\documentclass[11pt,a4paper,oneside,openany]{book}
\usepackage[spanish,es-nodecimaldot,es-noquoting]{babel}
\usepackage{fontspec}
\setmainfont{texgyrepagella-regular.otf}[Path=recursos/fuentes/,BoldFont=texgyrepagella-bold.otf,ItalicFont=texgyrepagella-italic.otf,BoldItalicFont=texgyrepagella-bolditalic.otf]
\setsansfont{texgyreheros-regular.otf}[Path=recursos/fuentes/,BoldFont=texgyreheros-bold.otf,ItalicFont=texgyreheros-italic.otf,BoldItalicFont=texgyreheros-bolditalic.otf]
\setmonofont{texgyrecursor-regular.otf}[Path=recursos/fuentes/,BoldFont=texgyrecursor-bold.otf,ItalicFont=texgyrecursor-italic.otf,BoldItalicFont=texgyrecursor-bolditalic.otf,Scale=MatchLowercase]
\usepackage[a4paper,margin=2.1cm,headheight=15pt,headsep=0.6cm,footskip=1cm]{geometry}
\usepackage{xcolor}
\definecolor{Azul}{HTML}{16324F}
\definecolor{Teal}{HTML}{007E80}
\definecolor{Gris}{HTML}{526171}
\usepackage{graphicx,booktabs,longtable,array,ragged2e,enumitem,pdflscape,needspace,titlesec,fancyhdr,microtype,seqsplit}
\usepackage[unicode,colorlinks=true,linkcolor=Azul,urlcolor=Teal,citecolor=Azul,bookmarksnumbered=true,pdfdisplaydoctitle=true]{hyperref}
\usepackage{bookmark}
\hypersetup{pdftitle={Manual de operación y aprendizaje de la plataforma BI con IA},pdfauthor={Proyecto académico BI - Universidad de Guayaquil},pdfsubject={Formación y operación supervisada de datamarts de ventas}}
\setcounter{tocdepth}{1}
\setcounter{secnumdepth}{1}
\renewcommand{\baselinestretch}{1.13}
\setlength{\parindent}{0pt}
\setlength{\parskip}{6pt plus 1pt minus 1pt}
\setlength{\emergencystretch}{3em}
\tolerance=1800
\widowpenalty=10000
\clubpenalty=10000
\displaywidowpenalty=10000
\raggedbottom
\setlist[itemize]{leftmargin=1.5em,itemsep=3pt,topsep=4pt,parsep=2pt}
\setlist[enumerate]{leftmargin=1.7em,label=\textcolor{Teal}{\bfseries\arabic*.},itemsep=5pt,topsep=5pt,parsep=3pt}
\titleformat{\chapter}[display]{\RaggedRight\sffamily\bfseries\color{Azul}\hyphenpenalty=10000\exhyphenpenalty=10000}{\fontsize{11}{13}\selectfont CAPÍTULO \thechapter}{5pt}{\fontsize{24}{28}\selectfont}
\titlespacing*{\chapter}{0pt}{-8pt}{18pt}
\titleformat{\section}{\RaggedRight\sffamily\bfseries\large\color{Azul}\hyphenpenalty=10000\exhyphenpenalty=10000}{\thesection}{0.7em}{}
\titleformat{\subsection}{\sffamily\bfseries\normalsize\color{Teal}}{}{0pt}{}
\titleformat{\subsubsection}{\sffamily\bfseries\normalsize\color{Azul}}{}{0pt}{}
\titleformat{\paragraph}[runin]{\bfseries}{}{0pt}{}
\titlespacing*{\section}{0pt}{16pt}{7pt}
\titlespacing*{\subsection}{0pt}{11pt}{5pt}
\titlespacing*{\subsubsection}{0pt}{9pt}{4pt}
\titlespacing*{\paragraph}{0pt}{6pt}{0.5em}
\usepackage{etoolbox}
\makeatletter
\renewcommand*{\l@section}{\@dottedtocline{1}{1.5em}{3.3em}}
\makeatother
\pretocmd{\section}{\Needspace{5\baselineskip}}{}{}
\pretocmd{\subsection}{\Needspace{4\baselineskip}}{}{}
\pretocmd{\subsubsection}{\Needspace{4\baselineskip}}{}{}
\pagestyle{fancy}
\fancyhf{}
\fancyhead[L]{\sffamily\footnotesize\color{Gris} PLATAFORMA BI + IA}
\fancyhead[R]{\sffamily\footnotesize\color{Gris} Operación y aprendizaje}
\fancyfoot[L]{\sffamily\footnotesize\color{Gris} Edición de formación 2.0 - 25 de noviembre de 2026}
\fancyfoot[R]{\sffamily\small\color{Azul}\thepage}
\renewcommand{\headrulewidth}{0.3pt}
\renewcommand{\footrulewidth}{0pt}
\fancypagestyle{plain}{\fancyhf{}\fancyfoot[L]{\sffamily\footnotesize\color{Gris} Manual de operación y aprendizaje}\fancyfoot[R]{\sffamily\small\thepage}\renewcommand{\headrulewidth}{0pt}}
\newcounter{pantallas}
\newcounter{diagramas}
\newcommand{\pantalla}[3]{%
  \clearpage\begin{landscape}\thispagestyle{empty}\begingroup
  \setlength{\parskip}{0pt}\phantomsection
  {\sffamily\bfseries\fontsize{14}{17}\selectfont\color{Azul}#2\par}
  \vspace{0.22cm}
  {\centering\includegraphics[width=25cm,height=13.8cm,keepaspectratio]{#1}\par}
  \vspace{0.18cm}\refstepcounter{pantallas}
  {\fontsize{10}{12}\selectfont\textbf{Pantalla \thepantallas.} #3\par}
  \vfill{\hfill\sffamily\footnotesize\thepage\par}
  \endgroup\end{landscape}%
}
\newcommand{\pantallaactual}[3]{%
  \clearpage\thispagestyle{plain}\begingroup
  {\sffamily\bfseries\large\color{Azul}#2\par}
  \vspace{0.3cm}
  {\centering\includegraphics[width=10.1cm,height=16.5cm,keepaspectratio]{#1}\par}
  \vspace{0.3cm}\refstepcounter{pantallas}
  {\small\textbf{Pantalla \thepantallas.} #3\par}
  \endgroup\clearpage
}
\newcommand{\diagramasistema}[3]{%
  \clearpage\begin{landscape}\thispagestyle{empty}\begingroup
  \setlength{\parskip}{0pt}\phantomsection
  {\sffamily\bfseries\fontsize{14}{17}\selectfont\color{Azul}#2\par}
  \vspace{0.22cm}
  {\centering\includegraphics[width=25cm,height=13.6cm,keepaspectratio]{#1}\par}
  \vspace{0.18cm}\refstepcounter{diagramas}
  {\fontsize{10}{12}\selectfont\textbf{Diagrama \thediagramas.} #3\par}
  \vfill{\hfill\sffamily\footnotesize\thepage\par}
  \endgroup\end{landscape}%
}
\begin{document}
\begin{titlepage}
\thispagestyle{empty}
{\sffamily\bfseries\color{Teal} UNIVERSIDAD DE GUAYAQUIL\par}
\vspace{0.35cm}
{\sffamily\color{Gris} Proyecto académico de inteligencia de negocios\par}
\vspace{2.0cm}
{\sffamily\bfseries\fontsize{30}{35}\selectfont\color{Azul} Manual de operación\par y aprendizaje\par}
\vspace{0.6cm}
{\sffamily\fontsize{20}{25}\selectfont\color{Teal} Plataforma BI con IA\par}
\vspace{0.8cm}
{\large Construcción supervisada de datamarts de ventas\par}
\vspace{0.5cm}
{\large Comprender, configurar, practicar y explicar\par}
\vspace{1cm}
{\color{Teal}\rule{\textwidth}{1.2pt}\par}
\vspace{0.6cm}
Para administración, analistas de BI, gerencia comercial\par
y programadoras del prototipo funcional.\par
\vfill
{\sffamily\color{Gris} Edición de formación 2.0\par 25 de noviembre de 2026\par}
\end{titlepage}
\frontmatter
\section*{Alcance de esta edición}
Esta edición de formación 2.0, del 25 de noviembre de 2026, integra las mejoras
operativas de navegación, sesiones, catálogo y cobertura de las propuestas,
así como la configuración de OpenAI y el criterio de costo de sus modelos.
Conserva los trece capítulos de la edición anterior, sus explicaciones,
sus 52 capturas de referencia del 18 de noviembre y sus ocho diagramas.
Añade tres vistas actuales, sin deformarlas ni presentarlas como capturas
de escritorio completo.

Las propuestas 129 y 131 y las ejecuciones 21 y 22 se mantienen como historia
de la campaña del 6 de noviembre. Los casos nuevos comprobados el 25 de
noviembre son AdventureWorks 136/24 y WideWorldImporters 135/23, con OpenAI
Sol y razonamiento Medio. La opción económica activa, OpenAI Luna con
razonamiento Bajo, sólo tiene una prueba mínima de conexión y salida
estructurada: su calidad BI completa sigue pendiente de una campaña propia.

Este manual permite estudiar y repetir el proceso, sin garantizar que
cualquier texto, modelo o motor produzca siempre una propuesta aprobable.
Las reglas verifican estructura y cálculos; la persona conserva la decisión
sobre significado comercial, alcance y uso del resultado. Los identificadores
son propios de cada instalación y cambian en un entorno limpio.

Las funciones de IA dependen de la disponibilidad del proveedor y pueden
consumir su cuota. El manual permite aprender con los expedientes conservados;
no exige generar otra propuesta ni repetir una carga para comprender el flujo.

\clearpage
\begingroup
\setlength{\parskip}{0pt}
\tableofcontents
\endgroup
\mainmatter
\chapter{Cómo aprender con este manual}
Este manual acompaña a las programadoras desde el primer ingreso hasta la consulta de un datamart de ventas. Cada recorrido explica qué se hace, por qué se hace, qué resultado debe aparecer y cómo reconocer un problema. Puede utilizarse para estudiar, practicar en el equipo de desarrollo y preparar una demostración ante el jurado.

La meta de aprendizaje es operar con criterio: identificar la fuente de cada cifra, explicar una medida, reconocer una relación que puede duplicar ventas y justificar una decisión. Memorizar una secuencia de botones no basta cuando cambian la base de datos, las preguntas comerciales o la respuesta del modelo de IA.

Primero lea los fundamentos y la arquitectura. Después complete la preparación con una cuenta administradora. Recorra el catálogo y el asistente como analista de BI. Finalmente revise los resultados como gerente comercial. En cada práctica registre la fuente, la versión de metadatos, la propuesta, el expediente, los filtros y el resultado observado.

Las pantallas proceden de la aplicación abierta en vista completa durante la preparación de esta edición. Los casos de referencia conservan pruebas anteriores verificadas en AdventureWorks2022 y WideWorldImporters. Consultar esos expedientes permite estudiar el resultado; no significa que se haya repetido su generación o su ETL al obtener una captura.

En una instalación limpia pueden no existir las propuestas 129 y 131 ni las ejecuciones 21 y 22. Esos números identifican los ejemplos de la instalación de referencia. Las prácticas explican también cómo iniciar un caso propio. Anote los identificadores que produzca su equipo y compare las condiciones y cifras, en lugar de intentar reproducir el mismo número de versión.

\Needspace{36\baselineskip}
\Needspace{8\baselineskip}
\begingroup
\setbox0=\vbox{\hsize=\linewidth
{\normalfont\normalsize\RaggedRight\sffamily\bfseries\large\color{Azul}\hyphenpenalty=10000\exhyphenpenalty=10000 Dónde encontrar cada opción\par}

Puede abrir la unidad desde la última columna. El menú visible depende de los permisos de su sesión; el enlace del manual no concede acceso a la aplicación.

\par
\begingroup
\fontsize{10}{12.5}\selectfont
\setlength{\tabcolsep}{5pt}
\renewcommand{\arraystretch}{1.15}
\begin{tabular}{@{}>{\RaggedRight\arraybackslash}p{0.25\linewidth}>{\RaggedRight\arraybackslash}p{0.44\linewidth}>{\RaggedRight\arraybackslash}p{0.22\linewidth}@{}}
\toprule
\textbf{Opción visible} & \textbf{Para qué acudir} & \textbf{Unidad del manual} \\
\midrule
Inicio & Reconocer la preparación y la fuente de trabajo. & \hyperref[unidad:operacion_02]{Abrir sección \ref*{unidad:operacion_02}} \\
\bottomrule
\end{tabular}
\endgroup
}
\Needspace{\dimexpr\ht0+\dp0+48pt\relax}
\endgroup
\section{Dónde encontrar cada opción}
Puede abrir la unidad desde la última columna. El menú visible depende de los permisos de su sesión; el enlace del manual no concede acceso a la aplicación.

\begingroup
\fontsize{10}{12.5}\selectfont
\setlength{\tabcolsep}{5pt}
\renewcommand{\arraystretch}{1.15}
\begin{longtable}{@{}>{\RaggedRight\arraybackslash}p{0.25\linewidth}>{\RaggedRight\arraybackslash}p{0.44\linewidth}>{\RaggedRight\arraybackslash}p{0.22\linewidth}@{}}
\toprule
\textbf{Opción visible} & \textbf{Para qué acudir} & \textbf{Unidad del manual} \\
\midrule
\endfirsthead
\toprule
\textbf{Opción visible} & \textbf{Para qué acudir} & \textbf{Unidad del manual} \\
\midrule
\endhead
\midrule
\multicolumn{3}{r}{\footnotesize Continúa en la página siguiente}\\
\endfoot
\bottomrule
\endlastfoot
Inicio & Reconocer la preparación y la fuente de trabajo. & \hyperref[unidad:operacion_02]{Abrir sección \ref*{unidad:operacion_02}} \\
\addlinespace[3pt]
Conexiones de datos & Registrar y comprobar el acceso a cada origen. & \hyperref[unidad:operacion_09]{Abrir sección \ref*{unidad:operacion_09}} \\
\addlinespace[3pt]
Parámetros & Comprender los valores generales de operación. & \hyperref[unidad:operacion_10]{Abrir sección \ref*{unidad:operacion_10}} \\
\addlinespace[3pt]
Configuración LLM & Preparar y comprobar el proveedor de IA. & \hyperref[unidad:operacion_11]{Abrir sección \ref*{unidad:operacion_11}} \\
\addlinespace[3pt]
Explorador de esquema & Revisar tablas, columnas y relaciones capturadas. & \hyperref[unidad:operacion_12]{Abrir sección \ref*{unidad:operacion_12}} \\
\addlinespace[3pt]
Catálogo analítico & Definir preguntas y opciones propias de cada fuente. & \hyperref[unidad:operacion_13]{Abrir sección \ref*{unidad:operacion_13}} \\
\addlinespace[3pt]
Asistente de datamart & Convertir la necesidad en una propuesta supervisada. & \hyperref[unidad:BI02]{Abrir sección \ref*{unidad:BI02}} \\
\addlinespace[3pt]
Datamart de ventas & Preparar, confirmar y conciliar una carga. & \hyperref[unidad:ETL01]{Abrir sección \ref*{unidad:ETL01}} \\
\addlinespace[3pt]
Analítica de ventas & Filtrar, interpretar, conversar y exportar. & \hyperref[unidad:AN01]{Abrir sección \ref*{unidad:AN01}} \\
\addlinespace[3pt]
Usuarios & Administrar cuentas y asignaciones de acceso. & \hyperref[unidad:operacion_04]{Abrir sección \ref*{unidad:operacion_04}} \\
\addlinespace[3pt]
Roles & Agrupar permisos de acuerdo con responsabilidades. & \hyperref[unidad:operacion_05]{Abrir sección \ref*{unidad:operacion_05}} \\
\addlinespace[3pt]
Permisos & Consultar qué autoriza cada capacidad. & \hyperref[unidad:operacion_06]{Abrir sección \ref*{unidad:operacion_06}} \\
\addlinespace[3pt]
Menús & Comprender la navegación y su visibilidad. & \hyperref[unidad:operacion_07]{Abrir sección \ref*{unidad:operacion_07}} \\
\addlinespace[3pt]
Auditoría & Consultar los eventos registrados por la aplicación. & \hyperref[unidad:operacion_08]{Abrir sección \ref*{unidad:operacion_08}} \\
\addlinespace[3pt]
\end{longtable}
\endgroup

\begingroup
\setbox0=\vbox{\hsize=\linewidth
{\normalfont\normalsize\RaggedRight\sffamily\bfseries\large\color{Azul}\hyphenpenalty=10000\exhyphenpenalty=10000 Responsabilidades durante la operación\par}

Una persona puede estudiar varios perfiles, pero la aplicación permite únicamente las acciones concedidas a su cuenta. Los papeles de Product Owner y Scrum Master pertenecen al trabajo del equipo del proyecto; no conceden por sí mismos permisos dentro de la plataforma.

No se aprende la seguridad dando permisos administrativos a todas las personas. Para una prueba de perfiles, la administradora prepara las cuentas que haya autorizado el responsable y cada estudiante ingresa con el perfil que le corresponda. Ver un rol en la tabla no demuestra que ya se haya probado una sesión con ese rol.

\par
\begingroup
\fontsize{10}{12.5}\selectfont
\setlength{\tabcolsep}{5pt}
\renewcommand{\arraystretch}{1.15}
\begin{tabular}{@{}>{\RaggedRight\arraybackslash}p{0.22\linewidth}>{\RaggedRight\arraybackslash}p{0.3\linewidth}>{\RaggedRight\arraybackslash}p{0.39\linewidth}@{}}
\toprule
\textbf{Perfil} & \textbf{Trabajo principal} & \textbf{Criterio que debe comprender} \\
\midrule
Administración & Preparar cuentas, permisos, conexiones, parámetros y proveedor de IA. & Un permiso habilita una acción; una conexión habilitada no combina sus datos con otra. \\
\bottomrule
\end{tabular}
\endgroup
}
\Needspace{\dimexpr\ht0+\dp0+48pt\relax}
\endgroup
\section{Responsabilidades durante la operación}
Una persona puede estudiar varios perfiles, pero la aplicación permite únicamente las acciones concedidas a su cuenta. Los papeles de Product Owner y Scrum Master pertenecen al trabajo del equipo del proyecto; no conceden por sí mismos permisos dentro de la plataforma.

No se aprende la seguridad dando permisos administrativos a todas las personas. Para una prueba de perfiles, la administradora prepara las cuentas que haya autorizado el responsable y cada estudiante ingresa con el perfil que le corresponda. Ver un rol en la tabla no demuestra que ya se haya probado una sesión con ese rol.

\begingroup
\fontsize{10}{12.5}\selectfont
\setlength{\tabcolsep}{5pt}
\renewcommand{\arraystretch}{1.15}
\begin{longtable}{@{}>{\RaggedRight\arraybackslash}p{0.22\linewidth}>{\RaggedRight\arraybackslash}p{0.3\linewidth}>{\RaggedRight\arraybackslash}p{0.39\linewidth}@{}}
\toprule
\textbf{Perfil} & \textbf{Trabajo principal} & \textbf{Criterio que debe comprender} \\
\midrule
\endfirsthead
\toprule
\textbf{Perfil} & \textbf{Trabajo principal} & \textbf{Criterio que debe comprender} \\
\midrule
\endhead
\midrule
\multicolumn{3}{r}{\footnotesize Continúa en la página siguiente}\\
\endfoot
\bottomrule
\endlastfoot
Administración & Preparar cuentas, permisos, conexiones, parámetros y proveedor de IA. & Un permiso habilita una acción; una conexión habilitada no combina sus datos con otra. \\
\addlinespace[3pt]
Analista de BI & Expresar la necesidad, comprobar su respaldo, revisar el modelo, aprobar y supervisar el ETL. & El significado comercial, el grano y las relaciones deben coincidir con lo solicitado. \\
\addlinespace[3pt]
Gerencia comercial & Consultar indicadores, comparar segmentos, formular preguntas y exportar resultados. & Una cifra necesita fuente, período, unidad y contexto antes de apoyar una decisión. \\
\addlinespace[3pt]
Programadora & Comprender el flujo, revisar la trazabilidad y reproducir una incidencia con evidencia. & El motor común valida y materializa contratos; no se programa una consulta distinta para cada demostración. \\
\addlinespace[3pt]
\end{longtable}
\endgroup

\begingroup
\setbox0=\vbox{\hsize=\linewidth
{\normalfont\normalsize\RaggedRight\sffamily\bfseries\large\color{Azul}\hyphenpenalty=10000\exhyphenpenalty=10000 Tres niveles de práctica\par}

La lectura de un expediente existente permite aprender aun cuando no haya una API key disponible. Generar sugerencias, revisar una necesidad con IA, obtener una propuesta o conversar con el copiloto requiere que el proveedor correspondiente esté configurado y disponible.

El ETL y las aprobaciones de una práctica crean estado persistente. Realice esos ejercicios en el entorno de aprendizaje y con nombres que permitan reconocer el caso. Para estudiar los botones de retirar, restaurar o descartar, utilice un caso de práctica identificado; conserve las aprobaciones y expedientes elegidos para la demostración.

\par
\begingroup
\fontsize{10}{12.5}\selectfont
\setlength{\tabcolsep}{5pt}
\renewcommand{\arraystretch}{1.15}
\begin{tabular}{@{}>{\RaggedRight\arraybackslash}p{0.22\linewidth}>{\RaggedRight\arraybackslash}p{0.3\linewidth}>{\RaggedRight\arraybackslash}p{0.39\linewidth}@{}}
\toprule
\textbf{Nivel} & \textbf{Qué puede practicar} & \textbf{Condición de cierre} \\
\midrule
Lectura y explicación & Cambiar la fuente, explorar metadatos, consultar una propuesta y un expediente conservados, aplicar filtros y revisar reportes. & Explica la procedencia y el significado de los valores sin modificar la propuesta ni repetir una carga. \\
\bottomrule
\end{tabular}
\endgroup
}
\Needspace{\dimexpr\ht0+\dp0+48pt\relax}
\endgroup
\section{Tres niveles de práctica}
La lectura de un expediente existente permite aprender aun cuando no haya una API key disponible. Generar sugerencias, revisar una necesidad con IA, obtener una propuesta o conversar con el copiloto requiere que el proveedor correspondiente esté configurado y disponible.

El ETL y las aprobaciones de una práctica crean estado persistente. Realice esos ejercicios en el entorno de aprendizaje y con nombres que permitan reconocer el caso. Para estudiar los botones de retirar, restaurar o descartar, utilice un caso de práctica identificado; conserve las aprobaciones y expedientes elegidos para la demostración.

\begingroup
\fontsize{10}{12.5}\selectfont
\setlength{\tabcolsep}{5pt}
\renewcommand{\arraystretch}{1.15}
\begin{longtable}{@{}>{\RaggedRight\arraybackslash}p{0.22\linewidth}>{\RaggedRight\arraybackslash}p{0.3\linewidth}>{\RaggedRight\arraybackslash}p{0.39\linewidth}@{}}
\toprule
\textbf{Nivel} & \textbf{Qué puede practicar} & \textbf{Condición de cierre} \\
\midrule
\endfirsthead
\toprule
\textbf{Nivel} & \textbf{Qué puede practicar} & \textbf{Condición de cierre} \\
\midrule
\endhead
\midrule
\multicolumn{3}{r}{\footnotesize Continúa en la página siguiente}\\
\endfoot
\bottomrule
\endlastfoot
Lectura y explicación & Cambiar la fuente, explorar metadatos, consultar una propuesta y un expediente conservados, aplicar filtros y revisar reportes. & Explica la procedencia y el significado de los valores sin modificar la propuesta ni repetir una carga. \\
\addlinespace[3pt]
Diseño supervisado & Redactar una necesidad nueva, solicitar ayuda de IA, revisar viabilidad, conceptos, medidas, dimensiones y fórmulas. & Conserva un contrato coherente con el objetivo y distingue advertencias aceptables de errores bloqueantes. \\
\addlinespace[3pt]
Construcción y análisis & Aprobar una propuesta de práctica, revisar indicadores, confirmar el plan, ejecutar ETL, conciliar y consultar la analítica. & Relaciona necesidad, contrato, expediente y resultados; registra las comprobaciones y los límites del caso. \\
\addlinespace[3pt]
\end{longtable}
\endgroup

\section{Cómo reconocer que una práctica terminó correctamente}
\begin{enumerate}
\item Confirme la fuente seleccionada antes de leer una tabla o un indicador. AdventureWorks2022 y WideWorldImporters tienen estructuras y documentos comerciales diferentes.
\item Lea el mensaje de resultado. Un botón pulsado, una pantalla abierta o una respuesta con apariencia correcta no demuestran que la validación haya terminado.
\item Anote el objetivo, las preguntas, la periodicidad, el proveedor y los identificadores producidos. Registre qué advertencias comprendió y qué decisiones tomó.
\item Después del ETL contraste las filas de origen y destino, las diferencias y los indicadores conciliados. La validación estructural de una propuesta no sustituye esta conciliación.
\item En la analítica compruebe el expediente y los filtros. Verifique un valor exacto desde el detalle o la tabla de respaldo y contraste la exportación con la selección visible.
\item Explique una limitación del caso. Por ejemplo, una moneda no comprobada, atributos actuales que no representan historia o un período extremo con cobertura parcial.
\end{enumerate}


\IfFileExists{capturas/fundamentos_01.tex}{\input{capturas/fundamentos_01.tex}}{}

\chapter{Comprender una venta antes de calcular indicadores}
Una operación comercial puede incluir un pedido, una factura, una entrega, un pago y una devolución. Son hechos relacionados, pero no equivalentes. La pregunta de negocio debe señalar qué evento se analiza y qué fecha corresponde a ese evento. Un sistema puede conservar el pedido sin disponer de una prueba de entrega o de cobro.

Una cabecera representa el documento comercial y suele reunir su fecha y cliente. Sus líneas representan productos o servicios y cantidades. Un documento con tres líneas sigue siendo un documento. Esta diferencia determina la granularidad del datamart y la fórmula del número de transacciones.

Los ejemplos numéricos de este capítulo son didácticos e inventados para practicar en papel. Los resultados comprobados de AdventureWorks y WideWorldImporters se presentan por separado en los casos del sistema.

\begingroup
\setbox0=\vbox{\hsize=\linewidth
{\normalfont\normalsize\RaggedRight\sffamily\bfseries\large\color{Azul}\hyphenpenalty=10000\exhyphenpenalty=10000 Vocabulario básico para seguir el recorrido\par}

\par
\begingroup
\fontsize{10}{12.5}\selectfont
\setlength{\tabcolsep}{5pt}
\renewcommand{\arraystretch}{1.15}
\begin{tabular}{@{}>{\RaggedRight\arraybackslash}p{0.22\linewidth}>{\RaggedRight\arraybackslash}p{0.3\linewidth}>{\RaggedRight\arraybackslash}p{0.39\linewidth}@{}}
\toprule
\textbf{Concepto} & \textbf{Explicación práctica} & \textbf{Ejemplo de decisión} \\
\midrule
Inteligencia de negocios o BI & Proceso de transformar datos en información comprensible para analizar una actividad. & Comparar el importe y las unidades por producto para decidir qué resultados revisar. \\
\bottomrule
\end{tabular}
\endgroup
}
\Needspace{\dimexpr\ht0+\dp0+48pt\relax}
\endgroup
\section{Vocabulario básico para seguir el recorrido}
\begingroup
\fontsize{10}{12.5}\selectfont
\setlength{\tabcolsep}{5pt}
\renewcommand{\arraystretch}{1.15}
\begin{longtable}{@{}>{\RaggedRight\arraybackslash}p{0.22\linewidth}>{\RaggedRight\arraybackslash}p{0.3\linewidth}>{\RaggedRight\arraybackslash}p{0.39\linewidth}@{}}
\toprule
\textbf{Concepto} & \textbf{Explicación práctica} & \textbf{Ejemplo de decisión} \\
\midrule
\endfirsthead
\toprule
\textbf{Concepto} & \textbf{Explicación práctica} & \textbf{Ejemplo de decisión} \\
\midrule
\endhead
\midrule
\multicolumn{3}{r}{\footnotesize Continúa en la página siguiente}\\
\endfoot
\bottomrule
\endlastfoot
Inteligencia de negocios o BI & Proceso de transformar datos en información comprensible para analizar una actividad. & Comparar el importe y las unidades por producto para decidir qué resultados revisar. \\
\addlinespace[3pt]
Base operacional u OLTP & Base que registra las operaciones del negocio y sus relaciones. & Identificar qué tabla conserva una línea y cuál conserva la fecha de su documento. \\
\addlinespace[3pt]
Datamart & Conjunto de datos preparado para analizar un dominio concreto, aquí ventas. & Conservar las líneas y organizar fecha, producto, cliente y territorio para consultar indicadores. \\
\addlinespace[3pt]
Hecho & Registro del evento que se medirá. & Una línea de pedido o una línea de factura, según la necesidad. \\
\addlinespace[3pt]
Granularidad o grano & Significado exacto de una fila del hecho. & Una fila por línea evita confundir el número de líneas con el número de pedidos. \\
\addlinespace[3pt]
Dimensión & Eje que describe o agrupa el hecho. & Fecha, producto, cliente o territorio comercial. \\
\addlinespace[3pt]
Atributo & Característica de una dimensión. & Nombre, categoría o tipo, cuando la fuente y sus relaciones lo respaldan. \\
\addlinespace[3pt]
Medida & Dato numérico asociado al hecho y una regla de cálculo o agregación. & Cantidad, importe registrado y número de documentos distintos. \\
\addlinespace[3pt]
Indicador o KPI & Resultado calculado para responder una pregunta comercial. & Importe promedio por pedido calculado con el importe total y los pedidos distintos. \\
\addlinespace[3pt]
ETL & Extracción, transformación y carga de los datos hacia el modelo analítico. & Leer el origen, resolver claves y cargar tablas antes de conciliar los resultados. \\
\addlinespace[3pt]
Conciliación & Comparación de resultados equivalentes entre el origen y el destino. & Comprobar filas, unidades e importes bajo las mismas reglas. \\
\addlinespace[3pt]
Trazabilidad & Enlaces que permiten conocer de dónde procede un resultado y qué decisiones lo produjeron. & Relacionar una tarjeta con fuente, instantánea, propuesta y ejecución. \\
\addlinespace[3pt]
\end{longtable}
\endgroup

\begingroup
\setbox0=\vbox{\hsize=\linewidth
{\normalfont\normalsize\RaggedRight\sffamily\bfseries\large\color{Azul}\hyphenpenalty=10000\exhyphenpenalty=10000 Un ejemplo pequeño de tres líneas y dos pedidos\par}

En este ejemplo el descuento es una tasa. El valor 0,10 representa 10 \%, no diez unidades monetarias. El importe bruto de una línea es precio por cantidad; su descuento monetario es precio por cantidad por tasa. Si el origen ya conserva el importe después del descuento, restarlo otra vez produciría un resultado incorrecto.

La primera línea tiene importe bruto de 20,00, descuento de 2,00 e importe neto de 18,00. Las otras dos líneas tienen importes de 20,00 y 30,00. El total neto es 68,00 y la cantidad total es seis unidades. Existen tres líneas, dos pedidos distintos y dos productos.

\par
\begingroup
\fontsize{10}{12.5}\selectfont
\setlength{\tabcolsep}{5pt}
\renewcommand{\arraystretch}{1.15}
\begin{tabular}{@{}>{\RaggedRight\arraybackslash}p{0.22\linewidth}>{\RaggedRight\arraybackslash}p{0.3\linewidth}>{\RaggedRight\arraybackslash}p{0.39\linewidth}@{}}
\toprule
\textbf{Pedido y línea} & \textbf{Descripción} & \textbf{Datos del ejemplo} \\
\midrule
A, línea 1 & \textbf{Producto:} P1\par\smallskip\textbf{Cantidad:} 2 & \textbf{Precio unitario:} 10,00\par\smallskip\textbf{Tasa de descuento:} 0,10\par\smallskip\textbf{Costo unitario:} 6,00 \\
\bottomrule
\end{tabular}
\endgroup
}
\Needspace{\dimexpr\ht0+\dp0+48pt\relax}
\endgroup
\section{Un ejemplo pequeño de tres líneas y dos pedidos}
En este ejemplo el descuento es una tasa. El valor 0,10 representa 10 \%, no diez unidades monetarias. El importe bruto de una línea es precio por cantidad; su descuento monetario es precio por cantidad por tasa. Si el origen ya conserva el importe después del descuento, restarlo otra vez produciría un resultado incorrecto.

La primera línea tiene importe bruto de 20,00, descuento de 2,00 e importe neto de 18,00. Las otras dos líneas tienen importes de 20,00 y 30,00. El total neto es 68,00 y la cantidad total es seis unidades. Existen tres líneas, dos pedidos distintos y dos productos.

\begingroup
\fontsize{10}{12.5}\selectfont
\setlength{\tabcolsep}{5pt}
\renewcommand{\arraystretch}{1.15}
\begin{longtable}{@{}>{\RaggedRight\arraybackslash}p{0.22\linewidth}>{\RaggedRight\arraybackslash}p{0.3\linewidth}>{\RaggedRight\arraybackslash}p{0.39\linewidth}@{}}
\toprule
\textbf{Pedido y línea} & \textbf{Descripción} & \textbf{Datos del ejemplo} \\
\midrule
\endfirsthead
\toprule
\textbf{Pedido y línea} & \textbf{Descripción} & \textbf{Datos del ejemplo} \\
\midrule
\endhead
\midrule
\multicolumn{3}{r}{\footnotesize Continúa en la página siguiente}\\
\endfoot
\bottomrule
\endlastfoot
A, línea 1 & \textbf{Producto:} P1\par\smallskip\textbf{Cantidad:} 2 & \textbf{Precio unitario:} 10,00\par\smallskip\textbf{Tasa de descuento:} 0,10\par\smallskip\textbf{Costo unitario:} 6,00 \\
\addlinespace[3pt]
A, línea 2 & \textbf{Producto:} P2\par\smallskip\textbf{Cantidad:} 1 & \textbf{Precio unitario:} 20,00\par\smallskip\textbf{Tasa de descuento:} 0,00\par\smallskip\textbf{Costo unitario:} 12,00 \\
\addlinespace[3pt]
B, línea 1 & \textbf{Producto:} P1\par\smallskip\textbf{Cantidad:} 3 & \textbf{Precio unitario:} 10,00\par\smallskip\textbf{Tasa de descuento:} 0,00\par\smallskip\textbf{Costo unitario:} 6,00 \\
\addlinespace[3pt]
\end{longtable}
\endgroup

\begingroup
\setbox0=\vbox{\hsize=\linewidth
{\normalfont\normalsize\RaggedRight\sffamily\bfseries\large\color{Azul}\hyphenpenalty=10000\exhyphenpenalty=10000 Elegir la agregación y el denominador\par}

Sumar sirve para magnitudes aditivas como cantidades o importes compatibles. Conteo cuenta registros con la regla definida; conteo distinto cuenta identificadores sin repetirlos dentro del grupo o filtro. Promedio divide una suma entre una cantidad de observaciones. Antes de elegirlo debe explicar qué representa una observación.

Un promedio por documento o por unidad es una razón de agregados. No se obtiene tomando el promedio simple de los promedios de cada producto o territorio. Al cambiar un filtro, el sistema debe volver a calcular numerador y denominador sobre esa selección.

El conteo distinto no siempre puede sumarse entre categorías. Un pedido puede contener más de un producto y aparecer en varias categorías. Tampoco debe promediar porcentajes de margen de grupos con tamaños diferentes para obtener el margen global.

Si el denominador es cero, el indicador carece de un cociente válido para esa selección. No lo reemplace por uno ni confunda la ausencia de resultado con una venta cero. Revise qué información existe y cómo comunica ese estado la pantalla.

\par
\begingroup
\fontsize{10}{12.5}\selectfont
\setlength{\tabcolsep}{5pt}
\renewcommand{\arraystretch}{1.15}
\begin{tabular}{@{}>{\RaggedRight\arraybackslash}p{0.22\linewidth}>{\RaggedRight\arraybackslash}p{0.3\linewidth}>{\RaggedRight\arraybackslash}p{0.39\linewidth}@{}}
\toprule
\textbf{Pregunta} & \textbf{Cálculo del ejemplo} & \textbf{Resultado e interpretación} \\
\midrule
¿Cuánto se vendió según el importe neto del ejemplo? & 18 + 20 + 30 & 68,00 unidades monetarias. \\
\bottomrule
\end{tabular}
\endgroup
}
\Needspace{\dimexpr\ht0+\dp0+48pt\relax}
\endgroup
\section{Elegir la agregación y el denominador}
Sumar sirve para magnitudes aditivas como cantidades o importes compatibles. Conteo cuenta registros con la regla definida; conteo distinto cuenta identificadores sin repetirlos dentro del grupo o filtro. Promedio divide una suma entre una cantidad de observaciones. Antes de elegirlo debe explicar qué representa una observación.

Un promedio por documento o por unidad es una razón de agregados. No se obtiene tomando el promedio simple de los promedios de cada producto o territorio. Al cambiar un filtro, el sistema debe volver a calcular numerador y denominador sobre esa selección.

El conteo distinto no siempre puede sumarse entre categorías. Un pedido puede contener más de un producto y aparecer en varias categorías. Tampoco debe promediar porcentajes de margen de grupos con tamaños diferentes para obtener el margen global.

Si el denominador es cero, el indicador carece de un cociente válido para esa selección. No lo reemplace por uno ni confunda la ausencia de resultado con una venta cero. Revise qué información existe y cómo comunica ese estado la pantalla.

\begingroup
\fontsize{10}{12.5}\selectfont
\setlength{\tabcolsep}{5pt}
\renewcommand{\arraystretch}{1.15}
\begin{longtable}{@{}>{\RaggedRight\arraybackslash}p{0.22\linewidth}>{\RaggedRight\arraybackslash}p{0.3\linewidth}>{\RaggedRight\arraybackslash}p{0.39\linewidth}@{}}
\toprule
\textbf{Pregunta} & \textbf{Cálculo del ejemplo} & \textbf{Resultado e interpretación} \\
\midrule
\endfirsthead
\toprule
\textbf{Pregunta} & \textbf{Cálculo del ejemplo} & \textbf{Resultado e interpretación} \\
\midrule
\endhead
\midrule
\multicolumn{3}{r}{\footnotesize Continúa en la página siguiente}\\
\endfoot
\bottomrule
\endlastfoot
¿Cuánto se vendió según el importe neto del ejemplo? & 18 + 20 + 30 & 68,00 unidades monetarias. \\
\addlinespace[3pt]
¿Cuántas unidades se vendieron? & 2 + 1 + 3 & 6 unidades físicas. \\
\addlinespace[3pt]
¿Cuántos pedidos hubo? & Contar los identificadores distintos A y B & 2 pedidos; contar las filas daría 3 líneas. \\
\addlinespace[3pt]
¿Cuál es el importe promedio por pedido? & 68 / 2 & 34,00 por pedido. \\
\addlinespace[3pt]
¿Cuál es el importe promedio por línea? & 68 / 3 & 22,67 por línea; responde una pregunta diferente. \\
\addlinespace[3pt]
¿Cuál es la venta promedio por unidad física? & 68 / 6 & 11,33 por unidad; utiliza cantidades como ponderación. \\
\addlinespace[3pt]
¿Cuántos pedidos contienen cada producto? & P1 aparece en A y B; P2 aparece en A & 2 para P1 y 1 para P2. Su suma 3 no sustituye los 2 pedidos totales. \\
\addlinespace[3pt]
\end{longtable}
\endgroup

\section{Descuento, costo y margen necesitan significado comprobado}
Descuento porcentual y descuento monetario son medidas diferentes. Precio de venta y costo unitario también lo son. Una columna numérica no puede intercambiarse por otra únicamente porque permita una multiplicación.

En el ejemplo didáctico, el costo total sería 2 por 6,00 más 1 por 12,00 más 3 por 6,00: 42,00. Si esos costos fueran pertinentes para las ventas analizadas, el margen bruto sería 68,00 menos 42,00, es decir 26,00. El porcentaje de margen sería 26,00 dividido entre 68,00 por 100: aproximadamente 38,24 \%.

Este cálculo no autoriza a llamar utilidad al margen bruto: todavía pueden faltar gastos, impuestos, devoluciones u otros componentes. Tampoco convierte un costo estándar actual en el costo histórico de una venta antigua. Primero hay que comprobar la procedencia, la fecha y el significado del costo.

En una fuente distinta, un importe registrado puede incluir componentes diferentes. Mantenga un nombre neutral como importe de líneas mientras no se compruebe si corresponde a ventas netas, bruto, impuestos o descuento. La IA puede proponer una interpretación; el analista debe exigir respaldo.

\section{Relacionar tablas sin multiplicar las ventas}
Una clave primaria identifica una fila de una tabla. Una clave foránea declara que un valor apunta a otra tabla. Una clave compuesta necesita todas sus columnas. Si se omite parte de ella al relacionar tablas, pueden aparecer coincidencias múltiples.

Una línea normalmente puede apuntar a una cabecera y a un producto. Al unirla con esas entidades, se espera conservar la cantidad de líneas. Si una línea encuentra dos filas en una dimensión y se duplica, también pueden duplicarse su importe y su cantidad. El aspecto correcto de un gráfico no revela ese error por sí solo.

No elija una relación por la semejanza de dos nombres. Cliente comprador, cliente de facturación, ciudad de entrega y ciudad postal pueden representar papeles diferentes. Compare el objetivo comercial, la ruta declarada y el efecto sobre el grano.

La fecha del hecho indica cuándo ocurrió el evento elegido. Un atributo dimensional actual no demuestra cómo era el cliente o el territorio en esa fecha. Un domicilio de entrega registrado tampoco demuestra que la mercancía haya sido entregada.

\section{Periodicidad, filtros, moneda y ausencia de datos}
La periodicidad define cómo se agrupan los resultados en el tiempo: día, semana, mes, trimestre o año. Conservar una fila por línea y mostrar ventas mensuales es compatible: el detalle físico permanece y la consulta agrupa sus valores. Elegir mensual no autoriza a eliminar la trazabilidad de las líneas.

Un filtro cambia la selección sobre la que se consulta un indicador. Un ranking ordena el resultado de esa selección y puede mostrar sólo las primeras categorías. Un porcentaje calculado sobre las categorías visibles no debe presentarse como participación en toda la empresa si el universo es más amplio.

La moneda debe estar comprobada. El punto o la coma decimal, un nombre en inglés o la nacionalidad de la base de demostración no prueban que el importe esté en USD. Cuando la aplicación indica moneda de origen, conserve esa identificación y explique la limitación.

Un nulo significa un valor ausente o desconocido según el contexto; cero es un valor conocido. Una vista sin resultados puede deberse a filtros sin coincidencias, ausencia de una ejecución conciliada o falta de cobertura. Revise el mensaje antes de concluir que no hubo ventas.


\IfFileExists{capturas/fundamentos_02.tex}{\input{capturas/fundamentos_02.tex}}{}

\chapter{Entender la arquitectura y el papel de la IA}
El navegador presenta formularios, recorridos y resultados. El backend recibe las solicitudes, comprueba los permisos y aplica las reglas. SQL Server conserva las fuentes operacionales de demostración; PostgreSQL conserva la información interna del prototipo y los datos analíticos materializados. El proveedor de IA recibe únicamente el contexto permitido para cada función.

React organiza la interfaz. En desarrollo se sirve con Vite; la entrega compilada se sirve con Nginx. FastAPI concentra los servicios y la autorización. Docker ejecuta los componentes del entorno. Estos nombres describen piezas técnicas; durante la operación ordinaria la analista trabaja desde las pantallas.

La arquitectura emplea un proceso común y configuraciones por conexión. AdventureWorks y WideWorldImporters demuestran su aplicación sobre dos estructuras SQL Server distintas. Esto no significa que el corte actual ya disponga de conectores funcionales para todos los motores del mercado.

El sistema no ejecuta SQL libre redactado por el modelo. La propuesta debe ajustarse a contratos y referencias verificables. Una persona autorizada decide su aprobación, y el motor determinístico prepara la materialización. Después de la carga se comparan los resultados con el origen.

OpenAI se integra como un adaptador adicional: no cambia el recorrido de catálogo, viabilidad, contrato, aprobación ni ETL. El modelo propone o explica; el servidor sigue comprobando referencias, recetas y permisos. La fuente activa y el proveedor activo son decisiones distintas: cambiar el proveedor no cambia la base del análisis.

\begingroup
\setbox0=\vbox{\hsize=\linewidth
{\normalfont\normalsize\RaggedRight\sffamily\bfseries\large\color{Azul}\hyphenpenalty=10000\exhyphenpenalty=10000 Qué ocurre en cada etapa\par}

Un LLM es un modelo de lenguaje. Puede producir una respuesta convincente que no conserve el objetivo o que confunda conceptos. El sistema inteligente combina esa capacidad con reglas, datos estructurales, validaciones y supervisión. Esas piezas deben explicarse juntas al presentar el proyecto.

Una respuesta en JSON puede ser sintácticamente válida y aun así incumplir el contrato del prototipo. Un contrato aceptado tampoco demuestra por sí solo que todos los registros sean completos o que un costo tenga significado histórico. Las verificaciones se acumulan durante el proceso.

Cambiar de proveedor modifica el destino de las consultas de IA; no cambia automáticamente las reglas del motor. El tiempo de respuesta, la cuota, la disponibilidad y el cumplimiento del contrato pueden variar. Evalúe el resultado que recibió, no el prestigio del nombre del modelo.

\par
\begingroup
\fontsize{10}{12.5}\selectfont
\setlength{\tabcolsep}{5pt}
\renewcommand{\arraystretch}{1.15}
\begin{tabular}{@{}>{\RaggedRight\arraybackslash}p{0.22\linewidth}>{\RaggedRight\arraybackslash}p{0.3\linewidth}>{\RaggedRight\arraybackslash}p{0.39\linewidth}@{}}
\toprule
\textbf{Etapa} & \textbf{Participación de la IA} & \textbf{Control de la plataforma y decisión humana} \\
\midrule
Conexión y metadatos & No sustituye la lectura del esquema. & El conector consulta tablas, columnas, tipos y claves con permisos de lectura; se conserva una instantánea. \\
\bottomrule
\end{tabular}
\endgroup
}
\Needspace{\dimexpr\ht0+\dp0+48pt\relax}
\endgroup
\section{Qué ocurre en cada etapa}
Un LLM es un modelo de lenguaje. Puede producir una respuesta convincente que no conserve el objetivo o que confunda conceptos. El sistema inteligente combina esa capacidad con reglas, datos estructurales, validaciones y supervisión. Esas piezas deben explicarse juntas al presentar el proyecto.

Una respuesta en JSON puede ser sintácticamente válida y aun así incumplir el contrato del prototipo. Un contrato aceptado tampoco demuestra por sí solo que todos los registros sean completos o que un costo tenga significado histórico. Las verificaciones se acumulan durante el proceso.

Cambiar de proveedor modifica el destino de las consultas de IA; no cambia automáticamente las reglas del motor. El tiempo de respuesta, la cuota, la disponibilidad y el cumplimiento del contrato pueden variar. Evalúe el resultado que recibió, no el prestigio del nombre del modelo.

\begingroup
\fontsize{10}{12.5}\selectfont
\setlength{\tabcolsep}{5pt}
\renewcommand{\arraystretch}{1.15}
\begin{longtable}{@{}>{\RaggedRight\arraybackslash}p{0.22\linewidth}>{\RaggedRight\arraybackslash}p{0.3\linewidth}>{\RaggedRight\arraybackslash}p{0.39\linewidth}@{}}
\toprule
\textbf{Etapa} & \textbf{Participación de la IA} & \textbf{Control de la plataforma y decisión humana} \\
\midrule
\endfirsthead
\toprule
\textbf{Etapa} & \textbf{Participación de la IA} & \textbf{Control de la plataforma y decisión humana} \\
\midrule
\endhead
\midrule
\multicolumn{3}{r}{\footnotesize Continúa en la página siguiente}\\
\endfoot
\bottomrule
\endlastfoot
Conexión y metadatos & No sustituye la lectura del esquema. & El conector consulta tablas, columnas, tipos y claves con permisos de lectura; se conserva una instantánea. \\
\addlinespace[3pt]
Catálogo & Utiliza las preguntas elegidas como orientación al generar una propuesta. & El analista mantiene preguntas y periodicidades por fuente. Las preguntas iniciales del producto no son una generación nueva de IA. \\
\addlinespace[3pt]
Necesidad y viabilidad & Puede sugerir necesidades, mejorar su redacción e interpretar los requisitos frente a la estructura. & Se comprueba el contexto y el respaldo de referencias y derivaciones. La persona adopta el texto y revisa los límites. \\
\addlinespace[3pt]
Conceptos y modelo & Propone conceptos, grano, dimensiones, medidas, indicadores y un plan declarativo. & El contrato valida estructura y semántica admitida. La analista corrige, justifica y aprueba o rechaza. \\
\addlinespace[3pt]
ETL y conciliación & La interpretación semántica puede sugerir etiquetas para categorías aptas. & El motor extrae, transforma y carga con reglas fijadas; compara resultados y exige revisar las etiquetas publicadas. \\
\addlinespace[3pt]
Analítica & El chat interpreta preguntas y explica evidencia de agregados permitidos. & Los indicadores y los hallazgos del panel se calculan mediante reglas. Los filtros y las exportaciones conservan el contexto. \\
\addlinespace[3pt]
Decisión comercial & Puede ayudar a resumir y formular nuevas preguntas. & La gerencia contrasta los resultados y sus límites; una explicación no prueba causalidad ni efectos futuros. \\
\addlinespace[3pt]
\end{longtable}
\endgroup


\IfFileExists{capturas/fundamentos_03.tex}{\input{capturas/fundamentos_03.tex}}{}
\diagramasistema{recursos/diagramas/Arquitectura_plataforma.pdf}{Arquitectura de la plataforma}{Interfaz, servicios, fuentes de lectura, persistencia y proveedor de IA. La aplicación comprueba las reglas y una persona autoriza el diseño y confirma la carga por separado.}

\chapter{Ingresar, orientarse y elegir la fuente}
\section{Ingresar y reconocer una sesión autorizada}
\label{unidad:operacion_01}
\textbf{Para quién:} Todos los perfiles\par
\textbf{Opción:} Pantalla de acceso\par

Abra la dirección de la plataforma proporcionada por la administración. En el repositorio de desarrollo, el puerto predeterminado es 5173; una entrega release o portable puede publicar otro puerto. Use la dirección de la guía de la instalación que está encendiendo. El puerto identifica el servicio: no determina sus permisos ni la base de datos que analizará.

Una cuenta representa a la persona que realiza las acciones. Sus roles reúnen permisos; por eso dos personas pueden iniciar sesión y ver opciones diferentes. La contraseña permite autenticar la cuenta, mientras que los permisos determinan si puede consultar, modificar o ejecutar una operación. Conocer la contraseña de una cuenta no convierte automáticamente a esa cuenta en administradora.

Al ingresar, compruebe el nombre mostrado en Sesión activa y los menús disponibles. Trabaje con su cuenta asignada. La cuenta administrativa inicial está protegida porque permite recuperar la administración del entorno. Su contraseña se entrega por un canal de administración y no forma parte de este manual.

La sesión tiene una duración absoluta: se cuenta desde el ingreso y no vuelve a comenzar con cada clic. El valor inicial es 60 minutos. La administración puede cambiar Duración de la sesión (minutos), en Parámetros, entre 15 y 240. El cambio se aplica a los próximos inicios de sesión; no alarga una sesión que ya está abierta.

\subsection{Qué hacer si la sesión vence durante una práctica}
Vencimiento y Cerrar sesión son situaciones diferentes. El vencimiento solicita autenticarse otra vez y permite recuperar información mínima de la tarea. Cerrar sesión es una salida explícita: limpia el borrador local y los contextos de la pestaña. Ninguna de las dos acciones elimina los expedientes guardados en el servidor.

El tiempo máximo del asistente y la duración de la sesión tampoco son lo mismo. El primero limita la espera de una solicitud; el segundo limita el acceso autenticado. Aumentar uno no cambia el otro ni corrige un contrato inválido.

\begin{enumerate}
\item Lea el aviso de sesión vencida y vuelva a ingresar con su cuenta.
\item Compruebe la fuente y la versión recuperadas. Revise el objetivo antes de continuar.
\item Si había una operación de generación o carga en curso, consulte primero el listado y el estado. Una sesión vencida no demuestra que el servidor haya cancelado esa operación.
\item Si necesita más tiempo para una práctica siguiente, solicite ajustar el parámetro y vuelva a iniciar sesión. No comparta la cuenta de administración para evitar el vencimiento.
\end{enumerate}

\subsection{Conceptos para comprender la opción}
\paragraph{Autenticación} Comprobación de la identidad de una cuenta mediante sus credenciales.

\paragraph{Autorización} Comprobación de que la cuenta tiene permiso para la acción solicitada.

\paragraph{Sesión} Acceso temporal de una cuenta autenticada; puede vencer sin borrar los expedientes almacenados.

\subsection{Paso a paso}
\begin{enumerate}
\item Escriba el correo de su cuenta en Correo y la clave suministrada por la administración en Contraseña.
\item Pulse Iniciar sesión una sola vez y espere a que aparezca el espacio de trabajo.
\item Compruebe su nombre en Sesión activa. Si no coincide con la cuenta que esperaba usar, pulse Cerrar sesión e ingrese con la cuenta correcta.
\item Despliegue el grupo de menú correspondiente a su tarea. Si una opción no está disponible, solicite revisar los permisos en lugar de utilizar la cuenta de otra persona.
\item Si vence la sesión, ingrese otra vez y confirme el módulo, la fuente y el borrador recuperados. No pulse otra vez Cargar para resolver el vencimiento: primero consulte si la ejecución anterior terminó.
\item Cuando termine, pulse Cerrar sesión. Esta acción cierra su acceso y elimina el borrador local del asistente; los expedientes ya guardados permanecen en el servidor.
\end{enumerate}

\Needspace{8\baselineskip}
\begingroup
\setbox0=\vbox{\hsize=\linewidth
{\normalfont\normalsize\sffamily\bfseries\normalsize\color{Teal}\hyphenpenalty=10000\exhyphenpenalty=10000 Campos y botones\par}

\par
\begingroup
\fontsize{10}{12.5}\selectfont
\setlength{\tabcolsep}{5pt}
\renewcommand{\arraystretch}{1.15}
\begin{tabular}{@{}>{\RaggedRight\arraybackslash}p{0.22\linewidth}>{\RaggedRight\arraybackslash}p{0.3\linewidth}>{\RaggedRight\arraybackslash}p{0.39\linewidth}@{}}
\toprule
\textbf{Control} & \textbf{Para qué sirve} & \textbf{Ejemplo y resultado} \\
\midrule
Correo & Identifica la cuenta registrada. & \textbf{Ejemplo:} El correo asignado a la estudiante; no un nombre de rol.\par\smallskip\textbf{Resultado:} Se consulta esa cuenta al autenticar. \\
\bottomrule
\end{tabular}
\endgroup
}
\Needspace{\dimexpr\ht0+\dp0+48pt\relax}
\endgroup
\subsection{Campos y botones}
\begingroup
\fontsize{10}{12.5}\selectfont
\setlength{\tabcolsep}{5pt}
\renewcommand{\arraystretch}{1.15}
\begin{longtable}{@{}>{\RaggedRight\arraybackslash}p{0.22\linewidth}>{\RaggedRight\arraybackslash}p{0.3\linewidth}>{\RaggedRight\arraybackslash}p{0.39\linewidth}@{}}
\toprule
\textbf{Control} & \textbf{Para qué sirve} & \textbf{Ejemplo y resultado} \\
\midrule
\endfirsthead
\toprule
\textbf{Control} & \textbf{Para qué sirve} & \textbf{Ejemplo y resultado} \\
\midrule
\endhead
\midrule
\multicolumn{3}{r}{\footnotesize Continúa en la página siguiente}\\
\endfoot
\bottomrule
\endlastfoot
Correo & Identifica la cuenta registrada. & \textbf{Ejemplo:} El correo asignado a la estudiante; no un nombre de rol.\par\smallskip\textbf{Resultado:} Se consulta esa cuenta al autenticar. \\
\addlinespace[3pt]
Contraseña & Acredita el acceso; se presenta enmascarada. & \textbf{Ejemplo:} La clave entregada para esa cuenta.\par\smallskip\textbf{Resultado:} Se valida sin devolver la contraseña guardada. \\
\addlinespace[3pt]
Iniciar sesión & Solicita la autenticación. & \textbf{Ejemplo:} Acceder como Analista de BI con su propia cuenta.\par\smallskip\textbf{Resultado:} Muestra el espacio de trabajo permitido o una explicación del rechazo. \\
\addlinespace[3pt]
Cerrar sesión & Finaliza el acceso del navegador. & \textbf{Ejemplo:} Concluir una práctica en un equipo compartido.\par\smallskip\textbf{Resultado:} Regresa al formulario de acceso; no elimina propuestas ni datamarts. \\
\addlinespace[3pt]
\end{longtable}
\endgroup

\subsection{Reconocer la cuenta antes de operar}
Práctica propuesta para el entorno de capacitación; no se presenta como una nueva ejecución realizada durante la redacción.

\begin{enumerate}
\item Ingrese con una cuenta de práctica.
\item Anote el nombre mostrado y tres opciones disponibles.
\item Cierre sesión y compruebe que vuelve el formulario.
\end{enumerate}

\textbf{Qué debe obtener.} La cuenta mostrada coincide con la utilizada. No se crean propuestas ni cargas por iniciar o cerrar sesión.

\textbf{Cómo comprobarlo.} La estudiante explica la diferencia entre acceso a la plataforma y permiso para ejecutar ETL.

\Needspace{8\baselineskip}
\begingroup
\setbox0=\vbox{\hsize=\linewidth
{\normalfont\normalsize\sffamily\bfseries\normalsize\color{Teal}\hyphenpenalty=10000\exhyphenpenalty=10000 Si aparece un problema\par}

\par
\begingroup
\fontsize{10}{12.5}\selectfont
\setlength{\tabcolsep}{5pt}
\renewcommand{\arraystretch}{1.15}
\begin{tabular}{@{}>{\RaggedRight\arraybackslash}p{0.22\linewidth}>{\RaggedRight\arraybackslash}p{0.3\linewidth}>{\RaggedRight\arraybackslash}p{0.39\linewidth}@{}}
\toprule
\textbf{Qué observa} & \textbf{Causa posible} & \textbf{Qué hacer y comprobar} \\
\midrule
Credenciales inválidas o acceso rechazado. & Correo o clave incorrectos, cuenta desactivada o servicio no disponible. & Verifique la dirección de la plataforma y el correo; solicite a la administración revisar el estado de la cuenta. No cambie la configuración de las bases para resolver un problema de acceso.\par\smallskip\textbf{Comprobación:} La cuenta correcta aparece en Sesión activa. \\
\bottomrule
\end{tabular}
\endgroup
}
\Needspace{\dimexpr\ht0+\dp0+48pt\relax}
\endgroup
\subsection{Si aparece un problema}
\begingroup
\fontsize{10}{12.5}\selectfont
\setlength{\tabcolsep}{5pt}
\renewcommand{\arraystretch}{1.15}
\begin{longtable}{@{}>{\RaggedRight\arraybackslash}p{0.22\linewidth}>{\RaggedRight\arraybackslash}p{0.3\linewidth}>{\RaggedRight\arraybackslash}p{0.39\linewidth}@{}}
\toprule
\textbf{Qué observa} & \textbf{Causa posible} & \textbf{Qué hacer y comprobar} \\
\midrule
\endfirsthead
\toprule
\textbf{Qué observa} & \textbf{Causa posible} & \textbf{Qué hacer y comprobar} \\
\midrule
\endhead
\midrule
\multicolumn{3}{r}{\footnotesize Continúa en la página siguiente}\\
\endfoot
\bottomrule
\endlastfoot
Credenciales inválidas o acceso rechazado. & Correo o clave incorrectos, cuenta desactivada o servicio no disponible. & Verifique la dirección de la plataforma y el correo; solicite a la administración revisar el estado de la cuenta. No cambie la configuración de las bases para resolver un problema de acceso.\par\smallskip\textbf{Comprobación:} La cuenta correcta aparece en Sesión activa. \\
\addlinespace[3pt]
Token inválido o vencido. & La sesión temporal terminó. & Autentíquese nuevamente y revise fuente, borrador y versión recuperados antes de continuar.\par\smallskip\textbf{Comprobación:} La consulta del expediente funciona y no se ha repetido una ejecución. \\
\addlinespace[3pt]
La clave aparece dentro de otro campo por autocompletado. & El navegador insertó datos guardados en un formulario. & Limpie el dato que no corresponde y revise cada etiqueta antes de guardar. No capture ni comparta contraseñas.\par\smallskip\textbf{Comprobación:} Correo y nombre contienen los datos previstos; los campos de claves siguen enmascarados. \\
\addlinespace[3pt]
\end{longtable}
\endgroup

\subsection{Para explicar durante la sustentación}
\paragraph{¿La IA decide quién puede entrar?} No. El acceso y los permisos los comprueba el servidor mediante las cuentas y los roles configurados.

\paragraph{¿Cerrar sesión elimina un datamart?} No. Finaliza el acceso local; los resultados persistidos se conservan. Un borrador local no guardado sí puede perderse.

\pantalla{recursos/pantallas/00_acceso.jpg}{Ingreso con una cuenta autorizada}{Formulario vacío. Utilice las credenciales que entregue el responsable del entorno; no se reproducen contraseñas en el manual. Captura de referencia del 18 de noviembre de 2026; revise las instrucciones actualizadas de esta unidad.}

\IfFileExists{capturas/operacion_01.tex}{\input{capturas/operacion_01.tex}}{}
\section{Comenzar desde Inicio y completar la preparación}
\label{unidad:operacion_02}
\textbf{Para quién:} Administración, Analista de BI, Gerencia para reconocer su contexto\par
\textbf{Opción:} Inicio\par

Inicio es el punto de orientación. Para una cuenta con consulta de propuestas, resume tres componentes: Fuente de ventas, Metadatos y Asistente de IA. Los tres deben estar preparados para crear una propuesta nueva. Una fuente habilitada no sustituye a sus metadatos; un modelo activo no garantiza que la pregunta comercial sea factible.

Todo está listo significa que se cumplen los prerrequisitos del asistente en ese contexto. Todavía faltan la necesidad concreta, la evaluación de su viabilidad, el diseño, la revisión humana y la ejecución del ETL. El mensaje no acredita anticipadamente que cualquier necesidad pueda resolverse.

Si el perfil no tiene consulta de propuestas, Inicio puede presentar cantidades de roles, permisos y opciones visibles en lugar de las tarjetas del asistente. Es una diferencia de acceso prevista. Para gerencia, el recorrido principal se dirige a Analítica de ventas, donde se consulta un datamart existente.

\subsection{Conceptos para comprender la opción}
\paragraph{Prerrequisito} Condición que debe cumplirse antes de comenzar una etapa. Una instantánea de metadatos antes de interpretar el esquema.

\paragraph{Metadatos} Descripción de tablas, columnas, tipos y relaciones; no las filas comerciales.

\paragraph{Proveedor de IA} Servicio configurado para interpretar y explicar dentro de las reglas de la plataforma.

\subsection{Paso a paso}
\begin{enumerate}
\item Abra Inicio y compruebe Fuente en contexto.
\item Lea el estado de Fuente de ventas, Metadatos y Asistente de IA si su perfil muestra esas tarjetas.
\item Si aparece Faltan prerrequisitos, pulse Completar configuración inicial.
\item Recorra la preparación con Continuar y Atrás; use Completar este paso para abrir el módulo pendiente.
\item Si aparece Todo está listo, pulse Crear propuesta de datamart cuando vaya a iniciar un diseño. Para revisar una configuración ya válida, utilice Revisar configuración.
\item Para terminar la ayuda de preparación, utilice Cerrar preparación.
\end{enumerate}

\Needspace{8\baselineskip}
\begingroup
\setbox0=\vbox{\hsize=\linewidth
{\normalfont\normalsize\sffamily\bfseries\normalsize\color{Teal}\hyphenpenalty=10000\exhyphenpenalty=10000 Campos y botones\par}

\par
\begingroup
\fontsize{10}{12.5}\selectfont
\setlength{\tabcolsep}{5pt}
\renewcommand{\arraystretch}{1.15}
\begin{tabular}{@{}>{\RaggedRight\arraybackslash}p{0.22\linewidth}>{\RaggedRight\arraybackslash}p{0.3\linewidth}>{\RaggedRight\arraybackslash}p{0.39\linewidth}@{}}
\toprule
\textbf{Control} & \textbf{Para qué sirve} & \textbf{Ejemplo y resultado} \\
\midrule
Crear propuesta de datamart & Abre el asistente. & \textbf{Ejemplo:} Iniciar un diseño sobre WideWorldImporters.\par\smallskip\textbf{Resultado:} Abre la etapa Necesidad; no aprueba ni carga datos. \\
\bottomrule
\end{tabular}
\endgroup
}
\Needspace{\dimexpr\ht0+\dp0+48pt\relax}
\endgroup
\subsection{Campos y botones}
\begingroup
\fontsize{10}{12.5}\selectfont
\setlength{\tabcolsep}{5pt}
\renewcommand{\arraystretch}{1.15}
\begin{longtable}{@{}>{\RaggedRight\arraybackslash}p{0.22\linewidth}>{\RaggedRight\arraybackslash}p{0.3\linewidth}>{\RaggedRight\arraybackslash}p{0.39\linewidth}@{}}
\toprule
\textbf{Control} & \textbf{Para qué sirve} & \textbf{Ejemplo y resultado} \\
\midrule
\endfirsthead
\toprule
\textbf{Control} & \textbf{Para qué sirve} & \textbf{Ejemplo y resultado} \\
\midrule
\endhead
\midrule
\multicolumn{3}{r}{\footnotesize Continúa en la página siguiente}\\
\endfoot
\bottomrule
\endlastfoot
Crear propuesta de datamart & Abre el asistente. & \textbf{Ejemplo:} Iniciar un diseño sobre WideWorldImporters.\par\smallskip\textbf{Resultado:} Abre la etapa Necesidad; no aprueba ni carga datos. \\
\addlinespace[3pt]
Completar configuración inicial & Muestra la preparación guiada. & \textbf{Ejemplo:} Faltan los metadatos de la conexión elegida.\par\smallskip\textbf{Resultado:} Permite localizar el módulo pendiente. \\
\addlinespace[3pt]
Revisar configuración & Abre el repaso de la preparación. & \textbf{Ejemplo:} Confirmar qué proveedor está activo.\par\smallskip\textbf{Resultado:} Muestra cinco pasos de preparación sin modificar registros por sí mismo. \\
\addlinespace[3pt]
Continuar / Atrás & Recorre la ayuda inicial. & \textbf{Ejemplo:} Pasar de Proveedor de IA a Fuente de ventas.\par\smallskip\textbf{Resultado:} Cambia el paso informativo, no ejecuta una carga. \\
\addlinespace[3pt]
\end{longtable}
\endgroup

\subsection{Distinguir preparación de resultados}
Práctica propuesta para el entorno de capacitación; no se presenta como una nueva ejecución realizada durante la redacción.

\begin{enumerate}
\item Seleccione AdventureWorks y lea sus tarjetas.
\item Cambie a WideWorldImporters y repita la lectura.
\item Explique qué falta para disponer de un análisis nuevo.
\end{enumerate}

\textbf{Qué debe obtener.} El nombre y los metadatos se corresponden con la fuente seleccionada. La estudiante identifica que el panel analítico requiere una ejecución conciliada.

\textbf{Cómo comprobarlo.} No confunde Todo está listo con ETL ejecutado o propuesta aprobada.

\Needspace{8\baselineskip}
\begingroup
\setbox0=\vbox{\hsize=\linewidth
{\normalfont\normalsize\sffamily\bfseries\normalsize\color{Teal}\hyphenpenalty=10000\exhyphenpenalty=10000 Si aparece un problema\par}

\par
\begingroup
\fontsize{10}{12.5}\selectfont
\setlength{\tabcolsep}{5pt}
\renewcommand{\arraystretch}{1.15}
\begin{tabular}{@{}>{\RaggedRight\arraybackslash}p{0.22\linewidth}>{\RaggedRight\arraybackslash}p{0.3\linewidth}>{\RaggedRight\arraybackslash}p{0.39\linewidth}@{}}
\toprule
\textbf{Qué observa} & \textbf{Causa posible} & \textbf{Qué hacer y comprobar} \\
\midrule
No se muestran las tarjetas del asistente. & El perfil carece de consulta de propuestas o no hay fuente seleccionable. & Revise el perfil y la existencia de fuentes; gerencia puede operar el panel sin crear propuestas.\par\smallskip\textbf{Comprobación:} La pantalla corresponde a los permisos reales, no a una supuesta cuenta universal. \\
\bottomrule
\end{tabular}
\endgroup
}
\Needspace{\dimexpr\ht0+\dp0+48pt\relax}
\endgroup
\subsection{Si aparece un problema}
\begingroup
\fontsize{10}{12.5}\selectfont
\setlength{\tabcolsep}{5pt}
\renewcommand{\arraystretch}{1.15}
\begin{longtable}{@{}>{\RaggedRight\arraybackslash}p{0.22\linewidth}>{\RaggedRight\arraybackslash}p{0.3\linewidth}>{\RaggedRight\arraybackslash}p{0.39\linewidth}@{}}
\toprule
\textbf{Qué observa} & \textbf{Causa posible} & \textbf{Qué hacer y comprobar} \\
\midrule
\endfirsthead
\toprule
\textbf{Qué observa} & \textbf{Causa posible} & \textbf{Qué hacer y comprobar} \\
\midrule
\endhead
\midrule
\multicolumn{3}{r}{\footnotesize Continúa en la página siguiente}\\
\endfoot
\bottomrule
\endlastfoot
No se muestran las tarjetas del asistente. & El perfil carece de consulta de propuestas o no hay fuente seleccionable. & Revise el perfil y la existencia de fuentes; gerencia puede operar el panel sin crear propuestas.\par\smallskip\textbf{Comprobación:} La pantalla corresponde a los permisos reales, no a una supuesta cuenta universal. \\
\addlinespace[3pt]
Asistente de IA no preparado. & Configuración ausente, inactiva o sin prueba vigente. & La administración revisa Configuración LLM y realiza la prueba del registro correcto.\par\smallskip\textbf{Comprobación:} Inicio reconoce el proveedor preparado después de la prueba válida. \\
\addlinespace[3pt]
\end{longtable}
\endgroup

\subsection{Para explicar durante la sustentación}
\paragraph{¿Todo está listo significa que la IA ya conoce el negocio?} No. Solo indica preparación técnica. La necesidad y su compatibilidad con los metadatos se revisan en la etapa siguiente.

\pantalla{recursos/pantallas/01_inicio.jpg}{Inicio y fuente seleccionada}{Inicio identifica la fuente seleccionada y el estado de preparación. Cambiar la fuente cambia el contexto del recorrido. Captura de referencia del 18 de noviembre de 2026; revise las instrucciones actualizadas de esta unidad.}

\IfFileExists{capturas/operacion_02.tex}{\input{capturas/operacion_02.tex}}{}
\section{Navegar y conservar la fuente de trabajo}
\label{unidad:operacion_03}
\textbf{Para quién:} Todos los perfiles autorizados\par
\textbf{Opción:} Fuente en contexto y navegación lateral\par

La navegación organiza la plataforma por preparación del entorno, diseño y transformación BI, análisis y decisiones, y administración y seguridad. Siga ese orden al preparar una fuente nueva. Para revisar resultados existentes, puede ir directamente a la etapa de consulta correspondiente, siempre que tenga permiso.

Cambiar fuente selecciona el origen de trabajo de su sesión. No apaga las demás conexiones ni une sus registros. Las dos bases pueden permanecer habilitadas simultáneamente. La plataforma conserva por separado sus instantáneas, catálogos, propuestas, ejecuciones y resultados.

Antes de interpretar un importe, confirme fuente, propuesta y ejecución. El nombre del origen no se debe inferir a partir de un color del gráfico. En una demostración, verbalice esta identidad al cambiar de AdventureWorks a WideWorldImporters. Si el listado no muestra una fuente, no reutilice una propuesta de otra base para simular su disponibilidad.

Al recargar la pestaña, el sistema conserva el módulo autorizado. En Analítica recuerda por separado la ejecución, el indicador, el año, el territorio y la vista confirmados de cada fuente. Cambiar a otra fuente y regresar recupera su selección propia; no combina ambos análisis. Cerrar sesión limpia estos contextos locales.

\subsection{Volver al punto correcto sin perder el contexto}
Atrás aparece cuando hay un recorrido previo conocido. Por ejemplo, Inicio puede llevarle a revisar una configuración y permitirle volver al mismo paso de preparación. Desde Analítica, Consultar respaldo permite revisar el expediente y regresar al panel con la selección confirmada. Entrar directamente desde el menú lateral no inventa un paso anterior.

Los filtros se recuerdan dentro de la pestaña y por conexión. No se guardan allí las cifras, el chat ni un resultado inventado: al volver se consultan de nuevo los datos autorizados. Una selección todavía pendiente o fallida no se presenta como un análisis confirmado.

\begin{enumerate}
\item En Analítica, seleccione una ejecución de AdventureWorks y un año o territorio. Espere a que se actualicen los resultados.
\item Cambie a WideWorldImporters y confirme su encabezado antes de elegir sus propios filtros.
\item Regrese a AdventureWorks y compruebe que recupera la selección de esa fuente.
\item Abra Consultar respaldo y vuelva con Atrás. Compare ejecución, año, territorio e indicador.
\item Recargue la pestaña y compruebe que sigue en un módulo permitido. Si cierra sesión, los contextos locales se limpian.
\end{enumerate}

\subsection{Conceptos para comprender la opción}
\paragraph{Fuente} Conexión a una base de datos concreta. AdventureWorks2022 y WideWorldImporters pueden estar en el mismo servidor y seguir siendo fuentes distintas.

\paragraph{Contexto} Selección que identifica de qué fuente y ejecución proviene la información consultada.

\paragraph{Linaje} Relación entre origen, instantánea, propuesta, ejecución y resultado que permite explicar su procedencia.

\subsection{Paso a paso}
\begin{enumerate}
\item Despliegue el grupo de navegación de la tarea que realizará.
\item Use Cambiar fuente y elija el nombre de la conexión y su base.
\item Espere a que se cargue el nuevo contexto antes de leer tablas, tarjetas o gráficos.
\item En el Explorador, confirme la instantánea; en el asistente, la propuesta; en el datamart o panel, la ejecución.
\item Si necesita más espacio, utilice Ocultar navegación lateral. Para recuperarla, utilice Mostrar navegación lateral.
\item En una pantalla estrecha, el botón Menú muestra las opciones disponibles.
\item Cuando la pantalla muestre Atrás, úselo para regresar al punto que abrió la configuración o el respaldo. Si entró desde el menú lateral puede no aparecer: en ese caso continúe con el menú.
\end{enumerate}

\Needspace{8\baselineskip}
\begingroup
\setbox0=\vbox{\hsize=\linewidth
{\normalfont\normalsize\sffamily\bfseries\normalsize\color{Teal}\hyphenpenalty=10000\exhyphenpenalty=10000 Campos y botones\par}

\par
\begingroup
\fontsize{10}{12.5}\selectfont
\setlength{\tabcolsep}{5pt}
\renewcommand{\arraystretch}{1.15}
\begin{tabular}{@{}>{\RaggedRight\arraybackslash}p{0.22\linewidth}>{\RaggedRight\arraybackslash}p{0.3\linewidth}>{\RaggedRight\arraybackslash}p{0.39\linewidth}@{}}
\toprule
\textbf{Control} & \textbf{Para qué sirve} & \textbf{Ejemplo y resultado} \\
\midrule
Cambiar fuente & Selecciona una conexión habilitada para la sesión. & \textbf{Ejemplo:} Alternar AdventureWorks2022 y WideWorldImporters.\par\smallskip\textbf{Resultado:} Recarga los componentes dependientes de esa fuente. \\
\bottomrule
\end{tabular}
\endgroup
}
\Needspace{\dimexpr\ht0+\dp0+48pt\relax}
\endgroup
\subsection{Campos y botones}
\begingroup
\fontsize{10}{12.5}\selectfont
\setlength{\tabcolsep}{5pt}
\renewcommand{\arraystretch}{1.15}
\begin{longtable}{@{}>{\RaggedRight\arraybackslash}p{0.22\linewidth}>{\RaggedRight\arraybackslash}p{0.3\linewidth}>{\RaggedRight\arraybackslash}p{0.39\linewidth}@{}}
\toprule
\textbf{Control} & \textbf{Para qué sirve} & \textbf{Ejemplo y resultado} \\
\midrule
\endfirsthead
\toprule
\textbf{Control} & \textbf{Para qué sirve} & \textbf{Ejemplo y resultado} \\
\midrule
\endhead
\midrule
\multicolumn{3}{r}{\footnotesize Continúa en la página siguiente}\\
\endfoot
\bottomrule
\endlastfoot
Cambiar fuente & Selecciona una conexión habilitada para la sesión. & \textbf{Ejemplo:} Alternar AdventureWorks2022 y WideWorldImporters.\par\smallskip\textbf{Resultado:} Recarga los componentes dependientes de esa fuente. \\
\addlinespace[3pt]
Ocultar navegación lateral / Mostrar navegación lateral & Aprovecha el ancho de la pantalla. & \textbf{Ejemplo:} Examinar un modelo o gráfico amplio.\par\smallskip\textbf{Resultado:} Cambia la presentación; no cambia permisos ni resultados. \\
\addlinespace[3pt]
Menú & Abre la navegación en pantallas estrechas. & \textbf{Ejemplo:} Uso en un equipo con menor resolución.\par\smallskip\textbf{Resultado:} Presenta las opciones autorizadas. \\
\addlinespace[3pt]
Anterior / Siguiente & Recorre un listado paginado cuando tiene más registros. & \textbf{Ejemplo:} Buscar una instantánea más antigua.\par\smallskip\textbf{Resultado:} Cambia el tramo de resultados del listado, no la fuente por sí mismo. \\
\addlinespace[3pt]
\end{longtable}
\endgroup

\subsection{Comprobar el aislamiento visible}
Práctica propuesta para el entorno de capacitación; no se presenta como una nueva ejecución realizada durante la redacción.

\begin{enumerate}
\item Abra el Explorador con AdventureWorks y anote base e instantánea.
\item Cambie a WideWorldImporters y anote los mismos datos.
\item Abra el catálogo de cada fuente y compare su contexto.
\end{enumerate}

\textbf{Qué debe obtener.} Cada módulo identifica el origen elegido. No se presupone que los números de instantánea o preguntas sean iguales en otra instalación.

\textbf{Cómo comprobarlo.} La estudiante puede señalar la fuente y el artefacto correspondiente antes de describir un resultado.

\Needspace{8\baselineskip}
\begingroup
\setbox0=\vbox{\hsize=\linewidth
{\normalfont\normalsize\sffamily\bfseries\normalsize\color{Teal}\hyphenpenalty=10000\exhyphenpenalty=10000 Si aparece un problema\par}

\par
\begingroup
\fontsize{10}{12.5}\selectfont
\setlength{\tabcolsep}{5pt}
\renewcommand{\arraystretch}{1.15}
\begin{tabular}{@{}>{\RaggedRight\arraybackslash}p{0.22\linewidth}>{\RaggedRight\arraybackslash}p{0.3\linewidth}>{\RaggedRight\arraybackslash}p{0.39\linewidth}@{}}
\toprule
\textbf{Qué observa} & \textbf{Causa posible} & \textbf{Qué hacer y comprobar} \\
\midrule
Falta el selector o no aparecen fuentes. & No existen conexiones habilitadas o falta Consultar metadatos, requerido por el selector actual. & Solicite revisar las conexiones y el permiso Consultar metadatos de la cuenta. No conceda ejecución ETL para corregir un permiso de lectura.\par\smallskip\textbf{Comprobación:} La fuente puede seleccionarse sin habilitar acciones de escritura innecesarias. \\
\bottomrule
\end{tabular}
\endgroup
}
\Needspace{\dimexpr\ht0+\dp0+48pt\relax}
\endgroup
\subsection{Si aparece un problema}
\begingroup
\fontsize{10}{12.5}\selectfont
\setlength{\tabcolsep}{5pt}
\renewcommand{\arraystretch}{1.15}
\begin{longtable}{@{}>{\RaggedRight\arraybackslash}p{0.22\linewidth}>{\RaggedRight\arraybackslash}p{0.3\linewidth}>{\RaggedRight\arraybackslash}p{0.39\linewidth}@{}}
\toprule
\textbf{Qué observa} & \textbf{Causa posible} & \textbf{Qué hacer y comprobar} \\
\midrule
\endfirsthead
\toprule
\textbf{Qué observa} & \textbf{Causa posible} & \textbf{Qué hacer y comprobar} \\
\midrule
\endhead
\midrule
\multicolumn{3}{r}{\footnotesize Continúa en la página siguiente}\\
\endfoot
\bottomrule
\endlastfoot
Falta el selector o no aparecen fuentes. & No existen conexiones habilitadas o falta Consultar metadatos, requerido por el selector actual. & Solicite revisar las conexiones y el permiso Consultar metadatos de la cuenta. No conceda ejecución ETL para corregir un permiso de lectura.\par\smallskip\textbf{Comprobación:} La fuente puede seleccionarse sin habilitar acciones de escritura innecesarias. \\
\addlinespace[3pt]
Parece permanecer información de otra fuente. & Se leyó la pantalla antes de terminar la actualización, hay una incidencia de contexto o se abrió un expediente histórico distinto. & Espere la carga y compruebe fuente y expediente. Si persiste, registre pantalla y pasos y detenga la interpretación de esas cifras; informe la incidencia.\par\smallskip\textbf{Comprobación:} Los identificadores visibles y los resultados pertenecen a la misma fuente. \\
\addlinespace[3pt]
\end{longtable}
\endgroup

\subsection{Para explicar durante la sustentación}
\paragraph{¿Por qué una empresa puede tener dos fuentes activas?} Porque cada conexión representa un origen independiente. La selección de contexto permite trabajar con una sin deshabilitar la otra.

\paragraph{¿Universal significa que admite cualquier motor inmediatamente?} El flujo y los contratos se reutilizan por fuente. En este corte se implementó y validó SQL Server; otro motor necesita su conector y pruebas.


\IfFileExists{capturas/operacion_03.tex}{\input{capturas/operacion_03.tex}}{}

\chapter{Preparar las fuentes y el asistente de IA}
\section{Registrar, probar y habilitar fuentes de datos}
\label{unidad:operacion_09}
\textbf{Para quién:} Administración\par
\textbf{Opción:} Preparación del entorno > Conexiones de datos\par

Una conexión identifica dónde se encuentra una base de datos y con qué cuenta de lectura puede consultarla la aplicación. En este corte, Motor muestra SQL Server y no es una opción libre para escoger otros motores. El sistema reutiliza el mismo flujo para diferentes bases SQL Server; registrar el nombre de MySQL no incorpora un conector que no existe.

Servidor se interpreta desde el contenedor del backend. En una instalación Compose del proyecto, sqlserver es el nombre interno del servicio y 1433 su puerto interno. No lo sustituya por localhost ni por el puerto publicado del host por costumbre: localhost dentro de un contenedor se refiere a ese mismo contenedor. Para un servidor externo, la administración proporciona el host y puerto accesibles desde el backend.

AdventureWorks2022 y WideWorldImporters pueden compartir servidor y puerto, pero llevan conexiones distintas porque cambia Base de datos y pueden cambiar sus credenciales. Nunca use una cuenta administrativa de escritura para el flujo BI. La prueba verifica acceso de lectura y ausencia de privilegios de escritura relevantes; una prueba rechazada por tener escritura exige corregir el usuario del origen.

Registrar conserva una configuración deshabilitada; Probar conexión comprueba sus condiciones; Habilitar la ofrece en el selector. Son tres acciones diferentes. Varias fuentes pueden estar habilitadas a la vez. Una fila habilitada muestra Probar conexión, Leer metadatos y Deshabilitar cuando la cuenta tiene los permisos correspondientes; no muestra Editar. El código ofrece edición solo para una conexión deshabilitada. No elimine una conexión que tiene expedientes para intentar corregir una captura: el servidor controla sus dependencias.

\subsection{Conceptos para comprender la opción}
\paragraph{Servidor y base} El servidor aloja una o varias bases; la base selecciona el origen concreto.

\paragraph{Solo lectura} Cuenta destinada a consultar el origen sin modificar sus tablas.

\paragraph{Cifrar transporte} Protección de la conexión hacia el servidor.

\paragraph{Confiar en certificado local} Aceptación del certificado de la instalación local prevista; no equivale a validar públicamente su identidad.

\subsection{Paso a paso}
\begin{enumerate}
\item Abra Conexiones de datos y revise si la conexión ya está registrada.
\item Si necesita una nueva, complete Nombre visible, Servidor, Puerto, Base de datos y Usuario de sólo lectura con los datos de la instalación.
\item Introduzca la contraseña suministrada para ese usuario. Revise Cifrar transporte y Confiar en certificado local con la administración.
\item Pulse Registrar conexión y compruebe su fila.
\item Pulse Probar conexión; continúe cuando Validación indique Sólo lectura validada.
\item Pulse Habilitar para que aparezca en Cambiar fuente.
\item Pulse Leer metadatos si tiene permiso de actualización, o continúe desde el Explorador.
\item No busque Editar en una fila habilitada. Si la administración necesita cambiar esa conexión, debe planificar su deshabilitación; la edición está prevista para el estado deshabilitado. No deshabilite una fuente solo para reproducir una captura del manual.
\end{enumerate}

\Needspace{8\baselineskip}
\begingroup
\setbox0=\vbox{\hsize=\linewidth
{\normalfont\normalsize\sffamily\bfseries\normalsize\color{Teal}\hyphenpenalty=10000\exhyphenpenalty=10000 Campos y botones\par}

\par
\begingroup
\fontsize{10}{12.5}\selectfont
\setlength{\tabcolsep}{5pt}
\renewcommand{\arraystretch}{1.15}
\begin{tabular}{@{}>{\RaggedRight\arraybackslash}p{0.22\linewidth}>{\RaggedRight\arraybackslash}p{0.3\linewidth}>{\RaggedRight\arraybackslash}p{0.39\linewidth}@{}}
\toprule
\textbf{Control} & \textbf{Para qué sirve} & \textbf{Ejemplo y resultado} \\
\midrule
Nombre visible & Distingue la fuente en los módulos. & \textbf{Ejemplo:} AdventureWorks local o WideWorldImporters local.\par\smallskip\textbf{Resultado:} Se reconoce el origen sin depender del orden del selector. \\
\bottomrule
\end{tabular}
\endgroup
}
\Needspace{\dimexpr\ht0+\dp0+48pt\relax}
\endgroup
\subsection{Campos y botones}
\begingroup
\fontsize{10}{12.5}\selectfont
\setlength{\tabcolsep}{5pt}
\renewcommand{\arraystretch}{1.15}
\begin{longtable}{@{}>{\RaggedRight\arraybackslash}p{0.22\linewidth}>{\RaggedRight\arraybackslash}p{0.3\linewidth}>{\RaggedRight\arraybackslash}p{0.39\linewidth}@{}}
\toprule
\textbf{Control} & \textbf{Para qué sirve} & \textbf{Ejemplo y resultado} \\
\midrule
\endfirsthead
\toprule
\textbf{Control} & \textbf{Para qué sirve} & \textbf{Ejemplo y resultado} \\
\midrule
\endhead
\midrule
\multicolumn{3}{r}{\footnotesize Continúa en la página siguiente}\\
\endfoot
\bottomrule
\endlastfoot
Nombre visible & Distingue la fuente en los módulos. & \textbf{Ejemplo:} AdventureWorks local o WideWorldImporters local.\par\smallskip\textbf{Resultado:} Se reconoce el origen sin depender del orden del selector. \\
\addlinespace[3pt]
Motor & Indica el conector implementado. & \textbf{Ejemplo:} SQL Server, deshabilitado para edición.\par\smallskip\textbf{Resultado:} Evita presentar soporte de un motor inexistente. \\
\addlinespace[3pt]
Servidor / Puerto & Determina el destino de red visto por el backend. & \textbf{Ejemplo:} sqlserver y 1433 en Compose, cuando corresponde a esa instalación.\par\smallskip\textbf{Resultado:} Se intenta conectar al servidor definido. \\
\addlinespace[3pt]
Base de datos & Selecciona la base dentro del servidor. & \textbf{Ejemplo:} AdventureWorks2022 o WideWorldImporters.\par\smallskip\textbf{Resultado:} Los metadatos posteriores se capturan de esa base. \\
\addlinespace[3pt]
Usuario de sólo lectura / Contraseña & Acredita el acceso de consulta al origen. & \textbf{Ejemplo:} Usuario preparado por la administración de la instalación.\par\smallskip\textbf{Resultado:} La contraseña se guarda cifrada y no se devuelve al navegador. \\
\addlinespace[3pt]
Cifrar transporte / Confiar en certificado local & Configura las condiciones de transporte admitidas. & \textbf{Ejemplo:} Certificado local del contenedor de demostración.\par\smallskip\textbf{Resultado:} Permite la conexión prevista sin confundir certificado local con identidad pública. \\
\addlinespace[3pt]
Probar conexión & Comprueba lectura y permisos relevantes. & \textbf{Ejemplo:} Probar cada una de las dos fuentes.\par\smallskip\textbf{Resultado:} Registra estado y fecha de la prueba. \\
\addlinespace[3pt]
Habilitar / Deshabilitar & Ofrece o retira la fuente del contexto seleccionable. & \textbf{Ejemplo:} Mantener ambas habilitadas para alternar análisis.\par\smallskip\textbf{Resultado:} No fusiona bases ni borra sus expedientes. \\
\addlinespace[3pt]
Leer metadatos & Captura la estructura de una fuente habilitada. & \textbf{Ejemplo:} Crear la instantánea inicial de WideWorldImporters.\par\smallskip\textbf{Resultado:} Genera o reutiliza una instantánea, sin enviar filas a la IA. \\
\addlinespace[3pt]
Eliminar & Retira una conexión deshabilitada sin dependencias que lo impidan. & \textbf{Ejemplo:} Conexión temporal creada por error, con autorización.\par\smallskip\textbf{Resultado:} No sirve para borrar indiscriminadamente el historial BI. \\
\addlinespace[3pt]
\end{longtable}
\endgroup

\subsection{Reconocer dos bases en un mismo servidor}
Práctica propuesta para el entorno de capacitación; no se presenta como una nueva ejecución realizada durante la redacción.

\begin{enumerate}
\item Consulte las dos conexiones existentes.
\item Compare servidor, puerto y Base de datos.
\item Seleccione cada una desde Cambiar fuente y compruebe su nombre.
\end{enumerate}

\textbf{Qué debe obtener.} Pueden coincidir servidor y puerto, pero las bases y sus contextos son diferentes.

\textbf{Cómo comprobarlo.} La estudiante distingue conexión, servidor y base sin cambiar credenciales para la práctica.

\Needspace{8\baselineskip}
\begingroup
\setbox0=\vbox{\hsize=\linewidth
{\normalfont\normalsize\sffamily\bfseries\normalsize\color{Teal}\hyphenpenalty=10000\exhyphenpenalty=10000 Si aparece un problema\par}

\par
\begingroup
\fontsize{10}{12.5}\selectfont
\setlength{\tabcolsep}{5pt}
\renewcommand{\arraystretch}{1.15}
\begin{tabular}{@{}>{\RaggedRight\arraybackslash}p{0.22\linewidth}>{\RaggedRight\arraybackslash}p{0.3\linewidth}>{\RaggedRight\arraybackslash}p{0.39\linewidth}@{}}
\toprule
\textbf{Qué observa} & \textbf{Causa posible} & \textbf{Qué hacer y comprobar} \\
\midrule
No se puede conectar al servidor. & Host o puerto incorrectos desde el backend, servicio apagado o red no disponible. & Revise los datos con la guía de instalación. En Compose, compruebe el nombre interno del servicio; no cambie a localhost por ensayo sin criterio.\par\smallskip\textbf{Comprobación:} Probar conexión responde y registra una fecha actual. \\
\bottomrule
\end{tabular}
\endgroup
}
\Needspace{\dimexpr\ht0+\dp0+48pt\relax}
\endgroup
\subsection{Si aparece un problema}
\begingroup
\fontsize{10}{12.5}\selectfont
\setlength{\tabcolsep}{5pt}
\renewcommand{\arraystretch}{1.15}
\begin{longtable}{@{}>{\RaggedRight\arraybackslash}p{0.22\linewidth}>{\RaggedRight\arraybackslash}p{0.3\linewidth}>{\RaggedRight\arraybackslash}p{0.39\linewidth}@{}}
\toprule
\textbf{Qué observa} & \textbf{Causa posible} & \textbf{Qué hacer y comprobar} \\
\midrule
\endfirsthead
\toprule
\textbf{Qué observa} & \textbf{Causa posible} & \textbf{Qué hacer y comprobar} \\
\midrule
\endhead
\midrule
\multicolumn{3}{r}{\footnotesize Continúa en la página siguiente}\\
\endfoot
\bottomrule
\endlastfoot
No se puede conectar al servidor. & Host o puerto incorrectos desde el backend, servicio apagado o red no disponible. & Revise los datos con la guía de instalación. En Compose, compruebe el nombre interno del servicio; no cambie a localhost por ensayo sin criterio.\par\smallskip\textbf{Comprobación:} Probar conexión responde y registra una fecha actual. \\
\addlinespace[3pt]
La prueba detecta permisos de escritura. & El usuario del origen tiene privilegios superiores a los admitidos. & Solicite un usuario de lectura o ajuste sus permisos en SQL Server mediante la administración. No se corrige desmarcando Cifrar transporte.\par\smallskip\textbf{Comprobación:} La prueba acredita solo lectura antes de habilitar. \\
\addlinespace[3pt]
Habilitar no está disponible. & No hay prueba satisfactoria vigente. & Corrija la causa y use Probar conexión.\par\smallskip\textbf{Comprobación:} La fila pasa a Sólo lectura validada y permite habilitar. \\
\addlinespace[3pt]
Eliminar es rechazado. & La conexión está activa o tiene dependencias. & Deshabilite solo si corresponde y lea las dependencias. Mantenga los registros que respaldan propuestas o ejecuciones.\par\smallskip\textbf{Comprobación:} No se destruye el origen de los expedientes para resolver un problema de configuración. \\
\addlinespace[3pt]
\end{longtable}
\endgroup

\subsection{Para explicar durante la sustentación}
\paragraph{¿Registrar ya lee todas las ventas?} No. Conserva la conexión. La prueba consulta condiciones de acceso, la introspección lee estructura y la ejecución ETL posterior obtiene las filas autorizadas.

\paragraph{¿Habilitar WideWorldImporters apaga AdventureWorks?} No. Las fuentes se habilitan independientemente y la sesión elige con cuál trabajar.

\pantalla{recursos/pantallas/02_conexiones_formulario.jpg}{Registro y edición de una conexión SQL Server}{Formulario de conexión vacío. Host y puerto se refieren al acceso desde el backend, que puede estar dentro de Docker. Captura de referencia del 18 de noviembre de 2026; revise las instrucciones actualizadas de esta unidad.}
\pantalla{recursos/pantallas/03_conexiones_registradas.jpg}{Fuentes de ejemplo habilitadas y validadas como lectura}{Cada registro representa una fuente independiente. No deshabilite una fuente sólo para pasar a otra. Captura de referencia del 18 de noviembre de 2026; revise las instrucciones actualizadas de esta unidad.}

\IfFileExists{capturas/operacion_09.tex}{\input{capturas/operacion_09.tex}}{}
\section{Ajustar parámetros con un propósito concreto}
\label{unidad:operacion_10}
\textbf{Para quién:} Administración\par
\textbf{Opción:} Preparación del entorno > Parámetros\par

Los parámetros son ajustes operativos aprobados, no un lugar para guardar claves, contraseñas o fórmulas del negocio. La pantalla muestra el nombre, una explicación, el valor actual, el rango admitido y el predeterminado. Modifique un parámetro cuando pueda explicar qué problema operativo intenta resolver.

Duración de la sesión (minutos) está en la categoría Seguridad. Su valor predeterminado es 60 y admite de 15 a 240 minutos. Para una capacitación larga pueden acordarse 60 o 120 minutos. El permiso Administrar parámetros es necesario para modificarlo. Guardarlo no renueva las sesiones ya emitidas; ingrese otra vez para recibir la nueva duración.

Tiempo máximo del asistente expresa cuánto esperará la aplicación una solicitud; no concede más capacidad semántica al modelo. Si un proveedor devuelve un contrato inválido, aumentar el tiempo no vuelve correcto ese contrato. Elementos por bloque de metadatos controla la organización del contexto técnico; no selecciona cuántas filas de ventas se cargarán.

Cada tarjeta tiene su propio Guardar y Restaurar. Guardar una tarjeta no debe darse por entendido como guardar todas las demás. Anote el valor anterior y cambie una cosa a la vez para poder observar el efecto. Restaurar recupera el predeterminado de ese parámetro, no la configuración completa de la aplicación.

En el código de este corte, los listados de la pantalla principal solicitan diez registros de forma explícita. Aunque Registros por página puede guardarse, no se debe prometer que cambie todos los listados inmediatamente. La capacitación debe comprobar el efecto visible antes de utilizarlo como demostración de parametrización.

\subsection{Práctica guiada: preparar tiempo suficiente para la capacitación}
Con una cuenta autorizada, busque la tarjeta Duración de la sesión (minutos). La imagen muestra el valor 60 de la instalación revisada. En un equipo diferente puede verse otro valor válido. No cambie todos los parámetros para intentar resolver una sola incidencia.

\begin{enumerate}
\item Anote el valor actual. Para una demostración habitual puede conservar 60.
\item Si necesita una práctica más larga, acuerde 120 con la administración, escriba 120 y pulse Guardar en esa tarjeta.
\item Compruebe el mensaje y el valor guardado. El rango es de 15 a 240 minutos.
\item Vuelva a iniciar sesión antes de la práctica para utilizar la nueva duración. Registre el ajuste si forma parte de una prueba.
\end{enumerate}

\pantallaactual{recursos/pantallas/parametros-sesion-60.jpg}{Parámetros: duración de la sesión}{Captura real del 25 de noviembre de 2026. La sesión nueva dura 60 minutos por defecto. El cambio no renueva una sesión que ya estaba abierta.}
\subsection{Conceptos para comprender la opción}
\paragraph{Tiempo máximo} Límite de espera de una operación, expresado en segundos.

\paragraph{Predeterminado} Valor de referencia que se recupera al restaurar la tarjeta.

\paragraph{Bloque de metadatos} Grupo de elementos estructurales usado para organizar el contexto del asistente; no es un lote de ventas.

\subsection{Paso a paso}
\begin{enumerate}
\item Abra Parámetros y lea la tarjeta que desea ajustar.
\item Anote valor actual, unidad, rango y motivo del cambio.
\item Introduzca un entero dentro del rango.
\item Pulse Guardar en esa tarjeta y compruebe el mensaje de actualización.
\item Repita la lectura del valor. Si el cambio no se justifica, utilice Restaurar en esa misma tarjeta.
\item Verifique el efecto real en la operación correspondiente; no infiera eficacia por el solo hecho de que se guardó.
\end{enumerate}

\Needspace{8\baselineskip}
\begingroup
\setbox0=\vbox{\hsize=\linewidth
{\normalfont\normalsize\sffamily\bfseries\normalsize\color{Teal}\hyphenpenalty=10000\exhyphenpenalty=10000 Campos y botones\par}

\par
\begingroup
\fontsize{10}{12.5}\selectfont
\setlength{\tabcolsep}{5pt}
\renewcommand{\arraystretch}{1.15}
\begin{tabular}{@{}>{\RaggedRight\arraybackslash}p{0.22\linewidth}>{\RaggedRight\arraybackslash}p{0.3\linewidth}>{\RaggedRight\arraybackslash}p{0.39\linewidth}@{}}
\toprule
\textbf{Control} & \textbf{Para qué sirve} & \textbf{Ejemplo y resultado} \\
\midrule
Tiempo máximo de conexión & Limita la espera al probar una fuente. & \textbf{Ejemplo:} Predeterminado 8 segundos; rango 3-30.\par\smallskip\textbf{Resultado:} Afecta el límite de espera de la prueba, no los permisos SQL. \\
\bottomrule
\end{tabular}
\endgroup
}
\Needspace{\dimexpr\ht0+\dp0+48pt\relax}
\endgroup
\subsection{Campos y botones}
\begingroup
\fontsize{10}{12.5}\selectfont
\setlength{\tabcolsep}{5pt}
\renewcommand{\arraystretch}{1.15}
\begin{longtable}{@{}>{\RaggedRight\arraybackslash}p{0.22\linewidth}>{\RaggedRight\arraybackslash}p{0.3\linewidth}>{\RaggedRight\arraybackslash}p{0.39\linewidth}@{}}
\toprule
\textbf{Control} & \textbf{Para qué sirve} & \textbf{Ejemplo y resultado} \\
\midrule
\endfirsthead
\toprule
\textbf{Control} & \textbf{Para qué sirve} & \textbf{Ejemplo y resultado} \\
\midrule
\endhead
\midrule
\multicolumn{3}{r}{\footnotesize Continúa en la página siguiente}\\
\endfoot
\bottomrule
\endlastfoot
Tiempo máximo de conexión & Limita la espera al probar una fuente. & \textbf{Ejemplo:} Predeterminado 8 segundos; rango 3-30.\par\smallskip\textbf{Resultado:} Afecta el límite de espera de la prueba, no los permisos SQL. \\
\addlinespace[3pt]
Tiempo máximo del asistente & Limita la espera de la solicitud al proveedor. & \textbf{Ejemplo:} Predeterminado 600 segundos; rango 10-900.\par\smallskip\textbf{Resultado:} Permite acotar la espera sin garantizar éxito. \\
\addlinespace[3pt]
Elementos por bloque de metadatos & Organiza la cantidad de elementos técnicos por bloque. & \textbf{Ejemplo:} Predeterminado 100; rango 25-200.\par\smallskip\textbf{Resultado:} Cambia el tamaño operativo del contexto, no las filas del ETL. \\
\addlinespace[3pt]
Registros por página & Conserva el ajuste de tamaño predeterminado de listados. & \textbf{Ejemplo:} Predeterminado 10; rango 10-100.\par\smallskip\textbf{Resultado:} El efecto depende del consumidor; los paginadores actuales envían 10 explícitamente. \\
\addlinespace[3pt]
Guardar & Persiste el valor de esa tarjeta. & \textbf{Ejemplo:} Guardar un tiempo acordado para una prueba.\par\smallskip\textbf{Resultado:} Muestra un mensaje con el nombre del parámetro actualizado. \\
\addlinespace[3pt]
Restaurar & Recupera el predeterminado de la tarjeta. & \textbf{Ejemplo:} Volver Tiempo máximo de conexión a 8.\par\smallskip\textbf{Resultado:} Se deshabilita si ya coincide con el predeterminado. \\
\addlinespace[3pt]
Duración de la sesión (minutos) & Fija la duración absoluta de los siguientes inicios de sesión. & \textbf{Ejemplo:} 60 para el uso habitual; 120 si se acuerda una práctica extensa.\par\smallskip\textbf{Resultado:} Guardar aplica el valor a sesiones nuevas, dentro del rango 15-240. \\
\addlinespace[3pt]
\end{longtable}
\endgroup

\subsection{Cambiar y recuperar un parámetro sin llamar a la IA}
Práctica propuesta para el entorno de capacitación; no se presenta como una nueva ejecución realizada durante la redacción.

\begin{enumerate}
\item En un entorno de capacitación, anote Tiempo máximo de conexión.
\item Guarde un valor permitido distinto y compruebe el mensaje.
\item Pulse Restaurar y compruebe el predeterminado.
\end{enumerate}

\textbf{Qué debe obtener.} Se conserva el dato operativo y se puede volver al valor de referencia. No se cambian catálogos, fuentes ni propuestas.

\textbf{Cómo comprobarlo.} La estudiante explica qué cambió y qué no demuestra este ejercicio.

\Needspace{8\baselineskip}
\begingroup
\setbox0=\vbox{\hsize=\linewidth
{\normalfont\normalsize\sffamily\bfseries\normalsize\color{Teal}\hyphenpenalty=10000\exhyphenpenalty=10000 Si aparece un problema\par}

\par
\begingroup
\fontsize{10}{12.5}\selectfont
\setlength{\tabcolsep}{5pt}
\renewcommand{\arraystretch}{1.15}
\begin{tabular}{@{}>{\RaggedRight\arraybackslash}p{0.22\linewidth}>{\RaggedRight\arraybackslash}p{0.3\linewidth}>{\RaggedRight\arraybackslash}p{0.39\linewidth}@{}}
\toprule
\textbf{Qué observa} & \textbf{Causa posible} & \textbf{Qué hacer y comprobar} \\
\midrule
El valor es rechazado. & Está fuera del rango o no es un entero admitido. & Utilice el rango visible; no agregue una clave nueva.\par\smallskip\textbf{Comprobación:} Se guarda un valor válido y legible. \\
\bottomrule
\end{tabular}
\endgroup
}
\Needspace{\dimexpr\ht0+\dp0+48pt\relax}
\endgroup
\subsection{Si aparece un problema}
\begingroup
\fontsize{10}{12.5}\selectfont
\setlength{\tabcolsep}{5pt}
\renewcommand{\arraystretch}{1.15}
\begin{longtable}{@{}>{\RaggedRight\arraybackslash}p{0.22\linewidth}>{\RaggedRight\arraybackslash}p{0.3\linewidth}>{\RaggedRight\arraybackslash}p{0.39\linewidth}@{}}
\toprule
\textbf{Qué observa} & \textbf{Causa posible} & \textbf{Qué hacer y comprobar} \\
\midrule
\endfirsthead
\toprule
\textbf{Qué observa} & \textbf{Causa posible} & \textbf{Qué hacer y comprobar} \\
\midrule
\endhead
\midrule
\multicolumn{3}{r}{\footnotesize Continúa en la página siguiente}\\
\endfoot
\bottomrule
\endlastfoot
El valor es rechazado. & Está fuera del rango o no es un entero admitido. & Utilice el rango visible; no agregue una clave nueva.\par\smallskip\textbf{Comprobación:} Se guarda un valor válido y legible. \\
\addlinespace[3pt]
El modelo sigue fallando tras aumentar el tiempo. & El problema puede ser cuota, clave, modelo o contrato estructurado. & Revise el mensaje del proveedor y su prueba de configuración. No atribuya todos los fallos a tiempo insuficiente.\par\smallskip\textbf{Comprobación:} La reparación se dirige a la causa observada. \\
\addlinespace[3pt]
Guardar Registros por página no cambia el listado. & El componente utiliza su tamaño explícito del corte actual. & Registre esta limitación; conserve el valor de referencia si no hay una prueba de efecto.\par\smallskip\textbf{Comprobación:} La demostración no promete un cambio que el navegador no muestra. \\
\addlinespace[3pt]
\end{longtable}
\endgroup

\subsection{Para explicar durante la sustentación}
\paragraph{¿Subir el tiempo hace más inteligente a la IA?} No. Solo modifica el límite de espera. La calidad de una propuesta sigue dependiendo del contexto, del proveedor y de las reglas de validación.

\pantalla{recursos/pantallas/04_parametros.jpg}{Parámetros operativos y rangos admitidos}{Los parámetros explican límites operativos. Un plazo mayor no corrige una credencial o un contrato inválido. Captura de referencia del 18 de noviembre de 2026; revise las instrucciones actualizadas de esta unidad.}
\pantalla{recursos/pantallas/05_parametros_detalle.jpg}{Bloques de metadatos y paginación configurables}{Revise el alcance de cada parámetro antes de guardar. La paginación de algunos listados conserva un tamaño fijado por la interfaz. Captura de referencia del 18 de noviembre de 2026; revise las instrucciones actualizadas de esta unidad.}

\IfFileExists{capturas/operacion_10.tex}{\input{capturas/operacion_10.tex}}{}
\section{Configurar y comprobar el asistente de IA}
\label{unidad:operacion_11}
\textbf{Para quién:} Administración\par
\textbf{Opción:} Preparación del entorno > Configuración LLM\par

El modelo de lenguaje interpreta y explica; no autentica usuarios ni ejecuta SQL libre sobre las bases. La plataforma mantiene proveedores como configuraciones separadas y una sola configuración activa a la vez. Cambiar la activa afecta las nuevas solicitudes de la aplicación, por lo que debe coordinarse antes de una demostración compartida.

Seleccione el Proveedor para completar la URL y el modelo sugeridos por el producto. Son valores del corte de software, no una garantía permanente de disponibilidad comercial. El identificador debe existir y estar autorizado en la cuenta del proveedor. La URL se valida contra destinos permitidos; no se usa esta pantalla para enviar datos a un servicio arbitrario.

Para un proveedor cloud, el flujo consiste en crear el registro, registrar la credencial, probar la conexión y activar la configuración prevista. La prueba hace una solicitud mínima, sin el contexto comercial de la demostración. Un resultado satisfactorio acredita ese acceso básico, pero no demuestra todavía la calidad de una propuesta de BI.

Ollama local no necesita API key, pero sí requiere el servicio local y el modelo descargado en la instalación. La pantalla no descarga modelos por escribir su nombre. Si se cambia proveedor, URL, modelo o razonamiento, se invalida la prueba previa; vuelva a probar antes de usar el asistente. Cambiar de tipo de proveedor en un registro activo exige desactivarlo primero.

Nivel de razonamiento es una opción común que los adaptadores traducen a las capacidades admitidas. La modalidad de Claude de este corte ofrece Automático, Mínimo o Bajo; no habilita pensamiento extendido con Medio o Alto. Una opción más alta puede cambiar tiempo y consumo en proveedores compatibles, pero no garantiza mejores respuestas ni sustituye la supervisión.

La edición del 25 de noviembre incorpora OpenAI como proveedor propio. La recomendación económica actual es gpt-6-luna con razonamiento Bajo. Su conexión y una respuesta estructurada mínima fueron comprobadas. Los recorridos nuevos completos con ambas fuentes se comprobaron con gpt-6.2-sol y razonamiento Medio; esto no acredita todavía un recorrido completo con Luna. Una prueba de conexión satisfactoria no reemplaza la revisión de cada propuesta.

\subsection{Configurar OpenAI paso a paso y controlar el consumo}
Una suscripción de ChatGPT y el uso de la API son servicios con facturación y disponibilidad propias. Un crédito promocional sólo debe darse por confirmado al consultar la cuenta API que lo recibió. La plataforma utiliza la API key de esa cuenta; no emplea automáticamente el modelo o la suscripción de este chat.

OpenAI se conecta al servicio oficial https://api.openai.com/v1 mediante Responses API. La credencial se introduce en Registrar credencial y se almacena cifrada; no se vuelve a mostrar. No la copie en el objetivo comercial, el catálogo, una captura ni este manual. La configuración no permite elegir cualquier destino de Internet.

Groq y OpenAI son proveedores distintos. Un modelo cuyo nombre contenga openai dentro de Groq sigue consumiendo la cuenta y cuota de Groq. Para utilizar su crédito API de OpenAI debe elegir el proveedor OpenAI.

La plataforma solicita una respuesta estructurada y vuelve a comprobarla localmente. Si hay rechazo del modelo, respuesta incompleta, cuota agotada o contrato inválido, el fallo debe quedar visible. No se cambia a un modelo más caro de manera automática para ocultarlo.

\begin{enumerate}
\item Abra Configuración LLM y revise si ya existe una configuración OpenAI. Evite crear duplicados.
\item Para crearla, escriba un nombre reconocible, por ejemplo OpenAI capacitación, y elija Proveedor: OpenAI.
\item Revise la URL oficial. Para la opción económica escriba gpt-6-luna y seleccione razonamiento Bajo. El modelo debe estar disponible para su cuenta.
\item Pulse Crear registro. Luego, en su fila, abra Registrar credencial, introduzca la API key privadamente y guárdela.
\item Pulse Probar conexión una sola vez. Una respuesta válida demuestra acceso básico y salida estructurada mínima, no calidad de toda la propuesta BI.
\item Pulse Activar sobre la configuración que utilizará el grupo. Sólo una queda activa para las nuevas solicitudes de toda la aplicación.
\item Antes de enviar la necesidad, revise el aviso sobre el contexto que se transferirá. Para diseño se utilizan necesidad y metadatos permitidos; no se envían contraseñas de las bases.
\item Si cambia modelo, URL, proveedor, razonamiento o credencial, repita la prueba. Si falla por cuota no reintente sin revisar primero su cuenta.
\end{enumerate}

\Needspace{8\baselineskip}
\begingroup
\setbox0=\vbox{\hsize=\linewidth
{\normalfont\normalsize\sffamily\bfseries\normalsize\color{Teal}\hyphenpenalty=10000\exhyphenpenalty=10000 Opciones OpenAI y qué se comprobó realmente\par}

\par
\begingroup
\fontsize{10}{12.5}\selectfont
\setlength{\tabcolsep}{5pt}
\renewcommand{\arraystretch}{1.15}
\begin{tabular}{@{}>{\RaggedRight\arraybackslash}p{0.22\linewidth}>{\RaggedRight\arraybackslash}p{0.3\linewidth}>{\RaggedRight\arraybackslash}p{0.39\linewidth}@{}}
\toprule
\textbf{Modelo y nivel} & \textbf{Tarifa estándar por millón de tokens} & \textbf{Evidencia de esta edición} \\
\midrule
GPT-6 Luna; Bajo & Entrada USD 0,10; salida USD 0,50 & Opción económica actual. Una prueba real de conexión y JSON estructurado; sin recorrido BI completo con Luna. \\
\bottomrule
\end{tabular}
\endgroup
}
\Needspace{\dimexpr\ht0+\dp0+48pt\relax}
\endgroup
\subsection{Opciones OpenAI y qué se comprobó realmente}
\begingroup
\fontsize{10}{12.5}\selectfont
\setlength{\tabcolsep}{5pt}
\renewcommand{\arraystretch}{1.15}
\begin{longtable}{@{}>{\RaggedRight\arraybackslash}p{0.22\linewidth}>{\RaggedRight\arraybackslash}p{0.3\linewidth}>{\RaggedRight\arraybackslash}p{0.39\linewidth}@{}}
\toprule
\textbf{Modelo y nivel} & \textbf{Tarifa estándar por millón de tokens} & \textbf{Evidencia de esta edición} \\
\midrule
\endfirsthead
\toprule
\textbf{Modelo y nivel} & \textbf{Tarifa estándar por millón de tokens} & \textbf{Evidencia de esta edición} \\
\midrule
\endhead
\midrule
\multicolumn{3}{r}{\footnotesize Continúa en la página siguiente}\\
\endfoot
\bottomrule
\endlastfoot
GPT-6 Luna; Bajo & Entrada USD 0,10; salida USD 0,50 & Opción económica actual. Una prueba real de conexión y JSON estructurado; sin recorrido BI completo con Luna. \\
\addlinespace[3pt]
GPT-5.4 mini & Entrada USD 0,75; salida USD 4,50 & Alternativa disponible según cuenta. No ensayada en esta campaña. \\
\addlinespace[3pt]
GPT-6.1 Terra & Entrada USD 1; salida USD 5 & No ensayada en esta campaña. \\
\addlinespace[3pt]
GPT-6.2 Sol; Medio & Entrada USD 2; salida USD 10 & Modelo utilizado en los dos recorridos nuevos completos: AdventureWorks 136/24 y WideWorldImporters 135/23. \\
\addlinespace[3pt]
GPT-6.2 Astra & Entrada USD 8; salida USD 40 & No ensayada. No se recomienda para este presupuesto de demostración. \\
\addlinespace[3pt]
\end{longtable}
\endgroup

Tarifas de texto estándar consultadas el 25 de noviembre de 2026, sin caché ni otros descuentos. Con un supuesto idéntico de 10.000 tokens de entrada y 2.000 de salida, el costo base sería USD 0,002; 0,0165; 0,020; 0,040 y 0,160, respectivamente. Es una comparación estimada, no el costo medido del proyecto. El supuesto no incluye razonamiento adicional, reintentos, herramientas ni llamadas extras.

Las tarifas de entrada y salida de Luna son 95 \% menores que las de Sol. Esto no demuestra igual calidad ni un ahorro real idéntico por tarea: cambian la longitud, el razonamiento y el número de llamadas. No se conservó un registro de costos por solicitud de las pruebas, por lo que el consumo real debe consultarse en Usage/Billing de la cuenta API.

En modelos de razonamiento, la salida disponible comparte presupuesto con el razonamiento interno. El adaptador reserva un límite mínimo según el nivel: 4.096 tokens en Automático, Mínimo o Bajo; 8.192 en Medio; 16.384 en Alto, con techo de 32.768. Es un límite permitido, no una cantidad que siempre se facture. Un proceso puede realizar más de una llamada.

Para gastar menos durante la sustentación, consulte propuestas aprobadas, expedientes y paneles publicados. Abrirlos no vuelve a generar la propuesta con IA. Sugerir una necesidad, verificar viabilidad mediante IA, generar una propuesta o conversar con el copiloto sí puede consumir la API. Use sólo la operación que necesita y no repita clics mientras espera.

\pantallaactual{recursos/pantallas/luna-modelo-bajo.jpg}{Configuración activa: OpenAI económico}{Captura real del 25 de noviembre de 2026. La fila OpenAI BI - Luna económico muestra razonamiento bajo y credencial configurada. No se muestra la clave ni debe copiarse al manual.}
\subsection{Conceptos para comprender la opción}
\paragraph{Proveedor} Servicio que ofrece el modelo, local o cloud.

\paragraph{Modelo} Identificador concreto de la capacidad solicitada dentro de ese proveedor.

\paragraph{Credencial} Clave que autoriza al proveedor; se almacena cifrada y no se vuelve a mostrar.

\paragraph{Activa} Configuración seleccionada para nuevas solicitudes; no significa que todas sus respuestas serán válidas.

\subsection{Paso a paso}
\begin{enumerate}
\item Abra Configuración LLM y revise los registros existentes para no duplicarlos innecesariamente.
\item En Crear configuración LLM, escriba Nombre de configuración y seleccione Proveedor.
\item Revise URL del servicio, Modelo y Nivel de razonamiento. Use valores autorizados para la instalación.
\item Pulse Crear registro. En cloud, pulse Registrar credencial en su fila, introduzca la API key y pulse Guardar credencial.
\item Pulse Probar conexión y lea el resultado. Corrija clave, cuota, modelo o disponibilidad si falla.
\item Pulse Activar sobre la configuración prevista y compruebe cuál queda activa.
\item Vuelva a Inicio para revisar el asistente preparado. Antes de enviar una necesidad, atienda la autorización de transferencia del contexto que corresponda al destino elegido.
\item Para una clave renovada, use Reemplazar credencial y repita la prueba. No necesita borrar la propuesta histórica para cambiar una configuración de futuras solicitudes.
\end{enumerate}

\Needspace{8\baselineskip}
\begingroup
\setbox0=\vbox{\hsize=\linewidth
{\normalfont\normalsize\sffamily\bfseries\normalsize\color{Teal}\hyphenpenalty=10000\exhyphenpenalty=10000 Campos y botones\par}

\par
\begingroup
\fontsize{10}{12.5}\selectfont
\setlength{\tabcolsep}{5pt}
\renewcommand{\arraystretch}{1.15}
\begin{tabular}{@{}>{\RaggedRight\arraybackslash}p{0.22\linewidth}>{\RaggedRight\arraybackslash}p{0.3\linewidth}>{\RaggedRight\arraybackslash}p{0.39\linewidth}@{}}
\toprule
\textbf{Control} & \textbf{Para qué sirve} & \textbf{Ejemplo y resultado} \\
\midrule
Nombre de configuración & Distingue el registro. & \textbf{Ejemplo:} Claude demostración o Groq capacitación.\par\smallskip\textbf{Resultado:} Permite reconocer la configuración que se prueba y activa. \\
\bottomrule
\end{tabular}
\endgroup
}
\Needspace{\dimexpr\ht0+\dp0+48pt\relax}
\endgroup
\subsection{Campos y botones}
\begingroup
\fontsize{10}{12.5}\selectfont
\setlength{\tabcolsep}{5pt}
\renewcommand{\arraystretch}{1.15}
\begin{longtable}{@{}>{\RaggedRight\arraybackslash}p{0.22\linewidth}>{\RaggedRight\arraybackslash}p{0.3\linewidth}>{\RaggedRight\arraybackslash}p{0.39\linewidth}@{}}
\toprule
\textbf{Control} & \textbf{Para qué sirve} & \textbf{Ejemplo y resultado} \\
\midrule
\endfirsthead
\toprule
\textbf{Control} & \textbf{Para qué sirve} & \textbf{Ejemplo y resultado} \\
\midrule
\endhead
\midrule
\multicolumn{3}{r}{\footnotesize Continúa en la página siguiente}\\
\endfoot
\bottomrule
\endlastfoot
Nombre de configuración & Distingue el registro. & \textbf{Ejemplo:} Claude demostración o Groq capacitación.\par\smallskip\textbf{Resultado:} Permite reconocer la configuración que se prueba y activa. \\
\addlinespace[3pt]
Proveedor & Selecciona el adaptador implementado. & \textbf{Ejemplo:} OpenAI, Anthropic Claude, Groq Cloud, Gemini Cloud, Qwen Cloud u Ollama local.\par\smallskip\textbf{Resultado:} Completa valores recomendados de URL, modelo y nivel. \\
\addlinespace[3pt]
URL del servicio & Identifica el destino permitido. & \textbf{Ejemplo:} URL completada al seleccionar el proveedor.\par\smallskip\textbf{Resultado:} Se valida su correspondencia; no admite un destino arbitrario. \\
\addlinespace[3pt]
Modelo & Especifica el identificador que el servicio debe reconocer. & \textbf{Ejemplo:} Modelo autorizado de la cuenta y del adaptador.\par\smallskip\textbf{Resultado:} La prueba puede detectar ausencia, acceso o incompatibilidad. \\
\addlinespace[3pt]
Nivel de razonamiento & Elige una modalidad admitida. & \textbf{Ejemplo:} Mínimo para una solicitud breve en la modalidad de Claude disponible.\par\smallskip\textbf{Resultado:} Se conserva en la configuración y forma parte de lo que debe volver a probarse si cambia. \\
\addlinespace[3pt]
Registrar credencial / Reemplazar credencial & Abre el formulario de la clave cloud. & \textbf{Ejemplo:} Registrar una clave nueva sin imprimirla ni pegarla en documentación.\par\smallskip\textbf{Resultado:} El servidor conserva el secreto cifrado. \\
\addlinespace[3pt]
Guardar credencial / Cancelar & Confirma o abandona el envío de la clave. & \textbf{Ejemplo:} Cancelar si se abrió la fila equivocada.\par\smallskip\textbf{Resultado:} La clave no vuelve a mostrarse en el listado. \\
\addlinespace[3pt]
Probar conexión & Realiza la comprobación mínima del proveedor y modelo. & \textbf{Ejemplo:} Comprobar un modelo local instalado.\par\smallskip\textbf{Resultado:} Devuelve un resultado básico, sin generar un datamart. \\
\addlinespace[3pt]
Activar / Desactivar & Selecciona la configuración operativa. & \textbf{Ejemplo:} Cambiar de proveedor con conocimiento de quienes usan la plataforma.\par\smallskip\textbf{Resultado:} Solo una queda activa para nuevas solicitudes. \\
\addlinespace[3pt]
Eliminar & Retira una configuración inactiva cuando el servidor lo permite. & \textbf{Ejemplo:} Limpiar un registro temporal autorizado.\par\smallskip\textbf{Resultado:} No sirve para reparar automáticamente un proveedor indisponible. \\
\addlinespace[3pt]
\end{longtable}
\endgroup

\subsection{Identificar los componentes de una configuración existente}
Práctica propuesta para el entorno de capacitación; no se presenta como una nueva ejecución realizada durante la redacción.

\begin{enumerate}
\item Consulte un registro sin mostrar la clave.
\item Identifique proveedor, modelo, nivel y estado.
\item Explique cuál usaría la próxima solicitud y qué ocurriría al cambiar el modelo.
\end{enumerate}

\textbf{Qué debe obtener.} Se distingue registro, prueba, credencial y activación. No se envía una solicitud comercial por consultar el listado.

\textbf{Cómo comprobarlo.} La estudiante explica por qué una prueba básica exitosa no es una evaluación completa de las propuestas BI.

\Needspace{8\baselineskip}
\begingroup
\setbox0=\vbox{\hsize=\linewidth
{\normalfont\normalsize\sffamily\bfseries\normalsize\color{Teal}\hyphenpenalty=10000\exhyphenpenalty=10000 Si aparece un problema\par}

\par
\begingroup
\fontsize{10}{12.5}\selectfont
\setlength{\tabcolsep}{5pt}
\renewcommand{\arraystretch}{1.15}
\begin{tabular}{@{}>{\RaggedRight\arraybackslash}p{0.22\linewidth}>{\RaggedRight\arraybackslash}p{0.3\linewidth}>{\RaggedRight\arraybackslash}p{0.39\linewidth}@{}}
\toprule
\textbf{Qué observa} & \textbf{Causa posible} & \textbf{Qué hacer y comprobar} \\
\midrule
Credencial pendiente o rechazo de autenticación. & No se registró una clave válida para ese proveedor. & Registre o reemplace la clave en la fila correcta y pruebe otra vez. No la escriba en Parámetros ni en el objetivo de negocio.\par\smallskip\textbf{Comprobación:} La prueba básica resulta satisfactoria sin exponer el secreto. \\
\bottomrule
\end{tabular}
\endgroup
}
\Needspace{\dimexpr\ht0+\dp0+48pt\relax}
\endgroup
\subsection{Si aparece un problema}
\begingroup
\fontsize{10}{12.5}\selectfont
\setlength{\tabcolsep}{5pt}
\renewcommand{\arraystretch}{1.15}
\begin{longtable}{@{}>{\RaggedRight\arraybackslash}p{0.22\linewidth}>{\RaggedRight\arraybackslash}p{0.3\linewidth}>{\RaggedRight\arraybackslash}p{0.39\linewidth}@{}}
\toprule
\textbf{Qué observa} & \textbf{Causa posible} & \textbf{Qué hacer y comprobar} \\
\midrule
\endfirsthead
\toprule
\textbf{Qué observa} & \textbf{Causa posible} & \textbf{Qué hacer y comprobar} \\
\midrule
\endhead
\midrule
\multicolumn{3}{r}{\footnotesize Continúa en la página siguiente}\\
\endfoot
\bottomrule
\endlastfoot
Credencial pendiente o rechazo de autenticación. & No se registró una clave válida para ese proveedor. & Registre o reemplace la clave en la fila correcta y pruebe otra vez. No la escriba en Parámetros ni en el objetivo de negocio.\par\smallskip\textbf{Comprobación:} La prueba básica resulta satisfactoria sin exponer el secreto. \\
\addlinespace[3pt]
Modelo no disponible o sin acceso. & Identificador incorrecto, no instalado o no autorizado para la cuenta. & Revise el modelo con la administración; en Ollama compruebe servicio y descarga mediante la guía técnica.\par\smallskip\textbf{Comprobación:} El servicio reconoce el modelo configurado. \\
\addlinespace[3pt]
Cuota, límite de servicio o tiempo agotado. & Restricción externa o recursos insuficientes. & Lea la causa y espere o cambie a una configuración previamente preparada con autorización. No describa un cambio silencioso de proveedor como comportamiento del sistema.\par\smallskip\textbf{Comprobación:} La nueva configuración queda identificada y probada. \\
\addlinespace[3pt]
El proveedor devuelve JSON, pero el contrato es rechazado. & La salida no cumple los campos, tipos o reglas que exige la plataforma. & Mantenga la necesidad guardada y revise el diagnóstico de recuperación del asistente. Reintente de forma controlada o elija un proveedor probado; no apruebe una respuesta inválida para continuar.\par\smallskip\textbf{Comprobación:} Se obtiene una versión válida y revisable o un fallo explicado; no se materializa contenido inválido. \\
\addlinespace[3pt]
Activar no devuelve un estado de preparación correcto. & La clave puede existir, pero la prueba es inválida o quedó obsoleta tras un cambio. & Pruebe exactamente la combinación que está activa y revise Inicio.\par\smallskip\textbf{Comprobación:} La preparación reconoce proveedor y modelo comprobados. \\
\addlinespace[3pt]
\end{longtable}
\endgroup

\subsection{Para explicar durante la sustentación}
\paragraph{¿Se reemplaza la IA sin cambiar el proceso BI?} Los adaptadores permiten cambiar de proveedor manteniendo los contratos de la aplicación. Cada combinación requiere prueba y validación; no se garantiza que todos los modelos produzcan la misma calidad.

\paragraph{¿La prueba mínima autoriza enviar datos de ventas completos?} No. La prueba utiliza una solicitud mínima. Las solicitudes posteriores respetan su propósito y autorización: el diseño usa contexto estructural y el copiloto analítico utiliza agregados permitidos.

\pantalla{recursos/pantallas/06_llm_formulario.jpg}{Proveedor, URL, modelo y razonamiento}{Configuración del proveedor, servicio, modelo y razonamiento. Seleccione opciones compatibles con el servicio disponible. Captura de referencia del 18 de noviembre de 2026; revise las instrucciones actualizadas de esta unidad.}
\pantalla{recursos/pantallas/07_llm_registros.jpg}{Configuraciones disponibles y proveedor activo}{Configuraciones registradas: una aparece activa. No se promete una salida idéntica entre proveedores. Captura de referencia del 18 de noviembre de 2026; revise las instrucciones actualizadas de esta unidad.}
\pantalla{recursos/pantallas/08_llm_credencial_vacia.jpg}{Formulario de credencial vacío y cancelación}{Panel de sustitución abierto con el campo vacío y cancelado. La credencial se introduce de forma privada; nunca se copia al manual. Captura de referencia del 18 de noviembre de 2026; revise las instrucciones actualizadas de esta unidad.}

\IfFileExists{capturas/operacion_11.tex}{\input{capturas/operacion_11.tex}}{}
\begingroup
\setlength{\parskip}{4pt plus 1pt minus 1pt}
\setlist[enumerate]{itemsep=3pt,topsep=4pt}
\titlespacing*{\paragraph}{0pt}{4pt}{0.5em}
\section{Leer y comprender el esquema antes de diseñar}
\label{unidad:operacion_12}
\textbf{Para quién:} Analista de BI, Administración, Gerencia con consulta de metadatos\par
\textbf{Opción:} Preparación del entorno > Explorador de esquema\par

El Explorador muestra una fotografía estructural de la fuente: esquemas, tablas, columnas, tipos, claves principales y relaciones declaradas. Es el primer paso para comprobar qué información existe y cómo puede relacionarse. No muestra las filas del negocio ni demuestra por sí solo que un atributo tenga el significado comercial que usted necesita.

Un esquema es una agrupación de objetos dentro de la base. Una tabla reúne registros de una entidad o evento; una columna representa un atributo. La clave principal identifica una fila de la tabla. Una relación declarada vincula columnas de tablas mediante una clave foránea. Dos columnas con el mismo nombre no prueban una relación válida si el esquema no la respalda.

La instantánea conserva versión, fecha y huella de contenido. Si la estructura no cambió, una actualización puede reutilizar la captura existente; no es obligatorio que aumente el número de versión. Versionar metadatos permite saber con qué estructura se revisó una propuesta histórica. Actualizar metadatos no actualiza las filas del datamart ni repite el ETL.

En AdventureWorks, una línea de pedido y su cabecera permiten distinguir la cantidad de líneas de la cantidad de pedidos. En WideWorldImporters, la línea de factura corresponde a otro evento comercial. Estos son ejemplos de lectura del esquema: el programa no exige esos nombres a una fuente nueva, pero sí necesita relaciones y tipos compatibles con la necesidad.

Para la demostración, encuentre al menos el detalle comercial, su fecha de evento y una relación hacia producto, cliente o territorio. Explique qué significa cada enlace antes de usarlo. Una relación hacia la persona que creó o modificó un registro no demuestra quién compró ni dónde se realizó la operación.

\subsection{Conceptos para comprender la opción}
\paragraph{Tabla} Conjunto de filas del mismo tipo de entidad o evento.

\paragraph{Columna} Atributo con un nombre y un tipo de dato.

\paragraph{PK} Clave principal; identifica una fila de la tabla. Puede incluir más de una columna.

\paragraph{Relación saliente} Enlace desde la tabla consultada hacia otra mediante sus columnas.

\paragraph{Relación entrante} Enlace de otra tabla hacia la tabla que se consulta.

\paragraph{Huella} Resumen calculado del contenido estructural para comprobar si la captura cambió.

\subsection{Paso a paso}
\begin{enumerate}
\item Seleccione la fuente correcta en Cambiar fuente.
\item Si no hay metadatos, pulse Crear primera instantánea si tiene permiso de actualización.
\item Lea el resumen de esquemas, tablas, columnas y relaciones y la versión vigente.
\item Use Versión de metadatos para una captura histórica cuando exista más de una.
\item En Buscar tabla o columna escriba parte del nombre; opcionalmente escriba el nombre exacto del esquema en Esquema y pulse Buscar.
\item Seleccione una tabla de la lista y revise Columna, Tipo, Permite nulos y Clave.
\item Lea Relaciones salientes y Relaciones entrantes. Anote las columnas del enlace y la tabla relacionada.
\item Para quitar un filtro, vacíe los campos y pulse Buscar. Al cambiar de fuente, compruebe que se cargan su propia instantánea y tablas.
\end{enumerate}

\Needspace{8\baselineskip}
\begingroup
\setbox0=\vbox{\hsize=\linewidth
{\normalfont\normalsize\sffamily\bfseries\normalsize\color{Teal}\hyphenpenalty=10000\exhyphenpenalty=10000 Campos y botones\par}

\par
\begingroup
\fontsize{10}{12.5}\selectfont
\setlength{\tabcolsep}{5pt}
\renewcommand{\arraystretch}{1.15}
\begin{tabular}{@{}>{\RaggedRight\arraybackslash}p{0.22\linewidth}>{\RaggedRight\arraybackslash}p{0.3\linewidth}>{\RaggedRight\arraybackslash}p{0.39\linewidth}@{}}
\toprule
\textbf{Control} & \textbf{Para qué sirve} & \textbf{Ejemplo y resultado} \\
\midrule
Crear primera instantánea & Captura estructura cuando aún no existe una versión. & \textbf{Ejemplo:} Preparar una fuente SQL Server nueva.\par\smallskip\textbf{Resultado:} Conserva tablas, columnas y relaciones, sin filas ni credenciales. \\
\bottomrule
\end{tabular}
\endgroup
}
\Needspace{\dimexpr\ht0+\dp0+48pt\relax}
\endgroup
\subsection{Campos y botones}
\begingroup
\fontsize{10}{12.5}\selectfont
\setlength{\tabcolsep}{5pt}
\renewcommand{\arraystretch}{1.15}
\begin{longtable}{@{}>{\RaggedRight\arraybackslash}p{0.22\linewidth}>{\RaggedRight\arraybackslash}p{0.3\linewidth}>{\RaggedRight\arraybackslash}p{0.39\linewidth}@{}}
\toprule
\textbf{Control} & \textbf{Para qué sirve} & \textbf{Ejemplo y resultado} \\
\midrule
\endfirsthead
\toprule
\textbf{Control} & \textbf{Para qué sirve} & \textbf{Ejemplo y resultado} \\
\midrule
\endhead
\midrule
\multicolumn{3}{r}{\footnotesize Continúa en la página siguiente}\\
\endfoot
\bottomrule
\endlastfoot
Crear primera instantánea & Captura estructura cuando aún no existe una versión. & \textbf{Ejemplo:} Preparar una fuente SQL Server nueva.\par\smallskip\textbf{Resultado:} Conserva tablas, columnas y relaciones, sin filas ni credenciales. \\
\addlinespace[3pt]
Actualizar metadatos & Vuelve a consultar la estructura de la fuente. & \textbf{Ejemplo:} Se agregó una tabla o cambió una relación declarada.\par\smallskip\textbf{Resultado:} Crea una versión si hay cambios o reutiliza la equivalente. \\
\addlinespace[3pt]
Versión de metadatos & Permite revisar una captura conservada. & \textbf{Ejemplo:} Abrir la estructura de una propuesta histórica.\par\smallskip\textbf{Resultado:} Con una única versión, el selector está deshabilitado y lo explica. \\
\addlinespace[3pt]
Buscar tabla o columna & Filtra por un texto de estructura. & \textbf{Ejemplo:} SalesOrderDetail o InvoiceLines en los ejemplos instalados.\par\smallskip\textbf{Resultado:} Presenta coincidencias, no filas comerciales. \\
\addlinespace[3pt]
Esquema & Acota la búsqueda por agrupación estructural. & \textbf{Ejemplo:} Sales, si corresponde a la fuente.\par\smallskip\textbf{Resultado:} No es un selector de base de datos. \\
\addlinespace[3pt]
Buscar & Aplica los campos escritos. & \textbf{Ejemplo:} Buscar InvoiceLines dentro de Sales.\par\smallskip\textbf{Resultado:} Actualiza la lista desde la instantánea elegida. \\
\addlinespace[3pt]
Tablas & Permite elegir el objeto que se examina. & \textbf{Ejemplo:} Detalle de una venta o factura.\par\smallskip\textbf{Resultado:} Presenta columnas y relaciones de la tabla seleccionada. \\
\addlinespace[3pt]
Relaciones salientes / Relaciones entrantes & Explica los enlaces declarados. & \textbf{Ejemplo:} Detalle hacia cabecera; otra tabla hacia cliente.\par\smallskip\textbf{Resultado:} Permite reconocer el sentido y columnas de una relación sin inventarla. \\
\addlinespace[3pt]
\end{longtable}
\endgroup

\subsection{Comparar pedido y factura}
Ejercicio de consulta propuesto; los nombres corresponden a las dos fuentes públicas de demostración y deben comprobarse en la instalación.

\begin{enumerate}
\item Seleccione AdventureWorks y busque SalesOrderDetail; revise sus claves y un enlace con SalesOrderHeader.
\item Cambie a WideWorldImporters y busque InvoiceLines; revise sus claves y un enlace con Invoices.
\item Explique por qué una línea no es un pedido ni una factura completa.
\end{enumerate}

\textbf{Qué debe obtener.} Se distinguen las estructuras de dos eventos comerciales. La estudiante identifica las referencias de la fuente sin proponer todavía un ETL.

\textbf{Cómo comprobarlo.} Puede justificar un conteo distinto de documentos frente a un conteo de líneas.

\Needspace{8\baselineskip}
\begingroup
\setbox0=\vbox{\hsize=\linewidth
{\normalfont\normalsize\sffamily\bfseries\normalsize\color{Teal}\hyphenpenalty=10000\exhyphenpenalty=10000 Si aparece un problema\par}

\par
\begingroup
\fontsize{10}{12.5}\selectfont
\setlength{\tabcolsep}{5pt}
\renewcommand{\arraystretch}{1.15}
\begin{tabular}{@{}>{\RaggedRight\arraybackslash}p{0.22\linewidth}>{\RaggedRight\arraybackslash}p{0.3\linewidth}>{\RaggedRight\arraybackslash}p{0.39\linewidth}@{}}
\toprule
\textbf{Qué observa} & \textbf{Causa posible} & \textbf{Qué hacer y comprobar} \\
\midrule
No aparecen coincidencias. & Texto o esquema no coinciden con la instantánea. & Vacíe los filtros, pulse Buscar y localice una tabla de la lista. No concluya que faltan las ventas por una búsqueda mal escrita.\par\smallskip\textbf{Comprobación:} La lista vuelve a mostrar objetos de la fuente correcta. \\
\bottomrule
\end{tabular}
\endgroup
}
\Needspace{\dimexpr\ht0+\dp0+48pt\relax}
\endgroup
\subsection{Si aparece un problema}
\begingroup
\fontsize{10}{12.5}\selectfont
\setlength{\tabcolsep}{5pt}
\renewcommand{\arraystretch}{1.15}
\begin{longtable}{@{}>{\RaggedRight\arraybackslash}p{0.22\linewidth}>{\RaggedRight\arraybackslash}p{0.3\linewidth}>{\RaggedRight\arraybackslash}p{0.39\linewidth}@{}}
\toprule
\textbf{Qué observa} & \textbf{Causa posible} & \textbf{Qué hacer y comprobar} \\
\midrule
\endfirsthead
\toprule
\textbf{Qué observa} & \textbf{Causa posible} & \textbf{Qué hacer y comprobar} \\
\midrule
\endhead
\midrule
\multicolumn{3}{r}{\footnotesize Continúa en la página siguiente}\\
\endfoot
\bottomrule
\endlastfoot
No aparecen coincidencias. & Texto o esquema no coinciden con la instantánea. & Vacíe los filtros, pulse Buscar y localice una tabla de la lista. No concluya que faltan las ventas por una búsqueda mal escrita.\par\smallskip\textbf{Comprobación:} La lista vuelve a mostrar objetos de la fuente correcta. \\
\addlinespace[3pt]
Versión de metadatos aparece deshabilitada. & Solo existe una captura disponible. & Lea la nota y consulte esa versión. No necesita generar versiones artificiales.\par\smallskip\textbf{Comprobación:} Se reconoce la única instantánea sin interpretarla como un fallo. \\
\addlinespace[3pt]
No hay una relación de negocio que el analista esperaba. & El origen no la declara o enlaza otro rol del negocio. & Registre la limitación y formule una necesidad compatible. No una tablas por nombres parecidos para forzar cobertura.\par\smallskip\textbf{Comprobación:} La propuesta posterior usa enlaces respaldados o informa la insuficiencia. \\
\addlinespace[3pt]
Actualizar no cambia el número de versión. & La huella estructural es equivalente. & Compruebe el mensaje y la estructura; una captura reutilizada es un resultado válido.\par\smallskip\textbf{Comprobación:} No se confunde ausencia de cambio de estructura con error de actualización de ventas. \\
\addlinespace[3pt]
\end{longtable}
\endgroup

\subsection{Para explicar durante la sustentación}
\paragraph{¿La instantánea permite comprobar que los importes están bien?} Permite comprobar estructura y referencias. La comprobación numérica corresponde a la conciliación posterior a una carga.

\paragraph{¿Por qué conservar nombres técnicos en inglés?} Identifican las referencias reales del origen. La explicación puede estar en español, pero no debe cambiar los nombres que permiten rastrear tablas y columnas.

\endgroup
\pantalla{recursos/pantallas/09_metadatos_AdventureWorks.jpg}{Instantáneas y resumen de metadatos de AdventureWorks}{AdventureWorks: instantánea estructural con seis esquemas y 71 tablas. Los recuentos no son cantidades de ventas. Captura de referencia del 18 de noviembre de 2026; revise las instrucciones actualizadas de esta unidad.}
\pantalla{recursos/pantallas/10_metadatos_AW_columnas.jpg}{Tipos, nulos y clave del detalle de pedidos}{SalesOrderDetail: claves y tipos de las columnas. El nombre de una columna es una pista que debe contrastarse con su significado. Captura de referencia del 18 de noviembre de 2026; revise las instrucciones actualizadas de esta unidad.}
\pantalla{recursos/pantallas/11_metadatos_AW_relaciones.jpg}{Relaciones de la línea con pedido y producto}{Las relaciones declaradas sustentan rutas seguras. Observe las columnas de origen y destino, no sólo el nombre de la tabla. Captura de referencia del 18 de noviembre de 2026; revise las instrucciones actualizadas de esta unidad.}
\pantalla{recursos/pantallas/12_metadatos_WWI.jpg}{Instantánea independiente de WideWorldImporters}{WideWorldImporters: cuatro esquemas y 48 tablas. Una única instantánea deja el selector visible sin alternativas. Captura de referencia del 18 de noviembre de 2026; revise las instrucciones actualizadas de esta unidad.}
\pantalla{recursos/pantallas/13_metadatos_WWI_columnas.jpg}{Columnas de InvoiceLines y significado de sus importes}{InvoiceLines: detalle físico de las facturas. No confunda una línea con una factura completa. Captura de referencia del 18 de noviembre de 2026; revise las instrucciones actualizadas de esta unidad.}
\pantalla{recursos/pantallas/14_metadatos_WWI_relaciones.jpg}{Rutas declaradas de InvoiceLines a sus entidades}{Relaciones de InvoiceLines con sus entidades. La ruta comercial se revisa según la necesidad, aunque exista una FK. Captura de referencia del 18 de noviembre de 2026; revise las instrucciones actualizadas de esta unidad.}

\IfFileExists{capturas/operacion_12.tex}{\input{capturas/operacion_12.tex}}{}

\chapter{Configurar el catálogo con criterio comercial}
\section{Construir un catálogo de preguntas con criterio de negocio}
\label{unidad:operacion_13}
\textbf{Para quién:} Analista responsable del catálogo con permiso de administración, Administración, Consulta para comprender el alcance\par
\textbf{Opción:} Diseño y transformación BI > Catálogo analítico > Preguntas de negocio\par

El catálogo es un conjunto de orientaciones comerciales que el analista podrá seleccionar al formular una necesidad. No es el datamart, no es una consulta SQL y no determina automáticamente qué tablas se cargarán. Su función es describir qué decisiones interesa apoyar con esa fuente.

La pantalla separa dos planos. Cobertura técnica de la fuente muestra qué referencias se descubren en la instantánea. Preguntas de negocio contiene las orientaciones que una persona autorizada administra. La IA utiliza el alcance y los metadatos en las etapas de necesidad y propuesta; la pantalla del catálogo no tiene un botón que genere libremente su configuración mediante la IA.

La línea base del producto contiene cuatro preguntas: evolución temporal, productos, clientes y territorio. No fueron creadas por un LLM en cada base. En el entorno histórico de AdventureWorks se incorporaron otras dos orientaciones, sobre costos y rentabilidad y sobre ventas y costos por unidad, por decisión del analista a partir de una recomendación conversacional. Ver seis preguntas en una fuente no significa que la plataforma deba copiar seis a todas las demás.

Cada catálogo pertenece a una conexión y a un dominio. Antes de editar, confirme el nombre en Fuente en contexto. Guardar un catálogo de AdventureWorks no debe modificar WideWorldImporters. Una nueva fuente comienza con las orientaciones del perfil y requiere su propio criterio comercial; la evidencia técnica se obtiene de su estructura, no de la fuente anterior.

Verificable significa que el descubrimiento dispone de referencias que pueden revisarse; Por validar conserva una orientación cuya evidencia tendrá que comprobarse antes de aprobar una propuesta; No disponible informa insuficiencia técnica. Ninguno de esos rótulos demuestra que ya se calcularon ventas o que la interpretación comercial sea incuestionable. Abra Referencias detectadas y lea su alcance.

Redacte una pregunta para una decisión concreta. Productos con mayor desempeño es una orientación; una instrucción útil especifica si se comparan sumas de importe o de unidades y advierte que un ranking no prueba preferencias del consumidor. No introduzca nombres de columnas inventadas ni convierta una limitación en una orden para que la IA la ignore.

Habilitada conserva la pregunta y controla si se ofrece para nuevos análisis. Retirar pregunta la quita del borrador de catálogo y el cambio se hace persistente al Guardar catálogo. Deshabilitar es adecuado cuando quiere conservar la redacción para un uso futuro. No utilice retirar como si fuera una aprobación o rechazo de las propuestas que ya existen.

Restaurar catálogo validado aplica inmediatamente la configuración predeterminada del perfil a la fuente elegida; no es un historial de deshacer ni recupera automáticamente las últimas preguntas personalizadas. En este corte puede sustituir las dos preguntas adicionales de AdventureWorks por las cuatro iniciales. Antes de utilizarlo, conserve por escrito las orientaciones que desea volver a incorporar y confirme que está trabajando en la fuente adecuada.

\subsection{Si una pregunta fue retirada del catálogo}
Un borrador local puede conservar preguntas que ya no existen en el catálogo de esa fuente. Esto ocurrió con una necesidad restaurada de WideWorldImporters. El rechazo se produjo antes de llamar al modelo: no era evidencia de corrupción del datamart ni prueba de un fallo específico de Claude.

Ahora el sistema conserva el objetivo y las preguntas que siguen disponibles, retira del borrador las que ya no están y muestra un aviso. La viabilidad anterior deja de servir para esa selección modificada: revise otra vez el alcance y la autorización antes de generar. No sustituye una pregunta retirada por otra con un nombre parecido.

Consultar una propuesta histórica conserva sus preguntas y decisiones de entonces. No se modifica ni se retira su aprobación sólo porque el catálogo actual haya cambiado. Un borrador nuevo y una versión aprobada son objetos diferentes.

\begin{enumerate}
\item Si aparece un aviso de preguntas retiradas, lea cuáles cambiaron y compruebe la fuente.
\item Vuelva a la necesidad; conserve o ajuste el objetivo y elija sólo preguntas disponibles que realmente correspondan.
\item Repita la comprobación de viabilidad. Lea las limitaciones antes de continuar.
\item Si desea estudiar el caso aprobado anterior, ábralo en consulta: no lo regenere para intentar recuperar el catálogo.
\end{enumerate}

\subsection{Conceptos para comprender la opción}
\paragraph{Pregunta de negocio} Tema comercial que orienta una necesidad, sin fijar por sí mismo el diseño físico. Comparar el importe vendido por producto y mes.

\paragraph{Etiqueta} Nombre breve que el analista reconocerá en la selección de preguntas.

\paragraph{Explicación para el analista} Describe para qué decisión resulta útil la pregunta.

\paragraph{Instrucción funcional para la IA} Precisa el resultado y sus límites en términos comerciales; no habilita SQL ni referencias inexistentes.

\paragraph{Cobertura técnica} Respaldo estructural detectado para las orientaciones del dominio.

\paragraph{Personalización por fuente} Configuración independiente que pertenece a una conexión concreta.

\subsection{Paso a paso}
\begin{enumerate}
\item Abra Catálogo analítico y confirme Fuente en contexto y Dominio habilitado.
\item Lea Cobertura técnica de la fuente. Abra Referencias detectadas en las tarjetas que la ofrezcan.
\item Revise las preguntas existentes antes de añadir otra: puede bastar con mejorar su explicación.
\item Pulse Agregar pregunta. Sustituya la redacción de ejemplo de los tres campos por una orientación comercial concreta.
\item Escriba Etiqueta, Explicación para el analista e Instrucción funcional para la IA. Revise Habilitada.
\item Pulse Guardar catálogo y espere el mensaje de actualización.
\item Vuelva al catálogo y compruebe los textos. Después, abra un análisis nuevo para comprobar qué orientaciones ofrece el asistente.
\item Cambie a la otra fuente y confirme que sus preguntas se mantienen propias.
\item Si necesita retirar una orientación, considere primero desmarcar Habilitada. Si elige Retirar pregunta, revise el borrador restante y pulse Guardar catálogo.
\item No utilice Restaurar catálogo validado como prueba de aprendizaje sobre un catálogo con personalizaciones que necesite conservar.
\end{enumerate}

\Needspace{8\baselineskip}
\begingroup
\setbox0=\vbox{\hsize=\linewidth
{\normalfont\normalsize\sffamily\bfseries\normalsize\color{Teal}\hyphenpenalty=10000\exhyphenpenalty=10000 Campos y botones\par}

\par
\begingroup
\fontsize{10}{12.5}\selectfont
\setlength{\tabcolsep}{5pt}
\renewcommand{\arraystretch}{1.15}
\begin{tabular}{@{}>{\RaggedRight\arraybackslash}p{0.22\linewidth}>{\RaggedRight\arraybackslash}p{0.3\linewidth}>{\RaggedRight\arraybackslash}p{0.39\linewidth}@{}}
\toprule
\textbf{Control} & \textbf{Para qué sirve} & \textbf{Ejemplo y resultado} \\
\midrule
Dominio habilitado & Selecciona un perfil analítico implementado. & \textbf{Ejemplo:} Datamart de ventas, único dominio operativo del corte.\par\smallskip\textbf{Resultado:} No crea inventario, compras u otro dominio por cambiar una etiqueta. \\
\bottomrule
\end{tabular}
\endgroup
}
\Needspace{\dimexpr\ht0+\dp0+48pt\relax}
\endgroup
\subsection{Campos y botones}
\begingroup
\fontsize{10}{12.5}\selectfont
\setlength{\tabcolsep}{5pt}
\renewcommand{\arraystretch}{1.15}
\begin{longtable}{@{}>{\RaggedRight\arraybackslash}p{0.22\linewidth}>{\RaggedRight\arraybackslash}p{0.3\linewidth}>{\RaggedRight\arraybackslash}p{0.39\linewidth}@{}}
\toprule
\textbf{Control} & \textbf{Para qué sirve} & \textbf{Ejemplo y resultado} \\
\midrule
\endfirsthead
\toprule
\textbf{Control} & \textbf{Para qué sirve} & \textbf{Ejemplo y resultado} \\
\midrule
\endhead
\midrule
\multicolumn{3}{r}{\footnotesize Continúa en la página siguiente}\\
\endfoot
\bottomrule
\endlastfoot
Dominio habilitado & Selecciona un perfil analítico implementado. & \textbf{Ejemplo:} Datamart de ventas, único dominio operativo del corte.\par\smallskip\textbf{Resultado:} No crea inventario, compras u otro dominio por cambiar una etiqueta. \\
\addlinespace[3pt]
Cobertura técnica de la fuente & Presenta el descubrimiento de referencias de la instantánea vigente. & \textbf{Ejemplo:} Comprobar qué estructura respalda análisis temporal o territorial.\par\smallskip\textbf{Resultado:} Distingue evidencia técnica de una preferencia comercial escrita. \\
\addlinespace[3pt]
Verificable / Por validar / No disponible & Resume el respaldo de una orientación. & \textbf{Ejemplo:} Una pregunta personalizada puede permanecer Por validar.\par\smallskip\textbf{Resultado:} Guía la revisión; no acredita una cifra calculada. \\
\addlinespace[3pt]
Referencias detectadas & Despliega las referencias que sustentan una tarjeta. & \textbf{Ejemplo:} Leer las tablas o columnas citadas y su pertenencia al origen.\par\smallskip\textbf{Resultado:} Permite una comprobación explícita del respaldo. \\
\addlinespace[3pt]
Agregar pregunta & Añade un borrador de orientación. & \textbf{Ejemplo:} Una orientación sobre documentos e importe promedio.\par\smallskip\textbf{Resultado:} Genera un identificador interno; máximo doce preguntas en total, no doce adicionales. \\
\addlinespace[3pt]
Habilitada & Ofrece o retira de las selecciones futuras sin eliminar el texto. & \textbf{Ejemplo:} Deshabilitar temporalmente una orientación de costos.\par\smallskip\textbf{Resultado:} Conserva la redacción al guardar, pero no la ofrece como pregunta habilitada. \\
\addlinespace[3pt]
Etiqueta & Nombre reconocible de la orientación; de 3 a 120 caracteres. & \textbf{Ejemplo:} Costos y rentabilidad de las ventas.\par\smallskip\textbf{Resultado:} Se muestra en el catálogo y en las selecciones que correspondan. \\
\addlinespace[3pt]
Explicación para el analista & Explica utilidad y decisión; de 10 a 300 caracteres. & \textbf{Ejemplo:} Permite comparar el importe vendido con costos respaldados para revisar diferencias por producto.\par\smallskip\textbf{Resultado:} La persona comprende para qué seleccionar esa orientación. \\
\addlinespace[3pt]
Instrucción funcional para la IA & Delimita la interpretación comercial; de 10 a 500 caracteres. & \textbf{Ejemplo:} Usar únicamente costos verificables; declarar si son actuales y no atribuirles validez histórica.\par\smallskip\textbf{Resultado:} Orienta las solicitudes futuras sin sustituir los validadores. \\
\addlinespace[3pt]
Identificador interno & Mantiene la identidad de la pregunta. & \textbf{Ejemplo:} Código automático custom\_\allowbreak{}… de una pregunta nueva.\par\smallskip\textbf{Resultado:} No es un campo que deba inventar ni editar la operadora. \\
\addlinespace[3pt]
Retirar pregunta & Quita la tarjeta del borrador. & \textbf{Ejemplo:} Retirar una orientación duplicada con autorización.\par\smallskip\textbf{Resultado:} Necesita Guardar catálogo para persistir; no permite retirar la última tarjeta por ese botón. \\
\addlinespace[3pt]
Guardar catálogo & Persiste preguntas y periodicidades de la fuente. & \textbf{Ejemplo:} Guardar una pregunta aclarada en WideWorldImporters.\par\smallskip\textbf{Resultado:} Los nuevos análisis usan la configuración actualizada; no reescribe propuestas previas. \\
\addlinespace[3pt]
Restaurar catálogo validado & Reaplica la línea base del perfil a la fuente actual. & \textbf{Ejemplo:} Recuperar las cuatro orientaciones estándar en un entorno limpio.\par\smallskip\textbf{Resultado:} Sustituye personalizaciones del catálogo; no revierte una propuesta ni una ejecución. \\
\addlinespace[3pt]
\end{longtable}
\endgroup

\subsection{Examinar el catálogo sin modificarlo}
Práctica propuesta para el entorno de capacitación; no se presenta como una nueva ejecución realizada durante la redacción.

\begin{enumerate}
\item Seleccione AdventureWorks y cuente las orientaciones que realmente conserva la instalación.
\item Identifique las cuatro estándar y las personalizadas si existen.
\item Cambie a WideWorldImporters y compare etiquetas, instrucciones y cobertura.
\end{enumerate}

\textbf{Qué debe obtener.} Se reconoce la configuración propia de cada origen. No se asume que la cantidad seis sea obligatoria o generada por IA.

\textbf{Cómo comprobarlo.} Puede explicar quién redacta las preguntas y de dónde procede la evidencia técnica.

\subsection{Añadir una orientación comercial acotada}
Práctica propuesta para el entorno de capacitación; no se presenta como una nueva ejecución realizada durante la redacción.

\begin{enumerate}
\item En el entorno de capacitación, seleccione una fuente y conserve los textos iniciales.
\item Agregue una pregunta sobre importe por documento, indicando conteo distinto de pedidos o facturas según el evento real.
\item Guarde el catálogo y compruebe la nueva tarjeta.
\item Consulte la otra fuente para verificar que no recibió la tarjeta.
\item En la etapa de necesidad, someta cualquier nueva solicitud a su evaluación de viabilidad.
\end{enumerate}

\textbf{Qué debe obtener.} La pregunta queda guardada únicamente en la fuente elegida. La tarjeta no crea automáticamente una medida, propuesta o ETL. Una cobertura Por validar no se describe como fallo si aún no se ha construido una propuesta.

\textbf{Cómo comprobarlo.} La estudiante distingue guardar orientación de demostrar su factibilidad sobre metadatos.

\Needspace{8\baselineskip}
\begingroup
\setbox0=\vbox{\hsize=\linewidth
{\normalfont\normalsize\sffamily\bfseries\normalsize\color{Teal}\hyphenpenalty=10000\exhyphenpenalty=10000 Si aparece un problema\par}

\par
\begingroup
\fontsize{10}{12.5}\selectfont
\setlength{\tabcolsep}{5pt}
\renewcommand{\arraystretch}{1.15}
\begin{tabular}{@{}>{\RaggedRight\arraybackslash}p{0.22\linewidth}>{\RaggedRight\arraybackslash}p{0.3\linewidth}>{\RaggedRight\arraybackslash}p{0.39\linewidth}@{}}
\toprule
\textbf{Qué observa} & \textbf{Causa posible} & \textbf{Qué hacer y comprobar} \\
\midrule
La pregunta aparece Por validar. & No tiene evidencia técnica suficiente en el descubrimiento inicial o es una orientación personalizada. & Revise sus términos y referencias; formule una necesidad concreta y utilice la evaluación de viabilidad. No la fuerce como Verificable por cambiar su etiqueta.\par\smallskip\textbf{Comprobación:} La factibilidad se acredita mediante referencias y reglas, o se informa una limitación. \\
\bottomrule
\end{tabular}
\endgroup
}
\Needspace{\dimexpr\ht0+\dp0+48pt\relax}
\endgroup
\subsection{Si aparece un problema}
\begingroup
\fontsize{10}{12.5}\selectfont
\setlength{\tabcolsep}{5pt}
\renewcommand{\arraystretch}{1.15}
\begin{longtable}{@{}>{\RaggedRight\arraybackslash}p{0.22\linewidth}>{\RaggedRight\arraybackslash}p{0.3\linewidth}>{\RaggedRight\arraybackslash}p{0.39\linewidth}@{}}
\toprule
\textbf{Qué observa} & \textbf{Causa posible} & \textbf{Qué hacer y comprobar} \\
\midrule
\endfirsthead
\toprule
\textbf{Qué observa} & \textbf{Causa posible} & \textbf{Qué hacer y comprobar} \\
\midrule
\endhead
\midrule
\multicolumn{3}{r}{\footnotesize Continúa en la página siguiente}\\
\endfoot
\bottomrule
\endlastfoot
La pregunta aparece Por validar. & No tiene evidencia técnica suficiente en el descubrimiento inicial o es una orientación personalizada. & Revise sus términos y referencias; formule una necesidad concreta y utilice la evaluación de viabilidad. No la fuerce como Verificable por cambiar su etiqueta.\par\smallskip\textbf{Comprobación:} La factibilidad se acredita mediante referencias y reglas, o se informa una limitación. \\
\addlinespace[3pt]
No se puede guardar el catálogo. & Campos fuera de longitud, tarjetas inválidas, fuente no habilitada o falta de permiso. & Lea el mensaje, complete los tres campos de cada pregunta y revise la fuente. Compruebe también las periodicidades del mismo formulario.\par\smallskip\textbf{Comprobación:} Aparece El catálogo analítico fue actualizado y una nueva consulta conserva los textos. \\
\addlinespace[3pt]
Faltan las dos preguntas adicionales después de restaurar. & Restaurar sustituyó la personalización por la línea base de cuatro preguntas. & Recupere los textos conservados o una evidencia histórica y vuelva a incorporarlos en la fuente correcta; no anuncie una recuperación automática inexistente.\par\smallskip\textbf{Comprobación:} Las orientaciones adicionales vuelven a estar guardadas y habilitadas solo donde corresponde. \\
\addlinespace[3pt]
La pregunta pide costos, pero no hay un costo confiable. & El deseo comercial supera la evidencia del origen o solo existe un costo actual. & Acote la necesidad a las medidas disponibles o declare el alcance del costo actual. No llame costo histórico de la venta a un atributo actual sin respaldo.\par\smallskip\textbf{Comprobación:} La propuesta no publica un margen supuesto ni confunde factibilidad estructural con criterio contable. \\
\addlinespace[3pt]
Agregar pregunta está deshabilitado. & Ya existen doce tarjetas. & Revise duplicados y retire o reformule una orientación si corresponde. Deshabilitar una tarjeta no reduce el número total.\par\smallskip\textbf{Comprobación:} El catálogo cumple el máximo de doce, con las preguntas necesarias. \\
\addlinespace[3pt]
\end{longtable}
\endgroup

\subsection{Para explicar durante la sustentación}
\paragraph{¿La IA construyó las cuatro preguntas iniciales?} No. Proceden del perfil del producto. Una persona autorizada puede ampliarlas; la IA interviene al sugerir o revisar necesidades y al proponer el diseño a partir de los metadatos.

\paragraph{¿Por qué una pregunta guardada no demuestra que puede responderse?} Porque el catálogo expresa intención. La evaluación de viabilidad y los validadores deben comprobar referencias, relaciones, grano y cálculos para esa necesidad.

\paragraph{¿Por qué el catálogo no fija las tablas?} Para separar el propósito comercial del esquema físico. El mismo objetivo puede requerir estructuras diferentes en cada fuente.

\paragraph{¿Qué ocurre al restaurar?} Se aplica la configuración predeterminada del perfil en esa fuente; las personalizaciones del catálogo pueden sustituirse. Las propuestas y ejecuciones previas no se reescriben.

\pantalla{recursos/pantallas/15_catalogo_WWI_contexto.jpg}{Catálogo separado para WideWorldImporters}{Catálogo en el contexto de WideWorldImporters. Cada fuente conserva sus propias preguntas y periodicidades. Captura de referencia del 18 de noviembre de 2026; revise las instrucciones actualizadas de esta unidad.}
\pantalla{recursos/pantallas/16_catalogo_cobertura_AW.jpg}{Cobertura potencial de las preguntas según metadatos}{Cobertura técnica calculada sobre AdventureWorks. Esta comprobación orienta el diseño y no sustituye la conciliación numérica. Captura de referencia del 18 de noviembre de 2026; revise las instrucciones actualizadas de esta unidad.}
\pantalla{recursos/pantallas/17_catalogo_preguntas.jpg}{Etiqueta, explicación e instrucción funcional de cada pregunta}{Etiqueta, explicación e instrucción funcional tienen fines distintos. La pregunta orienta a la IA sin fijar SQL ni tablas inventadas. Captura de referencia del 18 de noviembre de 2026; revise las instrucciones actualizadas de esta unidad.}

\IfFileExists{capturas/operacion_13.tex}{\input{capturas/operacion_13.tex}}{}
\begingroup
\setlength{\parskip}{4pt plus 1pt minus 1pt}
\setlist[enumerate]{itemsep=3pt,topsep=4pt}
\titlespacing*{\paragraph}{0pt}{4pt}{0.5em}
\section{Elegir periodicidades sin confundirlas con el detalle de los datos}
\label{unidad:operacion_14}
\textbf{Para quién:} Analista responsable del catálogo, Administración, Personas en formación BI\par
\textbf{Opción:} Diseño y transformación BI > Catálogo analítico > Periodicidades\par

Periodicidad indica cómo se agrupan temporalmente los resultados: día, semana, mes, trimestre o año. Granularidad indica qué representa una fila del hecho, como una línea de pedido o una línea de factura. Frecuencia de actualización indica cuándo se ejecutaría otra carga. Son tres decisiones distintas y no deben utilizarse como sinónimos.

Un caso puede conservar cada línea de venta y presentar sus importes agrupados mensualmente. En ese caso, Mensual no reduce por sí sola el hecho a una fila por mes. Tampoco agenda una carga mensual: el prototipo realiza cargas completas confirmadas por una persona; no implementa un calendario de refresco automático por habilitar esta opción.

La etiqueta temporal debe corresponder al evento de negocio. En AdventureWorks puede usarse la fecha del pedido cuando esa es la elegida y validada; en WideWorldImporters puede ser la fecha de la factura. Una comparación mensual de pedidos no debe llamarse facturación mensual solo porque ambas usan meses. La propuesta y su expediente conservan la fecha de evento y el alcance.

Las estrategias están implementadas como un catálogo cerrado de cinco opciones. Se pueden ofrecer varias y retirar alguna, pero escribir Bimestral como etiqueta de Mensual no crea una agregación bimestral. La etiqueta ayuda a leer la opción; no cambia el código de cálculo que la plataforma admite.

En un análisis particular se selecciona una de las periodicidades habilitadas. Los KPI pueden declarar una periodicidad explícita o heredar la de la propuesta; el lector analítico resuelve esa configuración al agrupar los resultados. Cambiar la configuración de un catálogo no debe reetiquetar una ejecución histórica que ya conserva su propio contrato.

\subsection{Conceptos para comprender la opción}
\paragraph{Granularidad} Significado de una fila del hecho. Una línea de pedido puede contener cantidad e importe.

\paragraph{Periodicidad} Agrupación temporal usada para presentar un resultado. Sumar importes de las líneas correspondientes a un mes.

\paragraph{Fecha de evento} Fecha cuyo significado comercial organiza el análisis. Fecha del pedido frente a fecha de factura.

\paragraph{Actualización} Obtención de datos mediante otra ejecución. No es un efecto de habilitar Mensual en el catálogo.

\paragraph{Herencia de periodicidad} Un indicador utiliza la periodicidad de su propuesta cuando así lo declara su receta.

\subsection{Paso a paso}
\begin{enumerate}
\item Revise primero qué fecha comercial respalda la fuente y la necesidad.
\item En Periodicidades lea cada opción y su Explicación.
\item Conserve habilitadas las que desea ofrecer al analista y desmarque las que no correspondan.
\item Para reincorporar una estrategia retirada, selecciónela en Agregar periodicidad y pulse Agregar.
\item Revise la Etiqueta y Explicación sin cambiar su significado temporal.
\item Pulse Guardar catálogo y compruebe las opciones disponibles al iniciar una necesidad nueva.
\item En la propuesta, confirme aparte la granularidad. En el expediente y el panel, compruebe la periodicidad efectiva de la serie.
\end{enumerate}

\Needspace{8\baselineskip}
\begingroup
\setbox0=\vbox{\hsize=\linewidth
{\normalfont\normalsize\sffamily\bfseries\normalsize\color{Teal}\hyphenpenalty=10000\exhyphenpenalty=10000 Campos y botones\par}

\par
\begingroup
\fontsize{10}{12.5}\selectfont
\setlength{\tabcolsep}{5pt}
\renewcommand{\arraystretch}{1.15}
\begin{tabular}{@{}>{\RaggedRight\arraybackslash}p{0.22\linewidth}>{\RaggedRight\arraybackslash}p{0.3\linewidth}>{\RaggedRight\arraybackslash}p{0.39\linewidth}@{}}
\toprule
\textbf{Control} & \textbf{Para qué sirve} & \textbf{Ejemplo y resultado} \\
\midrule
Habilitada & Ofrece una periodicidad para nuevas selecciones. & \textbf{Ejemplo:} Mantener Mensual habilitada para una comparación comercial.\par\smallskip\textbf{Resultado:} Conserva la estrategia y controla su disponibilidad. \\
\bottomrule
\end{tabular}
\endgroup
}
\Needspace{\dimexpr\ht0+\dp0+48pt\relax}
\endgroup
\subsection{Campos y botones}
\begingroup
\fontsize{10}{12.5}\selectfont
\setlength{\tabcolsep}{5pt}
\renewcommand{\arraystretch}{1.15}
\begin{longtable}{@{}>{\RaggedRight\arraybackslash}p{0.22\linewidth}>{\RaggedRight\arraybackslash}p{0.3\linewidth}>{\RaggedRight\arraybackslash}p{0.39\linewidth}@{}}
\toprule
\textbf{Control} & \textbf{Para qué sirve} & \textbf{Ejemplo y resultado} \\
\midrule
\endfirsthead
\toprule
\textbf{Control} & \textbf{Para qué sirve} & \textbf{Ejemplo y resultado} \\
\midrule
\endhead
\midrule
\multicolumn{3}{r}{\footnotesize Continúa en la página siguiente}\\
\endfoot
\bottomrule
\endlastfoot
Habilitada & Ofrece una periodicidad para nuevas selecciones. & \textbf{Ejemplo:} Mantener Mensual habilitada para una comparación comercial.\par\smallskip\textbf{Resultado:} Conserva la estrategia y controla su disponibilidad. \\
\addlinespace[3pt]
Etiqueta & Nombre legible de la periodicidad; de 3 a 100 caracteres. & \textbf{Ejemplo:} Mensual.\par\smallskip\textbf{Resultado:} No altera por sí mismo la agrupación implementada. \\
\addlinespace[3pt]
Explicación & Describe cuándo y para qué usar la opción; de 10 a 240 caracteres. & \textbf{Ejemplo:} Agrupa por mes usando una fecha comercial verificable.\par\smallskip\textbf{Resultado:} Orienta la selección del analista. \\
\addlinespace[3pt]
Retirar periodicidad & Quita una estrategia del borrador de configuración. & \textbf{Ejemplo:} Retirar temporalmente Diaria si no se ofrecerá.\par\smallskip\textbf{Resultado:} No permite retirar la última tarjeta desde ese botón; el cambio necesita guardarse. \\
\addlinespace[3pt]
Agregar periodicidad & Permite elegir una de las estrategias implementadas ausentes. & \textbf{Ejemplo:} Seleccionar Trimestral cuando no figura entre las tarjetas.\par\smallskip\textbf{Resultado:} No acepta una frecuencia arbitraria. \\
\addlinespace[3pt]
Agregar & Inserta la estrategia seleccionada en el borrador. & \textbf{Ejemplo:} Reincorporar Anual.\par\smallskip\textbf{Resultado:} La tarjeta aparece y se conserva al Guardar catálogo. \\
\addlinespace[3pt]
Guardar catálogo & Persiste también las opciones temporales del formulario. & \textbf{Ejemplo:} Guardar Mensual y Trimestral habilitadas.\par\smallskip\textbf{Resultado:} Las nuevas necesidades ofrecen esa configuración. \\
\addlinespace[3pt]
\end{longtable}
\endgroup

\subsection{Explicar un datamart de líneas con análisis mensual}
Práctica propuesta para el entorno de capacitación; no se presenta como una nueva ejecución realizada durante la redacción.

\begin{enumerate}
\item Seleccione una propuesta histórica de líneas en el recorrido de formación.
\item Lea su granularidad y su periodicidad.
\item Explique qué se suma durante un mes y qué representa cada fila almacenada.
\item Repita la explicación para la otra fuente con su propio evento comercial.
\end{enumerate}

\textbf{Qué debe obtener.} Se distingue nivel de detalle, fecha de negocio y agrupación. La estudiante no afirma que habilitar Mensual ejecute un ETL mensual automático.

\textbf{Cómo comprobarlo.} Puede responder qué representa una fila y qué representa un punto mensual del gráfico sin confundirlos.

\Needspace{8\baselineskip}
\begingroup
\setbox0=\vbox{\hsize=\linewidth
{\normalfont\normalsize\sffamily\bfseries\normalsize\color{Teal}\hyphenpenalty=10000\exhyphenpenalty=10000 Si aparece un problema\par}

\par
\begingroup
\fontsize{10}{12.5}\selectfont
\setlength{\tabcolsep}{5pt}
\renewcommand{\arraystretch}{1.15}
\begin{tabular}{@{}>{\RaggedRight\arraybackslash}p{0.22\linewidth}>{\RaggedRight\arraybackslash}p{0.3\linewidth}>{\RaggedRight\arraybackslash}p{0.39\linewidth}@{}}
\toprule
\textbf{Qué observa} & \textbf{Causa posible} & \textbf{Qué hacer y comprobar} \\
\midrule
No aparece una periodicidad en el asistente. & Está retirada o deshabilitada en el catálogo de esa fuente. & Revise las tarjetas y guarde la configuración correspondiente; confirme que inicia una necesidad nueva de la misma fuente.\par\smallskip\textbf{Comprobación:} La opción habilitada se ofrece en el contexto correcto. \\
\bottomrule
\end{tabular}
\endgroup
}
\Needspace{\dimexpr\ht0+\dp0+48pt\relax}
\endgroup
\subsection{Si aparece un problema}
\begingroup
\fontsize{10}{12.5}\selectfont
\setlength{\tabcolsep}{5pt}
\renewcommand{\arraystretch}{1.15}
\begin{longtable}{@{}>{\RaggedRight\arraybackslash}p{0.22\linewidth}>{\RaggedRight\arraybackslash}p{0.3\linewidth}>{\RaggedRight\arraybackslash}p{0.39\linewidth}@{}}
\toprule
\textbf{Qué observa} & \textbf{Causa posible} & \textbf{Qué hacer y comprobar} \\
\midrule
\endfirsthead
\toprule
\textbf{Qué observa} & \textbf{Causa posible} & \textbf{Qué hacer y comprobar} \\
\midrule
\endhead
\midrule
\multicolumn{3}{r}{\footnotesize Continúa en la página siguiente}\\
\endfoot
\bottomrule
\endlastfoot
No aparece una periodicidad en el asistente. & Está retirada o deshabilitada en el catálogo de esa fuente. & Revise las tarjetas y guarde la configuración correspondiente; confirme que inicia una necesidad nueva de la misma fuente.\par\smallskip\textbf{Comprobación:} La opción habilitada se ofrece en el contexto correcto. \\
\addlinespace[3pt]
Se desea una frecuencia que no aparece. & La estrategia no está implementada. & Utilice una opción admitida si responde a la necesidad o registre una ampliación de desarrollo. No renombre Mensual para simular Bimestral.\par\smallskip\textbf{Comprobación:} El título temporal coincide con la agrupación realmente ejecutada. \\
\addlinespace[3pt]
No hay fecha de negocio verificable. & La estructura no respalda la agrupación solicitada o la fecha pertenece a otro evento. & Revise metadatos y delimite la necesidad. La etiqueta temporal no reemplaza una fecha ausente.\par\smallskip\textbf{Comprobación:} La propuesta usa una fecha respaldada o informa la limitación. \\
\addlinespace[3pt]
Se supone que cambiar el catálogo cambió el panel histórico. & Se confundió la configuración futura con el contrato de una ejecución ya materializada. & Consulte la periodicidad y la propuesta de esa ejecución. Cree o revise un nuevo diseño solo cuando se requiera y se autorice.\par\smallskip\textbf{Comprobación:} La procedencia histórica permanece identificada. \\
\addlinespace[3pt]
\end{longtable}
\endgroup

\subsection{Para explicar durante la sustentación}
\paragraph{¿Mensual es la granularidad?} No necesariamente. En los casos de líneas, el hecho conserva líneas y el análisis agrupa sus importes por mes.

\paragraph{¿Por qué no usar siempre la fecha de modificación del registro?} Porque suele indicar una acción administrativa. La fecha del evento comercial debe corresponder al pedido, factura u otro evento solicitado.

\paragraph{¿Se puede refrescar automáticamente al marcar Semanal?} No. La opción agrupa resultados por semana; no programa una nueva carga.

\endgroup
\pantalla{recursos/pantallas/18_catalogo_periodicidades.jpg}{Habilitar periodicidades y guardar el catálogo de la fuente}{Periodicidades habilitadas. Agrupar por mes no transforma una línea física en una fila mensual. Captura de referencia del 18 de noviembre de 2026; revise las instrucciones actualizadas de esta unidad.}

\IfFileExists{capturas/operacion_14.tex}{\input{capturas/operacion_14.tex}}{}
\section{Seis orientaciones comerciales y cómo redactarlas}
Estos textos son ejemplos de capacitación. No se aplicaron a ningún catálogo durante la redacción; revíselos según el respaldo de la fuente antes de incorporarlos.

\subsection{Evolución de ventas en el tiempo}
Orientación estándar, redacción de capacitación; revisar antes de sustituir el texto local.

\textbf{Explicación para el analista.} Permite comparar el importe registrado y las unidades entre periodos de un mismo evento comercial.

\textbf{Instrucción funcional para la IA.} Comparar sumas por periodos usando la fecha de negocio verificable. Distinguir pedido de factura; no afirmar una tendencia causal ni comparar periodos incompletos como equivalentes.

\subsection{Productos con mayor desempeño}
Orientación estándar, ejemplo de instrucción más concreta.

\textbf{Explicación para el analista.} Permite identificar productos con mayor importe o unidades dentro de la selección analizada.

\textbf{Instrucción funcional para la IA.} Ordenar productos por sumas de importe o unidades respaldadas. Mostrar el alcance de la selección; el ranking no demuestra preferencias, demanda futura ni causalidad.

\subsection{Desempeño por cliente}
Orientación estándar.

\textbf{Explicación para el analista.} Permite comparar ventas atribuidas al cliente mediante relaciones de negocio comprobadas.

\textbf{Instrucción funcional para la IA.} Usar la identidad descriptiva del cliente obtenida mediante rutas verificadas. No sustituir cliente por vendedor, empleado que editó el registro o contacto sin respaldo.

\subsection{Desempeño por territorio}
Orientación estándar; el territorio efectivo depende de la fuente y la propuesta.

\textbf{Explicación para el analista.} Permite comparar resultados según la ubicación o el territorio comercial respaldado por la fuente.

\textbf{Instrucción funcional para la IA.} Declarar el rol territorial usado y conservar la ruta revisada. Una dirección actual de cliente no acredita lugar de entrega realizado ni ubicación histórica del evento.

\subsection{Costos y rentabilidad de las ventas}
Orientación personalizada histórica de AdventureWorks; no se garantiza aplicabilidad a otra fuente.

\textbf{Explicación para el analista.} Permite revisar la diferencia entre ventas y costos cuando los componentes tienen respaldo y alcance explícitos.

\textbf{Instrucción funcional para la IA.} Usar únicamente costos e importes verificables y declarar su fecha y significado. No presentar un costo actual de producto como costo histórico de cada venta ni inventar margen cuando falten componentes.

\subsection{Ventas y costos por unidad}
Orientación personalizada histórica; precisa revisión de factibilidad y denominadores.

\textbf{Explicación para el analista.} Permite comparar importes por unidad conservando un denominador de unidades y el alcance de los costos disponibles.

\textbf{Instrucción funcional para la IA.} Calcular razones entre sumas compatibles de importe y unidades, con denominador verificable y tratamiento de cero. Distinguir promedio por unidad, por línea y por documento; declarar costos actuales cuando corresponda.

\chapter{Del objetivo comercial a una propuesta supervisada}
\section{Antes de empezar: piense qué decisión necesita apoyar}
\label{unidad:BI01}
\textbf{Para quién:} Analista de BI, Gerente comercial\par
\textbf{Al terminar podrá:} Distinguir la pregunta comercial, el evento y el nivel de detalle antes de pedir un datamart.

Un datamart reúne datos organizados para un análisis concreto. En este proyecto el punto de partida es una necesidad de ventas escrita por una persona. La plataforma ayuda a relacionarla con la estructura disponible y conserva las decisiones que llevan hasta el tablero.

Antes de abrir el asistente, escriba qué quiere comparar, con qué indicador, por qué entidades y en qué período. Elija si observa pedidos o facturas. Un pedido representa una solicitud de venta; una factura representa otro evento. Sus importes no demuestran por sí solos cobro, entrega realizada o utilidad.

La interfaz ofrece un proceso común para fuentes distintas. La demostración comprobada utiliza dos bases SQL Server. Cambian sus tablas, claves y significados; no debe copiarse la elección técnica de una base a otra sólo porque los botones son iguales.

\subsection{Conceptos para decidir con criterio}
\paragraph{Hecho} Registro del evento que se analizará, por ejemplo una línea de pedido o una línea de factura.

\paragraph{Granularidad o grano} Qué representa exactamente una fila del hecho. Decir una línea de factura permite comprobar qué se conserva y evita llamar factura a cada línea.

\paragraph{Dimensión} Información que permite agrupar y describir el hecho: fecha, producto, cliente o territorio. Una clave identifica; un nombre ayuda a interpretar.

\paragraph{Medida} Valor asociado al hecho, como cantidad, importe de línea o identificador del documento que se contará de forma distinta.

\paragraph{Indicador} Cálculo con una finalidad comercial: total de unidades, facturas distintas o importe por factura. Debe tener fórmula, unidad, alcance y procedencia.

\paragraph{Agregación} Forma de resumir varias filas. Sumar, contar y promediar responden preguntas distintas.

\paragraph{Procedencia} Cadena que identifica fuente, instantánea de metadatos, propuesta, ejecución y relaciones empleadas.

\Needspace{8\baselineskip}
\begingroup
\setbox0=\vbox{\hsize=\linewidth
{\normalfont\normalsize\sffamily\bfseries\normalsize\color{Teal}\hyphenpenalty=10000\exhyphenpenalty=10000 Campos y controles\par}

\par
\begingroup
\fontsize{10}{12.5}\selectfont
\setlength{\tabcolsep}{5pt}
\renewcommand{\arraystretch}{1.15}
\begin{tabular}{@{}>{\RaggedRight\arraybackslash}p{0.27\linewidth}>{\RaggedRight\arraybackslash}p{0.66\linewidth}@{}}
\toprule
\textbf{Campo o control} & \textbf{Significado y criterio} \\
\midrule
Fuente en contexto & Base elegida para trabajar. Debe coincidir con la necesidad, los metadatos y el resultado consultado. \\
\bottomrule
\end{tabular}
\endgroup
}
\Needspace{\dimexpr\ht0+\dp0+48pt\relax}
\endgroup
\subsection{Campos y controles}
\begingroup
\fontsize{10}{12.5}\selectfont
\setlength{\tabcolsep}{5pt}
\renewcommand{\arraystretch}{1.15}
\begin{longtable}{@{}>{\RaggedRight\arraybackslash}p{0.27\linewidth}>{\RaggedRight\arraybackslash}p{0.66\linewidth}@{}}
\toprule
\textbf{Campo o control} & \textbf{Significado y criterio} \\
\midrule
\endfirsthead
\toprule
\textbf{Campo o control} & \textbf{Significado y criterio} \\
\midrule
\endhead
\midrule
\multicolumn{2}{r}{\footnotesize Continúa en la página siguiente}\\
\endfoot
\bottomrule
\endlastfoot
Fuente en contexto & Base elegida para trabajar. Debe coincidir con la necesidad, los metadatos y el resultado consultado. \\
\addlinespace[3pt]
Dominio Ventas & Perfil funcional implementado. No convierte automáticamente en ventas todas las tablas de la base. \\
\addlinespace[3pt]
\end{longtable}
\endgroup

\subsection{Paso a paso}
\begin{enumerate}
\item Escriba una decisión concreta, por ejemplo priorizar productos según unidades registradas cada mes.\par\nopagebreak[4]\smallskip{\small\textbf{Compruebe:} Una necesidad que se puede contrastar con datos, sin pedir todavía una conclusión causal.}
\item Precise el evento: pedido o factura, y la fila mínima que desea conservar.\par\nopagebreak[4]\smallskip{\small\textbf{Compruebe:} Un grano explícito, como una línea de pedido.}
\item Anote los indicadores y sus unidades.\par\nopagebreak[4]\smallskip{\small\textbf{Compruebe:} Por ejemplo importe, unidades, pedidos distintos e importe por pedido, sin confundirlos.}
\item Identifique la fuente y revise su esquema antes de generar la propuesta.\par\nopagebreak[4]\smallskip{\small\textbf{Compruebe:} Tablas y relaciones de esa base, sin suponer que tener un nombre parecido demuestra el significado.}
\end{enumerate}
\subsection{Qué hace la IA}
En esta preparación la persona define el propósito. La IA podrá sugerir una redacción y un diseño, pero todavía no existe evidencia de que una propuesta haya sido aprobada o cargada.

\subsection{Qué comprueba la plataforma}
Los controles posteriores comprobarán referencias, relaciones y recetas. La disponibilidad de una columna no demuestra por sí sola calidad, vigencia histórica o suficiencia de registros.

\subsection{Qué decide la persona}
El analista decide el evento y los significados; la gerencia aporta la decisión comercial que desea apoyar.

\paragraph{Resultado esperado} Una ficha de necesidad que se pueda leer y explicar antes de pulsar Generar.

\Needspace{8\baselineskip}
\begingroup
\setbox0=\vbox{\hsize=\linewidth
{\normalfont\normalsize\sffamily\bfseries\normalsize\color{Teal}\hyphenpenalty=10000\exhyphenpenalty=10000 Errores y recuperación\par}

\par
\begingroup
\fontsize{10}{12.5}\selectfont
\setlength{\tabcolsep}{5pt}
\renewcommand{\arraystretch}{1.15}
\begin{tabular}{@{}>{\RaggedRight\arraybackslash}p{0.27\linewidth}>{\RaggedRight\arraybackslash}p{0.66\linewidth}@{}}
\toprule
\textbf{Situación} & \textbf{Qué hacer} \\
\midrule
El objetivo dice analizar todas las ventas para mejorar el negocio. & Acótelo: indique evento, indicador, comparación y límite. Evite pedir que el sistema demuestre por adelantado un aumento de ventas. \\
\bottomrule
\end{tabular}
\endgroup
}
\Needspace{\dimexpr\ht0+\dp0+48pt\relax}
\endgroup
\subsection{Errores y recuperación}
\begingroup
\fontsize{10}{12.5}\selectfont
\setlength{\tabcolsep}{5pt}
\renewcommand{\arraystretch}{1.15}
\begin{longtable}{@{}>{\RaggedRight\arraybackslash}p{0.27\linewidth}>{\RaggedRight\arraybackslash}p{0.66\linewidth}@{}}
\toprule
\textbf{Situación} & \textbf{Qué hacer} \\
\midrule
\endfirsthead
\toprule
\textbf{Situación} & \textbf{Qué hacer} \\
\midrule
\endhead
\midrule
\multicolumn{2}{r}{\footnotesize Continúa en la página siguiente}\\
\endfoot
\bottomrule
\endlastfoot
El objetivo dice analizar todas las ventas para mejorar el negocio. & Acótelo: indique evento, indicador, comparación y límite. Evite pedir que el sistema demuestre por adelantado un aumento de ventas. \\
\addlinespace[3pt]
\end{longtable}
\endgroup

\subsection{Ejercicio}
Con un pedido que tiene tres líneas, indique cuántas filas de hecho y cuántos pedidos distintos habría.

\textbf{Respuesta orientativa.} Tres filas a grano de línea y un pedido distinto. Contar filas no da el número de pedidos.

\subsection{Compruebe lo que aprendió}
\paragraph{¿Un gráfico de importe prueba rentabilidad?} No. Harían falta costos comparables y una definición de margen, entre otros datos.

\paragraph{¿Una base nueva exige programar un tablero específico con su nombre?} El recorrido es común; primero se configuran y comprueban sus metadatos y decisiones. La compatibilidad real requiere los adaptadores y operaciones implementados.

\pantalla{recursos/pantallas/19_dominio_ventas.jpg}{Elegir el dominio de ventas}{El dominio habilitado es ventas. La pantalla no anuncia dominios que no tengan perfil y validadores implementados. Captura de referencia del 18 de noviembre de 2026; revise las instrucciones actualizadas de esta unidad.}

\IfFileExists{capturas/BI01.tex}{\input{capturas/BI01.tex}}{}
\section{Asistente, paso 1: escribir o elegir una necesidad}
\label{unidad:BI02}
\textbf{Para quién:} Analista de BI\par
\textbf{Al terminar podrá:} Completar el objetivo, las preguntas, la periodicidad y la autorización de metadatos sin perder el propósito comercial.

Entre a Asistente de datamart con la fuente correcta y seleccione el dominio de ventas. El recorrido tiene cinco etapas: Necesidad, Conceptos, Propuesta, Personalización y Revisión.

Puede comenzar con su propio texto o pedir sugerencias con la estructura de la fuente. Las sugerencias no sustituyen el objetivo automáticamente: primero se leen su explicación, evidencia y límites; después se pulsa Usar esta necesidad.

\subsection{Conceptos para decidir con criterio}
\paragraph{Objetivo del análisis} Describe qué necesita conocer la persona. Debe conservar el evento y las limitaciones que importan al negocio.

\paragraph{Preguntas de negocio} Orientaciones del catálogo de esta fuente. Seleccionarlas no fija por sí solo tablas, dimensiones ni fórmulas.

\paragraph{Periodicidad} Forma temporal de resumir o presentar el análisis. Elegir Mensual no significa programar una recarga automática cada mes ni reducir el hecho necesariamente a una fila mensual.

\Needspace{8\baselineskip}
\begingroup
\setbox0=\vbox{\hsize=\linewidth
{\normalfont\normalsize\sffamily\bfseries\normalsize\color{Teal}\hyphenpenalty=10000\exhyphenpenalty=10000 Campos y controles\par}

\par
\begingroup
\fontsize{10}{12.5}\selectfont
\setlength{\tabcolsep}{5pt}
\renewcommand{\arraystretch}{1.15}
\begin{tabular}{@{}>{\RaggedRight\arraybackslash}p{0.27\linewidth}>{\RaggedRight\arraybackslash}p{0.66\linewidth}@{}}
\toprule
\textbf{Campo o control} & \textbf{Significado y criterio} \\
\midrule
Objetivo del análisis & Texto de 20 a 2.000 caracteres. Escriba qué comparar, indicadores, grano y límites. No escriba SQL. \\
\bottomrule
\end{tabular}
\endgroup
}
\Needspace{\dimexpr\ht0+\dp0+48pt\relax}
\endgroup
\subsection{Campos y controles}
\begingroup
\fontsize{10}{12.5}\selectfont
\setlength{\tabcolsep}{5pt}
\renewcommand{\arraystretch}{1.15}
\begin{longtable}{@{}>{\RaggedRight\arraybackslash}p{0.27\linewidth}>{\RaggedRight\arraybackslash}p{0.66\linewidth}@{}}
\toprule
\textbf{Campo o control} & \textbf{Significado y criterio} \\
\midrule
\endfirsthead
\toprule
\textbf{Campo o control} & \textbf{Significado y criterio} \\
\midrule
\endhead
\midrule
\multicolumn{2}{r}{\footnotesize Continúa en la página siguiente}\\
\endfoot
\bottomrule
\endlastfoot
Objetivo del análisis & Texto de 20 a 2.000 caracteres. Escriba qué comparar, indicadores, grano y límites. No escriba SQL. \\
\addlinespace[3pt]
Destino & Identifica el proveedor/configuración que recibirá la necesidad y los metadatos autorizados. \\
\addlinespace[3pt]
Autorizo enviar estos metadatos y la necesidad al proveedor indicado para esta revisión & Habilita la asistencia para ese destino. Si cambia el proveedor se debe revisar de nuevo el consentimiento. \\
\addlinespace[3pt]
Preguntas de negocio disponibles & Casillas con etiqueta y explicación. Una pregunta no disponible no se convierte en viable por forzar su selección. \\
\addlinespace[3pt]
Periodicidad & Seleccione una opción disponible del catálogo, por ejemplo Mensual. \\
\addlinespace[3pt]
Sugerir necesidades con esta fuente & Solicita posibles objetivos desde los metadatos; no necesita que ya haya una propuesta. \\
\addlinespace[3pt]
Ayúdame a formular la necesidad & Propone una redacción a partir del objetivo y preguntas existentes. Requiere un objetivo suficiente y preguntas seleccionadas. \\
\addlinespace[3pt]
Usar esta redacción / Mantener mi redacción & Permite aceptar la propuesta textual o conservar el texto propio. \\
\addlinespace[3pt]
\end{longtable}
\endgroup

\subsection{Paso a paso}
\begin{enumerate}
\item Compruebe nombre de fuente e instantánea visibles antes de redactar.\par\nopagebreak[4]\smallskip{\small\textbf{Compruebe:} El objetivo quedará asociado a esa estructura, no a otra base.}
\item Seleccione preguntas que respondan a la decisión, sin marcar todo por costumbre.\par\nopagebreak[4]\smallskip{\small\textbf{Compruebe:} Alcance entendible y coherente con el catálogo de esta fuente.}
\item Elija periodicidad y autorice el destino que aparece en pantalla.\par\nopagebreak[4]\smallskip{\small\textbf{Compruebe:} Asistencia habilitada para el proveedor indicado.}
\item Escriba la necesidad o pulse Sugerir necesidades con esta fuente.\par\nopagebreak[4]\smallskip{\small\textbf{Compruebe:} Texto propio o varias opciones con fundamento y limitaciones.}
\item Si usa Ayúdame a formular la necesidad, compare la propuesta con el original.\par\nopagebreak[4]\smallskip{\small\textbf{Compruebe:} La redacción conserva el evento, indicadores y límites; no añade pronósticos ni otra fuente.}
\item Use sólo el texto que represente su necesidad y compruebe su viabilidad.\par\nopagebreak[4]\smallskip{\small\textbf{Compruebe:} La plataforma inicia una revisión previa; todavía no se aprueba ni se ejecuta ETL.}
\end{enumerate}
\subsection{Qué hace la IA}
Sugiere objetivos y mejora redacciones con el contexto estructural autorizado. Su respuesta pasa controles de formato y de referencias antes de ofrecerse como utilizable.

\subsection{Qué comprueba la plataforma}
Existencia de instantánea, catálogo y proveedor utilizables; longitud y selección requeridas; destino del consentimiento; referencias y semántica admitida en la revisión de necesidad. Cambiar objetivo, preguntas o periodicidad invalida la revisión anterior.

\subsection{Qué decide la persona}
El analista elige si usa una sugerencia y confirma que el texto representa la necesidad. Autorizar metadatos no equivale a aprobar el diseño ni la carga.

\paragraph{Resultado esperado} Una necesidad vigente, entendida y lista para la comprobación previa.

\Needspace{8\baselineskip}
\begingroup
\setbox0=\vbox{\hsize=\linewidth
{\normalfont\normalsize\sffamily\bfseries\normalsize\color{Teal}\hyphenpenalty=10000\exhyphenpenalty=10000 Errores y recuperación\par}

\par
\begingroup
\fontsize{10}{12.5}\selectfont
\setlength{\tabcolsep}{5pt}
\renewcommand{\arraystretch}{1.15}
\begin{tabular}{@{}>{\RaggedRight\arraybackslash}p{0.27\linewidth}>{\RaggedRight\arraybackslash}p{0.66\linewidth}@{}}
\toprule
\textbf{Situación} & \textbf{Qué hacer} \\
\midrule
No se habilita el botón de asistencia. & Revise permiso de generación, proveedor activo y probado, autorización, instantánea y los campos requeridos. \\
\bottomrule
\end{tabular}
\endgroup
}
\Needspace{\dimexpr\ht0+\dp0+48pt\relax}
\endgroup
\subsection{Errores y recuperación}
\begingroup
\fontsize{10}{12.5}\selectfont
\setlength{\tabcolsep}{5pt}
\renewcommand{\arraystretch}{1.15}
\begin{longtable}{@{}>{\RaggedRight\arraybackslash}p{0.27\linewidth}>{\RaggedRight\arraybackslash}p{0.66\linewidth}@{}}
\toprule
\textbf{Situación} & \textbf{Qué hacer} \\
\midrule
\endfirsthead
\toprule
\textbf{Situación} & \textbf{Qué hacer} \\
\midrule
\endhead
\midrule
\multicolumn{2}{r}{\footnotesize Continúa en la página siguiente}\\
\endfoot
\bottomrule
\endlastfoot
No se habilita el botón de asistencia. & Revise permiso de generación, proveedor activo y probado, autorización, instantánea y los campos requeridos. \\
\addlinespace[3pt]
La IA cambió el evento o añadió costos no solicitados. & Mantenga su texto o corríjalo antes de continuar; no acepte una ampliación sólo por estar bien redactada. \\
\addlinespace[3pt]
No aparecen sugerencias utilizables. & Lea el aviso, redacte un alcance más sencillo y compruebe la viabilidad; no se inventan datos faltantes. \\
\addlinespace[3pt]
\end{longtable}
\endgroup

\subsection{Ejercicio}
Redacte dos versiones del mismo objetivo: una vaga y otra que precise evento, grano, indicador y comparación.

\textbf{Respuesta orientativa.} La segunda debe permitir reconocer qué salida se revisará; por ejemplo unidades por producto y mes conservando cada línea de pedido.

\subsection{Compruebe lo que aprendió}
\paragraph{¿Usar esta necesidad ejecuta una carga?} No. Sólo coloca el texto elegido; se revisan preguntas y viabilidad antes de generar.

\paragraph{¿Mensual programa el ETL?} No. Es periodicidad analítica; el inicio de la carga sigue siendo una acción explícita.

\pantalla{recursos/pantallas/20_necesidad_y_autorizacion.jpg}{Escribir el objetivo del análisis}{Objetivo del análisis en un formulario nuevo. Describa la decisión, los indicadores, el detalle y las comparaciones; no escriba SQL. Captura de referencia del 18 de noviembre de 2026; revise las instrucciones actualizadas de esta unidad.}
\pantalla{recursos/pantallas/20b_consentimiento_IA.jpg}{Autorizar la revisión con IA}{Destino y alcance del envío de metadatos, ayudas de redacción y preguntas del catálogo. El consentimiento está desmarcado; no se envió una nueva solicitud. Captura de referencia del 18 de noviembre de 2026; revise las instrucciones actualizadas de esta unidad.}

\IfFileExists{capturas/BI02.tex}{\input{capturas/BI02.tex}}{}
\section{Comprender la viabilidad antes de generar}
\label{unidad:BI03}
\textbf{Para quién:} Analista de BI\par
\textbf{Al terminar podrá:} Interpretar directo, derivable, ambiguo y no disponible sin confundir la estructura con datos ya comprobados.

La revisión actual combina un análisis de la IA con reglas de la plataforma. Su objetivo es relacionar la necesidad con nombres, tipos, claves y relaciones de la instantánea. No consulta todas las filas del negocio para demostrar resultados.

Una derivación inequívoca puede resolverse al construir la propuesta mediante una ruta declarada o una fórmula controlada. Una interpretación candidata de IA sigue pendiente de validación: su frase no se ejecuta como fórmula libre.

\subsection{Del respaldo estructural a una respuesta analítica útil}
Derivable significa que una relación declarada o una receta controlada puede resolver el requisito a partir de referencias existentes. No significa que deba escribir SQL, ni que la IA tenga permiso para inventar una columna o ejecutar una fórmula libre. Lea las entradas, el grano y la ruta de relación.

La propuesta nueva debe conservar todos los requisitos revisados de la necesidad. No basta con encontrar una columna Quantity u OrderQty si se pidió analizar unidades y después el panel no ofrece ese indicador. Medida física es el dato de la línea; indicador es la operación que el usuario podrá seleccionar y explicar.

En los contratos nuevos, la plataforma sólo completa un indicador requerido cuando existe una receta segura y única a partir de las medidas admitidas. Su procedencia queda diferenciada de la sugerencia del modelo. Si hay varias interpretaciones posibles, corresponde revisar el caso; no elegir una en silencio. Si la persona excluye después un indicador requerido, se conserva su decisión y se muestra cobertura parcial.

\begin{enumerate}
\item Identifique las palabras de la necesidad que definen resultados obligatorios: importe, unidades, pedidos o facturas, por ejemplo.
\item En la propuesta, contraste cada requisito con su indicador, fórmula y unidad.
\item Compruebe que las unidades no se confundan con número de líneas y que los pedidos o facturas utilicen el conteo documental adecuado.
\item Si un requisito queda parcial, aclárelo o acepte explícitamente un alcance menor defendible. No apruebe sólo porque no hay errores de JSON.
\end{enumerate}

\subsection{Conceptos para decidir con criterio}
\paragraph{Directo} Existe una referencia estructural adecuada comprobada para el requisito.

\paragraph{Derivable por relación} El requisito puede obtenerse recorriendo vínculos verificados sin cambiar indebidamente el grano.

\paragraph{Derivable por fórmula} Puede calcularse con entradas compatibles y una operación permitida, por ejemplo precio \(\times\) tasa de descuento \(\times\) cantidad.

\paragraph{Ambiguo} Existen opciones o significados que requieren una definición del analista; una relación válida no basta para decidir su rol comercial.

\paragraph{No disponible} La revisión no encuentra respaldo suficiente para esa capacidad dentro del alcance actual.

\paragraph{Limitación aceptada} Constancia de que se entiende una ausencia; no crea una fuente ni demuestra que el requisito se cumpla.

\Needspace{8\baselineskip}
\begingroup
\setbox0=\vbox{\hsize=\linewidth
{\normalfont\normalsize\sffamily\bfseries\normalsize\color{Teal}\hyphenpenalty=10000\exhyphenpenalty=10000 Campos y controles\par}

\par
\begingroup
\fontsize{10}{12.5}\selectfont
\setlength{\tabcolsep}{5pt}
\renewcommand{\arraystretch}{1.15}
\begin{tabular}{@{}>{\RaggedRight\arraybackslash}p{0.27\linewidth}>{\RaggedRight\arraybackslash}p{0.66\linewidth}@{}}
\toprule
\textbf{Campo o control} & \textbf{Significado y criterio} \\
\midrule
Viabilidad contra los metadatos & Resumen de requisitos y evidencia. Lea cada tarjeta además del total de estados. \\
\bottomrule
\end{tabular}
\endgroup
}
\Needspace{\dimexpr\ht0+\dp0+48pt\relax}
\endgroup
\subsection{Campos y controles}
\begingroup
\fontsize{10}{12.5}\selectfont
\setlength{\tabcolsep}{5pt}
\renewcommand{\arraystretch}{1.15}
\begin{longtable}{@{}>{\RaggedRight\arraybackslash}p{0.27\linewidth}>{\RaggedRight\arraybackslash}p{0.66\linewidth}@{}}
\toprule
\textbf{Campo o control} & \textbf{Significado y criterio} \\
\midrule
\endfirsthead
\toprule
\textbf{Campo o control} & \textbf{Significado y criterio} \\
\midrule
\endhead
\midrule
\multicolumn{2}{r}{\footnotesize Continúa en la página siguiente}\\
\endfoot
\bottomrule
\endlastfoot
Viabilidad contra los metadatos & Resumen de requisitos y evidencia. Lea cada tarjeta además del total de estados. \\
\addlinespace[3pt]
Ver evidencia técnica & Muestra las referencias/rutas que sustentan el requisito. \\
\addlinespace[3pt]
Fórmula controlada & Receta propuesta por las reglas con entradas identificadas. \\
\addlinespace[3pt]
Interpretación propuesta por IA no ejecutable & Descripción que todavía debe convertirse y validarse en un contrato permitido. \\
\addlinespace[3pt]
Comprendo esta limitación y acepto continuar sin que el sistema invente una solución & Aceptación individual de requisitos ambiguos o ausentes; no debe marcarse sin leer. \\
\addlinespace[3pt]
Generar conceptos y propuesta & Usa la revisión vigente y las limitaciones expresamente aceptadas para generar una propuesta. \\
\addlinespace[3pt]
\end{longtable}
\endgroup

\subsection{Paso a paso}
\begin{enumerate}
\item Lea el resumen de viabilidad y abra la evidencia de cada requisito importante.\par\nopagebreak[4]\smallskip{\small\textbf{Compruebe:} Identifica qué está respaldado y qué requiere una decisión.}
\item En las derivaciones, compruebe si se trata de relación o fórmula y qué entradas usa.\par\nopagebreak[4]\smallskip{\small\textbf{Compruebe:} No se pide completar SQL ni buscar una columna sólo para quitar el aviso.}
\item Si un requisito imprescindible es ambiguo o ausente, vuelva a la necesidad y aclárelo.\par\nopagebreak[4]\smallskip{\small\textbf{Compruebe:} El objetivo no promete una salida imposible.}
\item Acepte sólo limitaciones que permitan todavía responder la pregunta útil del caso.\par\nopagebreak[4]\smallskip{\small\textbf{Compruebe:} Queda un alcance válido y las ausencias no se ocultan.}
\item Genere y después revise De la necesidad al datamart.\par\nopagebreak[4]\smallskip{\small\textbf{Compruebe:} Cada requisito debe aparecer cubierto, resuelto, limitado o pendiente, sin descarte silencioso.}
\end{enumerate}
\subsection{Qué hace la IA}
Examina la necesidad dentro del contexto de la fuente y aporta requisitos y explicaciones. La respuesta se combina con el análisis estructural y los controles de referencias.

\subsection{Qué comprueba la plataforma}
Referencias, relaciones, semántica admitida, alcance del evento y huella de entrada. La aprobación posterior vuelve a comprobar el contrato; aceptar una limitación no omite los controles bloqueantes.

\subsection{Qué decide la persona}
El analista aclara ambigüedades, reduce alcance cuando corresponde y decide qué limitaciones son aceptables.

\paragraph{Resultado esperado} Una revisión vigente y un alcance defendible antes del diseño.

\Needspace{8\baselineskip}
\begingroup
\setbox0=\vbox{\hsize=\linewidth
{\normalfont\normalsize\sffamily\bfseries\normalsize\color{Teal}\hyphenpenalty=10000\exhyphenpenalty=10000 Errores y recuperación\par}

\par
\begingroup
\fontsize{10}{12.5}\selectfont
\setlength{\tabcolsep}{5pt}
\renewcommand{\arraystretch}{1.15}
\begin{tabular}{@{}>{\RaggedRight\arraybackslash}p{0.27\linewidth}>{\RaggedRight\arraybackslash}p{0.66\linewidth}@{}}
\toprule
\textbf{Situación} & \textbf{Qué hacer} \\
\midrule
La necesidad tiene respaldo estructural y luego falla la propuesta. & Son controles distintos. La viabilidad no garantiza que el proveedor entregue un contrato válido; revise el error y reintente o ajuste el diseño. \\
\bottomrule
\end{tabular}
\endgroup
}
\Needspace{\dimexpr\ht0+\dp0+48pt\relax}
\endgroup
\subsection{Errores y recuperación}
\begingroup
\fontsize{10}{12.5}\selectfont
\setlength{\tabcolsep}{5pt}
\renewcommand{\arraystretch}{1.15}
\begin{longtable}{@{}>{\RaggedRight\arraybackslash}p{0.27\linewidth}>{\RaggedRight\arraybackslash}p{0.66\linewidth}@{}}
\toprule
\textbf{Situación} & \textbf{Qué hacer} \\
\midrule
\endfirsthead
\toprule
\textbf{Situación} & \textbf{Qué hacer} \\
\midrule
\endhead
\midrule
\multicolumn{2}{r}{\footnotesize Continúa en la página siguiente}\\
\endfoot
\bottomrule
\endlastfoot
La necesidad tiene respaldo estructural y luego falla la propuesta. & Son controles distintos. La viabilidad no garantiza que el proveedor entregue un contrato válido; revise el error y reintente o ajuste el diseño. \\
\addlinespace[3pt]
El sistema pide validar otra vez. & Compruebe si cambió el objetivo, periodicidad, preguntas, fuente o destino del proveedor; una aceptación anterior no debe reutilizarse a ciegas. \\
\addlinespace[3pt]
Una casilla de limitación parece desbloquear todo. & El botón puede permitir continuar con un alcance parcial, pero no fabrica la capacidad ausente ni habilita SQL libre. \\
\addlinespace[3pt]
\end{longtable}
\endgroup

\subsection{Ejercicio}
Explique qué haría si se pide margen y sólo se encuentra importe de venta.

\textbf{Respuesta orientativa.} No afirmar margen; comprobar si hay costo comparable verificable, aclarar el alcance o aceptar que esa salida no estará disponible.

\subsection{Compruebe lo que aprendió}
\paragraph{¿Derivable significa que la alumna debe escribir una fórmula SQL?} No. Las rutas/fórmulas inequívocas se resuelven mediante operaciones controladas. Una ambigüedad se decide con evidencia.

\paragraph{¿Viabilidad demuestra que todas las filas están bien?} No. Acredita respaldo estructural; la carga y conciliación revisarán datos concretos.


\IfFileExists{capturas/BI03.tex}{\input{capturas/BI03.tex}}{}
\section{Asistente, paso 2: revisar los conceptos}
\label{unidad:BI04}
\textbf{Para quién:} Analista de BI\par
\textbf{Al terminar podrá:} Seleccionar conceptos a partir de evidencia y usar el copiloto como apoyo a una decisión explícita.

La pantalla Conceptos encontrados traduce la necesidad en conceptos vinculados a referencias técnicas. La confianza alta no reemplaza revisar qué tabla representa el evento. Los conceptos con evidencia débil pueden quedar excluidos preventivamente.

El copiloto de conceptos es distinto del chat del tablero: aquí explica la conveniencia de incluir una entidad en el diseño, no analiza todavía los importes cargados.

\subsection{Conceptos para decidir con criterio}
\paragraph{Concepto de negocio} Idea que la persona reconoce, como producto, cliente o territorio, asociada a columnas o tablas verificables.

\paragraph{Confianza} Orientación de revisión; no es una medición de precisión universal del LLM.

\paragraph{Inclusión excepcional} Decisión consciente de incluir un concepto cuya evidencia exige atención. Después debe superar la nueva validación.

\Needspace{8\baselineskip}
\begingroup
\setbox0=\vbox{\hsize=\linewidth
{\normalfont\normalsize\sffamily\bfseries\normalsize\color{Teal}\hyphenpenalty=10000\exhyphenpenalty=10000 Campos y controles\par}

\par
\begingroup
\fontsize{10}{12.5}\selectfont
\setlength{\tabcolsep}{5pt}
\renewcommand{\arraystretch}{1.15}
\begin{tabular}{@{}>{\RaggedRight\arraybackslash}p{0.27\linewidth}>{\RaggedRight\arraybackslash}p{0.66\linewidth}@{}}
\toprule
\textbf{Campo o control} & \textbf{Significado y criterio} \\
\midrule
Incluir & Casilla que conserva o excluye el concepto de la siguiente propuesta. \\
\bottomrule
\end{tabular}
\endgroup
}
\Needspace{\dimexpr\ht0+\dp0+48pt\relax}
\endgroup
\subsection{Campos y controles}
\begingroup
\fontsize{10}{12.5}\selectfont
\setlength{\tabcolsep}{5pt}
\renewcommand{\arraystretch}{1.15}
\begin{longtable}{@{}>{\RaggedRight\arraybackslash}p{0.27\linewidth}>{\RaggedRight\arraybackslash}p{0.66\linewidth}@{}}
\toprule
\textbf{Campo o control} & \textbf{Significado y criterio} \\
\midrule
\endfirsthead
\toprule
\textbf{Campo o control} & \textbf{Significado y criterio} \\
\midrule
\endhead
\midrule
\multicolumn{2}{r}{\footnotesize Continúa en la página siguiente}\\
\endfoot
\bottomrule
\endlastfoot
Incluir & Casilla que conserva o excluye el concepto de la siguiente propuesta. \\
\addlinespace[3pt]
Ver evidencia y cómo resolver & Expone origen técnico, comprobaciones y orientación. \\
\addlinespace[3pt]
Confirmar inclusión excepcional & Se usa si el concepto con evidencia débil se mantiene incluido; no inventa una relación. \\
\addlinespace[3pt]
Consultar al copiloto & Abre conversación del concepto con objetivo, referencias y expediente. \\
\addlinespace[3pt]
Pregunta para el copiloto & Entre 10 y 500 caracteres. Pregunte qué significa la relación o qué riesgo tiene incluirla. \\
\addlinespace[3pt]
Aplicar recomendación & Cambia la selección según la recomendación de incluir/excluir; exige una acción de la persona. \\
\addlinespace[3pt]
Generar propuesta BI & Continúa con la selección; si cambia el alcance de conceptos puede generar otra versión con IA. \\
\addlinespace[3pt]
\end{longtable}
\endgroup

\subsection{Paso a paso}
\begin{enumerate}
\item Lea nombre, descripción, confianza y estado de cada tarjeta.\par\nopagebreak[4]\smallskip{\small\textbf{Compruebe:} Identifica por qué una entidad está incluida o excluida.}
\item Abra la evidencia de los conceptos con revisión requerida.\par\nopagebreak[4]\smallskip{\small\textbf{Compruebe:} Relaciona el nombre comercial con las referencias reales.}
\item Pregunte al copiloto por riesgos, claves y consecuencias si aún no puede justificar la selección.\par\nopagebreak[4]\smallskip{\small\textbf{Compruebe:} Obtiene una explicación contextual con evidencia citada.}
\item Excluya lo que no responda al objetivo o confirme expresamente una inclusión excepcional sustentada.\par\nopagebreak[4]\smallskip{\small\textbf{Compruebe:} No quedan decisiones pendientes de inclusión.}
\item Pulse Generar propuesta BI.\par\nopagebreak[4]\smallskip{\small\textbf{Compruebe:} Se obtiene el contrato para revisar; se conserva trazabilidad si se generó otra versión.}
\end{enumerate}
\subsection{Qué hace la IA}
Explica riesgo, consecuencias de incluir/excluir y recomendación. La conversación no modifica selecciones por sí sola.

\subsection{Qué comprueba la plataforma}
Las referencias y el contrato resultante vuelven a validarse. No se permite continuar sin conceptos ni con decisiones de inclusión pendientes.

\subsection{Qué decide la persona}
El analista decide la selección y aplica o rechaza la recomendación.

\paragraph{Resultado esperado} Un conjunto de conceptos que la alumna pueda justificar en lenguaje de negocio.

\Needspace{8\baselineskip}
\begingroup
\setbox0=\vbox{\hsize=\linewidth
{\normalfont\normalsize\sffamily\bfseries\normalsize\color{Teal}\hyphenpenalty=10000\exhyphenpenalty=10000 Errores y recuperación\par}

\par
\begingroup
\fontsize{10}{12.5}\selectfont
\setlength{\tabcolsep}{5pt}
\renewcommand{\arraystretch}{1.15}
\begin{tabular}{@{}>{\RaggedRight\arraybackslash}p{0.27\linewidth}>{\RaggedRight\arraybackslash}p{0.66\linewidth}@{}}
\toprule
\textbf{Situación} & \textbf{Qué hacer} \\
\midrule
La IA recomienda incluir pero la evidencia no demuestra la relación. & No aplique la recomendación sólo por su tono. Excluya o defina el negocio y revise una nueva propuesta. \\
\bottomrule
\end{tabular}
\endgroup
}
\Needspace{\dimexpr\ht0+\dp0+48pt\relax}
\endgroup
\subsection{Errores y recuperación}
\begingroup
\fontsize{10}{12.5}\selectfont
\setlength{\tabcolsep}{5pt}
\renewcommand{\arraystretch}{1.15}
\begin{longtable}{@{}>{\RaggedRight\arraybackslash}p{0.27\linewidth}>{\RaggedRight\arraybackslash}p{0.66\linewidth}@{}}
\toprule
\textbf{Situación} & \textbf{Qué hacer} \\
\midrule
\endfirsthead
\toprule
\textbf{Situación} & \textbf{Qué hacer} \\
\midrule
\endhead
\midrule
\multicolumn{2}{r}{\footnotesize Continúa en la página siguiente}\\
\endfoot
\bottomrule
\endlastfoot
La IA recomienda incluir pero la evidencia no demuestra la relación. & No aplique la recomendación sólo por su tono. Excluya o defina el negocio y revise una nueva propuesta. \\
\addlinespace[3pt]
Quedan conceptos por confirmar. & Confirme únicamente los que entiende y puede sostener, o desmárquelos para continuar con un alcance menor. \\
\addlinespace[3pt]
\end{longtable}
\endgroup

\subsection{Ejercicio}
Pregunte: ¿qué riesgo tendría incluir este concepto en el datamart? Registre una evidencia y una decisión.

\textbf{Respuesta orientativa.} La decisión debe relacionar el concepto con la necesidad; la recomendación de IA no equivale a aprobación.

\subsection{Compruebe lo que aprendió}
\paragraph{¿Consultar al copiloto carga datos?} No. Esta conversación usa el expediente y metadatos del diseño; no ejecuta el ETL.

\paragraph{¿Excluir un concepto elimina información del origen?} No. Lo retira del alcance de la propuesta.

\pantalla{recursos/pantallas/22_conceptos.jpg}{Conceptos y niveles de evidencia}{Conceptos revisados de la propuesta 129. La confianza semántica y la existencia técnica son comprobaciones diferentes. Captura de referencia del 18 de noviembre de 2026; revise las instrucciones actualizadas de esta unidad.}

\IfFileExists{capturas/BI04.tex}{\input{capturas/BI04.tex}}{}
\section{Asistente, paso 3: leer el diseño propuesto}
\label{unidad:BI05}
\textbf{Para quién:} Analista de BI\par
\textbf{Al terminar podrá:} Explicar grano, hecho, medidas, dimensiones, indicadores y cobertura antes de personalizar o aprobar.

La propuesta es un contrato de diseño que responde a la necesidad. Empiece por el grano y el evento; después revise medidas y relaciones.

Una propuesta puede tener más indicadores que los cuatro de una práctica. La cantidad depende de la generación. Revise cuáles responden al objetivo y no memorice una cantidad universal de tarjetas.

\subsection{Conceptos para decidir con criterio}
\paragraph{Clave del grano} Columna o combinación que identifica la fila mínima. En una fuente puede ser una clave simple y en otra una combinación.

\paragraph{Proveniencia de una medida} Referencias y fórmula que explican de dónde sale el valor.

\paragraph{Identidad descriptiva} Nombre reconocible de la entidad y su tipo cuando corresponde. No basta mostrar sólo una clave técnica.

\paragraph{Cobertura} Correspondencia entre cada requisito del objetivo y las salidas que realmente incluye el contrato.

\Needspace{8\baselineskip}
\begingroup
\setbox0=\vbox{\hsize=\linewidth
{\normalfont\normalsize\sffamily\bfseries\normalsize\color{Teal}\hyphenpenalty=10000\exhyphenpenalty=10000 Campos y controles\par}

\par
\begingroup
\fontsize{10}{12.5}\selectfont
\setlength{\tabcolsep}{5pt}
\renewcommand{\arraystretch}{1.15}
\begin{tabular}{@{}>{\RaggedRight\arraybackslash}p{0.27\linewidth}>{\RaggedRight\arraybackslash}p{0.66\linewidth}@{}}
\toprule
\textbf{Campo o control} & \textbf{Significado y criterio} \\
\midrule
Granularidad y Clave del grano & Qué conserva el hecho; no confundir periodicidad mensual con grano físico de línea. \\
\bottomrule
\end{tabular}
\endgroup
}
\Needspace{\dimexpr\ht0+\dp0+48pt\relax}
\endgroup
\subsection{Campos y controles}
\begingroup
\fontsize{10}{12.5}\selectfont
\setlength{\tabcolsep}{5pt}
\renewcommand{\arraystretch}{1.15}
\begin{longtable}{@{}>{\RaggedRight\arraybackslash}p{0.27\linewidth}>{\RaggedRight\arraybackslash}p{0.66\linewidth}@{}}
\toprule
\textbf{Campo o control} & \textbf{Significado y criterio} \\
\midrule
\endfirsthead
\toprule
\textbf{Campo o control} & \textbf{Significado y criterio} \\
\midrule
\endhead
\midrule
\multicolumn{2}{r}{\footnotesize Continúa en la página siguiente}\\
\endfoot
\bottomrule
\endlastfoot
Granularidad y Clave del grano & Qué conserva el hecho; no confundir periodicidad mensual con grano físico de línea. \\
\addlinespace[3pt]
Hecho y medidas & Para cada medida se muestran nombre, agregación, origen y fórmula. \\
\addlinespace[3pt]
Dimensiones y Columnas de origen & Atributos con los que se describirán y agruparán los hechos. \\
\addlinespace[3pt]
Campos que creará el ETL & Nombres/tipos descriptivos derivados por rutas verificadas; no son columnas que ya existan necesariamente en la fuente. \\
\addlinespace[3pt]
KPIs propuestos & Funciones analíticas planteadas; compare su significado con medida y denominador. \\
\addlinespace[3pt]
De la necesidad al datamart & Muestra cobertura directa, resolución automática, limitación aceptada o decisión pendiente. \\
\addlinespace[3pt]
Comprobar nombres y cobertura & Con permiso apropiado realiza una muestra controlada de la fuente para revisar identidad descriptiva. La muestra permanece en la plataforma y no se envía al LLM. \\
\addlinespace[3pt]
Plan ETL declarativo & Explica operaciones permitidas; esta vista no ejecuta la carga. \\
\addlinespace[3pt]
Ajustes automáticos / Validaciones y advertencias & Lea qué se corrigió y qué sigue exigiendo revisión. \\
\addlinespace[3pt]
Detalles técnicos y trazabilidad & Instantánea, huella, proveedor/modelo y observaciones originales. Una observación del proveedor no es una comprobación de la plataforma. \\
\addlinespace[3pt]
\end{longtable}
\endgroup

\subsection{Paso a paso}
\begin{enumerate}
\item Lea resumen y explicación de negocio.\par\nopagebreak[4]\smallskip{\small\textbf{Compruebe:} El evento coincide con lo solicitado.}
\item Compruebe grano y claves.\par\nopagebreak[4]\smallskip{\small\textbf{Compruebe:} Puede explicar qué representa una fila y cómo se identifica.}
\item Compare cada medida con origen y fórmula.\par\nopagebreak[4]\smallskip{\small\textbf{Compruebe:} Una tasa no está presentada como dinero ni un promedio de línea como promedio por pedido.}
\item Revise dimensiones, rutas e identidad descriptiva.\par\nopagebreak[4]\smallskip{\small\textbf{Compruebe:} Cliente y territorio representan el rol solicitado, con nombre y origen claros.}
\item Lea la cobertura requisito por requisito.\par\nopagebreak[4]\smallskip{\small\textbf{Compruebe:} No se promete una salida omitida.}
\item Revise el plan, ajustes y errores/advertencias.\par\nopagebreak[4]\smallskip{\small\textbf{Compruebe:} Cero errores bloqueantes antes de aprobar; los avisos se entienden.}
\item Elija Personalizar propuesta o Continuar a revisión.\par\nopagebreak[4]\smallskip{\small\textbf{Compruebe:} Pasa a una decisión humana sustentada.}
\end{enumerate}
\subsection{Qué hace la IA}
Propone el modelo y una explicación. Las decisiones originales se conservan para trazabilidad aunque el contrato validado haya aplicado correcciones.

\subsection{Qué comprueba la plataforma}
Existencia de tablas/columnas, tipos, claves, rutas de relación, coherencia del contrato y operaciones calculables; cobertura de requisitos y reglas semánticas implementadas.

\subsection{Qué decide la persona}
El analista decide si el modelo representa el evento comercial y si necesita ajustes.

\paragraph{Resultado esperado} Un diseño comprensible y defendible, o una lista concreta de correcciones.

\Needspace{8\baselineskip}
\begingroup
\setbox0=\vbox{\hsize=\linewidth
{\normalfont\normalsize\sffamily\bfseries\normalsize\color{Teal}\hyphenpenalty=10000\exhyphenpenalty=10000 Errores y recuperación\par}

\par
\begingroup
\fontsize{10}{12.5}\selectfont
\setlength{\tabcolsep}{5pt}
\renewcommand{\arraystretch}{1.15}
\begin{tabular}{@{}>{\RaggedRight\arraybackslash}p{0.27\linewidth}>{\RaggedRight\arraybackslash}p{0.66\linewidth}@{}}
\toprule
\textbf{Situación} & \textbf{Qué hacer} \\
\midrule
Aparecen compras u otro evento en una necesidad de ventas. & No apruebe por tener referencias existentes. Revise alcance y genere una versión que corresponda al evento. \\
\bottomrule
\end{tabular}
\endgroup
}
\Needspace{\dimexpr\ht0+\dp0+48pt\relax}
\endgroup
\subsection{Errores y recuperación}
\begingroup
\fontsize{10}{12.5}\selectfont
\setlength{\tabcolsep}{5pt}
\renewcommand{\arraystretch}{1.15}
\begin{longtable}{@{}>{\RaggedRight\arraybackslash}p{0.27\linewidth}>{\RaggedRight\arraybackslash}p{0.66\linewidth}@{}}
\toprule
\textbf{Situación} & \textbf{Qué hacer} \\
\midrule
\endfirsthead
\toprule
\textbf{Situación} & \textbf{Qué hacer} \\
\midrule
\endhead
\midrule
\multicolumn{2}{r}{\footnotesize Continúa en la página siguiente}\\
\endfoot
\bottomrule
\endlastfoot
Aparecen compras u otro evento en una necesidad de ventas. & No apruebe por tener referencias existentes. Revise alcance y genere una versión que corresponda al evento. \\
\addlinespace[3pt]
Hay nombres como Cliente 123 en lugar de identidad comprobada. & Revise nombres y cobertura con el permiso necesario; no dé por demostrada una entidad descriptiva sólo por su clave. \\
\addlinespace[3pt]
Error bloqueante en una fórmula. & Abra Personalización y excluya o reasigne una receta compatible; si faltan entradas, genere una propuesta nueva con alcance sustentado. \\
\addlinespace[3pt]
\end{longtable}
\endgroup

\subsection{Ejercicio}
Explique al jurado una medida sin leer sólo su nombre: indique origen, operación, unidad y grano.

\textbf{Respuesta orientativa.} Por ejemplo importe registrado de línea, sumado sobre líneas de pedido; no denominarlo cobro ni utilidad.

\subsection{Compruebe lo que aprendió}
\paragraph{¿Todos los campos mostrados ya existen en el origen?} No. Algunos son campos descriptivos que el ETL creará a partir de relaciones verificadas.

\paragraph{¿Comprobar nombres y cobertura envía nombres de clientes al LLM?} La muestra controlada permanece en la plataforma según la interfaz; no se remite al modelo para esa comprobación.

\pantalla{recursos/pantallas/23_grano_y_medidas.jpg}{Granularidad, hecho y medidas verificables}{El grano se define antes de sumar. En este caso, la clave identifica cada línea de pedido. Captura de referencia del 18 de noviembre de 2026; revise las instrucciones actualizadas de esta unidad.}
\pantalla{recursos/pantallas/24_cobertura_necesidad.jpg}{Cobertura y decisiones de la necesidad}{La matriz relaciona requisitos con salidas, decisiones y limitaciones. Una limitación aceptada debe explicarse al presentar el análisis. Captura de referencia del 18 de noviembre de 2026; revise las instrucciones actualizadas de esta unidad.}

\IfFileExists{capturas/BI05.tex}{\input{capturas/BI05.tex}}{}
\section{Asistente, paso 4: ajustar medidas y relaciones con criterio}
\label{unidad:BI06}
\textbf{Para quién:} Analista de BI\par
\textbf{Al terminar podrá:} Crear una versión corregida sin inventar estructuras ni perder las decisiones anteriores.

Personalizar crea una versión nueva. El contrato original conserva su historia. Compare siempre la regla del contrato actual con la regla del borrador: marcar algo en el borrador todavía no lo vuelve ejecutable.

Hay dos acciones de guardado diferentes. Guardar como nueva versión aplica los ajustes generales; Crear versión con relación corregida registra la decisión guiada de relación de manera independiente. Guarde y revise cada resultado antes de encadenar otra corrección.

\subsection{Conceptos para decidir con criterio}
\paragraph{Suma} Combina importes o unidades aditivos del mismo significado y alcance.

\paragraph{Conteo} Cuenta valores no nulos de la columna seleccionada. A grano de línea no equivale automáticamente a documentos.

\paragraph{Conteo distinto} Cuenta una vez cada identificador distinto. Un pedido con varias líneas se cuenta una vez en el universo consultado.

\paragraph{Promedio} Media de los valores de las filas; AVG de importe de línea mide importe medio por línea, no por pedido.

\paragraph{Mínimo y máximo} Extremos de los valores del campo. La línea de mayor importe no identifica necesariamente al producto con mayor suma de ventas.

\paragraph{Relación muchos a uno} Varias filas del hecho apuntan a una fila única de dimensión; la unicidad ayuda a conservar el detalle.

\paragraph{Rol de una relación} Significado del vínculo. Comprador y destinatario de facturación pueden ser clientes válidos, pero responden preguntas diferentes.

\Needspace{8\baselineskip}
\begingroup
\setbox0=\vbox{\hsize=\linewidth
{\normalfont\normalsize\sffamily\bfseries\normalsize\color{Teal}\hyphenpenalty=10000\exhyphenpenalty=10000 Campos y controles\par}

\par
\begingroup
\fontsize{10}{12.5}\selectfont
\setlength{\tabcolsep}{5pt}
\renewcommand{\arraystretch}{1.15}
\begin{tabular}{@{}>{\RaggedRight\arraybackslash}p{0.27\linewidth}>{\RaggedRight\arraybackslash}p{0.66\linewidth}@{}}
\toprule
\textbf{Campo o control} & \textbf{Significado y criterio} \\
\midrule
Resumen de negocio & Hasta 160 caracteres; describe el ajuste sin alterar lo que las fuentes pueden demostrar. \\
\bottomrule
\end{tabular}
\endgroup
}
\Needspace{\dimexpr\ht0+\dp0+48pt\relax}
\endgroup
\subsection{Campos y controles}
\begingroup
\fontsize{10}{12.5}\selectfont
\setlength{\tabcolsep}{5pt}
\renewcommand{\arraystretch}{1.15}
\begin{longtable}{@{}>{\RaggedRight\arraybackslash}p{0.27\linewidth}>{\RaggedRight\arraybackslash}p{0.66\linewidth}@{}}
\toprule
\textbf{Campo o control} & \textbf{Significado y criterio} \\
\midrule
\endfirsthead
\toprule
\textbf{Campo o control} & \textbf{Significado y criterio} \\
\midrule
\endhead
\midrule
\multicolumn{2}{r}{\footnotesize Continúa en la página siguiente}\\
\endfoot
\bottomrule
\endlastfoot
Resumen de negocio & Hasta 160 caracteres; describe el ajuste sin alterar lo que las fuentes pueden demostrar. \\
\addlinespace[3pt]
Granularidad propuesta & Hasta 240 caracteres; el texto debe corresponder a claves y hechos reales, no basta renombrar el grano. \\
\addlinespace[3pt]
Dimensiones incluidas / Medidas incluidas & Conserve al menos las necesarias para la pregunta y sus dependencias. \\
\addlinespace[3pt]
Agregación & Suma, Conteo, Conteo distinto, Promedio, Mínimo o Máximo según el significado. \\
\addlinespace[3pt]
Definir cálculo controlado & Permite operaciones sumar, multiplicar, restar o dividir sobre columnas numéricas disponibles. \\
\addlinespace[3pt]
Primer valor / Segundo valor & En resta y división el orden cambia el resultado; no intercambie costo e importe o cantidad y precio. \\
\addlinespace[3pt]
Factores o componentes en orden & Para suma/multiplicación se seleccionan entradas permitidas; el editor limita la cantidad y el servidor comprueba tipos y referencias. \\
\addlinespace[3pt]
KPIs incluidos y medida asociada & La asociación requiere una medida de función semántica compatible y las entradas de la receta. \\
\addlinespace[3pt]
Dimensión que desea corregir / Ruta de relación comprobada & Permite elegir entre vínculos existentes. Las opciones bloqueadas explican riesgo de duplicación, tipos o unicidad. \\
\addlinespace[3pt]
Justificación de esta corrección / Justificación del ajuste & Entre 10 y 500 caracteres. Explique el concepto de negocio y por qué corresponde, no escriba sólo porque funciona. \\
\addlinespace[3pt]
\end{longtable}
\endgroup

\subsection{Paso a paso}
\begin{enumerate}
\item Abra Personalizar propuesta y anote qué problema concreto va a corregir.\par\nopagebreak[4]\smallskip{\small\textbf{Compruebe:} Ajuste motivado por la necesidad, no por conseguir una luz verde.}
\item Revise cada origen y agregación; marque o desmarque dimensiones, medidas e indicadores.\par\nopagebreak[4]\smallskip{\small\textbf{Compruebe:} El borrador mantiene las dependencias necesarias.}
\item Si corrige un cálculo, seleccione sólo entradas comprobadas y respete el orden.\par\nopagebreak[4]\smallskip{\small\textbf{Compruebe:} Fórmula legible y controlada, sin SQL libre.}
\item Para una ruta ambigua elija dimensión, vínculo y justificación.\par\nopagebreak[4]\smallskip{\small\textbf{Compruebe:} La evidencia muestra tipos, cardinalidad, destino único y riesgo.}
\item Pulse Crear versión con relación corregida cuando corresponda y revise la nueva versión antes de otro ajuste.\par\nopagebreak[4]\smallskip{\small\textbf{Compruebe:} La corrección queda registrada; una decisión intermedia puede precisar una ruta sin cambiar la entidad final.}
\item Para ajustes generales complete Justificación del ajuste y pulse Guardar como nueva versión.\par\nopagebreak[4]\smallskip{\small\textbf{Compruebe:} Una versión derivada pasa de nuevo las reglas; la fuente y la versión original no se reescriben.}
\item Revise cobertura, rutas y errores del resultado y vuelva a Revisión.\par\nopagebreak[4]\smallskip{\small\textbf{Compruebe:} Sólo se aprueba una versión corregida comprobada.}
\end{enumerate}
\subsection{Qué hace la IA}
Puede haber propuesto las opciones iniciales; el ajuste general y la elección de relaciones son decisiones del analista sujetas a validadores. No hace falta pedir a la IA SQL para modificar la estructura.

\subsection{Qué comprueba la plataforma}
Tablas/columnas y tipos, número de entradas, dependencias de KPI, función semántica, clave destino única y conservación del grano; una relación contradictoria o multiplicadora se bloquea.

\subsection{Qué decide la persona}
La persona elige inclusiones, operaciones y roles comerciales, y justifica por escrito la corrección.

\paragraph{Resultado esperado} Nueva versión identificable, revalidada y con decisiones acumuladas.

\Needspace{8\baselineskip}
\begingroup
\setbox0=\vbox{\hsize=\linewidth
{\normalfont\normalsize\sffamily\bfseries\normalsize\color{Teal}\hyphenpenalty=10000\exhyphenpenalty=10000 Errores y recuperación\par}

\par
\begingroup
\fontsize{10}{12.5}\selectfont
\setlength{\tabcolsep}{5pt}
\renewcommand{\arraystretch}{1.15}
\begin{tabular}{@{}>{\RaggedRight\arraybackslash}p{0.27\linewidth}>{\RaggedRight\arraybackslash}p{0.66\linewidth}@{}}
\toprule
\textbf{Situación} & \textbf{Qué hacer} \\
\midrule
Al retirar una medida desaparece un KPI. & Revise sus dependencias; no es pérdida de datos del origen. Restituya una medida compatible o acepte excluir el KPI. \\
\bottomrule
\end{tabular}
\endgroup
}
\Needspace{\dimexpr\ht0+\dp0+48pt\relax}
\endgroup
\subsection{Errores y recuperación}
\begingroup
\fontsize{10}{12.5}\selectfont
\setlength{\tabcolsep}{5pt}
\renewcommand{\arraystretch}{1.15}
\begin{longtable}{@{}>{\RaggedRight\arraybackslash}p{0.27\linewidth}>{\RaggedRight\arraybackslash}p{0.66\linewidth}@{}}
\toprule
\textbf{Situación} & \textbf{Qué hacer} \\
\midrule
\endfirsthead
\toprule
\textbf{Situación} & \textbf{Qué hacer} \\
\midrule
\endhead
\midrule
\multicolumn{2}{r}{\footnotesize Continúa en la página siguiente}\\
\endfoot
\bottomrule
\endlastfoot
Al retirar una medida desaparece un KPI. & Revise sus dependencias; no es pérdida de datos del origen. Restituya una medida compatible o acepte excluir el KPI. \\
\addlinespace[3pt]
El vínculo es válido pero representa facturación en lugar de comprador. & Use la decisión guiada y justifique el rol; que coincidan las filas no demuestra que el rol comercial sea correcto. \\
\addlinespace[3pt]
La nueva versión está bloqueada. & Lea el motivo; cancele o ajuste el borrador. No cambie reglas ni huellas a mano para forzar la aprobación. \\
\addlinespace[3pt]
\end{longtable}
\endgroup

\subsection{Ejercicio}
Un documento tiene líneas de 100 y 50 y otro una línea de 90. Calcule promedio por línea y por documento.

\textbf{Respuesta orientativa.} Importe total 240; tres líneas: 80 por línea; dos documentos distintos: 120 por documento. Son indicadores diferentes.

\subsection{Compruebe lo que aprendió}
\paragraph{¿Editar el texto de granularidad cambia por sí solo las claves?} No. El contrato debe conservar una estructura compatible y superar validación.

\paragraph{¿Una relación uno a muchos desde el hecho es segura por existir?} No necesariamente: puede multiplicar líneas e importes. Revise dirección, unicidad y motivo de bloqueo.

\pantalla{recursos/pantallas/25_medidas_agregaciones.jpg}{Medidas y reglas del contrato conservado}{Contrato conservado en modo de consulta. Se muestran función, origen y regla para comparar conceptos y agregaciones. Captura de referencia del 18 de noviembre de 2026; revise las instrucciones actualizadas de esta unidad.}
\pantalla{recursos/pantallas/27_calculo_controlado.jpg}{Editor de cálculo: ejemplo abierto sin guardar}{Editor abierto como borrador sin guardar: no se eligieron factores ni se creó una versión. Este ejemplo enseña el control, no una fórmula aprobada. Captura de referencia del 18 de noviembre de 2026; revise las instrucciones actualizadas de esta unidad.}
\pantalla{recursos/pantallas/28_relacion_guiada.jpg}{Corrección de relaciones declaradas}{Corrección guiada de una relación. La justificación debe expresar el rol comercial; una FK existente por sí sola no lo demuestra. Captura de referencia del 18 de noviembre de 2026; revise las instrucciones actualizadas de esta unidad.}

\IfFileExists{capturas/BI06.tex}{\input{capturas/BI06.tex}}{}
\section{Asistente, paso 5: aprobar y consultar versiones}
\label{unidad:BI07}
\textbf{Para quién:} Analista de BI con permiso de revisión\par
\textbf{Al terminar podrá:} Distinguir validación técnica, aprobación humana, retiro de aprobación y consulta histórica.

La aprobación es una decisión sobre una versión concreta. La plataforma vuelve a ejecutar sus reglas y registra responsable, fecha y comentario. No crea tablas ni carga datos al pulsar Aprobar propuesta.

Consultar una versión o Verificar evidencia no debe retirar su aprobación por sí solo. Una diferencia de reconstrucción causada por el motor puede quedar como advertencia de compatibilidad; su efecto sobre nuevas cargas se comprueba por separado.

\subsection{Versiones históricas y nuevas reglas de composición}
Los contratos nuevos de esta edición emplean la composición 6 y el prompt sales-bi-v9. Son identificadores técnicos para reproducir el proceso; no son una calificación de la propuesta. Las versiones históricas compuestas con reglas 1 a 5 conservan sus datos y huellas originales.

Una comprobación de compatibilidad no debe reescribir el expediente histórico ni borrar una aprobación al consultarlo. Para una nueva carga se revisan las reglas y compatibilidad que correspondan. Si necesita un diseño nuevo, cree y apruebe una versión nueva, conservando la evidencia anterior.

En la campaña intermedia existió una propuesta que tenía cero errores estructurales, pero omitía el indicador requerido de unidades. No se aprobó ni se presenta aquí como caso final exitoso. La ausencia llevó a reforzar la cobertura. Los casos actuales de referencia son 135 y 136.

\subsection{Conceptos para decidir con criterio}
\paragraph{Lista para revisión} Contrato que puede recibir decisión humana si conserva las condiciones requeridas.

\paragraph{Aprobada} Versión con decisión humana registrada; su aptitud para otra carga también depende de las reglas actuales y la fuente.

\paragraph{Retirada} Aprobación retirada expresamente para impedir nuevas ejecuciones; no borra datamarts ya generados.

\paragraph{Descartada} Fuera de las listas operativas; se conservan auditoría y trazabilidad.

\paragraph{Huella} Resumen de contenido usado para detectar cambios. No se corrige editando el valor para que coincida.

\Needspace{8\baselineskip}
\begingroup
\setbox0=\vbox{\hsize=\linewidth
{\normalfont\normalsize\sffamily\bfseries\normalsize\color{Teal}\hyphenpenalty=10000\exhyphenpenalty=10000 Campos y controles\par}

\par
\begingroup
\fontsize{10}{12.5}\selectfont
\setlength{\tabcolsep}{5pt}
\renewcommand{\arraystretch}{1.15}
\begin{tabular}{@{}>{\RaggedRight\arraybackslash}p{0.27\linewidth}>{\RaggedRight\arraybackslash}p{0.66\linewidth}@{}}
\toprule
\textbf{Campo o control} & \textbf{Significado y criterio} \\
\midrule
Comentario de revisión & Hasta 500 caracteres. Para rechazar se exige un motivo suficiente, al menos 10 caracteres en la interfaz. \\
\bottomrule
\end{tabular}
\endgroup
}
\Needspace{\dimexpr\ht0+\dp0+48pt\relax}
\endgroup
\subsection{Campos y controles}
\begingroup
\fontsize{10}{12.5}\selectfont
\setlength{\tabcolsep}{5pt}
\renewcommand{\arraystretch}{1.15}
\begin{longtable}{@{}>{\RaggedRight\arraybackslash}p{0.27\linewidth}>{\RaggedRight\arraybackslash}p{0.66\linewidth}@{}}
\toprule
\textbf{Campo o control} & \textbf{Significado y criterio} \\
\midrule
\endfirsthead
\toprule
\textbf{Campo o control} & \textbf{Significado y criterio} \\
\midrule
\endhead
\midrule
\multicolumn{2}{r}{\footnotesize Continúa en la página siguiente}\\
\endfoot
\bottomrule
\endlastfoot
Comentario de revisión & Hasta 500 caracteres. Para rechazar se exige un motivo suficiente, al menos 10 caracteres en la interfaz. \\
\addlinespace[3pt]
Leí y comprendí las advertencias y las acciones indicadas & Aceptación requerida si hay avisos. Cero errores no significa cero aspectos que revisar. \\
\addlinespace[3pt]
Aprobar propuesta / Rechazar propuesta & Decisión humana distinta de la ejecución física. \\
\addlinespace[3pt]
Verificar evidencia & Comprueba instantánea, referencias, contrato y reproducción sin llamar al modelo, sin consultar filas y sin ETL. \\
\addlinespace[3pt]
Mostrar en Versiones generadas & Filtra el historial; no cambia el estado de las versiones. \\
\addlinespace[3pt]
Abrir resultado & Recupera la versión guardada sin volver a generar con IA. \\
\addlinespace[3pt]
Crear versión corregida desde ésta & Permite derivar una corrección conservando la original. \\
\addlinespace[3pt]
Retirar aprobación / Descartar versión & Acciones explícitas con motivo; no se usan para limpiar un aviso sin entenderlo. \\
\addlinespace[3pt]
Restaurar aprobación & Disponible en los estados y permisos apropiados; exige justificación, advertencias revisadas y revalidación. \\
\addlinespace[3pt]
\end{longtable}
\endgroup

\subsection{Paso a paso}
\begin{enumerate}
\item Lea necesidad, fuente, dominio e instantánea en Revisión.\par\nopagebreak[4]\smallskip{\small\textbf{Compruebe:} Está decidiendo sobre la versión correcta.}
\item Confirme cero errores y comprenda cada advertencia.\par\nopagebreak[4]\smallskip{\small\textbf{Compruebe:} Puede explicar qué está comprobado y qué límite conserva.}
\item Escriba comentario y apruebe si corresponde; rechace con motivo si no cumple.\par\nopagebreak[4]\smallskip{\small\textbf{Compruebe:} Comprobante de estado, responsable y fecha.}
\item Abra después el historial y consulte la misma versión.\par\nopagebreak[4]\smallskip{\small\textbf{Compruebe:} Se puede recorrer el diseño en modo de lectura sin regenerarlo.}
\item Si se necesita corregir, cree una derivada y revise su aprobación aparte.\par\nopagebreak[4]\smallskip{\small\textbf{Compruebe:} No cambia silenciosamente la evidencia del caso anterior.}
\end{enumerate}
\subsection{Qué hace la IA}
La IA no aprueba. Verificar evidencia es una comprobación determinística, no una nueva opinión del proveedor.

\subsection{Qué comprueba la plataforma}
Permisos, estado y reglas de aprobación vigentes; integridad de instantánea y contrato. Los errores bloqueantes no se resuelven marcando el aviso de advertencias.

\subsection{Qué decide la persona}
Una persona autorizada aprueba/rechaza o retira/restaura con un motivo.

\paragraph{Resultado esperado} Una decisión trazable y versiones que conservan su historia.

\Needspace{8\baselineskip}
\begingroup
\setbox0=\vbox{\hsize=\linewidth
{\normalfont\normalsize\sffamily\bfseries\normalsize\color{Teal}\hyphenpenalty=10000\exhyphenpenalty=10000 Errores y recuperación\par}

\par
\begingroup
\fontsize{10}{12.5}\selectfont
\setlength{\tabcolsep}{5pt}
\renewcommand{\arraystretch}{1.15}
\begin{tabular}{@{}>{\RaggedRight\arraybackslash}p{0.27\linewidth}>{\RaggedRight\arraybackslash}p{0.66\linewidth}@{}}
\toprule
\textbf{Situación} & \textbf{Qué hacer} \\
\midrule
No ve una versión en el filtro actual. & Cambie Mostrar al estado correspondiente y compruebe la fuente. No suponga que se eliminó. \\
\bottomrule
\end{tabular}
\endgroup
}
\Needspace{\dimexpr\ht0+\dp0+48pt\relax}
\endgroup
\subsection{Errores y recuperación}
\begingroup
\fontsize{10}{12.5}\selectfont
\setlength{\tabcolsep}{5pt}
\renewcommand{\arraystretch}{1.15}
\begin{longtable}{@{}>{\RaggedRight\arraybackslash}p{0.27\linewidth}>{\RaggedRight\arraybackslash}p{0.66\linewidth}@{}}
\toprule
\textbf{Situación} & \textbf{Qué hacer} \\
\midrule
\endfirsthead
\toprule
\textbf{Situación} & \textbf{Qué hacer} \\
\midrule
\endhead
\midrule
\multicolumn{2}{r}{\footnotesize Continúa en la página siguiente}\\
\endfoot
\bottomrule
\endlastfoot
No ve una versión en el filtro actual. & Cambie Mostrar al estado correspondiente y compruebe la fuente. No suponga que se eliminó. \\
\addlinespace[3pt]
La propuesta aprobada no aparece habilitada para una nueva carga. & Consulte motivos en Datamart de ventas. Puede ser histórica o incompatible para otra ejecución aunque sus resultados anteriores sigan disponibles. \\
\addlinespace[3pt]
Verificar evidencia muestra compatibilidad. & Lea la explicación y distinga el registro histórico de la aptitud actual. No retire aprobación sin una decisión deliberada. \\
\addlinespace[3pt]
\end{longtable}
\endgroup

\subsection{Ejercicio}
Explique la diferencia entre aprobar propuesta, preparar materialización y materializar.

\textbf{Respuesta orientativa.} Aprobar acepta un diseño; preparar fija selección y huellas; materializar lee la fuente y crea/carga el destino.

\subsection{Compruebe lo que aprendió}
\paragraph{¿Una aprobada puede ser consultable pero no apta para otra carga?} Sí. Se conserva historia, pero nuevas ejecuciones deben superar las reglas actuales.

\paragraph{¿Abrir resultado vuelve a consumir una generación IA?} No. Recupera una versión persistida.

\pantalla{recursos/pantallas/21_revision_aprobada.jpg}{Aprobación humana y versión conservada}{Aprobación histórica de la propuesta 129. Consultar el resultado conserva la decisión y no ejecuta el ETL. Captura de referencia del 18 de noviembre de 2026; revise las instrucciones actualizadas de esta unidad.}
\pantalla{recursos/pantallas/26_evidencia_estructural.jpg}{Cuatro controles estructurales y huellas coincidentes}{Reverificación del 18 de noviembre: cuatro controles cumplen y las huellas coinciden. Es validación estructural, no una nueva carga. Captura de referencia del 18 de noviembre de 2026; revise las instrucciones actualizadas de esta unidad.}

\IfFileExists{capturas/BI07.tex}{\input{capturas/BI07.tex}}{}

\chapter{Materializar, conciliar y revisar los resultados}
\begingroup
\setlength{\parskip}{4pt plus 1pt minus 1pt}
\setlist[enumerate]{itemsep=3pt,topsep=4pt}
\titlespacing*{\paragraph}{0pt}{4pt}{0.5em}
\section{Materialización, paso 1: elegir una propuesta}
\label{unidad:ETL01}
\textbf{Para quién:} Analista de BI\par
\textbf{Al terminar podrá:} Escoger un contrato aprobado entendiendo qué puede ejecutarse y qué debe abrirse como histórico.

Entre a Datamart de ventas con la fuente que corresponde. El proceso muestra cinco pasos: Elegir propuesta, Revisar indicadores, Confirmar ejecución, Materializar y cargar, y Validar resultados.

La propuesta recomendada es una ayuda para elegir entre candidatas compatibles. No ejecuta nada automáticamente. Si el mismo contrato y selección ya tienen expediente, se debe abrir ese expediente en lugar de repetir la carga.

\subsection{Conceptos para decidir con criterio}
\paragraph{Propuesta elegible} Conserva aprobación, metadatos, fuente habilitada y probada, contrato compatible e indicadores calculables.

\paragraph{Expediente} Registro de una preparación o ejecución, con selección, huellas, plan, responsable y resultados.

\paragraph{Consulta histórica} Lectura de lo ocurrido en una ejecución anterior; no es una prueba nueva ni una recarga.

\Needspace{8\baselineskip}
\begingroup
\setbox0=\vbox{\hsize=\linewidth
{\normalfont\normalsize\sffamily\bfseries\normalsize\color{Teal}\hyphenpenalty=10000\exhyphenpenalty=10000 Campos y controles\par}

\par
\begingroup
\fontsize{10}{12.5}\selectfont
\setlength{\tabcolsep}{5pt}
\renewcommand{\arraystretch}{1.15}
\begin{tabular}{@{}>{\RaggedRight\arraybackslash}p{0.27\linewidth}>{\RaggedRight\arraybackslash}p{0.66\linewidth}@{}}
\toprule
\textbf{Campo o control} & \textbf{Significado y criterio} \\
\midrule
Tarjeta de versión & Resumen, necesidad, grano, cantidad de KPI y fecha de aprobación. \\
\bottomrule
\end{tabular}
\endgroup
}
\Needspace{\dimexpr\ht0+\dp0+48pt\relax}
\endgroup
\subsection{Campos y controles}
\begingroup
\fontsize{10}{12.5}\selectfont
\setlength{\tabcolsep}{5pt}
\renewcommand{\arraystretch}{1.15}
\begin{longtable}{@{}>{\RaggedRight\arraybackslash}p{0.27\linewidth}>{\RaggedRight\arraybackslash}p{0.66\linewidth}@{}}
\toprule
\textbf{Campo o control} & \textbf{Significado y criterio} \\
\midrule
\endfirsthead
\toprule
\textbf{Campo o control} & \textbf{Significado y criterio} \\
\midrule
\endhead
\midrule
\multicolumn{2}{r}{\footnotesize Continúa en la página siguiente}\\
\endfoot
\bottomrule
\endlastfoot
Tarjeta de versión & Resumen, necesidad, grano, cantidad de KPI y fecha de aprobación. \\
\addlinespace[3pt]
Seleccionar & Elige el contrato del recorrido actual. \\
\addlinespace[3pt]
Comparar & Permite contrastar hasta dos propuestas por grano, medidas e indicadores. \\
\addlinespace[3pt]
Versiones anteriores no habilitadas para una nueva carga & Explica bloqueos; no significa que se hayan borrado sus datos. \\
\addlinespace[3pt]
Abrir expediente & Recupera una ejecución previa sin repetir el ETL. \\
\addlinespace[3pt]
Revisar indicadores & Avanza con la propuesta seleccionada. \\
\addlinespace[3pt]
\end{longtable}
\endgroup

\subsection{Paso a paso}
\begin{enumerate}
\item Compruebe la fuente antes de leer las tarjetas.\par\nopagebreak[4]\smallskip{\small\textbf{Compruebe:} No mezcla contratos de otra base.}
\item Compare necesidad y grano con su ficha de análisis.\par\nopagebreak[4]\smallskip{\small\textbf{Compruebe:} Selecciona un modelo adecuado, no sólo la versión de número mayor.}
\item Use Comparar si existen dos alternativas relevantes.\par\nopagebreak[4]\smallskip{\small\textbf{Compruebe:} Reconoce diferencias que afectan al significado comercial.}
\item Si indica Ya ejecutada, abra el expediente para consultar la evidencia.\par\nopagebreak[4]\smallskip{\small\textbf{Compruebe:} No produce otra carga idéntica accidental.}
\item Si requiere una carga nueva y la candidata es compatible, pulse Revisar indicadores.\par\nopagebreak[4]\smallskip{\small\textbf{Compruebe:} Continúa con la versión aprobada elegida.}
\end{enumerate}
\subsection{Qué hace la IA}
Los indicadores fueron sugeridos en el diseño; esta selección no exige otra generación con IA.

\subsection{Qué comprueba la plataforma}
Reglas de compilación actuales, aprobación, instantánea existente, fuente activa en lectura, recetas y preparación dimensional. Contratos históricos anteriores pueden quedar excluidos de nuevas ejecuciones.

\subsection{Qué decide la persona}
La persona selecciona la versión y decide si consulta lo previo o necesita una derivada corregida.

\paragraph{Resultado esperado} Un contrato correcto seleccionado o un expediente histórico abierto de forma consciente.

\Needspace{8\baselineskip}
\begingroup
\setbox0=\vbox{\hsize=\linewidth
{\normalfont\normalsize\sffamily\bfseries\normalsize\color{Teal}\hyphenpenalty=10000\exhyphenpenalty=10000 Errores y recuperación\par}

\par
\begingroup
\fontsize{10}{12.5}\selectfont
\setlength{\tabcolsep}{5pt}
\renewcommand{\arraystretch}{1.15}
\begin{tabular}{@{}>{\RaggedRight\arraybackslash}p{0.27\linewidth}>{\RaggedRight\arraybackslash}p{0.66\linewidth}@{}}
\toprule
\textbf{Situación} & \textbf{Qué hacer} \\
\midrule
No hay propuestas habilitadas. & Lea la lista de bloqueos; corrija o genere una propuesta actual. El historial puede seguir disponible. \\
\bottomrule
\end{tabular}
\endgroup
}
\Needspace{\dimexpr\ht0+\dp0+48pt\relax}
\endgroup
\subsection{Errores y recuperación}
\begingroup
\fontsize{10}{12.5}\selectfont
\setlength{\tabcolsep}{5pt}
\renewcommand{\arraystretch}{1.15}
\begin{longtable}{@{}>{\RaggedRight\arraybackslash}p{0.27\linewidth}>{\RaggedRight\arraybackslash}p{0.66\linewidth}@{}}
\toprule
\textbf{Situación} & \textbf{Qué hacer} \\
\midrule
\endfirsthead
\toprule
\textbf{Situación} & \textbf{Qué hacer} \\
\midrule
\endhead
\midrule
\multicolumn{2}{r}{\footnotesize Continúa en la página siguiente}\\
\endfoot
\bottomrule
\endlastfoot
No hay propuestas habilitadas. & Lea la lista de bloqueos; corrija o genere una propuesta actual. El historial puede seguir disponible. \\
\addlinespace[3pt]
Sólo se ve Abrir expediente. & Es la protección frente a repetición de contrato. Consulte el existente; no trate de borrar la ejecución para forzar otra. \\
\addlinespace[3pt]
\end{longtable}
\endgroup

\subsection{Ejercicio}
Elija dos propuestas y explique una diferencia de grano o indicador sin ejecutarlas.

\textbf{Respuesta orientativa.} La comparación tiene sentido si cambia la pregunta o su representación; el número de versión por sí solo no mide calidad.

\subsection{Compruebe lo que aprendió}
\paragraph{¿Recomendada equivale a aprobada automáticamente por la IA?} No. Es una recomendación operativa entre contratos que conservan aprobación y compatibilidad.

\paragraph{¿Histórico bloqueado equivale a datos eliminados?} No. Puede consultarse el expediente conservado.

\endgroup
\pantalla{recursos/pantallas/31_etl_elegir_propuesta.jpg}{Selección y control de ejecuciones idénticas}{Propuesta 129 ya materializada en el expediente 21. Una carga idéntica no se vuelve a habilitar accidentalmente. Captura de referencia del 18 de noviembre de 2026; revise las instrucciones actualizadas de esta unidad.}

\IfFileExists{capturas/ETL01.tex}{\input{capturas/ETL01.tex}}{}
\section{Materialización, paso 2: elegir indicadores y recetas}
\label{unidad:ETL02}
\textbf{Para quién:} Analista de BI\par
\textbf{Al terminar podrá:} Seleccionar salidas calculables y explicar suma, diferencia, razón y participación.

En Revisar indicadores no se redactan fórmulas libres. Se decide cuáles de las recetas admitidas del contrato se necesitan para el objetivo. Debe quedar al menos una.

Desmarcar indicadores puede reducir la cobertura de la necesidad. La preparación registra si la selección es parcial. No presente un ejercicio de dos KPI como si hubiera materializado todos los requisitos de una necesidad avanzada.

\subsection{Conceptos para decidir con criterio}
\paragraph{Agregación} Suma, conteo, conteo distinto, promedio, mínimo o máximo de una medida comprobada.

\paragraph{Diferencia} Resta entre dos entradas verificadas, por ejemplo importe menos costo comparable.

\paragraph{Razón} Numerador dividido entre denominador. El importe por pedido usa importe total \(\div\) pedidos distintos del mismo alcance.

\paragraph{Participación} Razón multiplicada por 100. Siempre se debe explicar el universo del denominador.

\paragraph{Promedio ponderado por unidad} Total monetario \(\div\) total de unidades, distinto de promediar precios de filas sin ponderar.

\paragraph{Sin valor} Una razón sin denominador válido no debe representarse como una división infinita ni como una venta cero demostrada.

\Needspace{8\baselineskip}
\begingroup
\setbox0=\vbox{\hsize=\linewidth
{\normalfont\normalsize\sffamily\bfseries\normalsize\color{Teal}\hyphenpenalty=10000\exhyphenpenalty=10000 Campos y controles\par}

\par
\begingroup
\fontsize{10}{12.5}\selectfont
\setlength{\tabcolsep}{5pt}
\renewcommand{\arraystretch}{1.15}
\begin{tabular}{@{}>{\RaggedRight\arraybackslash}p{0.27\linewidth}>{\RaggedRight\arraybackslash}p{0.66\linewidth}@{}}
\toprule
\textbf{Campo o control} & \textbf{Significado y criterio} \\
\midrule
Incluir indicador & Conserva o retira una receta de esta ejecución. \\
\bottomrule
\end{tabular}
\endgroup
}
\Needspace{\dimexpr\ht0+\dp0+48pt\relax}
\endgroup
\subsection{Campos y controles}
\begingroup
\fontsize{10}{12.5}\selectfont
\setlength{\tabcolsep}{5pt}
\renewcommand{\arraystretch}{1.15}
\begin{longtable}{@{}>{\RaggedRight\arraybackslash}p{0.27\linewidth}>{\RaggedRight\arraybackslash}p{0.66\linewidth}@{}}
\toprule
\textbf{Campo o control} & \textbf{Significado y criterio} \\
\midrule
\endfirsthead
\toprule
\textbf{Campo o control} & \textbf{Significado y criterio} \\
\midrule
\endhead
\midrule
\multicolumn{2}{r}{\footnotesize Continúa en la página siguiente}\\
\endfoot
\bottomrule
\endlastfoot
Incluir indicador & Conserva o retira una receta de esta ejecución. \\
\addlinespace[3pt]
Cálculo controlado & Describe operación y entradas; compruebe el denominador y la unidad. \\
\addlinespace[3pt]
Unidad & Unidades, documentos, porcentaje o moneda de origen; un código monetario necesita evidencia. \\
\addlinespace[3pt]
Período & Hereda o precisa la periodicidad del contrato. \\
\addlinespace[3pt]
Ajuste de seguridad & Por ejemplo una tasa presentada como dinero o un promedio de línea mal nombrado. \\
\addlinespace[3pt]
Ver trazabilidad técnica & Entradas verificadas, versión de definición y unidad originalmente sugerida. \\
\addlinespace[3pt]
Revisar transformaciones & Avanza con las recetas seleccionadas. \\
\addlinespace[3pt]
\end{longtable}
\endgroup

\subsection{Paso a paso}
\begin{enumerate}
\item Lea nombre, descripción y cálculo de cada indicador.\par\nopagebreak[4]\smallskip{\small\textbf{Compruebe:} Reconoce qué pregunta responde.}
\item Compruebe que el conteo documental sea distinto y que la razón use ese conteo.\par\nopagebreak[4]\smallskip{\small\textbf{Compruebe:} No confunde líneas con pedidos o facturas.}
\item Seleccione las salidas necesarias y revise las que retira.\par\nopagebreak[4]\smallskip{\small\textbf{Compruebe:} Queda al menos una receta y el alcance parcial se reconoce.}
\item Lea las correcciones de nombre o unidad.\par\nopagebreak[4]\smallskip{\small\textbf{Compruebe:} No denomina neto a precio \(\times\) cantidad sin descontar una tasa ni promedio por pedido a AVG de línea.}
\item Pulse Revisar transformaciones.\par\nopagebreak[4]\smallskip{\small\textbf{Compruebe:} Continúa sin cambiar el contrato físico ni ejecutar todavía.}
\end{enumerate}
\subsection{Qué hace la IA}
Propuso indicadores anteriormente; el cálculo lo controla el motor de recetas. La IA no calcula a mano el total que se publicará.

\subsection{Qué comprueba la plataforma}
Operaciones soportadas y entradas comprobadas. Las diferencias, razones y participaciones requieren dos entradas válidas. La división con divisor cero conserva un resultado nulo/no calculable.

\subsection{Qué decide la persona}
El analista selecciona qué indicadores quiere materializar y acepta conscientemente un alcance parcial cuando corresponda.

\paragraph{Resultado esperado} Selección de indicadores útil, con fórmula y unidad entendidas.

\Needspace{8\baselineskip}
\begingroup
\setbox0=\vbox{\hsize=\linewidth
{\normalfont\normalsize\sffamily\bfseries\normalsize\color{Teal}\hyphenpenalty=10000\exhyphenpenalty=10000 Errores y recuperación\par}

\par
\begingroup
\fontsize{10}{12.5}\selectfont
\setlength{\tabcolsep}{5pt}
\renewcommand{\arraystretch}{1.15}
\begin{tabular}{@{}>{\RaggedRight\arraybackslash}p{0.27\linewidth}>{\RaggedRight\arraybackslash}p{0.66\linewidth}@{}}
\toprule
\textbf{Situación} & \textbf{Qué hacer} \\
\midrule
Un KPI no tiene medida compatible. & Vuelva al asistente para corregir o excluir; este paso no permite inventar la entrada. \\
\bottomrule
\end{tabular}
\endgroup
}
\Needspace{\dimexpr\ht0+\dp0+48pt\relax}
\endgroup
\subsection{Errores y recuperación}
\begingroup
\fontsize{10}{12.5}\selectfont
\setlength{\tabcolsep}{5pt}
\renewcommand{\arraystretch}{1.15}
\begin{longtable}{@{}>{\RaggedRight\arraybackslash}p{0.27\linewidth}>{\RaggedRight\arraybackslash}p{0.66\linewidth}@{}}
\toprule
\textbf{Situación} & \textbf{Qué hacer} \\
\midrule
\endfirsthead
\toprule
\textbf{Situación} & \textbf{Qué hacer} \\
\midrule
\endhead
\midrule
\multicolumn{2}{r}{\footnotesize Continúa en la página siguiente}\\
\endfoot
\bottomrule
\endlastfoot
Un KPI no tiene medida compatible. & Vuelva al asistente para corregir o excluir; este paso no permite inventar la entrada. \\
\addlinespace[3pt]
La razón aparece Sin valor. & Revise denominador y universo. No reemplace el resultado por cero para ocultar una división no definida. \\
\addlinespace[3pt]
Las facturas por producto no suman las facturas globales. & El conteo distinto no es aditivo: una factura puede contener varios productos. \\
\addlinespace[3pt]
\end{longtable}
\endgroup

\subsection{Ejercicio}
Tres líneas: pedido A 100 y 50; pedido B 90. Calcule total, conteo de líneas, pedidos distintos e importe por pedido.

\textbf{Respuesta orientativa.} 240; 3; 2; 120. Sumar los pedidos distintos de grupos de producto puede contar el mismo pedido varias veces.

\subsection{Compruebe lo que aprendió}
\paragraph{¿Por qué 10 \% y 40 \% no se promedian siempre como 25 \%?} Si tienen distintos denominadores, hay que usar totales compatibles o ponderar; el promedio simple puede cambiar el significado.

\paragraph{¿Cambiar la etiqueta corrige una receta incorrecta?} No. Debe corresponder fórmula, medida y unidad, no sólo el texto.

\pantalla{recursos/pantallas/32_etl_indicadores.jpg}{Indicadores y recetas del ETL}{Cuatro indicadores conservados del expediente. Una razón requiere un numerador y un denominador comercialmente compatibles. Captura de referencia del 18 de noviembre de 2026; revise las instrucciones actualizadas de esta unidad.}

\IfFileExists{capturas/ETL02.tex}{\input{capturas/ETL02.tex}}{}
\section{Materialización, paso 3: confirmar el plan}
\label{unidad:ETL03}
\textbf{Para quién:} Analista de BI con permiso de ejecución\par
\textbf{Al terminar podrá:} Leer extracción, limpieza, transformación, carga y validación antes de fijar el expediente.

Esta pantalla permite explicar el proceso sin escribir SQL. Revise la tabla de hechos, el grano, las dimensiones y las medidas. Luego lea cada operación de la secuencia.

Las operaciones opcionales marcadas Revisar en el paso 5 preparan candidatos de etiquetas en español. No existe una lista para aprobar etiquetas aquí: esa revisión aparecerá después de la carga si hay candidatos.

\subsection{Conceptos para decidir con criterio}
\paragraph{Extracción} Lectura controlada de columnas del origen autorizado.

\paragraph{Limpieza y calidad} Tratamiento previsto de claves, nulos, textos o duplicados según el plan; no es una promesa de corregir cualquier dato defectuoso.

\paragraph{Transformación} Derivaciones y enriquecimientos permitidos que respetan la granularidad.

\paragraph{Clave sustituta} Identificador interno de la dimensión analítica, separado de la clave comercial del origen.

\paragraph{Carga dimensional} Creación y carga del hecho y dimensiones del destino de la ejecución.

\paragraph{Conciliación} Contraste de cantidades e importes de origen y destino del mismo alcance.

\Needspace{8\baselineskip}
\begingroup
\setbox0=\vbox{\hsize=\linewidth
{\normalfont\normalsize\sffamily\bfseries\normalsize\color{Teal}\hyphenpenalty=10000\exhyphenpenalty=10000 Campos y controles\par}

\par
\begingroup
\fontsize{10}{12.5}\selectfont
\setlength{\tabcolsep}{5pt}
\renewcommand{\arraystretch}{1.15}
\begin{tabular}{@{}>{\RaggedRight\arraybackslash}p{0.27\linewidth}>{\RaggedRight\arraybackslash}p{0.66\linewidth}@{}}
\toprule
\textbf{Campo o control} & \textbf{Significado y criterio} \\
\midrule
Tabla de hechos / Granularidad / Dimensiones / Medidas & Resumen del modelo que se va a materializar. \\
\bottomrule
\end{tabular}
\endgroup
}
\Needspace{\dimexpr\ht0+\dp0+48pt\relax}
\endgroup
\subsection{Campos y controles}
\begingroup
\fontsize{10}{12.5}\selectfont
\setlength{\tabcolsep}{5pt}
\renewcommand{\arraystretch}{1.15}
\begin{longtable}{@{}>{\RaggedRight\arraybackslash}p{0.27\linewidth}>{\RaggedRight\arraybackslash}p{0.66\linewidth}@{}}
\toprule
\textbf{Campo o control} & \textbf{Significado y criterio} \\
\midrule
\endfirsthead
\toprule
\textbf{Campo o control} & \textbf{Significado y criterio} \\
\midrule
\endhead
\midrule
\multicolumn{2}{r}{\footnotesize Continúa en la página siguiente}\\
\endfoot
\bottomrule
\endlastfoot
Tabla de hechos / Granularidad / Dimensiones / Medidas & Resumen del modelo que se va a materializar. \\
\addlinespace[3pt]
Etapas del plan & Extracción, Limpieza y calidad, Transformación y enriquecimiento, Carga dimensional y Conciliación y evidencia. \\
\addlinespace[3pt]
Advertencias heredadas & Avisos de la propuesta que siguen siendo relevantes para la carga. \\
\addlinespace[3pt]
Registro de decisión & Comentario de 10 a 500 caracteres que explica la decisión sobre plan y KPI. \\
\addlinespace[3pt]
Confirmo el plan de ejecución & Acepta el contrato y entiende que las etiquetas candidatas se revisarán después. \\
\addlinespace[3pt]
Preparar materialización & Guarda selección, plan y huellas. No lee ni altera filas de la fuente en esta confirmación. \\
\addlinespace[3pt]
\end{longtable}
\endgroup

\subsection{Paso a paso}
\begin{enumerate}
\item Compare el resumen con la propuesta aprobada.\par\nopagebreak[4]\smallskip{\small\textbf{Compruebe:} Mismo evento y grano.}
\item Lea las operaciones de cada etapa y advertencias.\par\nopagebreak[4]\smallskip{\small\textbf{Compruebe:} Puede explicar qué se leerá y qué se calculará.}
\item Compruebe que no se confunda una etiqueta en español con un valor original nuevo.\par\nopagebreak[4]\smallskip{\small\textbf{Compruebe:} Entiende la política de conservar originales.}
\item Escriba una decisión concreta y marque la confirmación.\par\nopagebreak[4]\smallskip{\small\textbf{Compruebe:} Registro humano suficiente.}
\item Pulse Preparar materialización una sola vez y anote el número de ejecución.\par\nopagebreak[4]\smallskip{\small\textbf{Compruebe:} Expediente preparado con propuesta, instantánea, responsable y huellas; aún no hay nueva carga física.}
\end{enumerate}
\subsection{Qué hace la IA}
No genera SQL libre en esta confirmación. La interpretación posterior de categorías puede usar IA bajo la política de valores aptos y revisión humana.

\subsection{Qué comprueba la plataforma}
Permiso de escritura ETL, aprobación y compatibilidad actuales, recetas seleccionadas, cobertura y duplicidad de contrato/selección.

\subsection{Qué decide la persona}
Una persona autorizada confirma el plan y la selección que quedará registrada.

\paragraph{Resultado esperado} Un expediente preparado, trazable y aún separado de la ejecución física.

\Needspace{8\baselineskip}
\begingroup
\setbox0=\vbox{\hsize=\linewidth
{\normalfont\normalsize\sffamily\bfseries\normalsize\color{Teal}\hyphenpenalty=10000\exhyphenpenalty=10000 Errores y recuperación\par}

\par
\begingroup
\fontsize{10}{12.5}\selectfont
\setlength{\tabcolsep}{5pt}
\renewcommand{\arraystretch}{1.15}
\begin{tabular}{@{}>{\RaggedRight\arraybackslash}p{0.27\linewidth}>{\RaggedRight\arraybackslash}p{0.66\linewidth}@{}}
\toprule
\textbf{Situación} & \textbf{Qué hacer} \\
\midrule
La preparación informa que el contrato ya corresponde a otra ejecución. & Abra ese expediente; el conflicto evita repetir la carga idéntica. \\
\bottomrule
\end{tabular}
\endgroup
}
\Needspace{\dimexpr\ht0+\dp0+48pt\relax}
\endgroup
\subsection{Errores y recuperación}
\begingroup
\fontsize{10}{12.5}\selectfont
\setlength{\tabcolsep}{5pt}
\renewcommand{\arraystretch}{1.15}
\begin{longtable}{@{}>{\RaggedRight\arraybackslash}p{0.27\linewidth}>{\RaggedRight\arraybackslash}p{0.66\linewidth}@{}}
\toprule
\textbf{Situación} & \textbf{Qué hacer} \\
\midrule
\endfirsthead
\toprule
\textbf{Situación} & \textbf{Qué hacer} \\
\midrule
\endhead
\midrule
\multicolumn{2}{r}{\footnotesize Continúa en la página siguiente}\\
\endfoot
\bottomrule
\endlastfoot
La preparación informa que el contrato ya corresponde a otra ejecución. & Abra ese expediente; el conflicto evita repetir la carga idéntica. \\
\addlinespace[3pt]
El botón sigue deshabilitado. & Revise permiso, casilla de confirmación y longitud del comentario. \\
\addlinespace[3pt]
Faltan etiquetas para aprobar en esta pantalla. & Es normal: primero se materializa, después se revisan los candidatos en Validar resultados. \\
\addlinespace[3pt]
\end{longtable}
\endgroup

\subsection{Ejercicio}
Explique por qué preparar no significa que ya se pueda mostrar un nuevo tablero.

\textbf{Respuesta orientativa.} Se guardó el contrato; faltan lectura, carga, conciliación y condiciones de publicación.

\subsection{Compruebe lo que aprendió}
\paragraph{¿Qué se conserva al preparar?} Selección de KPI, comentario, plan, huellas de propuesta e instantánea y responsable.

\paragraph{¿Quién ejecuta la lectura?} El materializador controlado después del inicio explícito del paso 4.

\pantalla{recursos/pantallas/33_etl_plan.jpg}{Plan controlado previo a la ejecución}{Plan explicado por etapas. Esta consulta histórica permite aprender los controles sin repetir la carga. Captura de referencia del 18 de noviembre de 2026; revise las instrucciones actualizadas de esta unidad.}

\IfFileExists{capturas/ETL03.tex}{\input{capturas/ETL03.tex}}{}
\section{Materialización, paso 4: cargar sin repetir accidentalmente}
\label{unidad:ETL04}
\textbf{Para quién:} Analista de BI con permiso de ejecución\par
\textbf{Al terminar podrá:} Iniciar una carga entendiendo sus efectos y recuperar el expediente ante una interrupción.

Antes de pulsar Materializar y cargar datamart, compruebe fuente, propuesta y número de expediente. La operación leerá columnas aprobadas de la fuente y creará un destino analítico propio de esa ejecución.

El hecho y las dimensiones se cargan con controles transaccionales. Un fallo bloqueante no debe dejar un datamart parcial publicado ni sustituir los resultados anteriores. No pulse repetidamente el botón porque tarde la respuesta.

\subsection{Conceptos para decidir con criterio}
\paragraph{Esquema por ejecución} Espacio analítico identificado para separar resultados de distintas ejecuciones; la base de origen no se sobrescribe.

\paragraph{Transacción} Conjunto de operaciones que se confirma completo o se revierte ante un fallo.

\paragraph{Preparada / En ejecución / Validada / Revisión pendiente / Fallida} Estados del expediente que se deben leer antes de intentar otra acción.

\Needspace{8\baselineskip}
\begingroup
\setbox0=\vbox{\hsize=\linewidth
{\normalfont\normalsize\sffamily\bfseries\normalsize\color{Teal}\hyphenpenalty=10000\exhyphenpenalty=10000 Campos y controles\par}

\par
\begingroup
\fontsize{10}{12.5}\selectfont
\setlength{\tabcolsep}{5pt}
\renewcommand{\arraystretch}{1.15}
\begin{tabular}{@{}>{\RaggedRight\arraybackslash}p{0.27\linewidth}>{\RaggedRight\arraybackslash}p{0.66\linewidth}@{}}
\toprule
\textbf{Campo o control} & \textbf{Significado y criterio} \\
\midrule
Contrato preparado correctamente & Confirma el registro previo; revise huellas y selección. \\
\bottomrule
\end{tabular}
\endgroup
}
\Needspace{\dimexpr\ht0+\dp0+48pt\relax}
\endgroup
\subsection{Campos y controles}
\begingroup
\fontsize{10}{12.5}\selectfont
\setlength{\tabcolsep}{5pt}
\renewcommand{\arraystretch}{1.15}
\begin{longtable}{@{}>{\RaggedRight\arraybackslash}p{0.27\linewidth}>{\RaggedRight\arraybackslash}p{0.66\linewidth}@{}}
\toprule
\textbf{Campo o control} & \textbf{Significado y criterio} \\
\midrule
\endfirsthead
\toprule
\textbf{Campo o control} & \textbf{Significado y criterio} \\
\midrule
\endhead
\midrule
\multicolumn{2}{r}{\footnotesize Continúa en la página siguiente}\\
\endfoot
\bottomrule
\endlastfoot
Contrato preparado correctamente & Confirma el registro previo; revise huellas y selección. \\
\addlinespace[3pt]
Materializar y cargar datamart & Inicia extracción, transformación, creación y carga según el contrato validado. \\
\addlinespace[3pt]
Ver trazabilidad técnica & Constructor, huellas, responsable y fecha. \\
\addlinespace[3pt]
Expedientes recientes & Permite recuperar la ejecución para ver qué ocurrió, sin iniciar otra. \\
\addlinespace[3pt]
\end{longtable}
\endgroup

\subsection{Paso a paso}
\begin{enumerate}
\item Lea el comprobante y confirme que sigue siendo el expediente esperado.\par\nopagebreak[4]\smallskip{\small\textbf{Compruebe:} No carga un modelo de otra fuente o versión.}
\item Pulse Materializar y cargar datamart una vez y mantenga la pantalla abierta.\par\nopagebreak[4]\smallskip{\small\textbf{Compruebe:} Estado de procesamiento; no cambia la fuente a modo escritura.}
\item Espere el resultado.\par\nopagebreak[4]\smallskip{\small\textbf{Compruebe:} Avanza al paso 5 con validación, revisión pendiente o fallo.}
\item Si se interrumpe el navegador, recupere sesión y abra el expediente antes de reintentar.\par\nopagebreak[4]\smallskip{\small\textbf{Compruebe:} Reconoce si terminó, sigue preparada o falló; no crea cargas duplicadas por incertidumbre.}
\item Registre identificadores y estado final.\par\nopagebreak[4]\smallskip{\small\textbf{Compruebe:} Evidencia de una ejecución nueva o de una consulta histórica, claramente distinguidas.}
\end{enumerate}
\subsection{Qué hace la IA}
El diseño puede proceder de IA, pero la lectura y carga se realizan con el materializador y operaciones permitidas. La IA no recibe autorización para ejecutar instrucciones SQL libres.

\subsection{Qué comprueba la plataforma}
Contrato y fuente al ejecutar, trazabilidad, claves, relaciones exactas, tipos y condiciones de carga. El destino conserva el aislamiento por ejecución.

\subsection{Qué decide la persona}
La persona inicia la materialización; si hay fallo investiga y prepara una corrección en lugar de alterar a mano el expediente.

\paragraph{Resultado esperado} Expediente de carga con métricas o fallo controlado y resultados anteriores conservados.

\Needspace{8\baselineskip}
\begingroup
\setbox0=\vbox{\hsize=\linewidth
{\normalfont\normalsize\sffamily\bfseries\normalsize\color{Teal}\hyphenpenalty=10000\exhyphenpenalty=10000 Errores y recuperación\par}

\par
\begingroup
\fontsize{10}{12.5}\selectfont
\setlength{\tabcolsep}{5pt}
\renewcommand{\arraystretch}{1.15}
\begin{tabular}{@{}>{\RaggedRight\arraybackslash}p{0.27\linewidth}>{\RaggedRight\arraybackslash}p{0.66\linewidth}@{}}
\toprule
\textbf{Situación} & \textbf{Qué hacer} \\
\midrule
Error de conexión o sesión vencida. & Vuelva a entrar y revise Expedientes recientes antes de reenviar. El servidor puede haber terminado aunque el navegador perdiera la respuesta. \\
\bottomrule
\end{tabular}
\endgroup
}
\Needspace{\dimexpr\ht0+\dp0+48pt\relax}
\endgroup
\subsection{Errores y recuperación}
\begingroup
\fontsize{10}{12.5}\selectfont
\setlength{\tabcolsep}{5pt}
\renewcommand{\arraystretch}{1.15}
\begin{longtable}{@{}>{\RaggedRight\arraybackslash}p{0.27\linewidth}>{\RaggedRight\arraybackslash}p{0.66\linewidth}@{}}
\toprule
\textbf{Situación} & \textbf{Qué hacer} \\
\midrule
\endfirsthead
\toprule
\textbf{Situación} & \textbf{Qué hacer} \\
\midrule
\endhead
\midrule
\multicolumn{2}{r}{\footnotesize Continúa en la página siguiente}\\
\endfoot
\bottomrule
\endlastfoot
Error de conexión o sesión vencida. & Vuelva a entrar y revise Expedientes recientes antes de reenviar. El servidor puede haber terminado aunque el navegador perdiera la respuesta. \\
\addlinespace[3pt]
Fallida y mensaje de reversión. & Revise conexión y contrato con el administrador; si cambió la estructura capture metadatos y genere la versión correcta. Conserve el estado fallido hasta resolver la causa y comprobar la nueva ejecución. \\
\addlinespace[3pt]
La etapa dice Materialización ya ejecutada. & Es consulta. Pase a Validar resultados; esa pantalla no vuelve a cargar por recorrer los pasos. \\
\addlinespace[3pt]
\end{longtable}
\endgroup

\subsection{Ejercicio}
Simule sólo en la bitácora la decisión ante una pantalla que no responde: ¿qué revisa antes de volver a pulsar?

\textbf{Respuesta orientativa.} Estado del expediente existente, fuente, propuesta y sesión. No ejecutar de nuevo sin conocer si la carga terminó.

\subsection{Compruebe lo que aprendió}
\paragraph{¿Se reemplaza AdventureWorks por el datamart?} No. Se usa como fuente de lectura y el destino analítico queda separado.

\paragraph{¿Abrir un paso histórico repite la lectura de origen?} No. Permite consultar evidencia de la ejecución ya realizada.

\pantalla{recursos/pantallas/34_etl_carga_conservada.jpg}{Materialización histórica de sólo lectura}{Etapa de materialización histórica en sólo lectura. Una ejecución nueva requiere la confirmación explícita de su propio expediente. Captura de referencia del 18 de noviembre de 2026; revise las instrucciones actualizadas de esta unidad.}

\IfFileExists{capturas/ETL04.tex}{\input{capturas/ETL04.tex}}{}
\begingroup
\setlength{\parskip}{4pt plus 1pt minus 1pt}
\setlist[enumerate]{itemsep=3pt,topsep=4pt}
\titlespacing*{\paragraph}{0pt}{4pt}{0.5em}
\section{Materialización, paso 5: revisar conciliación y calidad}
\label{unidad:ETL05}
\textbf{Para quién:} Analista de BI, Gerente comercial para lectura autorizada\par
\textbf{Al terminar podrá:} Comprobar cantidades y valores y distinguir una conciliación numérica de una interpretación comercial correcta.

Una carga no se da por buena sólo por mostrar un mensaje verde. Lea Filas de origen, Filas cargadas, Diferencia, Tablas creadas y las métricas de cada tabla. Compare también indicadores y procedencia.

La conciliación de la ejecución contrasta filas y totales de medidas de origen/destino. Un resultado de diferencia cero demuestra consistencia de ese alcance; no demuestra por sí solo que se eligió el cliente, fecha o territorio comercial correctos.

\subsection{Comprobar la conciliación con la campaña vigente}
La conciliación compara cantidades equivalentes bajo el mismo alcance. Una tabla con una fila por línea no tiene el mismo conteo que una tabla con una fila por pedido. Sumar importes con uniones que duplican líneas puede dar una cifra convincente y equivocada; por eso se revisan grano, relaciones y receta antes de interpretar el panel.

\Needspace{8\baselineskip}
\begingroup
\setbox0=\vbox{\hsize=\linewidth
{\normalfont\normalsize\sffamily\bfseries\normalsize\color{Teal}\hyphenpenalty=10000\exhyphenpenalty=10000 Casos comprobados del 25 de noviembre de 2026\par}

\par
\begingroup
\fontsize{10}{12.5}\selectfont
\setlength{\tabcolsep}{5pt}
\renewcommand{\arraystretch}{1.15}
\begin{tabular}{@{}>{\RaggedRight\arraybackslash}p{0.22\linewidth}>{\RaggedRight\arraybackslash}p{0.3\linewidth}>{\RaggedRight\arraybackslash}p{0.39\linewidth}@{}}
\toprule
\textbf{Fuente y expediente} & \textbf{Alcance total de referencia} & \textbf{Filtro de control} \\
\midrule
AdventureWorks: propuesta 136; ejecución 24 & 121.317 líneas; 274.914 unidades; 31.465 pedidos distintos; importe 109.846.381,399888 USD. & Año 2013 y territorio Northwest: 7.757 líneas; 17.452 unidades; 2.053 pedidos; 6.015.174,735836 USD. \\
\bottomrule
\end{tabular}
\endgroup
}
\Needspace{\dimexpr\ht0+\dp0+48pt\relax}
\endgroup
\subsection{Casos comprobados del 25 de noviembre de 2026}
\begingroup
\fontsize{10}{12.5}\selectfont
\setlength{\tabcolsep}{5pt}
\renewcommand{\arraystretch}{1.15}
\begin{longtable}{@{}>{\RaggedRight\arraybackslash}p{0.22\linewidth}>{\RaggedRight\arraybackslash}p{0.3\linewidth}>{\RaggedRight\arraybackslash}p{0.39\linewidth}@{}}
\toprule
\textbf{Fuente y expediente} & \textbf{Alcance total de referencia} & \textbf{Filtro de control} \\
\midrule
\endfirsthead
\toprule
\textbf{Fuente y expediente} & \textbf{Alcance total de referencia} & \textbf{Filtro de control} \\
\midrule
\endhead
\midrule
\multicolumn{3}{r}{\footnotesize Continúa en la página siguiente}\\
\endfoot
\bottomrule
\endlastfoot
AdventureWorks: propuesta 136; ejecución 24 & 121.317 líneas; 274.914 unidades; 31.465 pedidos distintos; importe 109.846.381,399888 USD. & Año 2013 y territorio Northwest: 7.757 líneas; 17.452 unidades; 2.053 pedidos; 6.015.174,735836 USD. \\
\addlinespace[3pt]
WideWorldImporters: propuesta 135; ejecución 23 & 228.265 líneas; 8.950.628 unidades; 70.510 facturas distintas; importe 198.043.439,45 en moneda de origen. & Año 2015 y provincia California: 3.771 líneas; 142.247 unidades; 1.179 facturas; 3.191.769,11 en moneda de origen. \\
\addlinespace[3pt]
\end{longtable}
\endgroup

Estos valores proceden de consultas agregadas independientes al origen, contrastadas con los expedientes nuevos. Son controles de las copias de demostración revisadas, no parámetros incrustados en la plataforma. En otra copia, con otros filtros o un contrato distinto pueden cambiar.

Northwest se obtiene del territorio del encabezado del pedido. California se obtiene del cliente comprador y su ciudad de entrega registrada en el maestro. Una ubicación actual del cliente no demuestra por sí sola una dirección histórica de entrega. No atribuya al indicador una precisión temporal que la relación no conserva.

\subsection{Conceptos para decidir con criterio}
\paragraph{Duplicados controlados en dimensión} Repeticiones agrupadas al conformar una dimensión. Muchas facturas del mismo día pueden producir una sola fila de fecha sin eliminar líneas del hecho.

\paragraph{Diferencia cero} Coinciden los controles de ese universo. Hay que comprobar además evento, rutas, unidades y criterios comerciales.

\paragraph{Indicador conciliado} Valor calculado a partir de entradas y recetas del expediente; un resultado no calculable debe revisarse por separado.

\paragraph{Revisión pendiente} Puede significar diferencia de carga o interpretación semántica pendiente. Lea el motivo concreto, no sólo el color.

\Needspace{8\baselineskip}
\begingroup
\setbox0=\vbox{\hsize=\linewidth
{\normalfont\normalsize\sffamily\bfseries\normalsize\color{Teal}\hyphenpenalty=10000\exhyphenpenalty=10000 Campos y controles\par}

\par
\begingroup
\fontsize{10}{12.5}\selectfont
\setlength{\tabcolsep}{5pt}
\renewcommand{\arraystretch}{1.15}
\begin{tabular}{@{}>{\RaggedRight\arraybackslash}p{0.27\linewidth}>{\RaggedRight\arraybackslash}p{0.66\linewidth}@{}}
\toprule
\textbf{Campo o control} & \textbf{Significado y criterio} \\
\midrule
Conciliación OLTP-datamart & Comprobante de filas, diferencias y medidas de la ejecución. \\
\bottomrule
\end{tabular}
\endgroup
}
\Needspace{\dimexpr\ht0+\dp0+48pt\relax}
\endgroup
\subsection{Campos y controles}
\begingroup
\fontsize{10}{12.5}\selectfont
\setlength{\tabcolsep}{5pt}
\renewcommand{\arraystretch}{1.15}
\begin{longtable}{@{}>{\RaggedRight\arraybackslash}p{0.27\linewidth}>{\RaggedRight\arraybackslash}p{0.66\linewidth}@{}}
\toprule
\textbf{Campo o control} & \textbf{Significado y criterio} \\
\midrule
\endfirsthead
\toprule
\textbf{Campo o control} & \textbf{Significado y criterio} \\
\midrule
\endhead
\midrule
\multicolumn{2}{r}{\footnotesize Continúa en la página siguiente}\\
\endfoot
\bottomrule
\endlastfoot
Conciliación OLTP-datamart & Comprobante de filas, diferencias y medidas de la ejecución. \\
\addlinespace[3pt]
Calidad por tabla & Filas de origen/carga y duplicados controlados; distinga hecho y dimensión. \\
\addlinespace[3pt]
Indicadores calculados & Valores, unidad y estado. Algunos se muestran abreviados; el valor exacto y los reportes permiten contrastar decimales. \\
\addlinespace[3pt]
Interpretación semántica & Estado y revisión de candidatos de etiquetas; se explica en la siguiente unidad. \\
\addlinespace[3pt]
Comprobar divisa sin repetir el ETL & Acción aparte si falta evidencia del código monetario. \\
\addlinespace[3pt]
\end{longtable}
\endgroup

\subsection{Paso a paso}
\begin{enumerate}
\item Revise que origen y destino correspondan al mismo expediente y grano.\par\nopagebreak[4]\smallskip{\small\textbf{Compruebe:} No compara el filtro de un tablero con un total global de otro caso.}
\item Compare filas y diferencia.\par\nopagebreak[4]\smallskip{\small\textbf{Compruebe:} El hecho conserva el número esperado sin rechazo o multiplicación indebida.}
\item Revise por tabla las agrupaciones de valores repetidos en las dimensiones.\par\nopagebreak[4]\smallskip{\small\textbf{Compruebe:} Entiende por qué la fecha puede tener menos filas que documentos sin perder hechos.}
\item Compare indicadores, unidades y razones con su fórmula.\par\nopagebreak[4]\smallskip{\small\textbf{Compruebe:} No hay denominador cambiado o unidad inventada.}
\item Revise rutas de cliente/territorio además de números.\par\nopagebreak[4]\smallskip{\small\textbf{Compruebe:} Se conserva el rol comercial solicitado.}
\item Si la conciliación y publicación cumplen, abra Analítica de ventas.\par\nopagebreak[4]\smallskip{\small\textbf{Compruebe:} Tablero identificado con fuente, propuesta y ejecución correctas.}
\end{enumerate}
\subsection{Qué hace la IA}
Puede apoyar etiquetas y explicar después resultados; el contraste numérico no se obtiene por opinión del modelo.

\subsection{Qué comprueba la plataforma}
En el materializador actual la conciliación requiere cero filas rechazadas, mismo conteo origen/destino y diferencias de totales de medidas iguales a cero. La condición numérica no reemplaza la revisión semántica del analista.

\subsection{Qué decide la persona}
La persona revisa métricas, procedencia y alcance antes de usar los resultados para una lectura comercial.

\paragraph{Resultado esperado} Expediente comprendido y evidencia concreta para aceptar o investigar.

\Needspace{8\baselineskip}
\begingroup
\setbox0=\vbox{\hsize=\linewidth
{\normalfont\normalsize\sffamily\bfseries\normalsize\color{Teal}\hyphenpenalty=10000\exhyphenpenalty=10000 Errores y recuperación\par}

\par
\begingroup
\fontsize{10}{12.5}\selectfont
\setlength{\tabcolsep}{5pt}
\renewcommand{\arraystretch}{1.15}
\begin{tabular}{@{}>{\RaggedRight\arraybackslash}p{0.27\linewidth}>{\RaggedRight\arraybackslash}p{0.66\linewidth}@{}}
\toprule
\textbf{Situación} & \textbf{Qué hacer} \\
\midrule
Hay diferencia de filas o importes. & No presente el tablero como verificado sin resolverla. Registre magnitud y alcance, revise claves/uniones/entradas y corrija el contrato. \\
\bottomrule
\end{tabular}
\endgroup
}
\Needspace{\dimexpr\ht0+\dp0+48pt\relax}
\endgroup
\subsection{Errores y recuperación}
\begingroup
\fontsize{10}{12.5}\selectfont
\setlength{\tabcolsep}{5pt}
\renewcommand{\arraystretch}{1.15}
\begin{longtable}{@{}>{\RaggedRight\arraybackslash}p{0.27\linewidth}>{\RaggedRight\arraybackslash}p{0.66\linewidth}@{}}
\toprule
\textbf{Situación} & \textbf{Qué hacer} \\
\midrule
\endfirsthead
\toprule
\textbf{Situación} & \textbf{Qué hacer} \\
\midrule
\endhead
\midrule
\multicolumn{2}{r}{\footnotesize Continúa en la página siguiente}\\
\endfoot
\bottomrule
\endlastfoot
Hay diferencia de filas o importes. & No presente el tablero como verificado sin resolverla. Registre magnitud y alcance, revise claves/uniones/entradas y corrija el contrato. \\
\addlinespace[3pt]
Fecha muestra miles de duplicados controlados. & Compruebe si son días repetidos en documentos; en el caso WWI 22 se conservaron 1.069 días desde 70.510 documentos sin perder 228.265 líneas. \\
\addlinespace[3pt]
Los totales coinciden pero cambia California entre versiones. & Compare roles de comprador/facturación y ciudad de entrega/postal. Distinta clasificación puede explicar el cambio sin mezcla entre bases. \\
\addlinespace[3pt]
\end{longtable}
\endgroup

\subsection{Ejercicio}
Conserve 100 líneas de factura sobre 10 días. ¿Es obligatorio tener 100 filas en la dimensión fecha?

\textbf{Respuesta orientativa.} No. Puede haber 10 fechas únicas y 100 filas en el hecho; lo importante es que las claves y los controles expliquen esa conformación.

\subsection{Compruebe lo que aprendió}
\paragraph{¿Diferencia cero prueba que el territorio era históricamente el del cliente?} No. Haría falta una historia de vigencias; el prototipo lee los atributos del estado disponible durante el ETL.

\paragraph{¿Revisión pendiente siempre significa carga fallida?} No. Puede haber conciliación correcta con etiquetas aún pendientes; lea el detalle.

\endgroup
\pantalla{recursos/pantallas/29_conciliacion_AW.jpg}{Conciliación del expediente AdventureWorks}{Expediente 21 de AdventureWorks: 121.317 filas en origen y destino, diferencia cero y cinco tablas. La ejecución ocurrió el 6 de noviembre. Captura de referencia del 18 de noviembre de 2026; revise las instrucciones actualizadas de esta unidad.}
\pantalla{recursos/pantallas/35_conciliacion_WWI.jpg}{Conciliación de WideWorldImporters y divisa sin confirmar}{Expediente 22 de WideWorldImporters: 228.265 filas conciliadas. La divisa continúa sin comprobar y no se atribuye USD. Captura de referencia del 18 de noviembre de 2026; revise las instrucciones actualizadas de esta unidad.}

\IfFileExists{capturas/ETL05.tex}{\input{capturas/ETL05.tex}}{}
\begingroup
\setlength{\parskip}{4pt plus 1pt minus 1pt}
\setlist[enumerate]{itemsep=3pt,topsep=4pt}
\titlespacing*{\paragraph}{0pt}{4pt}{0.5em}
\section{Revisar etiquetas en español y comprobar la moneda}
\label{unidad:ETL06}
\textbf{Para quién:} Analista de BI\par
\textbf{Al terminar podrá:} Publicar sólo interpretaciones revisadas y conservar el significado original y la divisa demostrable.

La interpretación semántica añade etiquetas españolas a categorías aptas; conserva los valores originales. No traduce identificadores, nombres de personas, direcciones, texto libre ni categorías de alta cardinalidad.

La moneda es un control diferente. Un importe conciliado no autoriza a escribir USD. Si la fuente no demuestra un código único adecuado, el resultado correcto es moneda de origen. No confunda el país del cliente con la divisa del importe.

\subsection{Entender moneda y razones sin inventar conversiones}
Cuando la fuente demuestra USD u otra divisa admitida, una etiqueta monetaria como moneda de origen por transacción se presenta con ese código, conservando el denominador. Si la fuente no demuestra una divisa, sigue diciendo moneda de origen. Una etiqueta USD o EUR sugerida por el modelo no es prueba de moneda.

La corrección es de presentación: no convierte valores, no aplica tipos de cambio y no modifica el contrato guardado. Las unidades físicas, por ejemplo piezas o kilogramos, no se transforman en moneda. Un promedio por línea y un promedio por pedido siguen siendo indicadores diferentes aunque ambos se expresen en USD.

\subsection{Conceptos para decidir con criterio}
\paragraph{Mapeo} Correspondencia entre categoría original y etiqueta española revisada.

\paragraph{Publicación supervisada} Las etiquetas sólo se publican después de la revisión individual y la confirmación humana.

\paragraph{Divisa comprobada} Código respaldado por evidencia relacional de fuente/base; no elegido por el modelo por intuición.

\Needspace{8\baselineskip}
\begingroup
\setbox0=\vbox{\hsize=\linewidth
{\normalfont\normalsize\sffamily\bfseries\normalsize\color{Teal}\hyphenpenalty=10000\exhyphenpenalty=10000 Campos y controles\par}

\par
\begingroup
\fontsize{10}{12.5}\selectfont
\setlength{\tabcolsep}{5pt}
\renewcommand{\arraystretch}{1.15}
\begin{tabular}{@{}>{\RaggedRight\arraybackslash}p{0.27\linewidth}>{\RaggedRight\arraybackslash}p{0.66\linewidth}@{}}
\toprule
\textbf{Campo o control} & \textbf{Significado y criterio} \\
\midrule
Publicar & Casilla por mapeo para conservarlo o excluirlo. \\
\bottomrule
\end{tabular}
\endgroup
}
\Needspace{\dimexpr\ht0+\dp0+48pt\relax}
\endgroup
\subsection{Campos y controles}
\begingroup
\fontsize{10}{12.5}\selectfont
\setlength{\tabcolsep}{5pt}
\renewcommand{\arraystretch}{1.15}
\begin{longtable}{@{}>{\RaggedRight\arraybackslash}p{0.27\linewidth}>{\RaggedRight\arraybackslash}p{0.66\linewidth}@{}}
\toprule
\textbf{Campo o control} & \textbf{Significado y criterio} \\
\midrule
\endfirsthead
\toprule
\textbf{Campo o control} & \textbf{Significado y criterio} \\
\midrule
\endhead
\midrule
\multicolumn{2}{r}{\footnotesize Continúa en la página siguiente}\\
\endfoot
\bottomrule
\endlastfoot
Publicar & Casilla por mapeo para conservarlo o excluirlo. \\
\addlinespace[3pt]
Etiqueta en español & Texto editable, hasta 100 caracteres, que debe preservar el significado. \\
\addlinespace[3pt]
Comentario de revisión & Entre 10 y 500 caracteres sobre la decisión tomada. \\
\addlinespace[3pt]
Confirmo que revisé los mapeos que se publicarán & Confirmación distinta de la aprobación del diseño y de la preparación del ETL. \\
\addlinespace[3pt]
Publicar etiquetas revisadas & Publica sólo los mapeos aceptados como columnas adicionales auditables. \\
\addlinespace[3pt]
Reintentar interpretación sin repetir el ETL & Recupera una interpretación pendiente sin recargar el hecho. \\
\addlinespace[3pt]
Comprobar divisa sin repetir el ETL & Consulta evidencia de moneda y actualiza su contexto sin otra materialización. \\
\addlinespace[3pt]
\end{longtable}
\endgroup

\subsection{Paso a paso}
\begin{enumerate}
\item Lea el estado de Interpretación semántica.\par\nopagebreak[4]\smallskip{\small\textbf{Compruebe:} Reconoce si no hay candidatos, están pendientes, requieren revisión o ya fueron publicados.}
\item Compare cada original con la etiqueta, corrija lo necesario y excluya códigos que no convenga traducir.\par\nopagebreak[4]\smallskip{\small\textbf{Compruebe:} No cambia el valor original ni acepta una etiqueta por rutina.}
\item Escriba comentario, marque confirmación y publique.\par\nopagebreak[4]\smallskip{\small\textbf{Compruebe:} Recibo con responsable, fecha y mapeos revisados.}
\item Si está pendiente por proveedor, revise la configuración y use el reintento semántico.\par\nopagebreak[4]\smallskip{\small\textbf{Compruebe:} No se repite el ETL para corregir sólo etiquetas.}
\item Revise la divisa y compruébela si la acción está disponible.\par\nopagebreak[4]\smallskip{\small\textbf{Compruebe:} Código verificado o moneda de origen con el límite conservado.}
\end{enumerate}
\subsection{Qué hace la IA}
Propone etiquetas para categorías seguras dentro de la política implementada. No debería convertir códigos, nombres propios o información ausente en traducciones inventadas.

\subsection{Qué comprueba la plataforma}
Aptitud de categorías, formato de mapeos, conservación del original y publicación autorizada. La comprobación de moneda usa evidencia de la fuente, no una respuesta libre del LLM.

\subsection{Qué decide la persona}
El analista publica, corrige o excluye cada etiqueta y confirma la revisión. No elige una divisa sin evidencia para mejorar la apariencia.

\paragraph{Resultado esperado} Etiquetas comprensibles con original conservado y unidad monetaria honesta.

\Needspace{8\baselineskip}
\begingroup
\setbox0=\vbox{\hsize=\linewidth
{\normalfont\normalsize\sffamily\bfseries\normalsize\color{Teal}\hyphenpenalty=10000\exhyphenpenalty=10000 Errores y recuperación\par}

\par
\begingroup
\fontsize{10}{12.5}\selectfont
\setlength{\tabcolsep}{5pt}
\renewcommand{\arraystretch}{1.15}
\begin{tabular}{@{}>{\RaggedRight\arraybackslash}p{0.27\linewidth}>{\RaggedRight\arraybackslash}p{0.66\linewidth}@{}}
\toprule
\textbf{Situación} & \textbf{Qué hacer} \\
\midrule
No aparece ninguna etiqueta para aprobar. & Si no hay candidatos aptos no existe una tarea de publicación. No hay que fabricar un mapeo. \\
\bottomrule
\end{tabular}
\endgroup
}
\Needspace{\dimexpr\ht0+\dp0+48pt\relax}
\endgroup
\subsection{Errores y recuperación}
\begingroup
\fontsize{10}{12.5}\selectfont
\setlength{\tabcolsep}{5pt}
\renewcommand{\arraystretch}{1.15}
\begin{longtable}{@{}>{\RaggedRight\arraybackslash}p{0.27\linewidth}>{\RaggedRight\arraybackslash}p{0.66\linewidth}@{}}
\toprule
\textbf{Situación} & \textbf{Qué hacer} \\
\midrule
\endfirsthead
\toprule
\textbf{Situación} & \textbf{Qué hacer} \\
\midrule
\endhead
\midrule
\multicolumn{2}{r}{\footnotesize Continúa en la página siguiente}\\
\endfoot
\bottomrule
\endlastfoot
No aparece ninguna etiqueta para aprobar. & Si no hay candidatos aptos no existe una tarea de publicación. No hay que fabricar un mapeo. \\
\addlinespace[3pt]
El proveedor falla después de una carga conciliada. & Reintente sólo interpretación; antes compruebe el expediente para conservar la carga existente. \\
\addlinespace[3pt]
WWI no demuestra USD. & Conserve moneda de origen. Cambiar el rótulo manualmente no convierte el dato en dólares. \\
\addlinespace[3pt]
\end{longtable}
\endgroup

\subsection{Ejercicio}
En el caso AW 21 se propusieron Europe, North America, Pacific y códigos AU, CA, DE, FR, GB, US. Explique la revisión documentada.

\textbf{Respuesta orientativa.} Se publicaron Europa, América del Norte y Pacífico; los seis códigos se excluyeron de traducción. Se conservaron originales.

\subsection{Compruebe lo que aprendió}
\paragraph{¿Publicar etiquetas altera los importes?} No. Añade interpretación categórica auditada; no cambia las medidas originales.

\paragraph{¿Siempre debe resolverse un código monetario?} No. La ausencia de evidencia suficiente debe conservarse como límite.

\endgroup
\pantalla{recursos/pantallas/30_semantica_publicada.jpg}{Etiquetas semánticas publicadas sin reemplazar originales}{Etiquetas revisadas y publicadas como columnas adicionales. Los valores originales se conservan para trazabilidad. Captura de referencia del 18 de noviembre de 2026; revise las instrucciones actualizadas de esta unidad.}

\IfFileExists{capturas/ETL06.tex}{\input{capturas/ETL06.tex}}{}

\chapter{Analizar, conversar y exportar con contexto}
\section{Elegir el datamart y filtrar el análisis}
\label{unidad:AN01}
\textbf{Para quién:} Analista de BI, Gerente comercial\par
\textbf{Al terminar podrá:} Leer los identificadores y usar filtros sin confundir vistas, bases ni universos.

Analítica de ventas consulta los datos físicos de una ejecución conciliada, no vuelve a extraer cada vez la base transaccional. Compruebe la fuente y los identificadores del encabezado antes de interpretar cifras.

Vista ejecutiva resume la información para gerencia. Vista analítica añade controles y tablas de respaldo. Cambiar de vista no otorga un rol nuevo ni añade datos que no se hayan materializado.

\subsection{Ejercicio actual: recuperar filtros de cada fuente}
\begin{enumerate}
\item Abra AdventureWorks, seleccione la ejecución 24 si existe, el año 2013 y Northwest.
\item Compruebe 17.452 unidades y 2.053 pedidos distintos. Lea el importe exacto en el respaldo y compare su abreviación en las tarjetas.
\item Cambie a WideWorldImporters y seleccione la ejecución 23, año 2015 y California, si están disponibles. Compruebe 142.247 unidades y 1.179 facturas distintas.
\item Regrese a AdventureWorks. Deben reaparecer sus filtros confirmados, sin importar filtros ni chat de WideWorldImporters.
\item Entre al respaldo y vuelva con Atrás. Si consulta una ejecución histórica, anuncie su número y su fecha. No espere ver los mismos identificadores en un portable limpio.
\end{enumerate}

\pantallaactual{recursos/pantallas/aw24-filtros-restaurados.jpg}{AdventureWorks: ejecución 24 y filtros restaurados}{Captura real del 25 de noviembre de 2026. Se observa 2013, Northwest, 17.452 unidades y 2.053 pedidos. La tarjeta abrevia el importe; la tabla o exportación permite leer el valor exacto.}
\subsection{Conceptos para decidir con criterio}
\paragraph{Datamart analizado} Ejecución concreta cuyos datos físicos se consultan; puede haber varios expedientes de la misma fuente.

\paragraph{Filtro} Delimita las filas que participan en el cálculo. Año y territorio se combinan en el mismo alcance.

\paragraph{Indicador de los gráficos} Métrica que explica las visualizaciones. No confundirla con todos los KPI del encabezado.

\paragraph{Última actualización} Fecha de finalización del expediente o de su creación si no existe otra; no equivale a una recarga automática al abrir el panel.

\Needspace{8\baselineskip}
\begingroup
\setbox0=\vbox{\hsize=\linewidth
{\normalfont\normalsize\sffamily\bfseries\normalsize\color{Teal}\hyphenpenalty=10000\exhyphenpenalty=10000 Campos y controles\par}

\par
\begingroup
\fontsize{10}{12.5}\selectfont
\setlength{\tabcolsep}{5pt}
\renewcommand{\arraystretch}{1.15}
\begin{tabular}{@{}>{\RaggedRight\arraybackslash}p{0.27\linewidth}>{\RaggedRight\arraybackslash}p{0.66\linewidth}@{}}
\toprule
\textbf{Campo o control} & \textbf{Significado y criterio} \\
\midrule
Fuente en contexto & Selecciona la base de trabajo. Cambiarla limpia ejecución, filtros y conversación y vuelve a pedir resultados de esa fuente. \\
\bottomrule
\end{tabular}
\endgroup
}
\Needspace{\dimexpr\ht0+\dp0+48pt\relax}
\endgroup
\subsection{Campos y controles}
\begingroup
\fontsize{10}{12.5}\selectfont
\setlength{\tabcolsep}{5pt}
\renewcommand{\arraystretch}{1.15}
\begin{longtable}{@{}>{\RaggedRight\arraybackslash}p{0.27\linewidth}>{\RaggedRight\arraybackslash}p{0.66\linewidth}@{}}
\toprule
\textbf{Campo o control} & \textbf{Significado y criterio} \\
\midrule
\endfirsthead
\toprule
\textbf{Campo o control} & \textbf{Significado y criterio} \\
\midrule
\endhead
\midrule
\multicolumn{2}{r}{\footnotesize Continúa en la página siguiente}\\
\endfoot
\bottomrule
\endlastfoot
Fuente en contexto & Selecciona la base de trabajo. Cambiarla limpia ejecución, filtros y conversación y vuelve a pedir resultados de esa fuente. \\
\addlinespace[3pt]
Datamart analizado & Combobox cuando hay más de una ejecución con datos físicos consultables. Si sólo hay una se muestra Datamart activo. \\
\addlinespace[3pt]
Vista ejecutiva / Vista analítica & Presentación según necesidad de lectura. En la ejecutiva se omite el visual de clientes; en la analítica se muestran controles/datos adicionales. \\
\addlinespace[3pt]
Indicador de los gráficos & En vista analítica lista agregaciones que admiten visualización segura. Las recetas derivadas pueden existir como KPI sin aparecer en ese selector. \\
\addlinespace[3pt]
Año / Territorio & Opciones obtenidas del datamart disponible; no inventan años futuros. \\
\addlinespace[3pt]
Restablecer filtros & Retira año y territorio; no cambia la fuente ni ejecuta otro ETL. \\
\addlinespace[3pt]
Consultar respaldo / Ver expediente técnico & Abre la evidencia de la ejecución que respalda el panel. \\
\addlinespace[3pt]
\end{longtable}
\endgroup

\subsection{Paso a paso}
\begin{enumerate}
\item Compruebe fuente, base, propuesta, ejecución, moneda y período en el encabezado.\par\nopagebreak[4]\smallskip{\small\textbf{Compruebe:} Contexto correcto antes de leer tarjetas.}
\item Seleccione el datamart deseado si hay varios.\par\nopagebreak[4]\smallskip{\small\textbf{Compruebe:} El panel confirma la ejecución seleccionada. Las selecciones se conservan separadas por fuente y sólo se recuperan cuando esa ejecución sigue disponible.}
\item Abra Vista analítica y elija el indicador de los gráficos.\par\nopagebreak[4]\smallskip{\small\textbf{Compruebe:} Las gráficas representan la métrica elegida.}
\item Seleccione año y territorio y espere Actualizando resultados.\par\nopagebreak[4]\smallskip{\small\textbf{Compruebe:} KPI, gráficos y hallazgos se recalculan para ese alcance.}
\item Compare un valor exacto de tabla con una tarjeta abreviada.\par\nopagebreak[4]\smallskip{\small\textbf{Compruebe:} Comprende la diferencia de formato sin declarar un error de cálculo.}
\item Use Restablecer filtros y vuelva a seleccionar la otra fuente cuando termine.\par\nopagebreak[4]\smallskip{\small\textbf{Compruebe:} Restablecer limpia los filtros de la fuente actual. Al volver a la otra fuente se recupera su último contexto confirmado; no se cruza la conversación.}
\end{enumerate}
\subsection{Qué hace la IA}
No hace falta una generación LLM para consultar o filtrar el panel. Los valores se calculan mediante recetas y consultas controladas.

\subsection{Qué comprueba la plataforma}
Correspondencia de fuente y ejecución, disponibilidad física, recetas y filtros permitidos; el servidor comprueba procedencia además de los controles del navegador.

\subsection{Qué decide la persona}
La persona elige el universo e interpreta los resultados dentro del evento materializado.

\paragraph{Resultado esperado} Un panel cuyo contexto y alcance puedan explicarse antes de emitir una conclusión.

\Needspace{8\baselineskip}
\begingroup
\setbox0=\vbox{\hsize=\linewidth
{\normalfont\normalsize\sffamily\bfseries\normalsize\color{Teal}\hyphenpenalty=10000\exhyphenpenalty=10000 Errores y recuperación\par}

\par
\begingroup
\fontsize{10}{12.5}\selectfont
\setlength{\tabcolsep}{5pt}
\renewcommand{\arraystretch}{1.15}
\begin{tabular}{@{}>{\RaggedRight\arraybackslash}p{0.27\linewidth}>{\RaggedRight\arraybackslash}p{0.66\linewidth}@{}}
\toprule
\textbf{Situación} & \textbf{Qué hacer} \\
\midrule
Primero complete el expediente del datamart. & No hay ejecución conciliada con datos físicos disponibles. Vuelva a Datamart de ventas y revise preparación, carga y validación. \\
\bottomrule
\end{tabular}
\endgroup
}
\Needspace{\dimexpr\ht0+\dp0+48pt\relax}
\endgroup
\subsection{Errores y recuperación}
\begingroup
\fontsize{10}{12.5}\selectfont
\setlength{\tabcolsep}{5pt}
\renewcommand{\arraystretch}{1.15}
\begin{longtable}{@{}>{\RaggedRight\arraybackslash}p{0.27\linewidth}>{\RaggedRight\arraybackslash}p{0.66\linewidth}@{}}
\toprule
\textbf{Situación} & \textbf{Qué hacer} \\
\midrule
\endfirsthead
\toprule
\textbf{Situación} & \textbf{Qué hacer} \\
\midrule
\endhead
\midrule
\multicolumn{2}{r}{\footnotesize Continúa en la página siguiente}\\
\endfoot
\bottomrule
\endlastfoot
Primero complete el expediente del datamart. & No hay ejecución conciliada con datos físicos disponibles. Vuelva a Datamart de ventas y revise preparación, carga y validación. \\
\addlinespace[3pt]
Falta el combobox de datamart. & Puede ser correcto si sólo una ejecución tiene datos físicos; no se debe crear otra para obtener el selector. \\
\addlinespace[3pt]
Aparece un resultado anterior mientras carga el filtro. & Espere a que termine y confirme el contexto del encabezado antes de copiar cifras. Si persiste una fuente equivocada, registre incidencia. \\
\addlinespace[3pt]
El chat se vació tras cambiar año o fuente. & Es protección de contexto: la respuesta anterior no se reutiliza como si perteneciera a la nueva selección. \\
\addlinespace[3pt]
\end{longtable}
\endgroup

\subsection{Ejercicio}
Explique por qué los 109,8 millones globales de AW 21 no se comparan directamente con 6.015.174,74 del filtro 2013/Northwest.

\textbf{Respuesta orientativa.} Son universos diferentes. La comparación necesita la misma ejecución y filtros.

\subsection{Compruebe lo que aprendió}
\paragraph{¿Cambiar Año vuelve a cargar el origen?} No. Recalcula sobre el datamart.

\paragraph{¿Todas las recetas pueden graficarse por entidad?} No. El selector ofrece las agregaciones que admiten visualización agrupada segura en la implementación.

\pantalla{recursos/pantallas/36_analitica_WWI.jpg}{Panel ejecutivo y procedencia de WideWorldImporters}{Panel de WideWorldImporters sin filtros. El encabezado identifica propuesta 131 y ejecución 22. Captura de referencia del 18 de noviembre de 2026; revise las instrucciones actualizadas de esta unidad.}
\pantalla{recursos/pantallas/42_analitica_AW.jpg}{Panel ejecutivo AdventureWorks y procedencia}{AdventureWorks: propuesta 129, ejecución 21 y divisa USD comprobada. Los filtros visibles corresponden al contexto de esa captura histórica. Captura de referencia del 18 de noviembre de 2026; revise las instrucciones actualizadas de esta unidad.}
\pantalla{recursos/pantallas/37_WWI_2015_California.jpg}{Caso filtrado: 2015 y California}{Caso de referencia WideWorldImporters: año 2015 y California; 1.179 facturas y promedio 2.707,18 en moneda de origen. Captura de referencia del 18 de noviembre de 2026; revise las instrucciones actualizadas de esta unidad.}
\pantalla{recursos/pantallas/43_AW_2013_Northwest.jpg}{Caso filtrado: 2013 y Northwest}{Caso AdventureWorks: 2013 y Northwest; 17.452 unidades, 2.053 pedidos y promedio 2.929,94 por pedido. Captura de referencia del 18 de noviembre de 2026; revise las instrucciones actualizadas de esta unidad.}

\IfFileExists{capturas/AN01.tex}{\input{capturas/AN01.tex}}{}
\begingroup
\setlength{\parskip}{4pt plus 1pt minus 1pt}
\setlist[enumerate]{itemsep=3pt,topsep=4pt}
\titlespacing*{\paragraph}{0pt}{4pt}{0.5em}
\section{Leer gráficos y hallazgos con criterio comercial}
\label{unidad:AN02}
\textbf{Para quién:} Analista de BI, Gerente comercial\par
\textbf{Al terminar podrá:} Diferenciar tendencia, ranking y distribución y explicar sus denominadores y límites.

Una línea temporal compara períodos; las barras ordenan entidades; una distribución muestra participaciones. Antes de decir que un producto es el mejor, precise si lo ordenó por importe o unidades y qué filtros utilizó.

Los hallazgos del panel se construyen con reglas sobre los resultados visibles. Aunque la pantalla los llame Copiloto analítico, no todos sus textos proceden de una nueva llamada LLM. La conversación inferior sí utiliza el proveedor de IA.

\subsection{Conceptos para decidir con criterio}
\paragraph{Top o ranking} Lista limitada de entidades ordenadas por un indicador. El máximo de una línea no equivale a la mayor suma por producto.

\paragraph{Participación visible} Porcentaje respecto del total definido para esa visualización. En el panel los rankings limitados pueden usar la suma de lo mostrado; en consultas agregadas del copiloto se informa el universo filtrado completo.

\paragraph{Asociación} Patrón observado en datos. No prueba por sí solo causalidad, preferencia del cliente, demanda de mercado o eficacia de una campaña.

\paragraph{Detalle seleccionado} Valor y contexto de un punto, barra o categoría. Seleccionarlo no cambia filtros sin una acción adicional.

\Needspace{8\baselineskip}
\begingroup
\setbox0=\vbox{\hsize=\linewidth
{\normalfont\normalsize\sffamily\bfseries\normalsize\color{Teal}\hyphenpenalty=10000\exhyphenpenalty=10000 Campos y controles\par}

\par
\begingroup
\fontsize{10}{12.5}\selectfont
\setlength{\tabcolsep}{5pt}
\renewcommand{\arraystretch}{1.15}
\begin{tabular}{@{}>{\RaggedRight\arraybackslash}p{0.27\linewidth}>{\RaggedRight\arraybackslash}p{0.66\linewidth}@{}}
\toprule
\textbf{Campo o control} & \textbf{Significado y criterio} \\
\midrule
Punto, barra o categoría & Clic o teclado permiten ver valor exacto y participación del dato seleccionado. \\
\bottomrule
\end{tabular}
\endgroup
}
\Needspace{\dimexpr\ht0+\dp0+48pt\relax}
\endgroup
\subsection{Campos y controles}
\begingroup
\fontsize{10}{12.5}\selectfont
\setlength{\tabcolsep}{5pt}
\renewcommand{\arraystretch}{1.15}
\begin{longtable}{@{}>{\RaggedRight\arraybackslash}p{0.27\linewidth}>{\RaggedRight\arraybackslash}p{0.66\linewidth}@{}}
\toprule
\textbf{Campo o control} & \textbf{Significado y criterio} \\
\midrule
\endfirsthead
\toprule
\textbf{Campo o control} & \textbf{Significado y criterio} \\
\midrule
\endhead
\midrule
\multicolumn{2}{r}{\footnotesize Continúa en la página siguiente}\\
\endfoot
\bottomrule
\endlastfoot
Punto, barra o categoría & Clic o teclado permiten ver valor exacto y participación del dato seleccionado. \\
\addlinespace[3pt]
Filtrar tablero por territorio & Botón explícito en el detalle territorial. Abrir el detalle no aplica el filtro automáticamente. \\
\addlinespace[3pt]
Cerrar detalle & Cierra la selección puntual sin modificar el universo. \\
\addlinespace[3pt]
Tabla de respaldo en vista analítica & Permite contrastar cifras de la gráfica. \\
\addlinespace[3pt]
Cómo se generaron & Explica la procedencia y reglas de los hallazgos. \\
\addlinespace[3pt]
\end{longtable}
\endgroup

\subsection{Paso a paso}
\begin{enumerate}
\item Lea título, subtítulo, indicador, unidad y filtros de la gráfica.\par\nopagebreak[4]\smallskip{\small\textbf{Compruebe:} Identifica el universo que está comparando.}
\item Seleccione el dato que desee revisar.\par\nopagebreak[4]\smallskip{\small\textbf{Compruebe:} Detalle con valor exacto sin cambio inesperado del panel.}
\item Compare con la tabla de respaldo.\par\nopagebreak[4]\smallskip{\small\textbf{Compruebe:} La cifra representada tiene soporte verificable.}
\item Si desea restringir por territorio, pulse el botón de filtro explícito.\par\nopagebreak[4]\smallskip{\small\textbf{Compruebe:} Se recalcula el panel para ese alcance.}
\item Lea hallazgo, evidencia y Cómo se generaron.\par\nopagebreak[4]\smallskip{\small\textbf{Compruebe:} Distingue una descripción numérica de una conclusión sobre causas.}
\item Redacte una lectura y una comprobación comercial adicional.\par\nopagebreak[4]\smallskip{\small\textbf{Compruebe:} Por ejemplo revisar composición de productos antes de atribuir una caída a un proveedor.}
\end{enumerate}
\subsection{Qué hace la IA}
Los hallazgos automáticos del panel son determinísticos. La IA conversacional puede explicar la evidencia después, pero su prosa requiere revisión.

\subsection{Qué comprueba la plataforma}
Cálculo de agregados y visualizaciones sobre el datamart identificado. La gráfica no certifica hipótesis comerciales ajenas a los registros.

\subsection{Qué decide la persona}
La persona decide qué dato merece atención y qué comprobaciones adicionales necesita antes de actuar.

\paragraph{Resultado esperado} Lectura comercial con cifra, unidad, período y límite explícitos.

\Needspace{8\baselineskip}
\begingroup
\setbox0=\vbox{\hsize=\linewidth
{\normalfont\normalsize\sffamily\bfseries\normalsize\color{Teal}\hyphenpenalty=10000\exhyphenpenalty=10000 Errores y recuperación\par}

\par
\begingroup
\fontsize{10}{12.5}\selectfont
\setlength{\tabcolsep}{5pt}
\renewcommand{\arraystretch}{1.15}
\begin{tabular}{@{}>{\RaggedRight\arraybackslash}p{0.27\linewidth}>{\RaggedRight\arraybackslash}p{0.66\linewidth}@{}}
\toprule
\textbf{Situación} & \textbf{Qué hacer} \\
\midrule
El Top 8 no suma el total global. & Es un subconjunto; revise su denominador. No ajuste la gráfica para que aparente cubrir todo. \\
\bottomrule
\end{tabular}
\endgroup
}
\Needspace{\dimexpr\ht0+\dp0+48pt\relax}
\endgroup
\subsection{Errores y recuperación}
\begingroup
\fontsize{10}{12.5}\selectfont
\setlength{\tabcolsep}{5pt}
\renewcommand{\arraystretch}{1.15}
\begin{longtable}{@{}>{\RaggedRight\arraybackslash}p{0.27\linewidth}>{\RaggedRight\arraybackslash}p{0.66\linewidth}@{}}
\toprule
\textbf{Situación} & \textbf{Qué hacer} \\
\midrule
\endfirsthead
\toprule
\textbf{Situación} & \textbf{Qué hacer} \\
\midrule
\endhead
\midrule
\multicolumn{2}{r}{\footnotesize Continúa en la página siguiente}\\
\endfoot
\bottomrule
\endlastfoot
El Top 8 no suma el total global. & Es un subconjunto; revise su denominador. No ajuste la gráfica para que aparente cubrir todo. \\
\addlinespace[3pt]
Pedidos distintos por producto superan los pedidos globales al sumarlos. & Un pedido puede contener varios productos; el conteo distinto no es aditivo entre grupos. \\
\addlinespace[3pt]
La IA dice que un producto domina el mercado. & El ranking interno no demuestra cuota de mercado. Reformule como posición dentro del universo consultado. \\
\addlinespace[3pt]
\end{longtable}
\endgroup

\subsection{Ejercicio}
Un producto lidera por importe y otro por unidades. Escriba dos conclusiones válidas.

\textbf{Respuesta orientativa.} Liderazgo por importe y liderazgo por unidades son distintos; ninguno demuestra automáticamente mayor margen o preferencia del mercado.

\subsection{Compruebe lo que aprendió}
\paragraph{¿Un gráfico por provincia demuestra dónde ocurrió físicamente la entrega?} No, salvo evidencia específica. En WWI el caso usa provincia de la ciudad de entrega registrada en la ficha.

\paragraph{¿Abrir una barra filtra todo el tablero?} No. El filtro territorial exige su botón explícito.

\endgroup
\pantalla{recursos/pantallas/38_graficos_analiticos.jpg}{Visualizaciones con indicador y filtros}{Vista analítica: serie y rankings bajo el mismo indicador y filtros. Las categorías principales no siempre representan el total completo. Captura de referencia del 18 de noviembre de 2026; revise las instrucciones actualizadas de esta unidad.}
\pantalla{recursos/pantallas/39_respaldo_grafico.jpg}{Valores exactos de la serie visible}{Tabla de valores exactos de la serie de 2015 y California. Úsela para contrastar una lectura abreviada del gráfico. Captura de referencia del 18 de noviembre de 2026; revise las instrucciones actualizadas de esta unidad.}
\pantalla{recursos/pantallas/40_hallazgos.jpg}{Hallazgos y límites de interpretación}{Hallazgos automáticos derivados de los agregados. Comparar extremos visibles no demuestra causalidad ni una tendencia estadística. Captura de referencia del 18 de noviembre de 2026; revise las instrucciones actualizadas de esta unidad.}

\IfFileExists{capturas/AN02.tex}{\input{capturas/AN02.tex}}{}
\section{Conversar con el copiloto y revisar su respuesta}
\label{unidad:AN03}
\textbf{Para quién:} Analista de BI, Gerente comercial\par
\textbf{Al terminar podrá:} Formular preguntas respondibles y contrastar consulta interpretada, evidencia y conclusión.

El copiloto recibe contexto agregado y procedencia del datamart. Interpreta indicador, dimensión, filtros y Top N, y la plataforma puede ejecutar una consulta agregada permitida antes de pedir una explicación.

Una pregunta puede solicitar otro año o territorio disponible aunque no esté seleccionado en pantalla. Por eso revise Consulta interpretada: la respuesta puede corresponder a otro alcance solicitado y no al filtro actualmente visible.

Las reglas acotan formato, consultas y cautelas; no existe un validador universal de verdad para cada frase del modelo. La alumna debe verificar cifras, denominadores y límites antes de usar la respuesta.

Si cambia de fuente, ejecución, indicador o filtros mientras una respuesta está pendiente, la conversación anterior deja de corresponder al contexto visible. El sistema descarta la respuesta tardía para que no se presente como si fuera del análisis nuevo. Espere la actualización del panel y formule otra pregunta si necesita continuar.

\subsection{Formular preguntas útiles al copiloto sin perder el alcance}
Una pregunta recomendable dice qué indicador quiere, para qué período y con qué agrupación. Ejemplo: ¿Cuáles son los cinco productos con más unidades en el contexto actual? Después revise Consulta interpretada y su evidencia. La pregunta pide una comparación descriptiva; no demuestra por qué un cliente compró ni predice lo que hará.

Para evaluar coherencia, compruebe que la cifra, la unidad, el año, el territorio, la dimensión y el denominador de la respuesta coinciden con la consulta interpretada. Si la explicación excede esa evidencia, solicite aclaración o descarte la conclusión. Una redacción convincente no convierte una hipótesis en una causa probada.

En la campaña actual se probaron recorridos y conversaciones con Sol. La comprobación mínima de Luna no acredita todavía la misma calidad de explicación; si la usa en una práctica nueva, conserve la pregunta y revise la evidencia con los mismos criterios. No cambie automáticamente a otro modelo ante una respuesta débil.

\subsection{Conceptos para decidir con criterio}
\paragraph{Consulta interpretada} Indicador, agrupación, filtros, límite y orden que la plataforma entendió de la pregunta.

\paragraph{Agregado permitido} Cálculo cerrado sobre dimensiones materializadas de producto, cliente o territorio, con recetas aprobadas; no SQL libre.

\paragraph{Evidencia} Valores y alcance que respaldan la explicación; debe poder contrastarse con el panel o tablas.

\paragraph{Pregunta de seguimiento} Sugerencia ejecutable construida desde capacidades verificadas; no debe asumirse que abre un nuevo módulo de análisis.

\Needspace{8\baselineskip}
\begingroup
\setbox0=\vbox{\hsize=\linewidth
{\normalfont\normalsize\sffamily\bfseries\normalsize\color{Teal}\hyphenpenalty=10000\exhyphenpenalty=10000 Campos y controles\par}

\par
\begingroup
\fontsize{10}{12.5}\selectfont
\setlength{\tabcolsep}{5pt}
\renewcommand{\arraystretch}{1.15}
\begin{tabular}{@{}>{\RaggedRight\arraybackslash}p{0.27\linewidth}>{\RaggedRight\arraybackslash}p{0.66\linewidth}@{}}
\toprule
\textbf{Campo o control} & \textbf{Significado y criterio} \\
\midrule
Escriba su pregunta & Entre 5 y 500 caracteres. Indique indicador, comparación, período, territorio y qué quiere comprender. \\
\bottomrule
\end{tabular}
\endgroup
}
\Needspace{\dimexpr\ht0+\dp0+48pt\relax}
\endgroup
\subsection{Campos y controles}
\begingroup
\fontsize{10}{12.5}\selectfont
\setlength{\tabcolsep}{5pt}
\renewcommand{\arraystretch}{1.15}
\begin{longtable}{@{}>{\RaggedRight\arraybackslash}p{0.27\linewidth}>{\RaggedRight\arraybackslash}p{0.66\linewidth}@{}}
\toprule
\textbf{Campo o control} & \textbf{Significado y criterio} \\
\midrule
\endfirsthead
\toprule
\textbf{Campo o control} & \textbf{Significado y criterio} \\
\midrule
\endhead
\midrule
\multicolumn{2}{r}{\footnotesize Continúa en la página siguiente}\\
\endfoot
\bottomrule
\endlastfoot
Escriba su pregunta & Entre 5 y 500 caracteres. Indique indicador, comparación, período, territorio y qué quiere comprender. \\
\addlinespace[3pt]
Guía de preguntas sugeridas & Ejemplos adaptados al contexto; pueden usarse como punto de partida. \\
\addlinespace[3pt]
Preguntar al copiloto & Solicita interpretación/explicación. No crea una propuesta ni ejecuta otro ETL. \\
\addlinespace[3pt]
Consulta interpretada & Revise alcance y denominador antes de leer la narración. \\
\addlinespace[3pt]
Evidencia y cautela & Permiten reconocer lo calculado y lo que no se ha demostrado. \\
\addlinespace[3pt]
Proveedor y modelo & Identifican quién generó la respuesta; cambiar de proveedor no debe cambiar las reglas del datamart. \\
\addlinespace[3pt]
\end{longtable}
\endgroup

\subsection{Paso a paso}
\begin{enumerate}
\item Confirme ejecución, indicador y filtros antes de preguntar.\par\nopagebreak[4]\smallskip{\small\textbf{Compruebe:} Tiene un contexto inicial claro.}
\item Escriba una pregunta concreta, por ejemplo Top 5 de productos por unidades con los filtros actuales.\par\nopagebreak[4]\smallskip{\small\textbf{Compruebe:} Solicitud respondible con capacidades presentes.}
\item Espere la respuesta sin enviar varias veces.\par\nopagebreak[4]\smallskip{\small\textbf{Compruebe:} El copiloto muestra interpretación y evidencia si la operación cumple.}
\item Lea Consulta interpretada y compare indicador, año, territorio, límite y denominador.\par\nopagebreak[4]\smallskip{\small\textbf{Compruebe:} Reconoce cuándo la pregunta pidió un alcance distinto del panel.}
\item Contraste al menos una cifra con tabla o consulta agregada y revise los límites comerciales.\par\nopagebreak[4]\smallskip{\small\textbf{Compruebe:} No acepta una redacción coherente si contradice datos.}
\item Pregunte de nuevo para aclarar un denominador o límite, sin atribuir causas que no se han comprobado.\par\nopagebreak[4]\smallskip{\small\textbf{Compruebe:} Conversación útil y controlada.}
\end{enumerate}
\subsection{Qué hace la IA}
Interpreta la intención y redacta respuesta con evidencia agregada. Hay dos generaciones diferenciadas: intención estructurada y explicación. La plataforma aporta seguimiento y cautelas.

\subsection{Qué comprueba la plataforma}
Permisos, fuente/ejecución, indicadores y dimensiones disponibles, filtros, cantidades Top N entre 1 y 20, consultas agregadas e identificadores permitidos; no acepta SQL libre del modelo.

\subsection{Qué decide la persona}
La persona formula, revisa y decide si la explicación permite una conclusión. Una recomendación es una comprobación por realizar, no un hecho demostrado.

\paragraph{Resultado esperado} Respuesta contrastada con significado comercial y alcance reconocibles.

\Needspace{8\baselineskip}
\begingroup
\setbox0=\vbox{\hsize=\linewidth
{\normalfont\normalsize\sffamily\bfseries\normalsize\color{Teal}\hyphenpenalty=10000\exhyphenpenalty=10000 Errores y recuperación\par}

\par
\begingroup
\fontsize{10}{12.5}\selectfont
\setlength{\tabcolsep}{5pt}
\renewcommand{\arraystretch}{1.15}
\begin{tabular}{@{}>{\RaggedRight\arraybackslash}p{0.27\linewidth}>{\RaggedRight\arraybackslash}p{0.66\linewidth}@{}}
\toprule
\textbf{Situación} & \textbf{Qué hacer} \\
\midrule
El proveedor no responde o devuelve formato inválido. & Revise configuración, cuota/tiempo y credencial con el administrador. Un fallo de chat no requiere recargar el datamart. \\
\bottomrule
\end{tabular}
\endgroup
}
\Needspace{\dimexpr\ht0+\dp0+48pt\relax}
\endgroup
\subsection{Errores y recuperación}
\begingroup
\fontsize{10}{12.5}\selectfont
\setlength{\tabcolsep}{5pt}
\renewcommand{\arraystretch}{1.15}
\begin{longtable}{@{}>{\RaggedRight\arraybackslash}p{0.27\linewidth}>{\RaggedRight\arraybackslash}p{0.66\linewidth}@{}}
\toprule
\textbf{Situación} & \textbf{Qué hacer} \\
\midrule
\endfirsthead
\toprule
\textbf{Situación} & \textbf{Qué hacer} \\
\midrule
\endhead
\midrule
\multicolumn{2}{r}{\footnotesize Continúa en la página siguiente}\\
\endfoot
\bottomrule
\endlastfoot
El proveedor no responde o devuelve formato inválido. & Revise configuración, cuota/tiempo y credencial con el administrador. Un fallo de chat no requiere recargar el datamart. \\
\addlinespace[3pt]
La respuesta usa el año actual para un conjunto histórico. & Recuerde que el reloj del servidor no añade registros futuros; no acepte fechas ajenas a la disponibilidad. \\
\addlinespace[3pt]
Se pide un pronóstico del próximo período. & Es fuera del alcance. El resultado correcto es explicar el límite; no inventar un valor o año. \\
\addlinespace[3pt]
La respuesta llama facturación a los pedidos de AW. & Corrija la interpretación: no se demuestra factura emitida a partir del pedido. \\
\addlinespace[3pt]
La respuesta es fluida pero cambia la cifra. & No la use como evidencia. Registre pregunta, contexto y discrepancia, reintente aclaración y contraste con cálculos determinísticos. \\
\addlinespace[3pt]
\end{longtable}
\endgroup

\subsection{Ejercicio}
En AW 21 filtre 2013/Northwest, indicador unidades, y pregunte: Resume las unidades vendidas en 2013 para Northwest. Indica el total y el mes con mayor cantidad, sin atribuir causas ni mezclar importes.

\textbf{Respuesta orientativa.} El control histórico del 6 noviembre registró 17.452 unidades y junio 2013 con 2.494. Una sesión nueva debe contrastar sus propios resultados; no copiar esa cifra si el alcance o la fuente cambiaron.

\subsection{Compruebe lo que aprendió}
\paragraph{¿Se ejecuta SQL redactado por el LLM?} No. La intención se limita a consultas agregadas permitidas del servidor.

\paragraph{¿El copiloto sólo puede responder el filtro visible?} Puede interpretar otro filtro disponible solicitado. Debe leer la consulta interpretada y no asumir que los dos alcances coinciden.

\paragraph{¿Una respuesta correcta una vez garantiza que el proveedor nunca fallará?} No. Se deben revisar cada respuesta y sus controles.

\pantalla{recursos/pantallas/41_copiloto.jpg}{Preguntas contextuales al copiloto}{Conversación contextual: formule una pregunta concreta y compruebe la evidencia de la respuesta. La captura no representa una respuesta nueva del LLM. Captura de referencia del 18 de noviembre de 2026; revise las instrucciones actualizadas de esta unidad.}

\IfFileExists{capturas/AN03.tex}{\input{capturas/AN03.tex}}{}
\section{Exportar y comprobar antes de compartir}
\label{unidad:AN04}
\textbf{Para quién:} Analista de BI, Gerente comercial con permiso de exportación\par
\textbf{Al terminar podrá:} Obtener informes que conserven universo y procedencia y revisarlos contra el panel.

Exportar PDF produce un informe para lectura. Exportar Excel produce tablas y gráficos de respaldo de las visualizaciones visibles. Ninguno es un respaldo de contenedores, volúmenes o todas las filas de la fuente.

La exportación se genera desde el dashboard autorizado. No es una exportación automática de la conversación del chat; si desea citar una respuesta debe registrar la pregunta y su contexto como evidencia aparte.

\subsection{Conceptos para decidir con criterio}
\paragraph{Contexto exportado} Fuente, propuesta, ejecución, vista, métrica, filtros y unidad que permiten interpretar el archivo.

\paragraph{Trazabilidad} Información para reconocer qué expediente y selección respaldan las cifras.

\paragraph{Redondeo} Presentación con decimales o abreviaturas; no autoriza a cambiar valores para hacerlos coincidir visualmente.

\Needspace{8\baselineskip}
\begingroup
\setbox0=\vbox{\hsize=\linewidth
{\normalfont\normalsize\sffamily\bfseries\normalsize\color{Teal}\hyphenpenalty=10000\exhyphenpenalty=10000 Campos y controles\par}

\par
\begingroup
\fontsize{10}{12.5}\selectfont
\setlength{\tabcolsep}{5pt}
\renewcommand{\arraystretch}{1.15}
\begin{tabular}{@{}>{\RaggedRight\arraybackslash}p{0.27\linewidth}>{\RaggedRight\arraybackslash}p{0.66\linewidth}@{}}
\toprule
\textbf{Campo o control} & \textbf{Significado y criterio} \\
\midrule
Exportar PDF & Descarga el informe de la vista y selección actuales. \\
\bottomrule
\end{tabular}
\endgroup
}
\Needspace{\dimexpr\ht0+\dp0+48pt\relax}
\endgroup
\subsection{Campos y controles}
\begingroup
\fontsize{10}{12.5}\selectfont
\setlength{\tabcolsep}{5pt}
\renewcommand{\arraystretch}{1.15}
\begin{longtable}{@{}>{\RaggedRight\arraybackslash}p{0.27\linewidth}>{\RaggedRight\arraybackslash}p{0.66\linewidth}@{}}
\toprule
\textbf{Campo o control} & \textbf{Significado y criterio} \\
\midrule
\endfirsthead
\toprule
\textbf{Campo o control} & \textbf{Significado y criterio} \\
\midrule
\endhead
\midrule
\multicolumn{2}{r}{\footnotesize Continúa en la página siguiente}\\
\endfoot
\bottomrule
\endlastfoot
Exportar PDF & Descarga el informe de la vista y selección actuales. \\
\addlinespace[3pt]
Exportar Excel & Descarga el libro con datos que respaldan los gráficos exportados. \\
\addlinespace[3pt]
Vista ejecutiva/Vista analítica & Afecta qué visualizaciones se comparten; revise la vista junto a los filtros. \\
\addlinespace[3pt]
Carpeta Descargas & Abra el archivo recibido y revíselo, no basta con pulsar el botón. \\
\addlinespace[3pt]
\end{longtable}
\endgroup

\subsection{Paso a paso}
\begin{enumerate}
\item Fije fuente, ejecución, vista, indicador, año y territorio.\par\nopagebreak[4]\smallskip{\small\textbf{Compruebe:} Alcance exacto que se quiere compartir.}
\item Exporte PDF y espere la descarga.\par\nopagebreak[4]\smallskip{\small\textbf{Compruebe:} Archivo con nombre de ejecución y vista.}
\item Exporte Excel.\par\nopagebreak[4]\smallskip{\small\textbf{Compruebe:} Libro de revisión, no volcado completo del origen.}
\item Abra ambos archivos y compare un total, una entidad y el filtro con el panel.\par\nopagebreak[4]\smallskip{\small\textbf{Compruebe:} Cifras y contexto congruentes.}
\item Revise unidad, advertencias, denominadores y procedencia.\par\nopagebreak[4]\smallskip{\small\textbf{Compruebe:} El informe no pierde el límite demostrado en pantalla.}
\item Registre fecha, nombre y contexto del archivo revisado.\par\nopagebreak[4]\smallskip{\small\textbf{Compruebe:} Evidencia de revisión real, no sólo descarga.}
\end{enumerate}
\subsection{Qué hace la IA}
No se necesita otra llamada al LLM para generar el informe del dashboard. El chat no convierte el reporte en una conclusión validada por un modelo.

\subsection{Qué comprueba la plataforma}
Permiso de exportación además de lectura de analítica; servidor reconstruye el mismo dashboard con filtros autorizados y registra la acción.

\subsection{Qué decide la persona}
La persona elige selección, vista y formato y comprueba el archivo antes de compartirlo.

\paragraph{Resultado esperado} Dos archivos revisados cuyo alcance puede explicarse sin abrir o trabase.

\Needspace{8\baselineskip}
\begingroup
\setbox0=\vbox{\hsize=\linewidth
{\normalfont\normalsize\sffamily\bfseries\normalsize\color{Teal}\hyphenpenalty=10000\exhyphenpenalty=10000 Errores y recuperación\par}

\par
\begingroup
\fontsize{10}{12.5}\selectfont
\setlength{\tabcolsep}{5pt}
\renewcommand{\arraystretch}{1.15}
\begin{tabular}{@{}>{\RaggedRight\arraybackslash}p{0.27\linewidth}>{\RaggedRight\arraybackslash}p{0.66\linewidth}@{}}
\toprule
\textbf{Situación} & \textbf{Qué hacer} \\
\midrule
No hay botones de exportación. & Revise que el rol tenga permisos de lectura del panel y exportación. No se corrige cambiando de vista. \\
\bottomrule
\end{tabular}
\endgroup
}
\Needspace{\dimexpr\ht0+\dp0+48pt\relax}
\endgroup
\subsection{Errores y recuperación}
\begingroup
\fontsize{10}{12.5}\selectfont
\setlength{\tabcolsep}{5pt}
\renewcommand{\arraystretch}{1.15}
\begin{longtable}{@{}>{\RaggedRight\arraybackslash}p{0.27\linewidth}>{\RaggedRight\arraybackslash}p{0.66\linewidth}@{}}
\toprule
\textbf{Situación} & \textbf{Qué hacer} \\
\midrule
\endfirsthead
\toprule
\textbf{Situación} & \textbf{Qué hacer} \\
\midrule
\endhead
\midrule
\multicolumn{2}{r}{\footnotesize Continúa en la página siguiente}\\
\endfoot
\bottomrule
\endlastfoot
No hay botones de exportación. & Revise que el rol tenga permisos de lectura del panel y exportación. No se corrige cambiando de vista. \\
\addlinespace[3pt]
El PDF tiene menos gráficos que la Vista analítica. & Compruebe qué vista exportó; la ejecutiva omite el visual de clientes. \\
\addlinespace[3pt]
El porcentaje de Excel no coincide. & Compruebe si es escala 0 a 100 o fracción y qué denominador se usó; no multiplique otra vez por 100 por costumbre. \\
\addlinespace[3pt]
El archivo muestra moneda de origen aunque AW esté en USD. & La evidencia histórica AW 21 conservó una etiqueta genérica en una razón del PDF, pese al valor correcto. No presente ese antecedente como perfección semántica; revise el archivo nuevo. \\
\addlinespace[3pt]
\end{longtable}
\endgroup

\subsection{Ejercicio}
Explique por qué no es suficiente mostrar que se descargó un informe para probar que es correcto.

\textbf{Respuesta orientativa.} Debe abrirse y comparar cifras, contexto, filtros, unidades y trazabilidad con la selección autorizada.

\subsection{Compruebe lo que aprendió}
\paragraph{¿Excel incluye todas las filas del OLTP?} No. Incluye respaldo de las visualizaciones exportadas.

\paragraph{¿Se exigen siempre seis páginas y seis hojas?} No. Esas cantidades pertenecen a los casos históricos concretos; varían con la vista y las visualizaciones.


\IfFileExists{capturas/AN04.tex}{\input{capturas/AN04.tex}}{}

\chapter{Administrar cuentas, permisos y trazabilidad}
\section{Crear y mantener cuentas de usuario}
\label{unidad:operacion_04}
\textbf{Para quién:} Administración\par
\textbf{Opción:} Administración y seguridad > Usuarios\par

Una cuenta se crea para una persona, no para una base de datos. Los roles asignados determinan sus permisos efectivos en toda la plataforma. La misma cuenta de analista puede seleccionar las fuentes habilitadas para el flujo disponible; el sistema no crea automáticamente una cuenta por conexión.

En Usuarios, el formulario se encuentra en la propia pantalla. Para crear se completan Correo, Nombre completo, Contraseña y Roles asignados, y se pulsa Crear registro. En edición, el correo identifica a la cuenta y permanece en consulta. Una contraseña vacía conserva la existente; escribir una nueva la sustituye.

No es necesario escribir identificadores numéricos: los roles se eligen por nombre y descripción. Si falla la carga de roles, el formulario queda bloqueado y aparece Reintentar carga de opciones. Un error de red no debe resolverse guardando una cuenta sin sus roles previstos.

El estado true de la tabla significa activo; false significa inactivo. Una cuenta desactivada conserva su registro, pero no puede usar una sesión válida. Eliminar es una operación diferente y solicita confirmación. Los registros Protegido y la recuperación administrativa tienen controles adicionales; la administración inicial no debe desaparecer como parte de una práctica.

\subsection{Conceptos para comprender la opción}
\paragraph{Cuenta} Identidad de acceso con correo y nombre.

\paragraph{Rol} Conjunto de permisos asignable a una cuenta. Analista BI o Gerente comercial son nombres de perfiles configurados.

\paragraph{Permisos efectivos} Unión de permisos activos de los roles activos de la cuenta.

\subsection{Paso a paso}
\begin{enumerate}
\item Abra Usuarios y compruebe si va a crear una cuenta o editar una existente.
\item Para una cuenta de práctica, escriba un correo no utilizado y un nombre que identifique a la persona. Defina una clave de al menos 12 caracteres sin incluirla en capturas o documentos.
\item En Roles asignados, marque los perfiles necesarios y lea su descripción.
\item Pulse Crear registro y verifique la nueva fila.
\item Para editar, pulse Editar en la fila exacta. Cambie el nombre o roles; deje Contraseña vacía si no desea cambiarla.
\item Pulse Guardar cambios y vuelva a abrir la misma fila para comprobar las asociaciones. Utilice Cancelar si no desea guardar.
\item La prueba completa del perfil requiere iniciar sesión con esa cuenta y verificar operaciones permitidas y denegadas; la lista de casillas por sí sola no lo demuestra.
\end{enumerate}

\Needspace{8\baselineskip}
\begingroup
\setbox0=\vbox{\hsize=\linewidth
{\normalfont\normalsize\sffamily\bfseries\normalsize\color{Teal}\hyphenpenalty=10000\exhyphenpenalty=10000 Campos y botones\par}

\par
\begingroup
\fontsize{10}{12.5}\selectfont
\setlength{\tabcolsep}{5pt}
\renewcommand{\arraystretch}{1.15}
\begin{tabular}{@{}>{\RaggedRight\arraybackslash}p{0.22\linewidth}>{\RaggedRight\arraybackslash}p{0.3\linewidth}>{\RaggedRight\arraybackslash}p{0.39\linewidth}@{}}
\toprule
\textbf{Control} & \textbf{Para qué sirve} & \textbf{Ejemplo y resultado} \\
\midrule
Correo & Identifica una cuenta nueva y no se edita después desde esta pantalla. & \textbf{Ejemplo:} Correo asignado a la persona.\par\smallskip\textbf{Resultado:} Evita que se confunda el nombre visible con el identificador de acceso. \\
\bottomrule
\end{tabular}
\endgroup
}
\Needspace{\dimexpr\ht0+\dp0+48pt\relax}
\endgroup
\subsection{Campos y botones}
\begingroup
\fontsize{10}{12.5}\selectfont
\setlength{\tabcolsep}{5pt}
\renewcommand{\arraystretch}{1.15}
\begin{longtable}{@{}>{\RaggedRight\arraybackslash}p{0.22\linewidth}>{\RaggedRight\arraybackslash}p{0.3\linewidth}>{\RaggedRight\arraybackslash}p{0.39\linewidth}@{}}
\toprule
\textbf{Control} & \textbf{Para qué sirve} & \textbf{Ejemplo y resultado} \\
\midrule
\endfirsthead
\toprule
\textbf{Control} & \textbf{Para qué sirve} & \textbf{Ejemplo y resultado} \\
\midrule
\endhead
\midrule
\multicolumn{3}{r}{\footnotesize Continúa en la página siguiente}\\
\endfoot
\bottomrule
\endlastfoot
Correo & Identifica una cuenta nueva y no se edita después desde esta pantalla. & \textbf{Ejemplo:} Correo asignado a la persona.\par\smallskip\textbf{Resultado:} Evita que se confunda el nombre visible con el identificador de acceso. \\
\addlinespace[3pt]
Nombre completo & Permite reconocer a quien opera y atribuirle acciones. & \textbf{Ejemplo:} Nombre de la estudiante, no su correo.\par\smallskip\textbf{Resultado:} Se usa como identificación visible de la cuenta. \\
\addlinespace[3pt]
Contraseña & Establece o cambia la clave. & \textbf{Ejemplo:} Al editar, dejar vacío para conservarla.\par\smallskip\textbf{Resultado:} Se guarda protegida y no vuelve a mostrarse. \\
\addlinespace[3pt]
Roles asignados & Concede los roles seleccionados. & \textbf{Ejemplo:} Marcar Gerente comercial para consulta.\par\smallskip\textbf{Resultado:} La cuenta recibe los permisos de sus roles activos. \\
\addlinespace[3pt]
Desactivar / Activar & Cambia la disponibilidad de una cuenta no protegida. & \textbf{Ejemplo:} Desactivar una cuenta de práctica que ya no se usa.\par\smallskip\textbf{Resultado:} El registro permanece; el acceso se deniega si está inactivo. \\
\addlinespace[3pt]
Eliminar & Solicita eliminación definitiva de una cuenta no protegida. & \textbf{Ejemplo:} Retirar una cuenta temporal autorizada.\par\smallskip\textbf{Resultado:} El servidor valida dependencias y conserva la evidencia histórica segura que corresponda. \\
\addlinespace[3pt]
\end{longtable}
\endgroup

\subsection{Editar sin cambiar la contraseña}
Práctica propuesta para el entorno de capacitación; no se presenta como una nueva ejecución realizada durante la redacción.

\begin{enumerate}
\item En un entorno de capacitación, abra una cuenta no protegida.
\item Cambie únicamente su Nombre completo y deje Contraseña vacía.
\item Guarde y vuelva a consultar la fila.
\item Compruebe acceso con las credenciales conocidas de esa cuenta cuando corresponda.
\end{enumerate}

\textbf{Qué debe obtener.} El nombre se actualiza. No se simula un cambio de correo ni se exige reemplazar la clave para editar los roles.

\textbf{Cómo comprobarlo.} La cuenta conserva el correo y la contraseña previstos y recibe solo los roles seleccionados.

\Needspace{8\baselineskip}
\begingroup
\setbox0=\vbox{\hsize=\linewidth
{\normalfont\normalsize\sffamily\bfseries\normalsize\color{Teal}\hyphenpenalty=10000\exhyphenpenalty=10000 Si aparece un problema\par}

\par
\begingroup
\fontsize{10}{12.5}\selectfont
\setlength{\tabcolsep}{5pt}
\renewcommand{\arraystretch}{1.15}
\begin{tabular}{@{}>{\RaggedRight\arraybackslash}p{0.22\linewidth}>{\RaggedRight\arraybackslash}p{0.3\linewidth}>{\RaggedRight\arraybackslash}p{0.39\linewidth}@{}}
\toprule
\textbf{Qué observa} & \textbf{Causa posible} & \textbf{Qué hacer y comprobar} \\
\midrule
No se pueden cargar los roles. & Permiso Consultar roles ausente o fallo al consultar el catálogo. & Solicite revisar el permiso complementario y utilice Reintentar carga de opciones. No guarde hasta que la tabla sea la correcta.\par\smallskip\textbf{Comprobación:} El formulario vuelve a habilitarse y muestra los roles esperados. \\
\bottomrule
\end{tabular}
\endgroup
}
\Needspace{\dimexpr\ht0+\dp0+48pt\relax}
\endgroup
\subsection{Si aparece un problema}
\begingroup
\fontsize{10}{12.5}\selectfont
\setlength{\tabcolsep}{5pt}
\renewcommand{\arraystretch}{1.15}
\begin{longtable}{@{}>{\RaggedRight\arraybackslash}p{0.22\linewidth}>{\RaggedRight\arraybackslash}p{0.3\linewidth}>{\RaggedRight\arraybackslash}p{0.39\linewidth}@{}}
\toprule
\textbf{Qué observa} & \textbf{Causa posible} & \textbf{Qué hacer y comprobar} \\
\midrule
\endfirsthead
\toprule
\textbf{Qué observa} & \textbf{Causa posible} & \textbf{Qué hacer y comprobar} \\
\midrule
\endhead
\midrule
\multicolumn{3}{r}{\footnotesize Continúa en la página siguiente}\\
\endfoot
\bottomrule
\endlastfoot
No se pueden cargar los roles. & Permiso Consultar roles ausente o fallo al consultar el catálogo. & Solicite revisar el permiso complementario y utilice Reintentar carga de opciones. No guarde hasta que la tabla sea la correcta.\par\smallskip\textbf{Comprobación:} El formulario vuelve a habilitarse y muestra los roles esperados. \\
\addlinespace[3pt]
El correo no cambia al editar. & Es un campo de consulta en la edición vigente. & Mantenga el correo original. Si se necesita otra identidad de acceso, acuerde con la administración el tratamiento de la cuenta; no anuncie un cambio no admitido.\par\smallskip\textbf{Comprobación:} La fila conserva el correo que identifica la cuenta. \\
\addlinespace[3pt]
No se permite eliminar o desactivar. & Cuenta protegida, sesión propia o control de recuperación/dependencias. & Lea la causa y utilice una cuenta de práctica autorizada. No altere directamente PostgreSQL para eludir el control.\par\smallskip\textbf{Comprobación:} La vía de administración permanece disponible. \\
\addlinespace[3pt]
\end{longtable}
\endgroup

\subsection{Para explicar durante la sustentación}
\paragraph{¿Por qué no se cambia el correo al editar?} La cuenta se identifica con ese correo y el contrato de edición no admite modificarlo. La pantalla lo muestra en consulta para evitar simular un cambio.

\paragraph{¿Ver las casillas marcadas prueba la seguridad del rol?} No. Permite revisar la configuración; la verificación operativa exige autenticar la cuenta y probar permisos permitidos y denegados.

\pantalla{recursos/pantallas/44_usuarios.jpg}{Alta y asignación de roles: formulario sin enviar}{Formulario de usuario vacío y sin enviar. Asignar un rol concede permisos; utilice cuentas autorizadas para las prácticas. Captura de referencia del 18 de noviembre de 2026; revise las instrucciones actualizadas de esta unidad.}

\IfFileExists{capturas/operacion_04.tex}{\input{capturas/operacion_04.tex}}{}
\section{Definir roles según las responsabilidades}
\label{unidad:operacion_05}
\textbf{Para quién:} Administración\par
\textbf{Opción:} Administración y seguridad > Roles\par

Diseñe primero la responsabilidad y después sus permisos. La administración prepara el entorno y el acceso; el analista convierte necesidades en diseños y resultados verificados; la gerencia consulta, interpreta y exporta. El nombre del rol es una etiqueta: llamar a un rol Gerente comercial no le concede automáticamente permisos de lectura.

En Crear rol se escriben Nombre y Descripción y se marcan Permisos otorgados. El sistema genera la clave interna. Use los nombres y las explicaciones de la tabla; no necesita conocer códigos técnicos para operar. Al editar un rol, una casilla retirada deja de formar parte de sus permisos al guardar.

La instalación base garantiza el rol administrativo protegido; Analista BI, Gerente comercial y otros perfiles pueden ser configuraciones locales de la demostración. Una copia limpia puede requerir crearlos o asignarlos. Por eso este manual distingue la matriz recomendada de la evidencia de una instalación concreta.

Tras una modificación, vuelva a abrir el rol para comprobar nombre y casillas. La edición de datos y de permisos utiliza operaciones separadas: si aparece un error después de una actualización parcial, revise qué quedó guardado y complete la corrección. No dé por sentado que todo se guardó por haber pulsado el botón.

Eliminar y Desactivar tienen efectos distintos. Eliminar borra definitivamente un rol no protegido y exige que ya no tenga asignaciones a usuarios ni permisos asociados; también cuentan las asociaciones existentes aunque el usuario o el permiso estén inactivos. La plataforma no retira esas relaciones automáticamente para permitir el borrado. Si se busca suspender el uso del perfil y conservar su definición, debe evaluarse Desactivar en lugar de eliminarlo.

\begingroup
\setbox0=\vbox{\hsize=\linewidth
{\normalfont\normalsize\sffamily\bfseries\normalsize\color{Teal}\hyphenpenalty=10000\exhyphenpenalty=10000 Roles y permisos observados en la instalación de referencia\par}

La siguiente relación procede de una consulta de definiciones de roles y permisos, sin cuentas ni asignaciones de personas. Describe la configuración observada; no acredita una sesión autenticada ni una prueba integral bajo cada perfil. Una instalación limpia puede conservar únicamente el perfil administrativo base.

\par
\begingroup
\fontsize{10}{12.5}\selectfont
\setlength{\tabcolsep}{5pt}
\renewcommand{\arraystretch}{1.15}
\begin{tabular}{@{}>{\RaggedRight\arraybackslash}p{0.22\linewidth}>{\RaggedRight\arraybackslash}p{0.3\linewidth}>{\RaggedRight\arraybackslash}p{0.39\linewidth}@{}}
\toprule
\textbf{Perfil} & \textbf{Permisos activos} & \textbf{Alcance observado} \\
\midrule
Administrador & 29 & Perfil protegido de administración con los 29 permisos base. \\
\bottomrule
\end{tabular}
\endgroup
}
\Needspace{\dimexpr\ht0+\dp0+48pt\relax}
\endgroup
\subsection{Roles y permisos observados en la instalación de referencia}
La siguiente relación procede de una consulta de definiciones de roles y permisos, sin cuentas ni asignaciones de personas. Describe la configuración observada; no acredita una sesión autenticada ni una prueba integral bajo cada perfil. Una instalación limpia puede conservar únicamente el perfil administrativo base.

\begingroup
\fontsize{10}{12.5}\selectfont
\setlength{\tabcolsep}{5pt}
\renewcommand{\arraystretch}{1.15}
\begin{longtable}{@{}>{\RaggedRight\arraybackslash}p{0.22\linewidth}>{\RaggedRight\arraybackslash}p{0.3\linewidth}>{\RaggedRight\arraybackslash}p{0.39\linewidth}@{}}
\toprule
\textbf{Perfil} & \textbf{Permisos activos} & \textbf{Alcance observado} \\
\midrule
\endfirsthead
\toprule
\textbf{Perfil} & \textbf{Permisos activos} & \textbf{Alcance observado} \\
\midrule
\endhead
\midrule
\multicolumn{3}{r}{\footnotesize Continúa en la página siguiente}\\
\endfoot
\bottomrule
\endlastfoot
Administrador & 29 & Perfil protegido de administración con los 29 permisos base. \\
\addlinespace[3pt]
Operador & 4 & Configuración local de cuatro permisos específicos; no representa un perfil completo de operación ni acredita acceso a todos los módulos. \\
\addlinespace[3pt]
Analista BI & 15 & Configuración local para diseñar, revisar y materializar con consulta de contexto, catálogo, auditoría y analítica. \\
\addlinespace[3pt]
Gerente comercial & 4 & Consulta analítica, conexiones sin secretos, metadatos y exportación. No incluye ETL ni aprobación de propuestas. \\
\addlinespace[3pt]
\end{longtable}
\endgroup

\begingroup
\setbox0=\vbox{\hsize=\linewidth
{\normalfont\normalsize\normalfont\normalsize\bfseries\hyphenpenalty=10000\exhyphenpenalty=10000 Administrador\par}

\par
\begingroup
\fontsize{10}{12.5}\selectfont
\setlength{\tabcolsep}{5pt}
\renewcommand{\arraystretch}{1.15}
\begin{tabular}{@{}>{\RaggedRight\arraybackslash}p{0.27\linewidth}>{\RaggedRight\arraybackslash}p{0.66\linewidth}@{}}
\toprule
\textbf{Permiso concedido} & \textbf{Estado} \\
\midrule
Consultar análisis visual & Activo \\
\bottomrule
\end{tabular}
\endgroup
}
\Needspace{\dimexpr\ht0+\dp0+48pt\relax}
\endgroup
\paragraph{Administrador}
\begingroup
\fontsize{10}{12.5}\selectfont
\setlength{\tabcolsep}{5pt}
\renewcommand{\arraystretch}{1.15}
\begin{longtable}{@{}>{\RaggedRight\arraybackslash}p{0.27\linewidth}>{\RaggedRight\arraybackslash}p{0.66\linewidth}@{}}
\toprule
\textbf{Permiso concedido} & \textbf{Estado} \\
\midrule
\endfirsthead
\toprule
\textbf{Permiso concedido} & \textbf{Estado} \\
\midrule
\endhead
\midrule
\multicolumn{2}{r}{\footnotesize Continúa en la página siguiente}\\
\endfoot
\bottomrule
\endlastfoot
Consultar análisis visual & Activo \\
\addlinespace[3pt]
Consultar auditoría & Activo \\
\addlinespace[3pt]
Consultar conexiones & Activo \\
\addlinespace[3pt]
Probar conexiones & Activo \\
\addlinespace[3pt]
Administrar conexiones & Activo \\
\addlinespace[3pt]
Consultar catálogo analítico & Activo \\
\addlinespace[3pt]
Administrar catálogo analítico & Activo \\
\addlinespace[3pt]
Generar propuestas BI & Activo \\
\addlinespace[3pt]
Consultar propuestas BI & Activo \\
\addlinespace[3pt]
Revisar propuestas BI & Activo \\
\addlinespace[3pt]
Consultar materializaciones & Activo \\
\addlinespace[3pt]
Preparar materializaciones & Activo \\
\addlinespace[3pt]
Consultar metadatos & Activo \\
\addlinespace[3pt]
Actualizar metadatos & Activo \\
\addlinespace[3pt]
Resolver identidad de dimensiones & Activo \\
\addlinespace[3pt]
Revisar identidad de dimensiones & Activo \\
\addlinespace[3pt]
Consultar configuración LLM & Activo \\
\addlinespace[3pt]
Administrar configuración LLM & Activo \\
\addlinespace[3pt]
Consultar parámetros & Activo \\
\addlinespace[3pt]
Administrar parámetros & Activo \\
\addlinespace[3pt]
Exportar reportes analíticos & Activo \\
\addlinespace[3pt]
Consultar menús & Activo \\
\addlinespace[3pt]
Administrar menús & Activo \\
\addlinespace[3pt]
Consultar permisos & Activo \\
\addlinespace[3pt]
Administrar permisos & Activo \\
\addlinespace[3pt]
Consultar roles & Activo \\
\addlinespace[3pt]
Administrar roles & Activo \\
\addlinespace[3pt]
Consultar usuarios & Activo \\
\addlinespace[3pt]
Administrar usuarios & Activo \\
\addlinespace[3pt]
\end{longtable}
\endgroup

\begingroup
\setbox0=\vbox{\hsize=\linewidth
{\normalfont\normalsize\normalfont\normalsize\bfseries\hyphenpenalty=10000\exhyphenpenalty=10000 Operador\par}

\par
\begingroup
\fontsize{10}{12.5}\selectfont
\setlength{\tabcolsep}{5pt}
\renewcommand{\arraystretch}{1.15}
\begin{tabular}{@{}>{\RaggedRight\arraybackslash}p{0.27\linewidth}>{\RaggedRight\arraybackslash}p{0.66\linewidth}@{}}
\toprule
\textbf{Permiso concedido} & \textbf{Estado} \\
\midrule
Probar conexiones & Activo \\
\bottomrule
\end{tabular}
\endgroup
}
\Needspace{\dimexpr\ht0+\dp0+48pt\relax}
\endgroup
\paragraph{Operador}
\begingroup
\fontsize{10}{12.5}\selectfont
\setlength{\tabcolsep}{5pt}
\renewcommand{\arraystretch}{1.15}
\begin{longtable}{@{}>{\RaggedRight\arraybackslash}p{0.27\linewidth}>{\RaggedRight\arraybackslash}p{0.66\linewidth}@{}}
\toprule
\textbf{Permiso concedido} & \textbf{Estado} \\
\midrule
\endfirsthead
\toprule
\textbf{Permiso concedido} & \textbf{Estado} \\
\midrule
\endhead
\midrule
\multicolumn{2}{r}{\footnotesize Continúa en la página siguiente}\\
\endfoot
\bottomrule
\endlastfoot
Probar conexiones & Activo \\
\addlinespace[3pt]
Consultar configuración LLM & Activo \\
\addlinespace[3pt]
Administrar configuración LLM & Activo \\
\addlinespace[3pt]
Administrar parámetros & Activo \\
\addlinespace[3pt]
\end{longtable}
\endgroup

\begingroup
\setbox0=\vbox{\hsize=\linewidth
{\normalfont\normalsize\normalfont\normalsize\bfseries\hyphenpenalty=10000\exhyphenpenalty=10000 Analista BI\par}

\par
\begingroup
\fontsize{10}{12.5}\selectfont
\setlength{\tabcolsep}{5pt}
\renewcommand{\arraystretch}{1.15}
\begin{tabular}{@{}>{\RaggedRight\arraybackslash}p{0.27\linewidth}>{\RaggedRight\arraybackslash}p{0.66\linewidth}@{}}
\toprule
\textbf{Permiso concedido} & \textbf{Estado} \\
\midrule
Consultar análisis visual & Activo \\
\bottomrule
\end{tabular}
\endgroup
}
\Needspace{\dimexpr\ht0+\dp0+48pt\relax}
\endgroup
\paragraph{Analista BI}
\begingroup
\fontsize{10}{12.5}\selectfont
\setlength{\tabcolsep}{5pt}
\renewcommand{\arraystretch}{1.15}
\begin{longtable}{@{}>{\RaggedRight\arraybackslash}p{0.27\linewidth}>{\RaggedRight\arraybackslash}p{0.66\linewidth}@{}}
\toprule
\textbf{Permiso concedido} & \textbf{Estado} \\
\midrule
\endfirsthead
\toprule
\textbf{Permiso concedido} & \textbf{Estado} \\
\midrule
\endhead
\midrule
\multicolumn{2}{r}{\footnotesize Continúa en la página siguiente}\\
\endfoot
\bottomrule
\endlastfoot
Consultar análisis visual & Activo \\
\addlinespace[3pt]
Consultar auditoría & Activo \\
\addlinespace[3pt]
Consultar conexiones & Activo \\
\addlinespace[3pt]
Consultar catálogo analítico & Activo \\
\addlinespace[3pt]
Administrar catálogo analítico & Activo \\
\addlinespace[3pt]
Generar propuestas BI & Activo \\
\addlinespace[3pt]
Consultar propuestas BI & Activo \\
\addlinespace[3pt]
Revisar propuestas BI & Activo \\
\addlinespace[3pt]
Consultar materializaciones & Activo \\
\addlinespace[3pt]
Preparar materializaciones & Activo \\
\addlinespace[3pt]
Consultar metadatos & Activo \\
\addlinespace[3pt]
Actualizar metadatos & Activo \\
\addlinespace[3pt]
Resolver identidad de dimensiones & Activo \\
\addlinespace[3pt]
Revisar identidad de dimensiones & Activo \\
\addlinespace[3pt]
Exportar reportes analíticos & Activo \\
\addlinespace[3pt]
\end{longtable}
\endgroup

\begingroup
\setbox0=\vbox{\hsize=\linewidth
{\normalfont\normalsize\normalfont\normalsize\bfseries\hyphenpenalty=10000\exhyphenpenalty=10000 Gerente comercial\par}

\par
\begingroup
\fontsize{10}{12.5}\selectfont
\setlength{\tabcolsep}{5pt}
\renewcommand{\arraystretch}{1.15}
\begin{tabular}{@{}>{\RaggedRight\arraybackslash}p{0.27\linewidth}>{\RaggedRight\arraybackslash}p{0.66\linewidth}@{}}
\toprule
\textbf{Permiso concedido} & \textbf{Estado} \\
\midrule
Consultar análisis visual & Activo \\
\bottomrule
\end{tabular}
\endgroup
}
\Needspace{\dimexpr\ht0+\dp0+48pt\relax}
\endgroup
\paragraph{Gerente comercial}
\begingroup
\fontsize{10}{12.5}\selectfont
\setlength{\tabcolsep}{5pt}
\renewcommand{\arraystretch}{1.15}
\begin{longtable}{@{}>{\RaggedRight\arraybackslash}p{0.27\linewidth}>{\RaggedRight\arraybackslash}p{0.66\linewidth}@{}}
\toprule
\textbf{Permiso concedido} & \textbf{Estado} \\
\midrule
\endfirsthead
\toprule
\textbf{Permiso concedido} & \textbf{Estado} \\
\midrule
\endhead
\midrule
\multicolumn{2}{r}{\footnotesize Continúa en la página siguiente}\\
\endfoot
\bottomrule
\endlastfoot
Consultar análisis visual & Activo \\
\addlinespace[3pt]
Consultar conexiones & Activo \\
\addlinespace[3pt]
Consultar metadatos & Activo \\
\addlinespace[3pt]
Exportar reportes analíticos & Activo \\
\addlinespace[3pt]
\end{longtable}
\endgroup

Gerente comercial necesita Consultar metadatos porque el selector actual de fuentes utiliza esa consulta. También puede abrir el Explorador en lectura; ese permiso no concede actualización de metadatos ni ejecución ETL. En Operador, Probar conexiones y Administrar parámetros no incluyen por sí solos los permisos de consulta complementarios: no se asume que pueda recorrer sus pantallas completas.

\subsection{Conceptos para comprender la opción}
\paragraph{Mínimo privilegio} Conceder solo las operaciones necesarias para una responsabilidad.

\paragraph{Consulta y administración} Leer un recurso y modificarlo son permisos diferentes. Consultar catálogo no equivale a Administrar catálogo.

\paragraph{Rol protegido} Perfil del sistema con controles que conservan la recuperación administrativa.

\subsection{Paso a paso}
\begin{enumerate}
\item Liste las tareas del perfil que va a crear. Para gerencia: consultar el panel, seleccionar fuentes y exportar resultados.
\item Abra Roles y complete Nombre y Descripción.
\item Marque los permisos necesarios en Permisos otorgados. Para un rol con selector de fuentes, incluya Consultar metadatos.
\item Pulse Crear registro. Para modificar uno existente, abra Editar, ajuste las casillas y pulse Guardar cambios.
\item Vuelva a abrir el rol y coteje sus permisos con la matriz acordada.
\item Asigne el rol a una cuenta en Usuarios y compruebe el recorrido autenticado. Si se retiró un permiso, verifique también que la operación correspondiente se deniega.
\end{enumerate}

\Needspace{8\baselineskip}
\begingroup
\setbox0=\vbox{\hsize=\linewidth
{\normalfont\normalsize\sffamily\bfseries\normalsize\color{Teal}\hyphenpenalty=10000\exhyphenpenalty=10000 Campos y botones\par}

\par
\begingroup
\fontsize{10}{12.5}\selectfont
\setlength{\tabcolsep}{5pt}
\renewcommand{\arraystretch}{1.15}
\begin{tabular}{@{}>{\RaggedRight\arraybackslash}p{0.22\linewidth}>{\RaggedRight\arraybackslash}p{0.3\linewidth}>{\RaggedRight\arraybackslash}p{0.39\linewidth}@{}}
\toprule
\textbf{Control} & \textbf{Para qué sirve} & \textbf{Ejemplo y resultado} \\
\midrule
Nombre & Identifica el perfil en Usuarios y Roles. & \textbf{Ejemplo:} Gerente comercial.\par\smallskip\textbf{Resultado:} Produce una etiqueta legible y una clave interna generada. \\
\bottomrule
\end{tabular}
\endgroup
}
\Needspace{\dimexpr\ht0+\dp0+48pt\relax}
\endgroup
\subsection{Campos y botones}
\begingroup
\fontsize{10}{12.5}\selectfont
\setlength{\tabcolsep}{5pt}
\renewcommand{\arraystretch}{1.15}
\begin{longtable}{@{}>{\RaggedRight\arraybackslash}p{0.22\linewidth}>{\RaggedRight\arraybackslash}p{0.3\linewidth}>{\RaggedRight\arraybackslash}p{0.39\linewidth}@{}}
\toprule
\textbf{Control} & \textbf{Para qué sirve} & \textbf{Ejemplo y resultado} \\
\midrule
\endfirsthead
\toprule
\textbf{Control} & \textbf{Para qué sirve} & \textbf{Ejemplo y resultado} \\
\midrule
\endhead
\midrule
\multicolumn{3}{r}{\footnotesize Continúa en la página siguiente}\\
\endfoot
\bottomrule
\endlastfoot
Nombre & Identifica el perfil en Usuarios y Roles. & \textbf{Ejemplo:} Gerente comercial.\par\smallskip\textbf{Resultado:} Produce una etiqueta legible y una clave interna generada. \\
\addlinespace[3pt]
Descripción & Explica qué puede hacer ese perfil. & \textbf{Ejemplo:} Consulta resultados y exporta informes, sin ejecutar ETL.\par\smallskip\textbf{Resultado:} Facilita una asignación de cuentas con criterio. \\
\addlinespace[3pt]
Permisos otorgados & Define el alcance real de las acciones. & \textbf{Ejemplo:} Consultar análisis visual y Exportar reportes analíticos.\par\smallskip\textbf{Resultado:} El servidor utiliza estos permisos al autorizar. \\
\addlinespace[3pt]
Editar & Abre el perfil concreto. & \textbf{Ejemplo:} Revisar un rol existente antes de agregar otro similar.\par\smallskip\textbf{Resultado:} Muestra sus datos y permisos actuales. \\
\addlinespace[3pt]
Desactivar / Activar & Cambia el uso de un rol no protegido. & \textbf{Ejemplo:} Retirar temporalmente un perfil de capacitación.\par\smallskip\textbf{Resultado:} Los permisos de un rol inactivo no aportan acceso efectivo. \\
\addlinespace[3pt]
Protegido & Identifica controles especiales del registro. & \textbf{Ejemplo:} Rol Administrador inicial.\par\smallskip\textbf{Resultado:} No equivale a una acción que deba pulsarse ni a un error. \\
\addlinespace[3pt]
Eliminar & Borra físicamente un rol no protegido que no conserve asignaciones a usuarios ni permisos. Aparece en las filas no protegidas cuando la cuenta puede administrar roles. & \textbf{Ejemplo:} Un rol de práctica vacío y sin asignaciones; es un ejemplo propuesto, no un borrado realizado durante esta edición.\par\smallskip\textbf{Resultado:} Solicita confirmación: «Eliminar definitivamente… Esta acción no se puede deshacer». Cancelar conserva el registro. Tras confirmar, el servidor comprueba permisos, protección y dependencias; si autoriza el borrado, la fila desaparece al actualizarse la lista y se registra el evento en Auditoría. \\
\addlinespace[3pt]
\end{longtable}
\endgroup

\subsection{Explicar un perfil de gerencia}
Práctica propuesta para el entorno de capacitación; no se presenta como una nueva ejecución realizada durante la redacción.

\begin{enumerate}
\item Consulte el rol de gerencia si existe o prepare su matriz antes de crearlo.
\item Identifique consulta del panel, exportación y consulta de metadatos.
\item Confirme que no requiere preparar materializaciones ni revisar propuestas.
\end{enumerate}

\textbf{Qué debe obtener.} Se distingue lectura analítica de escritura técnica. La estudiante explica la dependencia del selector de fuentes.

\textbf{Cómo comprobarlo.} Puede justificar cada permiso del perfil sin conceder administración general por comodidad.

\Needspace{8\baselineskip}
\begingroup
\setbox0=\vbox{\hsize=\linewidth
{\normalfont\normalsize\sffamily\bfseries\normalsize\color{Teal}\hyphenpenalty=10000\exhyphenpenalty=10000 Si aparece un problema\par}

\par
\begingroup
\fontsize{10}{12.5}\selectfont
\setlength{\tabcolsep}{5pt}
\renewcommand{\arraystretch}{1.15}
\begin{tabular}{@{}>{\RaggedRight\arraybackslash}p{0.22\linewidth}>{\RaggedRight\arraybackslash}p{0.3\linewidth}>{\RaggedRight\arraybackslash}p{0.39\linewidth}@{}}
\toprule
\textbf{Qué observa} & \textbf{Causa posible} & \textbf{Qué hacer y comprobar} \\
\midrule
La tabla de permisos no carga. & Falta Consultar permisos o hay un fallo de red. & Revise el permiso complementario y pulse Reintentar carga de opciones.\par\smallskip\textbf{Comprobación:} No queda habilitado el guardado de una selección vacía por error. \\
\bottomrule
\end{tabular}
\endgroup
}
\Needspace{\dimexpr\ht0+\dp0+48pt\relax}
\endgroup
\subsection{Si aparece un problema}
\begingroup
\fontsize{10}{12.5}\selectfont
\setlength{\tabcolsep}{5pt}
\renewcommand{\arraystretch}{1.15}
\begin{longtable}{@{}>{\RaggedRight\arraybackslash}p{0.22\linewidth}>{\RaggedRight\arraybackslash}p{0.3\linewidth}>{\RaggedRight\arraybackslash}p{0.39\linewidth}@{}}
\toprule
\textbf{Qué observa} & \textbf{Causa posible} & \textbf{Qué hacer y comprobar} \\
\midrule
\endfirsthead
\toprule
\textbf{Qué observa} & \textbf{Causa posible} & \textbf{Qué hacer y comprobar} \\
\midrule
\endhead
\midrule
\multicolumn{3}{r}{\footnotesize Continúa en la página siguiente}\\
\endfoot
\bottomrule
\endlastfoot
La tabla de permisos no carga. & Falta Consultar permisos o hay un fallo de red. & Revise el permiso complementario y pulse Reintentar carga de opciones.\par\smallskip\textbf{Comprobación:} No queda habilitado el guardado de una selección vacía por error. \\
\addlinespace[3pt]
El rol parece correcto, pero la cuenta no accede. & El rol o la cuenta están inactivos, falta la asignación o no se concedió una dependencia de interfaz. & Revise estado, asignación y permisos; autentique otra vez para actualizar la navegación.\par\smallskip\textbf{Comprobación:} La operación se autoriza en el servidor y la navegación corresponde a la cuenta. \\
\addlinespace[3pt]
No se permite retirar ciertos permisos del administrador. & La modificación dañaría la recuperación del entorno. & Mantenga el rol protegido y cree un perfil separado para la práctica.\par\smallskip\textbf{Comprobación:} Permanece una administración activa y recuperable. \\
\addlinespace[3pt]
Eliminar informa que primero se resuelvan asignaciones a usuarios o de permisos. & El rol tiene relaciones persistidas. No basta con que esté inactivo ni con que no haya usuarios activos. & Lea la dependencia indicada y decida con la responsable de acceso si corresponde conservar, desactivar o retirar deliberadamente las asignaciones. No vacíe un perfil en uso sólo para completar una práctica. Si se autoriza su retiro, revise primero las cuentas afectadas.\par\smallskip\textbf{Comprobación:} Un rechazo conserva el rol y sus relaciones. Un borrado autorizado elimina la fila de la lista y deja en Auditoría el evento de eliminación de rol. \\
\addlinespace[3pt]
No aparece Eliminar o el servidor informa «El rol protegido no puede eliminarse». & Es el perfil administrativo protegido, o faltan permisos para administrar roles. & Conserve el rol protegido y use un perfil de práctica no protegido para estudiar sus opciones. La interfaz muestra Protegido y oculta Eliminar en ese caso.\par\smallskip\textbf{Comprobación:} El perfil de recuperación administrativo sigue disponible. La protección se comprueba también en el servidor. \\
\addlinespace[3pt]
\end{longtable}
\endgroup

\subsection{Para explicar durante la sustentación}
\paragraph{¿Scrum Master es un rol de la plataforma?} No. Es una responsabilidad del proyecto de desarrollo. Product Owner y Scrum Master no conceden permisos de la aplicación.

\paragraph{¿Es lo mismo ocultar ETL que denegar ETL?} No. El menú refleja permisos, pero la autorización de la operación se comprueba también en el servidor.

\paragraph{¿Eliminar un rol equivale a desactivarlo y se puede deshacer?} No. Desactivar conserva la definición. Eliminar ejecuta un borrado físico, después de confirmar y sólo si el rol no es protegido y no tiene asignaciones de usuarios o permisos. No existe una papelera ni un botón de restauración. Crear otro perfil con el mismo nombre produce otro registro y exige volver a configurar sus permisos y asignaciones.

\pantalla{recursos/pantallas/45_roles.jpg}{Perfiles disponibles y acciones de administración}{Roles registrados. El perfil funcional Gerente comercial se diferencia del papel académico de Product Owner o Scrum Master. Captura de referencia del 18 de noviembre de 2026; revise las instrucciones actualizadas de esta unidad.}
\pantalla{recursos/pantallas/46_rol_gerencia.jpg}{Permisos del rol Gerente comercial: consulta sin guardar}{Consulta del editor de Gerente comercial, cancelada sin guardar. El rol tiene lectura analítica, conexiones, metadatos y exportación. Captura de referencia del 18 de noviembre de 2026; revise las instrucciones actualizadas de esta unidad.}

\IfFileExists{capturas/operacion_05.tex}{\input{capturas/operacion_05.tex}}{}
\section{Interpretar el catálogo de permisos}
\label{unidad:operacion_06}
\textbf{Para quién:} Administración, Personas autorizadas para consulta\par
\textbf{Opción:} Administración y seguridad > Permisos\par

Permisos es el diccionario de operaciones que los módulos conocen. Consultar propuestas BI, Generar propuestas BI y Revisar propuestas BI son capacidades distintas: consultar una versión no permite generarla ni aprobarla. La separación evita conceder escritura por el simple hecho de necesitar leer un expediente.

La pantalla vigente es informativa. Presenta nombre, descripción y estado en un listado paginado; no ofrece crear un permiso libre. Escribir un nombre como Administrar inventario no implementaría inventario ni protegería una operación real. Los permisos técnicos se incorporan cuando se desarrolla la funcionalidad correspondiente.

Algunas descripciones históricas del catálogo utilizan verbos como crear permisos o menús, pero la operación actual de la pantalla se rige por sus controles y por el servidor. Para asignar privilegios se trabaja en Roles; para conocer el significado de un permiso se consulta esta lista.

\subsection{Conceptos para comprender la opción}
\paragraph{Permiso} Autorización concreta que protege una operación implementada.

\paragraph{Permiso de lectura} Permite consultar sin modificar el recurso.

\paragraph{Permiso declarado} Entrada del catálogo; no demuestra por sí misma que exista una nueva función utilizable.

\subsection{Paso a paso}
\begin{enumerate}
\item Abra Permisos.
\item Lea el nombre y la descripción de las operaciones que correspondan a su tarea.
\item Use Anterior o Siguiente si el listado tiene más de una página.
\item Regrese a Roles para conceder o retirar esos permisos del perfil correspondiente.
\item Compruebe el resultado con una cuenta de ese perfil; no modifique códigos internos ni datos directamente.
\end{enumerate}

\Needspace{8\baselineskip}
\begingroup
\setbox0=\vbox{\hsize=\linewidth
{\normalfont\normalsize\sffamily\bfseries\normalsize\color{Teal}\hyphenpenalty=10000\exhyphenpenalty=10000 Campos y botones\par}

\par
\begingroup
\fontsize{10}{12.5}\selectfont
\setlength{\tabcolsep}{5pt}
\renewcommand{\arraystretch}{1.15}
\begin{tabular}{@{}>{\RaggedRight\arraybackslash}p{0.22\linewidth}>{\RaggedRight\arraybackslash}p{0.3\linewidth}>{\RaggedRight\arraybackslash}p{0.39\linewidth}@{}}
\toprule
\textbf{Control} & \textbf{Para qué sirve} & \textbf{Ejemplo y resultado} \\
\midrule
Elemento & Presenta el nombre del permiso. & \textbf{Ejemplo:} Revisar propuestas BI.\par\smallskip\textbf{Resultado:} Permite identificar una capacidad con lenguaje comprensible. \\
\bottomrule
\end{tabular}
\endgroup
}
\Needspace{\dimexpr\ht0+\dp0+48pt\relax}
\endgroup
\subsection{Campos y botones}
\begingroup
\fontsize{10}{12.5}\selectfont
\setlength{\tabcolsep}{5pt}
\renewcommand{\arraystretch}{1.15}
\begin{longtable}{@{}>{\RaggedRight\arraybackslash}p{0.22\linewidth}>{\RaggedRight\arraybackslash}p{0.3\linewidth}>{\RaggedRight\arraybackslash}p{0.39\linewidth}@{}}
\toprule
\textbf{Control} & \textbf{Para qué sirve} & \textbf{Ejemplo y resultado} \\
\midrule
\endfirsthead
\toprule
\textbf{Control} & \textbf{Para qué sirve} & \textbf{Ejemplo y resultado} \\
\midrule
\endhead
\midrule
\multicolumn{3}{r}{\footnotesize Continúa en la página siguiente}\\
\endfoot
\bottomrule
\endlastfoot
Elemento & Presenta el nombre del permiso. & \textbf{Ejemplo:} Revisar propuestas BI.\par\smallskip\textbf{Resultado:} Permite identificar una capacidad con lenguaje comprensible. \\
\addlinespace[3pt]
Detalle & Describe el alcance funcional. & \textbf{Ejemplo:} Aprobar o rechazar propuestas validadas.\par\smallskip\textbf{Resultado:} Ayuda a diferenciar revisión de mera consulta. \\
\addlinespace[3pt]
Estado & Indica si el permiso está activo. & \textbf{Ejemplo:} true en el catálogo aprobado.\par\smallskip\textbf{Resultado:} Un permiso inactivo no aporta acceso efectivo. \\
\addlinespace[3pt]
Anterior / Siguiente & Recorre el catálogo vigente. & \textbf{Ejemplo:} Consultar la segunda página de los 29 permisos base.\par\smallskip\textbf{Resultado:} Muestra otro tramo del listado sin crear registros. \\
\addlinespace[3pt]
\end{longtable}
\endgroup

\subsection{Separar tres permisos del asistente}
Práctica propuesta para el entorno de capacitación; no se presenta como una nueva ejecución realizada durante la redacción.

\begin{enumerate}
\item Localice Consultar propuestas BI, Generar propuestas BI y Revisar propuestas BI.
\item Asocie a cada uno una tarea.
\item Explique cuál requiere un analista revisor y cuál puede necesitar una persona que solo consulta.
\end{enumerate}

\textbf{Qué debe obtener.} Las tareas se vinculan con permisos diferentes.

\textbf{Cómo comprobarlo.} La estudiante no atribuye aprobación al permiso de lectura.

\Needspace{8\baselineskip}
\begingroup
\setbox0=\vbox{\hsize=\linewidth
{\normalfont\normalsize\sffamily\bfseries\normalsize\color{Teal}\hyphenpenalty=10000\exhyphenpenalty=10000 Si aparece un problema\par}

\par
\begingroup
\fontsize{10}{12.5}\selectfont
\setlength{\tabcolsep}{5pt}
\renewcommand{\arraystretch}{1.15}
\begin{tabular}{@{}>{\RaggedRight\arraybackslash}p{0.22\linewidth}>{\RaggedRight\arraybackslash}p{0.3\linewidth}>{\RaggedRight\arraybackslash}p{0.39\linewidth}@{}}
\toprule
\textbf{Qué observa} & \textbf{Causa posible} & \textbf{Qué hacer y comprobar} \\
\midrule
No existe el permiso que se desea escribir. & La capacidad no está implementada o se está buscando por un nombre diferente. & Revise el catálogo por su descripción y consulte la especificación; no intente crear códigos desde el navegador.\par\smallskip\textbf{Comprobación:} Se asigna una capacidad existente o se reconoce correctamente una ampliación pendiente de desarrollo. \\
\bottomrule
\end{tabular}
\endgroup
}
\Needspace{\dimexpr\ht0+\dp0+48pt\relax}
\endgroup
\subsection{Si aparece un problema}
\begingroup
\fontsize{10}{12.5}\selectfont
\setlength{\tabcolsep}{5pt}
\renewcommand{\arraystretch}{1.15}
\begin{longtable}{@{}>{\RaggedRight\arraybackslash}p{0.22\linewidth}>{\RaggedRight\arraybackslash}p{0.3\linewidth}>{\RaggedRight\arraybackslash}p{0.39\linewidth}@{}}
\toprule
\textbf{Qué observa} & \textbf{Causa posible} & \textbf{Qué hacer y comprobar} \\
\midrule
\endfirsthead
\toprule
\textbf{Qué observa} & \textbf{Causa posible} & \textbf{Qué hacer y comprobar} \\
\midrule
\endhead
\midrule
\multicolumn{3}{r}{\footnotesize Continúa en la página siguiente}\\
\endfoot
\bottomrule
\endlastfoot
No existe el permiso que se desea escribir. & La capacidad no está implementada o se está buscando por un nombre diferente. & Revise el catálogo por su descripción y consulte la especificación; no intente crear códigos desde el navegador.\par\smallskip\textbf{Comprobación:} Se asigna una capacidad existente o se reconoce correctamente una ampliación pendiente de desarrollo. \\
\addlinespace[3pt]
\end{longtable}
\endgroup

\subsection{Para explicar durante la sustentación}
\paragraph{¿Agregar un permiso implementa una función?} No. Un permiso protege una operación existente. La funcionalidad requiere código, controles y pruebas propios.

\pantalla{recursos/pantallas/47_permisos.jpg}{Permisos técnicos registrados}{El catálogo de permisos explica las capacidades existentes. Crear un nombre de permiso no implementa una función nueva. Captura de referencia del 18 de noviembre de 2026; revise las instrucciones actualizadas de esta unidad.}

\IfFileExists{capturas/operacion_06.tex}{\input{capturas/operacion_06.tex}}{}
\section{Ajustar los nombres y el orden de los menús}
\label{unidad:operacion_07}
\textbf{Para quién:} Administración\par
\textbf{Opción:} Administración y seguridad > Menús\par

Un menú es la entrada visible a una pantalla ya implementada. Cambiar su etiqueta ayuda a la navegación, pero no crea un nuevo módulo ni cambia el contrato de la operación. La aplicación conserva las rutas y claves internas; el formulario permite ajustar Etiqueta visible, Posición y Permisos requeridos.

Primero seleccione Editar en un menú registrado. No hay una operación normal para inventar una ruta o agregar una pantalla desde esta lista. En el catálogo base, los menús están protegidos: puede haber datos de presentación editables, pero no se permite eliminarlos o desactivarlos como si fueran registros de práctica.

Los grupos de navegación tienen un orden común y las opciones se organizan dentro de ellos. Un número de posición menor sitúa antes una opción dentro del orden aplicable. Con varios permisos asociados, el servidor de navegación muestra actualmente el menú si la cuenta posee al menos uno de ellos; cada acción sigue exigiendo su permiso específico. Por ello, la configuración del menú no sustituye al control de acceso.

\subsection{Conceptos para comprender la opción}
\paragraph{Menú} Acceso de navegación a una pantalla registrada.

\paragraph{Etiqueta} Nombre que lee la persona usuaria; es distinto de la ruta técnica.

\paragraph{Posición} Valor de orden; no es un identificador de permiso ni una prioridad comercial.

\subsection{Paso a paso}
\begin{enumerate}
\item Abra Menús y localice la entrada que necesita ajustar.
\item Pulse Editar en esa fila.
\item Cambie Etiqueta visible solo si mejora la comprensión y mantiene la función real.
\item Ajuste Posición y revise Permisos requeridos.
\item Pulse Guardar cambios o Cancelar.
\item Actualice la sesión para volver a consultar la navegación y compruebe la etiqueta con la cuenta destinataria.
\end{enumerate}

\Needspace{8\baselineskip}
\begingroup
\setbox0=\vbox{\hsize=\linewidth
{\normalfont\normalsize\sffamily\bfseries\normalsize\color{Teal}\hyphenpenalty=10000\exhyphenpenalty=10000 Campos y botones\par}

\par
\begingroup
\fontsize{10}{12.5}\selectfont
\setlength{\tabcolsep}{5pt}
\renewcommand{\arraystretch}{1.15}
\begin{tabular}{@{}>{\RaggedRight\arraybackslash}p{0.22\linewidth}>{\RaggedRight\arraybackslash}p{0.3\linewidth}>{\RaggedRight\arraybackslash}p{0.39\linewidth}@{}}
\toprule
\textbf{Control} & \textbf{Para qué sirve} & \textbf{Ejemplo y resultado} \\
\midrule
Etiqueta visible & Define el nombre de navegación. & \textbf{Ejemplo:} Conservar Analítica de ventas como nombre preciso.\par\smallskip\textbf{Resultado:} Modifica la presentación del menú registrado. \\
\bottomrule
\end{tabular}
\endgroup
}
\Needspace{\dimexpr\ht0+\dp0+48pt\relax}
\endgroup
\subsection{Campos y botones}
\begingroup
\fontsize{10}{12.5}\selectfont
\setlength{\tabcolsep}{5pt}
\renewcommand{\arraystretch}{1.15}
\begin{longtable}{@{}>{\RaggedRight\arraybackslash}p{0.22\linewidth}>{\RaggedRight\arraybackslash}p{0.3\linewidth}>{\RaggedRight\arraybackslash}p{0.39\linewidth}@{}}
\toprule
\textbf{Control} & \textbf{Para qué sirve} & \textbf{Ejemplo y resultado} \\
\midrule
\endfirsthead
\toprule
\textbf{Control} & \textbf{Para qué sirve} & \textbf{Ejemplo y resultado} \\
\midrule
\endhead
\midrule
\multicolumn{3}{r}{\footnotesize Continúa en la página siguiente}\\
\endfoot
\bottomrule
\endlastfoot
Etiqueta visible & Define el nombre de navegación. & \textbf{Ejemplo:} Conservar Analítica de ventas como nombre preciso.\par\smallskip\textbf{Resultado:} Modifica la presentación del menú registrado. \\
\addlinespace[3pt]
Posición & Ordena la entrada. & \textbf{Ejemplo:} Una posición anterior dentro de su grupo.\par\smallskip\textbf{Resultado:} Reordena la navegación aplicable, no concede permisos. \\
\addlinespace[3pt]
Permisos requeridos & Asocia condiciones de visibilidad del menú. & \textbf{Ejemplo:} Consultar análisis visual para la entrada del panel.\par\smallskip\textbf{Resultado:} Influye en la navegación; el servidor sigue validando la operación. \\
\addlinespace[3pt]
Guardar cambios / Cancelar & Confirma o abandona la edición. & \textbf{Ejemplo:} Ajustar una etiqueta sin cambiar una ruta.\par\smallskip\textbf{Resultado:} El registro se actualiza o se conserva como estaba. \\
\addlinespace[3pt]
\end{longtable}
\endgroup

\subsection{Reconocer menú y operación}
Práctica propuesta para el entorno de capacitación; no se presenta como una nueva ejecución realizada durante la redacción.

\begin{enumerate}
\item Consulte la entrada Datamart de ventas y la entrada Analítica de ventas.
\item Identifique sus funciones y los permisos de consulta asociados.
\item Explique por qué abrir el panel no permite ejecutar una carga.
\end{enumerate}

\textbf{Qué debe obtener.} Se distingue navegación, consulta y ejecución.

\textbf{Cómo comprobarlo.} Puede justificar la separación sin intentar crear rutas nuevas.

\Needspace{8\baselineskip}
\begingroup
\setbox0=\vbox{\hsize=\linewidth
{\normalfont\normalsize\sffamily\bfseries\normalsize\color{Teal}\hyphenpenalty=10000\exhyphenpenalty=10000 Si aparece un problema\par}

\par
\begingroup
\fontsize{10}{12.5}\selectfont
\setlength{\tabcolsep}{5pt}
\renewcommand{\arraystretch}{1.15}
\begin{tabular}{@{}>{\RaggedRight\arraybackslash}p{0.22\linewidth}>{\RaggedRight\arraybackslash}p{0.3\linewidth}>{\RaggedRight\arraybackslash}p{0.39\linewidth}@{}}
\toprule
\textbf{Qué observa} & \textbf{Causa posible} & \textbf{Qué hacer y comprobar} \\
\midrule
No aparece Crear menú. & La UI solo edita pantallas ya registradas. & Use Editar sobre una entrada existente. Para otra pantalla se necesita desarrollo e incorporación controlada.\par\smallskip\textbf{Comprobación:} No se presenta una ruta inventada como una función operativa. \\
\bottomrule
\end{tabular}
\endgroup
}
\Needspace{\dimexpr\ht0+\dp0+48pt\relax}
\endgroup
\subsection{Si aparece un problema}
\begingroup
\fontsize{10}{12.5}\selectfont
\setlength{\tabcolsep}{5pt}
\renewcommand{\arraystretch}{1.15}
\begin{longtable}{@{}>{\RaggedRight\arraybackslash}p{0.22\linewidth}>{\RaggedRight\arraybackslash}p{0.3\linewidth}>{\RaggedRight\arraybackslash}p{0.39\linewidth}@{}}
\toprule
\textbf{Qué observa} & \textbf{Causa posible} & \textbf{Qué hacer y comprobar} \\
\midrule
\endfirsthead
\toprule
\textbf{Qué observa} & \textbf{Causa posible} & \textbf{Qué hacer y comprobar} \\
\midrule
\endhead
\midrule
\multicolumn{3}{r}{\footnotesize Continúa en la página siguiente}\\
\endfoot
\bottomrule
\endlastfoot
No aparece Crear menú. & La UI solo edita pantallas ya registradas. & Use Editar sobre una entrada existente. Para otra pantalla se necesita desarrollo e incorporación controlada.\par\smallskip\textbf{Comprobación:} No se presenta una ruta inventada como una función operativa. \\
\addlinespace[3pt]
La etiqueta cambió, pero la cuenta no ve el menú. & La sesión conserva su navegación anterior o faltan los permisos efectivos. & Vuelva a autenticar y revise el rol; no retire la autorización del servidor para forzar la visualización.\par\smallskip\textbf{Comprobación:} La navegación actual refleja los permisos y la etiqueta guardada. \\
\addlinespace[3pt]
\end{longtable}
\endgroup

\subsection{Para explicar durante la sustentación}
\paragraph{¿Puedo agregar inventario escribiéndolo en Menús?} No. La pantalla, las reglas de inventario y sus pruebas deben desarrollarse antes de registrar una entrada operativa.

\pantalla{recursos/pantallas/48_menus.jpg}{Listado de menús y rutas protegidas}{Navegación registrada y protegida. Cambiar una etiqueta visible no cambia la ruta interna ni la autorización del backend. Captura de referencia del 18 de noviembre de 2026; revise las instrucciones actualizadas de esta unidad.}
\pantalla{recursos/pantallas/49_menu_edicion.jpg}{Edición de navegación sin guardar cambios}{Editor de menú abierto y cancelado. Etiqueta, posición y permisos afectan a la navegación de los usuarios. Captura de referencia del 18 de noviembre de 2026; revise las instrucciones actualizadas de esta unidad.}

\IfFileExists{capturas/operacion_07.tex}{\input{capturas/operacion_07.tex}}{}
\begingroup
\setlength{\parskip}{4pt plus 1pt minus 1pt}
\setlist[enumerate]{itemsep=3pt,topsep=4pt}
\titlespacing*{\paragraph}{0pt}{4pt}{0.5em}
\section{Consultar la auditoría y reportar una incidencia}
\label{unidad:operacion_08}
\textbf{Para quién:} Administración, Personas con Consultar auditoría\par
\textbf{Opción:} Administración y seguridad > Auditoría\par

La auditoría registra acciones críticas y ayuda a reconstruir quién actuó sobre un recurso. Sirve para investigar una modificación o acompañar un expediente. No reemplaza la conciliación de importes: un evento de ejecución informa que hubo una acción, mientras que el expediente muestra las comprobaciones y sus resultados.

La interfaz actual presenta una lista paginada de eventos recientes, con acción, actor y tipo de recurso. Algunas acciones aparecen con claves técnicas, como copilot.catalog.update o etl.execution.run. Deben interpretarse como actualización del catálogo y ejecución de carga, respectivamente. No son instrucciones que la estudiante deba ejecutar.

Esta pantalla es de consulta. No se crean, corrigen ni eliminan eventos para que una demostración parezca exitosa. Si una cuenta fue retirada, la evidencia puede conservar su etiqueta histórica. Para un reporte útil, acompañe la pantalla con la fuente, la propuesta o ejecución afectada y los pasos que reprodujeron la incidencia.

\subsection{Conceptos para comprender la opción}
\paragraph{Evento de auditoría} Registro de una acción realizada sobre un recurso.

\paragraph{Actor} Cuenta atribuida a la acción.

\paragraph{Recurso} Objeto sobre el que se actuó. Una propuesta, un catálogo, una conexión o una ejecución.

\subsection{Paso a paso}
\begin{enumerate}
\item Abra Auditoría.
\item Lea los eventos más recientes y distinga acción, actor y recurso.
\item Use Siguiente o Anterior para revisar eventos anteriores.
\item Si investiga un caso, abra por separado su expediente para comprobar estado y cifras.
\item Describa la incidencia con pasos, fuente e identificadores; adjunte la captura sin claves ni tokens.
\end{enumerate}

\Needspace{8\baselineskip}
\begingroup
\setbox0=\vbox{\hsize=\linewidth
{\normalfont\normalsize\sffamily\bfseries\normalsize\color{Teal}\hyphenpenalty=10000\exhyphenpenalty=10000 Campos y botones\par}

\par
\begingroup
\fontsize{10}{12.5}\selectfont
\setlength{\tabcolsep}{5pt}
\renewcommand{\arraystretch}{1.15}
\begin{tabular}{@{}>{\RaggedRight\arraybackslash}p{0.22\linewidth}>{\RaggedRight\arraybackslash}p{0.3\linewidth}>{\RaggedRight\arraybackslash}p{0.39\linewidth}@{}}
\toprule
\textbf{Control} & \textbf{Para qué sirve} & \textbf{Ejemplo y resultado} \\
\midrule
Elemento & Muestra la acción registrada. & \textbf{Ejemplo:} Inicio de sesión o copilot.catalog.update.\par\smallskip\textbf{Resultado:} Identifica qué tipo de operación ocurrió. \\
\bottomrule
\end{tabular}
\endgroup
}
\Needspace{\dimexpr\ht0+\dp0+48pt\relax}
\endgroup
\subsection{Campos y botones}
\begingroup
\fontsize{10}{12.5}\selectfont
\setlength{\tabcolsep}{5pt}
\renewcommand{\arraystretch}{1.15}
\begin{longtable}{@{}>{\RaggedRight\arraybackslash}p{0.22\linewidth}>{\RaggedRight\arraybackslash}p{0.3\linewidth}>{\RaggedRight\arraybackslash}p{0.39\linewidth}@{}}
\toprule
\textbf{Control} & \textbf{Para qué sirve} & \textbf{Ejemplo y resultado} \\
\midrule
\endfirsthead
\toprule
\textbf{Control} & \textbf{Para qué sirve} & \textbf{Ejemplo y resultado} \\
\midrule
\endhead
\midrule
\multicolumn{3}{r}{\footnotesize Continúa en la página siguiente}\\
\endfoot
\bottomrule
\endlastfoot
Elemento & Muestra la acción registrada. & \textbf{Ejemplo:} Inicio de sesión o copilot.catalog.update.\par\smallskip\textbf{Resultado:} Identifica qué tipo de operación ocurrió. \\
\addlinespace[3pt]
Detalle & Permite reconocer al actor histórico. & \textbf{Ejemplo:} Nombre y correo de la cuenta de práctica.\par\smallskip\textbf{Resultado:} Ayuda a atribuir la acción sin mostrar credenciales. \\
\addlinespace[3pt]
Estado & En esta tabla puede representar el tipo de recurso. & \textbf{Ejemplo:} etl\_\allowbreak{}execution.\par\smallskip\textbf{Resultado:} No debe interpretarse automáticamente como éxito o fracaso del ETL. \\
\addlinespace[3pt]
Anterior / Siguiente & Recorre los eventos. & \textbf{Ejemplo:} Revisar operaciones anteriores de una campaña.\par\smallskip\textbf{Resultado:} Cambia la página de consulta; no filtra por actor automáticamente. \\
\addlinespace[3pt]
\end{longtable}
\endgroup

\subsection{Relacionar un evento con su expediente}
Práctica propuesta para el entorno de capacitación; no se presenta como una nueva ejecución realizada durante la redacción.

\begin{enumerate}
\item Consulte un evento de ejecución existente.
\item Identifique su tipo de recurso.
\item Abra el expediente correspondiente desde Datamart de ventas y revise conciliación.
\end{enumerate}

\textbf{Qué debe obtener.} Se usan dos evidencias complementarias: acción auditada y resultado cuantitativo.

\textbf{Cómo comprobarlo.} La estudiante no utiliza la mera presencia del evento como prueba de igualdad de importes.

\Needspace{8\baselineskip}
\begingroup
\setbox0=\vbox{\hsize=\linewidth
{\normalfont\normalsize\sffamily\bfseries\normalsize\color{Teal}\hyphenpenalty=10000\exhyphenpenalty=10000 Si aparece un problema\par}

\par
\begingroup
\fontsize{10}{12.5}\selectfont
\setlength{\tabcolsep}{5pt}
\renewcommand{\arraystretch}{1.15}
\begin{tabular}{@{}>{\RaggedRight\arraybackslash}p{0.22\linewidth}>{\RaggedRight\arraybackslash}p{0.3\linewidth}>{\RaggedRight\arraybackslash}p{0.39\linewidth}@{}}
\toprule
\textbf{Qué observa} & \textbf{Causa posible} & \textbf{Qué hacer y comprobar} \\
\midrule
No encuentra filtros por fecha o actor. & La UI de este corte solo ofrece paginación. & Revise las páginas disponibles o solicite apoyo técnico para una consulta autorizada. No siga instrucciones de un manual anterior que anuncien filtros inexistentes.\par\smallskip\textbf{Comprobación:} La investigación documenta el alcance real de la consulta. \\
\bottomrule
\end{tabular}
\endgroup
}
\Needspace{\dimexpr\ht0+\dp0+48pt\relax}
\endgroup
\subsection{Si aparece un problema}
\begingroup
\fontsize{10}{12.5}\selectfont
\setlength{\tabcolsep}{5pt}
\renewcommand{\arraystretch}{1.15}
\begin{longtable}{@{}>{\RaggedRight\arraybackslash}p{0.22\linewidth}>{\RaggedRight\arraybackslash}p{0.3\linewidth}>{\RaggedRight\arraybackslash}p{0.39\linewidth}@{}}
\toprule
\textbf{Qué observa} & \textbf{Causa posible} & \textbf{Qué hacer y comprobar} \\
\midrule
\endfirsthead
\toprule
\textbf{Qué observa} & \textbf{Causa posible} & \textbf{Qué hacer y comprobar} \\
\midrule
\endhead
\midrule
\multicolumn{3}{r}{\footnotesize Continúa en la página siguiente}\\
\endfoot
\bottomrule
\endlastfoot
No encuentra filtros por fecha o actor. & La UI de este corte solo ofrece paginación. & Revise las páginas disponibles o solicite apoyo técnico para una consulta autorizada. No siga instrucciones de un manual anterior que anuncien filtros inexistentes.\par\smallskip\textbf{Comprobación:} La investigación documenta el alcance real de la consulta. \\
\addlinespace[3pt]
Aparece una acción técnica difícil de comprender. & El evento conserva el nombre interno de la operación. & Asóciela con el módulo y consulte el expediente; registre la clave literal si necesita reportar una incidencia.\par\smallskip\textbf{Comprobación:} El reporte identifica qué se hizo, sin inventar su resultado. \\
\addlinespace[3pt]
\end{longtable}
\endgroup

\subsection{Para explicar durante la sustentación}
\paragraph{¿Se puede borrar un error de la auditoría?} No desde la interfaz. La auditoría conserva acciones y aporta trazabilidad; la solución del problema se registra mediante nuevas acciones válidas.

\endgroup
\pantalla{recursos/pantallas/50_auditoria.jpg}{Auditoría consultable de acciones críticas}{La auditoría registra acciones críticas. La pantalla es de consulta y no ofrece edición de los eventos. Captura de referencia del 18 de noviembre de 2026; revise las instrucciones actualizadas de esta unidad.}

\IfFileExists{capturas/operacion_08.tex}{\input{capturas/operacion_08.tex}}{}

\chapter{Practicar, resolver incidencias y preparar la defensa}
\section{Practicar con las dos fuentes sin memorizar una respuesta del modelo}
\label{unidad:PR01}
\textbf{Para quién:} Analista de BI, Gerente comercial\par
\textbf{Al terminar podrá:} Repetir un recorrido básico y uno avanzado en cada fuente con criterios de aceptación observables.

Las cuatro fichas siguientes permiten aprender desde un alcance sencillo hasta documentos distintos, razones y roles de relación. Las básicas son ejercicios propuestos. Las avanzadas conservan objetivos y resultados de la campaña comprobada del 6 de noviembre de 2026.

Las fichas avanzadas conservan como historia la campaña del 6 de noviembre: propuestas 129 y 131, ejecuciones 21 y 22. Esta edición añade una campaña del 25 de noviembre, con las propuestas 136 y 135 y las ejecuciones 24 y 23, respectivamente. Una instalación limpia recibe sus propios identificadores: no tiene que reproducir esos números.

Si desea demostrar con un caso ya conservado, anuncie que consulta un expediente histórico. Si genera desde cero, diga que está realizando una ejecución nueva y registre los resultados que realmente obtenga.

La campaña nueva valida recorridos concretos con OpenAI Sol; las fichas básicas siguen siendo ejercicios para practicar. No equivale a garantizar cualquier texto, cualquier modelo, cualquier motor ni éxito idéntico en una solicitud futura. Aprenda a reconocer también un rechazo correcto y a recuperar una operación sin duplicarla.

\subsection{Ruta recomendada de estudio con los casos actuales}
Primero aprenda a consultar un expediente comprobado. Después repita su formulación en una propuesta nueva, sin exigir que el modelo use las mismas palabras o que el sistema asigne el mismo número. Evalúe el significado, los requisitos y la conciliación. Los casos históricos del 6 de noviembre se conservan a continuación como material de comparación.

\begin{enumerate}
\item Fase 1, lectura sin nueva generación: consulte 136/24 de AdventureWorks y explique grano, unidades, pedidos e importe.
\item Fase 2, contraste: consulte 135/23 de WideWorldImporters y explique la diferencia entre factura y pedido, y entre moneda de origen y USD.
\item Fase 3, navegación: reproduzca el cambio de fuentes, filtros y retorno al respaldo.
\item Fase 4, práctica nueva: escriba una necesidad breve con importe, unidades y documentos distintos por período y territorio. Confirme que el catálogo de esa fuente respalda el objetivo.
\item Fase 5, supervisión: revise viabilidad, conceptos, diseño, cobertura e indicadores. Apruebe únicamente la nueva versión que cumpla.
\item Fase 6, carga: prepare y confirme el ETL una vez; compruebe estado final, conciliación, publicación y panel. Registre los identificadores reales que obtenga.
\item Fase 7, explicación: formule una pregunta concreta al copiloto, contraste su evidencia y exporte el mismo contexto. Registre cualquier rechazo correcto o limitación.
\end{enumerate}

\Needspace{8\baselineskip}
\begingroup
\setbox0=\vbox{\hsize=\linewidth
{\normalfont\normalsize\sffamily\bfseries\normalsize\color{Teal}\hyphenpenalty=10000\exhyphenpenalty=10000 Preparación de los dos ejercicios vigentes\par}

\par
\begingroup
\fontsize{10}{12.5}\selectfont
\setlength{\tabcolsep}{5pt}
\renewcommand{\arraystretch}{1.15}
\begin{tabular}{@{}>{\RaggedRight\arraybackslash}p{0.22\linewidth}>{\RaggedRight\arraybackslash}p{0.3\linewidth}>{\RaggedRight\arraybackslash}p{0.39\linewidth}@{}}
\toprule
\textbf{Opción} & \textbf{AdventureWorks} & \textbf{WideWorldImporters} \\
\midrule
Periodicidad & Mensual & Mensual \\
\bottomrule
\end{tabular}
\endgroup
}
\Needspace{\dimexpr\ht0+\dp0+48pt\relax}
\endgroup
\subsection{Preparación de los dos ejercicios vigentes}
\begingroup
\fontsize{10}{12.5}\selectfont
\setlength{\tabcolsep}{5pt}
\renewcommand{\arraystretch}{1.15}
\begin{longtable}{@{}>{\RaggedRight\arraybackslash}p{0.22\linewidth}>{\RaggedRight\arraybackslash}p{0.3\linewidth}>{\RaggedRight\arraybackslash}p{0.39\linewidth}@{}}
\toprule
\textbf{Opción} & \textbf{AdventureWorks} & \textbf{WideWorldImporters} \\
\midrule
\endfirsthead
\toprule
\textbf{Opción} & \textbf{AdventureWorks} & \textbf{WideWorldImporters} \\
\midrule
\endhead
\midrule
\multicolumn{3}{r}{\footnotesize Continúa en la página siguiente}\\
\endfoot
\bottomrule
\endlastfoot
Periodicidad & Mensual & Mensual \\
\addlinespace[3pt]
Preguntas de negocio & Evolución temporal, productos, clientes y territorio. & Evolución temporal, productos, clientes y territorio. \\
\addlinespace[3pt]
Evento y granularidad & Una línea de pedido, conservando identificador del pedido. & Una línea de factura, conservando identificador de factura. \\
\addlinespace[3pt]
Indicadores imprescindibles & Importe total; unidades totales; pedidos distintos; importe promedio por pedido. & Importe total; unidades totales; facturas distintas; importe promedio por factura. \\
\addlinespace[3pt]
Límites del ejercicio & Sin costos, margen ni impuestos. USD sólo con evidencia de la fuente. & Sin costos, impuestos ni rentabilidad. Moneda de origen; ubicación actual del maestro del cliente. \\
\addlinespace[3pt]
\end{longtable}
\endgroup

Necesidad comprobada de AdventureWorks, para copiar y revisar: Analizar mensualmente el importe registrado y las unidades de las líneas de pedidos por producto, cliente y territorio comercial. Los indicadores deben incluir importe total, unidades totales, pedidos distintos e importe promedio por pedido. Usar la fecha del pedido y el territorio comercial asociado al pedido. Identificar los clientes con nombres descriptivos de personas u organizaciones mediante relaciones declaradas. Conservar el detalle de cada línea y la trazabilidad del pedido. No calcular costos, margen ni impuestos.

Necesidad comprobada de WideWorldImporters, para copiar y revisar: Analizar mensualmente el importe registrado y las unidades de las líneas de factura por producto y cliente, con el número de facturas distintas y el importe promedio por factura. Usar la fecha de emisión y, para el territorio, la provincia de la ciudad de entrega registrada del cliente. Mantener la trazabilidad de cada línea y aclarar que esa ubicación corresponde a la ficha actual del cliente. Conservar la moneda de origen sin asumir impuestos, costos ni rentabilidad.

Antes de aprobar, compruebe los cuatro indicadores. Las recetas esperadas son suma del importe, suma de unidades, conteo distinto del documento e importe dividido por documentos distintos. En la campaña vigente, la plataforma completó de manera trazable el KPI de unidades a partir de una receta inequívoca. No copie un número de advertencias como criterio de aprobación: lea y justifique las que presente su nueva versión.

Las rutas y los campos pueden cambiar en otra fuente. Estas fichas prueban dos aplicaciones del mismo proceso; no autorizan a fijar sus tablas, nombres ni fórmulas como lógica exclusiva de la plataforma. Para una nueva base, vuelva al esquema y al catálogo antes de reutilizar una necesidad.

\subsection{Conceptos para decidir con criterio}
\paragraph{Referencia histórica} Dato documentado de una campaña anterior, útil para comparar si copia y alcance coinciden.

\paragraph{Práctica nueva} Actividad que debe ejecutarse y registrarse; no se presenta como realizada por este manual.

\Needspace{8\baselineskip}
\begingroup
\setbox0=\vbox{\hsize=\linewidth
{\normalfont\normalsize\sffamily\bfseries\normalsize\color{Teal}\hyphenpenalty=10000\exhyphenpenalty=10000 Campos y controles\par}

\par
\begingroup
\fontsize{10}{12.5}\selectfont
\setlength{\tabcolsep}{5pt}
\renewcommand{\arraystretch}{1.15}
\begin{tabular}{@{}>{\RaggedRight\arraybackslash}p{0.27\linewidth}>{\RaggedRight\arraybackslash}p{0.66\linewidth}@{}}
\toprule
\textbf{Campo o control} & \textbf{Significado y criterio} \\
\midrule
Ficha de práctica & Necesidad copiable, preguntas, periodicidad, grano, indicadores y criterios de éxito. \\
\bottomrule
\end{tabular}
\endgroup
}
\Needspace{\dimexpr\ht0+\dp0+48pt\relax}
\endgroup
\subsection{Campos y controles}
\begingroup
\fontsize{10}{12.5}\selectfont
\setlength{\tabcolsep}{5pt}
\renewcommand{\arraystretch}{1.15}
\begin{longtable}{@{}>{\RaggedRight\arraybackslash}p{0.27\linewidth}>{\RaggedRight\arraybackslash}p{0.66\linewidth}@{}}
\toprule
\textbf{Campo o control} & \textbf{Significado y criterio} \\
\midrule
\endfirsthead
\toprule
\textbf{Campo o control} & \textbf{Significado y criterio} \\
\midrule
\endhead
\midrule
\multicolumn{2}{r}{\footnotesize Continúa en la página siguiente}\\
\endfoot
\bottomrule
\endlastfoot
Ficha de práctica & Necesidad copiable, preguntas, periodicidad, grano, indicadores y criterios de éxito. \\
\addlinespace[3pt]
Bitácora de la alumna & Registre fecha, fuente, instantánea, proveedor/modelo, propuesta, ejecución, acciones y resultados propios. \\
\addlinespace[3pt]
\end{longtable}
\endgroup

\subsection{Paso a paso}
\begin{enumerate}
\item Comience por la ruta básica de la primera fuente.\par\nopagebreak[4]\smallskip{\small\textbf{Compruebe:} Aprende el recorrido con pocas medidas.}
\item Repita la ruta básica de la segunda fuente.\par\nopagebreak[4]\smallskip{\small\textbf{Compruebe:} Comprueba qué se mantiene y qué cambió de evento y claves.}
\item Use el caso avanzado AW para conteo distinto y razones.\par\nopagebreak[4]\smallskip{\small\textbf{Compruebe:} Puede explicar diferencia entre línea y pedido.}
\item Use el caso avanzado WWI para roles de relación.\par\nopagebreak[4]\smallskip{\small\textbf{Compruebe:} Puede justificar comprador y ciudad de entrega y conservar su procedencia.}
\item Compare las dos lecturas sin sumar monedas ni eventos diferentes.\par\nopagebreak[4]\smallskip{\small\textbf{Compruebe:} Comprende la arquitectura común y los significados propios de cada fuente.}
\end{enumerate}
\subsection{Qué hace la IA}
La IA puede ofrecer diseños distintos en cada intento. Aprenderá a revisar cada resultado, justificar las decisiones y corregir lo necesario.

\subsection{Qué comprueba la plataforma}
Cada práctica vuelve a pasar autorizaciones, viabilidad, contrato, aprobación, preparación, carga y conciliación que correspondan.

\subsection{Qué decide la persona}
La alumna ejecuta y registra; el analista supervisa; la gerencia explica utilidad y límites de la lectura.

\paragraph{Resultado esperado} Cuatro fichas de trabajo con resultados propios y referencias históricas separadas.

\Needspace{8\baselineskip}
\begingroup
\setbox0=\vbox{\hsize=\linewidth
{\normalfont\normalsize\sffamily\bfseries\normalsize\color{Teal}\hyphenpenalty=10000\exhyphenpenalty=10000 Errores y recuperación\par}

\par
\begingroup
\fontsize{10}{12.5}\selectfont
\setlength{\tabcolsep}{5pt}
\renewcommand{\arraystretch}{1.15}
\begin{tabular}{@{}>{\RaggedRight\arraybackslash}p{0.27\linewidth}>{\RaggedRight\arraybackslash}p{0.66\linewidth}@{}}
\toprule
\textbf{Situación} & \textbf{Qué hacer} \\
\midrule
Los identificadores difieren del manual. & Es normal en otra instalación. Compruebe fuente, objetivo y evento; no cambie números a mano. \\
\bottomrule
\end{tabular}
\endgroup
}
\Needspace{\dimexpr\ht0+\dp0+48pt\relax}
\endgroup
\subsection{Errores y recuperación}
\begingroup
\fontsize{10}{12.5}\selectfont
\setlength{\tabcolsep}{5pt}
\renewcommand{\arraystretch}{1.15}
\begin{longtable}{@{}>{\RaggedRight\arraybackslash}p{0.27\linewidth}>{\RaggedRight\arraybackslash}p{0.66\linewidth}@{}}
\toprule
\textbf{Situación} & \textbf{Qué hacer} \\
\midrule
\endfirsthead
\toprule
\textbf{Situación} & \textbf{Qué hacer} \\
\midrule
\endhead
\midrule
\multicolumn{2}{r}{\footnotesize Continúa en la página siguiente}\\
\endfoot
\bottomrule
\endlastfoot
Los identificadores difieren del manual. & Es normal en otra instalación. Compruebe fuente, objetivo y evento; no cambie números a mano. \\
\addlinespace[3pt]
El LLM no entrega la cantidad de KPI esperada. & Revise entradas y recetas; personalice o excluya extras y valide otra versión. No existe un número universal de indicadores. \\
\addlinespace[3pt]
\end{longtable}
\endgroup

\subsection{Ejercicio}
Registre la evidencia de una práctica y clasifique qué es generación nueva y qué es consulta histórica.

\textbf{Respuesta orientativa.} La evidencia debe permitir reconocer fecha, fuente, propuesta y expediente, con resultado real y no sólo una captura antigua.

\subsection{Compruebe lo que aprendió}
\paragraph{¿Se deben forzar las cifras del manual?} No. Son referencias de una copia y alcance concretos; su ejecución debe conciliar datos reales.

\paragraph{¿La universalidad obliga a usar la misma relación en ambas bases?} No. El proceso es común; las relaciones y significados dependen de metadatos y decisiones verificadas.

\subsection{Casos para practicar}
\begingroup
\setbox0=\vbox{\hsize=\linewidth
{\normalfont\normalsize\sffamily\bfseries\normalsize\color{Azul}\hyphenpenalty=10000\exhyphenpenalty=10000 Ruta básica AdventureWorks unidades por mes y producto\par}

\par
\begingroup
\fontsize{10}{12.5}\selectfont
\setlength{\tabcolsep}{5pt}
\renewcommand{\arraystretch}{1.15}
\begin{tabular}{@{}>{\RaggedRight\arraybackslash}p{0.27\linewidth}>{\RaggedRight\arraybackslash}p{0.66\linewidth}@{}}
\toprule
\textbf{Dato} & \textbf{Contenido} \\
\midrule
Identificación de práctica & P-AW-B \\
\bottomrule
\end{tabular}
\endgroup
}
\Needspace{\dimexpr\ht0+\dp0+48pt\relax}
\endgroup
\subsubsection{Ruta básica AdventureWorks unidades por mes y producto}
\begingroup
\fontsize{10}{12.5}\selectfont
\setlength{\tabcolsep}{5pt}
\renewcommand{\arraystretch}{1.15}
\begin{longtable}{@{}>{\RaggedRight\arraybackslash}p{0.27\linewidth}>{\RaggedRight\arraybackslash}p{0.66\linewidth}@{}}
\toprule
\textbf{Dato} & \textbf{Contenido} \\
\midrule
\endfirsthead
\toprule
\textbf{Dato} & \textbf{Contenido} \\
\midrule
\endhead
\midrule
\multicolumn{2}{r}{\footnotesize Continúa en la página siguiente}\\
\endfoot
\bottomrule
\endlastfoot
Identificación de práctica & P-AW-B \\
\addlinespace[3pt]
Nivel & Inicio \\
\addlinespace[3pt]
Fuente & AdventureWorks 2022 \\
\addlinespace[3pt]
Estado de la evidencia & Ejercicio nuevo propuesto para las alumnas; no ejecutado en esta revisión documental. \\
\addlinespace[3pt]
Preguntas de negocio & Evolución de ventas en el tiempo y desempeño por producto, según las etiquetas disponibles del catálogo de esa fuente. \\
\addlinespace[3pt]
Periodicidad & Mensual \\
\addlinespace[3pt]
Granularidad & Línea de pedido \\
\addlinespace[3pt]
Indicadores & Suma de unidades; las salidas adicionales se revisan y excluyen si no son necesarias. \\
\addlinespace[3pt]
\end{longtable}
\endgroup

\paragraph{Necesidad para la práctica} Conservar cada línea de pedido de venta para comparar mensualmente la suma de unidades por producto. Usar la fecha del pedido y un nombre de producto comprobable. Mantener la referencia de cada línea a su pedido. Limitar esta práctica al total de unidades y su evolución; no interpretar estos registros como facturas, cobros o pronósticos.

\paragraph{Referencia numérica} No se asigna un resultado predeterminado a una ejecución no realizada. Si la copia y alcance equivalen al caso histórico, el total global de referencia es 274.914 unidades; debe conciliarse el caso propio.

\begingroup
\setbox0=\vbox{\hsize=\linewidth
{\normalfont\normalsize\bfseries Criterios de éxito\par}
\begin{itemize}
\item Objetivo conservado sin ampliar el evento.
\end{itemize}
}
\Needspace{\dimexpr\ht0+\dp0+24pt\relax}
\endgroup
\paragraph{Criterios de éxito}
\begin{itemize}
\item Objetivo conservado sin ampliar el evento.
\item Grano de línea y fecha del pedido verificados.
\item Producto descriptivo y ruta sin multiplicación.
\item Conciliación del universo nuevo.
\item Panel por unidades y mes; no se exige que la IA produzca una cantidad fija de tarjetas.
\end{itemize}

\begingroup
\setbox0=\vbox{\hsize=\linewidth
{\normalfont\normalsize\sffamily\bfseries\normalsize\color{Azul}\hyphenpenalty=10000\exhyphenpenalty=10000 Ruta avanzada AdventureWorks documentos distintos y razón por pedido\par}

\par
\begingroup
\fontsize{10}{12.5}\selectfont
\setlength{\tabcolsep}{5pt}
\renewcommand{\arraystretch}{1.15}
\begin{tabular}{@{}>{\RaggedRight\arraybackslash}p{0.27\linewidth}>{\RaggedRight\arraybackslash}p{0.66\linewidth}@{}}
\toprule
\textbf{Dato} & \textbf{Contenido} \\
\midrule
Identificación de práctica & P-AW-A \\
\bottomrule
\end{tabular}
\endgroup
}
\Needspace{\dimexpr\ht0+\dp0+48pt\relax}
\endgroup
\subsubsection{Ruta avanzada AdventureWorks documentos distintos y razón por pedido}
\begingroup
\fontsize{10}{12.5}\selectfont
\setlength{\tabcolsep}{5pt}
\renewcommand{\arraystretch}{1.15}
\begin{longtable}{@{}>{\RaggedRight\arraybackslash}p{0.27\linewidth}>{\RaggedRight\arraybackslash}p{0.66\linewidth}@{}}
\toprule
\textbf{Dato} & \textbf{Contenido} \\
\midrule
\endfirsthead
\toprule
\textbf{Dato} & \textbf{Contenido} \\
\midrule
\endhead
\midrule
\multicolumn{2}{r}{\footnotesize Continúa en la página siguiente}\\
\endfoot
\bottomrule
\endlastfoot
Identificación de práctica & P-AW-A \\
\addlinespace[3pt]
Nivel & Intermedio \\
\addlinespace[3pt]
Fuente & AdventureWorks 2022 \\
\addlinespace[3pt]
Estado de la evidencia & Caso histórico comprobado el 6 de noviembre de 2026: propuesta 129, ejecución 21. Repetición por las alumnas pendiente. \\
\addlinespace[3pt]
Preguntas de negocio & Tiempo, producto, cliente y territorio. Revise el catálogo propio antes de seleccionarlas. \\
\addlinespace[3pt]
Periodicidad & Mensual \\
\addlinespace[3pt]
Granularidad & Línea de pedido \\
\addlinespace[3pt]
Indicadores & Importe registrado total; unidades totales; pedidos distintos; importe total/pedidos distintos. \\
\addlinespace[3pt]
Pregunta para el copiloto & Resume las unidades vendidas en 2013 para Northwest. Indica el total y el mes con mayor cantidad, sin atribuir causas ni mezclar importes. \\
\addlinespace[3pt]
\end{longtable}
\endgroup

\paragraph{Necesidad para la práctica} Conservar cada línea de pedido de venta en un datamart. Para priorizar productos y territorios, comparar mensualmente cuatro indicadores: suma del importe registrado de las líneas, suma de unidades, conteo de identificadores de pedido distintos e importe total dividido entre pedidos distintos. Usar fecha y territorio comercial del pedido, producto y cliente con nombre verificable. Los rankings comparan sumas por entidad. Mantener la trazabilidad de cada línea al pedido. La moneda se reportará sólo si se comprueba en la fuente.

\begingroup
\setbox0=\vbox{\hsize=\linewidth
{\normalfont\normalsize\bfseries Rutas que deben comprobarse\par}
\begin{itemize}
\item Fecha y territorio comercial del pedido, con territorio procedente de Sales.SalesTerritory.
\end{itemize}
}
\Needspace{\dimexpr\ht0+\dp0+24pt\relax}
\endgroup
\paragraph{Rutas que deben comprobarse}
\begin{itemize}
\item Fecha y territorio comercial del pedido, con territorio procedente de Sales.SalesTerritory.
\item Producto y cliente con nombre verificable; clave de línea y pedido conservadas.
\end{itemize}

\begingroup
\setbox0=\vbox{\hsize=\linewidth
{\normalfont\normalsize\bfseries Pasos específicos\par}
\begin{itemize}
\item Use el objetivo exacto y pase por los cinco pasos del asistente.
\end{itemize}
}
\Needspace{\dimexpr\ht0+\dp0+24pt\relax}
\endgroup
\paragraph{Pasos específicos}
\begin{itemize}
\item Use el objetivo exacto y pase por los cinco pasos del asistente.
\item Revise que la cuenta del documento use conteo distinto y que la razón lo use como denominador.
\item No agregue costos o margen para igualar un número de tarjetas.
\item Apruebe únicamente cero errores y avisos comprendidos; complete los cinco pasos de materialización.
\item Consulte primero global, luego 2013/Northwest.
\item En Vista analítica elija unidades; lea tabla mensual; pregunte al chat por total y mes líder.
\item Exporte y abra PDF y Excel con ese mismo alcance.
\end{itemize}

\begingroup
\setbox0=\vbox{\hsize=\linewidth
{\normalfont\normalsize\normalfont\normalsize\bfseries\hyphenpenalty=10000\exhyphenpenalty=10000 Resultados del caso histórico\par}

\par
\begingroup
\fontsize{10}{12.5}\selectfont
\setlength{\tabcolsep}{5pt}
\renewcommand{\arraystretch}{1.15}
\begin{tabular}{@{}>{\RaggedRight\arraybackslash}p{0.27\linewidth}>{\RaggedRight\arraybackslash}p{0.66\linewidth}@{}}
\toprule
\textbf{Dato} & \textbf{Contenido} \\
\midrule
Fecha & 2026-11-06 \\
\bottomrule
\end{tabular}
\endgroup
}
\Needspace{\dimexpr\ht0+\dp0+48pt\relax}
\endgroup
\paragraph{Resultados del caso histórico}
\begingroup
\fontsize{10}{12.5}\selectfont
\setlength{\tabcolsep}{5pt}
\renewcommand{\arraystretch}{1.15}
\begin{longtable}{@{}>{\RaggedRight\arraybackslash}p{0.27\linewidth}>{\RaggedRight\arraybackslash}p{0.66\linewidth}@{}}
\toprule
\textbf{Dato} & \textbf{Contenido} \\
\midrule
\endfirsthead
\toprule
\textbf{Dato} & \textbf{Contenido} \\
\midrule
\endhead
\midrule
\multicolumn{2}{r}{\footnotesize Continúa en la página siguiente}\\
\endfoot
\bottomrule
\endlastfoot
Fecha & 2026-11-06 \\
\addlinespace[3pt]
Propuesta & 129 \\
\addlinespace[3pt]
Ejecución & 21 \\
\addlinespace[3pt]
Filas de origen & 121317 \\
\addlinespace[3pt]
Filas de destino & 121317 \\
\addlinespace[3pt]
Diferencia de filas & 0 \\
\addlinespace[3pt]
Tablas & 5 \\
\addlinespace[3pt]
Interpretación de etiquetas & Europe, North America y Pacific publicadas como Europa, América del Norte y Pacífico; AU, CA, DE, FR, GB, US excluidos de traducción. \\
\addlinespace[3pt]
Exportación & PDF de 6 páginas y Excel de 6 hojas de la Vista analítica por unidades 2013/Northwest; observación histórica de etiqueta genérica de razón monetaria conservada. \\
\addlinespace[3pt]
\end{longtable}
\endgroup

\begingroup
\setbox0=\vbox{\hsize=\linewidth
{\normalfont\normalsize\normalfont\normalsize\bfseries\hyphenpenalty=10000\exhyphenpenalty=10000 Selección completa\par}

\par
\begingroup
\fontsize{10}{12.5}\selectfont
\setlength{\tabcolsep}{5pt}
\renewcommand{\arraystretch}{1.15}
\begin{tabular}{@{}>{\RaggedRight\arraybackslash}p{0.27\linewidth}>{\RaggedRight\arraybackslash}p{0.66\linewidth}@{}}
\toprule
\textbf{Dato} & \textbf{Contenido} \\
\midrule
Importe exacto & 109846381.399888 \\
\bottomrule
\end{tabular}
\endgroup
}
\Needspace{\dimexpr\ht0+\dp0+48pt\relax}
\endgroup
\paragraph{Selección completa}
\begingroup
\fontsize{10}{12.5}\selectfont
\setlength{\tabcolsep}{5pt}
\renewcommand{\arraystretch}{1.15}
\begin{longtable}{@{}>{\RaggedRight\arraybackslash}p{0.27\linewidth}>{\RaggedRight\arraybackslash}p{0.66\linewidth}@{}}
\toprule
\textbf{Dato} & \textbf{Contenido} \\
\midrule
\endfirsthead
\toprule
\textbf{Dato} & \textbf{Contenido} \\
\midrule
\endhead
\midrule
\multicolumn{2}{r}{\footnotesize Continúa en la página siguiente}\\
\endfoot
\bottomrule
\endlastfoot
Importe exacto & 109846381.399888 \\
\addlinespace[3pt]
Moneda & USD comprobado \\
\addlinespace[3pt]
Unidades & 274914 \\
\addlinespace[3pt]
Pedidos distintos & 31465 \\
\addlinespace[3pt]
Importe por pedido exacto & 3491.065672966407 \\
\addlinespace[3pt]
\end{longtable}
\endgroup

\begingroup
\setbox0=\vbox{\hsize=\linewidth
{\normalfont\normalsize\normalfont\normalsize\bfseries\hyphenpenalty=10000\exhyphenpenalty=10000 Selección filtrada\par}

\par
\begingroup
\fontsize{10}{12.5}\selectfont
\setlength{\tabcolsep}{5pt}
\renewcommand{\arraystretch}{1.15}
\begin{tabular}{@{}>{\RaggedRight\arraybackslash}p{0.27\linewidth}>{\RaggedRight\arraybackslash}p{0.66\linewidth}@{}}
\toprule
\textbf{Dato} & \textbf{Contenido} \\
\midrule
Año & 2013 \\
\bottomrule
\end{tabular}
\endgroup
}
\Needspace{\dimexpr\ht0+\dp0+48pt\relax}
\endgroup
\paragraph{Selección filtrada}
\begingroup
\fontsize{10}{12.5}\selectfont
\setlength{\tabcolsep}{5pt}
\renewcommand{\arraystretch}{1.15}
\begin{longtable}{@{}>{\RaggedRight\arraybackslash}p{0.27\linewidth}>{\RaggedRight\arraybackslash}p{0.66\linewidth}@{}}
\toprule
\textbf{Dato} & \textbf{Contenido} \\
\midrule
\endfirsthead
\toprule
\textbf{Dato} & \textbf{Contenido} \\
\midrule
\endhead
\midrule
\multicolumn{2}{r}{\footnotesize Continúa en la página siguiente}\\
\endfoot
\bottomrule
\endlastfoot
Año & 2013 \\
\addlinespace[3pt]
Territorio & Northwest \\
\addlinespace[3pt]
Importe exacto & 6015174.735836 \\
\addlinespace[3pt]
Unidades & 17452 \\
\addlinespace[3pt]
Pedidos distintos & 2053 \\
\addlinespace[3pt]
Importe por pedido exacto & 2929.943855740867 \\
\addlinespace[3pt]
Mes de mayor cantidad & junio de 2013 \\
\addlinespace[3pt]
Unidades del mes de mayor cantidad & 2494 \\
\addlinespace[3pt]
\end{longtable}
\endgroup

\begingroup
\setbox0=\vbox{\hsize=\linewidth
{\normalfont\normalsize\bfseries Criterios de éxito\par}
\begin{itemize}
\item Mismo evento, grano y roles del caso.
\end{itemize}
}
\Needspace{\dimexpr\ht0+\dp0+24pt\relax}
\endgroup
\paragraph{Criterios de éxito}
\begin{itemize}
\item Mismo evento, grano y roles del caso.
\item Conciliación sin diferencia.
\item Cuatro indicadores revisados; no promedio de línea rotulado como promedio por pedido.
\item Unidad demostrada y límites de pedido conservados.
\item Valores del panel, chat y archivos contrastados con la selección.
\end{itemize}

\begingroup
\setbox0=\vbox{\hsize=\linewidth
{\normalfont\normalsize\sffamily\bfseries\normalsize\color{Azul}\hyphenpenalty=10000\exhyphenpenalty=10000 Ruta básica WideWorldImporters facturas y unidades por mes\par}

\par
\begingroup
\fontsize{10}{12.5}\selectfont
\setlength{\tabcolsep}{5pt}
\renewcommand{\arraystretch}{1.15}
\begin{tabular}{@{}>{\RaggedRight\arraybackslash}p{0.27\linewidth}>{\RaggedRight\arraybackslash}p{0.66\linewidth}@{}}
\toprule
\textbf{Dato} & \textbf{Contenido} \\
\midrule
Identificación de práctica & P-WWI-B \\
\bottomrule
\end{tabular}
\endgroup
}
\Needspace{\dimexpr\ht0+\dp0+48pt\relax}
\endgroup
\subsubsection{Ruta básica WideWorldImporters facturas y unidades por mes}
\begingroup
\fontsize{10}{12.5}\selectfont
\setlength{\tabcolsep}{5pt}
\renewcommand{\arraystretch}{1.15}
\begin{longtable}{@{}>{\RaggedRight\arraybackslash}p{0.27\linewidth}>{\RaggedRight\arraybackslash}p{0.66\linewidth}@{}}
\toprule
\textbf{Dato} & \textbf{Contenido} \\
\midrule
\endfirsthead
\toprule
\textbf{Dato} & \textbf{Contenido} \\
\midrule
\endhead
\midrule
\multicolumn{2}{r}{\footnotesize Continúa en la página siguiente}\\
\endfoot
\bottomrule
\endlastfoot
Identificación de práctica & P-WWI-B \\
\addlinespace[3pt]
Nivel & Inicio \\
\addlinespace[3pt]
Fuente & WideWorldImporters \\
\addlinespace[3pt]
Estado de la evidencia & Ejercicio nuevo propuesto para las alumnas; no ejecutado en esta revisión documental. \\
\addlinespace[3pt]
Preguntas de negocio & Tiempo y producto según el catálogo propio. \\
\addlinespace[3pt]
Periodicidad & Mensual \\
\addlinespace[3pt]
Granularidad & Línea de factura \\
\addlinespace[3pt]
Indicadores & Suma de unidades; otras salidas se revisan antes de materializar. \\
\addlinespace[3pt]
\end{longtable}
\endgroup

\paragraph{Necesidad para la práctica} Conservar cada línea de factura para comparar mensualmente la suma de unidades facturadas por producto. Usar la fecha de emisión y el producto con nombre verificable. Mantener el identificador de factura y el de línea para conservar la trazabilidad. Limitar la práctica a unidades y evolución; no afirmar moneda, cobro, utilidad ni entrega realizada.

\paragraph{Referencia numérica} No se inventa un resultado para esta ejecución nueva. El caso histórico completo WWI 22 conservó 8.950.628 unidades; sólo sirve de contraste si copia y alcance son equivalentes.

\begingroup
\setbox0=\vbox{\hsize=\linewidth
{\normalfont\normalsize\bfseries Criterios de éxito\par}
\begin{itemize}
\item La propuesta usa facturación y fecha de emisión, no pedidos como sustituto.
\end{itemize}
}
\Needspace{\dimexpr\ht0+\dp0+24pt\relax}
\endgroup
\paragraph{Criterios de éxito}
\begin{itemize}
\item La propuesta usa facturación y fecha de emisión, no pedidos como sustituto.
\item Producto descriptivo y línea identificable.
\item No se añade moneda ni utilidad por intuición.
\item Carga conciliada y panel en unidades.
\end{itemize}

\begingroup
\setbox0=\vbox{\hsize=\linewidth
{\normalfont\normalsize\sffamily\bfseries\normalsize\color{Azul}\hyphenpenalty=10000\exhyphenpenalty=10000 Ruta avanzada WideWorldImporters roles de cliente y territorio\par}

\par
\begingroup
\fontsize{10}{12.5}\selectfont
\setlength{\tabcolsep}{5pt}
\renewcommand{\arraystretch}{1.15}
\begin{tabular}{@{}>{\RaggedRight\arraybackslash}p{0.27\linewidth}>{\RaggedRight\arraybackslash}p{0.66\linewidth}@{}}
\toprule
\textbf{Dato} & \textbf{Contenido} \\
\midrule
Identificación de práctica & P-WWI-A \\
\bottomrule
\end{tabular}
\endgroup
}
\Needspace{\dimexpr\ht0+\dp0+48pt\relax}
\endgroup
\subsubsection{Ruta avanzada WideWorldImporters roles de cliente y territorio}
\begingroup
\fontsize{10}{12.5}\selectfont
\setlength{\tabcolsep}{5pt}
\renewcommand{\arraystretch}{1.15}
\begin{longtable}{@{}>{\RaggedRight\arraybackslash}p{0.27\linewidth}>{\RaggedRight\arraybackslash}p{0.66\linewidth}@{}}
\toprule
\textbf{Dato} & \textbf{Contenido} \\
\midrule
\endfirsthead
\toprule
\textbf{Dato} & \textbf{Contenido} \\
\midrule
\endhead
\midrule
\multicolumn{2}{r}{\footnotesize Continúa en la página siguiente}\\
\endfoot
\bottomrule
\endlastfoot
Identificación de práctica & P-WWI-A \\
\addlinespace[3pt]
Nivel & Intermedio/avanzado \\
\addlinespace[3pt]
Fuente & WideWorldImporters \\
\addlinespace[3pt]
Estado de la evidencia & Caso histórico comprobado el 6 de noviembre de 2026: propuesta 131, ejecución 22; dos correcciones guiadas sucesivas. Repetición por las alumnas pendiente. \\
\addlinespace[3pt]
Preguntas de negocio & Tiempo, producto, cliente y territorio en el catálogo WWI. \\
\addlinespace[3pt]
Periodicidad & Mensual \\
\addlinespace[3pt]
Granularidad & Línea de factura \\
\addlinespace[3pt]
Indicadores & Importe registrado; unidades; facturas distintas; importe total/facturas distintas. \\
\addlinespace[3pt]
\end{longtable}
\endgroup

\paragraph{Necesidad para la práctica} Diseñar un datamart cuyo grano físico sea una línea de factura. Sobre ese detalle, resumir mensualmente el importe registrado de la línea, las unidades y las facturas distintas por producto y cliente. Usar la fecha de emisión. Para el territorio, usar la provincia de la ciudad de entrega registrada en la ficha del cliente y aclarar que no representa necesariamente su ubicación histórica. Conservar el importe de origen sin asumir moneda, utilidad ni composición fiscal. Limitar los indicadores a importe total, unidades totales, facturas distintas e importe promedio por factura. Los rankings deben comparar totales por producto; el máximo de una sola línea no identifica al producto líder.

\paragraph{Pregunta para el copiloto} Resume el importe, las unidades y las facturas distintas en 2015 para California. Explica el denominador del importe por factura y qué significa la provincia usada, sin asumir moneda ni entrega realizada.

\begingroup
\setbox0=\vbox{\hsize=\linewidth
{\normalfont\normalsize\bfseries Rutas que deben comprobarse\par}
\begin{itemize}
\item Sales.InvoiceLines\(\rightarrow\)Sales.Invoices por la relación de factura.
\end{itemize}
}
\Needspace{\dimexpr\ht0+\dp0+24pt\relax}
\endgroup
\paragraph{Rutas que deben comprobarse}
\begin{itemize}
\item Sales.InvoiceLines\(\rightarrow\)Sales.Invoices por la relación de factura.
\item Cliente comprador: Invoices.CustomerID\(\rightarrow\)Customers, no BillToCustomerID para esta necesidad.
\item Territorio: Customers.DeliveryCityID\(\rightarrow\)Cities.StateProvinceID\(\rightarrow\)StateProvinces, no PostalCityID para esta necesidad.
\item La provincia es el destino dimensional; una elección intermedia puede precisar la ruta sin reemplazarla por ciudad.
\end{itemize}

\begingroup
\setbox0=\vbox{\hsize=\linewidth
{\normalfont\normalsize\bfseries Pasos específicos\par}
\begin{itemize}
\item Use el objetivo exacto y revise viabilidad.
\end{itemize}
}
\Needspace{\dimexpr\ht0+\dp0+24pt\relax}
\endgroup
\paragraph{Pasos específicos}
\begin{itemize}
\item Use el objetivo exacto y revise viabilidad.
\item Compruebe grano de factura, fecha de emisión, medidas y cuatro indicadores.
\item Si la generación usa otro rol, abra Personalización y Resolver una relación con evidencia.
\item Elija Cliente y Invoices.CustomerID\(\rightarrow\)Customers; justifique que se analiza al comprador; cree una versión con relación corregida.
\item Revise el nuevo resultado; elija Territorio y Customers.DeliveryCityID\(\rightarrow\)Cities; justifique ubicación registrada de entrega; cree otra versión.
\item Confirme que se conserva CustomerID y la provincia final de StateProvinces. No acepte una vía postal o de facturación sólo porque mantiene filas.
\item Apruebe y materialice si cumple; revise conciliación y divisa.
\item En analítica elija 2015/California, contraste cuatro indicadores; pregunte al chat sobre alcance, denominador y límites; exporte y abra archivos.
\end{itemize}

\begingroup
\setbox0=\vbox{\hsize=\linewidth
{\normalfont\normalsize\normalfont\normalsize\bfseries\hyphenpenalty=10000\exhyphenpenalty=10000 Resultados del caso histórico\par}

\par
\begingroup
\fontsize{10}{12.5}\selectfont
\setlength{\tabcolsep}{5pt}
\renewcommand{\arraystretch}{1.15}
\begin{tabular}{@{}>{\RaggedRight\arraybackslash}p{0.27\linewidth}>{\RaggedRight\arraybackslash}p{0.66\linewidth}@{}}
\toprule
\textbf{Dato} & \textbf{Contenido} \\
\midrule
Fecha & 2026-11-06 \\
\bottomrule
\end{tabular}
\endgroup
}
\Needspace{\dimexpr\ht0+\dp0+48pt\relax}
\endgroup
\paragraph{Resultados del caso histórico}
\begingroup
\fontsize{10}{12.5}\selectfont
\setlength{\tabcolsep}{5pt}
\renewcommand{\arraystretch}{1.15}
\begin{longtable}{@{}>{\RaggedRight\arraybackslash}p{0.27\linewidth}>{\RaggedRight\arraybackslash}p{0.66\linewidth}@{}}
\toprule
\textbf{Dato} & \textbf{Contenido} \\
\midrule
\endfirsthead
\toprule
\textbf{Dato} & \textbf{Contenido} \\
\midrule
\endhead
\midrule
\multicolumn{2}{r}{\footnotesize Continúa en la página siguiente}\\
\endfoot
\bottomrule
\endlastfoot
Fecha & 2026-11-06 \\
\addlinespace[3pt]
Propuesta & 131 \\
\addlinespace[3pt]
Ejecución & 22 \\
\addlinespace[3pt]
Filas de origen & 228265 \\
\addlinespace[3pt]
Filas de destino & 228265 \\
\addlinespace[3pt]
Diferencia de filas & 0 \\
\addlinespace[3pt]
Tablas & 5 \\
\addlinespace[3pt]
Exportación & PDF de 6 páginas y Excel de 6 hojas/4 gráficos de la vista filtrada; es un ejemplo histórico, no formato fijo. \\
\addlinespace[3pt]
\end{longtable}
\endgroup

\begingroup
\setbox0=\vbox{\hsize=\linewidth
{\normalfont\normalsize\normalfont\normalsize\bfseries\hyphenpenalty=10000\exhyphenpenalty=10000 Selección completa\par}

\par
\begingroup
\fontsize{10}{12.5}\selectfont
\setlength{\tabcolsep}{5pt}
\renewcommand{\arraystretch}{1.15}
\begin{tabular}{@{}>{\RaggedRight\arraybackslash}p{0.27\linewidth}>{\RaggedRight\arraybackslash}p{0.66\linewidth}@{}}
\toprule
\textbf{Dato} & \textbf{Contenido} \\
\midrule
Importe exacto & 198043439.45 \\
\bottomrule
\end{tabular}
\endgroup
}
\Needspace{\dimexpr\ht0+\dp0+48pt\relax}
\endgroup
\paragraph{Selección completa}
\begingroup
\fontsize{10}{12.5}\selectfont
\setlength{\tabcolsep}{5pt}
\renewcommand{\arraystretch}{1.15}
\begin{longtable}{@{}>{\RaggedRight\arraybackslash}p{0.27\linewidth}>{\RaggedRight\arraybackslash}p{0.66\linewidth}@{}}
\toprule
\textbf{Dato} & \textbf{Contenido} \\
\midrule
\endfirsthead
\toprule
\textbf{Dato} & \textbf{Contenido} \\
\midrule
\endhead
\midrule
\multicolumn{2}{r}{\footnotesize Continúa en la página siguiente}\\
\endfoot
\bottomrule
\endlastfoot
Importe exacto & 198043439.45 \\
\addlinespace[3pt]
Moneda & moneda de origen; sin código único demostrado \\
\addlinespace[3pt]
Unidades & 8950628 \\
\addlinespace[3pt]
Facturas distintas & 70510 \\
\addlinespace[3pt]
Importe por factura exacto & 2808.7283995177986 \\
\addlinespace[3pt]
\end{longtable}
\endgroup

\begingroup
\setbox0=\vbox{\hsize=\linewidth
{\normalfont\normalsize\normalfont\normalsize\bfseries\hyphenpenalty=10000\exhyphenpenalty=10000 Selección filtrada\par}

\par
\begingroup
\fontsize{10}{12.5}\selectfont
\setlength{\tabcolsep}{5pt}
\renewcommand{\arraystretch}{1.15}
\begin{tabular}{@{}>{\RaggedRight\arraybackslash}p{0.27\linewidth}>{\RaggedRight\arraybackslash}p{0.66\linewidth}@{}}
\toprule
\textbf{Dato} & \textbf{Contenido} \\
\midrule
Año & 2015 \\
\bottomrule
\end{tabular}
\endgroup
}
\Needspace{\dimexpr\ht0+\dp0+48pt\relax}
\endgroup
\paragraph{Selección filtrada}
\begingroup
\fontsize{10}{12.5}\selectfont
\setlength{\tabcolsep}{5pt}
\renewcommand{\arraystretch}{1.15}
\begin{longtable}{@{}>{\RaggedRight\arraybackslash}p{0.27\linewidth}>{\RaggedRight\arraybackslash}p{0.66\linewidth}@{}}
\toprule
\textbf{Dato} & \textbf{Contenido} \\
\midrule
\endfirsthead
\toprule
\textbf{Dato} & \textbf{Contenido} \\
\midrule
\endhead
\midrule
\multicolumn{2}{r}{\footnotesize Continúa en la página siguiente}\\
\endfoot
\bottomrule
\endlastfoot
Año & 2015 \\
\addlinespace[3pt]
Territorio & California \\
\addlinespace[3pt]
Lineas & 3771 \\
\addlinespace[3pt]
Importe exacto & 3191769.11 \\
\addlinespace[3pt]
Unidades & 142247 \\
\addlinespace[3pt]
Facturas distintas & 1179 \\
\addlinespace[3pt]
Importe por factura exacto & 2707.1832994062765 \\
\addlinespace[3pt]
\end{longtable}
\endgroup

\begingroup
\setbox0=\vbox{\hsize=\linewidth
{\normalfont\normalsize\normalfont\normalsize\bfseries\hyphenpenalty=10000\exhyphenpenalty=10000 Fecha dimension\par}

\par
\begingroup
\fontsize{10}{12.5}\selectfont
\setlength{\tabcolsep}{5pt}
\renewcommand{\arraystretch}{1.15}
\begin{tabular}{@{}>{\RaggedRight\arraybackslash}p{0.27\linewidth}>{\RaggedRight\arraybackslash}p{0.66\linewidth}@{}}
\toprule
\textbf{Dato} & \textbf{Contenido} \\
\midrule
Fechas unicas & 1069 \\
\bottomrule
\end{tabular}
\endgroup
}
\Needspace{\dimexpr\ht0+\dp0+48pt\relax}
\endgroup
\paragraph{Fecha dimension}
\begingroup
\fontsize{10}{12.5}\selectfont
\setlength{\tabcolsep}{5pt}
\renewcommand{\arraystretch}{1.15}
\begin{longtable}{@{}>{\RaggedRight\arraybackslash}p{0.27\linewidth}>{\RaggedRight\arraybackslash}p{0.66\linewidth}@{}}
\toprule
\textbf{Dato} & \textbf{Contenido} \\
\midrule
\endfirsthead
\toprule
\textbf{Dato} & \textbf{Contenido} \\
\midrule
\endhead
\midrule
\multicolumn{2}{r}{\footnotesize Continúa en la página siguiente}\\
\endfoot
\bottomrule
\endlastfoot
Fechas unicas & 1069 \\
\addlinespace[3pt]
Documentos origen & 70510 \\
\addlinespace[3pt]
Duplicados controlados & 69441 \\
\addlinespace[3pt]
\end{longtable}
\endgroup

\begingroup
\setbox0=\vbox{\hsize=\linewidth
{\normalfont\normalsize\bfseries Criterios de éxito\par}
\begin{itemize}
\item Evento de facturación con fecha de emisión.
\end{itemize}
}
\Needspace{\dimexpr\ht0+\dp0+24pt\relax}
\endgroup
\paragraph{Criterios de éxito}
\begin{itemize}
\item Evento de facturación con fecha de emisión.
\item Rutas exactas de comprador y ciudad de entrega registradas.
\item Facturas distintas, no conteo de líneas.
\item La moneda no se convierte en USD sin evidencia.
\item No se atribuyen ubicación histórica, entrega efectiva o composición fiscal no demostradas.
\item No se reutiliza el filtro del expediente legado 19 como oráculo del 22.
\end{itemize}


\IfFileExists{capturas/PR01.tex}{\input{capturas/PR01.tex}}{}
\section{Costos, descuentos y márgenes: una ampliación que exige evidencia}
\label{unidad:BI08}
\textbf{Para quién:} Analista de BI, Gerente comercial\par
\textbf{Al terminar podrá:} Reconocer cuándo se puede ampliar el análisis financiero y cuándo debe limitarse.

Los casos finales anteriores se centraron en importe, unidades, documentos y razones. No acreditan automáticamente un margen histórico. Si se necesita rentabilidad, primero se deben comprobar costos compatibles con la venta y su fecha.

Un costo unitario actual del catálogo puede servir de referencia, pero no se convierte por su nombre en el costo real histórico de cada línea vendida. Comparar ese costo con un importe de otro alcance puede producir un cálculo numérico correcto y una interpretación comercial incorrecta.

\subsection{Conceptos para decidir con criterio}
\paragraph{Descuento monetario} Importe derivado de una base y una tasa; la tasa por sí sola no es dinero.

\paragraph{Venta bruta} Precio por cantidad antes del descuento cuando esa es la receta comprobada.

\paragraph{Venta neta de un descuento} Requiere identificar qué deducción se aplica y si el importe de origen ya la contiene; no descontar dos veces.

\paragraph{Costo total} Costo unitario comparable \(\times\) unidades, o concepto total directo verificado; requiere una base coherente.

\paragraph{Margen bruto} Diferencia entre importe de venta y costo comparable del mismo alcance.

\paragraph{Porcentaje de margen} Margen \(\div\) importe de venta del mismo alcance \(\times\) 100, con divisor cero protegido.

\paragraph{Comparabilidad} Mismo evento, unidad, moneda, período y criterio de valoración para los componentes.

\Needspace{8\baselineskip}
\begingroup
\setbox0=\vbox{\hsize=\linewidth
{\normalfont\normalsize\sffamily\bfseries\normalsize\color{Teal}\hyphenpenalty=10000\exhyphenpenalty=10000 Campos y controles\par}

\par
\begingroup
\fontsize{10}{12.5}\selectfont
\setlength{\tabcolsep}{5pt}
\renewcommand{\arraystretch}{1.15}
\begin{tabular}{@{}>{\RaggedRight\arraybackslash}p{0.27\linewidth}>{\RaggedRight\arraybackslash}p{0.66\linewidth}@{}}
\toprule
\textbf{Campo o control} & \textbf{Significado y criterio} \\
\midrule
Fórmula y Origen & Deben mostrar entradas reales; no normalizar por nombre únicamente. \\
\bottomrule
\end{tabular}
\endgroup
}
\Needspace{\dimexpr\ht0+\dp0+48pt\relax}
\endgroup
\subsection{Campos y controles}
\begingroup
\fontsize{10}{12.5}\selectfont
\setlength{\tabcolsep}{5pt}
\renewcommand{\arraystretch}{1.15}
\begin{longtable}{@{}>{\RaggedRight\arraybackslash}p{0.27\linewidth}>{\RaggedRight\arraybackslash}p{0.66\linewidth}@{}}
\toprule
\textbf{Campo o control} & \textbf{Significado y criterio} \\
\midrule
\endfirsthead
\toprule
\textbf{Campo o control} & \textbf{Significado y criterio} \\
\midrule
\endhead
\midrule
\multicolumn{2}{r}{\footnotesize Continúa en la página siguiente}\\
\endfoot
\bottomrule
\endlastfoot
Fórmula y Origen & Deben mostrar entradas reales; no normalizar por nombre únicamente. \\
\addlinespace[3pt]
Cobertura de necesidad & Reconoce si costo o margen se resuelve o queda fuera. \\
\addlinespace[3pt]
Ajuste de seguridad & Puede corregir una tasa usada como importe, pero no inventa la evidencia económica ausente. \\
\addlinespace[3pt]
\end{longtable}
\endgroup

\subsection{Paso a paso}
\begin{enumerate}
\item Formule el objetivo aclarando qué costo desea usar y en qué condición temporal.\par\nopagebreak[4]\smallskip{\small\textbf{Compruebe:} Se distingue costo real, costo estándar actual u otra referencia.}
\item Revise viabilidad y entradas para costo, importe, cantidad y descuento.\par\nopagebreak[4]\smallskip{\small\textbf{Compruebe:} Si faltan datos, la salida se limita.}
\item Compruebe cálculos por fila y agregados.\par\nopagebreak[4]\smallskip{\small\textbf{Compruebe:} Evita sumar tasas como dinero o promediar márgenes sin ponderación.}
\item Apruebe sólo si el significado comparable puede defenderse.\par\nopagebreak[4]\smallskip{\small\textbf{Compruebe:} No se maquilla un bloqueo de calidad comercial.}
\end{enumerate}
\subsection{Qué hace la IA}
Sugiere recetas y explicaciones dentro del contexto. No puede demostrar costos históricos no disponibles a través de un texto convincente.

\subsection{Qué comprueba la plataforma}
Existencia, tipo, semántica implementada, recetas y relaciones. La revisión humana sigue siendo necesaria para condiciones económicas y temporalidad no acreditadas por el modelo.

\subsection{Qué decide la persona}
El analista decide qué costo es aceptable y qué límite debe declararse; la gerencia no debe recibir como utilidad real un indicador referencial.

\paragraph{Resultado esperado} Una ampliación justificable o un alcance financiero expresamente limitado.

\Needspace{8\baselineskip}
\begingroup
\setbox0=\vbox{\hsize=\linewidth
{\normalfont\normalsize\sffamily\bfseries\normalsize\color{Teal}\hyphenpenalty=10000\exhyphenpenalty=10000 Errores y recuperación\par}

\par
\begingroup
\fontsize{10}{12.5}\selectfont
\setlength{\tabcolsep}{5pt}
\renewcommand{\arraystretch}{1.15}
\begin{tabular}{@{}>{\RaggedRight\arraybackslash}p{0.27\linewidth}>{\RaggedRight\arraybackslash}p{0.66\linewidth}@{}}
\toprule
\textbf{Situación} & \textbf{Qué hacer} \\
\midrule
Se afirma margen sólo porque existen columnas Sales y Cost. & Compare procedencia, fecha, moneda y evento antes de usarlo. \\
\bottomrule
\end{tabular}
\endgroup
}
\Needspace{\dimexpr\ht0+\dp0+48pt\relax}
\endgroup
\subsection{Errores y recuperación}
\begingroup
\fontsize{10}{12.5}\selectfont
\setlength{\tabcolsep}{5pt}
\renewcommand{\arraystretch}{1.15}
\begin{longtable}{@{}>{\RaggedRight\arraybackslash}p{0.27\linewidth}>{\RaggedRight\arraybackslash}p{0.66\linewidth}@{}}
\toprule
\textbf{Situación} & \textbf{Qué hacer} \\
\midrule
\endfirsthead
\toprule
\textbf{Situación} & \textbf{Qué hacer} \\
\midrule
\endhead
\midrule
\multicolumn{2}{r}{\footnotesize Continúa en la página siguiente}\\
\endfoot
\bottomrule
\endlastfoot
Se afirma margen sólo porque existen columnas Sales y Cost. & Compare procedencia, fecha, moneda y evento antes de usarlo. \\
\addlinespace[3pt]
El importe original ya incluye descuento. & No aplique la tasa dos veces. Revise la fórmula y definición de origen. \\
\addlinespace[3pt]
Se solicitan impuestos o utilidad fiscal sin evidencia. & Limite la necesidad; no infiera composición fiscal por el nombre importe. \\
\addlinespace[3pt]
\end{longtable}
\endgroup

\subsection{Ejercicio}
Dos líneas tienen márgenes de 10 \% y 50 \% con importes de 900 y 100. Calcule el margen agregado.

\textbf{Respuesta orientativa.} 90+50=140 sobre 1.000; 14 \%. El promedio simple 30 \% no representa el porcentaje global.

\subsection{Compruebe lo que aprendió}
\paragraph{¿La existencia de un costo comprueba su vigencia histórica?} No. Requiere evidencia temporal y criterio de valoración.

\paragraph{¿Qué debe hacer si no puede justificarlo?} Declarar el límite y excluir rentabilidad o usar un indicador referencial correctamente rotulado.


\IfFileExists{capturas/BI08.tex}{\input{capturas/BI08.tex}}{}
\section{Resolver problemas y reconocer un rechazo correcto}
\label{unidad:QA01}
\textbf{Para quién:} Todos según permisos\par
\textbf{Al terminar podrá:} Recuperar el trabajo sin saltarse los controles y registrar incidencias útiles.

Que el sistema rechace una petición no siempre es un defecto. Debe rechazar estructuras inventadas, divisiones no definidas, mezclas de fuente o relaciones que contradigan el objetivo. La recuperación correcta conserva la evidencia y explica qué debe cambiar.

Ante un problema registre pantalla, fecha, fuente, propuesta, expediente, acción y mensaje. No copie contraseñas, claves API o credenciales a capturas o informes.

\subsection{Incidencias de esta edición: diagnóstico sencillo}
\Needspace{8\baselineskip}
\begingroup
\setbox0=\vbox{\hsize=\linewidth
{\normalfont\normalsize\sffamily\bfseries\normalsize\color{Teal}\hyphenpenalty=10000\exhyphenpenalty=10000 Cómo actuar sin crear otro problema\par}

\par
\begingroup
\fontsize{10}{12.5}\selectfont
\setlength{\tabcolsep}{5pt}
\renewcommand{\arraystretch}{1.15}
\begin{tabular}{@{}>{\RaggedRight\arraybackslash}p{0.22\linewidth}>{\RaggedRight\arraybackslash}p{0.3\linewidth}>{\RaggedRight\arraybackslash}p{0.39\linewidth}@{}}
\toprule
\textbf{Situación} & \textbf{Qué significa o revisar} & \textbf{Acción segura} \\
\midrule
Preguntas retiradas en un borrador & El catálogo de esa fuente cambió. No es una prueba de fallo del modelo. & Revise selección, repita viabilidad y confirme el alcance antes de generar. \\
\bottomrule
\end{tabular}
\endgroup
}
\Needspace{\dimexpr\ht0+\dp0+48pt\relax}
\endgroup
\subsection{Cómo actuar sin crear otro problema}
\begingroup
\fontsize{10}{12.5}\selectfont
\setlength{\tabcolsep}{5pt}
\renewcommand{\arraystretch}{1.15}
\begin{longtable}{@{}>{\RaggedRight\arraybackslash}p{0.22\linewidth}>{\RaggedRight\arraybackslash}p{0.3\linewidth}>{\RaggedRight\arraybackslash}p{0.39\linewidth}@{}}
\toprule
\textbf{Situación} & \textbf{Qué significa o revisar} & \textbf{Acción segura} \\
\midrule
\endfirsthead
\toprule
\textbf{Situación} & \textbf{Qué significa o revisar} & \textbf{Acción segura} \\
\midrule
\endhead
\midrule
\multicolumn{3}{r}{\footnotesize Continúa en la página siguiente}\\
\endfoot
\bottomrule
\endlastfoot
Preguntas retiradas en un borrador & El catálogo de esa fuente cambió. No es una prueba de fallo del modelo. & Revise selección, repita viabilidad y confirme el alcance antes de generar. \\
\addlinespace[3pt]
Sesión vencida & Terminó su duración absoluta. La tarea del servidor puede haber continuado. & Vuelva a ingresar; consulte estado y expediente antes de repetir una carga. \\
\addlinespace[3pt]
Respuesta estructurada inválida o incompleta & No cumple el contrato o alcanzó un límite. No se debe usar como propuesta válida. & Conserve la necesidad; revise aviso, modelo y presupuesto. Un reintento es una nueva llamada con posible costo. \\
\addlinespace[3pt]
Cuota o saldo insuficiente & La cuenta del proveedor no admite otra solicitud. & Revise Usage/Billing y el proyecto API. No cambie las bases ni repita solicitudes en bucle. \\
\addlinespace[3pt]
Una cifra no coincide con el otro panel & Puede haber distinto origen, ejecución, año, territorio, unidad o denominador. & Compare primero el encabezado y la tabla exacta; no sume ambas fuentes para ocultar la diferencia. \\
\addlinespace[3pt]
No aparece un KPI obligatorio & Encontrar la columna no demuestra cobertura analítica. & Revise cobertura y receta antes de aprobar. Si quedó parcial, corrija el alcance o el diseño. \\
\addlinespace[3pt]
Respuesta de chat pendiente al cambiar contexto & La respuesta anterior ya no corresponde al panel actual. & Espere los datos del contexto nuevo y formule la pregunta allí. No copie la respuesta vieja como si perteneciera a la nueva fuente. \\
\addlinespace[3pt]
\end{longtable}
\endgroup

\subsection{Conceptos para decidir con criterio}
\paragraph{Rechazo esperado} Control que impide una operación fuera de sus condiciones.

\paragraph{Regresión} Un comportamiento previamente correcto deja de cumplirse después de un cambio; necesita evidencia comparable.

\paragraph{Caso no ejecutado} Prueba planificada que no debe marcarse aprobada sin realizarla.

\Needspace{8\baselineskip}
\begingroup
\setbox0=\vbox{\hsize=\linewidth
{\normalfont\normalsize\sffamily\bfseries\normalsize\color{Teal}\hyphenpenalty=10000\exhyphenpenalty=10000 Campos y controles\par}

\par
\begingroup
\fontsize{10}{12.5}\selectfont
\setlength{\tabcolsep}{5pt}
\renewcommand{\arraystretch}{1.15}
\begin{tabular}{@{}>{\RaggedRight\arraybackslash}p{0.27\linewidth}>{\RaggedRight\arraybackslash}p{0.66\linewidth}@{}}
\toprule
\textbf{Campo o control} & \textbf{Significado y criterio} \\
\midrule
Reintentar la misma solicitud & Para una falla de proveedor, con revisión/consentimiento vigentes; no es un atajo del contrato. \\
\bottomrule
\end{tabular}
\endgroup
}
\Needspace{\dimexpr\ht0+\dp0+48pt\relax}
\endgroup
\subsection{Campos y controles}
\begingroup
\fontsize{10}{12.5}\selectfont
\setlength{\tabcolsep}{5pt}
\renewcommand{\arraystretch}{1.15}
\begin{longtable}{@{}>{\RaggedRight\arraybackslash}p{0.27\linewidth}>{\RaggedRight\arraybackslash}p{0.66\linewidth}@{}}
\toprule
\textbf{Campo o control} & \textbf{Significado y criterio} \\
\midrule
\endfirsthead
\toprule
\textbf{Campo o control} & \textbf{Significado y criterio} \\
\midrule
\endhead
\midrule
\multicolumn{2}{r}{\footnotesize Continúa en la página siguiente}\\
\endfoot
\bottomrule
\endlastfoot
Reintentar la misma solicitud & Para una falla de proveedor, con revisión/consentimiento vigentes; no es un atajo del contrato. \\
\addlinespace[3pt]
Revisar necesidad & Permite aclarar alcance sin modificar los datos de origen. \\
\addlinespace[3pt]
Abrir resultado / Abrir expediente & Recuperan evidencia persistida antes de repetir operaciones. \\
\addlinespace[3pt]
Motivo / Justificación & Explicación de una decisión humana, no código para cambiar estados. \\
\addlinespace[3pt]
\end{longtable}
\endgroup

\subsection{Paso a paso}
\begin{enumerate}
\item Determine en qué etapa ocurrió el problema: asistencia, diseño, aprobación, preparación, carga, semántica, panel, chat o exportación.\par\nopagebreak[4]\smallskip{\small\textbf{Compruebe:} No aplica el mismo remedio a todos los fallos.}
\item Lea el mensaje y compruebe fuente, versión y estado.\par\nopagebreak[4]\smallskip{\small\textbf{Compruebe:} Separa error de proveedor de una inconsistencia del contrato.}
\item Recupere el resultado persistido antes de reenviar.\par\nopagebreak[4]\smallskip{\small\textbf{Compruebe:} Evita duplicidad y consumo innecesario.}
\item Aplique la corrección mínima de la tabla de incidencias.\par\nopagebreak[4]\smallskip{\small\textbf{Compruebe:} No altera huellas, reglas o datos del origen.}
\item Repita el control afectado y registre resultado.\par\nopagebreak[4]\smallskip{\small\textbf{Compruebe:} La incidencia sólo se cierra con evidencia real.}
\end{enumerate}
\subsection{Qué hace la IA}
Puede fallar por formato, tiempo, cuota, capacidad o interpretación. Cambiar de proveedor puede ser una alternativa operativa, pero no debe cambiar las reglas de validación ni aprobar una respuesta inválida.

\subsection{Qué comprueba la plataforma}
Formato estructurado, requisitos de estado y permiso, coherencia de contrato y pertenencia de fuente. Los mensajes indican qué resultado no se usó o cuál acción quedó bloqueada.

\subsection{Qué decide la persona}
El operador recupera; el administrador revisa configuración; el analista corrige alcance y decisiones; la gerencia suspende conclusiones no sustentadas.

\paragraph{Resultado esperado} Trabajo recuperado o incidencia documentada con un punto de corrección identificable.

\Needspace{8\baselineskip}
\begingroup
\setbox0=\vbox{\hsize=\linewidth
{\normalfont\normalsize\sffamily\bfseries\normalsize\color{Teal}\hyphenpenalty=10000\exhyphenpenalty=10000 Errores y recuperación\par}

\par
\begingroup
\fontsize{10}{12.5}\selectfont
\setlength{\tabcolsep}{5pt}
\renewcommand{\arraystretch}{1.15}
\begin{tabular}{@{}>{\RaggedRight\arraybackslash}p{0.27\linewidth}>{\RaggedRight\arraybackslash}p{0.66\linewidth}@{}}
\toprule
\textbf{Situación} & \textbf{Qué hacer} \\
\midrule
El proveedor devolvió JSON pero no cumplió el contrato. & Se descarta su uso. Reintente con la misma revisión válida o coordine proveedor probado; no cargue a mano la respuesta inválida. \\
\bottomrule
\end{tabular}
\endgroup
}
\Needspace{\dimexpr\ht0+\dp0+48pt\relax}
\endgroup
\subsection{Errores y recuperación}
\begingroup
\fontsize{10}{12.5}\selectfont
\setlength{\tabcolsep}{5pt}
\renewcommand{\arraystretch}{1.15}
\begin{longtable}{@{}>{\RaggedRight\arraybackslash}p{0.27\linewidth}>{\RaggedRight\arraybackslash}p{0.66\linewidth}@{}}
\toprule
\textbf{Situación} & \textbf{Qué hacer} \\
\midrule
\endfirsthead
\toprule
\textbf{Situación} & \textbf{Qué hacer} \\
\midrule
\endhead
\midrule
\multicolumn{2}{r}{\footnotesize Continúa en la página siguiente}\\
\endfoot
\bottomrule
\endlastfoot
El proveedor devolvió JSON pero no cumplió el contrato. & Se descarta su uso. Reintente con la misma revisión válida o coordine proveedor probado; no cargue a mano la respuesta inválida. \\
\addlinespace[3pt]
La necesidad estaba respaldada pero el LLM no conserva el objetivo. & No apruebe; mantenga el texto original y reintente o corrija. Separar fallo de generación de falta de viabilidad. \\
\addlinespace[3pt]
Se pierde la respuesta del navegador. & Revise historial/versiones y expedientes antes de repetir; el servidor puede haber terminado. \\
\addlinespace[3pt]
No aparece una opción por rol. & Revise permisos con administrador; no haga que su compañera use otra identidad para simular un rol. \\
\addlinespace[3pt]
El filtro no da cifras. & Confirme que la combinación de año y territorio exista en el datamart; no cambie el cero a otro valor. \\
\addlinespace[3pt]
La respuesta del copiloto no es coherente. & Conserve pregunta, consulta interpretada, evidencia y contexto. Compare con el motor numérico y registre la corrección requerida. \\
\addlinespace[3pt]
\end{longtable}
\endgroup

\subsection{Ejercicio}
Escriba una incidencia que otra persona pueda reproducir sin conocer su contraseña.

\textbf{Respuesta orientativa.} Fuente, versión, expediente, condiciones previas, pasos, resultado esperado/obtenido y captura sin secretos.

\subsection{Compruebe lo que aprendió}
\paragraph{¿Cambiar el modelo resuelve una ruta ambigua por sí solo?} No. La elección del rol comercial requiere evidencia y decisión humana.

\paragraph{¿Un rechazo por facturación cuando se pidió pedido es siempre un defecto?} No. Puede ser un control correcto que protege el evento solicitado.

\begingroup
\setbox0=\vbox{\hsize=\linewidth
{\normalfont\normalsize\sffamily\bfseries\normalsize\color{Teal}\hyphenpenalty=10000\exhyphenpenalty=10000 Escenarios de comprobación y rechazo esperado\par}

\par
\begingroup
\fontsize{10}{12.5}\selectfont
\setlength{\tabcolsep}{5pt}
\renewcommand{\arraystretch}{1.15}
\begin{tabular}{@{}>{\RaggedRight\arraybackslash}p{0.27\linewidth}>{\RaggedRight\arraybackslash}p{0.66\linewidth}@{}}
\toprule
\textbf{Dato} & \textbf{Contenido} \\
\midrule
Caso & Cambiar AW\(\rightarrow\)WWI\(\rightarrow\)AW \\
\bottomrule
\end{tabular}
\endgroup
}
\Needspace{\dimexpr\ht0+\dp0+48pt\relax}
\endgroup
\subsection{Escenarios de comprobación y rechazo esperado}
\begingroup
\fontsize{10}{12.5}\selectfont
\setlength{\tabcolsep}{5pt}
\renewcommand{\arraystretch}{1.15}
\begin{longtable}{@{}>{\RaggedRight\arraybackslash}p{0.27\linewidth}>{\RaggedRight\arraybackslash}p{0.66\linewidth}@{}}
\toprule
\textbf{Dato} & \textbf{Contenido} \\
\midrule
\endfirsthead
\toprule
\textbf{Dato} & \textbf{Contenido} \\
\midrule
\endhead
\midrule
\multicolumn{2}{r}{\footnotesize Continúa en la página siguiente}\\
\endfoot
\bottomrule
\endlastfoot
Caso & Cambiar AW\(\rightarrow\)WWI\(\rightarrow\)AW \\
\addlinespace[3pt]
Resultado esperado & Fuente, ejecución, filtros, chat y panel corresponden al contexto nuevo; no se conservan datos de la otra fuente. \\
\addlinespace[3pt]
Tipo de comprobación & Aislamiento \\
\addlinespace[3pt]
Estado de la evidencia & Planificada para repetición por alumnas; existen antecedentes documentados. \\
\addlinespace[3pt]
\end{longtable}
\endgroup

\begingroup
\fontsize{10}{12.5}\selectfont
\setlength{\tabcolsep}{5pt}
\renewcommand{\arraystretch}{1.15}
\begin{longtable}{@{}>{\RaggedRight\arraybackslash}p{0.27\linewidth}>{\RaggedRight\arraybackslash}p{0.66\linewidth}@{}}
\toprule
\textbf{Dato} & \textbf{Contenido} \\
\midrule
\endfirsthead
\toprule
\textbf{Dato} & \textbf{Contenido} \\
\midrule
\endhead
\midrule
\multicolumn{2}{r}{\footnotesize Continúa en la página siguiente}\\
\endfoot
\bottomrule
\endlastfoot
Caso & Abrir versión aprobada y Verificar evidencia \\
\addlinespace[3pt]
Resultado esperado & Lectura y verificación sin cambiar silenciosamente estado, huellas o datos. \\
\addlinespace[3pt]
Tipo de comprobación & Historial \\
\addlinespace[3pt]
Estado de la evidencia & Planificada para repetición por alumnas. \\
\addlinespace[3pt]
\end{longtable}
\endgroup

\begingroup
\fontsize{10}{12.5}\selectfont
\setlength{\tabcolsep}{5pt}
\renewcommand{\arraystretch}{1.15}
\begin{longtable}{@{}>{\RaggedRight\arraybackslash}p{0.27\linewidth}>{\RaggedRight\arraybackslash}p{0.66\linewidth}@{}}
\toprule
\textbf{Dato} & \textbf{Contenido} \\
\midrule
\endfirsthead
\toprule
\textbf{Dato} & \textbf{Contenido} \\
\midrule
\endhead
\midrule
\multicolumn{2}{r}{\footnotesize Continúa en la página siguiente}\\
\endfoot
\bottomrule
\endlastfoot
Caso & Elegir contrato histórico sin rutas actuales para nueva carga \\
\addlinespace[3pt]
Resultado esperado & Nueva carga bloqueada con motivo; histórico sigue consultable. \\
\addlinespace[3pt]
Tipo de comprobación & Rechazo esperado \\
\addlinespace[3pt]
Estado de la evidencia & Regla contrastada con código. \\
\addlinespace[3pt]
\end{longtable}
\endgroup

\begingroup
\fontsize{10}{12.5}\selectfont
\setlength{\tabcolsep}{5pt}
\renewcommand{\arraystretch}{1.15}
\begin{longtable}{@{}>{\RaggedRight\arraybackslash}p{0.27\linewidth}>{\RaggedRight\arraybackslash}p{0.66\linewidth}@{}}
\toprule
\textbf{Dato} & \textbf{Contenido} \\
\midrule
\endfirsthead
\toprule
\textbf{Dato} & \textbf{Contenido} \\
\midrule
\endhead
\midrule
\multicolumn{2}{r}{\footnotesize Continúa en la página siguiente}\\
\endfoot
\bottomrule
\endlastfoot
Caso & Ruta WWI de BillToCustomerID en objetivo comprador \\
\addlinespace[3pt]
Resultado esperado & No se acepta como sustituto; usar decisión guiada con CustomerID. \\
\addlinespace[3pt]
Tipo de comprobación & Semántica \\
\addlinespace[3pt]
Estado de la evidencia & Antecedente histórico 6 noviembre; repetición nueva pendiente. \\
\addlinespace[3pt]
\end{longtable}
\endgroup

\begingroup
\fontsize{10}{12.5}\selectfont
\setlength{\tabcolsep}{5pt}
\renewcommand{\arraystretch}{1.15}
\begin{longtable}{@{}>{\RaggedRight\arraybackslash}p{0.27\linewidth}>{\RaggedRight\arraybackslash}p{0.66\linewidth}@{}}
\toprule
\textbf{Dato} & \textbf{Contenido} \\
\midrule
\endfirsthead
\toprule
\textbf{Dato} & \textbf{Contenido} \\
\midrule
\endhead
\midrule
\multicolumn{2}{r}{\footnotesize Continúa en la página siguiente}\\
\endfoot
\bottomrule
\endlastfoot
Caso & Ruta WWI de PostalCityID en objetivo ciudad de entrega \\
\addlinespace[3pt]
Resultado esperado & No se acepta como sustituto; DeliveryCityID con provincia conservada. \\
\addlinespace[3pt]
Tipo de comprobación & Semántica \\
\addlinespace[3pt]
Estado de la evidencia & Antecedente histórico 6 noviembre; repetición nueva pendiente. \\
\addlinespace[3pt]
\end{longtable}
\endgroup

\begingroup
\fontsize{10}{12.5}\selectfont
\setlength{\tabcolsep}{5pt}
\renewcommand{\arraystretch}{1.15}
\begin{longtable}{@{}>{\RaggedRight\arraybackslash}p{0.27\linewidth}>{\RaggedRight\arraybackslash}p{0.66\linewidth}@{}}
\toprule
\textbf{Dato} & \textbf{Contenido} \\
\midrule
\endfirsthead
\toprule
\textbf{Dato} & \textbf{Contenido} \\
\midrule
\endhead
\midrule
\multicolumn{2}{r}{\footnotesize Continúa en la página siguiente}\\
\endfoot
\bottomrule
\endlastfoot
Caso & Volver a preparar contrato/selección ya ejecutados \\
\addlinespace[3pt]
Resultado esperado & Abre expediente o conflicto de duplicidad; no otra carga idéntica accidental. \\
\addlinespace[3pt]
Tipo de comprobación & Rechazo esperado \\
\addlinespace[3pt]
Estado de la evidencia & Regla contrastada con código. \\
\addlinespace[3pt]
\end{longtable}
\endgroup

\begingroup
\fontsize{10}{12.5}\selectfont
\setlength{\tabcolsep}{5pt}
\renewcommand{\arraystretch}{1.15}
\begin{longtable}{@{}>{\RaggedRight\arraybackslash}p{0.27\linewidth}>{\RaggedRight\arraybackslash}p{0.66\linewidth}@{}}
\toprule
\textbf{Dato} & \textbf{Contenido} \\
\midrule
\endfirsthead
\toprule
\textbf{Dato} & \textbf{Contenido} \\
\midrule
\endhead
\midrule
\multicolumn{2}{r}{\footnotesize Continúa en la página siguiente}\\
\endfoot
\bottomrule
\endlastfoot
Caso & Divisor de una razón cero \\
\addlinespace[3pt]
Resultado esperado & Sin valor/no calculable; no infinito ni porcentaje inventado. \\
\addlinespace[3pt]
Tipo de comprobación & Cálculo \\
\addlinespace[3pt]
Estado de la evidencia & Regla contrastada con código; no se modifican datos para provocarla aquí. \\
\addlinespace[3pt]
\end{longtable}
\endgroup

\begingroup
\fontsize{10}{12.5}\selectfont
\setlength{\tabcolsep}{5pt}
\renewcommand{\arraystretch}{1.15}
\begin{longtable}{@{}>{\RaggedRight\arraybackslash}p{0.27\linewidth}>{\RaggedRight\arraybackslash}p{0.66\linewidth}@{}}
\toprule
\textbf{Dato} & \textbf{Contenido} \\
\midrule
\endfirsthead
\toprule
\textbf{Dato} & \textbf{Contenido} \\
\midrule
\endhead
\midrule
\multicolumn{2}{r}{\footnotesize Continúa en la página siguiente}\\
\endfoot
\bottomrule
\endlastfoot
Caso & Divisa WWI no demostrada \\
\addlinespace[3pt]
Resultado esperado & Moneda de origen, con importes conservados. \\
\addlinespace[3pt]
Tipo de comprobación & Unidad \\
\addlinespace[3pt]
Estado de la evidencia & Caso 22 documentado 6 noviembre. \\
\addlinespace[3pt]
\end{longtable}
\endgroup

\begingroup
\fontsize{10}{12.5}\selectfont
\setlength{\tabcolsep}{5pt}
\renewcommand{\arraystretch}{1.15}
\begin{longtable}{@{}>{\RaggedRight\arraybackslash}p{0.27\linewidth}>{\RaggedRight\arraybackslash}p{0.66\linewidth}@{}}
\toprule
\textbf{Dato} & \textbf{Contenido} \\
\midrule
\endfirsthead
\toprule
\textbf{Dato} & \textbf{Contenido} \\
\midrule
\endhead
\midrule
\multicolumn{2}{r}{\footnotesize Continúa en la página siguiente}\\
\endfoot
\bottomrule
\endlastfoot
Caso & Pedir pronóstico para próximo período \\
\addlinespace[3pt]
Resultado esperado & Explica fuera del alcance sin inventar cifra futura. \\
\addlinespace[3pt]
Tipo de comprobación & Límite de IA \\
\addlinespace[3pt]
Estado de la evidencia & Ejercicio supervisado; no se promete que toda respuesta libre sea correcta. \\
\addlinespace[3pt]
\end{longtable}
\endgroup

\begingroup
\fontsize{10}{12.5}\selectfont
\setlength{\tabcolsep}{5pt}
\renewcommand{\arraystretch}{1.15}
\begin{longtable}{@{}>{\RaggedRight\arraybackslash}p{0.27\linewidth}>{\RaggedRight\arraybackslash}p{0.66\linewidth}@{}}
\toprule
\textbf{Dato} & \textbf{Contenido} \\
\midrule
\endfirsthead
\toprule
\textbf{Dato} & \textbf{Contenido} \\
\midrule
\endhead
\midrule
\multicolumn{2}{r}{\footnotesize Continúa en la página siguiente}\\
\endfoot
\bottomrule
\endlastfoot
Caso & Exportar sin lectura del panel o sin permiso de exportación \\
\addlinespace[3pt]
Resultado esperado & Operación rechazada por servidor; no archivo con datos no autorizados. \\
\addlinespace[3pt]
Tipo de comprobación & Permisos \\
\addlinespace[3pt]
Estado de la evidencia & Regla contrastada con código; sesión manual por rol debe registrarse aparte. \\
\addlinespace[3pt]
\end{longtable}
\endgroup


\IfFileExists{capturas/QA01.tex}{\input{capturas/QA01.tex}}{}
\section{Demostrar y defender el recorrido completo}
\label{unidad:DEMO01}
\textbf{Para quién:} Analista de BI, Gerente comercial, Programadoras\par
\textbf{Al terminar podrá:} Explicar cada etapa con evidencia y responder preguntas del jurado sin prometer capacidades no implementadas.

Una demostración sólida no consiste en pulsar botones rápido. Explique la decisión de negocio, muestre su traducción al modelo y compruebe el resultado. Haga pausas en los puntos donde una persona acepta una decisión.

Practique primero la consulta del caso histórico y después una ejecución nueva cuando tenga tiempo y proveedor disponible. Lleve el PDF y Excel ya revisados como referencia, identificados con su fecha; no los presente como una descarga nueva si no lo es.

\subsection{Conceptos para decidir con criterio}
\paragraph{Eficacia observable} Cumplimiento de criterios funcionales y numéricos del caso.

\paragraph{Eficiencia comparativa} Requeriría tareas equivalentes y medidas de tiempo/esfuerzo/costo frente a una referencia. No se deduce de que el prototipo funcione.

\paragraph{Supervisión} Decisiones humanas explícitas entre asistencia de IA, reglas determinísticas y ejecución controlada.

\Needspace{8\baselineskip}
\begingroup
\setbox0=\vbox{\hsize=\linewidth
{\normalfont\normalsize\sffamily\bfseries\normalsize\color{Teal}\hyphenpenalty=10000\exhyphenpenalty=10000 Campos y controles\par}

\par
\begingroup
\fontsize{10}{12.5}\selectfont
\setlength{\tabcolsep}{5pt}
\renewcommand{\arraystretch}{1.15}
\begin{tabular}{@{}>{\RaggedRight\arraybackslash}p{0.27\linewidth}>{\RaggedRight\arraybackslash}p{0.66\linewidth}@{}}
\toprule
\textbf{Campo o control} & \textbf{Significado y criterio} \\
\midrule
Ficha de caso & Propósito, fuente, objetivo, grano, indicadores, rutas, unidad y filtros. \\
\bottomrule
\end{tabular}
\endgroup
}
\Needspace{\dimexpr\ht0+\dp0+48pt\relax}
\endgroup
\subsection{Campos y controles}
\begingroup
\fontsize{10}{12.5}\selectfont
\setlength{\tabcolsep}{5pt}
\renewcommand{\arraystretch}{1.15}
\begin{longtable}{@{}>{\RaggedRight\arraybackslash}p{0.27\linewidth}>{\RaggedRight\arraybackslash}p{0.66\linewidth}@{}}
\toprule
\textbf{Campo o control} & \textbf{Significado y criterio} \\
\midrule
\endfirsthead
\toprule
\textbf{Campo o control} & \textbf{Significado y criterio} \\
\midrule
\endhead
\midrule
\multicolumn{2}{r}{\footnotesize Continúa en la página siguiente}\\
\endfoot
\bottomrule
\endlastfoot
Ficha de caso & Propósito, fuente, objetivo, grano, indicadores, rutas, unidad y filtros. \\
\addlinespace[3pt]
Comprobantes & Aprobación, preparación, carga, conciliación y publicación de etiquetas. \\
\addlinespace[3pt]
Bitácora & Estado observado, hora, identificadores y archivos realmente revisados. \\
\addlinespace[3pt]
\end{longtable}
\endgroup

\subsection{Paso a paso}
\begin{enumerate}
\item Presente el problema como correspondencia entre necesidad comercial y datos disponibles.\par\nopagebreak[4]\smallskip{\small\textbf{Compruebe:} No afirma que faltaba IA en todas las empresas ni que la herramienta aumentó ventas.}
\item Muestre esquema y catálogo de cada fuente.\par\nopagebreak[4]\smallskip{\small\textbf{Compruebe:} Reconoce diferencias sin cambiar el proceso común.}
\item Explique necesidad y viabilidad.\par\nopagebreak[4]\smallskip{\small\textbf{Compruebe:} Distingue estructura, candidatos y límites.}
\item Lea modelo, personalización y aprobación.\par\nopagebreak[4]\smallskip{\small\textbf{Compruebe:} Se ven las decisiones humanas y las reglas que no pueden saltarse.}
\item Muestre preparación, carga y conciliación, nueva o histórica identificada.\par\nopagebreak[4]\smallskip{\small\textbf{Compruebe:} La audiencia reconoce de dónde salen los datos.}
\item Interprete el panel, use un filtro y pida una explicación al chat.\par\nopagebreak[4]\smallskip{\small\textbf{Compruebe:} Compara cifras y consulta interpretada sin atribuir causas.}
\item Exporte y abra los archivos, o muestre los históricos como tales.\par\nopagebreak[4]\smallskip{\small\textbf{Compruebe:} Cierra con contexto y trazabilidad.}
\item Explique una restricción: pronóstico fuera del alcance, moneda no comprobada o costo no comparable.\par\nopagebreak[4]\smallskip{\small\textbf{Compruebe:} Demuestra criterio y no infalibilidad universal.}
\end{enumerate}
\subsection{Qué hace la IA}
Asiste en formulación, sugerencias, diseño, consejo semántico, etiquetas aptas y conversación analítica. No reemplaza todos los cálculos ni la aprobación humana.

\subsection{Qué comprueba la plataforma}
Puertas de entrada, requisitos, contrato, relaciones, recetas, permisos, procedencia, carga y conciliación según el alcance implementado.

\subsection{Qué decide la persona}
La persona define, selecciona, justifica, aprueba, confirma, revisa etiquetas e interpreta el significado comercial del resultado.

\paragraph{Resultado esperado} Demostración que diferencia lo propuesto, lo comprobado, lo decidido y lo ejecutado.

\Needspace{8\baselineskip}
\begingroup
\setbox0=\vbox{\hsize=\linewidth
{\normalfont\normalsize\sffamily\bfseries\normalsize\color{Teal}\hyphenpenalty=10000\exhyphenpenalty=10000 Errores y recuperación\par}

\par
\begingroup
\fontsize{10}{12.5}\selectfont
\setlength{\tabcolsep}{5pt}
\renewcommand{\arraystretch}{1.15}
\begin{tabular}{@{}>{\RaggedRight\arraybackslash}p{0.27\linewidth}>{\RaggedRight\arraybackslash}p{0.66\linewidth}@{}}
\toprule
\textbf{Situación} & \textbf{Qué hacer} \\
\midrule
El proveedor falla frente al jurado. & Explique el rechazo seguro, conserve el estado y consulte el caso histórico identificado; no diga que una ejecución nueva terminó si no ocurrió. \\
\bottomrule
\end{tabular}
\endgroup
}
\Needspace{\dimexpr\ht0+\dp0+48pt\relax}
\endgroup
\subsection{Errores y recuperación}
\begingroup
\fontsize{10}{12.5}\selectfont
\setlength{\tabcolsep}{5pt}
\renewcommand{\arraystretch}{1.15}
\begin{longtable}{@{}>{\RaggedRight\arraybackslash}p{0.27\linewidth}>{\RaggedRight\arraybackslash}p{0.66\linewidth}@{}}
\toprule
\textbf{Situación} & \textbf{Qué hacer} \\
\midrule
\endfirsthead
\toprule
\textbf{Situación} & \textbf{Qué hacer} \\
\midrule
\endhead
\midrule
\multicolumn{2}{r}{\footnotesize Continúa en la página siguiente}\\
\endfoot
\bottomrule
\endlastfoot
El proveedor falla frente al jurado. & Explique el rechazo seguro, conserve el estado y consulte el caso histórico identificado; no diga que una ejecución nueva terminó si no ocurrió. \\
\addlinespace[3pt]
El jurado pregunta por otro motor o idioma. & Explique arquitectura adaptable y límites comprobados. Dos fuentes SQL Server y variantes de pruebas no certifican cualquier motor, idioma o modelo. \\
\addlinespace[3pt]
\end{longtable}
\endgroup

\subsection{Ejercicio}
Prepare una exposición de 10 minutos y otra de 25 minutos; en ambas identifique qué pasos consulta y qué pasos ejecuta.

\textbf{Respuesta orientativa.} La corta puede usar expedientes históricos verificados; la larga puede generar un caso nuevo con margen para revisión sin prometer una respuesta fija.

\subsection{Compruebe lo que aprendió}
\paragraph{¿Qué hace inteligente al sistema?} Combina interpretación y propuestas de IA con contexto, reglas, memoria de expedientes y supervisión; no que toda acción la realice un LLM.

\paragraph{¿Qué evidencia usaría para demostrar una cifra?} Fuente, propuesta, ejecución, alcance, fórmula, conciliación y tabla/reporte del mismo universo.

\paragraph{¿Qué no se puede prometer?} Infalibilidad LLM, pronóstico, eficiencia comparativa no medida, compatibilidad de todos los motores ni efectos económicos reales no investigados.


\IfFileExists{capturas/DEMO01.tex}{\input{capturas/DEMO01.tex}}{}
\section{Preguntas para explicarlo con sus propias palabras}
\textbf{Al terminar podrá:} Defender el propósito y las decisiones del sistema sin memorizar frases.

Conteste primero sin mirar la respuesta; luego señale una pantalla, cálculo o diagrama que sostenga su explicación. Las respuestas son orientativas y deben relacionarse con el contexto realmente observado.

\paragraph{¿Qué problema atiende el proyecto y cómo demostraría que lo atiende?} Ayuda a relacionar una necesidad comercial con fuentes, grano y cálculos verificables mediante un recorrido supervisado. Se demuestra siguiendo un caso y contrastando sus resultados. Que funcione no demuestra que aumente ventas ni que sea más rápido que otro método; esa comparación exigiría mediciones equivalentes.

\textbf{Cómo sostener la respuesta:} Explique una decisión comercial y señale propuesta, ejecución y conciliación del caso correspondiente.

\paragraph{Si un pedido tiene tres líneas, ¿cuántas filas del hecho y cuántos pedidos distintos representa?} A grano de línea son tres filas y un pedido distinto. Contar filas o sumar el identificador del pedido no da la cantidad de pedidos.

\textbf{Cómo sostener la respuesta:} Dibuje un pedido con tres productos y escriba ambos conteos.

\paragraph{¿Elegir Mensual cambia el detalle físico o programa una carga mensual?} Mensual indica cómo agrupar el análisis. Las filas pueden seguir siendo líneas de pedido o factura y la carga se confirma por separado. No significa que exista un programador automático.

\textbf{Cómo sostener la respuesta:} Compare el grano declarado en la propuesta con la agrupación temporal del panel.

\paragraph{¿Qué aprende la IA al recibir metadatos y qué todavía no está demostrado?} Recibe estructura autorizada: nombres, tipos, claves y relaciones. Esto puede respaldar una ruta o receta, pero no prueba que los valores estén completos, que un costo sea histórico o que una dirección represente una entrega realizada.

\textbf{Cómo sostener la respuesta:} Abra una tabla y una relación del explorador; explique su alcance sin consultar filas de ventas.

\paragraph{¿La IA arma por sí sola el catálogo de preguntas y lo copia a todas las bases?} El perfil del producto aporta cuatro orientaciones base; la persona autorizada administra preguntas y periodicidades por conexión y dominio. Las dos preguntas de costos de AdventureWorks fueron adiciones humanas históricas. Las sugerencias de necesidades del asistente no sustituyen automáticamente ese catálogo.

\textbf{Cómo sostener la respuesta:} Compare ambos catálogos y justifique una pregunta apropiada para cada fuente.

\paragraph{¿Qué haría con un requisito marcado como derivable?} Leer si depende de una relación declarada o de una fórmula controlada, comprobar sus entradas y conservar el grano. Una explicación candidata de IA no se convierte por sí sola en cálculo ejecutable. Si falta una decisión de negocio, hay que aclararla antes de aprobar.

\textbf{Cómo sostener la respuesta:} Explique una derivación posible y sus entradas; no escriba SQL libre para quitar el aviso.

\paragraph{¿Que la necesidad tenga respaldo estructural garantiza que la propuesta se genere correctamente?} No. La revisión de viabilidad y la generación son controles distintos. El proveedor puede fallar en formato, cuota o interpretación. Aceptar una limitación tampoco crea datos ausentes ni elimina los errores bloqueantes.

\textbf{Cómo sostener la respuesta:} Distinga una revisión válida de una generación rechazada y registre cada estado por separado.

\paragraph{¿Qué hacen la IA, la plataforma y la persona en el mismo caso?} La IA puede sugerir, interpretar y explicar. La plataforma valida referencias y contratos, construye consultas controladas y calcula. La persona define el propósito, resuelve ambigüedades, aprueba el diseño, confirma la carga y revisa el resultado.

\textbf{Cómo sostener la respuesta:} Señale una acción concreta de cada participante en el asistente y en el panel.

\paragraph{¿Puede llamar facturas, cobros o utilidades a las cifras de un datamart de pedidos?} Debe conservar el evento demostrado. Un pedido no prueba facturación ni cobro; un importe de venta no prueba utilidad. En WideWorldImporters el caso de referencia usa líneas de factura y en AdventureWorks líneas de pedido.

\textbf{Cómo sostener la respuesta:} Lea el evento y el grano de ambas necesidades antes de comparar sus tarjetas.

\paragraph{¿Por qué importan comprador, destinatario de facturación y ciudad de entrega?} Pueden representar funciones distintas aunque las relaciones sean válidas y los totales coincidan. El caso WWI pide comprador y provincia de la ciudad de entrega de su ficha actual: no sustituir por destinatario de facturación o dirección postal ni afirmar ubicación histórica.

\textbf{Cómo sostener la respuesta:} Explique la ruta revisada de WWI: CustomerID, DeliveryCityID y StateProvinceID, usando las etiquetas de la interfaz.

\paragraph{¿Cómo obtiene el importe promedio por pedido o factura?} Divide el importe total entre documentos distintos del mismo filtro. No usa el promedio de importes de línea ni un denominador de unidades. Si el denominador es cero, debe quedar sin valor calculable; no inventar cero, infinito o un porcentaje.

\textbf{Cómo sostener la respuesta:} Compruebe numerador, denominador, unidad y filtro de una razón visible.

\paragraph{¿Por qué el importe de WWI no se presenta automáticamente como dólares?} El caso conservado no demuestra un código de moneda único. Debe decir moneda de origen y conservar el importe. En AdventureWorks el caso de referencia sí comprobó USD. No se suman ambos importes como si compartieran unidad y evento.

\textbf{Cómo sostener la respuesta:} Lea la unidad de las dos ejecuciones y explique el límite antes de emitir una comparación monetaria.

\paragraph{¿Qué debe revisar antes de interpretar un producto líder o una participación?} Fuente, ejecución, indicador y filtros. Un ranking por sumas del conjunto no equivale al máximo de una sola línea. La participación utiliza el universo filtrado completo; el Top N no convierte al subconjunto mostrado en el total del negocio.

\textbf{Cómo sostener la respuesta:} Cambie un filtro y compruebe que tarjeta, gráfico, tabla y explicación correspondan al nuevo alcance.

\paragraph{¿Qué se obtiene al aprobar y qué exige cargar el datamart?} La aprobación conserva un contrato revisado. La carga necesita preparar la selección, revisar el plan, confirmar la ejecución y comprobar la conciliación. Tener una propuesta aprobada no significa que ya existan datos analíticos.

\textbf{Cómo sostener la respuesta:} Muestre las identidades separadas de propuesta y ejecución sin volver a ejecutar un expediente histórico.

\paragraph{¿Puede consultar o verificar una versión aprobada sin retirarla?} La consulta y la verificación deben conservar el estado. Una decisión de retiro o restauración es explícita y auditada. Si las reglas actuales requieren una versión nueva para otra carga, el histórico no se corrige sobrescribiéndolo.

\textbf{Cómo sostener la respuesta:} Anote el estado antes y después de una consulta autorizada; si cambia sin decisión, registre la incidencia.

\paragraph{¿Qué debe ocurrir al pasar de AdventureWorks a WideWorldImporters?} Debe cambiar el contexto de fuente y seleccionarse un resultado de esa base. Filtros, ejecución y conversación anteriores no deben interpretarse como pertenecientes a la nueva fuente. Cambiar fuente no mezcla ni desactiva ambas conexiones.

\textbf{Cómo sostener la respuesta:} Repita el recorrido de ida y vuelta y lea los identificadores del encabezado antes de copiar cifras.

\paragraph{¿Un menú visible concede permisos? ¿Qué puede hacer el Gerente comercial observado?} FastAPI autoriza cada operación; el menú sólo refleja permisos. El rol observado puede consultar el panel, conexiones y metadatos y exportar análisis. Consultar metadatos ayuda al selector y permite el explorador; no da escritura, aprobación ni ETL. Una cuenta con varios roles puede acumular permisos.

\textbf{Cómo sostener la respuesta:} Revise la matriz de roles. Si prueba una cuenta, registre su identidad de práctica y sus roles reales; no simule gerencia usando una sesión administrativa.

\paragraph{¿Qué distingue una FK real de un enlace dentro de JSON o del datamart?} La FK es una restricción del motor. Los documentos JSONB guardan diseños y evidencia y requieren validación de aplicación. En e21/e22 las dimensiones tienen PK, pero fact\_\allowbreak{}ventas no tiene PK ni FK físicas; sus enlaces a dimensiones y etl\_\allowbreak{}execution\_\allowbreak{}id son lógicos. No dibuje restricciones inexistentes.

\textbf{Cómo sostener la respuesta:} Explique una línea continua de app y una discontinua de los datamarts en los diagramas.

\paragraph{¿Una instalación limpia debe tener 83 tablas y las propuestas 129 y 131?} No. Ese inventario describe historial acumulado local. La estructura inicial de app se crea con migraciones y las tablas analíticas aparecen al cargar. Las propuestas y ejecuciones reciben identificadores propios; los casos 129/21 y 131/22 son referencias de una instalación, no valores que haya que forzar.

\textbf{Cómo sostener la respuesta:} Registre los identificadores propios y diferencie consulta histórica de ejecución nueva.

\paragraph{¿Qué dice y hace si falla el proveedor frente al jurado?} Conserva el mensaje, comprueba si la solicitud terminó y evita reenviarla mientras siga pendiente. No aprueba una respuesta inválida. Puede consultar evidencia histórica anunciándola como tal o continuar después con un proveedor autorizado y nueva revisión del destino. Un rechazo correcto no es una carga exitosa ni una prueba de que cualquier modelo funcionará.

\textbf{Cómo sostener la respuesta:} Ensaye una explicación verbal de contingencia y complete un registro con estado real, sin secretos ni resultados inventados.

\section{Diez minutos: leer y explicar ventas como gerente}
\textbf{Al terminar podrá:} Explicar una decisión comercial con resultados de ambas fuentes, mostrando el contexto y sus límites.

Guion de aprendizaje propuesto; no ejecutado durante esta redacción.

Consulta de datos conservados y, si se realiza, chat y exportación nuevos. No genera propuestas ni ejecuta ETL.

Plataforma ya preparada, con acceso autorizado y datamarts conciliados físicamente disponibles de ambas fuentes. Esta ruta no parte de una aplicación vacía.

Identificar los expedientes de referencia disponibles o sus equivalentes comprobados. En la instalación documentada son AW 129/21 y WWI 131/22; en otra pueden cambiar.

Usar una cuenta de demostración con rol y permisos realmente asignados. La matriz observada describe el rol; este guion no acredita una prueba previa de login con esa cuenta.

Para chat nuevo, proveedor ya activo y probado y consumo autorizado. Conservar como respaldo archivos históricos revisados y rotulados con su fecha.

\Needspace{8\baselineskip}
\begingroup
\setbox0=\vbox{\hsize=\linewidth
{\normalfont\normalsize\sffamily\bfseries\normalsize\color{Teal}\hyphenpenalty=10000\exhyphenpenalty=10000 Tiempo previsto y alternativas\par}

\par
\begingroup
\fontsize{10}{12.5}\selectfont
\setlength{\tabcolsep}{5pt}
\renewcommand{\arraystretch}{1.15}
\begin{tabular}{@{}>{\RaggedRight\arraybackslash}p{0.27\linewidth}>{\RaggedRight\arraybackslash}p{0.66\linewidth}@{}}
\toprule
\textbf{Condición} & \textbf{Cómo organizarse} \\
\midrule
Tiempo total & 10 minutos de exposición; no garantiza disponibilidad del servicio externo. \\
\bottomrule
\end{tabular}
\endgroup
}
\Needspace{\dimexpr\ht0+\dp0+48pt\relax}
\endgroup
\subsection{Tiempo previsto y alternativas}
\begingroup
\fontsize{10}{12.5}\selectfont
\setlength{\tabcolsep}{5pt}
\renewcommand{\arraystretch}{1.15}
\begin{longtable}{@{}>{\RaggedRight\arraybackslash}p{0.27\linewidth}>{\RaggedRight\arraybackslash}p{0.66\linewidth}@{}}
\toprule
\textbf{Condición} & \textbf{Cómo organizarse} \\
\midrule
\endfirsthead
\toprule
\textbf{Condición} & \textbf{Cómo organizarse} \\
\midrule
\endhead
\midrule
\multicolumn{2}{r}{\footnotesize Continúa en la página siguiente}\\
\endfoot
\bottomrule
\endlastfoot
Tiempo total & 10 minutos de exposición; no garantiza disponibilidad del servicio externo. \\
\addlinespace[3pt]
Intentos chat manuales & 1 \\
\addlinespace[3pt]
Espera chat dentro del bloque & Como máximo el tiempo del bloque de dos minutos; no modificar el timeout técnico ni reenviar mientras siga pendiente. \\
\addlinespace[3pt]
Alternativa & Explicar la falla y consultar un archivo histórico identificado. Eso no completa una consulta de IA nueva. \\
\addlinespace[3pt]
\end{longtable}
\endgroup

\begingroup
\setbox0=\vbox{\hsize=\dimexpr\linewidth-1.7em\relax
\everypar{}\csname @newlistfalse\endcsname
\textbf{Minutos 0-1.} Presente la decisión: comparar unidades y documentos dentro de cada fuente. Lea la cuenta y el contexto de trabajo.

{\small\textbf{Tiempo máximo previsto de este bloque:} 1 minuto.\par}
\textbf{Qué debe observar:} La audiencia entiende el propósito y sabe con qué rol real se trabaja.

\textbf{Qué registrar:} Fuente y rol observados; indicar que los datos cargados son de un expediente conservado.

\textbf{Si no se cumple:} Si sólo hay una sesión administrativa disponible, diga que usa administración para la lectura; no la presente como validación del rol gerente.

\par
}
\dimen0=\dimexpr\ht0+\dp0+4\baselineskip\relax
\ifdim\dimen0<\textheight\Needspace{\dimen0}\fi
\endgroup
\subsection{Guion paso a paso}
\begin{enumerate}
\begingroup
\setbox0=\vbox{\hsize=\linewidth
\everypar{}\csname @newlistfalse\endcsname
\textbf{Minutos 0-1.} Presente la decisión: comparar unidades y documentos dentro de cada fuente. Lea la cuenta y el contexto de trabajo.

{\small\textbf{Tiempo máximo previsto de este bloque:} 1 minuto.\par}
\textbf{Qué debe observar:} La audiencia entiende el propósito y sabe con qué rol real se trabaja.

\textbf{Qué registrar:} Fuente y rol observados; indicar que los datos cargados son de un expediente conservado.

\textbf{Si no se cumple:} Si sólo hay una sesión administrativa disponible, diga que usa administración para la lectura; no la presente como validación del rol gerente.

\par
}
\dimen0=\dimexpr\ht0+\dp0+2\baselineskip\relax
\ifdim\dimen0<\textheight
\Needspace{\dimen0}
\endgroup
\item \begin{samepage}
\textbf{Minutos 0-1.} Presente la decisión: comparar unidades y documentos dentro de cada fuente. Lea la cuenta y el contexto de trabajo.

{\small\textbf{Tiempo máximo previsto de este bloque:} 1 minuto.\par}
\textbf{Qué debe observar:} La audiencia entiende el propósito y sabe con qué rol real se trabaja.

\textbf{Qué registrar:} Fuente y rol observados; indicar que los datos cargados son de un expediente conservado.

\textbf{Si no se cumple:} Si sólo hay una sesión administrativa disponible, diga que usa administración para la lectura; no la presente como validación del rol gerente.

\end{samepage}
\else
\endgroup
\item \textbf{Minutos 0-1.} Presente la decisión: comparar unidades y documentos dentro de cada fuente. Lea la cuenta y el contexto de trabajo.

{\small\textbf{Tiempo máximo previsto de este bloque:} 1 minuto.\par}
\textbf{Qué debe observar:} La audiencia entiende el propósito y sabe con qué rol real se trabaja.

\textbf{Qué registrar:} Fuente y rol observados; indicar que los datos cargados son de un expediente conservado.

\textbf{Si no se cumple:} Si sólo hay una sesión administrativa disponible, diga que usa administración para la lectura; no la presente como validación del rol gerente.

\fi
\begingroup
\setbox0=\vbox{\hsize=\linewidth
\everypar{}\csname @newlistfalse\endcsname
\textbf{Minutos 1-3.} Abra Analítica de ventas de AdventureWorks. Seleccione el expediente correcto y el filtro 2013/Northwest; espere que termine la actualización.

{\small\textbf{Tiempo máximo previsto de este bloque:} 2 minutos.\par}
\textbf{Qué debe observar:} En el caso histórico equivalente: 17.452 unidades y 2.053 pedidos; el importe se expresa como USD comprobado. Las tarjetas abreviadas se contrastan con su respaldo exacto.

\textbf{Qué registrar:} Captura del contexto y al menos una cifra exacta. Referencia histórica: 6.015.174,74 USD; 2.929,94 USD por pedido.

\textbf{Si no se cumple:} Si el alcance o la copia no coincide, no fuerce esas cifras. Registre la diferencia o limite la lectura a lo que realmente esté disponible.

\par
}
\dimen0=\dimexpr\ht0+\dp0+2\baselineskip\relax
\ifdim\dimen0<\textheight
\Needspace{\dimen0}
\endgroup
\item \begin{samepage}
\textbf{Minutos 1-3.} Abra Analítica de ventas de AdventureWorks. Seleccione el expediente correcto y el filtro 2013/Northwest; espere que termine la actualización.

{\small\textbf{Tiempo máximo previsto de este bloque:} 2 minutos.\par}
\textbf{Qué debe observar:} En el caso histórico equivalente: 17.452 unidades y 2.053 pedidos; el importe se expresa como USD comprobado. Las tarjetas abreviadas se contrastan con su respaldo exacto.

\textbf{Qué registrar:} Captura del contexto y al menos una cifra exacta. Referencia histórica: 6.015.174,74 USD; 2.929,94 USD por pedido.

\textbf{Si no se cumple:} Si el alcance o la copia no coincide, no fuerce esas cifras. Registre la diferencia o limite la lectura a lo que realmente esté disponible.

\end{samepage}
\else
\endgroup
\item \textbf{Minutos 1-3.} Abra Analítica de ventas de AdventureWorks. Seleccione el expediente correcto y el filtro 2013/Northwest; espere que termine la actualización.

{\small\textbf{Tiempo máximo previsto de este bloque:} 2 minutos.\par}
\textbf{Qué debe observar:} En el caso histórico equivalente: 17.452 unidades y 2.053 pedidos; el importe se expresa como USD comprobado. Las tarjetas abreviadas se contrastan con su respaldo exacto.

\textbf{Qué registrar:} Captura del contexto y al menos una cifra exacta. Referencia histórica: 6.015.174,74 USD; 2.929,94 USD por pedido.

\textbf{Si no se cumple:} Si el alcance o la copia no coincide, no fuerce esas cifras. Registre la diferencia o limite la lectura a lo que realmente esté disponible.

\fi
\begingroup
\setbox0=\vbox{\hsize=\linewidth
\everypar{}\csname @newlistfalse\endcsname
\textbf{Minutos 3-5.} Cambie a WideWorldImporters y seleccione la ejecución de esa fuente. Aplique 2015/California y lea la unidad.

{\small\textbf{Tiempo máximo previsto de este bloque:} 2 minutos.\par}
\textbf{Qué debe observar:} En el caso histórico equivalente: 142.247 unidades y 1.179 facturas; 3.191.769,11 en moneda de origen y 2.707,18 por factura. El territorio procede de la ficha actual del comprador.

\textbf{Qué registrar:} Fuente, ejecución y filtros nuevos. Señalar que pedido y factura son eventos distintos.

\textbf{Si no se cumple:} Si sigue mostrando el contexto anterior al terminar la carga de pantalla, suspenda la conclusión y registre incidencia de aislamiento.

\par
}
\dimen0=\dimexpr\ht0+\dp0+2\baselineskip\relax
\ifdim\dimen0<\textheight
\Needspace{\dimen0}
\endgroup
\item \begin{samepage}
\textbf{Minutos 3-5.} Cambie a WideWorldImporters y seleccione la ejecución de esa fuente. Aplique 2015/California y lea la unidad.

{\small\textbf{Tiempo máximo previsto de este bloque:} 2 minutos.\par}
\textbf{Qué debe observar:} En el caso histórico equivalente: 142.247 unidades y 1.179 facturas; 3.191.769,11 en moneda de origen y 2.707,18 por factura. El territorio procede de la ficha actual del comprador.

\textbf{Qué registrar:} Fuente, ejecución y filtros nuevos. Señalar que pedido y factura son eventos distintos.

\textbf{Si no se cumple:} Si sigue mostrando el contexto anterior al terminar la carga de pantalla, suspenda la conclusión y registre incidencia de aislamiento.

\end{samepage}
\else
\endgroup
\item \textbf{Minutos 3-5.} Cambie a WideWorldImporters y seleccione la ejecución de esa fuente. Aplique 2015/California y lea la unidad.

{\small\textbf{Tiempo máximo previsto de este bloque:} 2 minutos.\par}
\textbf{Qué debe observar:} En el caso histórico equivalente: 142.247 unidades y 1.179 facturas; 3.191.769,11 en moneda de origen y 2.707,18 por factura. El territorio procede de la ficha actual del comprador.

\textbf{Qué registrar:} Fuente, ejecución y filtros nuevos. Señalar que pedido y factura son eventos distintos.

\textbf{Si no se cumple:} Si sigue mostrando el contexto anterior al terminar la carga de pantalla, suspenda la conclusión y registre incidencia de aislamiento.

\fi
\begingroup
\setbox0=\vbox{\hsize=\linewidth
\everypar{}\csname @newlistfalse\endcsname
\textbf{Minutos 5-7.} Muestre una tabla de respaldo o exporte el alcance visible y abra el archivo. Explique por qué los datos no son una predicción ni demuestran utilidad.

{\small\textbf{Tiempo máximo previsto de este bloque:} 2 minutos.\par}
\textbf{Qué debe observar:} El archivo y la lectura tienen la misma fuente, ejecución, filtro, indicador y unidad.

\textbf{Qué registrar:} Nombre del archivo realmente abierto y su contexto. Si utiliza un archivo anterior, indicarlo como histórico.

\textbf{Si no se cumple:} Si no se puede exportar o abrir, marque el paso no completado. Una captura antigua no se anuncia como una exportación nueva.

\par
}
\dimen0=\dimexpr\ht0+\dp0+2\baselineskip\relax
\ifdim\dimen0<\textheight
\Needspace{\dimen0}
\endgroup
\item \begin{samepage}
\textbf{Minutos 5-7.} Muestre una tabla de respaldo o exporte el alcance visible y abra el archivo. Explique por qué los datos no son una predicción ni demuestran utilidad.

{\small\textbf{Tiempo máximo previsto de este bloque:} 2 minutos.\par}
\textbf{Qué debe observar:} El archivo y la lectura tienen la misma fuente, ejecución, filtro, indicador y unidad.

\textbf{Qué registrar:} Nombre del archivo realmente abierto y su contexto. Si utiliza un archivo anterior, indicarlo como histórico.

\textbf{Si no se cumple:} Si no se puede exportar o abrir, marque el paso no completado. Una captura antigua no se anuncia como una exportación nueva.

\end{samepage}
\else
\endgroup
\item \textbf{Minutos 5-7.} Muestre una tabla de respaldo o exporte el alcance visible y abra el archivo. Explique por qué los datos no son una predicción ni demuestran utilidad.

{\small\textbf{Tiempo máximo previsto de este bloque:} 2 minutos.\par}
\textbf{Qué debe observar:} El archivo y la lectura tienen la misma fuente, ejecución, filtro, indicador y unidad.

\textbf{Qué registrar:} Nombre del archivo realmente abierto y su contexto. Si utiliza un archivo anterior, indicarlo como histórico.

\textbf{Si no se cumple:} Si no se puede exportar o abrir, marque el paso no completado. Una captura antigua no se anuncia como una exportación nueva.

\fi
\begingroup
\setbox0=\vbox{\hsize=\linewidth
\everypar{}\csname @newlistfalse\endcsname
\textbf{Minutos 7-9.} Haga una pregunta concreta al copiloto sobre WWI 2015/California. Lea Consulta interpretada y contraste una cifra con el panel.

{\small\textbf{Tiempo máximo previsto de este bloque:} 2 minutos.\par}
\textbf{Qué debe observar:} La explicación conserva alcance, denominador y cautelas; el texto no atribuye cobro, entrega efectiva ni causalidad no demostrada.

\textbf{Qué registrar:} Pregunta, estado de respuesta y una comparación numérica. Puede usar la pregunta de referencia del caso WWI.

\textbf{Si no se cumple:} Un intento manual. Si falla o permanece pendiente al terminar el bloque, registre el estado y explique la contingencia; el panel sigue siendo una evidencia independiente del chat.

\par
}
\dimen0=\dimexpr\ht0+\dp0+2\baselineskip\relax
\ifdim\dimen0<\textheight
\Needspace{\dimen0}
\endgroup
\item \begin{samepage}
\textbf{Minutos 7-9.} Haga una pregunta concreta al copiloto sobre WWI 2015/California. Lea Consulta interpretada y contraste una cifra con el panel.

{\small\textbf{Tiempo máximo previsto de este bloque:} 2 minutos.\par}
\textbf{Qué debe observar:} La explicación conserva alcance, denominador y cautelas; el texto no atribuye cobro, entrega efectiva ni causalidad no demostrada.

\textbf{Qué registrar:} Pregunta, estado de respuesta y una comparación numérica. Puede usar la pregunta de referencia del caso WWI.

\textbf{Si no se cumple:} Un intento manual. Si falla o permanece pendiente al terminar el bloque, registre el estado y explique la contingencia; el panel sigue siendo una evidencia independiente del chat.

\end{samepage}
\else
\endgroup
\item \textbf{Minutos 7-9.} Haga una pregunta concreta al copiloto sobre WWI 2015/California. Lea Consulta interpretada y contraste una cifra con el panel.

{\small\textbf{Tiempo máximo previsto de este bloque:} 2 minutos.\par}
\textbf{Qué debe observar:} La explicación conserva alcance, denominador y cautelas; el texto no atribuye cobro, entrega efectiva ni causalidad no demostrada.

\textbf{Qué registrar:} Pregunta, estado de respuesta y una comparación numérica. Puede usar la pregunta de referencia del caso WWI.

\textbf{Si no se cumple:} Un intento manual. Si falla o permanece pendiente al terminar el bloque, registre el estado y explique la contingencia; el panel sigue siendo una evidencia independiente del chat.

\fi
\begingroup
\setbox0=\vbox{\hsize=\linewidth
\everypar{}\csname @newlistfalse\endcsname
\textbf{Minutos 9-10.} Cierre con una conclusión comercial acotada y un límite. Explique qué tendría que pedir al analista para otra pregunta.

{\small\textbf{Tiempo máximo previsto de este bloque:} 1 minuto.\par}
\textbf{Qué debe observar:} Se interpreta el dato observado sin afirmar aumento de ventas, moneda común o compatibilidad universal.

\textbf{Qué registrar:} Conclusión breve y una limitación explícita.

\textbf{Si no se cumple:} No convierta una impresión en un resultado medido ni sume importes de ambas fuentes.

\par
}
\dimen0=\dimexpr\ht0+\dp0+2\baselineskip\relax
\ifdim\dimen0<\textheight
\Needspace{\dimen0}
\endgroup
\item \begin{samepage}
\textbf{Minutos 9-10.} Cierre con una conclusión comercial acotada y un límite. Explique qué tendría que pedir al analista para otra pregunta.

{\small\textbf{Tiempo máximo previsto de este bloque:} 1 minuto.\par}
\textbf{Qué debe observar:} Se interpreta el dato observado sin afirmar aumento de ventas, moneda común o compatibilidad universal.

\textbf{Qué registrar:} Conclusión breve y una limitación explícita.

\textbf{Si no se cumple:} No convierta una impresión en un resultado medido ni sume importes de ambas fuentes.

\end{samepage}
\else
\endgroup
\item \textbf{Minutos 9-10.} Cierre con una conclusión comercial acotada y un límite. Explique qué tendría que pedir al analista para otra pregunta.

{\small\textbf{Tiempo máximo previsto de este bloque:} 1 minuto.\par}
\textbf{Qué debe observar:} Se interpreta el dato observado sin afirmar aumento de ventas, moneda común o compatibilidad universal.

\textbf{Qué registrar:} Conclusión breve y una limitación explícita.

\textbf{Si no se cumple:} No convierta una impresión en un resultado medido ni sume importes de ambas fuentes.

\fi
\end{enumerate}
\subsection{Antes de dar por terminado el ejercicio}
\begin{itemize}
\item Se identificaron las dos fuentes y sus contextos.
\item Se contrastó al menos una cifra por fuente con su alcance.
\item Se diferenció una operación nueva de un expediente o archivo histórico.
\item Las fallas o pasos omitidos quedaron registrados con su estado real.
\end{itemize}

\section{Veinticinco minutos: construir y justificar un caso como analista}
\textbf{Al terminar podrá:} Recorrer un caso nuevo de una fuente, desde la lectura del esquema hasta la conciliación y su lectura analítica.

Guion de aprendizaje propuesto; no ejecutado durante esta redacción.

Una generación y una carga nuevas de una fuente si superan sus controles; la consulta histórica es una alternativa de exposición y no reemplaza su resultado.

Administrador preparó conexiones, permisos y proveedor activo probado. El analista no necesita ni recibe claves técnicas para operar.

Elegir una sola fuente para la ejecución nueva. Repetir el mismo guion en otra sesión para la segunda; no prometer dos cargas nuevas dentro de 25 minutos.

Utilizar la necesidad avanzada de referencia de AW o WWI, con sus límites, o acordar un alcance sencillo explícito. Las cifras históricas sólo sirven si copia, filtros y contrato son equivalentes.

Disponer de permiso de generación, revisión, materialización, consulta y exportación; aceptar el envío al destino mostrado cuando el flujo lo solicite.

Tener un expediente histórico comprobado como respaldo, identificado como consulta y no como salida de la generación actual.

\Needspace{8\baselineskip}
\begingroup
\setbox0=\vbox{\hsize=\linewidth
{\normalfont\normalsize\sffamily\bfseries\normalsize\color{Teal}\hyphenpenalty=10000\exhyphenpenalty=10000 Tiempo previsto y alternativas\par}

\par
\begingroup
\fontsize{10}{12.5}\selectfont
\setlength{\tabcolsep}{5pt}
\renewcommand{\arraystretch}{1.15}
\begin{tabular}{@{}>{\RaggedRight\arraybackslash}p{0.27\linewidth}>{\RaggedRight\arraybackslash}p{0.66\linewidth}@{}}
\toprule
\textbf{Condición} & \textbf{Cómo organizarse} \\
\midrule
Tiempo total & 25 minutos de práctica guiada; las demoras de proveedor o carga pueden obligar a continuar en otra sesión. \\
\bottomrule
\end{tabular}
\endgroup
}
\Needspace{\dimexpr\ht0+\dp0+48pt\relax}
\endgroup
\subsection{Tiempo previsto y alternativas}
\begingroup
\fontsize{10}{12.5}\selectfont
\setlength{\tabcolsep}{5pt}
\renewcommand{\arraystretch}{1.15}
\begin{longtable}{@{}>{\RaggedRight\arraybackslash}p{0.27\linewidth}>{\RaggedRight\arraybackslash}p{0.66\linewidth}@{}}
\toprule
\textbf{Condición} & \textbf{Cómo organizarse} \\
\midrule
\endfirsthead
\toprule
\textbf{Condición} & \textbf{Cómo organizarse} \\
\midrule
\endhead
\midrule
\multicolumn{2}{r}{\footnotesize Continúa en la página siguiente}\\
\endfoot
\bottomrule
\endlastfoot
Tiempo total & 25 minutos de práctica guiada; las demoras de proveedor o carga pueden obligar a continuar en otra sesión. \\
\addlinespace[3pt]
Generaciones nuevas planificadas & 1 \\
\addlinespace[3pt]
Reintentos manuales & Como máximo uno después de un error definitivo y si hay tiempo; nunca reenviar una operación pendiente. \\
\addlinespace[3pt]
Regla de corte & Al agotar el bloque, anotar estado, conservar evidencia y consultar el historial antes de decidir un nuevo intento. El reloj de la exposición no cambia el timeout del producto. \\
\addlinespace[3pt]
Segunda fuente & Repetir el guion en otra sesión con su propio catálogo y metadatos. Las referencias de ambas fuentes se incluyen para preparar esa repetición. \\
\addlinespace[3pt]
\end{longtable}
\endgroup

\begingroup
\setbox0=\vbox{\hsize=\dimexpr\linewidth-1.7em\relax
\everypar{}\csname @newlistfalse\endcsname
\textbf{Minutos 0-2.} Explique evento y grano. En el explorador compruebe fuente, instantánea, fecha de negocio y una relación relevante.

{\small\textbf{Tiempo máximo previsto de este bloque:} 2 minutos.\par}
\textbf{Qué debe observar:} El caso distingue línea de pedido o factura y conoce su procedencia.

\textbf{Qué registrar:} Fuente, instantánea y relación elegida; no es necesario escribir SQL.

\textbf{Si no se cumple:} Si la estructura no respalda una capacidad imprescindible, aclare el objetivo o detenga ese caso.

\par
}
\dimen0=\dimexpr\ht0+\dp0+4\baselineskip\relax
\ifdim\dimen0<\textheight\Needspace{\dimen0}\fi
\endgroup
\subsection{Guion paso a paso}
\begin{enumerate}
\begingroup
\setbox0=\vbox{\hsize=\linewidth
\everypar{}\csname @newlistfalse\endcsname
\textbf{Minutos 0-2.} Explique evento y grano. En el explorador compruebe fuente, instantánea, fecha de negocio y una relación relevante.

{\small\textbf{Tiempo máximo previsto de este bloque:} 2 minutos.\par}
\textbf{Qué debe observar:} El caso distingue línea de pedido o factura y conoce su procedencia.

\textbf{Qué registrar:} Fuente, instantánea y relación elegida; no es necesario escribir SQL.

\textbf{Si no se cumple:} Si la estructura no respalda una capacidad imprescindible, aclare el objetivo o detenga ese caso.

\par
}
\dimen0=\dimexpr\ht0+\dp0+2\baselineskip\relax
\ifdim\dimen0<\textheight
\Needspace{\dimen0}
\endgroup
\item \begin{samepage}
\textbf{Minutos 0-2.} Explique evento y grano. En el explorador compruebe fuente, instantánea, fecha de negocio y una relación relevante.

{\small\textbf{Tiempo máximo previsto de este bloque:} 2 minutos.\par}
\textbf{Qué debe observar:} El caso distingue línea de pedido o factura y conoce su procedencia.

\textbf{Qué registrar:} Fuente, instantánea y relación elegida; no es necesario escribir SQL.

\textbf{Si no se cumple:} Si la estructura no respalda una capacidad imprescindible, aclare el objetivo o detenga ese caso.

\end{samepage}
\else
\endgroup
\item \textbf{Minutos 0-2.} Explique evento y grano. En el explorador compruebe fuente, instantánea, fecha de negocio y una relación relevante.

{\small\textbf{Tiempo máximo previsto de este bloque:} 2 minutos.\par}
\textbf{Qué debe observar:} El caso distingue línea de pedido o factura y conoce su procedencia.

\textbf{Qué registrar:} Fuente, instantánea y relación elegida; no es necesario escribir SQL.

\textbf{Si no se cumple:} Si la estructura no respalda una capacidad imprescindible, aclare el objetivo o detenga ese caso.

\fi
\begingroup
\setbox0=\vbox{\hsize=\linewidth
\everypar{}\csname @newlistfalse\endcsname
\textbf{Minutos 2-4.} Revise el catálogo propio y escriba la necesidad. Seleccione sólo las preguntas pertinentes y Mensual.

{\small\textbf{Tiempo máximo previsto de este bloque:} 2 minutos.\par}
\textbf{Qué debe observar:} Texto, preguntas y período representan el mismo propósito; no se restauran catálogos sólo para igualar una captura.

\textbf{Qué registrar:} Objetivo definitivo y opciones seleccionadas.

\textbf{Si no se cumple:} Si falta una orientación, revise con una persona autorizada. No copie relaciones de la otra base por parecido de nombres.

\par
}
\dimen0=\dimexpr\ht0+\dp0+2\baselineskip\relax
\ifdim\dimen0<\textheight
\Needspace{\dimen0}
\endgroup
\item \begin{samepage}
\textbf{Minutos 2-4.} Revise el catálogo propio y escriba la necesidad. Seleccione sólo las preguntas pertinentes y Mensual.

{\small\textbf{Tiempo máximo previsto de este bloque:} 2 minutos.\par}
\textbf{Qué debe observar:} Texto, preguntas y período representan el mismo propósito; no se restauran catálogos sólo para igualar una captura.

\textbf{Qué registrar:} Objetivo definitivo y opciones seleccionadas.

\textbf{Si no se cumple:} Si falta una orientación, revise con una persona autorizada. No copie relaciones de la otra base por parecido de nombres.

\end{samepage}
\else
\endgroup
\item \textbf{Minutos 2-4.} Revise el catálogo propio y escriba la necesidad. Seleccione sólo las preguntas pertinentes y Mensual.

{\small\textbf{Tiempo máximo previsto de este bloque:} 2 minutos.\par}
\textbf{Qué debe observar:} Texto, preguntas y período representan el mismo propósito; no se restauran catálogos sólo para igualar una captura.

\textbf{Qué registrar:} Objetivo definitivo y opciones seleccionadas.

\textbf{Si no se cumple:} Si falta una orientación, revise con una persona autorizada. No copie relaciones de la otra base por parecido de nombres.

\fi
\begingroup
\setbox0=\vbox{\hsize=\linewidth
\everypar{}\csname @newlistfalse\endcsname
\textbf{Minutos 4-8.} Autorice el destino mostrado y revise la viabilidad. Abra la evidencia de requisitos directos o derivables; aclare las ambigüedades importantes.

{\small\textbf{Tiempo máximo previsto de este bloque:} 4 minutos.\par}
\textbf{Qué debe observar:} Existe una revisión vigente del texto y contexto definitivo. Aceptar un límite no fabrica una capacidad.

\textbf{Qué registrar:} Estado de viabilidad, límites leídos y decisión humana.

\textbf{Si no se cumple:} Si la IA cambia el propósito, conserve el texto original. Ante falla use el protocolo; no continúe como si la revisión hubiera terminado.

\par
}
\dimen0=\dimexpr\ht0+\dp0+2\baselineskip\relax
\ifdim\dimen0<\textheight
\Needspace{\dimen0}
\endgroup
\item \begin{samepage}
\textbf{Minutos 4-8.} Autorice el destino mostrado y revise la viabilidad. Abra la evidencia de requisitos directos o derivables; aclare las ambigüedades importantes.

{\small\textbf{Tiempo máximo previsto de este bloque:} 4 minutos.\par}
\textbf{Qué debe observar:} Existe una revisión vigente del texto y contexto definitivo. Aceptar un límite no fabrica una capacidad.

\textbf{Qué registrar:} Estado de viabilidad, límites leídos y decisión humana.

\textbf{Si no se cumple:} Si la IA cambia el propósito, conserve el texto original. Ante falla use el protocolo; no continúe como si la revisión hubiera terminado.

\end{samepage}
\else
\endgroup
\item \textbf{Minutos 4-8.} Autorice el destino mostrado y revise la viabilidad. Abra la evidencia de requisitos directos o derivables; aclare las ambigüedades importantes.

{\small\textbf{Tiempo máximo previsto de este bloque:} 4 minutos.\par}
\textbf{Qué debe observar:} Existe una revisión vigente del texto y contexto definitivo. Aceptar un límite no fabrica una capacidad.

\textbf{Qué registrar:} Estado de viabilidad, límites leídos y decisión humana.

\textbf{Si no se cumple:} Si la IA cambia el propósito, conserve el texto original. Ante falla use el protocolo; no continúe como si la revisión hubiera terminado.

\fi
\begingroup
\setbox0=\vbox{\hsize=\linewidth
\everypar{}\csname @newlistfalse\endcsname
\textbf{Minutos 8-12.} Genere conceptos y propuesta. Revise las cinco etapas, los nombres descriptivos, las medidas y las rutas que preservan el grano.

{\small\textbf{Tiempo máximo previsto de este bloque:} 4 minutos.\par}
\textbf{Qué debe observar:} Contrato revisable que conserva el objetivo. En WWI se distingue comprador de destinatario de facturación y entrega de dirección postal.

\textbf{Qué registrar:} Identificador propio de la propuesta, proveedor/modelo y aspectos comprobados.

\textbf{Si no se cumple:} Si no aparece un contrato válido dentro del bloque, marque generación pendiente o bloqueada y pase a consultar un histórico anunciado. No lo apruebe para ahorrar tiempo.

\par
}
\dimen0=\dimexpr\ht0+\dp0+2\baselineskip\relax
\ifdim\dimen0<\textheight
\Needspace{\dimen0}
\endgroup
\item \begin{samepage}
\textbf{Minutos 8-12.} Genere conceptos y propuesta. Revise las cinco etapas, los nombres descriptivos, las medidas y las rutas que preservan el grano.

{\small\textbf{Tiempo máximo previsto de este bloque:} 4 minutos.\par}
\textbf{Qué debe observar:} Contrato revisable que conserva el objetivo. En WWI se distingue comprador de destinatario de facturación y entrega de dirección postal.

\textbf{Qué registrar:} Identificador propio de la propuesta, proveedor/modelo y aspectos comprobados.

\textbf{Si no se cumple:} Si no aparece un contrato válido dentro del bloque, marque generación pendiente o bloqueada y pase a consultar un histórico anunciado. No lo apruebe para ahorrar tiempo.

\end{samepage}
\else
\endgroup
\item \textbf{Minutos 8-12.} Genere conceptos y propuesta. Revise las cinco etapas, los nombres descriptivos, las medidas y las rutas que preservan el grano.

{\small\textbf{Tiempo máximo previsto de este bloque:} 4 minutos.\par}
\textbf{Qué debe observar:} Contrato revisable que conserva el objetivo. En WWI se distingue comprador de destinatario de facturación y entrega de dirección postal.

\textbf{Qué registrar:} Identificador propio de la propuesta, proveedor/modelo y aspectos comprobados.

\textbf{Si no se cumple:} Si no aparece un contrato válido dentro del bloque, marque generación pendiente o bloqueada y pase a consultar un histórico anunciado. No lo apruebe para ahorrar tiempo.

\fi
\begingroup
\setbox0=\vbox{\hsize=\linewidth
\everypar{}\csname @newlistfalse\endcsname
\textbf{Minutos 12-15.} Personalice sólo decisiones permitidas y justifique el cambio si hace falta. Apruebe únicamente cuando no haya errores bloqueantes y comprenda las advertencias.

{\small\textbf{Tiempo máximo previsto de este bloque:} 3 minutos.\par}
\textbf{Qué debe observar:} La aprobación conserva una versión; una personalización crea otra ligada a su origen sin sobrescribirla.

\textbf{Qué registrar:} Propuesta aprobada o motivo real de no aprobación; justificar las advertencias aceptadas.

\textbf{Si no se cumple:} No marque advertencias a ciegas. Si el significado no se puede justificar, mantenga la propuesta sin aprobar.

\par
}
\dimen0=\dimexpr\ht0+\dp0+2\baselineskip\relax
\ifdim\dimen0<\textheight
\Needspace{\dimen0}
\endgroup
\item \begin{samepage}
\textbf{Minutos 12-15.} Personalice sólo decisiones permitidas y justifique el cambio si hace falta. Apruebe únicamente cuando no haya errores bloqueantes y comprenda las advertencias.

{\small\textbf{Tiempo máximo previsto de este bloque:} 3 minutos.\par}
\textbf{Qué debe observar:} La aprobación conserva una versión; una personalización crea otra ligada a su origen sin sobrescribirla.

\textbf{Qué registrar:} Propuesta aprobada o motivo real de no aprobación; justificar las advertencias aceptadas.

\textbf{Si no se cumple:} No marque advertencias a ciegas. Si el significado no se puede justificar, mantenga la propuesta sin aprobar.

\end{samepage}
\else
\endgroup
\item \textbf{Minutos 12-15.} Personalice sólo decisiones permitidas y justifique el cambio si hace falta. Apruebe únicamente cuando no haya errores bloqueantes y comprenda las advertencias.

{\small\textbf{Tiempo máximo previsto de este bloque:} 3 minutos.\par}
\textbf{Qué debe observar:} La aprobación conserva una versión; una personalización crea otra ligada a su origen sin sobrescribirla.

\textbf{Qué registrar:} Propuesta aprobada o motivo real de no aprobación; justificar las advertencias aceptadas.

\textbf{Si no se cumple:} No marque advertencias a ciegas. Si el significado no se puede justificar, mantenga la propuesta sin aprobar.

\fi
\begingroup
\setbox0=\vbox{\hsize=\linewidth
\everypar{}\csname @newlistfalse\endcsname
\textbf{Minutos 15-19.} En Datamart de ventas elija esa propuesta, revise indicadores y denominadores, confirme el plan e inicie una sola carga.

{\small\textbf{Tiempo máximo previsto de este bloque:} 4 minutos.\par}
\textbf{Qué debe observar:} Se prepara y ejecuta un expediente propio. Generación, aprobación e inicio de carga son decisiones distintas.

\textbf{Qué registrar:} Identificador de ejecución, estado y selección; registrar hora sin reenviar por una respuesta lenta.

\textbf{Si no se cumple:} Si la pantalla pierde respuesta, consulte el expediente antes de repetir. Si sigue en curso al terminar el bloque, no anuncie carga completada.

\par
}
\dimen0=\dimexpr\ht0+\dp0+2\baselineskip\relax
\ifdim\dimen0<\textheight
\Needspace{\dimen0}
\endgroup
\item \begin{samepage}
\textbf{Minutos 15-19.} En Datamart de ventas elija esa propuesta, revise indicadores y denominadores, confirme el plan e inicie una sola carga.

{\small\textbf{Tiempo máximo previsto de este bloque:} 4 minutos.\par}
\textbf{Qué debe observar:} Se prepara y ejecuta un expediente propio. Generación, aprobación e inicio de carga son decisiones distintas.

\textbf{Qué registrar:} Identificador de ejecución, estado y selección; registrar hora sin reenviar por una respuesta lenta.

\textbf{Si no se cumple:} Si la pantalla pierde respuesta, consulte el expediente antes de repetir. Si sigue en curso al terminar el bloque, no anuncie carga completada.

\end{samepage}
\else
\endgroup
\item \textbf{Minutos 15-19.} En Datamart de ventas elija esa propuesta, revise indicadores y denominadores, confirme el plan e inicie una sola carga.

{\small\textbf{Tiempo máximo previsto de este bloque:} 4 minutos.\par}
\textbf{Qué debe observar:} Se prepara y ejecuta un expediente propio. Generación, aprobación e inicio de carga son decisiones distintas.

\textbf{Qué registrar:} Identificador de ejecución, estado y selección; registrar hora sin reenviar por una respuesta lenta.

\textbf{Si no se cumple:} Si la pantalla pierde respuesta, consulte el expediente antes de repetir. Si sigue en curso al terminar el bloque, no anuncie carga completada.

\fi
\begingroup
\setbox0=\vbox{\hsize=\linewidth
\everypar{}\csname @newlistfalse\endcsname
\textbf{Minutos 19-22.} Revise conciliación, cobertura y estado de publicación. Abra el panel con esa ejecución y aplique el filtro de contraste.

{\small\textbf{Tiempo máximo previsto de este bloque:} 3 minutos.\par}
\textbf{Qué debe observar:} Filas y cálculos se contrastan con el mismo universo; no basta que se vea un gráfico. Moneda y fecha tienen fundamento.

\textbf{Qué registrar:} Resultado de conciliación y al menos una suma, conteo distinto y razón si existen en la selección.

\textbf{Si no se cumple:} Si hay diferencias sin explicar, el resultado no se declara aprobado. Un filtro nuevo se compara con su propio denominador.

\par
}
\dimen0=\dimexpr\ht0+\dp0+2\baselineskip\relax
\ifdim\dimen0<\textheight
\Needspace{\dimen0}
\endgroup
\item \begin{samepage}
\textbf{Minutos 19-22.} Revise conciliación, cobertura y estado de publicación. Abra el panel con esa ejecución y aplique el filtro de contraste.

{\small\textbf{Tiempo máximo previsto de este bloque:} 3 minutos.\par}
\textbf{Qué debe observar:} Filas y cálculos se contrastan con el mismo universo; no basta que se vea un gráfico. Moneda y fecha tienen fundamento.

\textbf{Qué registrar:} Resultado de conciliación y al menos una suma, conteo distinto y razón si existen en la selección.

\textbf{Si no se cumple:} Si hay diferencias sin explicar, el resultado no se declara aprobado. Un filtro nuevo se compara con su propio denominador.

\end{samepage}
\else
\endgroup
\item \textbf{Minutos 19-22.} Revise conciliación, cobertura y estado de publicación. Abra el panel con esa ejecución y aplique el filtro de contraste.

{\small\textbf{Tiempo máximo previsto de este bloque:} 3 minutos.\par}
\textbf{Qué debe observar:} Filas y cálculos se contrastan con el mismo universo; no basta que se vea un gráfico. Moneda y fecha tienen fundamento.

\textbf{Qué registrar:} Resultado de conciliación y al menos una suma, conteo distinto y razón si existen en la selección.

\textbf{Si no se cumple:} Si hay diferencias sin explicar, el resultado no se declara aprobado. Un filtro nuevo se compara con su propio denominador.

\fi
\begingroup
\setbox0=\vbox{\hsize=\linewidth
\everypar{}\csname @newlistfalse\endcsname
\textbf{Minutos 22-24.} Pida una explicación breve al copiloto y, si el tiempo lo permite, exporte y abra el archivo con el mismo alcance.

{\small\textbf{Tiempo máximo previsto de este bloque:} 2 minutos.\par}
\textbf{Qué debe observar:} Se revisa la interpretación y una cifra; sólo se marca lo que realmente se obtuvo y abrió.

\textbf{Qué registrar:} Respuesta o error del chat; archivo y contexto comprobados, o paso no ejecutado.

\textbf{Si no se cumple:} Cuando no alcance el tiempo, registre el paso pendiente para otra práctica; no use un archivo histórico como si fuera una descarga nueva.

\par
}
\dimen0=\dimexpr\ht0+\dp0+2\baselineskip\relax
\ifdim\dimen0<\textheight
\Needspace{\dimen0}
\endgroup
\item \begin{samepage}
\textbf{Minutos 22-24.} Pida una explicación breve al copiloto y, si el tiempo lo permite, exporte y abra el archivo con el mismo alcance.

{\small\textbf{Tiempo máximo previsto de este bloque:} 2 minutos.\par}
\textbf{Qué debe observar:} Se revisa la interpretación y una cifra; sólo se marca lo que realmente se obtuvo y abrió.

\textbf{Qué registrar:} Respuesta o error del chat; archivo y contexto comprobados, o paso no ejecutado.

\textbf{Si no se cumple:} Cuando no alcance el tiempo, registre el paso pendiente para otra práctica; no use un archivo histórico como si fuera una descarga nueva.

\end{samepage}
\else
\endgroup
\item \textbf{Minutos 22-24.} Pida una explicación breve al copiloto y, si el tiempo lo permite, exporte y abra el archivo con el mismo alcance.

{\small\textbf{Tiempo máximo previsto de este bloque:} 2 minutos.\par}
\textbf{Qué debe observar:} Se revisa la interpretación y una cifra; sólo se marca lo que realmente se obtuvo y abrió.

\textbf{Qué registrar:} Respuesta o error del chat; archivo y contexto comprobados, o paso no ejecutado.

\textbf{Si no se cumple:} Cuando no alcance el tiempo, registre el paso pendiente para otra práctica; no use un archivo histórico como si fuera una descarga nueva.

\fi
\begingroup
\setbox0=\vbox{\hsize=\linewidth
\everypar{}\csname @newlistfalse\endcsname
\textbf{Minutos 24-25.} Explique quién propuso, quién validó, quién aprobó y qué cálculo respalda la conclusión. Complete el registro.

{\small\textbf{Tiempo máximo previsto de este bloque:} 1 minuto.\par}
\textbf{Qué debe observar:} La persona puede defender el caso con evidencia y sus límites.

\textbf{Qué registrar:} Conclusión, estado de cada etapa y siguiente paso.

\textbf{Si no se cumple:} Si se usó la contingencia histórica, separar claramente la ejecución nueva incompleta del resultado conservado.

\par
}
\dimen0=\dimexpr\ht0+\dp0+2\baselineskip\relax
\ifdim\dimen0<\textheight
\Needspace{\dimen0}
\endgroup
\item \begin{samepage}
\textbf{Minutos 24-25.} Explique quién propuso, quién validó, quién aprobó y qué cálculo respalda la conclusión. Complete el registro.

{\small\textbf{Tiempo máximo previsto de este bloque:} 1 minuto.\par}
\textbf{Qué debe observar:} La persona puede defender el caso con evidencia y sus límites.

\textbf{Qué registrar:} Conclusión, estado de cada etapa y siguiente paso.

\textbf{Si no se cumple:} Si se usó la contingencia histórica, separar claramente la ejecución nueva incompleta del resultado conservado.

\end{samepage}
\else
\endgroup
\item \textbf{Minutos 24-25.} Explique quién propuso, quién validó, quién aprobó y qué cálculo respalda la conclusión. Complete el registro.

{\small\textbf{Tiempo máximo previsto de este bloque:} 1 minuto.\par}
\textbf{Qué debe observar:} La persona puede defender el caso con evidencia y sus límites.

\textbf{Qué registrar:} Conclusión, estado de cada etapa y siguiente paso.

\textbf{Si no se cumple:} Si se usó la contingencia histórica, separar claramente la ejecución nueva incompleta del resultado conservado.

\fi
\end{enumerate}
\subsection{Antes de dar por terminado el ejercicio}
\begin{itemize}
\item Objetivo, evento, grano y fuente conservados.
\item Aprobación y carga identificadas de forma separada.
\item Conciliación y contexto numérico revisados antes de concluir.
\item Ninguna propuesta inválida fue aprobada para cumplir el tiempo.
\item Las etapas que no terminaron se registran y no se convierten en éxitos por la existencia de un histórico.
\end{itemize}

\section{Sesenta minutos: empezar en una aplicación limpia y aprender el recorrido}
\textbf{Al terminar podrá:} Configurar desde la web un entorno de práctica ya instalado y realizar dos ejercicios básicos, uno por fuente, registrando el resultado real de cada etapa.

Guion de aprendizaje propuesto; no ejecutado durante esta redacción.

Configuración funcional desde la web y dos ejecuciones básicas nuevas condicionadas a sus controles. Sin instalación, parada, limpieza ni migración de infraestructura durante la práctica.

La infraestructura ya fue levantada por el responsable: aplicación accesible, contenedores saludables, migraciones y administrador inicial disponibles y las dos bases demostrativas restauradas. No se instalan ni se eliminan contenedores en este guion.

Usar una instancia de práctica limpia dispuesta por el responsable. No borrar una instalación con trabajo anterior para simular limpieza.

Tener acceso administrativo autorizado y, por un canal privado, credenciales de lectura de ambas fuentes y una credencial de IA propia autorizada cuando corresponda. No escribirlas en el manual ni en la hoja de práctica.

Acordar un presupuesto de consumo y tiempo. Las dos rutas básicas son ejercicios propuestos; su repetición no se presenta como ya realizada por este manual.

Conservar las necesidades de referencia y la hoja de registro. Los IDs de las nuevas propuestas y ejecuciones serán los que genere esta instalación.

\Needspace{8\baselineskip}
\begingroup
\setbox0=\vbox{\hsize=\linewidth
{\normalfont\normalsize\sffamily\bfseries\normalsize\color{Teal}\hyphenpenalty=10000\exhyphenpenalty=10000 Tiempo previsto y alternativas\par}

\par
\begingroup
\fontsize{10}{12.5}\selectfont
\setlength{\tabcolsep}{5pt}
\renewcommand{\arraystretch}{1.15}
\begin{tabular}{@{}>{\RaggedRight\arraybackslash}p{0.27\linewidth}>{\RaggedRight\arraybackslash}p{0.66\linewidth}@{}}
\toprule
\textbf{Condición} & \textbf{Cómo organizarse} \\
\midrule
Tiempo total & 60 minutos de aprendizaje. Es una agenda orientativa, no una garantía de duración del proveedor o de las cargas. \\
\bottomrule
\end{tabular}
\endgroup
}
\Needspace{\dimexpr\ht0+\dp0+48pt\relax}
\endgroup
\subsection{Tiempo previsto y alternativas}
\begingroup
\fontsize{10}{12.5}\selectfont
\setlength{\tabcolsep}{5pt}
\renewcommand{\arraystretch}{1.15}
\begin{longtable}{@{}>{\RaggedRight\arraybackslash}p{0.27\linewidth}>{\RaggedRight\arraybackslash}p{0.66\linewidth}@{}}
\toprule
\textbf{Condición} & \textbf{Cómo organizarse} \\
\midrule
\endfirsthead
\toprule
\textbf{Condición} & \textbf{Cómo organizarse} \\
\midrule
\endhead
\midrule
\multicolumn{2}{r}{\footnotesize Continúa en la página siguiente}\\
\endfoot
\bottomrule
\endlastfoot
Tiempo total & 60 minutos de aprendizaje. Es una agenda orientativa, no una garantía de duración del proveedor o de las cargas. \\
\addlinespace[3pt]
Dos fuentes & Doce minutos de trabajo guiado por fuente; ampliar a otra sesión cuando una fase siga pendiente. \\
\addlinespace[3pt]
Reintentos manuales & Máximo uno tras un error definitivo dentro del bloque; no repetir una carga en curso ni consumir cuota para forzar un caso. \\
\addlinespace[3pt]
Si falta ia & Completar preparación, catálogo, esquema, roles y explicación. Registrar las generaciones como no ejecutadas y retomarlas cuando haya un destino autorizado utilizable. \\
\addlinespace[3pt]
\end{longtable}
\endgroup

\begingroup
\setbox0=\vbox{\hsize=\dimexpr\linewidth-1.7em\relax
\everypar{}\csname @newlistfalse\endcsname
\textbf{Minutos 0-4.} Abra Inicio e identifique lo que falta. Explique navegador, servidor, SQL Server de origen y PostgreSQL de configuración y resultados.

{\small\textbf{Tiempo máximo previsto de este bloque:} 4 minutos.\par}
\textbf{Qué debe observar:} Se reconoce una aplicación preparada para configurar, no un panel de ventas ya cargado.

\textbf{Qué registrar:} Estado inicial, entorno y rol real. Explicar que las 83 tablas del inventario histórico no son un requisito del portable limpio.

\textbf{Si no se cumple:} Si no inicia sesión o la aplicación no está disponible, registrar bloqueo y volver con el responsable de infraestructura; no gastar la sesión inventando resultados.

\par
}
\dimen0=\dimexpr\ht0+\dp0+4\baselineskip\relax
\ifdim\dimen0<\textheight\Needspace{\dimen0}\fi
\endgroup
\subsection{Guion paso a paso}
\begin{enumerate}
\begingroup
\setbox0=\vbox{\hsize=\linewidth
\everypar{}\csname @newlistfalse\endcsname
\textbf{Minutos 0-4.} Abra Inicio e identifique lo que falta. Explique navegador, servidor, SQL Server de origen y PostgreSQL de configuración y resultados.

{\small\textbf{Tiempo máximo previsto de este bloque:} 4 minutos.\par}
\textbf{Qué debe observar:} Se reconoce una aplicación preparada para configurar, no un panel de ventas ya cargado.

\textbf{Qué registrar:} Estado inicial, entorno y rol real. Explicar que las 83 tablas del inventario histórico no son un requisito del portable limpio.

\textbf{Si no se cumple:} Si no inicia sesión o la aplicación no está disponible, registrar bloqueo y volver con el responsable de infraestructura; no gastar la sesión inventando resultados.

\par
}
\dimen0=\dimexpr\ht0+\dp0+2\baselineskip\relax
\ifdim\dimen0<\textheight
\Needspace{\dimen0}
\endgroup
\item \begin{samepage}
\textbf{Minutos 0-4.} Abra Inicio e identifique lo que falta. Explique navegador, servidor, SQL Server de origen y PostgreSQL de configuración y resultados.

{\small\textbf{Tiempo máximo previsto de este bloque:} 4 minutos.\par}
\textbf{Qué debe observar:} Se reconoce una aplicación preparada para configurar, no un panel de ventas ya cargado.

\textbf{Qué registrar:} Estado inicial, entorno y rol real. Explicar que las 83 tablas del inventario histórico no son un requisito del portable limpio.

\textbf{Si no se cumple:} Si no inicia sesión o la aplicación no está disponible, registrar bloqueo y volver con el responsable de infraestructura; no gastar la sesión inventando resultados.

\end{samepage}
\else
\endgroup
\item \textbf{Minutos 0-4.} Abra Inicio e identifique lo que falta. Explique navegador, servidor, SQL Server de origen y PostgreSQL de configuración y resultados.

{\small\textbf{Tiempo máximo previsto de este bloque:} 4 minutos.\par}
\textbf{Qué debe observar:} Se reconoce una aplicación preparada para configurar, no un panel de ventas ya cargado.

\textbf{Qué registrar:} Estado inicial, entorno y rol real. Explicar que las 83 tablas del inventario histórico no son un requisito del portable limpio.

\textbf{Si no se cumple:} Si no inicia sesión o la aplicación no está disponible, registrar bloqueo y volver con el responsable de infraestructura; no gastar la sesión inventando resultados.

\fi
\begingroup
\setbox0=\vbox{\hsize=\linewidth
\everypar{}\csname @newlistfalse\endcsname
\textbf{Minutos 4-12.} Registre, pruebe y habilite cada fuente con sus datos privados. Capture metadatos de AW y WWI y observe que cada explorador cambia con el selector.

{\small\textbf{Tiempo máximo previsto de este bloque:} 8 minutos.\par}
\textbf{Qué debe observar:} Dos conexiones de solo lectura identificadas y sus instantáneas propias; obtener metadatos no extrae ventas.

\textbf{Qué registrar:} Nombre de base, conector y resultado de prueba; no guardar contraseña ni nombre técnico de usuario en una captura.

\textbf{Si no se cumple:} Si una prueba no confirma lectura autorizada, no habilitar por omisión ni usar una cuenta administrativa de SQL para saltar el control.

\par
}
\dimen0=\dimexpr\ht0+\dp0+2\baselineskip\relax
\ifdim\dimen0<\textheight
\Needspace{\dimen0}
\endgroup
\item \begin{samepage}
\textbf{Minutos 4-12.} Registre, pruebe y habilite cada fuente con sus datos privados. Capture metadatos de AW y WWI y observe que cada explorador cambia con el selector.

{\small\textbf{Tiempo máximo previsto de este bloque:} 8 minutos.\par}
\textbf{Qué debe observar:} Dos conexiones de solo lectura identificadas y sus instantáneas propias; obtener metadatos no extrae ventas.

\textbf{Qué registrar:} Nombre de base, conector y resultado de prueba; no guardar contraseña ni nombre técnico de usuario en una captura.

\textbf{Si no se cumple:} Si una prueba no confirma lectura autorizada, no habilitar por omisión ni usar una cuenta administrativa de SQL para saltar el control.

\end{samepage}
\else
\endgroup
\item \textbf{Minutos 4-12.} Registre, pruebe y habilite cada fuente con sus datos privados. Capture metadatos de AW y WWI y observe que cada explorador cambia con el selector.

{\small\textbf{Tiempo máximo previsto de este bloque:} 8 minutos.\par}
\textbf{Qué debe observar:} Dos conexiones de solo lectura identificadas y sus instantáneas propias; obtener metadatos no extrae ventas.

\textbf{Qué registrar:} Nombre de base, conector y resultado de prueba; no guardar contraseña ni nombre técnico de usuario en una captura.

\textbf{Si no se cumple:} Si una prueba no confirma lectura autorizada, no habilitar por omisión ni usar una cuenta administrativa de SQL para saltar el control.

\fi
\begingroup
\setbox0=\vbox{\hsize=\linewidth
\everypar{}\csname @newlistfalse\endcsname
\textbf{Minutos 12-18.} Configure el proveedor y modelo acordados, registre la credencial de manera privada si hace falta y ejecute su prueba.

{\small\textbf{Tiempo máximo previsto de este bloque:} 6 minutos.\par}
\textbf{Qué debe observar:} Proveedor probado con credencial configurada, sin mostrar su valor. Configurar un proveedor no valida todavía cada propuesta que produzca.

\textbf{Qué registrar:} Proveedor/modelo y resultado de prueba sin secretos; registro de consumo sólo si se informa realmente.

\textbf{Si no se cumple:} Si no hay cuota o conectividad, registrar la limitación y continuar con esquema, catálogo y explicación del flujo. Una práctica sin proveedor utilizable no se declara ciclo BI completo.

\par
}
\dimen0=\dimexpr\ht0+\dp0+2\baselineskip\relax
\ifdim\dimen0<\textheight
\Needspace{\dimen0}
\endgroup
\item \begin{samepage}
\textbf{Minutos 12-18.} Configure el proveedor y modelo acordados, registre la credencial de manera privada si hace falta y ejecute su prueba.

{\small\textbf{Tiempo máximo previsto de este bloque:} 6 minutos.\par}
\textbf{Qué debe observar:} Proveedor probado con credencial configurada, sin mostrar su valor. Configurar un proveedor no valida todavía cada propuesta que produzca.

\textbf{Qué registrar:} Proveedor/modelo y resultado de prueba sin secretos; registro de consumo sólo si se informa realmente.

\textbf{Si no se cumple:} Si no hay cuota o conectividad, registrar la limitación y continuar con esquema, catálogo y explicación del flujo. Una práctica sin proveedor utilizable no se declara ciclo BI completo.

\end{samepage}
\else
\endgroup
\item \textbf{Minutos 12-18.} Configure el proveedor y modelo acordados, registre la credencial de manera privada si hace falta y ejecute su prueba.

{\small\textbf{Tiempo máximo previsto de este bloque:} 6 minutos.\par}
\textbf{Qué debe observar:} Proveedor probado con credencial configurada, sin mostrar su valor. Configurar un proveedor no valida todavía cada propuesta que produzca.

\textbf{Qué registrar:} Proveedor/modelo y resultado de prueba sin secretos; registro de consumo sólo si se informa realmente.

\textbf{Si no se cumple:} Si no hay cuota o conectividad, registrar la limitación y continuar con esquema, catálogo y explicación del flujo. Una práctica sin proveedor utilizable no se declara ciclo BI completo.

\fi
\begingroup
\setbox0=\vbox{\hsize=\linewidth
\everypar{}\csname @newlistfalse\endcsname
\textbf{Minutos 18-24.} Abra el catálogo de cada fuente. Lea las cuatro orientaciones base y elija tiempo y producto para las rutas básicas. Revise las periodicidades.

{\small\textbf{Tiempo máximo previsto de este bloque:} 6 minutos.\par}
\textbf{Qué debe observar:} Preguntas y períodos pertenecen a cada conexión. El catálogo no fija por adelantado el hecho y las dimensiones ni se genera automáticamente con una llamada de IA.

\textbf{Qué registrar:} Etiquetas e instrucciones utilizadas y Mensual habilitado por fuente.

\textbf{Si no se cumple:} No pulsar Restaurar catálogo si se quieren conservar cambios. No copiar las preguntas de costos a otra base sin comprobar su sentido.

\par
}
\dimen0=\dimexpr\ht0+\dp0+2\baselineskip\relax
\ifdim\dimen0<\textheight
\Needspace{\dimen0}
\endgroup
\item \begin{samepage}
\textbf{Minutos 18-24.} Abra el catálogo de cada fuente. Lea las cuatro orientaciones base y elija tiempo y producto para las rutas básicas. Revise las periodicidades.

{\small\textbf{Tiempo máximo previsto de este bloque:} 6 minutos.\par}
\textbf{Qué debe observar:} Preguntas y períodos pertenecen a cada conexión. El catálogo no fija por adelantado el hecho y las dimensiones ni se genera automáticamente con una llamada de IA.

\textbf{Qué registrar:} Etiquetas e instrucciones utilizadas y Mensual habilitado por fuente.

\textbf{Si no se cumple:} No pulsar Restaurar catálogo si se quieren conservar cambios. No copiar las preguntas de costos a otra base sin comprobar su sentido.

\end{samepage}
\else
\endgroup
\item \textbf{Minutos 18-24.} Abra el catálogo de cada fuente. Lea las cuatro orientaciones base y elija tiempo y producto para las rutas básicas. Revise las periodicidades.

{\small\textbf{Tiempo máximo previsto de este bloque:} 6 minutos.\par}
\textbf{Qué debe observar:} Preguntas y períodos pertenecen a cada conexión. El catálogo no fija por adelantado el hecho y las dimensiones ni se genera automáticamente con una llamada de IA.

\textbf{Qué registrar:} Etiquetas e instrucciones utilizadas y Mensual habilitado por fuente.

\textbf{Si no se cumple:} No pulsar Restaurar catálogo si se quieren conservar cambios. No copiar las preguntas de costos a otra base sin comprobar su sentido.

\fi
\begingroup
\setbox0=\vbox{\hsize=\linewidth
\everypar{}\csname @newlistfalse\endcsname
\textbf{Minutos 24-29.} Revise roles y permisos con la matriz. Si el ejercicio incluye crearlos, prepare Analista BI y Gerente comercial y asigne sólo cuentas de práctica autorizadas.

{\small\textbf{Tiempo máximo previsto de este bloque:} 5 minutos.\par}
\textbf{Qué debe observar:} Se distingue crear un rol, asignarle permisos y asociarlo a una cuenta. Gerencia puede leer y exportar, sin aprobar o ejecutar ETL.

\textbf{Qué registrar:} Definiciones y asignaciones realizadas o paso no ejecutado. Registrar aparte cualquier prueba de login realmente realizada.

\textbf{Si no se cumple:} Si sigue usando la cuenta administrativa, indicarlo. Revisar definiciones no acredita haber operado toda la plataforma como analista o gerente.

\par
}
\dimen0=\dimexpr\ht0+\dp0+2\baselineskip\relax
\ifdim\dimen0<\textheight
\Needspace{\dimen0}
\endgroup
\item \begin{samepage}
\textbf{Minutos 24-29.} Revise roles y permisos con la matriz. Si el ejercicio incluye crearlos, prepare Analista BI y Gerente comercial y asigne sólo cuentas de práctica autorizadas.

{\small\textbf{Tiempo máximo previsto de este bloque:} 5 minutos.\par}
\textbf{Qué debe observar:} Se distingue crear un rol, asignarle permisos y asociarlo a una cuenta. Gerencia puede leer y exportar, sin aprobar o ejecutar ETL.

\textbf{Qué registrar:} Definiciones y asignaciones realizadas o paso no ejecutado. Registrar aparte cualquier prueba de login realmente realizada.

\textbf{Si no se cumple:} Si sigue usando la cuenta administrativa, indicarlo. Revisar definiciones no acredita haber operado toda la plataforma como analista o gerente.

\end{samepage}
\else
\endgroup
\item \textbf{Minutos 24-29.} Revise roles y permisos con la matriz. Si el ejercicio incluye crearlos, prepare Analista BI y Gerente comercial y asigne sólo cuentas de práctica autorizadas.

{\small\textbf{Tiempo máximo previsto de este bloque:} 5 minutos.\par}
\textbf{Qué debe observar:} Se distingue crear un rol, asignarle permisos y asociarlo a una cuenta. Gerencia puede leer y exportar, sin aprobar o ejecutar ETL.

\textbf{Qué registrar:} Definiciones y asignaciones realizadas o paso no ejecutado. Registrar aparte cualquier prueba de login realmente realizada.

\textbf{Si no se cumple:} Si sigue usando la cuenta administrativa, indicarlo. Revisar definiciones no acredita haber operado toda la plataforma como analista o gerente.

\fi
\begingroup
\setbox0=\vbox{\hsize=\linewidth
\everypar{}\csname @newlistfalse\endcsname
\textbf{Minutos 29-41.} En AdventureWorks utilice la necesidad básica de unidades por producto y mes. Revise viabilidad, conceptos y modelo; apruebe sólo lo válido, confirme la carga y revise conciliación.

{\small\textbf{Tiempo máximo previsto de este bloque:} 12 minutos.\par}
\textbf{Qué debe observar:} Un caso propio de línea de pedido y unidades. El proveedor puede proponer otras etiquetas; se revisan significado y receta, no una captura idéntica.

\textbf{Qué registrar:} Texto usado, instantánea, propuesta, ejecución, avisos comprendidos y resultado de conciliación, o la etapa precisa donde se interrumpió.

\textbf{Si no se cumple:} Si falla un control, detener ese caso y registrarlo. No sustituir el diseño por SQL libre ni completar a mano la hoja con las cifras de referencia.

\par
}
\dimen0=\dimexpr\ht0+\dp0+2\baselineskip\relax
\ifdim\dimen0<\textheight
\Needspace{\dimen0}
\endgroup
\item \begin{samepage}
\textbf{Minutos 29-41.} En AdventureWorks utilice la necesidad básica de unidades por producto y mes. Revise viabilidad, conceptos y modelo; apruebe sólo lo válido, confirme la carga y revise conciliación.

{\small\textbf{Tiempo máximo previsto de este bloque:} 12 minutos.\par}
\textbf{Qué debe observar:} Un caso propio de línea de pedido y unidades. El proveedor puede proponer otras etiquetas; se revisan significado y receta, no una captura idéntica.

\textbf{Qué registrar:} Texto usado, instantánea, propuesta, ejecución, avisos comprendidos y resultado de conciliación, o la etapa precisa donde se interrumpió.

\textbf{Si no se cumple:} Si falla un control, detener ese caso y registrarlo. No sustituir el diseño por SQL libre ni completar a mano la hoja con las cifras de referencia.

\end{samepage}
\else
\endgroup
\item \textbf{Minutos 29-41.} En AdventureWorks utilice la necesidad básica de unidades por producto y mes. Revise viabilidad, conceptos y modelo; apruebe sólo lo válido, confirme la carga y revise conciliación.

{\small\textbf{Tiempo máximo previsto de este bloque:} 12 minutos.\par}
\textbf{Qué debe observar:} Un caso propio de línea de pedido y unidades. El proveedor puede proponer otras etiquetas; se revisan significado y receta, no una captura idéntica.

\textbf{Qué registrar:} Texto usado, instantánea, propuesta, ejecución, avisos comprendidos y resultado de conciliación, o la etapa precisa donde se interrumpió.

\textbf{Si no se cumple:} Si falla un control, detener ese caso y registrarlo. No sustituir el diseño por SQL libre ni completar a mano la hoja con las cifras de referencia.

\fi
\begingroup
\setbox0=\vbox{\hsize=\linewidth
\everypar{}\csname @newlistfalse\endcsname
\textbf{Minutos 41-53.} Cambie a WideWorldImporters y repita el recorrido básico de unidades por producto y mes, conservando línea de factura y fecha de emisión.

{\small\textbf{Tiempo máximo previsto de este bloque:} 12 minutos.\par}
\textbf{Qué debe observar:} Segundo caso con contexto propio. No reutiliza las claves de AW ni afirma cobro, margen, moneda o entrega efectiva.

\textbf{Qué registrar:} Los identificadores propios de WWI y el resultado de cada control; contraste entre ambos eventos.

\textbf{Si no se cumple:} Si no alcanza el tiempo o el proveedor falla, dejar la etapa como pendiente o bloqueada y continuar en otra sesión. No completar una carga ficticia para cerrar la hora.

\par
}
\dimen0=\dimexpr\ht0+\dp0+2\baselineskip\relax
\ifdim\dimen0<\textheight
\Needspace{\dimen0}
\endgroup
\item \begin{samepage}
\textbf{Minutos 41-53.} Cambie a WideWorldImporters y repita el recorrido básico de unidades por producto y mes, conservando línea de factura y fecha de emisión.

{\small\textbf{Tiempo máximo previsto de este bloque:} 12 minutos.\par}
\textbf{Qué debe observar:} Segundo caso con contexto propio. No reutiliza las claves de AW ni afirma cobro, margen, moneda o entrega efectiva.

\textbf{Qué registrar:} Los identificadores propios de WWI y el resultado de cada control; contraste entre ambos eventos.

\textbf{Si no se cumple:} Si no alcanza el tiempo o el proveedor falla, dejar la etapa como pendiente o bloqueada y continuar en otra sesión. No completar una carga ficticia para cerrar la hora.

\end{samepage}
\else
\endgroup
\item \textbf{Minutos 41-53.} Cambie a WideWorldImporters y repita el recorrido básico de unidades por producto y mes, conservando línea de factura y fecha de emisión.

{\small\textbf{Tiempo máximo previsto de este bloque:} 12 minutos.\par}
\textbf{Qué debe observar:} Segundo caso con contexto propio. No reutiliza las claves de AW ni afirma cobro, margen, moneda o entrega efectiva.

\textbf{Qué registrar:} Los identificadores propios de WWI y el resultado de cada control; contraste entre ambos eventos.

\textbf{Si no se cumple:} Si no alcanza el tiempo o el proveedor falla, dejar la etapa como pendiente o bloqueada y continuar en otra sesión. No completar una carga ficticia para cerrar la hora.

\fi
\begingroup
\setbox0=\vbox{\hsize=\linewidth
\everypar{}\csname @newlistfalse\endcsname
\textbf{Minutos 53-58.} Consulte los resultados que sí quedaron conciliados. Cambie de fuente, aplique un filtro, compruebe una cifra y formule una pregunta breve al copiloto. Abra una exportación si se realizó.

{\small\textbf{Tiempo máximo previsto de este bloque:} 5 minutos.\par}
\textbf{Qué debe observar:} Se entiende el alcance de los datos propios. La referencia de 274.914 unidades AW o 8.950.628 WWI sólo sirve si la copia y selección equivalen al histórico y la conciliación lo confirma.

\textbf{Qué registrar:} Fuente, filtros, cálculo propio y respuesta o fallo del chat. Archivos realmente abiertos con su contexto.

\textbf{Si no se cumple:} No compare como equivalentes dos eventos o monedas diferentes. Si sólo una fuente quedó completada, describa ese resultado parcial.

\par
}
\dimen0=\dimexpr\ht0+\dp0+2\baselineskip\relax
\ifdim\dimen0<\textheight
\Needspace{\dimen0}
\endgroup
\item \begin{samepage}
\textbf{Minutos 53-58.} Consulte los resultados que sí quedaron conciliados. Cambie de fuente, aplique un filtro, compruebe una cifra y formule una pregunta breve al copiloto. Abra una exportación si se realizó.

{\small\textbf{Tiempo máximo previsto de este bloque:} 5 minutos.\par}
\textbf{Qué debe observar:} Se entiende el alcance de los datos propios. La referencia de 274.914 unidades AW o 8.950.628 WWI sólo sirve si la copia y selección equivalen al histórico y la conciliación lo confirma.

\textbf{Qué registrar:} Fuente, filtros, cálculo propio y respuesta o fallo del chat. Archivos realmente abiertos con su contexto.

\textbf{Si no se cumple:} No compare como equivalentes dos eventos o monedas diferentes. Si sólo una fuente quedó completada, describa ese resultado parcial.

\end{samepage}
\else
\endgroup
\item \textbf{Minutos 53-58.} Consulte los resultados que sí quedaron conciliados. Cambie de fuente, aplique un filtro, compruebe una cifra y formule una pregunta breve al copiloto. Abra una exportación si se realizó.

{\small\textbf{Tiempo máximo previsto de este bloque:} 5 minutos.\par}
\textbf{Qué debe observar:} Se entiende el alcance de los datos propios. La referencia de 274.914 unidades AW o 8.950.628 WWI sólo sirve si la copia y selección equivalen al histórico y la conciliación lo confirma.

\textbf{Qué registrar:} Fuente, filtros, cálculo propio y respuesta o fallo del chat. Archivos realmente abiertos con su contexto.

\textbf{Si no se cumple:} No compare como equivalentes dos eventos o monedas diferentes. Si sólo una fuente quedó completada, describa ese resultado parcial.

\fi
\begingroup
\setbox0=\vbox{\hsize=\linewidth
\everypar{}\csname @newlistfalse\endcsname
\textbf{Minutos 58-60.} Complete la hoja y responda qué hizo la IA, qué controló el servidor y qué aprobó una persona. Elija la próxima práctica.

{\small\textbf{Tiempo máximo previsto de este bloque:} 2 minutos.\par}
\textbf{Qué debe observar:} Queda una bitácora reproducible con avances, límites e incidencias, incluso si no concluyeron las dos cargas.

\textbf{Qué registrar:} Estado de cada fuente y tareas pendientes específicas.

\textbf{Si no se cumple:} No calificar todo el recorrido como exitoso si falta evidencia de alguna etapa.

\par
}
\dimen0=\dimexpr\ht0+\dp0+2\baselineskip\relax
\ifdim\dimen0<\textheight
\Needspace{\dimen0}
\endgroup
\item \begin{samepage}
\textbf{Minutos 58-60.} Complete la hoja y responda qué hizo la IA, qué controló el servidor y qué aprobó una persona. Elija la próxima práctica.

{\small\textbf{Tiempo máximo previsto de este bloque:} 2 minutos.\par}
\textbf{Qué debe observar:} Queda una bitácora reproducible con avances, límites e incidencias, incluso si no concluyeron las dos cargas.

\textbf{Qué registrar:} Estado de cada fuente y tareas pendientes específicas.

\textbf{Si no se cumple:} No calificar todo el recorrido como exitoso si falta evidencia de alguna etapa.

\end{samepage}
\else
\endgroup
\item \textbf{Minutos 58-60.} Complete la hoja y responda qué hizo la IA, qué controló el servidor y qué aprobó una persona. Elija la próxima práctica.

{\small\textbf{Tiempo máximo previsto de este bloque:} 2 minutos.\par}
\textbf{Qué debe observar:} Queda una bitácora reproducible con avances, límites e incidencias, incluso si no concluyeron las dos cargas.

\textbf{Qué registrar:} Estado de cada fuente y tareas pendientes específicas.

\textbf{Si no se cumple:} No calificar todo el recorrido como exitoso si falta evidencia de alguna etapa.

\fi
\end{enumerate}
\subsection{Antes de dar por terminado el ejercicio}
\begin{itemize}
\item Las fuentes e instantáneas se mantuvieron separadas.
\item Cada generación/carga realizada tiene contexto e identificadores propios.
\item Se registró cuál de los dos ejercicios terminó y cuál no.
\item La configuración de roles se distinguió de una prueba de sesión por perfil.
\item No se usaron números del histórico como resultados obtenidos de una instalación limpia.
\end{itemize}

\subsection{Necesidades básicas para la práctica}
\paragraph{AdventureWorks 2022} Conservar cada línea de pedido de venta para comparar mensualmente la suma de unidades por producto. Usar la fecha del pedido y un nombre de producto comprobable. Mantener la referencia de cada línea a su pedido. Limitar esta práctica al total de unidades y su evolución; no interpretar estos registros como facturas, cobros o pronósticos.

Ejercicio nuevo propuesto; resultado propio pendiente de ejecución.

\paragraph{WideWorldImporters} Conservar cada línea de factura para comparar mensualmente la suma de unidades facturadas por producto. Usar la fecha de emisión y el producto con nombre verificable. Mantener el identificador de factura y el de línea para conservar la trazabilidad. Limitar la práctica a unidades y evolución; no afirmar moneda, cobro, utilidad ni entrega realizada.

Ejercicio nuevo propuesto; resultado propio pendiente de ejecución.

\section{Qué hacer cuando una etapa no termina}
El protocolo protege el aprendizaje y la evidencia. No requiere que la demostración parezca perfecta. Hay que distinguir un rechazo correcto, un fallo externo y una inconsistencia del producto.

\Needspace{8\baselineskip}
\begingroup
\setbox0=\vbox{\hsize=\linewidth
{\normalfont\normalsize\sffamily\bfseries\normalsize\color{Teal}\hyphenpenalty=10000\exhyphenpenalty=10000 Reconocer la situación y actuar con criterio\par}

\par
\begingroup
\fontsize{10}{12.5}\selectfont
\setlength{\tabcolsep}{5pt}
\renewcommand{\arraystretch}{1.15}
\begin{tabular}{@{}>{\RaggedRight\arraybackslash}p{0.22\linewidth}>{\RaggedRight\arraybackslash}p{0.3\linewidth}>{\RaggedRight\arraybackslash}p{0.39\linewidth}@{}}
\toprule
\textbf{Qué ocurre} & \textbf{Qué hacer} & \textbf{Qué no debe afirmarse} \\
\midrule
Respuesta de IA pendiente o demasiado lenta & Conservar la pantalla y el contexto. No pulsar varias veces. Si termina el tiempo del bloque, anotar pendiente y comprobar historial antes de reintentar. & No decir que se generó, aprobó o cargó algo sin el comprobante correspondiente. \\
\bottomrule
\end{tabular}
\endgroup
}
\Needspace{\dimexpr\ht0+\dp0+48pt\relax}
\endgroup
\subsection{Reconocer la situación y actuar con criterio}
\begingroup
\fontsize{10}{12.5}\selectfont
\setlength{\tabcolsep}{5pt}
\renewcommand{\arraystretch}{1.15}
\begin{longtable}{@{}>{\RaggedRight\arraybackslash}p{0.22\linewidth}>{\RaggedRight\arraybackslash}p{0.3\linewidth}>{\RaggedRight\arraybackslash}p{0.39\linewidth}@{}}
\toprule
\textbf{Qué ocurre} & \textbf{Qué hacer} & \textbf{Qué no debe afirmarse} \\
\midrule
\endfirsthead
\toprule
\textbf{Qué ocurre} & \textbf{Qué hacer} & \textbf{Qué no debe afirmarse} \\
\midrule
\endhead
\midrule
\multicolumn{3}{r}{\footnotesize Continúa en la página siguiente}\\
\endfoot
\bottomrule
\endlastfoot
Respuesta de IA pendiente o demasiado lenta & Conservar la pantalla y el contexto. No pulsar varias veces. Si termina el tiempo del bloque, anotar pendiente y comprobar historial antes de reintentar. & No decir que se generó, aprobó o cargó algo sin el comprobante correspondiente. \\
\addlinespace[3pt]
El proveedor informa cuota, capacidad o conectividad insuficiente & Registrar el mensaje y la etapa. Una persona autorizada puede revisar la configuración o acordar otro proveedor; si cambia el destino, revisar su prueba y autorización de envío. & No cambiar silenciosamente de modelo ni atribuir éxito a un intento interrumpido. \\
\addlinespace[3pt]
La respuesta no cumple el formato o contrato & Aceptar el rechazo de la plataforma; revisar la solicitud o hacer un único reintento manual si el error es definitivo y queda tiempo. & No pegar una respuesta inválida a mano ni desactivar validaciones para continuar. \\
\addlinespace[3pt]
La propuesta cambia el objetivo, evento o rol de una relación & Mantener el propósito original, aclarar la ambigüedad y revisar una versión nueva cuando corresponda. Si no puede justificarse, no aprobar. & No afirmar que una FK válida demuestra por sí sola la interpretación comercial. \\
\addlinespace[3pt]
Advertencia o requisito no disponible & Leer la limitación y decidir si el alcance restante sigue siendo útil. Si falta un requisito esencial, corregir la necesidad. & Aceptar el aviso no hace que exista la capacidad ausente. \\
\addlinespace[3pt]
Se pierde la respuesta durante una carga & Recuperar el expediente y leer su estado antes de iniciar cualquier otra operación. Registrar si terminó, falló o sigue en curso. & No reenviar por reflejo ni anunciar conciliación por el mero paso del tiempo. \\
\addlinespace[3pt]
Hay diferencia numérica, fuente equivocada o denominador incoherente & Suspender la conclusión, guardar contexto y comparar universos equivalentes. Registrar esperado y obtenido sin alterar filas ni la evidencia anterior. & No ajustar las cifras para que coincidan con la guía ni sumar monedas sin evidencia. \\
\addlinespace[3pt]
Se muestra una explicación o un archivo histórico como respaldo & Anunciar fuente, fecha, propuesta y ejecución histórica. Separarlo del estado del intento nuevo. & Un respaldo histórico no prueba que la operación nueva se completó. \\
\addlinespace[3pt]
El rol no tiene permiso o no aparece una opción & Pedir revisión a administración y registrar el permiso faltante. Continuar sólo las operaciones autorizadas. & No presentar una sesión administrativa como prueba de un perfil con menos permisos. \\
\addlinespace[3pt]
El chat explica algo que las cifras no sostienen & Conservar pregunta, interpretación, evidencia, respuesta y filtros. Contrastar números y cautelas antes de usar la narración. & Una frase fluida no demuestra causalidad, rentabilidad, pronóstico o precisión universal. \\
\addlinespace[3pt]
\end{longtable}
\endgroup

\begingroup
\setbox0=\vbox{\hsize=\linewidth
{\normalfont\normalsize\RaggedRight\sffamily\bfseries\large\color{Azul}\hyphenpenalty=10000\exhyphenpenalty=10000 Rúbrica de aprendizaje: veinte puntos\par}

Guía orientativa para autoevaluación o acompañamiento docente; no es un instrumento validado de usabilidad ni una medición de la eficacia del producto.

\textbf{Puntuación máxima:} 20 puntos.

\par
\begingroup
\fontsize{10}{12.5}\selectfont
\setlength{\tabcolsep}{5pt}
\renewcommand{\arraystretch}{1.15}
\begin{tabular}{@{}>{\RaggedRight\arraybackslash}p{0.22\linewidth}>{\RaggedRight\arraybackslash}p{0.3\linewidth}>{\RaggedRight\arraybackslash}p{0.39\linewidth}@{}}
\toprule
\textbf{Criterio} & \textbf{Puntos máximos} & \textbf{Evidencia que se espera} \\
\midrule
Comprender la decisión comercial & 4 & Explica evento, grano, indicadores, unidades y comparación sin confundir pedido, factura o cobro. \\
\bottomrule
\end{tabular}
\endgroup
}
\Needspace{\dimexpr\ht0+\dp0+48pt\relax}
\endgroup
\section{Rúbrica de aprendizaje: veinte puntos}
Guía orientativa para autoevaluación o acompañamiento docente; no es un instrumento validado de usabilidad ni una medición de la eficacia del producto.

\textbf{Puntuación máxima:} 20 puntos.

\begingroup
\fontsize{10}{12.5}\selectfont
\setlength{\tabcolsep}{5pt}
\renewcommand{\arraystretch}{1.15}
\begin{longtable}{@{}>{\RaggedRight\arraybackslash}p{0.22\linewidth}>{\RaggedRight\arraybackslash}p{0.3\linewidth}>{\RaggedRight\arraybackslash}p{0.39\linewidth}@{}}
\toprule
\textbf{Criterio} & \textbf{Puntos máximos} & \textbf{Evidencia que se espera} \\
\midrule
\endfirsthead
\toprule
\textbf{Criterio} & \textbf{Puntos máximos} & \textbf{Evidencia que se espera} \\
\midrule
\endhead
\midrule
\multicolumn{3}{r}{\footnotesize Continúa en la página siguiente}\\
\endfoot
\bottomrule
\endlastfoot
Comprender la decisión comercial & 4 & Explica evento, grano, indicadores, unidades y comparación sin confundir pedido, factura o cobro. \\
\addlinespace[3pt]
Preparar el contexto de cada fuente & 4 & Distingue conexión, instantánea y catálogo; conserva sus identidades y usa preguntas y período apropiados. \\
\addlinespace[3pt]
Supervisar y recuperar con criterio & 4 & Distingue IA, reglas y decisión humana; revisa advertencias, no aprueba errores y registra con honestidad una interrupción. \\
\addlinespace[3pt]
Comprobar e interpretar los resultados & 4 & Contrasta conciliación, filtros, moneda, sumas, documentos distintos y denominadores de una razón; o explica correctamente qué control no pudo ejecutar. \\
\addlinespace[3pt]
Cuidar acceso, evidencia y comunicación & 4 & Trabaja con la identidad autorizada, evita exponer secretos y presenta contexto, archivos y procedencia sin confundir lo histórico con lo nuevo. \\
\addlinespace[3pt]
\end{longtable}
\endgroup

\Needspace{8\baselineskip}
\begingroup
\setbox0=\vbox{\hsize=\linewidth
{\normalfont\normalsize\sffamily\bfseries\normalsize\color{Teal}\hyphenpenalty=10000\exhyphenpenalty=10000 Cómo asignar los puntos\par}

\par
\begingroup
\fontsize{10}{12.5}\selectfont
\setlength{\tabcolsep}{5pt}
\renewcommand{\arraystretch}{1.15}
\begin{tabular}{@{}>{\RaggedRight\arraybackslash}p{0.27\linewidth}>{\RaggedRight\arraybackslash}p{0.66\linewidth}@{}}
\toprule
\textbf{Puntos} & \textbf{Descripción} \\
\midrule
4 & Lo realiza o explica con autonomía, evidencia suficiente y sin errores sustantivos. \\
\bottomrule
\end{tabular}
\endgroup
}
\Needspace{\dimexpr\ht0+\dp0+48pt\relax}
\endgroup
\subsection{Cómo asignar los puntos}
\begingroup
\fontsize{10}{12.5}\selectfont
\setlength{\tabcolsep}{5pt}
\renewcommand{\arraystretch}{1.15}
\begin{longtable}{@{}>{\RaggedRight\arraybackslash}p{0.27\linewidth}>{\RaggedRight\arraybackslash}p{0.66\linewidth}@{}}
\toprule
\textbf{Puntos} & \textbf{Descripción} \\
\midrule
\endfirsthead
\toprule
\textbf{Puntos} & \textbf{Descripción} \\
\midrule
\endhead
\midrule
\multicolumn{2}{r}{\footnotesize Continúa en la página siguiente}\\
\endfoot
\bottomrule
\endlastfoot
4 & Lo realiza o explica con autonomía, evidencia suficiente y sin errores sustantivos. \\
\addlinespace[3pt]
3 & Lo realiza de forma correcta con una aclaración menor. \\
\addlinespace[3pt]
2 & Comprensión parcial; necesita apoyo para completar o justificar una decisión. \\
\addlinespace[3pt]
1 & Reconoce parte del concepto, pero requiere acompañamiento continuo. \\
\addlinespace[3pt]
0 & No presenta evidencia de aprendizaje o mantiene una interpretación incorrecta después de revisarla. \\
\addlinespace[3pt]
\end{longtable}
\endgroup

Marcar No evaluado cuando no hubo oportunidad de observar el criterio. No asignar cuatro puntos por un paso no realizado ni publicar un total sobre veinte como si todos se hubieran evaluado. Puede registrarse por separado una explicación verbal correcta y una operación bloqueada.

Con los cinco criterios evaluados, dieciséis puntos o más puede utilizarse como orientación para repetir el caso con menor ayuda. El umbral es didáctico y no representa una validación estadística del producto.

\subsection{Condiciones que obligan a detenerse y corregir}
\begin{itemize}
\item Exponer una credencial, aprobar una propuesta con error no resuelto, mezclar deliberadamente fuentes/monedas o declarar una carga inexistente requiere detener y corregir la práctica, aunque la suma de otros criterios sea alta.
\item Un fallo externo de cuota no se convierte en una falta de aprendizaje si se interpreta y registra correctamente; la operación continúa como bloqueada o no ejecutada, sin calificarse como ciclo funcional terminado.
\end{itemize}

\section{Hoja de registro de una práctica nueva}
Copiar una hoja por fuente y ejecución. Completar el resultado esperado antes de ejecutar y el obtenido después.

Dejar sin valor lo que no se obtuvo y marcar su estado. No copiar las cifras históricas en la columna de resultado obtenido.

Usar alias de práctica y el rol real; no escribir contraseñas, API keys, tokens ni datos personales innecesarios.

\Needspace{8\baselineskip}
\begingroup
\setbox0=\vbox{\hsize=\linewidth
{\normalfont\normalsize\sffamily\bfseries\normalsize\color{Teal}\hyphenpenalty=10000\exhyphenpenalty=10000 Datos de contexto de su práctica\par}

\par
\begingroup
\fontsize{10}{12.5}\selectfont
\setlength{\tabcolsep}{5pt}
\renewcommand{\arraystretch}{1.15}
\begin{tabular}{@{}>{\RaggedRight\arraybackslash}p{0.27\linewidth}>{\RaggedRight\arraybackslash}p{0.66\linewidth}@{}}
\toprule
\textbf{Dato que debe registrar} & \textbf{Valor de su ejecución} \\
\midrule
Fecha y hora &          \\
\bottomrule
\end{tabular}
\endgroup
}
\Needspace{\dimexpr\ht0+\dp0+48pt\relax}
\endgroup
\subsection{Datos de contexto de su práctica}
\begingroup
\fontsize{10}{12.5}\selectfont
\setlength{\tabcolsep}{5pt}
\renewcommand{\arraystretch}{1.15}
\begin{longtable}{@{}>{\RaggedRight\arraybackslash}p{0.27\linewidth}>{\RaggedRight\arraybackslash}p{0.66\linewidth}@{}}
\toprule
\textbf{Dato que debe registrar} & \textbf{Valor de su ejecución} \\
\midrule
\endfirsthead
\toprule
\textbf{Dato que debe registrar} & \textbf{Valor de su ejecución} \\
\midrule
\endhead
\midrule
\multicolumn{2}{r}{\footnotesize Continúa en la página siguiente}\\
\endfoot
\bottomrule
\endlastfoot
Fecha y hora &          \\
\addlinespace[3pt]
Participante o alias &          \\
\addlinespace[3pt]
Persona que acompana &          \\
\addlinespace[3pt]
Objetivo de aprendizaje &          \\
\addlinespace[3pt]
Tipo de actividad &          \\
\addlinespace[3pt]
Entorno o version visible &          \\
\addlinespace[3pt]
Rol y permisos reales &          \\
\addlinespace[3pt]
Fuente y base &          \\
\addlinespace[3pt]
Conexion local &          \\
\addlinespace[3pt]
Instantanea local &          \\
\addlinespace[3pt]
Proveedor modelo si se usa &          \\
\addlinespace[3pt]
Destino de envio autorizado &          \\
\addlinespace[3pt]
Necesidad textual &          \\
\addlinespace[3pt]
Preguntas de catalogo &          \\
\addlinespace[3pt]
Periodicidad &          \\
\addlinespace[3pt]
Evento y grano &          \\
\addlinespace[3pt]
Indicadores y unidades &          \\
\addlinespace[3pt]
Limitaciones aceptadas y motivo &          \\
\addlinespace[3pt]
Propuesta local &          \\
\addlinespace[3pt]
Propuesta origen si hay derivacion &          \\
\addlinespace[3pt]
Estado de revision &          \\
\addlinespace[3pt]
Ejecucion local &          \\
\addlinespace[3pt]
Estado de ejecucion &          \\
\addlinespace[3pt]
Contexto del filtro &          \\
\addlinespace[3pt]
Formula y denominador a comprobar &          \\
\addlinespace[3pt]
Resultado esperado y fuente del criterio &          \\
\addlinespace[3pt]
Resultado obtenido &          \\
\addlinespace[3pt]
Diferencia y explicacion &          \\
\addlinespace[3pt]
Archivo o captura sin secretos &          \\
\addlinespace[3pt]
Naturaleza y fecha de la evidencia &          \\
\addlinespace[3pt]
Incidencia y recuperacion &          \\
\addlinespace[3pt]
Pasos no ejecutados o pendientes &          \\
\addlinespace[3pt]
Conclusion con limites &          \\
\addlinespace[3pt]
Proxima accion &          \\
\addlinespace[3pt]
\end{longtable}
\endgroup

\subsection{Tipos de actividad}
\begin{itemize}
\item Consulta de expediente histórico
\item Generación y ejecución nuevas
\item Ejercicio propuesto aún no ejecutado
\item Prueba de acceso bajo un perfil realmente autenticado
\end{itemize}

\subsection{Estados que puede registrar}
\begin{itemize}
\item No ejecutado
\item En curso
\item Terminado y comprobado para el alcance indicado
\item Rechazo esperado observado
\item Bloqueado por proveedor
\item Bloqueado por permisos o preparación
\item Interrumpido
\item Resultado inconsistente por revisar
\end{itemize}

\Needspace{8\baselineskip}
\begingroup
\setbox0=\vbox{\hsize=\linewidth
{\normalfont\normalsize\sffamily\bfseries\normalsize\color{Teal}\hyphenpenalty=10000\exhyphenpenalty=10000 Registro de comprobaciones\par}

\par
\begingroup
\fontsize{10}{12.5}\selectfont
\setlength{\tabcolsep}{5pt}
\renewcommand{\arraystretch}{1.15}
\begin{tabular}{@{}>{\RaggedRight\arraybackslash}p{0.22\linewidth}>{\RaggedRight\arraybackslash}p{0.3\linewidth}>{\RaggedRight\arraybackslash}p{0.39\linewidth}@{}}
\toprule
\textbf{Etapa o indicador} & \textbf{Descripción} & \textbf{Datos del ejemplo} \\
\midrule
Viabilidad de la necesidad & \textbf{Fuente/propuesta/ejecución y filtros:} \par\smallskip\textbf{Fórmula o criterio esperado:} \par\smallskip\textbf{Resultado esperado antes de ejecutar:}  & \textbf{Resultado obtenido:} \par\smallskip\textbf{Estado:} \par\smallskip\textbf{Evidencia:}  \\
\bottomrule
\end{tabular}
\endgroup
}
\Needspace{\dimexpr\ht0+\dp0+48pt\relax}
\endgroup
\subsection{Registro de comprobaciones}
\begingroup
\fontsize{10}{12.5}\selectfont
\setlength{\tabcolsep}{5pt}
\renewcommand{\arraystretch}{1.15}
\begin{longtable}{@{}>{\RaggedRight\arraybackslash}p{0.22\linewidth}>{\RaggedRight\arraybackslash}p{0.3\linewidth}>{\RaggedRight\arraybackslash}p{0.39\linewidth}@{}}
\toprule
\textbf{Etapa o indicador} & \textbf{Descripción} & \textbf{Datos del ejemplo} \\
\midrule
\endfirsthead
\toprule
\textbf{Etapa o indicador} & \textbf{Descripción} & \textbf{Datos del ejemplo} \\
\midrule
\endhead
\midrule
\multicolumn{3}{r}{\footnotesize Continúa en la página siguiente}\\
\endfoot
\bottomrule
\endlastfoot
Viabilidad de la necesidad & \textbf{Fuente/propuesta/ejecución y filtros:} \par\smallskip\textbf{Fórmula o criterio esperado:} \par\smallskip\textbf{Resultado esperado antes de ejecutar:}  & \textbf{Resultado obtenido:} \par\smallskip\textbf{Estado:} \par\smallskip\textbf{Evidencia:}  \\
\addlinespace[3pt]
Contrato y revisión humana & \textbf{Fuente/propuesta/ejecución y filtros:} \par\smallskip\textbf{Fórmula o criterio esperado:} \par\smallskip\textbf{Resultado esperado antes de ejecutar:}  & \textbf{Resultado obtenido:} \par\smallskip\textbf{Estado:} \par\smallskip\textbf{Evidencia:}  \\
\addlinespace[3pt]
Preparación y ejecución de carga & \textbf{Fuente/propuesta/ejecución y filtros:} \par\smallskip\textbf{Fórmula o criterio esperado:} \par\smallskip\textbf{Resultado esperado antes de ejecutar:}  & \textbf{Resultado obtenido:} \par\smallskip\textbf{Estado:} \par\smallskip\textbf{Evidencia:}  \\
\addlinespace[3pt]
Conciliación de filas e importes & \textbf{Fuente/propuesta/ejecución y filtros:} \par\smallskip\textbf{Fórmula o criterio esperado:} \par\smallskip\textbf{Resultado esperado antes de ejecutar:}  & \textbf{Resultado obtenido:} \par\smallskip\textbf{Estado:} \par\smallskip\textbf{Evidencia:}  \\
\addlinespace[3pt]
Panel y filtro & \textbf{Fuente/propuesta/ejecución y filtros:} \par\smallskip\textbf{Fórmula o criterio esperado:} \par\smallskip\textbf{Resultado esperado antes de ejecutar:}  & \textbf{Resultado obtenido:} \par\smallskip\textbf{Estado:} \par\smallskip\textbf{Evidencia:}  \\
\addlinespace[3pt]
Conteo distinto y razón, si aplican & \textbf{Fuente/propuesta/ejecución y filtros:} \par\smallskip\textbf{Fórmula o criterio esperado:} \par\smallskip\textbf{Resultado esperado antes de ejecutar:}  & \textbf{Resultado obtenido:} \par\smallskip\textbf{Estado:} \par\smallskip\textbf{Evidencia:}  \\
\addlinespace[3pt]
Copiloto y cautelas & \textbf{Fuente/propuesta/ejecución y filtros:} \par\smallskip\textbf{Fórmula o criterio esperado:} \par\smallskip\textbf{Resultado esperado antes de ejecutar:}  & \textbf{Resultado obtenido:} \par\smallskip\textbf{Estado:} \par\smallskip\textbf{Evidencia:}  \\
\addlinespace[3pt]
Exportación abierta & \textbf{Fuente/propuesta/ejecución y filtros:} \par\smallskip\textbf{Fórmula o criterio esperado:} \par\smallskip\textbf{Resultado esperado antes de ejecutar:}  & \textbf{Resultado obtenido:} \par\smallskip\textbf{Estado:} \par\smallskip\textbf{Evidencia:}  \\
\addlinespace[3pt]
Acceso con el rol de práctica & \textbf{Fuente/propuesta/ejecución y filtros:} \par\smallskip\textbf{Fórmula o criterio esperado:} \par\smallskip\textbf{Resultado esperado antes de ejecutar:}  & \textbf{Resultado obtenido:} \par\smallskip\textbf{Estado:} \par\smallskip\textbf{Evidencia:}  \\
\addlinespace[3pt]
\end{longtable}
\endgroup

\Needspace{8\baselineskip}
\begingroup
\setbox0=\vbox{\hsize=\linewidth
{\normalfont\normalsize\sffamily\bfseries\normalsize\color{Teal}\hyphenpenalty=10000\exhyphenpenalty=10000 Registro de autoevaluación\par}

\par
\begingroup
\fontsize{10}{12.5}\selectfont
\setlength{\tabcolsep}{5pt}
\renewcommand{\arraystretch}{1.15}
\begin{tabular}{@{}>{\RaggedRight\arraybackslash}p{0.22\linewidth}>{\RaggedRight\arraybackslash}p{0.3\linewidth}>{\RaggedRight\arraybackslash}p{0.39\linewidth}@{}}
\toprule
\textbf{Criterio} & \textbf{Puntuación 0-4 o No evaluado} & \textbf{Evidencia o apoyo requerido} \\
\midrule
Comprender la decisión comercial &  &  \\
\bottomrule
\end{tabular}
\endgroup
}
\Needspace{\dimexpr\ht0+\dp0+48pt\relax}
\endgroup
\subsection{Registro de autoevaluación}
\begingroup
\fontsize{10}{12.5}\selectfont
\setlength{\tabcolsep}{5pt}
\renewcommand{\arraystretch}{1.15}
\begin{longtable}{@{}>{\RaggedRight\arraybackslash}p{0.22\linewidth}>{\RaggedRight\arraybackslash}p{0.3\linewidth}>{\RaggedRight\arraybackslash}p{0.39\linewidth}@{}}
\toprule
\textbf{Criterio} & \textbf{Puntuación 0-4 o No evaluado} & \textbf{Evidencia o apoyo requerido} \\
\midrule
\endfirsthead
\toprule
\textbf{Criterio} & \textbf{Puntuación 0-4 o No evaluado} & \textbf{Evidencia o apoyo requerido} \\
\midrule
\endhead
\midrule
\multicolumn{3}{r}{\footnotesize Continúa en la página siguiente}\\
\endfoot
\bottomrule
\endlastfoot
Comprender la decisión comercial &  &  \\
\addlinespace[3pt]
Preparar el contexto de cada fuente &  &  \\
\addlinespace[3pt]
Supervisar y recuperar con criterio &  &  \\
\addlinespace[3pt]
Comprobar e interpretar los resultados &  &  \\
\addlinespace[3pt]
Cuidar acceso, evidencia y comunicación &  &  \\
\addlinespace[3pt]
\end{longtable}
\endgroup

\subsection{Conclusión de la práctica}
Un registro con fases bloqueadas o no ejecutadas sigue siendo útil si permite reconocer qué se hizo y qué falta. No demuestra un ciclo completo hasta contar con sus evidencias.

\section{Referencias históricas para contrastar los casos}
Estos valores corresponden a la campaña del 6 de noviembre de 2026. Los IDs y resultados no se fuerzan en una instalación nueva. Antes de comparar cifras se comprueban fuente, copia, contrato, período, filtros y receta. En WWI la unidad permanece como moneda de origen.

\begingroup
\setbox0=\vbox{\hsize=\linewidth
{\normalfont\normalsize\sffamily\bfseries\normalsize\color{Teal}\hyphenpenalty=10000\exhyphenpenalty=10000 AdventureWorks 2022\par}

\par
\begingroup
\fontsize{10}{12.5}\selectfont
\setlength{\tabcolsep}{5pt}
\renewcommand{\arraystretch}{1.15}
\begin{tabular}{@{}>{\RaggedRight\arraybackslash}p{0.27\linewidth}>{\RaggedRight\arraybackslash}p{0.66\linewidth}@{}}
\toprule
\textbf{Dato} & \textbf{Contenido} \\
\midrule
Fecha evidencia & 2026-11-06 \\
\bottomrule
\end{tabular}
\endgroup
}
\Needspace{\dimexpr\ht0+\dp0+48pt\relax}
\endgroup
\subsection{AdventureWorks 2022}
\begingroup
\fontsize{10}{12.5}\selectfont
\setlength{\tabcolsep}{5pt}
\renewcommand{\arraystretch}{1.15}
\begin{longtable}{@{}>{\RaggedRight\arraybackslash}p{0.27\linewidth}>{\RaggedRight\arraybackslash}p{0.66\linewidth}@{}}
\toprule
\textbf{Dato} & \textbf{Contenido} \\
\midrule
\endfirsthead
\toprule
\textbf{Dato} & \textbf{Contenido} \\
\midrule
\endhead
\midrule
\multicolumn{2}{r}{\footnotesize Continúa en la página siguiente}\\
\endfoot
\bottomrule
\endlastfoot
Fecha evidencia & 2026-11-06 \\
\addlinespace[3pt]
Propuesta local & 129 \\
\addlinespace[3pt]
Ejecucion local & 21 \\
\addlinespace[3pt]
Granularidad & Línea de pedido \\
\addlinespace[3pt]
Filas de origen & 121317 \\
\addlinespace[3pt]
Filas de destino & 121317 \\
\addlinespace[3pt]
Diferencia de filas & 0 \\
\addlinespace[3pt]
Pregunta chat & Resume las unidades vendidas en 2013 para Northwest. Indica el total y el mes con mayor cantidad, sin atribuir causas ni mezclar importes. \\
\addlinespace[3pt]
Alcance & Evidencia histórica comprobada en esa campaña. No es un resultado nuevo obtenido al redactar este material ni el contenido garantizado de un portable limpio. \\
\addlinespace[3pt]
\end{longtable}
\endgroup

\begingroup
\setbox0=\vbox{\hsize=\linewidth
{\normalfont\normalsize\normalfont\normalsize\bfseries\hyphenpenalty=10000\exhyphenpenalty=10000 Necesidad de referencia\par}
 Conservar cada línea de pedido de venta en un datamart. Para priorizar productos y territorios, comparar mensualmente cuatro indicadores: suma del importe registrado de las líneas, suma de unidades, conteo de identificadores de pedido distintos e importe total dividido entre pedidos distintos. Usar fecha y territorio comercial del pedido, producto y cliente con nombre verificable. Los rankings comparan sumas por entidad. Mantener la trazabilidad de cada línea al pedido. La moneda se reportará sólo si se comprueba en la fuente.

\par\vspace{6pt}\noindent\textbf{Selección completa}\hspace{0.5em}
\par
\begingroup
\fontsize{10}{12.5}\selectfont
\setlength{\tabcolsep}{5pt}
\renewcommand{\arraystretch}{1.15}
\begin{tabular}{@{}>{\RaggedRight\arraybackslash}p{0.27\linewidth}>{\RaggedRight\arraybackslash}p{0.66\linewidth}@{}}
\toprule
\textbf{Dato} & \textbf{Contenido} \\
\midrule
Importe exacto & 109846381.399888 \\
\bottomrule
\end{tabular}
\endgroup
}
\Needspace{\dimexpr\ht0+\dp0+48pt\relax}
\endgroup
\paragraph{Necesidad de referencia} Conservar cada línea de pedido de venta en un datamart. Para priorizar productos y territorios, comparar mensualmente cuatro indicadores: suma del importe registrado de las líneas, suma de unidades, conteo de identificadores de pedido distintos e importe total dividido entre pedidos distintos. Usar fecha y territorio comercial del pedido, producto y cliente con nombre verificable. Los rankings comparan sumas por entidad. Mantener la trazabilidad de cada línea al pedido. La moneda se reportará sólo si se comprueba en la fuente.

\paragraph{Selección completa}
\begingroup
\fontsize{10}{12.5}\selectfont
\setlength{\tabcolsep}{5pt}
\renewcommand{\arraystretch}{1.15}
\begin{longtable}{@{}>{\RaggedRight\arraybackslash}p{0.27\linewidth}>{\RaggedRight\arraybackslash}p{0.66\linewidth}@{}}
\toprule
\textbf{Dato} & \textbf{Contenido} \\
\midrule
\endfirsthead
\toprule
\textbf{Dato} & \textbf{Contenido} \\
\midrule
\endhead
\midrule
\multicolumn{2}{r}{\footnotesize Continúa en la página siguiente}\\
\endfoot
\bottomrule
\endlastfoot
Importe exacto & 109846381.399888 \\
\addlinespace[3pt]
Moneda & USD comprobado \\
\addlinespace[3pt]
Unidades & 274914 \\
\addlinespace[3pt]
Pedidos distintos & 31465 \\
\addlinespace[3pt]
Importe por pedido exacto & 3491.065672966407 \\
\addlinespace[3pt]
\end{longtable}
\endgroup

\begingroup
\setbox0=\vbox{\hsize=\linewidth
{\normalfont\normalsize\normalfont\normalsize\bfseries\hyphenpenalty=10000\exhyphenpenalty=10000 Selección filtrada\par}

\par
\begingroup
\fontsize{10}{12.5}\selectfont
\setlength{\tabcolsep}{5pt}
\renewcommand{\arraystretch}{1.15}
\begin{tabular}{@{}>{\RaggedRight\arraybackslash}p{0.27\linewidth}>{\RaggedRight\arraybackslash}p{0.66\linewidth}@{}}
\toprule
\textbf{Dato} & \textbf{Contenido} \\
\midrule
Año & 2013 \\
\bottomrule
\end{tabular}
\endgroup
}
\Needspace{\dimexpr\ht0+\dp0+48pt\relax}
\endgroup
\paragraph{Selección filtrada}
\begingroup
\fontsize{10}{12.5}\selectfont
\setlength{\tabcolsep}{5pt}
\renewcommand{\arraystretch}{1.15}
\begin{longtable}{@{}>{\RaggedRight\arraybackslash}p{0.27\linewidth}>{\RaggedRight\arraybackslash}p{0.66\linewidth}@{}}
\toprule
\textbf{Dato} & \textbf{Contenido} \\
\midrule
\endfirsthead
\toprule
\textbf{Dato} & \textbf{Contenido} \\
\midrule
\endhead
\midrule
\multicolumn{2}{r}{\footnotesize Continúa en la página siguiente}\\
\endfoot
\bottomrule
\endlastfoot
Año & 2013 \\
\addlinespace[3pt]
Territorio & Northwest \\
\addlinespace[3pt]
Importe exacto & 6015174.735836 \\
\addlinespace[3pt]
Unidades & 17452 \\
\addlinespace[3pt]
Pedidos distintos & 2053 \\
\addlinespace[3pt]
Importe por pedido exacto & 2929.943855740867 \\
\addlinespace[3pt]
Mes de mayor cantidad & junio de 2013 \\
\addlinespace[3pt]
Unidades del mes de mayor cantidad & 2494 \\
\addlinespace[3pt]
\end{longtable}
\endgroup

\begingroup
\setbox0=\vbox{\hsize=\linewidth
{\normalfont\normalsize\sffamily\bfseries\normalsize\color{Teal}\hyphenpenalty=10000\exhyphenpenalty=10000 WideWorldImporters\par}

\par
\begingroup
\fontsize{10}{12.5}\selectfont
\setlength{\tabcolsep}{5pt}
\renewcommand{\arraystretch}{1.15}
\begin{tabular}{@{}>{\RaggedRight\arraybackslash}p{0.27\linewidth}>{\RaggedRight\arraybackslash}p{0.66\linewidth}@{}}
\toprule
\textbf{Dato} & \textbf{Contenido} \\
\midrule
Fecha evidencia & 2026-11-06 \\
\bottomrule
\end{tabular}
\endgroup
}
\Needspace{\dimexpr\ht0+\dp0+48pt\relax}
\endgroup
\subsection{WideWorldImporters}
\begingroup
\fontsize{10}{12.5}\selectfont
\setlength{\tabcolsep}{5pt}
\renewcommand{\arraystretch}{1.15}
\begin{longtable}{@{}>{\RaggedRight\arraybackslash}p{0.27\linewidth}>{\RaggedRight\arraybackslash}p{0.66\linewidth}@{}}
\toprule
\textbf{Dato} & \textbf{Contenido} \\
\midrule
\endfirsthead
\toprule
\textbf{Dato} & \textbf{Contenido} \\
\midrule
\endhead
\midrule
\multicolumn{2}{r}{\footnotesize Continúa en la página siguiente}\\
\endfoot
\bottomrule
\endlastfoot
Fecha evidencia & 2026-11-06 \\
\addlinespace[3pt]
Propuesta local & 131 \\
\addlinespace[3pt]
Ejecucion local & 22 \\
\addlinespace[3pt]
Granularidad & Línea de factura \\
\addlinespace[3pt]
Filas de origen & 228265 \\
\addlinespace[3pt]
Filas de destino & 228265 \\
\addlinespace[3pt]
Diferencia de filas & 0 \\
\addlinespace[3pt]
Alcance & Evidencia histórica comprobada en esa campaña. No es un resultado nuevo obtenido al redactar este material ni el contenido garantizado de un portable limpio. \\
\addlinespace[3pt]
\end{longtable}
\endgroup

\begingroup
\setbox0=\vbox{\hsize=\linewidth
{\normalfont\normalsize\normalfont\normalsize\bfseries\hyphenpenalty=10000\exhyphenpenalty=10000 Necesidad de referencia\par}
 Diseñar un datamart cuyo grano físico sea una línea de factura. Sobre ese detalle, resumir mensualmente el importe registrado de la línea, las unidades y las facturas distintas por producto y cliente. Usar la fecha de emisión. Para el territorio, usar la provincia de la ciudad de entrega registrada en la ficha del cliente y aclarar que no representa necesariamente su ubicación histórica. Conservar el importe de origen sin asumir moneda, utilidad ni composición fiscal. Limitar los indicadores a importe total, unidades totales, facturas distintas e importe promedio por factura. Los rankings deben comparar totales por producto; el máximo de una sola línea no identifica al producto líder.

\par\vspace{6pt}\noindent\textbf{Pregunta chat}\hspace{0.5em} Resume el importe, las unidades y las facturas distintas en 2015 para California. Explica el denominador del importe por factura y qué significa la provincia usada, sin asumir moneda ni entrega realizada.

\par\vspace{6pt}\noindent\textbf{Selección completa}\hspace{0.5em}
\par
\begingroup
\fontsize{10}{12.5}\selectfont
\setlength{\tabcolsep}{5pt}
\renewcommand{\arraystretch}{1.15}
\begin{tabular}{@{}>{\RaggedRight\arraybackslash}p{0.27\linewidth}>{\RaggedRight\arraybackslash}p{0.66\linewidth}@{}}
\toprule
\textbf{Dato} & \textbf{Contenido} \\
\midrule
Importe exacto & 198043439.45 \\
\bottomrule
\end{tabular}
\endgroup
}
\Needspace{\dimexpr\ht0+\dp0+48pt\relax}
\endgroup
\paragraph{Necesidad de referencia} Diseñar un datamart cuyo grano físico sea una línea de factura. Sobre ese detalle, resumir mensualmente el importe registrado de la línea, las unidades y las facturas distintas por producto y cliente. Usar la fecha de emisión. Para el territorio, usar la provincia de la ciudad de entrega registrada en la ficha del cliente y aclarar que no representa necesariamente su ubicación histórica. Conservar el importe de origen sin asumir moneda, utilidad ni composición fiscal. Limitar los indicadores a importe total, unidades totales, facturas distintas e importe promedio por factura. Los rankings deben comparar totales por producto; el máximo de una sola línea no identifica al producto líder.

\paragraph{Pregunta chat} Resume el importe, las unidades y las facturas distintas en 2015 para California. Explica el denominador del importe por factura y qué significa la provincia usada, sin asumir moneda ni entrega realizada.

\paragraph{Selección completa}
\begingroup
\fontsize{10}{12.5}\selectfont
\setlength{\tabcolsep}{5pt}
\renewcommand{\arraystretch}{1.15}
\begin{longtable}{@{}>{\RaggedRight\arraybackslash}p{0.27\linewidth}>{\RaggedRight\arraybackslash}p{0.66\linewidth}@{}}
\toprule
\textbf{Dato} & \textbf{Contenido} \\
\midrule
\endfirsthead
\toprule
\textbf{Dato} & \textbf{Contenido} \\
\midrule
\endhead
\midrule
\multicolumn{2}{r}{\footnotesize Continúa en la página siguiente}\\
\endfoot
\bottomrule
\endlastfoot
Importe exacto & 198043439.45 \\
\addlinespace[3pt]
Moneda & moneda de origen; sin código único demostrado \\
\addlinespace[3pt]
Unidades & 8950628 \\
\addlinespace[3pt]
Facturas distintas & 70510 \\
\addlinespace[3pt]
Importe por factura exacto & 2808.7283995177986 \\
\addlinespace[3pt]
\end{longtable}
\endgroup

\begingroup
\setbox0=\vbox{\hsize=\linewidth
{\normalfont\normalsize\normalfont\normalsize\bfseries\hyphenpenalty=10000\exhyphenpenalty=10000 Selección filtrada\par}

\par
\begingroup
\fontsize{10}{12.5}\selectfont
\setlength{\tabcolsep}{5pt}
\renewcommand{\arraystretch}{1.15}
\begin{tabular}{@{}>{\RaggedRight\arraybackslash}p{0.27\linewidth}>{\RaggedRight\arraybackslash}p{0.66\linewidth}@{}}
\toprule
\textbf{Dato} & \textbf{Contenido} \\
\midrule
Año & 2015 \\
\bottomrule
\end{tabular}
\endgroup
}
\Needspace{\dimexpr\ht0+\dp0+48pt\relax}
\endgroup
\paragraph{Selección filtrada}
\begingroup
\fontsize{10}{12.5}\selectfont
\setlength{\tabcolsep}{5pt}
\renewcommand{\arraystretch}{1.15}
\begin{longtable}{@{}>{\RaggedRight\arraybackslash}p{0.27\linewidth}>{\RaggedRight\arraybackslash}p{0.66\linewidth}@{}}
\toprule
\textbf{Dato} & \textbf{Contenido} \\
\midrule
\endfirsthead
\toprule
\textbf{Dato} & \textbf{Contenido} \\
\midrule
\endhead
\midrule
\multicolumn{2}{r}{\footnotesize Continúa en la página siguiente}\\
\endfoot
\bottomrule
\endlastfoot
Año & 2015 \\
\addlinespace[3pt]
Territorio & California \\
\addlinespace[3pt]
Lineas & 3771 \\
\addlinespace[3pt]
Importe exacto & 3191769.11 \\
\addlinespace[3pt]
Unidades & 142247 \\
\addlinespace[3pt]
Facturas distintas & 1179 \\
\addlinespace[3pt]
Importe por factura exacto & 2707.1832994062765 \\
\addlinespace[3pt]
\end{longtable}
\endgroup

\chapter{Comprender la base de datos y el linaje}
\section{Dos bases de datos con trabajos diferentes}
SQL Server conserva las bases de demostración de origen. PostgreSQL conserva cuentas, configuración, metadatos, propuestas, expedientes y las tablas analíticas cargadas. Registrar una conexión no copia toda la base origen; obtener metadatos tampoco copia ventas. La aplicación lee datos autorizados cuando se confirma la carga. Cada ejecución nueva conserva su propio destino analítico.

\section{Base de datos, esquema, tabla y fila}
bi\_\allowbreak{}prototype es la base PostgreSQL auditada. Un esquema es un espacio de nombres dentro de esa base: app agrupa el estado de la plataforma; mart\_\allowbreak{}ventas\_\allowbreak{}e23 y mart\_\allowbreak{}ventas\_\allowbreak{}e24 separan los resultados de las dos ejecuciones nuevas de esta edición. Una tabla contiene columnas y filas. Un esquema llamado mart\_\allowbreak{}ventas\_\allowbreak{}e23 no es otra base SQL Server ni un catálogo de preguntas. Las ejecuciones anteriores, entre ellas e21 y e22, conservan su propio esquema y no deben confundirse con una carga nueva.

\section{Qué se comprobó realmente}
Se volvieron a leer los catálogos PostgreSQL el 25 de noviembre de 2026 en modo de solo lectura, sin filas comerciales, credenciales ni documentos JSON de las propuestas. El corte vigente incluye 93 tablas y 860 columnas: 18 tablas de app con 174 columnas, 74 tablas analíticas y public.alembic\_\allowbreak{}version. Hay 24 FK físicas, todas en app. Los nombres de las 174 columnas de las 18 tablas de app coinciden con el ORM actual. Se cotejaron 20 archivos de migración; la revisión aplicada y la única cabeza de la cadena son 20261125\_\allowbreak{}20. Estos datos describen la estructura observada; por sí solos no prueban que todas las cargas históricas sean correctas.

El inventario anterior, del 18 de noviembre, registraba 83 tablas y 782 columnas, con la revisión 20261003\_\allowbreak{}19. La diferencia corresponde a las diez tablas y 78 columnas de los nuevos destinos e23 y e24. La tabla app.parameters recibió una fila de configuración; las tablas, columnas y FK de app no cambiaron. Los diagramas de app siguen representando esas mismas 18 entidades y 24 relaciones físicas.

\section{Qué cambia con la duración configurable de la sesión}
La opción \textbf{Duración de la sesión (minutos)} se guarda como una fila de app.parameters, con la clave técnica \texttt{AUTH\_SESSION\_MINUTES}. No representa una nueva tabla de usuarios ni una nueva entidad del negocio. En este corte su valor es 60 minutos, con un intervalo permitido de 15 a 240 minutos. La migración 20261125\_\allowbreak{}20 añade el parámetro cuando falta y conserva un valor que ya estuviera configurado.

La persona administradora cambia esta duración desde \textbf{Parámetros}; no necesita editar PostgreSQL. El valor se aplica al siguiente inicio de sesión. Cambiarlo no amplía una sesión que ya está abierta ni convierte el tiempo de sesión en un contador de inactividad. Si se modifica durante una demostración, se debe cerrar la sesión y volver a entrar antes de comprobar la nueva duración. Las propuestas y los expedientes guardados permanecen en sus tablas cuando termina la sesión.

\section{Claves primarias y claves foráneas}
Una PK identifica una fila dentro de una tabla. Las tablas user\_\allowbreak{}roles, role\_\allowbreak{}permissions y menu\_\allowbreak{}permissions usan PK de dos columnas: la pareja es única. Una FK comprueba que un valor hijo exista en una tabla padre. Por ejemplo, una ejecución debe referir una propuesta existente. No se deduce una FK porque una columna termine en \_\allowbreak{}id.

\section{Cómo leer las cardinalidades}
1 junto a la tabla padre significa que cada fila hija necesita un padre; 0..1 permite que la referencia sea nula. 0..* junto a la hija indica que un padre puede tener ninguna o muchas hijas. No se deduce que todos los padres ya tengan hijos. Un índice único completo sobre la FK limita a una hija por padre. Por eso secret\_\allowbreak{}id es uno a uno opcional por cada tabla de configuración, pero la unicidad de data\_\allowbreak{}connections y de llm\_\allowbreak{}configurations no es una exclusión global entre ambas tablas.

\section{Autoría, roles y permisos}
app.users identifica a la persona de la aplicación. user\_\allowbreak{}roles la relaciona con uno o varios roles, y role\_\allowbreak{}permissions relaciona esos roles con operaciones. menu\_\allowbreak{}permissions vincula navegación con permisos. No existe una FK directa entre roles y menus. La persona puede ver una opción en React únicamente según sus permisos; FastAPI vuelve a comprobarlos antes de operar. Las FK de autoría son opcionales y usan SET NULL para preservar evidencia si se retira una cuenta; las reglas de sesión del servidor siguen siendo obligatorias.

\section{JSON y JSONB no son tablas separadas}
JSONB almacena un documento estructurado en una columna: por ejemplo la definición de un catálogo o el contrato de una propuesta. JSON almacena el documento de detalle de auditoría. Estos documentos no se descomponen automáticamente en tablas con FK. Los nombres de tablas, rutas, identificadores o esquemas guardados dentro del documento se comprueban mediante reglas del servidor. No deben dibujarse como FK físicas. El inventario registra únicamente el tipo y nombre de esas columnas, sin leer su contenido.

\section{La revisión previa no está unida por una FK directa a la propuesta}
business\_\allowbreak{}need\_\allowbreak{}reviews y bi\_\allowbreak{}proposals apuntan a metadata\_\allowbreak{}snapshots. El servidor comprueba entradas, huellas, actor y contexto al generar. No existe review\_\allowbreak{}id en bi\_\allowbreak{}proposals ni una FK que materialice la correspondencia entre ambos documentos. Esa parte de la trazabilidad es una regla de aplicación, complementaria a las restricciones SQL.

\section{Una propuesta y una ejecución no son lo mismo}
bi\_\allowbreak{}proposals guarda el diseño conservado y su decisión. etl\_\allowbreak{}executions guarda un expediente que apunta por FK a esa propuesta y a sus metadatos. Una propuesta puede fundamentar cero o varias ejecuciones. Aprobar el diseño no llena fact\_\allowbreak{}ventas; la materialización requiere la confirmación posterior y controles de conciliación.

\section{Cómo se conservan los resultados físicos}
Los destinos actuales mart\_\allowbreak{}ventas\_\allowbreak{}e23 y mart\_\allowbreak{}ventas\_\allowbreak{}e24 existen en el catálogo y contienen dim\_\allowbreak{}fecha, dim\_\allowbreak{}producto, dim\_\allowbreak{}cliente, dim\_\allowbreak{}territorio y fact\_\allowbreak{}ventas. Los históricos e21 y e22 conservan esas cinco tablas con sus propias columnas. Las dimensiones tienen PK surrogate\_\allowbreak{}key. En los cuatro destinos observados, fact\_\allowbreak{}ventas no tiene PK ni FK declaradas. Sus dim\_\allowbreak{}*\_\allowbreak{}sk se unen lógicamente con surrogate\_\allowbreak{}key, y etl\_\allowbreak{}execution\_\allowbreak{}id es una etiqueta de ejecución sin FK a app.etl\_\allowbreak{}executions. Las líneas discontinuas de sus diagramas no prometen restricciones que PostgreSQL no impone. El constructor y la conciliación deben comprobar esas correspondencias; la PK de una dimensión no demuestra por sí sola que todos los hechos tengan miembro válido.

\section{Por qué las columnas de ventas cambian entre fuentes}
Los nombres y atributos físicos dependen del contrato aprobado. AdventureWorks conserva la identidad de pedido y detalle; WideWorldImporters conserva la identidad de línea de factura. Las dos demostraciones comparten el motor general, sin exigir nombres de tablas fuente iguales. Los esquemas físicos de ejecuciones anteriores son evidencia histórica; su mera existencia no garantiza que estén publicados o disponibles en el selector analítico.

\section{Cómo interpretar medidas y conteos}
Una columna física llamada numero de pedidos o facturas distintas puede conservar el identificador documental de cada línea para que el lector calcule COUNT DISTINCT en el conjunto filtrado. No se debe sumar ese identificador ni sumar conteos de grupos que pueden repetir documentos. Para el importe promedio por documento, se divide el importe total entre los documentos distintos del filtro. Agrupar por mes no transforma la fila física en un resumen mensual.

\section{Asesoría del diseño y conversación analítica}
semantic\_\allowbreak{}advice conserva consultas sobre conceptos de una propuesta. El copiloto analítico conversa sobre agregados de una ejecución y sus filtros. La arquitectura documenta que no conserva el historial completo de chat en una tabla de mensajes: la auditoría registra huella y contexto mínimo. No debe presentarse semantic\_\allowbreak{}advice como un almacén de todos los mensajes analíticos.

\section{Qué garantiza y qué no garantiza este inventario}
La coincidencia entre ORM y catálogo acredita los nombres de columnas de app observados. Las FK físicas acreditan integridad referencial de esas relaciones SQL, no el significado comercial de cada importe o territorio. Los documentos JSON, el plan, la selección y los enlaces analíticos necesitan controles de aplicación. El manual debe enseñar a revisar tanto estructura como conciliación, contexto, moneda y denominadores, sin pedir a la persona usuaria que modifique estas tablas con SQL.

\section{Una fotografía local, no un requisito del portable}
Las 93 tablas del corte vigente describen esta instalación con resultados históricos acumulados. Las 74 tablas analíticas no son un catálogo fijo que deba venir con el portable limpio: aparecen al materializar ejecuciones. El portable limpio conserva la estructura de app y el control de migraciones; los esquemas analíticos se generan después. Los identificadores de propuestas, instantáneas y ejecuciones también se asignan en cada instalación. Por eso una réplica correcta no depende de obtener los números 135, 136, 23 o 24: depende de conservar la correspondencia entre la fuente, la propuesta que se aprobó y su propia ejecución.

\begingroup
\setbox0=\vbox{\hsize=\linewidth
{\normalfont\normalsize\RaggedRight\sffamily\bfseries\large\color{Azul}\hyphenpenalty=10000\exhyphenpenalty=10000 Roles observados y alcance de la comprobación\par}

El inventario de roles conservado del 18 de noviembre consultó únicamente sus definiciones y permisos, sin cuentas ni asignaciones de personas. En ese corte constan Administrador, Operador, Analista BI y Gerente comercial. Gerente comercial tiene lectura analítica, lectura de conexiones, consulta de metadatos y exportación analítica. Consultar metadatos permite cargar el selector de fuentes y también abrir el explorador; no concede escritura ni ETL. La nueva inspección estructural de este capítulo no vuelve a evaluar recorridos por perfil. Una lectura de permisos no acredita que se haya iniciado sesión y completado el recorrido con cada uno.

\begingroup
\setbox0=\vbox{\hsize=\linewidth
{\normalfont\normalsize\RaggedRight\sffamily\bfseries\large\color{Azul}\hyphenpenalty=10000\exhyphenpenalty=10000 Esquemas conservados en la instalación de referencia\par}

\par
\par
\begingroup
\fontsize{10}{12.5}\selectfont
\setlength{\tabcolsep}{5pt}
\renewcommand{\arraystretch}{1.15}
\begin{tabular}{@{}>{\RaggedRight\arraybackslash}p{0.27\linewidth}>{\RaggedRight\arraybackslash}p{0.66\linewidth}@{}}
\toprule
\textbf{Esquema} & \textbf{Tablas observadas} \\
\midrule
app & 18 \\
\bottomrule
\end{tabular}
\endgroup
}
\Needspace{\dimexpr\ht0+\dp0+48pt\relax}
\endgroup
\section{Esquemas conservados en la instalación de referencia}
La tabla corresponde al catálogo leído el 25 de noviembre de 2026. La presencia de un esquema conserva un resultado físico; no significa que deba aparecer como carga vigente en todos los perfiles o que todas sus cargas históricas hayan sido reevaluadas.
\begingroup
\fontsize{10}{12.5}\selectfont
\setlength{\tabcolsep}{5pt}
\renewcommand{\arraystretch}{1.15}
\begin{tabular}{@{}>{\RaggedRight\arraybackslash}p{0.27\linewidth}>{\RaggedRight\arraybackslash}p{0.66\linewidth}@{}}
\toprule
\textbf{Esquema} & \textbf{Tablas observadas} \\
\midrule
app & 18 \\
\bottomrule
\end{tabular}
\endgroup
}
\Needspace{\dimexpr\ht0+\dp0+48pt\relax}
\endgroup
\section{Roles observados y alcance de la comprobación}
El inventario de roles conservado del 18 de noviembre consultó únicamente sus definiciones y permisos, sin cuentas ni asignaciones de personas. En ese corte constan Administrador, Operador, Analista BI y Gerente comercial. Gerente comercial tiene lectura analítica, lectura de conexiones, consulta de metadatos y exportación analítica. Consultar metadatos permite cargar el selector de fuentes y también abrir el explorador; no concede escritura ni ETL. La nueva inspección estructural de este capítulo no vuelve a evaluar recorridos por perfil. Una lectura de permisos no acredita que se haya iniciado sesión y completado el recorrido con cada uno.

\begingroup
\setbox0=\vbox{\hsize=\linewidth
{\normalfont\normalsize\RaggedRight\sffamily\bfseries\large\color{Azul}\hyphenpenalty=10000\exhyphenpenalty=10000 Esquemas conservados en la instalación de referencia\par}

\par
\begingroup
\fontsize{10}{12.5}\selectfont
\setlength{\tabcolsep}{5pt}
\renewcommand{\arraystretch}{1.15}
\begin{tabular}{@{}>{\RaggedRight\arraybackslash}p{0.27\linewidth}>{\RaggedRight\arraybackslash}p{0.66\linewidth}@{}}
\toprule
\textbf{Esquema} & \textbf{Tablas observadas} \\
\midrule
app & 18 \\
\bottomrule
\end{tabular}
\endgroup
}
\Needspace{\dimexpr\ht0+\dp0+48pt\relax}
\endgroup
\section{Esquemas conservados en la instalación de referencia}
La tabla corresponde al catálogo leído el 25 de noviembre de 2026. La presencia de un esquema conserva un resultado físico; no significa que deba aparecer como carga vigente en todos los perfiles o que todas sus cargas históricas hayan sido reevaluadas.
\begingroup
\fontsize{10}{12.5}\selectfont
\setlength{\tabcolsep}{5pt}
\renewcommand{\arraystretch}{1.15}
\begin{longtable}{@{}>{\RaggedRight\arraybackslash}p{0.27\linewidth}>{\RaggedRight\arraybackslash}p{0.66\linewidth}@{}}
\toprule
\textbf{Esquema} & \textbf{Tablas observadas} \\
\midrule
\endfirsthead
\toprule
\textbf{Esquema} & \textbf{Tablas observadas} \\
\midrule
\endhead
\midrule
\multicolumn{2}{r}{\footnotesize Continúa en la página siguiente}\\
\endfoot
\bottomrule
\endlastfoot
app & 18 \\
\addlinespace[3pt]
mart & 0 \\
\addlinespace[3pt]
mart\_\allowbreak{}ventas & 5 \\
\addlinespace[3pt]
mart\_\allowbreak{}ventas\_\allowbreak{}e11 & 5 \\
\addlinespace[3pt]
mart\_\allowbreak{}ventas\_\allowbreak{}e12 & 5 \\
\addlinespace[3pt]
mart\_\allowbreak{}ventas\_\allowbreak{}e13 & 5 \\
\addlinespace[3pt]
mart\_\allowbreak{}ventas\_\allowbreak{}e14 & 5 \\
\addlinespace[3pt]
mart\_\allowbreak{}ventas\_\allowbreak{}e15 & 5 \\
\addlinespace[3pt]
mart\_\allowbreak{}ventas\_\allowbreak{}e16 & 5 \\
\addlinespace[3pt]
mart\_\allowbreak{}ventas\_\allowbreak{}e17 & 4 \\
\addlinespace[3pt]
mart\_\allowbreak{}ventas\_\allowbreak{}e18 & 5 \\
\addlinespace[3pt]
mart\_\allowbreak{}ventas\_\allowbreak{}e19 & 5 \\
\addlinespace[3pt]
mart\_\allowbreak{}ventas\_\allowbreak{}e20 & 5 \\
\addlinespace[3pt]
mart\_\allowbreak{}ventas\_\allowbreak{}e21 & 5 \\
\addlinespace[3pt]
mart\_\allowbreak{}ventas\_\allowbreak{}e22 & 5 \\
\addlinespace[3pt]
mart\_\allowbreak{}ventas\_\allowbreak{}e23 & 5 \\
\addlinespace[3pt]
mart\_\allowbreak{}ventas\_\allowbreak{}e24 & 5 \\
\addlinespace[3pt]
public & 1 \\
\addlinespace[3pt]
\end{longtable}
\endgroup

\begingroup
\setbox0=\vbox{\hsize=\linewidth
{\normalfont\normalsize\RaggedRight\sffamily\bfseries\large\color{Azul}\hyphenpenalty=10000\exhyphenpenalty=10000 Relación comprobada entre los casos de demostración\par}

\par
\begingroup
\fontsize{10}{12.5}\selectfont
\setlength{\tabcolsep}{5pt}
\renewcommand{\arraystretch}{1.15}
\begin{tabular}{@{}>{\RaggedRight\arraybackslash}p{0.22\linewidth}>{\RaggedRight\arraybackslash}p{0.3\linewidth}>{\RaggedRight\arraybackslash}p{0.39\linewidth}@{}}
\toprule
\textbf{Fuente} & \textbf{Diseño y metadatos} & \textbf{Ejecución} \\
\midrule
AdventureWorks2022 & Propuesta 136 aprobada; instantánea 2. & Expediente 24: carga completada. Vigente en esta edición. \\
\bottomrule
\end{tabular}
\endgroup
}
\Needspace{\dimexpr\ht0+\dp0+48pt\relax}
\endgroup
\section{Relación comprobada entre los casos de demostración}
La lectura técnica confirma que cada expediente apunta a su propuesta y a la instantánea indicada. La edición actual usa los casos 135/23 y 136/24 para la nueva demostración. Se conservan 129/21 y 131/22 como referencias históricas del manual del 18 de noviembre; no se ha sustituido su contenido por el de una ejecución posterior. En otra instalación se debe reconocer la fuente y el diseño, porque sus números locales pueden ser diferentes.
\begingroup
\fontsize{10}{12.5}\selectfont
\setlength{\tabcolsep}{5pt}
\renewcommand{\arraystretch}{1.15}
\begin{longtable}{@{}>{\RaggedRight\arraybackslash}p{0.22\linewidth}>{\RaggedRight\arraybackslash}p{0.3\linewidth}>{\RaggedRight\arraybackslash}p{0.39\linewidth}@{}}
\toprule
\textbf{Fuente} & \textbf{Diseño y metadatos} & \textbf{Ejecución} \\
\midrule
\endfirsthead
\toprule
\textbf{Fuente} & \textbf{Diseño y metadatos} & \textbf{Ejecución} \\
\midrule
\endhead
\midrule
\multicolumn{3}{r}{\footnotesize Continúa en la página siguiente}\\
\endfoot
\bottomrule
\endlastfoot
AdventureWorks2022 & Propuesta 136 aprobada; instantánea 2. & Expediente 24: carga completada. Vigente en esta edición. \\
\addlinespace[3pt]
WideWorldImporters & Propuesta 135 aprobada; instantánea 3. & Expediente 23: carga completada. Vigente en esta edición. \\
\addlinespace[3pt]
AdventureWorks2022 & Propuesta 129 aprobada; instantánea 2. & Expediente 21: carga completada. Referencia histórica del 18 de noviembre. \\
\addlinespace[3pt]
WideWorldImporters & Propuesta 131 aprobada; instantánea 3. & Expediente 22: carga completada. Referencia histórica del 18 de noviembre. \\
\addlinespace[3pt]
\end{longtable}
\endgroup

\diagramasistema{recursos/diagramas/ER_00_recorrido_identidades.pdf}{Identidades principales que conectan fuente, diseño y ejecución}{Vista inicial de las identidades principales; no incluye todas las FK. El inventario y los cuatro esquemas de área contienen las 18 entidades de app y las 24 relaciones físicas, sin deducir relaciones por parecido de nombres.}
\diagramasistema{recursos/diagramas/ER_01_seguridad.pdf}{Modelo físico de seguridad y navegación}{Se muestran las ocho tablas de seguridad y siete relaciones físicas. Campos principales; el diccionario enumera todas las columnas. La auditoría conserva actor opcional y una etiqueta histórica separada.}
\diagramasistema{recursos/diagramas/ER_02_configuracion_metadatos.pdf}{Modelo físico de configuración y metadatos}{Seis tablas de configuración y metadatos. Cuatro relaciones físicas; las dos FK de autoría aparecen en el esquema separado de responsables. Catálogo único por conexión y dominio; instantánea única por conexión, versión de contrato y huella.}
\diagramasistema{recursos/diagramas/ER_03_propuestas_ejecuciones.pdf}{Modelo físico de necesidades, propuestas y ejecuciones}{Cuatro entidades operativas y referencia a metadatos; seis FK físicas, incluida la relación de versiones derivadas de bi\_\allowbreak{}proposals. Las revisiones previas se comprueban con contexto y huellas por lógica del servidor; no existe una FK review\_\allowbreak{}id en bi\_\allowbreak{}proposals.}
\diagramasistema{recursos/diagramas/ER_04_autoria.pdf}{Relaciones físicas con la persona que realiza cada acción}{Siete FK de autoría hacia users. La opcionalidad se obtiene del catálogo real: la eliminación de una cuenta preserva el expediente anulando la referencia. Las reglas del servidor y las etiquetas históricas complementan la FK.}
\diagramasistema{recursos/diagramas/ER_05_datamart_e21.pdf}{Referencia histórica: estructura física del datamart, ejecución 21}{Diagrama del inventario del 18 de noviembre; no representa la nueva ejecución 24. Conserva todas las columnas de las cinco tablas físicas observadas entonces. Los nombres y atributos corresponden a esta ejecución, no son un modelo fijo universal. Las líneas discontinuas representan la unión que efectúa el motor entre dim\_\allowbreak{}*\_\allowbreak{}sk y surrogate\_\allowbreak{}key; la base no las impone como FK. No se atribuyen cardinalidades garantizadas a esas líneas.}
\diagramasistema{recursos/diagramas/ER_05_datamart_e22.pdf}{Referencia histórica: estructura física del datamart, ejecución 22}{Diagrama del inventario del 18 de noviembre; no representa la nueva ejecución 23. Conserva todas las columnas de las cinco tablas físicas observadas entonces. Los nombres y atributos corresponden a esta ejecución, no son un modelo fijo universal. Las líneas discontinuas representan la unión que efectúa el motor entre dim\_\allowbreak{}*\_\allowbreak{}sk y surrogate\_\allowbreak{}key; la base no las impone como FK. No se atribuyen cardinalidades garantizadas a esas líneas.}
\section{Cómo reconocer los destinos nuevos sin mezclar su información}
El destino de WideWorldImporters, ejecución 23, tiene 36 columnas entre sus cinco tablas. Su tabla de hechos conserva la identidad de la línea de factura, el importe registrado, el identificador usado para facturas distintas, las unidades y las claves de las dimensiones. El destino de AdventureWorks, ejecución 24, tiene 42 columnas entre sus cinco tablas. Su tabla de hechos conserva la identidad de pedido y detalle, el importe registrado, el identificador usado para pedidos distintos, las unidades y las claves de las dimensiones. Que ambos destinos tengan una tabla llamada fact\_\allowbreak{}ventas no significa que sus filas o atributos deban combinarse.

Para revisar el linaje desde la aplicación: primero se elige la fuente; después se comprueba la propuesta del expediente; por último se abre el resultado de esa ejecución en Analítica de ventas. Antes de interpretar un valor, se verifican también el año, el territorio y los demás filtros. La fecha de un pedido, la fecha de una factura y la provincia de entrega del cliente son decisiones de negocio diferentes: se conserva la definición del contrato correspondiente. No se debe identificar un territorio solo por el parecido de su nombre.

Las reglas nuevas de composición no reescriben las propuestas históricas. Una receta que explica un conteo de documentos distintos sigue siendo una receta de conteo, aunque la tabla física guarde el identificador documental repetido en varias líneas. Revisar un diseño o un resultado histórico no debe retirar su aprobación ni sustituirlo por otro diseño. Cuando se necesita otro alcance, se genera y aprueba una propuesta nueva y se conserva su expediente separado.
\begingroup
\setbox0=\vbox{\hsize=\linewidth
{\normalfont\normalsize\RaggedRight\sffamily\bfseries\large\color{Azul}\hyphenpenalty=10000\exhyphenpenalty=10000 Diccionario de las tablas de la aplicación\par}

El nombre de cada tabla y columna coincide con el inventario observado. Estas fichas ayudan a comprender lo que guarda la plataforma; las tareas de operación se realizan desde la interfaz, sin modificar estas tablas con SQL.

\begingroup
\setbox0=\vbox{\hsize=\linewidth
{\normalfont\normalsize\sffamily\bfseries\normalsize\color{Teal}\hyphenpenalty=10000\exhyphenpenalty=10000 app.analysis\_\allowbreak{}catalogs\par}

\par\vspace{6pt}\noindent\textbf{Área}\hspace{0.5em} Configuración y metadatos

\par\vspace{6pt}\noindent\textbf{Qué guarda}\hspace{0.5em} Guarda orientación funcional por conexión y dominio, con versión.

\par\vspace{6pt}\noindent\textbf{Cómo se utiliza}\hspace{0.5em} Catálogo analítico: preguntas de negocio y periodicidades. No define un modelo dimensional físico fijo. La combinación conexión/dominio es única.

\par
\par
\begingroup
\fontsize{10}{12.5}\selectfont
\setlength{\tabcolsep}{5pt}
\renewcommand{\arraystretch}{1.15}
\begin{tabular}{@{}>{\RaggedRight\arraybackslash}p{0.23\linewidth}>{\RaggedRight\arraybackslash}p{0.23\linewidth}>{\RaggedRight\arraybackslash}p{0.45\linewidth}@{}}
\toprule
\textbf{Columna} & \textbf{Tipo y condición} & \textbf{Qué representa} \\
\midrule
\texttt{\seqsplit{id}} & integer\par\smallskip Obligatorio. Clave primaria. Con valor predeterminado. & Identificador interno de esta entidad; su PK distingue cada registro. \\
\bottomrule
\end{tabular}
\endgroup
}
\Needspace{\dimexpr\ht0+\dp0+48pt\relax}
\endgroup
\subsection{app.analysis\_\allowbreak{}catalogs}
\paragraph{Área} Configuración y metadatos

\paragraph{Qué guarda} Guarda orientación funcional por conexión y dominio, con versión.

\paragraph{Cómo se utiliza} Catálogo analítico: preguntas de negocio y periodicidades. No define un modelo dimensional físico fijo. La combinación conexión/dominio es única.

\begingroup
\fontsize{10}{12.5}\selectfont
\setlength{\tabcolsep}{5pt}
\renewcommand{\arraystretch}{1.15}
\begin{tabular}{@{}>{\RaggedRight\arraybackslash}p{0.23\linewidth}>{\RaggedRight\arraybackslash}p{0.23\linewidth}>{\RaggedRight\arraybackslash}p{0.45\linewidth}@{}}
\toprule
\textbf{Columna} & \textbf{Tipo y condición} & \textbf{Qué representa} \\
\midrule
\texttt{\seqsplit{id}} & integer\par\smallskip Obligatorio. Clave primaria. Con valor predeterminado. & Identificador interno de esta entidad; su PK distingue cada registro. \\
\bottomrule
\end{tabular}
\endgroup
}
\Needspace{\dimexpr\ht0+\dp0+48pt\relax}
\endgroup
\section{Diccionario de las tablas de la aplicación}
El nombre de cada tabla y columna coincide con el inventario observado. Estas fichas ayudan a comprender lo que guarda la plataforma; las tareas de operación se realizan desde la interfaz, sin modificar estas tablas con SQL.

\begingroup
\setbox0=\vbox{\hsize=\linewidth
{\normalfont\normalsize\sffamily\bfseries\normalsize\color{Teal}\hyphenpenalty=10000\exhyphenpenalty=10000 app.analysis\_\allowbreak{}catalogs\par}

\par\vspace{6pt}\noindent\textbf{Área}\hspace{0.5em} Configuración y metadatos

\par\vspace{6pt}\noindent\textbf{Qué guarda}\hspace{0.5em} Guarda orientación funcional por conexión y dominio, con versión.

\par\vspace{6pt}\noindent\textbf{Cómo se utiliza}\hspace{0.5em} Catálogo analítico: preguntas de negocio y periodicidades. No define un modelo dimensional físico fijo. La combinación conexión/dominio es única.

\par
\begingroup
\fontsize{10}{12.5}\selectfont
\setlength{\tabcolsep}{5pt}
\renewcommand{\arraystretch}{1.15}
\begin{tabular}{@{}>{\RaggedRight\arraybackslash}p{0.23\linewidth}>{\RaggedRight\arraybackslash}p{0.23\linewidth}>{\RaggedRight\arraybackslash}p{0.45\linewidth}@{}}
\toprule
\textbf{Columna} & \textbf{Tipo y condición} & \textbf{Qué representa} \\
\midrule
\texttt{\seqsplit{id}} & integer\par\smallskip Obligatorio. Clave primaria. Con valor predeterminado. & Identificador interno de esta entidad; su PK distingue cada registro. \\
\bottomrule
\end{tabular}
\endgroup
}
\Needspace{\dimexpr\ht0+\dp0+48pt\relax}
\endgroup
\subsection{app.analysis\_\allowbreak{}catalogs}
\paragraph{Área} Configuración y metadatos

\paragraph{Qué guarda} Guarda orientación funcional por conexión y dominio, con versión.

\paragraph{Cómo se utiliza} Catálogo analítico: preguntas de negocio y periodicidades. No define un modelo dimensional físico fijo. La combinación conexión/dominio es única.

\begingroup
\fontsize{10}{12.5}\selectfont
\setlength{\tabcolsep}{5pt}
\renewcommand{\arraystretch}{1.15}
\begin{longtable}{@{}>{\RaggedRight\arraybackslash}p{0.23\linewidth}>{\RaggedRight\arraybackslash}p{0.23\linewidth}>{\RaggedRight\arraybackslash}p{0.45\linewidth}@{}}
\toprule
\textbf{Columna} & \textbf{Tipo y condición} & \textbf{Qué representa} \\
\midrule
\endfirsthead
\toprule
\textbf{Columna} & \textbf{Tipo y condición} & \textbf{Qué representa} \\
\midrule
\endhead
\midrule
\multicolumn{3}{r}{\footnotesize Continúa en la página siguiente}\\
\endfoot
\bottomrule
\endlastfoot
\texttt{\seqsplit{id}} & integer\par\smallskip Obligatorio. Clave primaria. Con valor predeterminado. & Identificador interno de esta entidad; su PK distingue cada registro. \\
\addlinespace[3pt]
\texttt{\seqsplit{data\_connection\_id}} & integer\par\smallskip Obligatorio. & FK a app.data\_\allowbreak{}connections.id; preserva el contexto de fuente. \\
\addlinespace[3pt]
\texttt{\seqsplit{domain\_code}} & character varying(40)\par\smallskip Obligatorio. & Dominio analítico; ventas es el dominio implementado. \\
\addlinespace[3pt]
\texttt{\seqsplit{version}} & integer\par\smallskip Obligatorio. Con valor predeterminado. & Versión de configuración del catálogo; no representa un expediente de ejecución. \\
\addlinespace[3pt]
\texttt{\seqsplit{configuration\_document}} & jsonb\par\smallskip Obligatorio. & JSONB de preguntas de negocio y periodicidades administradas para esta conexión y dominio. \\
\addlinespace[3pt]
\texttt{\seqsplit{updated\_by\_user\_id}} & integer\par\smallskip Admite nulos. & FK opcional al usuario que actualizó el catálogo. \\
\addlinespace[3pt]
\texttt{\seqsplit{updated\_by\_label}} & character varying(320)\par\smallskip Obligatorio. & Etiqueta histórica de la persona que actualizó el catálogo. \\
\addlinespace[3pt]
\texttt{\seqsplit{updated\_at}} & timestamp with time zone\par\smallskip Obligatorio. Con valor predeterminado. & Fecha de actualización del registro; su opcionalidad física consta en el inventario. \\
\addlinespace[3pt]
\end{longtable}
\endgroup

\begingroup
\setbox0=\vbox{\hsize=\linewidth
{\normalfont\normalsize\bfseries Restricciones físicas observadas\par}
\begin{itemize}
\item Clave foránea: data\_\allowbreak{}connection\_\allowbreak{}id hacia app.data\_\allowbreak{}connections (id); al eliminar: cascada SQL.
\end{itemize}
}
\Needspace{\dimexpr\ht0+\dp0+24pt\relax}
\endgroup
\begingroup
\setbox0=\vbox{\hsize=\linewidth
{\normalfont\normalsize\bfseries Restricciones físicas observadas\par}
\begin{itemize}
\item Clave foránea: data\_\allowbreak{}connection\_\allowbreak{}id hacia app.data\_\allowbreak{}connections (id); al eliminar: cascada SQL.
\end{itemize}
}
\Needspace{\dimexpr\ht0+\dp0+24pt\relax}
\endgroup
\paragraph{Restricciones físicas observadas} \begin{itemize}
\item Clave foránea: data\_\allowbreak{}connection\_\allowbreak{}id hacia app.data\_\allowbreak{}connections (id); al eliminar: cascada SQL.
\item Clave foránea: updated\_\allowbreak{}by\_\allowbreak{}user\_\allowbreak{}id hacia app.users (id); al eliminar: referencia nula.
\item Clave primaria: id.
\item Unicidad: data\_\allowbreak{}connection\_\allowbreak{}id, domain\_\allowbreak{}code.
\end{itemize}

\begingroup
\setbox0=\vbox{\hsize=\linewidth
{\normalfont\normalsize\bfseries Índices únicos\par}
\begin{itemize}
\item id; índice completo.
\end{itemize}
}
\Needspace{\dimexpr\ht0+\dp0+24pt\relax}
\endgroup
\begingroup
\setbox0=\vbox{\hsize=\linewidth
{\normalfont\normalsize\bfseries Índices únicos\par}
\begin{itemize}
\item id; índice completo.
\end{itemize}
}
\Needspace{\dimexpr\ht0+\dp0+24pt\relax}
\endgroup
\paragraph{Índices únicos} \begin{itemize}
\item id; índice completo.
\item data\_\allowbreak{}connection\_\allowbreak{}id, domain\_\allowbreak{}code; índice completo.
\end{itemize}

\begingroup
\setbox0=\vbox{\hsize=\linewidth
{\normalfont\normalsize\sffamily\bfseries\normalsize\color{Teal}\hyphenpenalty=10000\exhyphenpenalty=10000 app.audit\_\allowbreak{}events\par}

\par\vspace{6pt}\noindent\textbf{Área}\hspace{0.5em} Seguridad

\par\vspace{6pt}\noindent\textbf{Qué guarda}\hspace{0.5em} Registra quién realizó una acción y sobre qué recurso, conservando contexto técnico.

\par\vspace{6pt}\noindent\textbf{Cómo se utiliza}\hspace{0.5em} Auditoría de operaciones. La columna detail es JSON, no JSONB; actor\_\allowbreak{}user\_\allowbreak{}id es la FK de usuario, pero resource\_\allowbreak{}id es una referencia textual, no una FK universal.

\par
\begingroup
\fontsize{10}{12.5}\selectfont
\setlength{\tabcolsep}{5pt}
\renewcommand{\arraystretch}{1.15}
\begin{tabular}{@{}>{\RaggedRight\arraybackslash}p{0.23\linewidth}>{\RaggedRight\arraybackslash}p{0.23\linewidth}>{\RaggedRight\arraybackslash}p{0.45\linewidth}@{}}
\toprule
\textbf{Columna} & \textbf{Tipo y condición} & \textbf{Qué representa} \\
\midrule
\texttt{\seqsplit{id}} & integer\par\smallskip Obligatorio. Clave primaria. Con valor predeterminado. & Identificador interno de esta entidad; su PK distingue cada registro. \\
\bottomrule
\end{tabular}
\endgroup
}
\Needspace{\dimexpr\ht0+\dp0+48pt\relax}
\endgroup
\subsection{app.audit\_\allowbreak{}events}
\paragraph{Área} Seguridad

\paragraph{Qué guarda} Registra quién realizó una acción y sobre qué recurso, conservando contexto técnico.

\paragraph{Cómo se utiliza} Auditoría de operaciones. La columna detail es JSON, no JSONB; actor\_\allowbreak{}user\_\allowbreak{}id es la FK de usuario, pero resource\_\allowbreak{}id es una referencia textual, no una FK universal.

\begingroup
\fontsize{10}{12.5}\selectfont
\setlength{\tabcolsep}{5pt}
\renewcommand{\arraystretch}{1.15}
\begin{longtable}{@{}>{\RaggedRight\arraybackslash}p{0.23\linewidth}>{\RaggedRight\arraybackslash}p{0.23\linewidth}>{\RaggedRight\arraybackslash}p{0.45\linewidth}@{}}
\toprule
\textbf{Columna} & \textbf{Tipo y condición} & \textbf{Qué representa} \\
\midrule
\endfirsthead
\toprule
\textbf{Columna} & \textbf{Tipo y condición} & \textbf{Qué representa} \\
\midrule
\endhead
\midrule
\multicolumn{3}{r}{\footnotesize Continúa en la página siguiente}\\
\endfoot
\bottomrule
\endlastfoot
\texttt{\seqsplit{id}} & integer\par\smallskip Obligatorio. Clave primaria. Con valor predeterminado. & Identificador interno de esta entidad; su PK distingue cada registro. \\
\addlinespace[3pt]
\texttt{\seqsplit{actor\_user\_id}} & integer\par\smallskip Admite nulos. & FK opcional al usuario que realizó la acción; ON DELETE SET NULL conserva el evento si se retira la cuenta. \\
\addlinespace[3pt]
\texttt{\seqsplit{action}} & character varying(120)\par\smallskip Obligatorio. & Código de la operación registrada en auditoría. \\
\addlinespace[3pt]
\texttt{\seqsplit{resource\_type}} & character varying(100)\par\smallskip Obligatorio. & Clase de recurso afectado por la operación. \\
\addlinespace[3pt]
\texttt{\seqsplit{resource\_id}} & character varying(100)\par\smallskip Admite nulos. & Identidad del recurso guardada como texto; PostgreSQL no la impone como FK a todas las tablas. \\
\addlinespace[3pt]
\texttt{\seqsplit{detail}} & json\par\smallskip Obligatorio. Con valor predeterminado. & Documento JSON con detalle de auditoría; no se leyó su contenido. \\
\addlinespace[3pt]
\texttt{\seqsplit{created\_at}} & timestamp with time zone\par\smallskip Admite nulos. Con valor predeterminado. & Fecha de creación del registro; su opcionalidad física consta en el inventario. \\
\addlinespace[3pt]
\texttt{\seqsplit{actor\_label}} & character varying(320)\par\smallskip Admite nulos. & Etiqueta histórica del actor. Es texto de evidencia y no otra FK. \\
\addlinespace[3pt]
\end{longtable}
\endgroup

\begingroup
\setbox0=\vbox{\hsize=\linewidth
{\normalfont\normalsize\bfseries Restricciones físicas observadas\par}
\begin{itemize}
\item Clave foránea: actor\_\allowbreak{}user\_\allowbreak{}id hacia app.users (id); al eliminar: referencia nula.
\end{itemize}
}
\Needspace{\dimexpr\ht0+\dp0+24pt\relax}
\endgroup
\begingroup
\setbox0=\vbox{\hsize=\linewidth
{\normalfont\normalsize\bfseries Restricciones físicas observadas\par}
\begin{itemize}
\item Clave foránea: actor\_\allowbreak{}user\_\allowbreak{}id hacia app.users (id); al eliminar: referencia nula.
\end{itemize}
}
\Needspace{\dimexpr\ht0+\dp0+24pt\relax}
\endgroup
\paragraph{Restricciones físicas observadas} \begin{itemize}
\item Clave foránea: actor\_\allowbreak{}user\_\allowbreak{}id hacia app.users (id); al eliminar: referencia nula.
\item Clave primaria: id.
\end{itemize}

\begingroup
\setbox0=\vbox{\hsize=\linewidth
{\normalfont\normalsize\bfseries Índices únicos\par}
\begin{itemize}
\item id; índice completo.
\end{itemize}
}
\Needspace{\dimexpr\ht0+\dp0+24pt\relax}
\endgroup
\begingroup
\setbox0=\vbox{\hsize=\linewidth
{\normalfont\normalsize\bfseries Índices únicos\par}
\begin{itemize}
\item id; índice completo.
\end{itemize}
}
\Needspace{\dimexpr\ht0+\dp0+24pt\relax}
\endgroup
\paragraph{Índices únicos} \begin{itemize}
\item id; índice completo.
\end{itemize}

\begingroup
\setbox0=\vbox{\hsize=\linewidth
{\normalfont\normalsize\sffamily\bfseries\normalsize\color{Teal}\hyphenpenalty=10000\exhyphenpenalty=10000 app.bi\_\allowbreak{}proposals\par}

\par\vspace{6pt}\noindent\textbf{Área}\hspace{0.5em} Propuestas y ejecución

\par\vspace{6pt}\noindent\textbf{Qué guarda}\hspace{0.5em} Conserva versiones del diseño BI, sus documentos, proveedor, validación y revisión humana.

\par\vspace{6pt}\noindent\textbf{Cómo se utiliza}\hspace{0.5em} Asistente de datamart: conceptos, propuesta, personalización y revisión. Aprobar no ejecuta la carga. Una versión derivada enlaza con su origen mediante source\_\allowbreak{}proposal\_\allowbreak{}id.

\par
\begingroup
\fontsize{10}{12.5}\selectfont
\setlength{\tabcolsep}{5pt}
\renewcommand{\arraystretch}{1.15}
\begin{tabular}{@{}>{\RaggedRight\arraybackslash}p{0.23\linewidth}>{\RaggedRight\arraybackslash}p{0.23\linewidth}>{\RaggedRight\arraybackslash}p{0.45\linewidth}@{}}
\toprule
\textbf{Columna} & \textbf{Tipo y condición} & \textbf{Qué representa} \\
\midrule
\texttt{\seqsplit{id}} & integer\par\smallskip Obligatorio. Clave primaria. Con valor predeterminado. & Identificador interno de esta entidad; su PK distingue cada registro. \\
\bottomrule
\end{tabular}
\endgroup
}
\Needspace{\dimexpr\ht0+\dp0+48pt\relax}
\endgroup
\subsection{app.bi\_\allowbreak{}proposals}
\paragraph{Área} Propuestas y ejecución

\paragraph{Qué guarda} Conserva versiones del diseño BI, sus documentos, proveedor, validación y revisión humana.

\paragraph{Cómo se utiliza} Asistente de datamart: conceptos, propuesta, personalización y revisión. Aprobar no ejecuta la carga. Una versión derivada enlaza con su origen mediante source\_\allowbreak{}proposal\_\allowbreak{}id.

\begingroup
\fontsize{10}{12.5}\selectfont
\setlength{\tabcolsep}{5pt}
\renewcommand{\arraystretch}{1.15}
\begin{longtable}{@{}>{\RaggedRight\arraybackslash}p{0.23\linewidth}>{\RaggedRight\arraybackslash}p{0.23\linewidth}>{\RaggedRight\arraybackslash}p{0.45\linewidth}@{}}
\toprule
\textbf{Columna} & \textbf{Tipo y condición} & \textbf{Qué representa} \\
\midrule
\endfirsthead
\toprule
\textbf{Columna} & \textbf{Tipo y condición} & \textbf{Qué representa} \\
\midrule
\endhead
\midrule
\multicolumn{3}{r}{\footnotesize Continúa en la página siguiente}\\
\endfoot
\bottomrule
\endlastfoot
\texttt{\seqsplit{id}} & integer\par\smallskip Obligatorio. Clave primaria. Con valor predeterminado. & Identificador interno de esta entidad; su PK distingue cada registro. \\
\addlinespace[3pt]
\texttt{\seqsplit{metadata\_snapshot\_id}} & integer\par\smallskip Obligatorio. & FK a app.metadata\_\allowbreak{}snapshots.id; enlaza el artefacto a una estructura de origen conservada. \\
\addlinespace[3pt]
\texttt{\seqsplit{business\_goal}} & text\par\smallskip Obligatorio. & Necesidad comercial redactada o adoptada por la persona analista. \\
\addlinespace[3pt]
\texttt{\seqsplit{business\_questions}} & jsonb\par\smallskip Obligatorio. & JSONB de las preguntas seleccionadas para esa necesidad. \\
\addlinespace[3pt]
\texttt{\seqsplit{requested\_dimensions}} & jsonb\par\smallskip Obligatorio. & Campo JSONB conservado en el contrato histórico; no implica que el catálogo actual fije dimensiones predeterminadas. \\
\addlinespace[3pt]
\texttt{\seqsplit{periodicity}} & character varying(20)\par\smallskip Obligatorio. & Agrupación temporal solicitada; no cambia por sí sola el grano físico de las filas. \\
\addlinespace[3pt]
\texttt{\seqsplit{scope\_document}} & jsonb\par\smallskip Obligatorio. & JSONB del alcance y contexto verificados para el diseño. \\
\addlinespace[3pt]
\texttt{\seqsplit{semantic\_map\_document}} & jsonb\par\smallskip Obligatorio. & JSONB de conceptos, explicaciones y correspondencias con la estructura. \\
\addlinespace[3pt]
\texttt{\seqsplit{status}} & character varying(40)\par\smallskip Obligatorio. & Estado del ciclo de vida del artefacto. Una propuesta aprobada no equivale a una ejecución terminada. \\
\addlinespace[3pt]
\texttt{\seqsplit{input\_hash}} & character varying(64)\par\smallskip Obligatorio. & Huella de entradas y contexto de la solicitud; no es el hash de contraseña de una persona. \\
\addlinespace[3pt]
\texttt{\seqsplit{prompt\_version}} & character varying(30)\par\smallskip Obligatorio. & Versión de las instrucciones utilizadas para preparar el artefacto. \\
\addlinespace[3pt]
\texttt{\seqsplit{contract\_version}} & integer\par\smallskip Obligatorio. & Versión del contrato estructural guardado; no es el número de propuesta ni de ejecución. \\
\addlinespace[3pt]
\texttt{\seqsplit{provider\_kind}} & character varying(40)\par\smallskip Obligatorio. & Tipo de proveedor de IA. En los expedientes conserva la procedencia como texto, sin FK a la configuración que después puede cambiar. \\
\addlinespace[3pt]
\texttt{\seqsplit{model\_id}} & character varying(160)\par\smallskip Obligatorio. & Identificador del modelo utilizado o configurado, distinto de una API key. \\
\addlinespace[3pt]
\texttt{\seqsplit{proposal\_document}} & jsonb\par\smallskip Obligatorio. & JSONB del diseño BI declarativo: hechos, dimensiones, medidas, indicadores y plan propuesto. No es el datamart cargado. \\
\addlinespace[3pt]
\texttt{\seqsplit{validation\_document}} & jsonb\par\smallskip Obligatorio. & JSONB de comprobaciones, errores o advertencias del artefacto o ejecución. \\
\addlinespace[3pt]
\texttt{\seqsplit{review\_comment}} & text\par\smallskip Admite nulos. & Justificación o comentario de la revisión humana. \\
\addlinespace[3pt]
\texttt{\seqsplit{warnings\_confirmed}} & boolean\par\smallskip Obligatorio. Con valor predeterminado. & Indica si se confirmaron las advertencias sujetas a decisión humana. \\
\addlinespace[3pt]
\texttt{\seqsplit{created\_by\_user\_id}} & integer\par\smallskip Admite nulos. & FK opcional a la persona que creó el artefacto o expediente. \\
\addlinespace[3pt]
\texttt{\seqsplit{created\_by\_label}} & character varying(320)\par\smallskip Obligatorio. & Etiqueta histórica de la persona creadora. \\
\addlinespace[3pt]
\texttt{\seqsplit{reviewed\_by\_user\_id}} & integer\par\smallskip Admite nulos. & FK opcional a la persona que revisó la propuesta; puede diferir de quien la creó. \\
\addlinespace[3pt]
\texttt{\seqsplit{reviewed\_by\_label}} & character varying(320)\par\smallskip Admite nulos. & Etiqueta histórica de la persona revisora. \\
\addlinespace[3pt]
\texttt{\seqsplit{created\_at}} & timestamp with time zone\par\smallskip Obligatorio. Con valor predeterminado. & Fecha de creación del registro; su opcionalidad física consta en el inventario. \\
\addlinespace[3pt]
\texttt{\seqsplit{reviewed\_at}} & timestamp with time zone\par\smallskip Admite nulos. & Fecha de la decisión de revisión. \\
\addlinespace[3pt]
\texttt{\seqsplit{domain\_code}} & character varying(40)\par\smallskip Obligatorio. & Dominio analítico; ventas es el dominio implementado. \\
\addlinespace[3pt]
\texttt{\seqsplit{source\_proposal\_id}} & integer\par\smallskip Admite nulos. & FK opcional a otra fila de app.bi\_\allowbreak{}proposals; registra el origen de una versión derivada. \\
\addlinespace[3pt]
\end{longtable}
\endgroup

\begingroup
\setbox0=\vbox{\hsize=\linewidth
{\normalfont\normalsize\bfseries Restricciones físicas observadas\par}
\begin{itemize}
\item Clave foránea: created\_\allowbreak{}by\_\allowbreak{}user\_\allowbreak{}id hacia app.users (id); al eliminar: referencia nula.
\end{itemize}
}
\Needspace{\dimexpr\ht0+\dp0+24pt\relax}
\endgroup
\begingroup
\setbox0=\vbox{\hsize=\linewidth
{\normalfont\normalsize\bfseries Restricciones físicas observadas\par}
\begin{itemize}
\item Clave foránea: created\_\allowbreak{}by\_\allowbreak{}user\_\allowbreak{}id hacia app.users (id); al eliminar: referencia nula.
\end{itemize}
}
\Needspace{\dimexpr\ht0+\dp0+24pt\relax}
\endgroup
\paragraph{Restricciones físicas observadas} \begin{itemize}
\item Clave foránea: created\_\allowbreak{}by\_\allowbreak{}user\_\allowbreak{}id hacia app.users (id); al eliminar: referencia nula.
\item Clave foránea: metadata\_\allowbreak{}snapshot\_\allowbreak{}id hacia app.metadata\_\allowbreak{}snapshots (id); al eliminar: restricción.
\item Clave foránea: reviewed\_\allowbreak{}by\_\allowbreak{}user\_\allowbreak{}id hacia app.users (id); al eliminar: referencia nula.
\item Clave foránea: source\_\allowbreak{}proposal\_\allowbreak{}id hacia app.bi\_\allowbreak{}proposals (id); al eliminar: referencia nula.
\item Clave primaria: id.
\end{itemize}

\begingroup
\setbox0=\vbox{\hsize=\linewidth
{\normalfont\normalsize\bfseries Índices únicos\par}
\begin{itemize}
\item id; índice completo.
\end{itemize}
}
\Needspace{\dimexpr\ht0+\dp0+24pt\relax}
\endgroup
\begingroup
\setbox0=\vbox{\hsize=\linewidth
{\normalfont\normalsize\bfseries Índices únicos\par}
\begin{itemize}
\item id; índice completo.
\end{itemize}
}
\Needspace{\dimexpr\ht0+\dp0+24pt\relax}
\endgroup
\paragraph{Índices únicos} \begin{itemize}
\item id; índice completo.
\end{itemize}

\begingroup
\setbox0=\vbox{\hsize=\linewidth
{\normalfont\normalsize\sffamily\bfseries\normalsize\color{Teal}\hyphenpenalty=10000\exhyphenpenalty=10000 app.business\_\allowbreak{}need\_\allowbreak{}reviews\par}

\par\vspace{6pt}\noindent\textbf{Área}\hspace{0.5em} Propuestas y ejecución

\par\vspace{6pt}\noindent\textbf{Qué guarda}\hspace{0.5em} Conserva la evaluación previa de una necesidad, contexto y huellas.

\par\vspace{6pt}\noindent\textbf{Cómo se utiliza}\hspace{0.5em} Paso Necesidad. La revisión evalúa viabilidad antes de generar; no crea tablas ni garantiza la calidad de todos los datos.

\par
\begingroup
\fontsize{10}{12.5}\selectfont
\setlength{\tabcolsep}{5pt}
\renewcommand{\arraystretch}{1.15}
\begin{tabular}{@{}>{\RaggedRight\arraybackslash}p{0.23\linewidth}>{\RaggedRight\arraybackslash}p{0.23\linewidth}>{\RaggedRight\arraybackslash}p{0.45\linewidth}@{}}
\toprule
\textbf{Columna} & \textbf{Tipo y condición} & \textbf{Qué representa} \\
\midrule
\texttt{\seqsplit{id}} & integer\par\smallskip Obligatorio. Clave primaria. Con valor predeterminado. & Identificador interno de esta entidad; su PK distingue cada registro. \\
\bottomrule
\end{tabular}
\endgroup
}
\Needspace{\dimexpr\ht0+\dp0+48pt\relax}
\endgroup
\subsection{app.business\_\allowbreak{}need\_\allowbreak{}reviews}
\paragraph{Área} Propuestas y ejecución

\paragraph{Qué guarda} Conserva la evaluación previa de una necesidad, contexto y huellas.

\paragraph{Cómo se utiliza} Paso Necesidad. La revisión evalúa viabilidad antes de generar; no crea tablas ni garantiza la calidad de todos los datos.

\begingroup
\fontsize{10}{12.5}\selectfont
\setlength{\tabcolsep}{5pt}
\renewcommand{\arraystretch}{1.15}
\begin{longtable}{@{}>{\RaggedRight\arraybackslash}p{0.23\linewidth}>{\RaggedRight\arraybackslash}p{0.23\linewidth}>{\RaggedRight\arraybackslash}p{0.45\linewidth}@{}}
\toprule
\textbf{Columna} & \textbf{Tipo y condición} & \textbf{Qué representa} \\
\midrule
\endfirsthead
\toprule
\textbf{Columna} & \textbf{Tipo y condición} & \textbf{Qué representa} \\
\midrule
\endhead
\midrule
\multicolumn{3}{r}{\footnotesize Continúa en la página siguiente}\\
\endfoot
\bottomrule
\endlastfoot
\texttt{\seqsplit{id}} & integer\par\smallskip Obligatorio. Clave primaria. Con valor predeterminado. & Identificador interno de esta entidad; su PK distingue cada registro. \\
\addlinespace[3pt]
\texttt{\seqsplit{metadata\_snapshot\_id}} & integer\par\smallskip Obligatorio. & FK a app.metadata\_\allowbreak{}snapshots.id; enlaza el artefacto a una estructura de origen conservada. \\
\addlinespace[3pt]
\texttt{\seqsplit{input\_hash}} & character varying(64)\par\smallskip Obligatorio. & Huella de entradas y contexto de la solicitud; no es el hash de contraseña de una persona. \\
\addlinespace[3pt]
\texttt{\seqsplit{assessment\_hash}} & character varying(64)\par\smallskip Obligatorio. & Huella del resultado de la revisión previa. \\
\addlinespace[3pt]
\texttt{\seqsplit{assessment\_document}} & jsonb\par\smallskip Obligatorio. & JSONB de la evaluación de viabilidad, sus requisitos y límites; no se ejecuta como SQL. \\
\addlinespace[3pt]
\texttt{\seqsplit{provider\_kind}} & character varying(40)\par\smallskip Obligatorio. & Tipo de proveedor de IA. En los expedientes conserva la procedencia como texto, sin FK a la configuración que después puede cambiar. \\
\addlinespace[3pt]
\texttt{\seqsplit{model\_id}} & character varying(160)\par\smallskip Obligatorio. & Identificador del modelo utilizado o configurado, distinto de una API key. \\
\addlinespace[3pt]
\texttt{\seqsplit{created\_by\_user\_id}} & integer\par\smallskip Admite nulos. & FK opcional a la persona que creó el artefacto o expediente. \\
\addlinespace[3pt]
\texttt{\seqsplit{created\_at}} & timestamp with time zone\par\smallskip Obligatorio. Con valor predeterminado. & Fecha de creación del registro; su opcionalidad física consta en el inventario. \\
\addlinespace[3pt]
\end{longtable}
\endgroup

\begingroup
\setbox0=\vbox{\hsize=\linewidth
{\normalfont\normalsize\bfseries Restricciones físicas observadas\par}
\begin{itemize}
\item Clave foránea: created\_\allowbreak{}by\_\allowbreak{}user\_\allowbreak{}id hacia app.users (id); al eliminar: referencia nula.
\end{itemize}
}
\Needspace{\dimexpr\ht0+\dp0+24pt\relax}
\endgroup
\begingroup
\setbox0=\vbox{\hsize=\linewidth
{\normalfont\normalsize\bfseries Restricciones físicas observadas\par}
\begin{itemize}
\item Clave foránea: created\_\allowbreak{}by\_\allowbreak{}user\_\allowbreak{}id hacia app.users (id); al eliminar: referencia nula.
\end{itemize}
}
\Needspace{\dimexpr\ht0+\dp0+24pt\relax}
\endgroup
\paragraph{Restricciones físicas observadas} \begin{itemize}
\item Clave foránea: created\_\allowbreak{}by\_\allowbreak{}user\_\allowbreak{}id hacia app.users (id); al eliminar: referencia nula.
\item Clave foránea: metadata\_\allowbreak{}snapshot\_\allowbreak{}id hacia app.metadata\_\allowbreak{}snapshots (id); al eliminar: restricción.
\item Clave primaria: id.
\end{itemize}

\begingroup
\setbox0=\vbox{\hsize=\linewidth
{\normalfont\normalsize\bfseries Índices únicos\par}
\begin{itemize}
\item id; índice completo.
\end{itemize}
}
\Needspace{\dimexpr\ht0+\dp0+24pt\relax}
\endgroup
\begingroup
\setbox0=\vbox{\hsize=\linewidth
{\normalfont\normalsize\bfseries Índices únicos\par}
\begin{itemize}
\item id; índice completo.
\end{itemize}
}
\Needspace{\dimexpr\ht0+\dp0+24pt\relax}
\endgroup
\paragraph{Índices únicos} \begin{itemize}
\item id; índice completo.
\end{itemize}

\begingroup
\setbox0=\vbox{\hsize=\linewidth
{\normalfont\normalsize\sffamily\bfseries\normalsize\color{Teal}\hyphenpenalty=10000\exhyphenpenalty=10000 app.data\_\allowbreak{}connections\par}

\par\vspace{6pt}\noindent\textbf{Área}\hspace{0.5em} Configuración y metadatos

\par\vspace{6pt}\noindent\textbf{Qué guarda}\hspace{0.5em} Identifica una fuente de datos, sus propiedades de conexión y la referencia a una credencial cifrada.

\par\vspace{6pt}\noindent\textbf{Cómo se utiliza}\hspace{0.5em} Pantalla Conexiones de datos y selector de fuente. Ambas conexiones pueden estar habilitadas; seleccionar una no fusiona sus datos.

\par
\begingroup
\fontsize{10}{12.5}\selectfont
\setlength{\tabcolsep}{5pt}
\renewcommand{\arraystretch}{1.15}
\begin{tabular}{@{}>{\RaggedRight\arraybackslash}p{0.23\linewidth}>{\RaggedRight\arraybackslash}p{0.23\linewidth}>{\RaggedRight\arraybackslash}p{0.45\linewidth}@{}}
\toprule
\textbf{Columna} & \textbf{Tipo y condición} & \textbf{Qué representa} \\
\midrule
\texttt{\seqsplit{id}} & integer\par\smallskip Obligatorio. Clave primaria. Con valor predeterminado. & Identificador interno de esta entidad; su PK distingue cada registro. \\
\bottomrule
\end{tabular}
\endgroup
}
\Needspace{\dimexpr\ht0+\dp0+48pt\relax}
\endgroup
\subsection{app.data\_\allowbreak{}connections}
\paragraph{Área} Configuración y metadatos

\paragraph{Qué guarda} Identifica una fuente de datos, sus propiedades de conexión y la referencia a una credencial cifrada.

\paragraph{Cómo se utiliza} Pantalla Conexiones de datos y selector de fuente. Ambas conexiones pueden estar habilitadas; seleccionar una no fusiona sus datos.

\begingroup
\fontsize{10}{12.5}\selectfont
\setlength{\tabcolsep}{5pt}
\renewcommand{\arraystretch}{1.15}
\begin{longtable}{@{}>{\RaggedRight\arraybackslash}p{0.23\linewidth}>{\RaggedRight\arraybackslash}p{0.23\linewidth}>{\RaggedRight\arraybackslash}p{0.45\linewidth}@{}}
\toprule
\textbf{Columna} & \textbf{Tipo y condición} & \textbf{Qué representa} \\
\midrule
\endfirsthead
\toprule
\textbf{Columna} & \textbf{Tipo y condición} & \textbf{Qué representa} \\
\midrule
\endhead
\midrule
\multicolumn{3}{r}{\footnotesize Continúa en la página siguiente}\\
\endfoot
\bottomrule
\endlastfoot
\texttt{\seqsplit{id}} & integer\par\smallskip Obligatorio. Clave primaria. Con valor predeterminado. & Identificador interno de esta entidad; su PK distingue cada registro. \\
\addlinespace[3pt]
\texttt{\seqsplit{name}} & character varying(120)\par\smallskip Obligatorio. & Nombre o denominación visible en el módulo. \\
\addlinespace[3pt]
\texttt{\seqsplit{connector\_kind}} & character varying(40)\par\smallskip Obligatorio. & Tipo de conector de la fuente; las dos fuentes ensayadas usan el adaptador SQL Server. \\
\addlinespace[3pt]
\texttt{\seqsplit{host}} & character varying(255)\par\smallskip Obligatorio. & Servidor de origen configurado; no se extrajo su valor. \\
\addlinespace[3pt]
\texttt{\seqsplit{port}} & integer\par\smallskip Obligatorio. & Puerto de la conexión de origen; no se extrajo su valor. \\
\addlinespace[3pt]
\texttt{\seqsplit{database\_name}} & character varying(128)\par\smallskip Obligatorio. & Nombre de la base de origen. En una instantánea conserva la identidad de la captura. \\
\addlinespace[3pt]
\texttt{\seqsplit{username}} & character varying(128)\par\smallskip Obligatorio. & Cuenta técnica de lectura de la fuente; no es la cuenta con la que una persona entra a la aplicación. No se extrajo su valor. \\
\addlinespace[3pt]
\texttt{\seqsplit{encrypt}} & boolean\par\smallskip Obligatorio. & Opción de cifrado de transporte hacia la fuente. \\
\addlinespace[3pt]
\texttt{\seqsplit{trust\_server\_certificate}} & boolean\par\smallskip Obligatorio. & Opción de validación del certificado de la conexión de origen. \\
\addlinespace[3pt]
\texttt{\seqsplit{secret\_id}} & integer\par\smallskip Obligatorio. & FK a app.secrets.id. La credencial se recupera y descifra sólo bajo controles del servidor. \\
\addlinespace[3pt]
\texttt{\seqsplit{is\_active}} & boolean\par\smallskip Obligatorio. Con valor predeterminado. & Estado de habilitación del registro; su efecto depende de las reglas del módulo. \\
\addlinespace[3pt]
\texttt{\seqsplit{last\_test\_status}} & character varying(40)\par\smallskip Admite nulos. & Resultado de la última prueba de conexión del registro. \\
\addlinespace[3pt]
\texttt{\seqsplit{last\_test\_message}} & text\par\smallskip Admite nulos. & Descripción del resultado de la última prueba; no se extrajo su texto. \\
\addlinespace[3pt]
\texttt{\seqsplit{last\_tested\_at}} & timestamp with time zone\par\smallskip Admite nulos. & Fecha de la última prueba de conectividad. \\
\addlinespace[3pt]
\texttt{\seqsplit{created\_at}} & timestamp with time zone\par\smallskip Admite nulos. Con valor predeterminado. & Fecha de creación del registro; su opcionalidad física consta en el inventario. \\
\addlinespace[3pt]
\texttt{\seqsplit{updated\_at}} & timestamp with time zone\par\smallskip Admite nulos. Con valor predeterminado. & Fecha de actualización del registro; su opcionalidad física consta en el inventario. \\
\addlinespace[3pt]
\end{longtable}
\endgroup

\begingroup
\setbox0=\vbox{\hsize=\linewidth
{\normalfont\normalsize\bfseries Restricciones físicas observadas\par}
\begin{itemize}
\item Clave foránea: secret\_\allowbreak{}id hacia app.secrets (id); al eliminar: restricción.
\end{itemize}
}
\Needspace{\dimexpr\ht0+\dp0+24pt\relax}
\endgroup
\begingroup
\setbox0=\vbox{\hsize=\linewidth
{\normalfont\normalsize\bfseries Restricciones físicas observadas\par}
\begin{itemize}
\item Clave foránea: secret\_\allowbreak{}id hacia app.secrets (id); al eliminar: restricción.
\end{itemize}
}
\Needspace{\dimexpr\ht0+\dp0+24pt\relax}
\endgroup
\paragraph{Restricciones físicas observadas} \begin{itemize}
\item Clave foránea: secret\_\allowbreak{}id hacia app.secrets (id); al eliminar: restricción.
\item Clave primaria: id.
\item Unicidad: secret\_\allowbreak{}id.
\end{itemize}

\begingroup
\setbox0=\vbox{\hsize=\linewidth
{\normalfont\normalsize\bfseries Índices únicos\par}
\begin{itemize}
\item id; índice completo.
\end{itemize}
}
\Needspace{\dimexpr\ht0+\dp0+24pt\relax}
\endgroup
\begingroup
\setbox0=\vbox{\hsize=\linewidth
{\normalfont\normalsize\bfseries Índices únicos\par}
\begin{itemize}
\item id; índice completo.
\end{itemize}
}
\Needspace{\dimexpr\ht0+\dp0+24pt\relax}
\endgroup
\paragraph{Índices únicos} \begin{itemize}
\item id; índice completo.
\item secret\_\allowbreak{}id; índice completo.
\end{itemize}

\begingroup
\setbox0=\vbox{\hsize=\linewidth
{\normalfont\normalsize\sffamily\bfseries\normalsize\color{Teal}\hyphenpenalty=10000\exhyphenpenalty=10000 app.etl\_\allowbreak{}executions\par}

\par\vspace{6pt}\noindent\textbf{Área}\hspace{0.5em} Propuestas y ejecución

\par\vspace{6pt}\noindent\textbf{Qué guarda}\hspace{0.5em} Conserva el expediente de preparación y ejecución: selección, plan, validación, métricas y fechas.

\par\vspace{6pt}\noindent\textbf{Cómo se utiliza}\hspace{0.5em} Datamart de ventas e historial de ejecuciones. El registro del expediente no es la tabla de hechos; los datos cargados viven en el esquema analítico de esa ejecución.

\par
\begingroup
\fontsize{10}{12.5}\selectfont
\setlength{\tabcolsep}{5pt}
\renewcommand{\arraystretch}{1.15}
\begin{tabular}{@{}>{\RaggedRight\arraybackslash}p{0.23\linewidth}>{\RaggedRight\arraybackslash}p{0.23\linewidth}>{\RaggedRight\arraybackslash}p{0.45\linewidth}@{}}
\toprule
\textbf{Columna} & \textbf{Tipo y condición} & \textbf{Qué representa} \\
\midrule
\texttt{\seqsplit{id}} & integer\par\smallskip Obligatorio. Clave primaria. Con valor predeterminado. & Identificador interno de esta entidad; su PK distingue cada registro. \\
\bottomrule
\end{tabular}
\endgroup
}
\Needspace{\dimexpr\ht0+\dp0+48pt\relax}
\endgroup
\subsection{app.etl\_\allowbreak{}executions}
\paragraph{Área} Propuestas y ejecución

\paragraph{Qué guarda} Conserva el expediente de preparación y ejecución: selección, plan, validación, métricas y fechas.

\paragraph{Cómo se utiliza} Datamart de ventas e historial de ejecuciones. El registro del expediente no es la tabla de hechos; los datos cargados viven en el esquema analítico de esa ejecución.

\begingroup
\fontsize{10}{12.5}\selectfont
\setlength{\tabcolsep}{5pt}
\renewcommand{\arraystretch}{1.15}
\begin{longtable}{@{}>{\RaggedRight\arraybackslash}p{0.23\linewidth}>{\RaggedRight\arraybackslash}p{0.23\linewidth}>{\RaggedRight\arraybackslash}p{0.45\linewidth}@{}}
\toprule
\textbf{Columna} & \textbf{Tipo y condición} & \textbf{Qué representa} \\
\midrule
\endfirsthead
\toprule
\textbf{Columna} & \textbf{Tipo y condición} & \textbf{Qué representa} \\
\midrule
\endhead
\midrule
\multicolumn{3}{r}{\footnotesize Continúa en la página siguiente}\\
\endfoot
\bottomrule
\endlastfoot
\texttt{\seqsplit{id}} & integer\par\smallskip Obligatorio. Clave primaria. Con valor predeterminado. & Identificador interno de esta entidad; su PK distingue cada registro. \\
\addlinespace[3pt]
\texttt{\seqsplit{proposal\_id}} & integer\par\smallskip Obligatorio. & FK a app.bi\_\allowbreak{}proposals.id; vincula la asesoría o ejecución con una propuesta concreta. \\
\addlinespace[3pt]
\texttt{\seqsplit{metadata\_snapshot\_id}} & integer\par\smallskip Obligatorio. & FK a app.metadata\_\allowbreak{}snapshots.id; enlaza el artefacto a una estructura de origen conservada. \\
\addlinespace[3pt]
\texttt{\seqsplit{status}} & character varying(40)\par\smallskip Obligatorio. & Estado del ciclo de vida del artefacto. Una propuesta aprobada no equivale a una ejecución terminada. \\
\addlinespace[3pt]
\texttt{\seqsplit{domain\_code}} & character varying(40)\par\smallskip Obligatorio. & Dominio analítico; ventas es el dominio implementado. \\
\addlinespace[3pt]
\texttt{\seqsplit{builder\_version}} & character varying(40)\par\smallskip Obligatorio. & Versión del constructor determinístico que preparó el plan de ejecución. \\
\addlinespace[3pt]
\texttt{\seqsplit{proposal\_hash}} & character varying(64)\par\smallskip Obligatorio. & Huella del contrato de propuesta usado para la ejecución. \\
\addlinespace[3pt]
\texttt{\seqsplit{snapshot\_hash}} & character varying(64)\par\smallskip Obligatorio. & Huella de los metadatos conservados utilizados para la ejecución. \\
\addlinespace[3pt]
\texttt{\seqsplit{selection\_document}} & jsonb\par\smallskip Obligatorio. & JSONB de indicadores y decisiones seleccionados para materializar. \\
\addlinespace[3pt]
\texttt{\seqsplit{plan\_document}} & jsonb\par\smallskip Obligatorio. & JSONB del plan compilado y su destino; el nombre de un esquema en un documento no constituye una FK de PostgreSQL. \\
\addlinespace[3pt]
\texttt{\seqsplit{validation\_document}} & jsonb\par\smallskip Obligatorio. & JSONB de comprobaciones, errores o advertencias del artefacto o ejecución. \\
\addlinespace[3pt]
\texttt{\seqsplit{metrics\_document}} & jsonb\par\smallskip Obligatorio. & JSONB de métricas y comprobaciones de la carga; no contiene la tabla física de hechos. \\
\addlinespace[3pt]
\texttt{\seqsplit{created\_by\_user\_id}} & integer\par\smallskip Admite nulos. & FK opcional a la persona que creó el artefacto o expediente. \\
\addlinespace[3pt]
\texttt{\seqsplit{created\_by\_label}} & character varying(320)\par\smallskip Obligatorio. & Etiqueta histórica de la persona creadora. \\
\addlinespace[3pt]
\texttt{\seqsplit{created\_at}} & timestamp with time zone\par\smallskip Obligatorio. Con valor predeterminado. & Fecha de creación del registro; su opcionalidad física consta en el inventario. \\
\addlinespace[3pt]
\texttt{\seqsplit{started\_at}} & timestamp with time zone\par\smallskip Admite nulos. & Fecha de comienzo de la ejecución. \\
\addlinespace[3pt]
\texttt{\seqsplit{finished\_at}} & timestamp with time zone\par\smallskip Admite nulos. & Fecha de finalización de la ejecución. \\
\addlinespace[3pt]
\end{longtable}
\endgroup

\begingroup
\setbox0=\vbox{\hsize=\linewidth
{\normalfont\normalsize\bfseries Restricciones físicas observadas\par}
\begin{itemize}
\item Clave foránea: created\_\allowbreak{}by\_\allowbreak{}user\_\allowbreak{}id hacia app.users (id); al eliminar: referencia nula.
\end{itemize}
}
\Needspace{\dimexpr\ht0+\dp0+24pt\relax}
\endgroup
\begingroup
\setbox0=\vbox{\hsize=\linewidth
{\normalfont\normalsize\bfseries Restricciones físicas observadas\par}
\begin{itemize}
\item Clave foránea: created\_\allowbreak{}by\_\allowbreak{}user\_\allowbreak{}id hacia app.users (id); al eliminar: referencia nula.
\end{itemize}
}
\Needspace{\dimexpr\ht0+\dp0+24pt\relax}
\endgroup
\paragraph{Restricciones físicas observadas} \begin{itemize}
\item Clave foránea: created\_\allowbreak{}by\_\allowbreak{}user\_\allowbreak{}id hacia app.users (id); al eliminar: referencia nula.
\item Clave foránea: metadata\_\allowbreak{}snapshot\_\allowbreak{}id hacia app.metadata\_\allowbreak{}snapshots (id); al eliminar: restricción.
\item Clave foránea: proposal\_\allowbreak{}id hacia app.bi\_\allowbreak{}proposals (id); al eliminar: restricción.
\item Clave primaria: id.
\end{itemize}

\begingroup
\setbox0=\vbox{\hsize=\linewidth
{\normalfont\normalsize\bfseries Índices únicos\par}
\begin{itemize}
\item id; índice completo.
\end{itemize}
}
\Needspace{\dimexpr\ht0+\dp0+24pt\relax}
\endgroup
\begingroup
\setbox0=\vbox{\hsize=\linewidth
{\normalfont\normalsize\bfseries Índices únicos\par}
\begin{itemize}
\item id; índice completo.
\end{itemize}
}
\Needspace{\dimexpr\ht0+\dp0+24pt\relax}
\endgroup
\paragraph{Índices únicos} \begin{itemize}
\item id; índice completo.
\end{itemize}

\begingroup
\setbox0=\vbox{\hsize=\linewidth
{\normalfont\normalsize\sffamily\bfseries\normalsize\color{Teal}\hyphenpenalty=10000\exhyphenpenalty=10000 app.llm\_\allowbreak{}configurations\par}

\par\vspace{6pt}\noindent\textbf{Área}\hspace{0.5em} Configuración y metadatos

\par\vspace{6pt}\noindent\textbf{Qué guarda}\hspace{0.5em} Configura proveedor, modelo, endpoint, razonamiento y prueba de conectividad.

\par\vspace{6pt}\noindent\textbf{Cómo se utiliza}\hspace{0.5em} Pantalla Configuración IA. secret\_\allowbreak{}id es opcional y único dentro de esta tabla; un proveedor local puede no necesitar API key.

\par
\begingroup
\fontsize{10}{12.5}\selectfont
\setlength{\tabcolsep}{5pt}
\renewcommand{\arraystretch}{1.15}
\begin{tabular}{@{}>{\RaggedRight\arraybackslash}p{0.23\linewidth}>{\RaggedRight\arraybackslash}p{0.23\linewidth}>{\RaggedRight\arraybackslash}p{0.45\linewidth}@{}}
\toprule
\textbf{Columna} & \textbf{Tipo y condición} & \textbf{Qué representa} \\
\midrule
\texttt{\seqsplit{id}} & integer\par\smallskip Obligatorio. Clave primaria. Con valor predeterminado. & Identificador interno de esta entidad; su PK distingue cada registro. \\
\bottomrule
\end{tabular}
\endgroup
}
\Needspace{\dimexpr\ht0+\dp0+48pt\relax}
\endgroup
\subsection{app.llm\_\allowbreak{}configurations}
\paragraph{Área} Configuración y metadatos

\paragraph{Qué guarda} Configura proveedor, modelo, endpoint, razonamiento y prueba de conectividad.

\paragraph{Cómo se utiliza} Pantalla Configuración IA. secret\_\allowbreak{}id es opcional y único dentro de esta tabla; un proveedor local puede no necesitar API key.

\begingroup
\fontsize{10}{12.5}\selectfont
\setlength{\tabcolsep}{5pt}
\renewcommand{\arraystretch}{1.15}
\begin{longtable}{@{}>{\RaggedRight\arraybackslash}p{0.23\linewidth}>{\RaggedRight\arraybackslash}p{0.23\linewidth}>{\RaggedRight\arraybackslash}p{0.45\linewidth}@{}}
\toprule
\textbf{Columna} & \textbf{Tipo y condición} & \textbf{Qué representa} \\
\midrule
\endfirsthead
\toprule
\textbf{Columna} & \textbf{Tipo y condición} & \textbf{Qué representa} \\
\midrule
\endhead
\midrule
\multicolumn{3}{r}{\footnotesize Continúa en la página siguiente}\\
\endfoot
\bottomrule
\endlastfoot
\texttt{\seqsplit{id}} & integer\par\smallskip Obligatorio. Clave primaria. Con valor predeterminado. & Identificador interno de esta entidad; su PK distingue cada registro. \\
\addlinespace[3pt]
\texttt{\seqsplit{name}} & character varying(120)\par\smallskip Obligatorio. & Nombre o denominación visible en el módulo. \\
\addlinespace[3pt]
\texttt{\seqsplit{provider\_kind}} & character varying(40)\par\smallskip Obligatorio. & Tipo de proveedor de IA. En los expedientes conserva la procedencia como texto, sin FK a la configuración que después puede cambiar. \\
\addlinespace[3pt]
\texttt{\seqsplit{base\_url}} & character varying(300)\par\smallskip Obligatorio. & Dirección del servicio de IA configurado; no se extrajo su valor. \\
\addlinespace[3pt]
\texttt{\seqsplit{model\_id}} & character varying(160)\par\smallskip Obligatorio. & Identificador del modelo utilizado o configurado, distinto de una API key. \\
\addlinespace[3pt]
\texttt{\seqsplit{credential\_reference}} & character varying(80)\par\smallskip Obligatorio. & Referencia técnica de la credencial; no equivale a una clave en texto. \\
\addlinespace[3pt]
\texttt{\seqsplit{is\_active}} & boolean\par\smallskip Obligatorio. Con valor predeterminado. & Estado de habilitación del registro; su efecto depende de las reglas del módulo. \\
\addlinespace[3pt]
\texttt{\seqsplit{last\_test\_status}} & character varying(40)\par\smallskip Admite nulos. & Resultado de la última prueba de conexión del registro. \\
\addlinespace[3pt]
\texttt{\seqsplit{last\_test\_message}} & text\par\smallskip Admite nulos. & Descripción del resultado de la última prueba; no se extrajo su texto. \\
\addlinespace[3pt]
\texttt{\seqsplit{last\_tested\_at}} & timestamp with time zone\par\smallskip Admite nulos. & Fecha de la última prueba de conectividad. \\
\addlinespace[3pt]
\texttt{\seqsplit{created\_at}} & timestamp with time zone\par\smallskip Admite nulos. Con valor predeterminado. & Fecha de creación del registro; su opcionalidad física consta en el inventario. \\
\addlinespace[3pt]
\texttt{\seqsplit{updated\_at}} & timestamp with time zone\par\smallskip Admite nulos. Con valor predeterminado. & Fecha de actualización del registro; su opcionalidad física consta en el inventario. \\
\addlinespace[3pt]
\texttt{\seqsplit{secret\_id}} & integer\par\smallskip Admite nulos. & FK a app.secrets.id. La credencial se recupera y descifra sólo bajo controles del servidor. \\
\addlinespace[3pt]
\texttt{\seqsplit{reasoning\_level}} & character varying(20)\par\smallskip Obligatorio. Con valor predeterminado. & Nivel de razonamiento admitido por la configuración y el adaptador. \\
\addlinespace[3pt]
\end{longtable}
\endgroup

\begingroup
\setbox0=\vbox{\hsize=\linewidth
{\normalfont\normalsize\bfseries Restricciones físicas observadas\par}
\begin{itemize}
\item Clave foránea: secret\_\allowbreak{}id hacia app.secrets (id); al eliminar: referencia nula.
\end{itemize}
}
\Needspace{\dimexpr\ht0+\dp0+24pt\relax}
\endgroup
\begingroup
\setbox0=\vbox{\hsize=\linewidth
{\normalfont\normalsize\bfseries Restricciones físicas observadas\par}
\begin{itemize}
\item Clave foránea: secret\_\allowbreak{}id hacia app.secrets (id); al eliminar: referencia nula.
\end{itemize}
}
\Needspace{\dimexpr\ht0+\dp0+24pt\relax}
\endgroup
\paragraph{Restricciones físicas observadas} \begin{itemize}
\item Clave foránea: secret\_\allowbreak{}id hacia app.secrets (id); al eliminar: referencia nula.
\item Clave primaria: id.
\item Unicidad: secret\_\allowbreak{}id.
\end{itemize}

\begingroup
\setbox0=\vbox{\hsize=\linewidth
{\normalfont\normalsize\bfseries Índices únicos\par}
\begin{itemize}
\item id; índice completo.
\end{itemize}
}
\Needspace{\dimexpr\ht0+\dp0+24pt\relax}
\endgroup
\begingroup
\setbox0=\vbox{\hsize=\linewidth
{\normalfont\normalsize\bfseries Índices únicos\par}
\begin{itemize}
\item id; índice completo.
\end{itemize}
}
\Needspace{\dimexpr\ht0+\dp0+24pt\relax}
\endgroup
\paragraph{Índices únicos} \begin{itemize}
\item id; índice completo.
\item secret\_\allowbreak{}id; índice completo.
\end{itemize}

\begingroup
\setbox0=\vbox{\hsize=\linewidth
{\normalfont\normalsize\sffamily\bfseries\normalsize\color{Teal}\hyphenpenalty=10000\exhyphenpenalty=10000 app.menu\_\allowbreak{}permissions\par}

\par\vspace{6pt}\noindent\textbf{Área}\hspace{0.5em} Seguridad

\par\vspace{6pt}\noindent\textbf{Qué guarda}\hspace{0.5em} Relaciona entradas de navegación con permisos.

\par\vspace{6pt}\noindent\textbf{Cómo se utiliza}\hspace{0.5em} Determina qué requisito de permiso acompaña a una entrada del menú; no es una asignación directa de menú a rol.

\par
\begingroup
\fontsize{10}{12.5}\selectfont
\setlength{\tabcolsep}{5pt}
\renewcommand{\arraystretch}{1.15}
\begin{tabular}{@{}>{\RaggedRight\arraybackslash}p{0.23\linewidth}>{\RaggedRight\arraybackslash}p{0.23\linewidth}>{\RaggedRight\arraybackslash}p{0.45\linewidth}@{}}
\toprule
\textbf{Columna} & \textbf{Tipo y condición} & \textbf{Qué representa} \\
\midrule
\texttt{\seqsplit{menu\_id}} & integer\par\smallskip Obligatorio. Clave primaria. & FK a app.menus.id; identifica una entrada de navegación. \\
\bottomrule
\end{tabular}
\endgroup
}
\Needspace{\dimexpr\ht0+\dp0+48pt\relax}
\endgroup
\subsection{app.menu\_\allowbreak{}permissions}
\paragraph{Área} Seguridad

\paragraph{Qué guarda} Relaciona entradas de navegación con permisos.

\paragraph{Cómo se utiliza} Determina qué requisito de permiso acompaña a una entrada del menú; no es una asignación directa de menú a rol.

\begingroup
\fontsize{10}{12.5}\selectfont
\setlength{\tabcolsep}{5pt}
\renewcommand{\arraystretch}{1.15}
\begin{longtable}{@{}>{\RaggedRight\arraybackslash}p{0.23\linewidth}>{\RaggedRight\arraybackslash}p{0.23\linewidth}>{\RaggedRight\arraybackslash}p{0.45\linewidth}@{}}
\toprule
\textbf{Columna} & \textbf{Tipo y condición} & \textbf{Qué representa} \\
\midrule
\endfirsthead
\toprule
\textbf{Columna} & \textbf{Tipo y condición} & \textbf{Qué representa} \\
\midrule
\endhead
\midrule
\multicolumn{3}{r}{\footnotesize Continúa en la página siguiente}\\
\endfoot
\bottomrule
\endlastfoot
\texttt{\seqsplit{menu\_id}} & integer\par\smallskip Obligatorio. Clave primaria. & FK a app.menus.id; identifica una entrada de navegación. \\
\addlinespace[3pt]
\texttt{\seqsplit{permission\_id}} & integer\par\smallskip Obligatorio. Clave primaria. & FK a app.permissions.id; identifica una operación autorizable. \\
\addlinespace[3pt]
\end{longtable}
\endgroup

\begingroup
\setbox0=\vbox{\hsize=\linewidth
{\normalfont\normalsize\bfseries Restricciones físicas observadas\par}
\begin{itemize}
\item Clave foránea: menu\_\allowbreak{}id hacia app.menus (id); al eliminar: cascada SQL.
\end{itemize}
}
\Needspace{\dimexpr\ht0+\dp0+24pt\relax}
\endgroup
\begingroup
\setbox0=\vbox{\hsize=\linewidth
{\normalfont\normalsize\bfseries Restricciones físicas observadas\par}
\begin{itemize}
\item Clave foránea: menu\_\allowbreak{}id hacia app.menus (id); al eliminar: cascada SQL.
\end{itemize}
}
\Needspace{\dimexpr\ht0+\dp0+24pt\relax}
\endgroup
\paragraph{Restricciones físicas observadas} \begin{itemize}
\item Clave foránea: menu\_\allowbreak{}id hacia app.menus (id); al eliminar: cascada SQL.
\item Clave foránea: permission\_\allowbreak{}id hacia app.permissions (id); al eliminar: cascada SQL.
\item Clave primaria: menu\_\allowbreak{}id, permission\_\allowbreak{}id.
\end{itemize}

\begingroup
\setbox0=\vbox{\hsize=\linewidth
{\normalfont\normalsize\bfseries Índices únicos\par}
\begin{itemize}
\item menu\_\allowbreak{}id, permission\_\allowbreak{}id; índice completo.
\end{itemize}
}
\Needspace{\dimexpr\ht0+\dp0+24pt\relax}
\endgroup
\begingroup
\setbox0=\vbox{\hsize=\linewidth
{\normalfont\normalsize\bfseries Índices únicos\par}
\begin{itemize}
\item menu\_\allowbreak{}id, permission\_\allowbreak{}id; índice completo.
\end{itemize}
}
\Needspace{\dimexpr\ht0+\dp0+24pt\relax}
\endgroup
\paragraph{Índices únicos} \begin{itemize}
\item menu\_\allowbreak{}id, permission\_\allowbreak{}id; índice completo.
\end{itemize}

\begingroup
\setbox0=\vbox{\hsize=\linewidth
{\normalfont\normalsize\sffamily\bfseries\normalsize\color{Teal}\hyphenpenalty=10000\exhyphenpenalty=10000 app.menus\par}

\par\vspace{6pt}\noindent\textbf{Área}\hspace{0.5em} Seguridad

\par\vspace{6pt}\noindent\textbf{Qué guarda}\hspace{0.5em} Define entradas de navegación, ruta, agrupación y orden.

\par\vspace{6pt}\noindent\textbf{Cómo se utiliza}\hspace{0.5em} La barra lateral muestra las opciones permitidas. La autorización efectiva se comprueba en FastAPI.

\par
\begingroup
\fontsize{10}{12.5}\selectfont
\setlength{\tabcolsep}{5pt}
\renewcommand{\arraystretch}{1.15}
\begin{tabular}{@{}>{\RaggedRight\arraybackslash}p{0.23\linewidth}>{\RaggedRight\arraybackslash}p{0.23\linewidth}>{\RaggedRight\arraybackslash}p{0.45\linewidth}@{}}
\toprule
\textbf{Columna} & \textbf{Tipo y condición} & \textbf{Qué representa} \\
\midrule
\texttt{\seqsplit{id}} & integer\par\smallskip Obligatorio. Clave primaria. Con valor predeterminado. & Identificador interno de esta entidad; su PK distingue cada registro. \\
\bottomrule
\end{tabular}
\endgroup
}
\Needspace{\dimexpr\ht0+\dp0+48pt\relax}
\endgroup
\subsection{app.menus}
\paragraph{Área} Seguridad

\paragraph{Qué guarda} Define entradas de navegación, ruta, agrupación y orden.

\paragraph{Cómo se utiliza} La barra lateral muestra las opciones permitidas. La autorización efectiva se comprueba en FastAPI.

\begingroup
\fontsize{10}{12.5}\selectfont
\setlength{\tabcolsep}{5pt}
\renewcommand{\arraystretch}{1.15}
\begin{longtable}{@{}>{\RaggedRight\arraybackslash}p{0.23\linewidth}>{\RaggedRight\arraybackslash}p{0.23\linewidth}>{\RaggedRight\arraybackslash}p{0.45\linewidth}@{}}
\toprule
\textbf{Columna} & \textbf{Tipo y condición} & \textbf{Qué representa} \\
\midrule
\endfirsthead
\toprule
\textbf{Columna} & \textbf{Tipo y condición} & \textbf{Qué representa} \\
\midrule
\endhead
\midrule
\multicolumn{3}{r}{\footnotesize Continúa en la página siguiente}\\
\endfoot
\bottomrule
\endlastfoot
\texttt{\seqsplit{id}} & integer\par\smallskip Obligatorio. Clave primaria. Con valor predeterminado. & Identificador interno de esta entidad; su PK distingue cada registro. \\
\addlinespace[3pt]
\texttt{\seqsplit{code}} & character varying(100)\par\smallskip Obligatorio. & Código estable de la entidad; no depende de su etiqueta visible. \\
\addlinespace[3pt]
\texttt{\seqsplit{label}} & character varying(120)\par\smallskip Obligatorio. & Etiqueta de navegación que se presenta a la persona usuaria. \\
\addlinespace[3pt]
\texttt{\seqsplit{path}} & character varying(160)\par\smallskip Obligatorio. & Ruta de la pantalla; no es una credencial ni una regla de autorización. \\
\addlinespace[3pt]
\texttt{\seqsplit{position}} & integer\par\smallskip Obligatorio. Con valor predeterminado. & Orden de presentación dentro de la navegación. \\
\addlinespace[3pt]
\texttt{\seqsplit{is\_active}} & boolean\par\smallskip Obligatorio. Con valor predeterminado. & Estado de habilitación del registro; su efecto depende de las reglas del módulo. \\
\addlinespace[3pt]
\texttt{\seqsplit{module\_code}} & character varying(80)\par\smallskip Obligatorio. Con valor predeterminado. & Código del módulo o agrupación funcional. \\
\addlinespace[3pt]
\texttt{\seqsplit{module\_label}} & character varying(120)\par\smallskip Obligatorio. Con valor predeterminado. & Nombre visible de la agrupación de menús. \\
\addlinespace[3pt]
\texttt{\seqsplit{is\_system\_protected}} & boolean\par\smallskip Obligatorio. Con valor predeterminado. & Marca de un registro base protegido para conservar administración y navegación esenciales. \\
\addlinespace[3pt]
\end{longtable}
\endgroup

\begingroup
\setbox0=\vbox{\hsize=\linewidth
{\normalfont\normalsize\bfseries Restricciones físicas observadas\par}
\begin{itemize}
\item Clave primaria: id.
\end{itemize}
}
\Needspace{\dimexpr\ht0+\dp0+24pt\relax}
\endgroup
\begingroup
\setbox0=\vbox{\hsize=\linewidth
{\normalfont\normalsize\bfseries Restricciones físicas observadas\par}
\begin{itemize}
\item Clave primaria: id.
\end{itemize}
}
\Needspace{\dimexpr\ht0+\dp0+24pt\relax}
\endgroup
\paragraph{Restricciones físicas observadas} \begin{itemize}
\item Clave primaria: id.
\item Unicidad: code.
\item Unicidad: path.
\end{itemize}

\begingroup
\setbox0=\vbox{\hsize=\linewidth
{\normalfont\normalsize\bfseries Índices únicos\par}
\begin{itemize}
\item code; índice completo.
\end{itemize}
}
\Needspace{\dimexpr\ht0+\dp0+24pt\relax}
\endgroup
\begingroup
\setbox0=\vbox{\hsize=\linewidth
{\normalfont\normalsize\bfseries Índices únicos\par}
\begin{itemize}
\item code; índice completo.
\end{itemize}
}
\Needspace{\dimexpr\ht0+\dp0+24pt\relax}
\endgroup
\paragraph{Índices únicos} \begin{itemize}
\item code; índice completo.
\item id; índice completo.
\item code; índice completo.
\item path; índice completo.
\end{itemize}

\begingroup
\setbox0=\vbox{\hsize=\linewidth
{\normalfont\normalsize\sffamily\bfseries\normalsize\color{Teal}\hyphenpenalty=10000\exhyphenpenalty=10000 app.metadata\_\allowbreak{}snapshots\par}

\par\vspace{6pt}\noindent\textbf{Área}\hspace{0.5em} Configuración y metadatos

\par\vspace{6pt}\noindent\textbf{Qué guarda}\hspace{0.5em} Guarda una captura estructural inmutable de una fuente y sus huellas.

\par\vspace{6pt}\noindent\textbf{Cómo se utiliza}\hspace{0.5em} Explorador de esquema. El JSONB describe tablas, columnas y relaciones declaradas de la fuente, no contiene filas de ventas.

\par
\begingroup
\fontsize{10}{12.5}\selectfont
\setlength{\tabcolsep}{5pt}
\renewcommand{\arraystretch}{1.15}
\begin{tabular}{@{}>{\RaggedRight\arraybackslash}p{0.23\linewidth}>{\RaggedRight\arraybackslash}p{0.23\linewidth}>{\RaggedRight\arraybackslash}p{0.45\linewidth}@{}}
\toprule
\textbf{Columna} & \textbf{Tipo y condición} & \textbf{Qué representa} \\
\midrule
\texttt{\seqsplit{id}} & integer\par\smallskip Obligatorio. Clave primaria. Con valor predeterminado. & Identificador interno de esta entidad; su PK distingue cada registro. \\
\bottomrule
\end{tabular}
\endgroup
}
\Needspace{\dimexpr\ht0+\dp0+48pt\relax}
\endgroup
\subsection{app.metadata\_\allowbreak{}snapshots}
\paragraph{Área} Configuración y metadatos

\paragraph{Qué guarda} Guarda una captura estructural inmutable de una fuente y sus huellas.

\paragraph{Cómo se utiliza} Explorador de esquema. El JSONB describe tablas, columnas y relaciones declaradas de la fuente, no contiene filas de ventas.

\begingroup
\fontsize{10}{12.5}\selectfont
\setlength{\tabcolsep}{5pt}
\renewcommand{\arraystretch}{1.15}
\begin{longtable}{@{}>{\RaggedRight\arraybackslash}p{0.23\linewidth}>{\RaggedRight\arraybackslash}p{0.23\linewidth}>{\RaggedRight\arraybackslash}p{0.45\linewidth}@{}}
\toprule
\textbf{Columna} & \textbf{Tipo y condición} & \textbf{Qué representa} \\
\midrule
\endfirsthead
\toprule
\textbf{Columna} & \textbf{Tipo y condición} & \textbf{Qué representa} \\
\midrule
\endhead
\midrule
\multicolumn{3}{r}{\footnotesize Continúa en la página siguiente}\\
\endfoot
\bottomrule
\endlastfoot
\texttt{\seqsplit{id}} & integer\par\smallskip Obligatorio. Clave primaria. Con valor predeterminado. & Identificador interno de esta entidad; su PK distingue cada registro. \\
\addlinespace[3pt]
\texttt{\seqsplit{data\_connection\_id}} & integer\par\smallskip Obligatorio. & FK a app.data\_\allowbreak{}connections.id; preserva el contexto de fuente. \\
\addlinespace[3pt]
\texttt{\seqsplit{connector\_code}} & character varying(40)\par\smallskip Obligatorio. & Código del adaptador que realizó la captura estructural. \\
\addlinespace[3pt]
\texttt{\seqsplit{database\_name}} & character varying(128)\par\smallskip Obligatorio. & Nombre de la base de origen. En una instantánea conserva la identidad de la captura. \\
\addlinespace[3pt]
\texttt{\seqsplit{contract\_version}} & integer\par\smallskip Obligatorio. & Versión del contrato estructural guardado; no es el número de propuesta ni de ejecución. \\
\addlinespace[3pt]
\texttt{\seqsplit{content\_hash}} & character varying(64)\par\smallskip Obligatorio. & Huella del documento estructural canónico; permite detectar cambios y evitar duplicar capturas iguales. \\
\addlinespace[3pt]
\texttt{\seqsplit{schema\_document}} & jsonb\par\smallskip Obligatorio. & JSONB de los metadatos de origen: estructura, tipos, claves y relaciones declaradas. No son filas de negocio. \\
\addlinespace[3pt]
\texttt{\seqsplit{schema\_count}} & integer\par\smallskip Obligatorio. & Número de esquemas descritos en la instantánea de origen. \\
\addlinespace[3pt]
\texttt{\seqsplit{table\_count}} & integer\par\smallskip Obligatorio. & Número de tablas descritas en la instantánea de origen. \\
\addlinespace[3pt]
\texttt{\seqsplit{column\_count}} & integer\par\smallskip Obligatorio. & Número de columnas descritas en la instantánea de origen. \\
\addlinespace[3pt]
\texttt{\seqsplit{relationship\_count}} & integer\par\smallskip Obligatorio. & Número de relaciones declaradas descritas; no cuenta ventas ni coincidencias de datos. \\
\addlinespace[3pt]
\texttt{\seqsplit{captured\_by\_user\_id}} & integer\par\smallskip Admite nulos. & FK opcional al usuario que obtuvo la instantánea. \\
\addlinespace[3pt]
\texttt{\seqsplit{captured\_by\_label}} & character varying(320)\par\smallskip Obligatorio. & Etiqueta histórica de la persona que obtuvo la instantánea. \\
\addlinespace[3pt]
\texttt{\seqsplit{captured\_at}} & timestamp with time zone\par\smallskip Obligatorio. Con valor predeterminado. & Fecha de captura de los metadatos. \\
\addlinespace[3pt]
\end{longtable}
\endgroup

\begingroup
\setbox0=\vbox{\hsize=\linewidth
{\normalfont\normalsize\bfseries Restricciones físicas observadas\par}
\begin{itemize}
\item Clave foránea: captured\_\allowbreak{}by\_\allowbreak{}user\_\allowbreak{}id hacia app.users (id); al eliminar: referencia nula.
\end{itemize}
}
\Needspace{\dimexpr\ht0+\dp0+24pt\relax}
\endgroup
\begingroup
\setbox0=\vbox{\hsize=\linewidth
{\normalfont\normalsize\bfseries Restricciones físicas observadas\par}
\begin{itemize}
\item Clave foránea: captured\_\allowbreak{}by\_\allowbreak{}user\_\allowbreak{}id hacia app.users (id); al eliminar: referencia nula.
\end{itemize}
}
\Needspace{\dimexpr\ht0+\dp0+24pt\relax}
\endgroup
\paragraph{Restricciones físicas observadas} \begin{itemize}
\item Clave foránea: captured\_\allowbreak{}by\_\allowbreak{}user\_\allowbreak{}id hacia app.users (id); al eliminar: referencia nula.
\item Clave foránea: data\_\allowbreak{}connection\_\allowbreak{}id hacia app.data\_\allowbreak{}connections (id); al eliminar: restricción.
\item Clave primaria: id.
\item Unicidad: data\_\allowbreak{}connection\_\allowbreak{}id, contract\_\allowbreak{}version, content\_\allowbreak{}hash.
\end{itemize}

\begingroup
\setbox0=\vbox{\hsize=\linewidth
{\normalfont\normalsize\bfseries Índices únicos\par}
\begin{itemize}
\item id; índice completo.
\end{itemize}
}
\Needspace{\dimexpr\ht0+\dp0+24pt\relax}
\endgroup
\begingroup
\setbox0=\vbox{\hsize=\linewidth
{\normalfont\normalsize\bfseries Índices únicos\par}
\begin{itemize}
\item id; índice completo.
\end{itemize}
}
\Needspace{\dimexpr\ht0+\dp0+24pt\relax}
\endgroup
\paragraph{Índices únicos} \begin{itemize}
\item id; índice completo.
\item data\_\allowbreak{}connection\_\allowbreak{}id, contract\_\allowbreak{}version, content\_\allowbreak{}hash; índice completo.
\end{itemize}

\begingroup
\setbox0=\vbox{\hsize=\linewidth
{\normalfont\normalsize\sffamily\bfseries\normalsize\color{Teal}\hyphenpenalty=10000\exhyphenpenalty=10000 app.parameters\par}

\par\vspace{6pt}\noindent\textbf{Área}\hspace{0.5em} Configuración y metadatos

\par\vspace{6pt}\noindent\textbf{Qué guarda}\hspace{0.5em} Conserva parámetros del producto, su tipo y límites configurados.

\par\vspace{6pt}\noindent\textbf{Cómo se utiliza}\hspace{0.5em} Pantalla Parámetros. value se conserva como texto; el servidor interpreta el tipo y valida los límites.

\par
\begingroup
\fontsize{10}{12.5}\selectfont
\setlength{\tabcolsep}{5pt}
\renewcommand{\arraystretch}{1.15}
\begin{tabular}{@{}>{\RaggedRight\arraybackslash}p{0.23\linewidth}>{\RaggedRight\arraybackslash}p{0.23\linewidth}>{\RaggedRight\arraybackslash}p{0.45\linewidth}@{}}
\toprule
\textbf{Columna} & \textbf{Tipo y condición} & \textbf{Qué representa} \\
\midrule
\texttt{\seqsplit{id}} & integer\par\smallskip Obligatorio. Clave primaria. Con valor predeterminado. & Identificador interno de esta entidad; su PK distingue cada registro. \\
\bottomrule
\end{tabular}
\endgroup
}
\Needspace{\dimexpr\ht0+\dp0+48pt\relax}
\endgroup
\subsection{app.parameters}
\paragraph{Área} Configuración y metadatos

\paragraph{Qué guarda} Conserva parámetros del producto, su tipo y límites configurados.

\paragraph{Cómo se utiliza} Pantalla Parámetros. value se conserva como texto; el servidor interpreta el tipo y valida los límites.

\paragraph{Ejemplo vigente} Duración de la sesión (minutos) utiliza esta tabla existente. La clave \texttt{AUTH\_SESSION\_MINUTES} identifica una fila; su valor inicial es 60 y sus límites son 15 y 240. La operación se realiza desde la pantalla, sin añadir columnas ni editar valores por SQL. Este ejemplo es una configuración técnica del acceso y no una medida de ventas.

\begingroup
\fontsize{10}{12.5}\selectfont
\setlength{\tabcolsep}{5pt}
\renewcommand{\arraystretch}{1.15}
\begin{longtable}{@{}>{\RaggedRight\arraybackslash}p{0.23\linewidth}>{\RaggedRight\arraybackslash}p{0.23\linewidth}>{\RaggedRight\arraybackslash}p{0.45\linewidth}@{}}
\toprule
\textbf{Columna} & \textbf{Tipo y condición} & \textbf{Qué representa} \\
\midrule
\endfirsthead
\toprule
\textbf{Columna} & \textbf{Tipo y condición} & \textbf{Qué representa} \\
\midrule
\endhead
\midrule
\multicolumn{3}{r}{\footnotesize Continúa en la página siguiente}\\
\endfoot
\bottomrule
\endlastfoot
\texttt{\seqsplit{id}} & integer\par\smallskip Obligatorio. Clave primaria. Con valor predeterminado. & Identificador interno de esta entidad; su PK distingue cada registro. \\
\addlinespace[3pt]
\texttt{\seqsplit{key}} & character varying(100)\par\smallskip Obligatorio. & Clave única de un parámetro del producto. \\
\addlinespace[3pt]
\texttt{\seqsplit{value}} & text\par\smallskip Obligatorio. & Valor del parámetro guardado como texto; no se extrajo su contenido. \\
\addlinespace[3pt]
\texttt{\seqsplit{description}} & text\par\smallskip Admite nulos. & Explicación de la entidad o parámetro. \\
\addlinespace[3pt]
\texttt{\seqsplit{is\_active}} & boolean\par\smallskip Obligatorio. Con valor predeterminado. & Estado de habilitación del registro; su efecto depende de las reglas del módulo. \\
\addlinespace[3pt]
\texttt{\seqsplit{updated\_at}} & timestamp with time zone\par\smallskip Admite nulos. Con valor predeterminado. & Fecha de actualización del registro; su opcionalidad física consta en el inventario. \\
\addlinespace[3pt]
\texttt{\seqsplit{name}} & character varying(160)\par\smallskip Obligatorio. Con valor predeterminado. & Nombre o denominación visible en el módulo. \\
\addlinespace[3pt]
\texttt{\seqsplit{module\_code}} & character varying(80)\par\smallskip Obligatorio. Con valor predeterminado. & Código del módulo o agrupación funcional. \\
\addlinespace[3pt]
\texttt{\seqsplit{value\_type}} & character varying(30)\par\smallskip Obligatorio. Con valor predeterminado. & Tipo lógico con el que el servidor interpreta el valor textual. \\
\addlinespace[3pt]
\texttt{\seqsplit{default\_value}} & text\par\smallskip Obligatorio. Con valor predeterminado. & Valor previsto por defecto para el parámetro; no se extrajo su contenido. \\
\addlinespace[3pt]
\texttt{\seqsplit{min\_value}} & integer\par\smallskip Admite nulos. & Límite inferior configurado para valores numéricos compatibles. \\
\addlinespace[3pt]
\texttt{\seqsplit{max\_value}} & integer\par\smallskip Admite nulos. & Límite superior configurado para valores numéricos compatibles. \\
\addlinespace[3pt]
\end{longtable}
\endgroup

\begingroup
\setbox0=\vbox{\hsize=\linewidth
{\normalfont\normalsize\bfseries Restricciones físicas observadas\par}
\begin{itemize}
\item Clave primaria: id.
\end{itemize}
}
\Needspace{\dimexpr\ht0+\dp0+24pt\relax}
\endgroup
\begingroup
\setbox0=\vbox{\hsize=\linewidth
{\normalfont\normalsize\bfseries Restricciones físicas observadas\par}
\begin{itemize}
\item Clave primaria: id.
\end{itemize}
}
\Needspace{\dimexpr\ht0+\dp0+24pt\relax}
\endgroup
\paragraph{Restricciones físicas observadas} \begin{itemize}
\item Clave primaria: id.
\item Unicidad: key.
\end{itemize}

\begingroup
\setbox0=\vbox{\hsize=\linewidth
{\normalfont\normalsize\bfseries Índices únicos\par}
\begin{itemize}
\item key; índice completo.
\end{itemize}
}
\Needspace{\dimexpr\ht0+\dp0+24pt\relax}
\endgroup
\begingroup
\setbox0=\vbox{\hsize=\linewidth
{\normalfont\normalsize\bfseries Índices únicos\par}
\begin{itemize}
\item key; índice completo.
\end{itemize}
}
\Needspace{\dimexpr\ht0+\dp0+24pt\relax}
\endgroup
\paragraph{Índices únicos} \begin{itemize}
\item key; índice completo.
\item id; índice completo.
\item key; índice completo.
\end{itemize}

\begingroup
\setbox0=\vbox{\hsize=\linewidth
{\normalfont\normalsize\sffamily\bfseries\normalsize\color{Teal}\hyphenpenalty=10000\exhyphenpenalty=10000 app.permissions\par}

\par\vspace{6pt}\noindent\textbf{Área}\hspace{0.5em} Seguridad

\par\vspace{6pt}\noindent\textbf{Qué guarda}\hspace{0.5em} Catálogo de operaciones que el servidor autoriza.

\par\vspace{6pt}\noindent\textbf{Cómo se utiliza}\hspace{0.5em} Permisos de lectura, escritura, revisión o ejecución según el módulo. Un menú visible no concede el permiso.

\par
\begingroup
\fontsize{10}{12.5}\selectfont
\setlength{\tabcolsep}{5pt}
\renewcommand{\arraystretch}{1.15}
\begin{tabular}{@{}>{\RaggedRight\arraybackslash}p{0.23\linewidth}>{\RaggedRight\arraybackslash}p{0.23\linewidth}>{\RaggedRight\arraybackslash}p{0.45\linewidth}@{}}
\toprule
\textbf{Columna} & \textbf{Tipo y condición} & \textbf{Qué representa} \\
\midrule
\texttt{\seqsplit{id}} & integer\par\smallskip Obligatorio. Clave primaria. Con valor predeterminado. & Identificador interno de esta entidad; su PK distingue cada registro. \\
\bottomrule
\end{tabular}
\endgroup
}
\Needspace{\dimexpr\ht0+\dp0+48pt\relax}
\endgroup
\subsection{app.permissions}
\paragraph{Área} Seguridad

\paragraph{Qué guarda} Catálogo de operaciones que el servidor autoriza.

\paragraph{Cómo se utiliza} Permisos de lectura, escritura, revisión o ejecución según el módulo. Un menú visible no concede el permiso.

\begingroup
\fontsize{10}{12.5}\selectfont
\setlength{\tabcolsep}{5pt}
\renewcommand{\arraystretch}{1.15}
\begin{longtable}{@{}>{\RaggedRight\arraybackslash}p{0.23\linewidth}>{\RaggedRight\arraybackslash}p{0.23\linewidth}>{\RaggedRight\arraybackslash}p{0.45\linewidth}@{}}
\toprule
\textbf{Columna} & \textbf{Tipo y condición} & \textbf{Qué representa} \\
\midrule
\endfirsthead
\toprule
\textbf{Columna} & \textbf{Tipo y condición} & \textbf{Qué representa} \\
\midrule
\endhead
\midrule
\multicolumn{3}{r}{\footnotesize Continúa en la página siguiente}\\
\endfoot
\bottomrule
\endlastfoot
\texttt{\seqsplit{id}} & integer\par\smallskip Obligatorio. Clave primaria. Con valor predeterminado. & Identificador interno de esta entidad; su PK distingue cada registro. \\
\addlinespace[3pt]
\texttt{\seqsplit{code}} & character varying(120)\par\smallskip Obligatorio. & Código estable de la entidad; no depende de su etiqueta visible. \\
\addlinespace[3pt]
\texttt{\seqsplit{name}} & character varying(160)\par\smallskip Obligatorio. & Nombre o denominación visible en el módulo. \\
\addlinespace[3pt]
\texttt{\seqsplit{description}} & text\par\smallskip Admite nulos. & Explicación de la entidad o parámetro. \\
\addlinespace[3pt]
\texttt{\seqsplit{is\_active}} & boolean\par\smallskip Obligatorio. Con valor predeterminado. & Estado de habilitación del registro; su efecto depende de las reglas del módulo. \\
\addlinespace[3pt]
\texttt{\seqsplit{is\_system\_protected}} & boolean\par\smallskip Obligatorio. Con valor predeterminado. & Marca de un registro base protegido para conservar administración y navegación esenciales. \\
\addlinespace[3pt]
\end{longtable}
\endgroup

\begingroup
\setbox0=\vbox{\hsize=\linewidth
{\normalfont\normalsize\bfseries Restricciones físicas observadas\par}
\begin{itemize}
\item Clave primaria: id.
\end{itemize}
}
\Needspace{\dimexpr\ht0+\dp0+24pt\relax}
\endgroup
\begingroup
\setbox0=\vbox{\hsize=\linewidth
{\normalfont\normalsize\bfseries Restricciones físicas observadas\par}
\begin{itemize}
\item Clave primaria: id.
\end{itemize}
}
\Needspace{\dimexpr\ht0+\dp0+24pt\relax}
\endgroup
\paragraph{Restricciones físicas observadas} \begin{itemize}
\item Clave primaria: id.
\item Unicidad: code.
\end{itemize}

\begingroup
\setbox0=\vbox{\hsize=\linewidth
{\normalfont\normalsize\bfseries Índices únicos\par}
\begin{itemize}
\item code; índice completo.
\end{itemize}
}
\Needspace{\dimexpr\ht0+\dp0+24pt\relax}
\endgroup
\begingroup
\setbox0=\vbox{\hsize=\linewidth
{\normalfont\normalsize\bfseries Índices únicos\par}
\begin{itemize}
\item code; índice completo.
\end{itemize}
}
\Needspace{\dimexpr\ht0+\dp0+24pt\relax}
\endgroup
\paragraph{Índices únicos} \begin{itemize}
\item code; índice completo.
\item id; índice completo.
\item code; índice completo.
\end{itemize}

\begingroup
\setbox0=\vbox{\hsize=\linewidth
{\normalfont\normalsize\sffamily\bfseries\normalsize\color{Teal}\hyphenpenalty=10000\exhyphenpenalty=10000 app.role\_\allowbreak{}permissions\par}

\par\vspace{6pt}\noindent\textbf{Área}\hspace{0.5em} Seguridad

\par\vspace{6pt}\noindent\textbf{Qué guarda}\hspace{0.5em} Relaciona roles con permisos mediante dos FK y una PK compuesta.

\par\vspace{6pt}\noindent\textbf{Cómo se utiliza}\hspace{0.5em} Conceder a un rol las operaciones autorizadas.

\par
\begingroup
\fontsize{10}{12.5}\selectfont
\setlength{\tabcolsep}{5pt}
\renewcommand{\arraystretch}{1.15}
\begin{tabular}{@{}>{\RaggedRight\arraybackslash}p{0.23\linewidth}>{\RaggedRight\arraybackslash}p{0.23\linewidth}>{\RaggedRight\arraybackslash}p{0.45\linewidth}@{}}
\toprule
\textbf{Columna} & \textbf{Tipo y condición} & \textbf{Qué representa} \\
\midrule
\texttt{\seqsplit{role\_id}} & integer\par\smallskip Obligatorio. Clave primaria. & FK a app.roles.id; identifica un rol en la tabla de asociación. \\
\bottomrule
\end{tabular}
\endgroup
}
\Needspace{\dimexpr\ht0+\dp0+48pt\relax}
\endgroup
\subsection{app.role\_\allowbreak{}permissions}
\paragraph{Área} Seguridad

\paragraph{Qué guarda} Relaciona roles con permisos mediante dos FK y una PK compuesta.

\paragraph{Cómo se utiliza} Conceder a un rol las operaciones autorizadas.

\begingroup
\fontsize{10}{12.5}\selectfont
\setlength{\tabcolsep}{5pt}
\renewcommand{\arraystretch}{1.15}
\begin{longtable}{@{}>{\RaggedRight\arraybackslash}p{0.23\linewidth}>{\RaggedRight\arraybackslash}p{0.23\linewidth}>{\RaggedRight\arraybackslash}p{0.45\linewidth}@{}}
\toprule
\textbf{Columna} & \textbf{Tipo y condición} & \textbf{Qué representa} \\
\midrule
\endfirsthead
\toprule
\textbf{Columna} & \textbf{Tipo y condición} & \textbf{Qué representa} \\
\midrule
\endhead
\midrule
\multicolumn{3}{r}{\footnotesize Continúa en la página siguiente}\\
\endfoot
\bottomrule
\endlastfoot
\texttt{\seqsplit{role\_id}} & integer\par\smallskip Obligatorio. Clave primaria. & FK a app.roles.id; identifica un rol en la tabla de asociación. \\
\addlinespace[3pt]
\texttt{\seqsplit{permission\_id}} & integer\par\smallskip Obligatorio. Clave primaria. & FK a app.permissions.id; identifica una operación autorizable. \\
\addlinespace[3pt]
\end{longtable}
\endgroup

\begingroup
\setbox0=\vbox{\hsize=\linewidth
{\normalfont\normalsize\bfseries Restricciones físicas observadas\par}
\begin{itemize}
\item Clave foránea: permission\_\allowbreak{}id hacia app.permissions (id); al eliminar: cascada SQL.
\end{itemize}
}
\Needspace{\dimexpr\ht0+\dp0+24pt\relax}
\endgroup
\begingroup
\setbox0=\vbox{\hsize=\linewidth
{\normalfont\normalsize\bfseries Restricciones físicas observadas\par}
\begin{itemize}
\item Clave foránea: permission\_\allowbreak{}id hacia app.permissions (id); al eliminar: cascada SQL.
\end{itemize}
}
\Needspace{\dimexpr\ht0+\dp0+24pt\relax}
\endgroup
\paragraph{Restricciones físicas observadas} \begin{itemize}
\item Clave foránea: permission\_\allowbreak{}id hacia app.permissions (id); al eliminar: cascada SQL.
\item Clave foránea: role\_\allowbreak{}id hacia app.roles (id); al eliminar: cascada SQL.
\item Clave primaria: role\_\allowbreak{}id, permission\_\allowbreak{}id.
\end{itemize}

\begingroup
\setbox0=\vbox{\hsize=\linewidth
{\normalfont\normalsize\bfseries Índices únicos\par}
\begin{itemize}
\item role\_\allowbreak{}id, permission\_\allowbreak{}id; índice completo.
\end{itemize}
}
\Needspace{\dimexpr\ht0+\dp0+24pt\relax}
\endgroup
\begingroup
\setbox0=\vbox{\hsize=\linewidth
{\normalfont\normalsize\bfseries Índices únicos\par}
\begin{itemize}
\item role\_\allowbreak{}id, permission\_\allowbreak{}id; índice completo.
\end{itemize}
}
\Needspace{\dimexpr\ht0+\dp0+24pt\relax}
\endgroup
\paragraph{Índices únicos} \begin{itemize}
\item role\_\allowbreak{}id, permission\_\allowbreak{}id; índice completo.
\end{itemize}

\begingroup
\setbox0=\vbox{\hsize=\linewidth
{\normalfont\normalsize\sffamily\bfseries\normalsize\color{Teal}\hyphenpenalty=10000\exhyphenpenalty=10000 app.roles\par}

\par\vspace{6pt}\noindent\textbf{Área}\hspace{0.5em} Seguridad

\par\vspace{6pt}\noindent\textbf{Qué guarda}\hspace{0.5em} Agrupa responsabilidades de acceso, como administración, análisis o gerencia.

\par\vspace{6pt}\noindent\textbf{Cómo se utiliza}\hspace{0.5em} Administración de roles; el nombre visible no sustituye los permisos asignados.

\par
\begingroup
\fontsize{10}{12.5}\selectfont
\setlength{\tabcolsep}{5pt}
\renewcommand{\arraystretch}{1.15}
\begin{tabular}{@{}>{\RaggedRight\arraybackslash}p{0.23\linewidth}>{\RaggedRight\arraybackslash}p{0.23\linewidth}>{\RaggedRight\arraybackslash}p{0.45\linewidth}@{}}
\toprule
\textbf{Columna} & \textbf{Tipo y condición} & \textbf{Qué representa} \\
\midrule
\texttt{\seqsplit{id}} & integer\par\smallskip Obligatorio. Clave primaria. Con valor predeterminado. & Identificador interno de esta entidad; su PK distingue cada registro. \\
\bottomrule
\end{tabular}
\endgroup
}
\Needspace{\dimexpr\ht0+\dp0+48pt\relax}
\endgroup
\subsection{app.roles}
\paragraph{Área} Seguridad

\paragraph{Qué guarda} Agrupa responsabilidades de acceso, como administración, análisis o gerencia.

\paragraph{Cómo se utiliza} Administración de roles; el nombre visible no sustituye los permisos asignados.

\begingroup
\fontsize{10}{12.5}\selectfont
\setlength{\tabcolsep}{5pt}
\renewcommand{\arraystretch}{1.15}
\begin{longtable}{@{}>{\RaggedRight\arraybackslash}p{0.23\linewidth}>{\RaggedRight\arraybackslash}p{0.23\linewidth}>{\RaggedRight\arraybackslash}p{0.45\linewidth}@{}}
\toprule
\textbf{Columna} & \textbf{Tipo y condición} & \textbf{Qué representa} \\
\midrule
\endfirsthead
\toprule
\textbf{Columna} & \textbf{Tipo y condición} & \textbf{Qué representa} \\
\midrule
\endhead
\midrule
\multicolumn{3}{r}{\footnotesize Continúa en la página siguiente}\\
\endfoot
\bottomrule
\endlastfoot
\texttt{\seqsplit{id}} & integer\par\smallskip Obligatorio. Clave primaria. Con valor predeterminado. & Identificador interno de esta entidad; su PK distingue cada registro. \\
\addlinespace[3pt]
\texttt{\seqsplit{code}} & character varying(80)\par\smallskip Obligatorio. & Código estable de la entidad; no depende de su etiqueta visible. \\
\addlinespace[3pt]
\texttt{\seqsplit{name}} & character varying(120)\par\smallskip Obligatorio. & Nombre o denominación visible en el módulo. \\
\addlinespace[3pt]
\texttt{\seqsplit{description}} & text\par\smallskip Admite nulos. & Explicación de la entidad o parámetro. \\
\addlinespace[3pt]
\texttt{\seqsplit{is\_active}} & boolean\par\smallskip Obligatorio. Con valor predeterminado. & Estado de habilitación del registro; su efecto depende de las reglas del módulo. \\
\addlinespace[3pt]
\texttt{\seqsplit{is\_system\_protected}} & boolean\par\smallskip Obligatorio. Con valor predeterminado. & Marca de un registro base protegido para conservar administración y navegación esenciales. \\
\addlinespace[3pt]
\end{longtable}
\endgroup

\begingroup
\setbox0=\vbox{\hsize=\linewidth
{\normalfont\normalsize\bfseries Restricciones físicas observadas\par}
\begin{itemize}
\item Clave primaria: id.
\end{itemize}
}
\Needspace{\dimexpr\ht0+\dp0+24pt\relax}
\endgroup
\begingroup
\setbox0=\vbox{\hsize=\linewidth
{\normalfont\normalsize\bfseries Restricciones físicas observadas\par}
\begin{itemize}
\item Clave primaria: id.
\end{itemize}
}
\Needspace{\dimexpr\ht0+\dp0+24pt\relax}
\endgroup
\paragraph{Restricciones físicas observadas} \begin{itemize}
\item Clave primaria: id.
\item Unicidad: code.
\item Unicidad: name.
\end{itemize}

\begingroup
\setbox0=\vbox{\hsize=\linewidth
{\normalfont\normalsize\bfseries Índices únicos\par}
\begin{itemize}
\item code; índice completo.
\end{itemize}
}
\Needspace{\dimexpr\ht0+\dp0+24pt\relax}
\endgroup
\begingroup
\setbox0=\vbox{\hsize=\linewidth
{\normalfont\normalsize\bfseries Índices únicos\par}
\begin{itemize}
\item code; índice completo.
\end{itemize}
}
\Needspace{\dimexpr\ht0+\dp0+24pt\relax}
\endgroup
\paragraph{Índices únicos} \begin{itemize}
\item code; índice completo.
\item id; índice completo.
\item code; índice completo.
\item name; índice completo.
\end{itemize}

\begingroup
\setbox0=\vbox{\hsize=\linewidth
{\normalfont\normalsize\sffamily\bfseries\normalsize\color{Teal}\hyphenpenalty=10000\exhyphenpenalty=10000 app.secrets\par}

\par\vspace{6pt}\noindent\textbf{Área}\hspace{0.5em} Configuración y metadatos

\par\vspace{6pt}\noindent\textbf{Qué guarda}\hspace{0.5em} Conserva valores cifrados de credenciales mediante ciphertext.

\par\vspace{6pt}\noindent\textbf{Cómo se utiliza}\hspace{0.5em} No se consultan sus valores desde la interfaz ni desde este inventario. La raíz de cifrado está en un volumen separado de PostgreSQL.

\par
\begingroup
\fontsize{10}{12.5}\selectfont
\setlength{\tabcolsep}{5pt}
\renewcommand{\arraystretch}{1.15}
\begin{tabular}{@{}>{\RaggedRight\arraybackslash}p{0.23\linewidth}>{\RaggedRight\arraybackslash}p{0.23\linewidth}>{\RaggedRight\arraybackslash}p{0.45\linewidth}@{}}
\toprule
\textbf{Columna} & \textbf{Tipo y condición} & \textbf{Qué representa} \\
\midrule
\texttt{\seqsplit{id}} & integer\par\smallskip Obligatorio. Clave primaria. Con valor predeterminado. & Identificador interno de esta entidad; su PK distingue cada registro. \\
\bottomrule
\end{tabular}
\endgroup
}
\Needspace{\dimexpr\ht0+\dp0+48pt\relax}
\endgroup
\subsection{app.secrets}
\paragraph{Área} Configuración y metadatos

\paragraph{Qué guarda} Conserva valores cifrados de credenciales mediante ciphertext.

\paragraph{Cómo se utiliza} No se consultan sus valores desde la interfaz ni desde este inventario. La raíz de cifrado está en un volumen separado de PostgreSQL.

\begingroup
\fontsize{10}{12.5}\selectfont
\setlength{\tabcolsep}{5pt}
\renewcommand{\arraystretch}{1.15}
\begin{longtable}{@{}>{\RaggedRight\arraybackslash}p{0.23\linewidth}>{\RaggedRight\arraybackslash}p{0.23\linewidth}>{\RaggedRight\arraybackslash}p{0.45\linewidth}@{}}
\toprule
\textbf{Columna} & \textbf{Tipo y condición} & \textbf{Qué representa} \\
\midrule
\endfirsthead
\toprule
\textbf{Columna} & \textbf{Tipo y condición} & \textbf{Qué representa} \\
\midrule
\endhead
\midrule
\multicolumn{3}{r}{\footnotesize Continúa en la página siguiente}\\
\endfoot
\bottomrule
\endlastfoot
\texttt{\seqsplit{id}} & integer\par\smallskip Obligatorio. Clave primaria. Con valor predeterminado. & Identificador interno de esta entidad; su PK distingue cada registro. \\
\addlinespace[3pt]
\texttt{\seqsplit{kind}} & character varying(40)\par\smallskip Obligatorio. & Clase de credencial cifrada. \\
\addlinespace[3pt]
\texttt{\seqsplit{ciphertext}} & text\par\smallskip Obligatorio. & Contenido cifrado del secreto; no se leyó ni se incluye en los documentos. \\
\addlinespace[3pt]
\texttt{\seqsplit{created\_at}} & timestamp with time zone\par\smallskip Admite nulos. Con valor predeterminado. & Fecha de creación del registro; su opcionalidad física consta en el inventario. \\
\addlinespace[3pt]
\texttt{\seqsplit{updated\_at}} & timestamp with time zone\par\smallskip Admite nulos. Con valor predeterminado. & Fecha de actualización del registro; su opcionalidad física consta en el inventario. \\
\addlinespace[3pt]
\end{longtable}
\endgroup

\begingroup
\setbox0=\vbox{\hsize=\linewidth
{\normalfont\normalsize\bfseries Restricciones físicas observadas\par}
\begin{itemize}
\item Clave primaria: id.
\end{itemize}
}
\Needspace{\dimexpr\ht0+\dp0+24pt\relax}
\endgroup
\begingroup
\setbox0=\vbox{\hsize=\linewidth
{\normalfont\normalsize\bfseries Restricciones físicas observadas\par}
\begin{itemize}
\item Clave primaria: id.
\end{itemize}
}
\Needspace{\dimexpr\ht0+\dp0+24pt\relax}
\endgroup
\paragraph{Restricciones físicas observadas} \begin{itemize}
\item Clave primaria: id.
\end{itemize}

\begingroup
\setbox0=\vbox{\hsize=\linewidth
{\normalfont\normalsize\bfseries Índices únicos\par}
\begin{itemize}
\item id; índice completo.
\end{itemize}
}
\Needspace{\dimexpr\ht0+\dp0+24pt\relax}
\endgroup
\begingroup
\setbox0=\vbox{\hsize=\linewidth
{\normalfont\normalsize\bfseries Índices únicos\par}
\begin{itemize}
\item id; índice completo.
\end{itemize}
}
\Needspace{\dimexpr\ht0+\dp0+24pt\relax}
\endgroup
\paragraph{Índices únicos} \begin{itemize}
\item id; índice completo.
\end{itemize}

\begingroup
\setbox0=\vbox{\hsize=\linewidth
{\normalfont\normalsize\sffamily\bfseries\normalsize\color{Teal}\hyphenpenalty=10000\exhyphenpenalty=10000 app.semantic\_\allowbreak{}advice\par}

\par\vspace{6pt}\noindent\textbf{Área}\hspace{0.5em} Propuestas y ejecución

\par\vspace{6pt}\noindent\textbf{Qué guarda}\hspace{0.5em} Conserva asesoría contextual sobre un concepto de una propuesta.

\par\vspace{6pt}\noindent\textbf{Cómo se utiliza}\hspace{0.5em} Ayuda para una decisión semántica del diseño. No es una tabla de conversaciones completas del copiloto analítico.

\par
\begingroup
\fontsize{10}{12.5}\selectfont
\setlength{\tabcolsep}{5pt}
\renewcommand{\arraystretch}{1.15}
\begin{tabular}{@{}>{\RaggedRight\arraybackslash}p{0.23\linewidth}>{\RaggedRight\arraybackslash}p{0.23\linewidth}>{\RaggedRight\arraybackslash}p{0.45\linewidth}@{}}
\toprule
\textbf{Columna} & \textbf{Tipo y condición} & \textbf{Qué representa} \\
\midrule
\texttt{\seqsplit{id}} & integer\par\smallskip Obligatorio. Clave primaria. Con valor predeterminado. & Identificador interno de esta entidad; su PK distingue cada registro. \\
\bottomrule
\end{tabular}
\endgroup
}
\Needspace{\dimexpr\ht0+\dp0+48pt\relax}
\endgroup
\subsection{app.semantic\_\allowbreak{}advice}
\paragraph{Área} Propuestas y ejecución

\paragraph{Qué guarda} Conserva asesoría contextual sobre un concepto de una propuesta.

\paragraph{Cómo se utiliza} Ayuda para una decisión semántica del diseño. No es una tabla de conversaciones completas del copiloto analítico.

\begingroup
\fontsize{10}{12.5}\selectfont
\setlength{\tabcolsep}{5pt}
\renewcommand{\arraystretch}{1.15}
\begin{longtable}{@{}>{\RaggedRight\arraybackslash}p{0.23\linewidth}>{\RaggedRight\arraybackslash}p{0.23\linewidth}>{\RaggedRight\arraybackslash}p{0.45\linewidth}@{}}
\toprule
\textbf{Columna} & \textbf{Tipo y condición} & \textbf{Qué representa} \\
\midrule
\endfirsthead
\toprule
\textbf{Columna} & \textbf{Tipo y condición} & \textbf{Qué representa} \\
\midrule
\endhead
\midrule
\multicolumn{3}{r}{\footnotesize Continúa en la página siguiente}\\
\endfoot
\bottomrule
\endlastfoot
\texttt{\seqsplit{id}} & integer\par\smallskip Obligatorio. Clave primaria. Con valor predeterminado. & Identificador interno de esta entidad; su PK distingue cada registro. \\
\addlinespace[3pt]
\texttt{\seqsplit{proposal\_id}} & integer\par\smallskip Obligatorio. & FK a app.bi\_\allowbreak{}proposals.id; vincula la asesoría o ejecución con una propuesta concreta. \\
\addlinespace[3pt]
\texttt{\seqsplit{concept\_code}} & character varying(80)\par\smallskip Obligatorio. & Identificador del concepto sobre el que se pidió asesoría. \\
\addlinespace[3pt]
\texttt{\seqsplit{question}} & text\par\smallskip Obligatorio. & Consulta contextual sobre un concepto del diseño. \\
\addlinespace[3pt]
\texttt{\seqsplit{response\_document}} & jsonb\par\smallskip Obligatorio. & JSONB de la asesoría semántica; no se leyó su contenido. \\
\addlinespace[3pt]
\texttt{\seqsplit{provider\_kind}} & character varying(40)\par\smallskip Obligatorio. & Tipo de proveedor de IA. En los expedientes conserva la procedencia como texto, sin FK a la configuración que después puede cambiar. \\
\addlinespace[3pt]
\texttt{\seqsplit{model\_id}} & character varying(160)\par\smallskip Obligatorio. & Identificador del modelo utilizado o configurado, distinto de una API key. \\
\addlinespace[3pt]
\texttt{\seqsplit{created\_by\_user\_id}} & integer\par\smallskip Admite nulos. & FK opcional a la persona que creó el artefacto o expediente. \\
\addlinespace[3pt]
\texttt{\seqsplit{created\_by\_label}} & character varying(320)\par\smallskip Obligatorio. & Etiqueta histórica de la persona creadora. \\
\addlinespace[3pt]
\texttt{\seqsplit{created\_at}} & timestamp with time zone\par\smallskip Obligatorio. Con valor predeterminado. & Fecha de creación del registro; su opcionalidad física consta en el inventario. \\
\addlinespace[3pt]
\end{longtable}
\endgroup

\begingroup
\setbox0=\vbox{\hsize=\linewidth
{\normalfont\normalsize\bfseries Restricciones físicas observadas\par}
\begin{itemize}
\item Clave foránea: created\_\allowbreak{}by\_\allowbreak{}user\_\allowbreak{}id hacia app.users (id); al eliminar: referencia nula.
\end{itemize}
}
\Needspace{\dimexpr\ht0+\dp0+24pt\relax}
\endgroup
\begingroup
\setbox0=\vbox{\hsize=\linewidth
{\normalfont\normalsize\bfseries Restricciones físicas observadas\par}
\begin{itemize}
\item Clave foránea: created\_\allowbreak{}by\_\allowbreak{}user\_\allowbreak{}id hacia app.users (id); al eliminar: referencia nula.
\end{itemize}
}
\Needspace{\dimexpr\ht0+\dp0+24pt\relax}
\endgroup
\paragraph{Restricciones físicas observadas} \begin{itemize}
\item Clave foránea: created\_\allowbreak{}by\_\allowbreak{}user\_\allowbreak{}id hacia app.users (id); al eliminar: referencia nula.
\item Clave foránea: proposal\_\allowbreak{}id hacia app.bi\_\allowbreak{}proposals (id); al eliminar: cascada SQL.
\item Clave primaria: id.
\end{itemize}

\begingroup
\setbox0=\vbox{\hsize=\linewidth
{\normalfont\normalsize\bfseries Índices únicos\par}
\begin{itemize}
\item id; índice completo.
\end{itemize}
}
\Needspace{\dimexpr\ht0+\dp0+24pt\relax}
\endgroup
\begingroup
\setbox0=\vbox{\hsize=\linewidth
{\normalfont\normalsize\bfseries Índices únicos\par}
\begin{itemize}
\item id; índice completo.
\end{itemize}
}
\Needspace{\dimexpr\ht0+\dp0+24pt\relax}
\endgroup
\paragraph{Índices únicos} \begin{itemize}
\item id; índice completo.
\end{itemize}

\begingroup
\setbox0=\vbox{\hsize=\linewidth
{\normalfont\normalsize\sffamily\bfseries\normalsize\color{Teal}\hyphenpenalty=10000\exhyphenpenalty=10000 app.user\_\allowbreak{}roles\par}

\par\vspace{6pt}\noindent\textbf{Área}\hspace{0.5em} Seguridad

\par\vspace{6pt}\noindent\textbf{Qué guarda}\hspace{0.5em} Relaciona cuentas con roles mediante dos FK y una PK compuesta.

\par\vspace{6pt}\noindent\textbf{Cómo se utiliza}\hspace{0.5em} Asignar uno o varios roles a una cuenta; un rol puede pertenecer a varias cuentas.

\par
\begingroup
\fontsize{10}{12.5}\selectfont
\setlength{\tabcolsep}{5pt}
\renewcommand{\arraystretch}{1.15}
\begin{tabular}{@{}>{\RaggedRight\arraybackslash}p{0.23\linewidth}>{\RaggedRight\arraybackslash}p{0.23\linewidth}>{\RaggedRight\arraybackslash}p{0.45\linewidth}@{}}
\toprule
\textbf{Columna} & \textbf{Tipo y condición} & \textbf{Qué representa} \\
\midrule
\texttt{\seqsplit{user\_id}} & integer\par\smallskip Obligatorio. Clave primaria. & FK a app.users.id; identifica una cuenta en la tabla de asociación. \\
\bottomrule
\end{tabular}
\endgroup
}
\Needspace{\dimexpr\ht0+\dp0+48pt\relax}
\endgroup
\subsection{app.user\_\allowbreak{}roles}
\paragraph{Área} Seguridad

\paragraph{Qué guarda} Relaciona cuentas con roles mediante dos FK y una PK compuesta.

\paragraph{Cómo se utiliza} Asignar uno o varios roles a una cuenta; un rol puede pertenecer a varias cuentas.

\begingroup
\fontsize{10}{12.5}\selectfont
\setlength{\tabcolsep}{5pt}
\renewcommand{\arraystretch}{1.15}
\begin{longtable}{@{}>{\RaggedRight\arraybackslash}p{0.23\linewidth}>{\RaggedRight\arraybackslash}p{0.23\linewidth}>{\RaggedRight\arraybackslash}p{0.45\linewidth}@{}}
\toprule
\textbf{Columna} & \textbf{Tipo y condición} & \textbf{Qué representa} \\
\midrule
\endfirsthead
\toprule
\textbf{Columna} & \textbf{Tipo y condición} & \textbf{Qué representa} \\
\midrule
\endhead
\midrule
\multicolumn{3}{r}{\footnotesize Continúa en la página siguiente}\\
\endfoot
\bottomrule
\endlastfoot
\texttt{\seqsplit{user\_id}} & integer\par\smallskip Obligatorio. Clave primaria. & FK a app.users.id; identifica una cuenta en la tabla de asociación. \\
\addlinespace[3pt]
\texttt{\seqsplit{role\_id}} & integer\par\smallskip Obligatorio. Clave primaria. & FK a app.roles.id; identifica un rol en la tabla de asociación. \\
\addlinespace[3pt]
\end{longtable}
\endgroup

\begingroup
\setbox0=\vbox{\hsize=\linewidth
{\normalfont\normalsize\bfseries Restricciones físicas observadas\par}
\begin{itemize}
\item Clave foránea: role\_\allowbreak{}id hacia app.roles (id); al eliminar: cascada SQL.
\end{itemize}
}
\Needspace{\dimexpr\ht0+\dp0+24pt\relax}
\endgroup
\begingroup
\setbox0=\vbox{\hsize=\linewidth
{\normalfont\normalsize\bfseries Restricciones físicas observadas\par}
\begin{itemize}
\item Clave foránea: role\_\allowbreak{}id hacia app.roles (id); al eliminar: cascada SQL.
\end{itemize}
}
\Needspace{\dimexpr\ht0+\dp0+24pt\relax}
\endgroup
\paragraph{Restricciones físicas observadas} \begin{itemize}
\item Clave foránea: role\_\allowbreak{}id hacia app.roles (id); al eliminar: cascada SQL.
\item Clave foránea: user\_\allowbreak{}id hacia app.users (id); al eliminar: cascada SQL.
\item Clave primaria: user\_\allowbreak{}id, role\_\allowbreak{}id.
\end{itemize}

\begingroup
\setbox0=\vbox{\hsize=\linewidth
{\normalfont\normalsize\bfseries Índices únicos\par}
\begin{itemize}
\item user\_\allowbreak{}id, role\_\allowbreak{}id; índice completo.
\end{itemize}
}
\Needspace{\dimexpr\ht0+\dp0+24pt\relax}
\endgroup
\begingroup
\setbox0=\vbox{\hsize=\linewidth
{\normalfont\normalsize\bfseries Índices únicos\par}
\begin{itemize}
\item user\_\allowbreak{}id, role\_\allowbreak{}id; índice completo.
\end{itemize}
}
\Needspace{\dimexpr\ht0+\dp0+24pt\relax}
\endgroup
\paragraph{Índices únicos} \begin{itemize}
\item user\_\allowbreak{}id, role\_\allowbreak{}id; índice completo.
\end{itemize}

\begingroup
\setbox0=\vbox{\hsize=\linewidth
{\normalfont\normalsize\sffamily\bfseries\normalsize\color{Teal}\hyphenpenalty=10000\exhyphenpenalty=10000 app.users\par}

\par\vspace{6pt}\noindent\textbf{Área}\hspace{0.5em} Seguridad

\par\vspace{6pt}\noindent\textbf{Qué guarda}\hspace{0.5em} La cuenta de una persona: identidad, estado de acceso y hash de contraseña.

\par\vspace{6pt}\noindent\textbf{Cómo se utiliza}\hspace{0.5em} Administración y acceso a la plataforma. No confundir al usuario de la aplicación con el usuario técnico de PostgreSQL o SQL Server.

\par
\begingroup
\fontsize{10}{12.5}\selectfont
\setlength{\tabcolsep}{5pt}
\renewcommand{\arraystretch}{1.15}
\begin{tabular}{@{}>{\RaggedRight\arraybackslash}p{0.23\linewidth}>{\RaggedRight\arraybackslash}p{0.23\linewidth}>{\RaggedRight\arraybackslash}p{0.45\linewidth}@{}}
\toprule
\textbf{Columna} & \textbf{Tipo y condición} & \textbf{Qué representa} \\
\midrule
\texttt{\seqsplit{id}} & integer\par\smallskip Obligatorio. Clave primaria. Con valor predeterminado. & Identificador interno de esta entidad; su PK distingue cada registro. \\
\bottomrule
\end{tabular}
\endgroup
}
\Needspace{\dimexpr\ht0+\dp0+48pt\relax}
\endgroup
\subsection{app.users}
\paragraph{Área} Seguridad

\paragraph{Qué guarda} La cuenta de una persona: identidad, estado de acceso y hash de contraseña.

\paragraph{Cómo se utiliza} Administración y acceso a la plataforma. No confundir al usuario de la aplicación con el usuario técnico de PostgreSQL o SQL Server.

\begingroup
\fontsize{10}{12.5}\selectfont
\setlength{\tabcolsep}{5pt}
\renewcommand{\arraystretch}{1.15}
\begin{longtable}{@{}>{\RaggedRight\arraybackslash}p{0.23\linewidth}>{\RaggedRight\arraybackslash}p{0.23\linewidth}>{\RaggedRight\arraybackslash}p{0.45\linewidth}@{}}
\toprule
\textbf{Columna} & \textbf{Tipo y condición} & \textbf{Qué representa} \\
\midrule
\endfirsthead
\toprule
\textbf{Columna} & \textbf{Tipo y condición} & \textbf{Qué representa} \\
\midrule
\endhead
\midrule
\multicolumn{3}{r}{\footnotesize Continúa en la página siguiente}\\
\endfoot
\bottomrule
\endlastfoot
\texttt{\seqsplit{id}} & integer\par\smallskip Obligatorio. Clave primaria. Con valor predeterminado. & Identificador interno de esta entidad; su PK distingue cada registro. \\
\addlinespace[3pt]
\texttt{\seqsplit{email}} & character varying(255)\par\smallskip Obligatorio. & Identificador de inicio de sesión; su índice único impide cuentas duplicadas con el mismo valor. \\
\addlinespace[3pt]
\texttt{\seqsplit{full\_name}} & character varying(160)\par\smallskip Obligatorio. & Nombre visible de la persona de la cuenta. \\
\addlinespace[3pt]
\texttt{\seqsplit{password\_hash}} & character varying(255)\par\smallskip Obligatorio. & Hash de contraseña, no contraseña en texto ni credencial recuperable. Este inventario sólo lee el nombre y tipo de la columna. \\
\addlinespace[3pt]
\texttt{\seqsplit{is\_active}} & boolean\par\smallskip Obligatorio. Con valor predeterminado. & Estado de habilitación del registro; su efecto depende de las reglas del módulo. \\
\addlinespace[3pt]
\texttt{\seqsplit{created\_at}} & timestamp with time zone\par\smallskip Admite nulos. Con valor predeterminado. & Fecha de creación del registro; su opcionalidad física consta en el inventario. \\
\addlinespace[3pt]
\texttt{\seqsplit{updated\_at}} & timestamp with time zone\par\smallskip Admite nulos. Con valor predeterminado. & Fecha de actualización del registro; su opcionalidad física consta en el inventario. \\
\addlinespace[3pt]
\texttt{\seqsplit{is\_system\_protected}} & boolean\par\smallskip Obligatorio. Con valor predeterminado. & Marca de un registro base protegido para conservar administración y navegación esenciales. \\
\addlinespace[3pt]
\end{longtable}
\endgroup

\begingroup
\setbox0=\vbox{\hsize=\linewidth
{\normalfont\normalsize\bfseries Restricciones físicas observadas\par}
\begin{itemize}
\item Clave primaria: id.
\end{itemize}
}
\Needspace{\dimexpr\ht0+\dp0+24pt\relax}
\endgroup
\begingroup
\setbox0=\vbox{\hsize=\linewidth
{\normalfont\normalsize\bfseries Restricciones físicas observadas\par}
\begin{itemize}
\item Clave primaria: id.
\end{itemize}
}
\Needspace{\dimexpr\ht0+\dp0+24pt\relax}
\endgroup
\paragraph{Restricciones físicas observadas} \begin{itemize}
\item Clave primaria: id.
\item Unicidad: email.
\end{itemize}

\begingroup
\setbox0=\vbox{\hsize=\linewidth
{\normalfont\normalsize\bfseries Índices únicos\par}
\begin{itemize}
\item email; índice completo.
\end{itemize}
}
\Needspace{\dimexpr\ht0+\dp0+24pt\relax}
\endgroup
\begingroup
\setbox0=\vbox{\hsize=\linewidth
{\normalfont\normalsize\bfseries Índices únicos\par}
\begin{itemize}
\item email; índice completo.
\end{itemize}
}
\Needspace{\dimexpr\ht0+\dp0+24pt\relax}
\endgroup
\paragraph{Índices únicos} \begin{itemize}
\item email; índice completo.
\item email; índice completo.
\item id; índice completo.
\end{itemize}

\begingroup
\setbox0=\vbox{\hsize=\linewidth
{\normalfont\normalsize\RaggedRight\sffamily\bfseries\large\color{Azul}\hyphenpenalty=10000\exhyphenpenalty=10000 Diccionario histórico de dos datamarts de demostración\par}

Las fichas siguientes conservan las columnas de los expedientes 21 y 22 del inventario del 18 de noviembre, para estudiar ejemplos concretos sin alterar la historia. No son el diccionario físico de los nuevos expedientes 23 y 24. Un diseño nuevo puede producir otro conjunto de atributos, como se explica en el apartado de destinos nuevos. Las 174 columnas de app descritas antes siguen vigentes al corte del 25 de noviembre. Una relación lógica del materializador no debe confundirse con una clave foránea física que la base no ha declarado.

\begingroup
\setbox0=\vbox{\hsize=\linewidth
{\normalfont\normalsize\sffamily\bfseries\normalsize\color{Teal}\hyphenpenalty=10000\exhyphenpenalty=10000 mart\_\allowbreak{}ventas\_\allowbreak{}e21.dim\_\allowbreak{}cliente\par}

\par
\par
\begingroup
\fontsize{10}{12.5}\selectfont
\setlength{\tabcolsep}{5pt}
\renewcommand{\arraystretch}{1.15}
\begin{tabular}{@{}>{\RaggedRight\arraybackslash}p{0.23\linewidth}>{\RaggedRight\arraybackslash}p{0.23\linewidth}>{\RaggedRight\arraybackslash}p{0.45\linewidth}@{}}
\toprule
\textbf{Columna} & \textbf{Tipo y condición} & \textbf{Qué representa} \\
\midrule
\texttt{\seqsplit{surrogate\_key}} & bigint\par\smallskip Obligatorio. & Clave sustituta interna de la dimensión; PK física. \\
\bottomrule
\end{tabular}
\endgroup
}
\Needspace{\dimexpr\ht0+\dp0+48pt\relax}
\endgroup
\subsection{mart\_\allowbreak{}ventas\_\allowbreak{}e21.dim\_\allowbreak{}cliente}
\begingroup
\fontsize{10}{12.5}\selectfont
\setlength{\tabcolsep}{5pt}
\renewcommand{\arraystretch}{1.15}
\begin{tabular}{@{}>{\RaggedRight\arraybackslash}p{0.23\linewidth}>{\RaggedRight\arraybackslash}p{0.23\linewidth}>{\RaggedRight\arraybackslash}p{0.45\linewidth}@{}}
\toprule
\textbf{Columna} & \textbf{Tipo y condición} & \textbf{Qué representa} \\
\midrule
\texttt{\seqsplit{surrogate\_key}} & bigint\par\smallskip Obligatorio. & Clave sustituta interna de la dimensión; PK física. \\
\bottomrule
\end{tabular}
\endgroup
}
\Needspace{\dimexpr\ht0+\dp0+48pt\relax}
\endgroup
\section{Diccionario histórico de dos datamarts de demostración}
Las fichas siguientes conservan las columnas de los expedientes 21 y 22 del inventario del 18 de noviembre, para estudiar ejemplos concretos sin alterar la historia. No son el diccionario físico de los nuevos expedientes 23 y 24. Un diseño nuevo puede producir otro conjunto de atributos, como se explica en el apartado de destinos nuevos. Las 174 columnas de app descritas antes siguen vigentes al corte del 25 de noviembre. Una relación lógica del materializador no debe confundirse con una clave foránea física que la base no ha declarado.

\begingroup
\setbox0=\vbox{\hsize=\linewidth
{\normalfont\normalsize\sffamily\bfseries\normalsize\color{Teal}\hyphenpenalty=10000\exhyphenpenalty=10000 mart\_\allowbreak{}ventas\_\allowbreak{}e21.dim\_\allowbreak{}cliente\par}

\par
\begingroup
\fontsize{10}{12.5}\selectfont
\setlength{\tabcolsep}{5pt}
\renewcommand{\arraystretch}{1.15}
\begin{tabular}{@{}>{\RaggedRight\arraybackslash}p{0.23\linewidth}>{\RaggedRight\arraybackslash}p{0.23\linewidth}>{\RaggedRight\arraybackslash}p{0.45\linewidth}@{}}
\toprule
\textbf{Columna} & \textbf{Tipo y condición} & \textbf{Qué representa} \\
\midrule
\texttt{\seqsplit{surrogate\_key}} & bigint\par\smallskip Obligatorio. & Clave sustituta interna de la dimensión; PK física. \\
\bottomrule
\end{tabular}
\endgroup
}
\Needspace{\dimexpr\ht0+\dp0+48pt\relax}
\endgroup
\subsection{mart\_\allowbreak{}ventas\_\allowbreak{}e21.dim\_\allowbreak{}cliente}
\begingroup
\fontsize{10}{12.5}\selectfont
\setlength{\tabcolsep}{5pt}
\renewcommand{\arraystretch}{1.15}
\begin{longtable}{@{}>{\RaggedRight\arraybackslash}p{0.23\linewidth}>{\RaggedRight\arraybackslash}p{0.23\linewidth}>{\RaggedRight\arraybackslash}p{0.45\linewidth}@{}}
\toprule
\textbf{Columna} & \textbf{Tipo y condición} & \textbf{Qué representa} \\
\midrule
\endfirsthead
\toprule
\textbf{Columna} & \textbf{Tipo y condición} & \textbf{Qué representa} \\
\midrule
\endhead
\midrule
\multicolumn{3}{r}{\footnotesize Continúa en la página siguiente}\\
\endfoot
\bottomrule
\endlastfoot
\texttt{\seqsplit{surrogate\_key}} & bigint\par\smallskip Obligatorio. & Clave sustituta interna de la dimensión; PK física. \\
\addlinespace[3pt]
\texttt{\seqsplit{customerid}} & bigint\par\smallskip Admite nulos. & Atributo de origen o etiqueta descriptiva seleccionada para esta dimensión. La propuesta y la evidencia definen su procedencia; no se deduce significado histórico únicamente del nombre. \\
\addlinespace[3pt]
\texttt{\seqsplit{accountnumber}} & text\par\smallskip Admite nulos. & Atributo de origen o etiqueta descriptiva seleccionada para esta dimensión. La propuesta y la evidencia definen su procedencia; no se deduce significado histórico únicamente del nombre. \\
\addlinespace[3pt]
\texttt{\seqsplit{nombre\_cliente}} & text\par\smallskip Admite nulos. & Atributo de origen o etiqueta descriptiva seleccionada para esta dimensión. La propuesta y la evidencia definen su procedencia; no se deduce significado histórico únicamente del nombre. \\
\addlinespace[3pt]
\texttt{\seqsplit{tipo\_cliente}} & text\par\smallskip Admite nulos. & Atributo de origen o etiqueta descriptiva seleccionada para esta dimensión. La propuesta y la evidencia definen su procedencia; no se deduce significado histórico únicamente del nombre. \\
\addlinespace[3pt]
\texttt{\seqsplit{etl\_execution\_id}} & bigint\par\smallskip Obligatorio. & Etiqueta de ejecución obligatoria; no tiene FK física a app.etl\_\allowbreak{}executions. \\
\addlinespace[3pt]
\end{longtable}
\endgroup

\begingroup
\setbox0=\vbox{\hsize=\linewidth
{\normalfont\normalsize\bfseries Restricciones físicas observadas\par}
\begin{itemize}
\item Clave primaria: surrogate\_\allowbreak{}key.
\end{itemize}
}
\Needspace{\dimexpr\ht0+\dp0+24pt\relax}
\endgroup
\begingroup
\setbox0=\vbox{\hsize=\linewidth
{\normalfont\normalsize\bfseries Restricciones físicas observadas\par}
\begin{itemize}
\item Clave primaria: surrogate\_\allowbreak{}key.
\end{itemize}
}
\Needspace{\dimexpr\ht0+\dp0+24pt\relax}
\endgroup
\paragraph{Restricciones físicas observadas} \begin{itemize}
\item Clave primaria: surrogate\_\allowbreak{}key.
\end{itemize}

\begingroup
\setbox0=\vbox{\hsize=\linewidth
{\normalfont\normalsize\sffamily\bfseries\normalsize\color{Teal}\hyphenpenalty=10000\exhyphenpenalty=10000 mart\_\allowbreak{}ventas\_\allowbreak{}e21.dim\_\allowbreak{}fecha\par}

\par
\begingroup
\fontsize{10}{12.5}\selectfont
\setlength{\tabcolsep}{5pt}
\renewcommand{\arraystretch}{1.15}
\begin{tabular}{@{}>{\RaggedRight\arraybackslash}p{0.23\linewidth}>{\RaggedRight\arraybackslash}p{0.23\linewidth}>{\RaggedRight\arraybackslash}p{0.45\linewidth}@{}}
\toprule
\textbf{Columna} & \textbf{Tipo y condición} & \textbf{Qué representa} \\
\midrule
\texttt{\seqsplit{surrogate\_key}} & bigint\par\smallskip Obligatorio. & Clave sustituta interna de la dimensión; PK física. \\
\bottomrule
\end{tabular}
\endgroup
}
\Needspace{\dimexpr\ht0+\dp0+48pt\relax}
\endgroup
\subsection{mart\_\allowbreak{}ventas\_\allowbreak{}e21.dim\_\allowbreak{}fecha}
\begingroup
\fontsize{10}{12.5}\selectfont
\setlength{\tabcolsep}{5pt}
\renewcommand{\arraystretch}{1.15}
\begin{longtable}{@{}>{\RaggedRight\arraybackslash}p{0.23\linewidth}>{\RaggedRight\arraybackslash}p{0.23\linewidth}>{\RaggedRight\arraybackslash}p{0.45\linewidth}@{}}
\toprule
\textbf{Columna} & \textbf{Tipo y condición} & \textbf{Qué representa} \\
\midrule
\endfirsthead
\toprule
\textbf{Columna} & \textbf{Tipo y condición} & \textbf{Qué representa} \\
\midrule
\endhead
\midrule
\multicolumn{3}{r}{\footnotesize Continúa en la página siguiente}\\
\endfoot
\bottomrule
\endlastfoot
\texttt{\seqsplit{surrogate\_key}} & bigint\par\smallskip Obligatorio. & Clave sustituta interna de la dimensión; PK física. \\
\addlinespace[3pt]
\texttt{\seqsplit{orderdate}} & timestamp without time zone\par\smallskip Admite nulos. & Atributo de origen o etiqueta descriptiva seleccionada para esta dimensión. La propuesta y la evidencia definen su procedencia; no se deduce significado histórico únicamente del nombre. \\
\addlinespace[3pt]
\texttt{\seqsplit{anio}} & bigint\par\smallskip Obligatorio. & Componente de calendario derivado de la fecha seleccionada. \\
\addlinespace[3pt]
\texttt{\seqsplit{trimestre}} & bigint\par\smallskip Obligatorio. & Componente de calendario derivado de la fecha seleccionada. \\
\addlinespace[3pt]
\texttt{\seqsplit{mes}} & bigint\par\smallskip Obligatorio. & Componente de calendario derivado de la fecha seleccionada. \\
\addlinespace[3pt]
\texttt{\seqsplit{dia}} & bigint\par\smallskip Obligatorio. & Componente de calendario derivado de la fecha seleccionada. \\
\addlinespace[3pt]
\texttt{\seqsplit{etl\_execution\_id}} & bigint\par\smallskip Obligatorio. & Etiqueta de ejecución obligatoria; no tiene FK física a app.etl\_\allowbreak{}executions. \\
\addlinespace[3pt]
\end{longtable}
\endgroup

\begingroup
\setbox0=\vbox{\hsize=\linewidth
{\normalfont\normalsize\bfseries Restricciones físicas observadas\par}
\begin{itemize}
\item Clave primaria: surrogate\_\allowbreak{}key.
\end{itemize}
}
\Needspace{\dimexpr\ht0+\dp0+24pt\relax}
\endgroup
\begingroup
\setbox0=\vbox{\hsize=\linewidth
{\normalfont\normalsize\bfseries Restricciones físicas observadas\par}
\begin{itemize}
\item Clave primaria: surrogate\_\allowbreak{}key.
\end{itemize}
}
\Needspace{\dimexpr\ht0+\dp0+24pt\relax}
\endgroup
\paragraph{Restricciones físicas observadas} \begin{itemize}
\item Clave primaria: surrogate\_\allowbreak{}key.
\end{itemize}

\begingroup
\setbox0=\vbox{\hsize=\linewidth
{\normalfont\normalsize\sffamily\bfseries\normalsize\color{Teal}\hyphenpenalty=10000\exhyphenpenalty=10000 mart\_\allowbreak{}ventas\_\allowbreak{}e21.dim\_\allowbreak{}producto\par}

\par
\begingroup
\fontsize{10}{12.5}\selectfont
\setlength{\tabcolsep}{5pt}
\renewcommand{\arraystretch}{1.15}
\begin{tabular}{@{}>{\RaggedRight\arraybackslash}p{0.23\linewidth}>{\RaggedRight\arraybackslash}p{0.23\linewidth}>{\RaggedRight\arraybackslash}p{0.45\linewidth}@{}}
\toprule
\textbf{Columna} & \textbf{Tipo y condición} & \textbf{Qué representa} \\
\midrule
\texttt{\seqsplit{surrogate\_key}} & bigint\par\smallskip Obligatorio. & Clave sustituta interna de la dimensión; PK física. \\
\bottomrule
\end{tabular}
\endgroup
}
\Needspace{\dimexpr\ht0+\dp0+48pt\relax}
\endgroup
\subsection{mart\_\allowbreak{}ventas\_\allowbreak{}e21.dim\_\allowbreak{}producto}
\begingroup
\fontsize{10}{12.5}\selectfont
\setlength{\tabcolsep}{5pt}
\renewcommand{\arraystretch}{1.15}
\begin{longtable}{@{}>{\RaggedRight\arraybackslash}p{0.23\linewidth}>{\RaggedRight\arraybackslash}p{0.23\linewidth}>{\RaggedRight\arraybackslash}p{0.45\linewidth}@{}}
\toprule
\textbf{Columna} & \textbf{Tipo y condición} & \textbf{Qué representa} \\
\midrule
\endfirsthead
\toprule
\textbf{Columna} & \textbf{Tipo y condición} & \textbf{Qué representa} \\
\midrule
\endhead
\midrule
\multicolumn{3}{r}{\footnotesize Continúa en la página siguiente}\\
\endfoot
\bottomrule
\endlastfoot
\texttt{\seqsplit{surrogate\_key}} & bigint\par\smallskip Obligatorio. & Clave sustituta interna de la dimensión; PK física. \\
\addlinespace[3pt]
\texttt{\seqsplit{productid}} & bigint\par\smallskip Admite nulos. & Atributo de origen o etiqueta descriptiva seleccionada para esta dimensión. La propuesta y la evidencia definen su procedencia; no se deduce significado histórico únicamente del nombre. \\
\addlinespace[3pt]
\texttt{\seqsplit{name}} & text\par\smallskip Admite nulos. & Atributo de origen o etiqueta descriptiva seleccionada para esta dimensión. La propuesta y la evidencia definen su procedencia; no se deduce significado histórico únicamente del nombre. \\
\addlinespace[3pt]
\texttt{\seqsplit{sizeunitmeasurecode}} & text\par\smallskip Admite nulos. & Atributo de origen o etiqueta descriptiva seleccionada para esta dimensión. La propuesta y la evidencia definen su procedencia; no se deduce significado histórico únicamente del nombre. \\
\addlinespace[3pt]
\texttt{\seqsplit{productnumber}} & text\par\smallskip Admite nulos. & Atributo de origen o etiqueta descriptiva seleccionada para esta dimensión. La propuesta y la evidencia definen su procedencia; no se deduce significado histórico únicamente del nombre. \\
\addlinespace[3pt]
\texttt{\seqsplit{weightunitmeasurecode}} & text\par\smallskip Admite nulos. & Atributo de origen o etiqueta descriptiva seleccionada para esta dimensión. La propuesta y la evidencia definen su procedencia; no se deduce significado histórico únicamente del nombre. \\
\addlinespace[3pt]
\texttt{\seqsplit{nombre\_producto}} & text\par\smallskip Admite nulos. & Atributo de origen o etiqueta descriptiva seleccionada para esta dimensión. La propuesta y la evidencia definen su procedencia; no se deduce significado histórico únicamente del nombre. \\
\addlinespace[3pt]
\texttt{\seqsplit{tipo\_producto}} & text\par\smallskip Admite nulos. & Atributo de origen o etiqueta descriptiva seleccionada para esta dimensión. La propuesta y la evidencia definen su procedencia; no se deduce significado histórico únicamente del nombre. \\
\addlinespace[3pt]
\texttt{\seqsplit{etl\_execution\_id}} & bigint\par\smallskip Obligatorio. & Etiqueta de ejecución obligatoria; no tiene FK física a app.etl\_\allowbreak{}executions. \\
\addlinespace[3pt]
\end{longtable}
\endgroup

\begingroup
\setbox0=\vbox{\hsize=\linewidth
{\normalfont\normalsize\bfseries Restricciones físicas observadas\par}
\begin{itemize}
\item Clave primaria: surrogate\_\allowbreak{}key.
\end{itemize}
}
\Needspace{\dimexpr\ht0+\dp0+24pt\relax}
\endgroup
\begingroup
\setbox0=\vbox{\hsize=\linewidth
{\normalfont\normalsize\bfseries Restricciones físicas observadas\par}
\begin{itemize}
\item Clave primaria: surrogate\_\allowbreak{}key.
\end{itemize}
}
\Needspace{\dimexpr\ht0+\dp0+24pt\relax}
\endgroup
\paragraph{Restricciones físicas observadas} \begin{itemize}
\item Clave primaria: surrogate\_\allowbreak{}key.
\end{itemize}

\begingroup
\setbox0=\vbox{\hsize=\linewidth
{\normalfont\normalsize\sffamily\bfseries\normalsize\color{Teal}\hyphenpenalty=10000\exhyphenpenalty=10000 mart\_\allowbreak{}ventas\_\allowbreak{}e21.dim\_\allowbreak{}territorio\par}

\par
\begingroup
\fontsize{10}{12.5}\selectfont
\setlength{\tabcolsep}{5pt}
\renewcommand{\arraystretch}{1.15}
\begin{tabular}{@{}>{\RaggedRight\arraybackslash}p{0.23\linewidth}>{\RaggedRight\arraybackslash}p{0.23\linewidth}>{\RaggedRight\arraybackslash}p{0.45\linewidth}@{}}
\toprule
\textbf{Columna} & \textbf{Tipo y condición} & \textbf{Qué representa} \\
\midrule
\texttt{\seqsplit{surrogate\_key}} & bigint\par\smallskip Obligatorio. & Clave sustituta interna de la dimensión; PK física. \\
\bottomrule
\end{tabular}
\endgroup
}
\Needspace{\dimexpr\ht0+\dp0+48pt\relax}
\endgroup
\subsection{mart\_\allowbreak{}ventas\_\allowbreak{}e21.dim\_\allowbreak{}territorio}
\begingroup
\fontsize{10}{12.5}\selectfont
\setlength{\tabcolsep}{5pt}
\renewcommand{\arraystretch}{1.15}
\begin{longtable}{@{}>{\RaggedRight\arraybackslash}p{0.23\linewidth}>{\RaggedRight\arraybackslash}p{0.23\linewidth}>{\RaggedRight\arraybackslash}p{0.45\linewidth}@{}}
\toprule
\textbf{Columna} & \textbf{Tipo y condición} & \textbf{Qué representa} \\
\midrule
\endfirsthead
\toprule
\textbf{Columna} & \textbf{Tipo y condición} & \textbf{Qué representa} \\
\midrule
\endhead
\midrule
\multicolumn{3}{r}{\footnotesize Continúa en la página siguiente}\\
\endfoot
\bottomrule
\endlastfoot
\texttt{\seqsplit{surrogate\_key}} & bigint\par\smallskip Obligatorio. & Clave sustituta interna de la dimensión; PK física. \\
\addlinespace[3pt]
\texttt{\seqsplit{territoryid}} & bigint\par\smallskip Admite nulos. & Atributo de origen o etiqueta descriptiva seleccionada para esta dimensión. La propuesta y la evidencia definen su procedencia; no se deduce significado histórico únicamente del nombre. \\
\addlinespace[3pt]
\texttt{\seqsplit{countryregioncode}} & text\par\smallskip Admite nulos. & Atributo de origen o etiqueta descriptiva seleccionada para esta dimensión. La propuesta y la evidencia definen su procedencia; no se deduce significado histórico únicamente del nombre. \\
\addlinespace[3pt]
\texttt{\seqsplit{group}} & text\par\smallskip Admite nulos. & Atributo de origen o etiqueta descriptiva seleccionada para esta dimensión. La propuesta y la evidencia definen su procedencia; no se deduce significado histórico únicamente del nombre. \\
\addlinespace[3pt]
\texttt{\seqsplit{name}} & text\par\smallskip Admite nulos. & Atributo de origen o etiqueta descriptiva seleccionada para esta dimensión. La propuesta y la evidencia definen su procedencia; no se deduce significado histórico únicamente del nombre. \\
\addlinespace[3pt]
\texttt{\seqsplit{nombre\_territory\_dimension}} & text\par\smallskip Admite nulos. & Atributo de origen o etiqueta descriptiva seleccionada para esta dimensión. La propuesta y la evidencia definen su procedencia; no se deduce significado histórico únicamente del nombre. \\
\addlinespace[3pt]
\texttt{\seqsplit{tipo\_territory\_dimension}} & text\par\smallskip Admite nulos. & Atributo de origen o etiqueta descriptiva seleccionada para esta dimensión. La propuesta y la evidencia definen su procedencia; no se deduce significado histórico únicamente del nombre. \\
\addlinespace[3pt]
\texttt{\seqsplit{etl\_execution\_id}} & bigint\par\smallskip Obligatorio. & Etiqueta de ejecución obligatoria; no tiene FK física a app.etl\_\allowbreak{}executions. \\
\addlinespace[3pt]
\texttt{\seqsplit{group\_es}} & text\par\smallskip Admite nulos. & Atributo de origen o etiqueta descriptiva seleccionada para esta dimensión. La propuesta y la evidencia definen su procedencia; no se deduce significado histórico únicamente del nombre. \\
\addlinespace[3pt]
\end{longtable}
\endgroup

\begingroup
\setbox0=\vbox{\hsize=\linewidth
{\normalfont\normalsize\bfseries Restricciones físicas observadas\par}
\begin{itemize}
\item Clave primaria: surrogate\_\allowbreak{}key.
\end{itemize}
}
\Needspace{\dimexpr\ht0+\dp0+24pt\relax}
\endgroup
\begingroup
\setbox0=\vbox{\hsize=\linewidth
{\normalfont\normalsize\bfseries Restricciones físicas observadas\par}
\begin{itemize}
\item Clave primaria: surrogate\_\allowbreak{}key.
\end{itemize}
}
\Needspace{\dimexpr\ht0+\dp0+24pt\relax}
\endgroup
\paragraph{Restricciones físicas observadas} \begin{itemize}
\item Clave primaria: surrogate\_\allowbreak{}key.
\end{itemize}

\begingroup
\setbox0=\vbox{\hsize=\linewidth
{\normalfont\normalsize\sffamily\bfseries\normalsize\color{Teal}\hyphenpenalty=10000\exhyphenpenalty=10000 mart\_\allowbreak{}ventas\_\allowbreak{}e21.fact\_\allowbreak{}ventas\par}

\par
\begingroup
\fontsize{10}{12.5}\selectfont
\setlength{\tabcolsep}{5pt}
\renewcommand{\arraystretch}{1.15}
\begin{tabular}{@{}>{\RaggedRight\arraybackslash}p{0.23\linewidth}>{\RaggedRight\arraybackslash}p{0.23\linewidth}>{\RaggedRight\arraybackslash}p{0.45\linewidth}@{}}
\toprule
\textbf{Columna} & \textbf{Tipo y condición} & \textbf{Qué representa} \\
\midrule
\texttt{\seqsplit{salesorderid}} & bigint\par\smallskip Admite nulos. & Clave de trazabilidad o medida de origen seleccionada por el contrato de esta ejecución. Su significado y agregación se revisan en la propuesta; el nombre físico no acredita su semántica. \\
\bottomrule
\end{tabular}
\endgroup
}
\Needspace{\dimexpr\ht0+\dp0+48pt\relax}
\endgroup
\subsection{mart\_\allowbreak{}ventas\_\allowbreak{}e21.fact\_\allowbreak{}ventas}
\begingroup
\fontsize{10}{12.5}\selectfont
\setlength{\tabcolsep}{5pt}
\renewcommand{\arraystretch}{1.15}
\begin{longtable}{@{}>{\RaggedRight\arraybackslash}p{0.23\linewidth}>{\RaggedRight\arraybackslash}p{0.23\linewidth}>{\RaggedRight\arraybackslash}p{0.45\linewidth}@{}}
\toprule
\textbf{Columna} & \textbf{Tipo y condición} & \textbf{Qué representa} \\
\midrule
\endfirsthead
\toprule
\textbf{Columna} & \textbf{Tipo y condición} & \textbf{Qué representa} \\
\midrule
\endhead
\midrule
\multicolumn{3}{r}{\footnotesize Continúa en la página siguiente}\\
\endfoot
\bottomrule
\endlastfoot
\texttt{\seqsplit{salesorderid}} & bigint\par\smallskip Admite nulos. & Clave de trazabilidad o medida de origen seleccionada por el contrato de esta ejecución. Su significado y agregación se revisan en la propuesta; el nombre físico no acredita su semántica. \\
\addlinespace[3pt]
\texttt{\seqsplit{salesorderdetailid}} & bigint\par\smallskip Admite nulos. & Clave de trazabilidad o medida de origen seleccionada por el contrato de esta ejecución. Su significado y agregación se revisan en la propuesta; el nombre físico no acredita su semántica. \\
\addlinespace[3pt]
\texttt{\seqsplit{importe\_total\_de\_l\_neas}} & numeric\par\smallskip Admite nulos. & Clave de trazabilidad o medida de origen seleccionada por el contrato de esta ejecución. Su significado y agregación se revisan en la propuesta; el nombre físico no acredita su semántica. \\
\addlinespace[3pt]
\texttt{\seqsplit{cantidad\_de\_unidades}} & bigint\par\smallskip Admite nulos. & Clave de trazabilidad o medida de origen seleccionada por el contrato de esta ejecución. Su significado y agregación se revisan en la propuesta; el nombre físico no acredita su semántica. \\
\addlinespace[3pt]
\texttt{\seqsplit{n\_mero\_de\_pedidos\_distintos}} & bigint\par\smallskip Admite nulos. & Campo materializado para la medida documental: su receta puede calcular COUNT DISTINCT. No debe sumarse como si cada valor ya fuera un conteo. \\
\addlinespace[3pt]
\texttt{\seqsplit{dim\_fecha\_sk}} & bigint\par\smallskip Admite nulos. & Referencia lógica a surrogate\_\allowbreak{}key de una dimensión; puede ser nula y no tiene FK física. \\
\addlinespace[3pt]
\texttt{\seqsplit{dim\_producto\_sk}} & bigint\par\smallskip Admite nulos. & Referencia lógica a surrogate\_\allowbreak{}key de una dimensión; puede ser nula y no tiene FK física. \\
\addlinespace[3pt]
\texttt{\seqsplit{dim\_cliente\_sk}} & bigint\par\smallskip Admite nulos. & Referencia lógica a surrogate\_\allowbreak{}key de una dimensión; puede ser nula y no tiene FK física. \\
\addlinespace[3pt]
\texttt{\seqsplit{dim\_territorio\_sk}} & bigint\par\smallskip Admite nulos. & Referencia lógica a surrogate\_\allowbreak{}key de una dimensión; puede ser nula y no tiene FK física. \\
\addlinespace[3pt]
\texttt{\seqsplit{etl\_execution\_id}} & bigint\par\smallskip Obligatorio. & Etiqueta de ejecución obligatoria; no tiene FK física a app.etl\_\allowbreak{}executions. \\
\addlinespace[3pt]
\end{longtable}
\endgroup

\begingroup
\setbox0=\vbox{\hsize=\linewidth
{\normalfont\normalsize\sffamily\bfseries\normalsize\color{Teal}\hyphenpenalty=10000\exhyphenpenalty=10000 mart\_\allowbreak{}ventas\_\allowbreak{}e22.dim\_\allowbreak{}cliente\par}

\par
\begingroup
\fontsize{10}{12.5}\selectfont
\setlength{\tabcolsep}{5pt}
\renewcommand{\arraystretch}{1.15}
\begin{tabular}{@{}>{\RaggedRight\arraybackslash}p{0.23\linewidth}>{\RaggedRight\arraybackslash}p{0.23\linewidth}>{\RaggedRight\arraybackslash}p{0.45\linewidth}@{}}
\toprule
\textbf{Columna} & \textbf{Tipo y condición} & \textbf{Qué representa} \\
\midrule
\texttt{\seqsplit{surrogate\_key}} & bigint\par\smallskip Obligatorio. & Clave sustituta interna de la dimensión; PK física. \\
\bottomrule
\end{tabular}
\endgroup
}
\Needspace{\dimexpr\ht0+\dp0+48pt\relax}
\endgroup
\subsection{mart\_\allowbreak{}ventas\_\allowbreak{}e22.dim\_\allowbreak{}cliente}
\begingroup
\fontsize{10}{12.5}\selectfont
\setlength{\tabcolsep}{5pt}
\renewcommand{\arraystretch}{1.15}
\begin{longtable}{@{}>{\RaggedRight\arraybackslash}p{0.23\linewidth}>{\RaggedRight\arraybackslash}p{0.23\linewidth}>{\RaggedRight\arraybackslash}p{0.45\linewidth}@{}}
\toprule
\textbf{Columna} & \textbf{Tipo y condición} & \textbf{Qué representa} \\
\midrule
\endfirsthead
\toprule
\textbf{Columna} & \textbf{Tipo y condición} & \textbf{Qué representa} \\
\midrule
\endhead
\midrule
\multicolumn{3}{r}{\footnotesize Continúa en la página siguiente}\\
\endfoot
\bottomrule
\endlastfoot
\texttt{\seqsplit{surrogate\_key}} & bigint\par\smallskip Obligatorio. & Clave sustituta interna de la dimensión; PK física. \\
\addlinespace[3pt]
\texttt{\seqsplit{customerid}} & bigint\par\smallskip Admite nulos. & Atributo de origen o etiqueta descriptiva seleccionada para esta dimensión. La propuesta y la evidencia definen su procedencia; no se deduce significado histórico únicamente del nombre. \\
\addlinespace[3pt]
\texttt{\seqsplit{customername}} & text\par\smallskip Admite nulos. & Atributo de origen o etiqueta descriptiva seleccionada para esta dimensión. La propuesta y la evidencia definen su procedencia; no se deduce significado histórico únicamente del nombre. \\
\addlinespace[3pt]
\texttt{\seqsplit{nombre\_cliente}} & text\par\smallskip Admite nulos. & Atributo de origen o etiqueta descriptiva seleccionada para esta dimensión. La propuesta y la evidencia definen su procedencia; no se deduce significado histórico únicamente del nombre. \\
\addlinespace[3pt]
\texttt{\seqsplit{tipo\_cliente}} & text\par\smallskip Admite nulos. & Atributo de origen o etiqueta descriptiva seleccionada para esta dimensión. La propuesta y la evidencia definen su procedencia; no se deduce significado histórico únicamente del nombre. \\
\addlinespace[3pt]
\texttt{\seqsplit{etl\_execution\_id}} & bigint\par\smallskip Obligatorio. & Etiqueta de ejecución obligatoria; no tiene FK física a app.etl\_\allowbreak{}executions. \\
\addlinespace[3pt]
\end{longtable}
\endgroup

\begingroup
\setbox0=\vbox{\hsize=\linewidth
{\normalfont\normalsize\bfseries Restricciones físicas observadas\par}
\begin{itemize}
\item Clave primaria: surrogate\_\allowbreak{}key.
\end{itemize}
}
\Needspace{\dimexpr\ht0+\dp0+24pt\relax}
\endgroup
\begingroup
\setbox0=\vbox{\hsize=\linewidth
{\normalfont\normalsize\bfseries Restricciones físicas observadas\par}
\begin{itemize}
\item Clave primaria: surrogate\_\allowbreak{}key.
\end{itemize}
}
\Needspace{\dimexpr\ht0+\dp0+24pt\relax}
\endgroup
\paragraph{Restricciones físicas observadas} \begin{itemize}
\item Clave primaria: surrogate\_\allowbreak{}key.
\end{itemize}

\begingroup
\setbox0=\vbox{\hsize=\linewidth
{\normalfont\normalsize\sffamily\bfseries\normalsize\color{Teal}\hyphenpenalty=10000\exhyphenpenalty=10000 mart\_\allowbreak{}ventas\_\allowbreak{}e22.dim\_\allowbreak{}fecha\par}

\par
\begingroup
\fontsize{10}{12.5}\selectfont
\setlength{\tabcolsep}{5pt}
\renewcommand{\arraystretch}{1.15}
\begin{tabular}{@{}>{\RaggedRight\arraybackslash}p{0.23\linewidth}>{\RaggedRight\arraybackslash}p{0.23\linewidth}>{\RaggedRight\arraybackslash}p{0.45\linewidth}@{}}
\toprule
\textbf{Columna} & \textbf{Tipo y condición} & \textbf{Qué representa} \\
\midrule
\texttt{\seqsplit{surrogate\_key}} & bigint\par\smallskip Obligatorio. & Clave sustituta interna de la dimensión; PK física. \\
\bottomrule
\end{tabular}
\endgroup
}
\Needspace{\dimexpr\ht0+\dp0+48pt\relax}
\endgroup
\subsection{mart\_\allowbreak{}ventas\_\allowbreak{}e22.dim\_\allowbreak{}fecha}
\begingroup
\fontsize{10}{12.5}\selectfont
\setlength{\tabcolsep}{5pt}
\renewcommand{\arraystretch}{1.15}
\begin{longtable}{@{}>{\RaggedRight\arraybackslash}p{0.23\linewidth}>{\RaggedRight\arraybackslash}p{0.23\linewidth}>{\RaggedRight\arraybackslash}p{0.45\linewidth}@{}}
\toprule
\textbf{Columna} & \textbf{Tipo y condición} & \textbf{Qué representa} \\
\midrule
\endfirsthead
\toprule
\textbf{Columna} & \textbf{Tipo y condición} & \textbf{Qué representa} \\
\midrule
\endhead
\midrule
\multicolumn{3}{r}{\footnotesize Continúa en la página siguiente}\\
\endfoot
\bottomrule
\endlastfoot
\texttt{\seqsplit{surrogate\_key}} & bigint\par\smallskip Obligatorio. & Clave sustituta interna de la dimensión; PK física. \\
\addlinespace[3pt]
\texttt{\seqsplit{invoicedate}} & date\par\smallskip Admite nulos. & Atributo de origen o etiqueta descriptiva seleccionada para esta dimensión. La propuesta y la evidencia definen su procedencia; no se deduce significado histórico únicamente del nombre. \\
\addlinespace[3pt]
\texttt{\seqsplit{anio}} & bigint\par\smallskip Obligatorio. & Componente de calendario derivado de la fecha seleccionada. \\
\addlinespace[3pt]
\texttt{\seqsplit{trimestre}} & bigint\par\smallskip Obligatorio. & Componente de calendario derivado de la fecha seleccionada. \\
\addlinespace[3pt]
\texttt{\seqsplit{mes}} & bigint\par\smallskip Obligatorio. & Componente de calendario derivado de la fecha seleccionada. \\
\addlinespace[3pt]
\texttt{\seqsplit{dia}} & bigint\par\smallskip Obligatorio. & Componente de calendario derivado de la fecha seleccionada. \\
\addlinespace[3pt]
\texttt{\seqsplit{etl\_execution\_id}} & bigint\par\smallskip Obligatorio. & Etiqueta de ejecución obligatoria; no tiene FK física a app.etl\_\allowbreak{}executions. \\
\addlinespace[3pt]
\end{longtable}
\endgroup

\begingroup
\setbox0=\vbox{\hsize=\linewidth
{\normalfont\normalsize\bfseries Restricciones físicas observadas\par}
\begin{itemize}
\item Clave primaria: surrogate\_\allowbreak{}key.
\end{itemize}
}
\Needspace{\dimexpr\ht0+\dp0+24pt\relax}
\endgroup
\begingroup
\setbox0=\vbox{\hsize=\linewidth
{\normalfont\normalsize\bfseries Restricciones físicas observadas\par}
\begin{itemize}
\item Clave primaria: surrogate\_\allowbreak{}key.
\end{itemize}
}
\Needspace{\dimexpr\ht0+\dp0+24pt\relax}
\endgroup
\paragraph{Restricciones físicas observadas} \begin{itemize}
\item Clave primaria: surrogate\_\allowbreak{}key.
\end{itemize}

\begingroup
\setbox0=\vbox{\hsize=\linewidth
{\normalfont\normalsize\sffamily\bfseries\normalsize\color{Teal}\hyphenpenalty=10000\exhyphenpenalty=10000 mart\_\allowbreak{}ventas\_\allowbreak{}e22.dim\_\allowbreak{}producto\par}

\par
\begingroup
\fontsize{10}{12.5}\selectfont
\setlength{\tabcolsep}{5pt}
\renewcommand{\arraystretch}{1.15}
\begin{tabular}{@{}>{\RaggedRight\arraybackslash}p{0.23\linewidth}>{\RaggedRight\arraybackslash}p{0.23\linewidth}>{\RaggedRight\arraybackslash}p{0.45\linewidth}@{}}
\toprule
\textbf{Columna} & \textbf{Tipo y condición} & \textbf{Qué representa} \\
\midrule
\texttt{\seqsplit{surrogate\_key}} & bigint\par\smallskip Obligatorio. & Clave sustituta interna de la dimensión; PK física. \\
\bottomrule
\end{tabular}
\endgroup
}
\Needspace{\dimexpr\ht0+\dp0+48pt\relax}
\endgroup
\subsection{mart\_\allowbreak{}ventas\_\allowbreak{}e22.dim\_\allowbreak{}producto}
\begingroup
\fontsize{10}{12.5}\selectfont
\setlength{\tabcolsep}{5pt}
\renewcommand{\arraystretch}{1.15}
\begin{longtable}{@{}>{\RaggedRight\arraybackslash}p{0.23\linewidth}>{\RaggedRight\arraybackslash}p{0.23\linewidth}>{\RaggedRight\arraybackslash}p{0.45\linewidth}@{}}
\toprule
\textbf{Columna} & \textbf{Tipo y condición} & \textbf{Qué representa} \\
\midrule
\endfirsthead
\toprule
\textbf{Columna} & \textbf{Tipo y condición} & \textbf{Qué representa} \\
\midrule
\endhead
\midrule
\multicolumn{3}{r}{\footnotesize Continúa en la página siguiente}\\
\endfoot
\bottomrule
\endlastfoot
\texttt{\seqsplit{surrogate\_key}} & bigint\par\smallskip Obligatorio. & Clave sustituta interna de la dimensión; PK física. \\
\addlinespace[3pt]
\texttt{\seqsplit{stockitemid}} & bigint\par\smallskip Admite nulos. & Atributo de origen o etiqueta descriptiva seleccionada para esta dimensión. La propuesta y la evidencia definen su procedencia; no se deduce significado histórico únicamente del nombre. \\
\addlinespace[3pt]
\texttt{\seqsplit{stockitemname}} & text\par\smallskip Admite nulos. & Atributo de origen o etiqueta descriptiva seleccionada para esta dimensión. La propuesta y la evidencia definen su procedencia; no se deduce significado histórico únicamente del nombre. \\
\addlinespace[3pt]
\texttt{\seqsplit{nombre\_producto}} & text\par\smallskip Admite nulos. & Atributo de origen o etiqueta descriptiva seleccionada para esta dimensión. La propuesta y la evidencia definen su procedencia; no se deduce significado histórico únicamente del nombre. \\
\addlinespace[3pt]
\texttt{\seqsplit{tipo\_producto}} & text\par\smallskip Admite nulos. & Atributo de origen o etiqueta descriptiva seleccionada para esta dimensión. La propuesta y la evidencia definen su procedencia; no se deduce significado histórico únicamente del nombre. \\
\addlinespace[3pt]
\texttt{\seqsplit{etl\_execution\_id}} & bigint\par\smallskip Obligatorio. & Etiqueta de ejecución obligatoria; no tiene FK física a app.etl\_\allowbreak{}executions. \\
\addlinespace[3pt]
\end{longtable}
\endgroup

\begingroup
\setbox0=\vbox{\hsize=\linewidth
{\normalfont\normalsize\bfseries Restricciones físicas observadas\par}
\begin{itemize}
\item Clave primaria: surrogate\_\allowbreak{}key.
\end{itemize}
}
\Needspace{\dimexpr\ht0+\dp0+24pt\relax}
\endgroup
\begingroup
\setbox0=\vbox{\hsize=\linewidth
{\normalfont\normalsize\bfseries Restricciones físicas observadas\par}
\begin{itemize}
\item Clave primaria: surrogate\_\allowbreak{}key.
\end{itemize}
}
\Needspace{\dimexpr\ht0+\dp0+24pt\relax}
\endgroup
\paragraph{Restricciones físicas observadas} \begin{itemize}
\item Clave primaria: surrogate\_\allowbreak{}key.
\end{itemize}

\begingroup
\setbox0=\vbox{\hsize=\linewidth
{\normalfont\normalsize\sffamily\bfseries\normalsize\color{Teal}\hyphenpenalty=10000\exhyphenpenalty=10000 mart\_\allowbreak{}ventas\_\allowbreak{}e22.dim\_\allowbreak{}territorio\par}

\par
\begingroup
\fontsize{10}{12.5}\selectfont
\setlength{\tabcolsep}{5pt}
\renewcommand{\arraystretch}{1.15}
\begin{tabular}{@{}>{\RaggedRight\arraybackslash}p{0.23\linewidth}>{\RaggedRight\arraybackslash}p{0.23\linewidth}>{\RaggedRight\arraybackslash}p{0.45\linewidth}@{}}
\toprule
\textbf{Columna} & \textbf{Tipo y condición} & \textbf{Qué representa} \\
\midrule
\texttt{\seqsplit{surrogate\_key}} & bigint\par\smallskip Obligatorio. & Clave sustituta interna de la dimensión; PK física. \\
\bottomrule
\end{tabular}
\endgroup
}
\Needspace{\dimexpr\ht0+\dp0+48pt\relax}
\endgroup
\subsection{mart\_\allowbreak{}ventas\_\allowbreak{}e22.dim\_\allowbreak{}territorio}
\begingroup
\fontsize{10}{12.5}\selectfont
\setlength{\tabcolsep}{5pt}
\renewcommand{\arraystretch}{1.15}
\begin{longtable}{@{}>{\RaggedRight\arraybackslash}p{0.23\linewidth}>{\RaggedRight\arraybackslash}p{0.23\linewidth}>{\RaggedRight\arraybackslash}p{0.45\linewidth}@{}}
\toprule
\textbf{Columna} & \textbf{Tipo y condición} & \textbf{Qué representa} \\
\midrule
\endfirsthead
\toprule
\textbf{Columna} & \textbf{Tipo y condición} & \textbf{Qué representa} \\
\midrule
\endhead
\midrule
\multicolumn{3}{r}{\footnotesize Continúa en la página siguiente}\\
\endfoot
\bottomrule
\endlastfoot
\texttt{\seqsplit{surrogate\_key}} & bigint\par\smallskip Obligatorio. & Clave sustituta interna de la dimensión; PK física. \\
\addlinespace[3pt]
\texttt{\seqsplit{stateprovinceid}} & bigint\par\smallskip Admite nulos. & Atributo de origen o etiqueta descriptiva seleccionada para esta dimensión. La propuesta y la evidencia definen su procedencia; no se deduce significado histórico únicamente del nombre. \\
\addlinespace[3pt]
\texttt{\seqsplit{stateprovincename}} & text\par\smallskip Admite nulos. & Atributo de origen o etiqueta descriptiva seleccionada para esta dimensión. La propuesta y la evidencia definen su procedencia; no se deduce significado histórico únicamente del nombre. \\
\addlinespace[3pt]
\texttt{\seqsplit{stateprovincecode}} & text\par\smallskip Admite nulos. & Atributo de origen o etiqueta descriptiva seleccionada para esta dimensión. La propuesta y la evidencia definen su procedencia; no se deduce significado histórico únicamente del nombre. \\
\addlinespace[3pt]
\texttt{\seqsplit{salesterritory}} & text\par\smallskip Admite nulos. & Atributo de origen o etiqueta descriptiva seleccionada para esta dimensión. La propuesta y la evidencia definen su procedencia; no se deduce significado histórico únicamente del nombre. \\
\addlinespace[3pt]
\texttt{\seqsplit{nombre\_territorio}} & text\par\smallskip Admite nulos. & Atributo de origen o etiqueta descriptiva seleccionada para esta dimensión. La propuesta y la evidencia definen su procedencia; no se deduce significado histórico únicamente del nombre. \\
\addlinespace[3pt]
\texttt{\seqsplit{tipo\_territorio}} & text\par\smallskip Admite nulos. & Atributo de origen o etiqueta descriptiva seleccionada para esta dimensión. La propuesta y la evidencia definen su procedencia; no se deduce significado histórico únicamente del nombre. \\
\addlinespace[3pt]
\texttt{\seqsplit{etl\_execution\_id}} & bigint\par\smallskip Obligatorio. & Etiqueta de ejecución obligatoria; no tiene FK física a app.etl\_\allowbreak{}executions. \\
\addlinespace[3pt]
\end{longtable}
\endgroup

\begingroup
\setbox0=\vbox{\hsize=\linewidth
{\normalfont\normalsize\bfseries Restricciones físicas observadas\par}
\begin{itemize}
\item Clave primaria: surrogate\_\allowbreak{}key.
\end{itemize}
}
\Needspace{\dimexpr\ht0+\dp0+24pt\relax}
\endgroup
\begingroup
\setbox0=\vbox{\hsize=\linewidth
{\normalfont\normalsize\bfseries Restricciones físicas observadas\par}
\begin{itemize}
\item Clave primaria: surrogate\_\allowbreak{}key.
\end{itemize}
}
\Needspace{\dimexpr\ht0+\dp0+24pt\relax}
\endgroup
\paragraph{Restricciones físicas observadas} \begin{itemize}
\item Clave primaria: surrogate\_\allowbreak{}key.
\end{itemize}

\begingroup
\setbox0=\vbox{\hsize=\linewidth
{\normalfont\normalsize\sffamily\bfseries\normalsize\color{Teal}\hyphenpenalty=10000\exhyphenpenalty=10000 mart\_\allowbreak{}ventas\_\allowbreak{}e22.fact\_\allowbreak{}ventas\par}

\par
\begingroup
\fontsize{10}{12.5}\selectfont
\setlength{\tabcolsep}{5pt}
\renewcommand{\arraystretch}{1.15}
\begin{tabular}{@{}>{\RaggedRight\arraybackslash}p{0.23\linewidth}>{\RaggedRight\arraybackslash}p{0.23\linewidth}>{\RaggedRight\arraybackslash}p{0.45\linewidth}@{}}
\toprule
\textbf{Columna} & \textbf{Tipo y condición} & \textbf{Qué representa} \\
\midrule
\texttt{\seqsplit{invoicelineid}} & bigint\par\smallskip Admite nulos. & Clave de trazabilidad o medida de origen seleccionada por el contrato de esta ejecución. Su significado y agregación se revisan en la propuesta; el nombre físico no acredita su semántica. \\
\bottomrule
\end{tabular}
\endgroup
}
\Needspace{\dimexpr\ht0+\dp0+48pt\relax}
\endgroup
\subsection{mart\_\allowbreak{}ventas\_\allowbreak{}e22.fact\_\allowbreak{}ventas}
\begingroup
\fontsize{10}{12.5}\selectfont
\setlength{\tabcolsep}{5pt}
\renewcommand{\arraystretch}{1.15}
\begin{longtable}{@{}>{\RaggedRight\arraybackslash}p{0.23\linewidth}>{\RaggedRight\arraybackslash}p{0.23\linewidth}>{\RaggedRight\arraybackslash}p{0.45\linewidth}@{}}
\toprule
\textbf{Columna} & \textbf{Tipo y condición} & \textbf{Qué representa} \\
\midrule
\endfirsthead
\toprule
\textbf{Columna} & \textbf{Tipo y condición} & \textbf{Qué representa} \\
\midrule
\endhead
\midrule
\multicolumn{3}{r}{\footnotesize Continúa en la página siguiente}\\
\endfoot
\bottomrule
\endlastfoot
\texttt{\seqsplit{invoicelineid}} & bigint\par\smallskip Admite nulos. & Clave de trazabilidad o medida de origen seleccionada por el contrato de esta ejecución. Su significado y agregación se revisan en la propuesta; el nombre físico no acredita su semántica. \\
\addlinespace[3pt]
\texttt{\seqsplit{importe\_total\_de\_l\_nea}} & numeric\par\smallskip Admite nulos. & Clave de trazabilidad o medida de origen seleccionada por el contrato de esta ejecución. Su significado y agregación se revisan en la propuesta; el nombre físico no acredita su semántica. \\
\addlinespace[3pt]
\texttt{\seqsplit{unidades\_totales}} & bigint\par\smallskip Admite nulos. & Clave de trazabilidad o medida de origen seleccionada por el contrato de esta ejecución. Su significado y agregación se revisan en la propuesta; el nombre físico no acredita su semántica. \\
\addlinespace[3pt]
\texttt{\seqsplit{facturas\_distintas}} & bigint\par\smallskip Admite nulos. & Campo materializado para la medida documental: su receta puede calcular COUNT DISTINCT. No debe sumarse como si cada valor ya fuera un conteo. \\
\addlinespace[3pt]
\texttt{\seqsplit{dim\_fecha\_sk}} & bigint\par\smallskip Admite nulos. & Referencia lógica a surrogate\_\allowbreak{}key de una dimensión; puede ser nula y no tiene FK física. \\
\addlinespace[3pt]
\texttt{\seqsplit{dim\_producto\_sk}} & bigint\par\smallskip Admite nulos. & Referencia lógica a surrogate\_\allowbreak{}key de una dimensión; puede ser nula y no tiene FK física. \\
\addlinespace[3pt]
\texttt{\seqsplit{dim\_cliente\_sk}} & bigint\par\smallskip Admite nulos. & Referencia lógica a surrogate\_\allowbreak{}key de una dimensión; puede ser nula y no tiene FK física. \\
\addlinespace[3pt]
\texttt{\seqsplit{dim\_territorio\_sk}} & bigint\par\smallskip Admite nulos. & Referencia lógica a surrogate\_\allowbreak{}key de una dimensión; puede ser nula y no tiene FK física. \\
\addlinespace[3pt]
\texttt{\seqsplit{etl\_execution\_id}} & bigint\par\smallskip Obligatorio. & Etiqueta de ejecución obligatoria; no tiene FK física a app.etl\_\allowbreak{}executions. \\
\addlinespace[3pt]
\end{longtable}
\endgroup

\section{Alcance del inventario y límites}
La inspección estructural realizada para actualizar este capítulo no modificó código, base, configuración ni contenedores.

Se consultaron catálogos, identificadores y estados técnicos mínimos para el linaje y la revisión Alembic, además del parámetro público de duración de sesión. No se consultaron ventas individuales, credenciales ni documentos completos de las propuestas.

La existencia física y el estado técnico no sustituyen la prueba operativa ni la conciliación de una carga. Los casos nuevos 135/23 y 136/24 se explican y se prueban en los capítulos de operación de esta edición; este capítulo distingue su estructura de los ejemplos históricos 129/21 y 131/22. Los estados de todos ellos se cotejaron sin generar propuestas ni ejecutar cargas adicionales durante esta inspección.

\chapter{Responsabilidades, permisos y precisión de esta edición}
Catálogo base extraído literalmente del código; las columnas de Analista y Gerencia son una propuesta de capacitación, no una lectura de sus asignaciones actuales en PostgreSQL.

\section{Lectura de la matriz}
El rol de Analista responsable del catálogo incluye Administrar catálogo solo si esa responsabilidad se le asignó. Un perfil de analista que únicamente crea propuestas puede mantener el catálogo en consulta.

Actualizar metadatos puede reservarse a administración y retirarse del analista si la preparación de fuentes se realiza centralmente.

La propuesta de Gerente comercial incluye lectura de conexiones sin secretos, lectura de metadatos por dependencia del selector, panel y exportación. No incluye materialización, revisión ni administración del catálogo.

Administrador es el perfil base sembrado. Los nombres locales Analista BI, Gerente comercial u Operador no garantizan por sí mismos ninguna lista de permisos.

Product Owner y Scrum Master son responsabilidades del proyecto académico y no perfiles de autorización de la aplicación.

\section{Permisos de la plataforma}
\begingroup
\setbox0=\vbox{\hsize=\linewidth
{\normalfont\normalsize\sffamily\bfseries\normalsize\color{Teal}\hyphenpenalty=10000\exhyphenpenalty=10000 Consultar usuarios\par}

Permite consultar usuarios.

\textbf{Referencia técnica:} \texttt{\seqsplit{security.users.read}}\par
\par
\begingroup
\fontsize{10}{12.5}\selectfont
\setlength{\tabcolsep}{5pt}
\renewcommand{\arraystretch}{1.15}
\begin{tabular}{@{}>{\RaggedRight\arraybackslash}p{0.27\linewidth}>{\RaggedRight\arraybackslash}p{0.66\linewidth}@{}}
\toprule
\textbf{Perfil} & \textbf{Tratamiento en la formación} \\
\midrule
Administración & Incluido en el rol administrativo base protegido. \\
\bottomrule
\end{tabular}
\endgroup
}
\Needspace{\dimexpr\ht0+\dp0+48pt\relax}
\endgroup
\subsection{Consultar usuarios}
Permite consultar usuarios.

\textbf{Referencia técnica:} \texttt{\seqsplit{security.users.read}}\par
\begingroup
\fontsize{10}{12.5}\selectfont
\setlength{\tabcolsep}{5pt}
\renewcommand{\arraystretch}{1.15}
\begin{longtable}{@{}>{\RaggedRight\arraybackslash}p{0.27\linewidth}>{\RaggedRight\arraybackslash}p{0.66\linewidth}@{}}
\toprule
\textbf{Perfil} & \textbf{Tratamiento en la formación} \\
\midrule
\endfirsthead
\toprule
\textbf{Perfil} & \textbf{Tratamiento en la formación} \\
\midrule
\endhead
\midrule
\multicolumn{2}{r}{\footnotesize Continúa en la página siguiente}\\
\endfoot
\bottomrule
\endlastfoot
Administración & Incluido en el rol administrativo base protegido. \\
\addlinespace[3pt]
Analista de BI & No necesario para el recorrido propuesto. \\
\addlinespace[3pt]
Gerencia comercial & No necesario para el recorrido de gerencia propuesto. \\
\addlinespace[3pt]
\end{longtable}
\endgroup

\begingroup
\setbox0=\vbox{\hsize=\linewidth
{\normalfont\normalsize\sffamily\bfseries\normalsize\color{Teal}\hyphenpenalty=10000\exhyphenpenalty=10000 Administrar usuarios\par}

Permite crear y modificar usuarios.

\textbf{Referencia técnica:} \texttt{\seqsplit{security.users.write}}\par
\par
\begingroup
\fontsize{10}{12.5}\selectfont
\setlength{\tabcolsep}{5pt}
\renewcommand{\arraystretch}{1.15}
\begin{tabular}{@{}>{\RaggedRight\arraybackslash}p{0.27\linewidth}>{\RaggedRight\arraybackslash}p{0.66\linewidth}@{}}
\toprule
\textbf{Perfil} & \textbf{Tratamiento en la formación} \\
\midrule
Administración & Incluido en el rol administrativo base protegido. \\
\bottomrule
\end{tabular}
\endgroup
}
\Needspace{\dimexpr\ht0+\dp0+48pt\relax}
\endgroup
\subsection{Administrar usuarios}
Permite crear y modificar usuarios.

\textbf{Referencia técnica:} \texttt{\seqsplit{security.users.write}}\par
\begingroup
\fontsize{10}{12.5}\selectfont
\setlength{\tabcolsep}{5pt}
\renewcommand{\arraystretch}{1.15}
\begin{longtable}{@{}>{\RaggedRight\arraybackslash}p{0.27\linewidth}>{\RaggedRight\arraybackslash}p{0.66\linewidth}@{}}
\toprule
\textbf{Perfil} & \textbf{Tratamiento en la formación} \\
\midrule
\endfirsthead
\toprule
\textbf{Perfil} & \textbf{Tratamiento en la formación} \\
\midrule
\endhead
\midrule
\multicolumn{2}{r}{\footnotesize Continúa en la página siguiente}\\
\endfoot
\bottomrule
\endlastfoot
Administración & Incluido en el rol administrativo base protegido. \\
\addlinespace[3pt]
Analista de BI & No necesario para el recorrido propuesto. \\
\addlinespace[3pt]
Gerencia comercial & No necesario para el recorrido de gerencia propuesto. \\
\addlinespace[3pt]
\end{longtable}
\endgroup

\begingroup
\setbox0=\vbox{\hsize=\linewidth
{\normalfont\normalsize\sffamily\bfseries\normalsize\color{Teal}\hyphenpenalty=10000\exhyphenpenalty=10000 Consultar roles\par}

Permite consultar roles.

\textbf{Referencia técnica:} \texttt{\seqsplit{security.roles.read}}\par
\par
\begingroup
\fontsize{10}{12.5}\selectfont
\setlength{\tabcolsep}{5pt}
\renewcommand{\arraystretch}{1.15}
\begin{tabular}{@{}>{\RaggedRight\arraybackslash}p{0.27\linewidth}>{\RaggedRight\arraybackslash}p{0.66\linewidth}@{}}
\toprule
\textbf{Perfil} & \textbf{Tratamiento en la formación} \\
\midrule
Administración & Incluido en el rol administrativo base protegido. \\
\bottomrule
\end{tabular}
\endgroup
}
\Needspace{\dimexpr\ht0+\dp0+48pt\relax}
\endgroup
\subsection{Consultar roles}
Permite consultar roles.

\textbf{Referencia técnica:} \texttt{\seqsplit{security.roles.read}}\par
\begingroup
\fontsize{10}{12.5}\selectfont
\setlength{\tabcolsep}{5pt}
\renewcommand{\arraystretch}{1.15}
\begin{longtable}{@{}>{\RaggedRight\arraybackslash}p{0.27\linewidth}>{\RaggedRight\arraybackslash}p{0.66\linewidth}@{}}
\toprule
\textbf{Perfil} & \textbf{Tratamiento en la formación} \\
\midrule
\endfirsthead
\toprule
\textbf{Perfil} & \textbf{Tratamiento en la formación} \\
\midrule
\endhead
\midrule
\multicolumn{2}{r}{\footnotesize Continúa en la página siguiente}\\
\endfoot
\bottomrule
\endlastfoot
Administración & Incluido en el rol administrativo base protegido. \\
\addlinespace[3pt]
Analista de BI & No necesario para el recorrido propuesto. \\
\addlinespace[3pt]
Gerencia comercial & No necesario para el recorrido de gerencia propuesto. \\
\addlinespace[3pt]
\end{longtable}
\endgroup

\begingroup
\setbox0=\vbox{\hsize=\linewidth
{\normalfont\normalsize\sffamily\bfseries\normalsize\color{Teal}\hyphenpenalty=10000\exhyphenpenalty=10000 Administrar roles\par}

Permite crear roles y asignar permisos.

\textbf{Referencia técnica:} \texttt{\seqsplit{security.roles.write}}\par
\par
\begingroup
\fontsize{10}{12.5}\selectfont
\setlength{\tabcolsep}{5pt}
\renewcommand{\arraystretch}{1.15}
\begin{tabular}{@{}>{\RaggedRight\arraybackslash}p{0.27\linewidth}>{\RaggedRight\arraybackslash}p{0.66\linewidth}@{}}
\toprule
\textbf{Perfil} & \textbf{Tratamiento en la formación} \\
\midrule
Administración & Incluido en el rol administrativo base protegido. \\
\bottomrule
\end{tabular}
\endgroup
}
\Needspace{\dimexpr\ht0+\dp0+48pt\relax}
\endgroup
\subsection{Administrar roles}
Permite crear roles y asignar permisos.

\textbf{Referencia técnica:} \texttt{\seqsplit{security.roles.write}}\par
\begingroup
\fontsize{10}{12.5}\selectfont
\setlength{\tabcolsep}{5pt}
\renewcommand{\arraystretch}{1.15}
\begin{longtable}{@{}>{\RaggedRight\arraybackslash}p{0.27\linewidth}>{\RaggedRight\arraybackslash}p{0.66\linewidth}@{}}
\toprule
\textbf{Perfil} & \textbf{Tratamiento en la formación} \\
\midrule
\endfirsthead
\toprule
\textbf{Perfil} & \textbf{Tratamiento en la formación} \\
\midrule
\endhead
\midrule
\multicolumn{2}{r}{\footnotesize Continúa en la página siguiente}\\
\endfoot
\bottomrule
\endlastfoot
Administración & Incluido en el rol administrativo base protegido. \\
\addlinespace[3pt]
Analista de BI & No necesario para el recorrido propuesto. \\
\addlinespace[3pt]
Gerencia comercial & No necesario para el recorrido de gerencia propuesto. \\
\addlinespace[3pt]
\end{longtable}
\endgroup

\begingroup
\setbox0=\vbox{\hsize=\linewidth
{\normalfont\normalsize\sffamily\bfseries\normalsize\color{Teal}\hyphenpenalty=10000\exhyphenpenalty=10000 Consultar permisos\par}

Permite consultar permisos.

\textbf{Referencia técnica:} \texttt{\seqsplit{security.permissions.read}}\par
\par
\begingroup
\fontsize{10}{12.5}\selectfont
\setlength{\tabcolsep}{5pt}
\renewcommand{\arraystretch}{1.15}
\begin{tabular}{@{}>{\RaggedRight\arraybackslash}p{0.27\linewidth}>{\RaggedRight\arraybackslash}p{0.66\linewidth}@{}}
\toprule
\textbf{Perfil} & \textbf{Tratamiento en la formación} \\
\midrule
Administración & Incluido en el rol administrativo base protegido. \\
\bottomrule
\end{tabular}
\endgroup
}
\Needspace{\dimexpr\ht0+\dp0+48pt\relax}
\endgroup
\subsection{Consultar permisos}
Permite consultar permisos.

\textbf{Referencia técnica:} \texttt{\seqsplit{security.permissions.read}}\par
\begingroup
\fontsize{10}{12.5}\selectfont
\setlength{\tabcolsep}{5pt}
\renewcommand{\arraystretch}{1.15}
\begin{longtable}{@{}>{\RaggedRight\arraybackslash}p{0.27\linewidth}>{\RaggedRight\arraybackslash}p{0.66\linewidth}@{}}
\toprule
\textbf{Perfil} & \textbf{Tratamiento en la formación} \\
\midrule
\endfirsthead
\toprule
\textbf{Perfil} & \textbf{Tratamiento en la formación} \\
\midrule
\endhead
\midrule
\multicolumn{2}{r}{\footnotesize Continúa en la página siguiente}\\
\endfoot
\bottomrule
\endlastfoot
Administración & Incluido en el rol administrativo base protegido. \\
\addlinespace[3pt]
Analista de BI & No necesario para el recorrido propuesto. \\
\addlinespace[3pt]
Gerencia comercial & No necesario para el recorrido de gerencia propuesto. \\
\addlinespace[3pt]
\end{longtable}
\endgroup

\begingroup
\setbox0=\vbox{\hsize=\linewidth
{\normalfont\normalsize\sffamily\bfseries\normalsize\color{Teal}\hyphenpenalty=10000\exhyphenpenalty=10000 Administrar permisos\par}

Permite crear permisos.

\textbf{Referencia técnica:} \texttt{\seqsplit{security.permissions.write}}\par
\par
\begingroup
\fontsize{10}{12.5}\selectfont
\setlength{\tabcolsep}{5pt}
\renewcommand{\arraystretch}{1.15}
\begin{tabular}{@{}>{\RaggedRight\arraybackslash}p{0.27\linewidth}>{\RaggedRight\arraybackslash}p{0.66\linewidth}@{}}
\toprule
\textbf{Perfil} & \textbf{Tratamiento en la formación} \\
\midrule
Administración & Incluido en el rol administrativo base protegido. \\
\bottomrule
\end{tabular}
\endgroup
}
\Needspace{\dimexpr\ht0+\dp0+48pt\relax}
\endgroup
\subsection{Administrar permisos}
Permite crear permisos.

\textbf{Referencia técnica:} \texttt{\seqsplit{security.permissions.write}}\par
\begingroup
\fontsize{10}{12.5}\selectfont
\setlength{\tabcolsep}{5pt}
\renewcommand{\arraystretch}{1.15}
\begin{longtable}{@{}>{\RaggedRight\arraybackslash}p{0.27\linewidth}>{\RaggedRight\arraybackslash}p{0.66\linewidth}@{}}
\toprule
\textbf{Perfil} & \textbf{Tratamiento en la formación} \\
\midrule
\endfirsthead
\toprule
\textbf{Perfil} & \textbf{Tratamiento en la formación} \\
\midrule
\endhead
\midrule
\multicolumn{2}{r}{\footnotesize Continúa en la página siguiente}\\
\endfoot
\bottomrule
\endlastfoot
Administración & Incluido en el rol administrativo base protegido. \\
\addlinespace[3pt]
Analista de BI & No necesario para el recorrido propuesto. \\
\addlinespace[3pt]
Gerencia comercial & No necesario para el recorrido de gerencia propuesto. \\
\addlinespace[3pt]
\end{longtable}
\endgroup

\begingroup
\setbox0=\vbox{\hsize=\linewidth
{\normalfont\normalsize\sffamily\bfseries\normalsize\color{Teal}\hyphenpenalty=10000\exhyphenpenalty=10000 Consultar menús\par}

Permite consultar la navegación.

\textbf{Referencia técnica:} \texttt{\seqsplit{security.menus.read}}\par
\par
\begingroup
\fontsize{10}{12.5}\selectfont
\setlength{\tabcolsep}{5pt}
\renewcommand{\arraystretch}{1.15}
\begin{tabular}{@{}>{\RaggedRight\arraybackslash}p{0.27\linewidth}>{\RaggedRight\arraybackslash}p{0.66\linewidth}@{}}
\toprule
\textbf{Perfil} & \textbf{Tratamiento en la formación} \\
\midrule
Administración & Incluido en el rol administrativo base protegido. \\
\bottomrule
\end{tabular}
\endgroup
}
\Needspace{\dimexpr\ht0+\dp0+48pt\relax}
\endgroup
\subsection{Consultar menús}
Permite consultar la navegación.

\textbf{Referencia técnica:} \texttt{\seqsplit{security.menus.read}}\par
\begingroup
\fontsize{10}{12.5}\selectfont
\setlength{\tabcolsep}{5pt}
\renewcommand{\arraystretch}{1.15}
\begin{longtable}{@{}>{\RaggedRight\arraybackslash}p{0.27\linewidth}>{\RaggedRight\arraybackslash}p{0.66\linewidth}@{}}
\toprule
\textbf{Perfil} & \textbf{Tratamiento en la formación} \\
\midrule
\endfirsthead
\toprule
\textbf{Perfil} & \textbf{Tratamiento en la formación} \\
\midrule
\endhead
\midrule
\multicolumn{2}{r}{\footnotesize Continúa en la página siguiente}\\
\endfoot
\bottomrule
\endlastfoot
Administración & Incluido en el rol administrativo base protegido. \\
\addlinespace[3pt]
Analista de BI & No necesario para el recorrido propuesto. \\
\addlinespace[3pt]
Gerencia comercial & No necesario para el recorrido de gerencia propuesto. \\
\addlinespace[3pt]
\end{longtable}
\endgroup

\begingroup
\setbox0=\vbox{\hsize=\linewidth
{\normalfont\normalsize\sffamily\bfseries\normalsize\color{Teal}\hyphenpenalty=10000\exhyphenpenalty=10000 Administrar menús\par}

Permite crear y parametrizar menús.

\textbf{Referencia técnica:} \texttt{\seqsplit{security.menus.write}}\par
\par
\begingroup
\fontsize{10}{12.5}\selectfont
\setlength{\tabcolsep}{5pt}
\renewcommand{\arraystretch}{1.15}
\begin{tabular}{@{}>{\RaggedRight\arraybackslash}p{0.27\linewidth}>{\RaggedRight\arraybackslash}p{0.66\linewidth}@{}}
\toprule
\textbf{Perfil} & \textbf{Tratamiento en la formación} \\
\midrule
Administración & Incluido en el rol administrativo base protegido. \\
\bottomrule
\end{tabular}
\endgroup
}
\Needspace{\dimexpr\ht0+\dp0+48pt\relax}
\endgroup
\subsection{Administrar menús}
Permite crear y parametrizar menús.

\textbf{Referencia técnica:} \texttt{\seqsplit{security.menus.write}}\par
\begingroup
\fontsize{10}{12.5}\selectfont
\setlength{\tabcolsep}{5pt}
\renewcommand{\arraystretch}{1.15}
\begin{longtable}{@{}>{\RaggedRight\arraybackslash}p{0.27\linewidth}>{\RaggedRight\arraybackslash}p{0.66\linewidth}@{}}
\toprule
\textbf{Perfil} & \textbf{Tratamiento en la formación} \\
\midrule
\endfirsthead
\toprule
\textbf{Perfil} & \textbf{Tratamiento en la formación} \\
\midrule
\endhead
\midrule
\multicolumn{2}{r}{\footnotesize Continúa en la página siguiente}\\
\endfoot
\bottomrule
\endlastfoot
Administración & Incluido en el rol administrativo base protegido. \\
\addlinespace[3pt]
Analista de BI & No necesario para el recorrido propuesto. \\
\addlinespace[3pt]
Gerencia comercial & No necesario para el recorrido de gerencia propuesto. \\
\addlinespace[3pt]
\end{longtable}
\endgroup

\begingroup
\setbox0=\vbox{\hsize=\linewidth
{\normalfont\normalsize\sffamily\bfseries\normalsize\color{Teal}\hyphenpenalty=10000\exhyphenpenalty=10000 Consultar parámetros\par}

Permite consultar parámetros operativos.

\textbf{Referencia técnica:} \texttt{\seqsplit{parameters.read}}\par
\par
\begingroup
\fontsize{10}{12.5}\selectfont
\setlength{\tabcolsep}{5pt}
\renewcommand{\arraystretch}{1.15}
\begin{tabular}{@{}>{\RaggedRight\arraybackslash}p{0.27\linewidth}>{\RaggedRight\arraybackslash}p{0.66\linewidth}@{}}
\toprule
\textbf{Perfil} & \textbf{Tratamiento en la formación} \\
\midrule
Administración & Incluido en el rol administrativo base protegido. \\
\bottomrule
\end{tabular}
\endgroup
}
\Needspace{\dimexpr\ht0+\dp0+48pt\relax}
\endgroup
\subsection{Consultar parámetros}
Permite consultar parámetros operativos.

\textbf{Referencia técnica:} \texttt{\seqsplit{parameters.read}}\par
\begingroup
\fontsize{10}{12.5}\selectfont
\setlength{\tabcolsep}{5pt}
\renewcommand{\arraystretch}{1.15}
\begin{longtable}{@{}>{\RaggedRight\arraybackslash}p{0.27\linewidth}>{\RaggedRight\arraybackslash}p{0.66\linewidth}@{}}
\toprule
\textbf{Perfil} & \textbf{Tratamiento en la formación} \\
\midrule
\endfirsthead
\toprule
\textbf{Perfil} & \textbf{Tratamiento en la formación} \\
\midrule
\endhead
\midrule
\multicolumn{2}{r}{\footnotesize Continúa en la página siguiente}\\
\endfoot
\bottomrule
\endlastfoot
Administración & Incluido en el rol administrativo base protegido. \\
\addlinespace[3pt]
Analista de BI & No necesario para el recorrido propuesto. \\
\addlinespace[3pt]
Gerencia comercial & No necesario para el recorrido de gerencia propuesto. \\
\addlinespace[3pt]
\end{longtable}
\endgroup

\begingroup
\setbox0=\vbox{\hsize=\linewidth
{\normalfont\normalsize\sffamily\bfseries\normalsize\color{Teal}\hyphenpenalty=10000\exhyphenpenalty=10000 Administrar parámetros\par}

Permite crear y modificar parámetros.

\textbf{Referencia técnica:} \texttt{\seqsplit{parameters.write}}\par
\par
\begingroup
\fontsize{10}{12.5}\selectfont
\setlength{\tabcolsep}{5pt}
\renewcommand{\arraystretch}{1.15}
\begin{tabular}{@{}>{\RaggedRight\arraybackslash}p{0.27\linewidth}>{\RaggedRight\arraybackslash}p{0.66\linewidth}@{}}
\toprule
\textbf{Perfil} & \textbf{Tratamiento en la formación} \\
\midrule
Administración & Incluido en el rol administrativo base protegido. \\
\bottomrule
\end{tabular}
\endgroup
}
\Needspace{\dimexpr\ht0+\dp0+48pt\relax}
\endgroup
\subsection{Administrar parámetros}
Permite crear y modificar parámetros.

\textbf{Referencia técnica:} \texttt{\seqsplit{parameters.write}}\par
\begingroup
\fontsize{10}{12.5}\selectfont
\setlength{\tabcolsep}{5pt}
\renewcommand{\arraystretch}{1.15}
\begin{longtable}{@{}>{\RaggedRight\arraybackslash}p{0.27\linewidth}>{\RaggedRight\arraybackslash}p{0.66\linewidth}@{}}
\toprule
\textbf{Perfil} & \textbf{Tratamiento en la formación} \\
\midrule
\endfirsthead
\toprule
\textbf{Perfil} & \textbf{Tratamiento en la formación} \\
\midrule
\endhead
\midrule
\multicolumn{2}{r}{\footnotesize Continúa en la página siguiente}\\
\endfoot
\bottomrule
\endlastfoot
Administración & Incluido en el rol administrativo base protegido. \\
\addlinespace[3pt]
Analista de BI & No necesario para el recorrido propuesto. \\
\addlinespace[3pt]
Gerencia comercial & No necesario para el recorrido de gerencia propuesto. \\
\addlinespace[3pt]
\end{longtable}
\endgroup

\begingroup
\setbox0=\vbox{\hsize=\linewidth
{\normalfont\normalsize\sffamily\bfseries\normalsize\color{Teal}\hyphenpenalty=10000\exhyphenpenalty=10000 Consultar configuración LLM\par}

Permite consultar proveedores LLM.

\textbf{Referencia técnica:} \texttt{\seqsplit{parameters.llm.read}}\par
\par
\begingroup
\fontsize{10}{12.5}\selectfont
\setlength{\tabcolsep}{5pt}
\renewcommand{\arraystretch}{1.15}
\begin{tabular}{@{}>{\RaggedRight\arraybackslash}p{0.27\linewidth}>{\RaggedRight\arraybackslash}p{0.66\linewidth}@{}}
\toprule
\textbf{Perfil} & \textbf{Tratamiento en la formación} \\
\midrule
Administración & Incluido en el rol administrativo base protegido. \\
\bottomrule
\end{tabular}
\endgroup
}
\Needspace{\dimexpr\ht0+\dp0+48pt\relax}
\endgroup
\subsection{Consultar configuración LLM}
Permite consultar proveedores LLM.

\textbf{Referencia técnica:} \texttt{\seqsplit{parameters.llm.read}}\par
\begingroup
\fontsize{10}{12.5}\selectfont
\setlength{\tabcolsep}{5pt}
\renewcommand{\arraystretch}{1.15}
\begin{longtable}{@{}>{\RaggedRight\arraybackslash}p{0.27\linewidth}>{\RaggedRight\arraybackslash}p{0.66\linewidth}@{}}
\toprule
\textbf{Perfil} & \textbf{Tratamiento en la formación} \\
\midrule
\endfirsthead
\toprule
\textbf{Perfil} & \textbf{Tratamiento en la formación} \\
\midrule
\endhead
\midrule
\multicolumn{2}{r}{\footnotesize Continúa en la página siguiente}\\
\endfoot
\bottomrule
\endlastfoot
Administración & Incluido en el rol administrativo base protegido. \\
\addlinespace[3pt]
Analista de BI & No necesario para el recorrido propuesto. \\
\addlinespace[3pt]
Gerencia comercial & No necesario para el recorrido de gerencia propuesto. \\
\addlinespace[3pt]
\end{longtable}
\endgroup

\begingroup
\setbox0=\vbox{\hsize=\linewidth
{\normalfont\normalsize\sffamily\bfseries\normalsize\color{Teal}\hyphenpenalty=10000\exhyphenpenalty=10000 Administrar configuración LLM\par}

Permite parametrizar proveedores LLM.

\textbf{Referencia técnica:} \texttt{\seqsplit{parameters.llm.write}}\par
\par
\begingroup
\fontsize{10}{12.5}\selectfont
\setlength{\tabcolsep}{5pt}
\renewcommand{\arraystretch}{1.15}
\begin{tabular}{@{}>{\RaggedRight\arraybackslash}p{0.27\linewidth}>{\RaggedRight\arraybackslash}p{0.66\linewidth}@{}}
\toprule
\textbf{Perfil} & \textbf{Tratamiento en la formación} \\
\midrule
Administración & Incluido en el rol administrativo base protegido. \\
\bottomrule
\end{tabular}
\endgroup
}
\Needspace{\dimexpr\ht0+\dp0+48pt\relax}
\endgroup
\subsection{Administrar configuración LLM}
Permite parametrizar proveedores LLM.

\textbf{Referencia técnica:} \texttt{\seqsplit{parameters.llm.write}}\par
\begingroup
\fontsize{10}{12.5}\selectfont
\setlength{\tabcolsep}{5pt}
\renewcommand{\arraystretch}{1.15}
\begin{longtable}{@{}>{\RaggedRight\arraybackslash}p{0.27\linewidth}>{\RaggedRight\arraybackslash}p{0.66\linewidth}@{}}
\toprule
\textbf{Perfil} & \textbf{Tratamiento en la formación} \\
\midrule
\endfirsthead
\toprule
\textbf{Perfil} & \textbf{Tratamiento en la formación} \\
\midrule
\endhead
\midrule
\multicolumn{2}{r}{\footnotesize Continúa en la página siguiente}\\
\endfoot
\bottomrule
\endlastfoot
Administración & Incluido en el rol administrativo base protegido. \\
\addlinespace[3pt]
Analista de BI & No necesario para el recorrido propuesto. \\
\addlinespace[3pt]
Gerencia comercial & No necesario para el recorrido de gerencia propuesto. \\
\addlinespace[3pt]
\end{longtable}
\endgroup

\begingroup
\setbox0=\vbox{\hsize=\linewidth
{\normalfont\normalsize\sffamily\bfseries\normalsize\color{Teal}\hyphenpenalty=10000\exhyphenpenalty=10000 Consultar conexiones\par}

Permite consultar fuentes sin secretos.

\textbf{Referencia técnica:} \texttt{\seqsplit{connections.read}}\par
\par
\begingroup
\fontsize{10}{12.5}\selectfont
\setlength{\tabcolsep}{5pt}
\renewcommand{\arraystretch}{1.15}
\begin{tabular}{@{}>{\RaggedRight\arraybackslash}p{0.27\linewidth}>{\RaggedRight\arraybackslash}p{0.66\linewidth}@{}}
\toprule
\textbf{Perfil} & \textbf{Tratamiento en la formación} \\
\midrule
Administración & Incluido en el rol administrativo base protegido. \\
\bottomrule
\end{tabular}
\endgroup
}
\Needspace{\dimexpr\ht0+\dp0+48pt\relax}
\endgroup
\subsection{Consultar conexiones}
Permite consultar fuentes sin secretos.

\textbf{Referencia técnica:} \texttt{\seqsplit{connections.read}}\par
\begingroup
\fontsize{10}{12.5}\selectfont
\setlength{\tabcolsep}{5pt}
\renewcommand{\arraystretch}{1.15}
\begin{longtable}{@{}>{\RaggedRight\arraybackslash}p{0.27\linewidth}>{\RaggedRight\arraybackslash}p{0.66\linewidth}@{}}
\toprule
\textbf{Perfil} & \textbf{Tratamiento en la formación} \\
\midrule
\endfirsthead
\toprule
\textbf{Perfil} & \textbf{Tratamiento en la formación} \\
\midrule
\endhead
\midrule
\multicolumn{2}{r}{\footnotesize Continúa en la página siguiente}\\
\endfoot
\bottomrule
\endlastfoot
Administración & Incluido en el rol administrativo base protegido. \\
\addlinespace[3pt]
Analista de BI & No necesario para el recorrido propuesto. \\
\addlinespace[3pt]
Gerencia comercial & Incluir para el recorrido de consulta propuesto. \\
\addlinespace[3pt]
\end{longtable}
\endgroup

\begingroup
\setbox0=\vbox{\hsize=\linewidth
{\normalfont\normalsize\sffamily\bfseries\normalsize\color{Teal}\hyphenpenalty=10000\exhyphenpenalty=10000 Administrar conexiones\par}

Permite parametrizar fuentes de datos.

\textbf{Referencia técnica:} \texttt{\seqsplit{connections.write}}\par
\par
\begingroup
\fontsize{10}{12.5}\selectfont
\setlength{\tabcolsep}{5pt}
\renewcommand{\arraystretch}{1.15}
\begin{tabular}{@{}>{\RaggedRight\arraybackslash}p{0.27\linewidth}>{\RaggedRight\arraybackslash}p{0.66\linewidth}@{}}
\toprule
\textbf{Perfil} & \textbf{Tratamiento en la formación} \\
\midrule
Administración & Incluido en el rol administrativo base protegido. \\
\bottomrule
\end{tabular}
\endgroup
}
\Needspace{\dimexpr\ht0+\dp0+48pt\relax}
\endgroup
\subsection{Administrar conexiones}
Permite parametrizar fuentes de datos.

\textbf{Referencia técnica:} \texttt{\seqsplit{connections.write}}\par
\begingroup
\fontsize{10}{12.5}\selectfont
\setlength{\tabcolsep}{5pt}
\renewcommand{\arraystretch}{1.15}
\begin{longtable}{@{}>{\RaggedRight\arraybackslash}p{0.27\linewidth}>{\RaggedRight\arraybackslash}p{0.66\linewidth}@{}}
\toprule
\textbf{Perfil} & \textbf{Tratamiento en la formación} \\
\midrule
\endfirsthead
\toprule
\textbf{Perfil} & \textbf{Tratamiento en la formación} \\
\midrule
\endhead
\midrule
\multicolumn{2}{r}{\footnotesize Continúa en la página siguiente}\\
\endfoot
\bottomrule
\endlastfoot
Administración & Incluido en el rol administrativo base protegido. \\
\addlinespace[3pt]
Analista de BI & No necesario para el recorrido propuesto. \\
\addlinespace[3pt]
Gerencia comercial & No necesario para el recorrido de gerencia propuesto. \\
\addlinespace[3pt]
\end{longtable}
\endgroup

\begingroup
\setbox0=\vbox{\hsize=\linewidth
{\normalfont\normalsize\sffamily\bfseries\normalsize\color{Teal}\hyphenpenalty=10000\exhyphenpenalty=10000 Probar conexiones\par}

Permite validar fuentes de solo lectura.

\textbf{Referencia técnica:} \texttt{\seqsplit{connections.test}}\par
\par
\begingroup
\fontsize{10}{12.5}\selectfont
\setlength{\tabcolsep}{5pt}
\renewcommand{\arraystretch}{1.15}
\begin{tabular}{@{}>{\RaggedRight\arraybackslash}p{0.27\linewidth}>{\RaggedRight\arraybackslash}p{0.66\linewidth}@{}}
\toprule
\textbf{Perfil} & \textbf{Tratamiento en la formación} \\
\midrule
Administración & Incluido en el rol administrativo base protegido. \\
\bottomrule
\end{tabular}
\endgroup
}
\Needspace{\dimexpr\ht0+\dp0+48pt\relax}
\endgroup
\subsection{Probar conexiones}
Permite validar fuentes de solo lectura.

\textbf{Referencia técnica:} \texttt{\seqsplit{connections.test}}\par
\begingroup
\fontsize{10}{12.5}\selectfont
\setlength{\tabcolsep}{5pt}
\renewcommand{\arraystretch}{1.15}
\begin{longtable}{@{}>{\RaggedRight\arraybackslash}p{0.27\linewidth}>{\RaggedRight\arraybackslash}p{0.66\linewidth}@{}}
\toprule
\textbf{Perfil} & \textbf{Tratamiento en la formación} \\
\midrule
\endfirsthead
\toprule
\textbf{Perfil} & \textbf{Tratamiento en la formación} \\
\midrule
\endhead
\midrule
\multicolumn{2}{r}{\footnotesize Continúa en la página siguiente}\\
\endfoot
\bottomrule
\endlastfoot
Administración & Incluido en el rol administrativo base protegido. \\
\addlinespace[3pt]
Analista de BI & No necesario para el recorrido propuesto. \\
\addlinespace[3pt]
Gerencia comercial & No necesario para el recorrido de gerencia propuesto. \\
\addlinespace[3pt]
\end{longtable}
\endgroup

\begingroup
\setbox0=\vbox{\hsize=\linewidth
{\normalfont\normalsize\sffamily\bfseries\normalsize\color{Teal}\hyphenpenalty=10000\exhyphenpenalty=10000 Consultar metadatos\par}

Permite consultar instantáneas del esquema.

\textbf{Referencia técnica:} \texttt{\seqsplit{metadata.read}}\par
\par
\begingroup
\fontsize{10}{12.5}\selectfont
\setlength{\tabcolsep}{5pt}
\renewcommand{\arraystretch}{1.15}
\begin{tabular}{@{}>{\RaggedRight\arraybackslash}p{0.27\linewidth}>{\RaggedRight\arraybackslash}p{0.66\linewidth}@{}}
\toprule
\textbf{Perfil} & \textbf{Tratamiento en la formación} \\
\midrule
Administración & Incluido en el rol administrativo base protegido. \\
\bottomrule
\end{tabular}
\endgroup
}
\Needspace{\dimexpr\ht0+\dp0+48pt\relax}
\endgroup
\subsection{Consultar metadatos}
Permite consultar instantáneas del esquema.

\textbf{Referencia técnica:} \texttt{\seqsplit{metadata.read}}\par
\begingroup
\fontsize{10}{12.5}\selectfont
\setlength{\tabcolsep}{5pt}
\renewcommand{\arraystretch}{1.15}
\begin{longtable}{@{}>{\RaggedRight\arraybackslash}p{0.27\linewidth}>{\RaggedRight\arraybackslash}p{0.66\linewidth}@{}}
\toprule
\textbf{Perfil} & \textbf{Tratamiento en la formación} \\
\midrule
\endfirsthead
\toprule
\textbf{Perfil} & \textbf{Tratamiento en la formación} \\
\midrule
\endhead
\midrule
\multicolumn{2}{r}{\footnotesize Continúa en la página siguiente}\\
\endfoot
\bottomrule
\endlastfoot
Administración & Incluido en el rol administrativo base protegido. \\
\addlinespace[3pt]
Analista de BI & Incluir según la responsabilidad acordada. \\
\addlinespace[3pt]
Gerencia comercial & Incluir para el recorrido de consulta propuesto. \\
\addlinespace[3pt]
\end{longtable}
\endgroup

\begingroup
\setbox0=\vbox{\hsize=\linewidth
{\normalfont\normalsize\sffamily\bfseries\normalsize\color{Teal}\hyphenpenalty=10000\exhyphenpenalty=10000 Actualizar metadatos\par}

Permite leer la estructura de una fuente habilitada y seleccionada.

\textbf{Referencia técnica:} \texttt{\seqsplit{metadata.refresh}}\par
\par
\begingroup
\fontsize{10}{12.5}\selectfont
\setlength{\tabcolsep}{5pt}
\renewcommand{\arraystretch}{1.15}
\begin{tabular}{@{}>{\RaggedRight\arraybackslash}p{0.27\linewidth}>{\RaggedRight\arraybackslash}p{0.66\linewidth}@{}}
\toprule
\textbf{Perfil} & \textbf{Tratamiento en la formación} \\
\midrule
Administración & Incluido en el rol administrativo base protegido. \\
\bottomrule
\end{tabular}
\endgroup
}
\Needspace{\dimexpr\ht0+\dp0+48pt\relax}
\endgroup
\subsection{Actualizar metadatos}
Permite leer la estructura de una fuente habilitada y seleccionada.

\textbf{Referencia técnica:} \texttt{\seqsplit{metadata.refresh}}\par
\begingroup
\fontsize{10}{12.5}\selectfont
\setlength{\tabcolsep}{5pt}
\renewcommand{\arraystretch}{1.15}
\begin{longtable}{@{}>{\RaggedRight\arraybackslash}p{0.27\linewidth}>{\RaggedRight\arraybackslash}p{0.66\linewidth}@{}}
\toprule
\textbf{Perfil} & \textbf{Tratamiento en la formación} \\
\midrule
\endfirsthead
\toprule
\textbf{Perfil} & \textbf{Tratamiento en la formación} \\
\midrule
\endhead
\midrule
\multicolumn{2}{r}{\footnotesize Continúa en la página siguiente}\\
\endfoot
\bottomrule
\endlastfoot
Administración & Incluido en el rol administrativo base protegido. \\
\addlinespace[3pt]
Analista de BI & Incluir según la responsabilidad acordada. \\
\addlinespace[3pt]
Gerencia comercial & No necesario para el recorrido de gerencia propuesto. \\
\addlinespace[3pt]
\end{longtable}
\endgroup

\begingroup
\setbox0=\vbox{\hsize=\linewidth
{\normalfont\normalsize\sffamily\bfseries\normalsize\color{Teal}\hyphenpenalty=10000\exhyphenpenalty=10000 Revisar identidad de dimensiones\par}

Permite comprobar relaciones, cobertura y muestras descriptivas controladas.

\textbf{Referencia técnica:} \texttt{\seqsplit{metadata.semantic\_resolution.read}}\par
\par
\begingroup
\fontsize{10}{12.5}\selectfont
\setlength{\tabcolsep}{5pt}
\renewcommand{\arraystretch}{1.15}
\begin{tabular}{@{}>{\RaggedRight\arraybackslash}p{0.27\linewidth}>{\RaggedRight\arraybackslash}p{0.66\linewidth}@{}}
\toprule
\textbf{Perfil} & \textbf{Tratamiento en la formación} \\
\midrule
Administración & Incluido en el rol administrativo base protegido. \\
\bottomrule
\end{tabular}
\endgroup
}
\Needspace{\dimexpr\ht0+\dp0+48pt\relax}
\endgroup
\subsection{Revisar identidad de dimensiones}
Permite comprobar relaciones, cobertura y muestras descriptivas controladas.

\textbf{Referencia técnica:} \texttt{\seqsplit{metadata.semantic\_resolution.read}}\par
\begingroup
\fontsize{10}{12.5}\selectfont
\setlength{\tabcolsep}{5pt}
\renewcommand{\arraystretch}{1.15}
\begin{longtable}{@{}>{\RaggedRight\arraybackslash}p{0.27\linewidth}>{\RaggedRight\arraybackslash}p{0.66\linewidth}@{}}
\toprule
\textbf{Perfil} & \textbf{Tratamiento en la formación} \\
\midrule
\endfirsthead
\toprule
\textbf{Perfil} & \textbf{Tratamiento en la formación} \\
\midrule
\endhead
\midrule
\multicolumn{2}{r}{\footnotesize Continúa en la página siguiente}\\
\endfoot
\bottomrule
\endlastfoot
Administración & Incluido en el rol administrativo base protegido. \\
\addlinespace[3pt]
Analista de BI & Incluir según la responsabilidad acordada. \\
\addlinespace[3pt]
Gerencia comercial & No necesario para el recorrido de gerencia propuesto. \\
\addlinespace[3pt]
\end{longtable}
\endgroup

\begingroup
\setbox0=\vbox{\hsize=\linewidth
{\normalfont\normalsize\sffamily\bfseries\normalsize\color{Teal}\hyphenpenalty=10000\exhyphenpenalty=10000 Resolver identidad de dimensiones\par}

Permite crear una versión corregida desde rutas semánticas verificadas.

\textbf{Referencia técnica:} \texttt{\seqsplit{metadata.semantic\_resolution.manage}}\par
\par
\begingroup
\fontsize{10}{12.5}\selectfont
\setlength{\tabcolsep}{5pt}
\renewcommand{\arraystretch}{1.15}
\begin{tabular}{@{}>{\RaggedRight\arraybackslash}p{0.27\linewidth}>{\RaggedRight\arraybackslash}p{0.66\linewidth}@{}}
\toprule
\textbf{Perfil} & \textbf{Tratamiento en la formación} \\
\midrule
Administración & Incluido en el rol administrativo base protegido. \\
\bottomrule
\end{tabular}
\endgroup
}
\Needspace{\dimexpr\ht0+\dp0+48pt\relax}
\endgroup
\subsection{Resolver identidad de dimensiones}
Permite crear una versión corregida desde rutas semánticas verificadas.

\textbf{Referencia técnica:} \texttt{\seqsplit{metadata.semantic\_resolution.manage}}\par
\begingroup
\fontsize{10}{12.5}\selectfont
\setlength{\tabcolsep}{5pt}
\renewcommand{\arraystretch}{1.15}
\begin{longtable}{@{}>{\RaggedRight\arraybackslash}p{0.27\linewidth}>{\RaggedRight\arraybackslash}p{0.66\linewidth}@{}}
\toprule
\textbf{Perfil} & \textbf{Tratamiento en la formación} \\
\midrule
\endfirsthead
\toprule
\textbf{Perfil} & \textbf{Tratamiento en la formación} \\
\midrule
\endhead
\midrule
\multicolumn{2}{r}{\footnotesize Continúa en la página siguiente}\\
\endfoot
\bottomrule
\endlastfoot
Administración & Incluido en el rol administrativo base protegido. \\
\addlinespace[3pt]
Analista de BI & No necesario para el recorrido propuesto. \\
\addlinespace[3pt]
Gerencia comercial & No necesario para el recorrido de gerencia propuesto. \\
\addlinespace[3pt]
\end{longtable}
\endgroup

Declarado en el catálogo; no se identificó una operación autorizada con este código en el recorrido web auditado. No atribuir una capacidad nueva por su mera existencia.

\begingroup
\setbox0=\vbox{\hsize=\linewidth
{\normalfont\normalsize\sffamily\bfseries\normalsize\color{Teal}\hyphenpenalty=10000\exhyphenpenalty=10000 Consultar propuestas BI\par}

Permite consultar solicitudes, propuestas y decisiones del asistente.

\textbf{Referencia técnica:} \texttt{\seqsplit{copilot.proposals.read}}\par
\par
\begingroup
\fontsize{10}{12.5}\selectfont
\setlength{\tabcolsep}{5pt}
\renewcommand{\arraystretch}{1.15}
\begin{tabular}{@{}>{\RaggedRight\arraybackslash}p{0.27\linewidth}>{\RaggedRight\arraybackslash}p{0.66\linewidth}@{}}
\toprule
\textbf{Perfil} & \textbf{Tratamiento en la formación} \\
\midrule
Administración & Incluido en el rol administrativo base protegido. \\
\bottomrule
\end{tabular}
\endgroup
}
\Needspace{\dimexpr\ht0+\dp0+48pt\relax}
\endgroup
\subsection{Consultar propuestas BI}
Permite consultar solicitudes, propuestas y decisiones del asistente.

\textbf{Referencia técnica:} \texttt{\seqsplit{copilot.proposals.read}}\par
\begingroup
\fontsize{10}{12.5}\selectfont
\setlength{\tabcolsep}{5pt}
\renewcommand{\arraystretch}{1.15}
\begin{longtable}{@{}>{\RaggedRight\arraybackslash}p{0.27\linewidth}>{\RaggedRight\arraybackslash}p{0.66\linewidth}@{}}
\toprule
\textbf{Perfil} & \textbf{Tratamiento en la formación} \\
\midrule
\endfirsthead
\toprule
\textbf{Perfil} & \textbf{Tratamiento en la formación} \\
\midrule
\endhead
\midrule
\multicolumn{2}{r}{\footnotesize Continúa en la página siguiente}\\
\endfoot
\bottomrule
\endlastfoot
Administración & Incluido en el rol administrativo base protegido. \\
\addlinespace[3pt]
Analista de BI & Incluir según la responsabilidad acordada. \\
\addlinespace[3pt]
Gerencia comercial & No necesario para el recorrido de gerencia propuesto. \\
\addlinespace[3pt]
\end{longtable}
\endgroup

\begingroup
\setbox0=\vbox{\hsize=\linewidth
{\normalfont\normalsize\sffamily\bfseries\normalsize\color{Teal}\hyphenpenalty=10000\exhyphenpenalty=10000 Generar propuestas BI\par}

Permite solicitar interpretación semántica y propuestas de ventas.

\textbf{Referencia técnica:} \texttt{\seqsplit{copilot.proposals.generate}}\par
\par
\begingroup
\fontsize{10}{12.5}\selectfont
\setlength{\tabcolsep}{5pt}
\renewcommand{\arraystretch}{1.15}
\begin{tabular}{@{}>{\RaggedRight\arraybackslash}p{0.27\linewidth}>{\RaggedRight\arraybackslash}p{0.66\linewidth}@{}}
\toprule
\textbf{Perfil} & \textbf{Tratamiento en la formación} \\
\midrule
Administración & Incluido en el rol administrativo base protegido. \\
\bottomrule
\end{tabular}
\endgroup
}
\Needspace{\dimexpr\ht0+\dp0+48pt\relax}
\endgroup
\subsection{Generar propuestas BI}
Permite solicitar interpretación semántica y propuestas de ventas.

\textbf{Referencia técnica:} \texttt{\seqsplit{copilot.proposals.generate}}\par
\begingroup
\fontsize{10}{12.5}\selectfont
\setlength{\tabcolsep}{5pt}
\renewcommand{\arraystretch}{1.15}
\begin{longtable}{@{}>{\RaggedRight\arraybackslash}p{0.27\linewidth}>{\RaggedRight\arraybackslash}p{0.66\linewidth}@{}}
\toprule
\textbf{Perfil} & \textbf{Tratamiento en la formación} \\
\midrule
\endfirsthead
\toprule
\textbf{Perfil} & \textbf{Tratamiento en la formación} \\
\midrule
\endhead
\midrule
\multicolumn{2}{r}{\footnotesize Continúa en la página siguiente}\\
\endfoot
\bottomrule
\endlastfoot
Administración & Incluido en el rol administrativo base protegido. \\
\addlinespace[3pt]
Analista de BI & Incluir según la responsabilidad acordada. \\
\addlinespace[3pt]
Gerencia comercial & No necesario para el recorrido de gerencia propuesto. \\
\addlinespace[3pt]
\end{longtable}
\endgroup

\begingroup
\setbox0=\vbox{\hsize=\linewidth
{\normalfont\normalsize\sffamily\bfseries\normalsize\color{Teal}\hyphenpenalty=10000\exhyphenpenalty=10000 Revisar propuestas BI\par}

Permite aprobar o rechazar propuestas validadas.

\textbf{Referencia técnica:} \texttt{\seqsplit{copilot.proposals.review}}\par
\par
\begingroup
\fontsize{10}{12.5}\selectfont
\setlength{\tabcolsep}{5pt}
\renewcommand{\arraystretch}{1.15}
\begin{tabular}{@{}>{\RaggedRight\arraybackslash}p{0.27\linewidth}>{\RaggedRight\arraybackslash}p{0.66\linewidth}@{}}
\toprule
\textbf{Perfil} & \textbf{Tratamiento en la formación} \\
\midrule
Administración & Incluido en el rol administrativo base protegido. \\
\bottomrule
\end{tabular}
\endgroup
}
\Needspace{\dimexpr\ht0+\dp0+48pt\relax}
\endgroup
\subsection{Revisar propuestas BI}
Permite aprobar o rechazar propuestas validadas.

\textbf{Referencia técnica:} \texttt{\seqsplit{copilot.proposals.review}}\par
\begingroup
\fontsize{10}{12.5}\selectfont
\setlength{\tabcolsep}{5pt}
\renewcommand{\arraystretch}{1.15}
\begin{longtable}{@{}>{\RaggedRight\arraybackslash}p{0.27\linewidth}>{\RaggedRight\arraybackslash}p{0.66\linewidth}@{}}
\toprule
\textbf{Perfil} & \textbf{Tratamiento en la formación} \\
\midrule
\endfirsthead
\toprule
\textbf{Perfil} & \textbf{Tratamiento en la formación} \\
\midrule
\endhead
\midrule
\multicolumn{2}{r}{\footnotesize Continúa en la página siguiente}\\
\endfoot
\bottomrule
\endlastfoot
Administración & Incluido en el rol administrativo base protegido. \\
\addlinespace[3pt]
Analista de BI & Incluir según la responsabilidad acordada. \\
\addlinespace[3pt]
Gerencia comercial & No necesario para el recorrido de gerencia propuesto. \\
\addlinespace[3pt]
\end{longtable}
\endgroup

\begingroup
\setbox0=\vbox{\hsize=\linewidth
{\normalfont\normalsize\sffamily\bfseries\normalsize\color{Teal}\hyphenpenalty=10000\exhyphenpenalty=10000 Consultar catálogo analítico\par}

Permite consultar preguntas de negocio y periodicidades por dominio.

\textbf{Referencia técnica:} \texttt{\seqsplit{copilot.catalog.read}}\par
\par
\begingroup
\fontsize{10}{12.5}\selectfont
\setlength{\tabcolsep}{5pt}
\renewcommand{\arraystretch}{1.15}
\begin{tabular}{@{}>{\RaggedRight\arraybackslash}p{0.27\linewidth}>{\RaggedRight\arraybackslash}p{0.66\linewidth}@{}}
\toprule
\textbf{Perfil} & \textbf{Tratamiento en la formación} \\
\midrule
Administración & Incluido en el rol administrativo base protegido. \\
\bottomrule
\end{tabular}
\endgroup
}
\Needspace{\dimexpr\ht0+\dp0+48pt\relax}
\endgroup
\subsection{Consultar catálogo analítico}
Permite consultar preguntas de negocio y periodicidades por dominio.

\textbf{Referencia técnica:} \texttt{\seqsplit{copilot.catalog.read}}\par
\begingroup
\fontsize{10}{12.5}\selectfont
\setlength{\tabcolsep}{5pt}
\renewcommand{\arraystretch}{1.15}
\begin{longtable}{@{}>{\RaggedRight\arraybackslash}p{0.27\linewidth}>{\RaggedRight\arraybackslash}p{0.66\linewidth}@{}}
\toprule
\textbf{Perfil} & \textbf{Tratamiento en la formación} \\
\midrule
\endfirsthead
\toprule
\textbf{Perfil} & \textbf{Tratamiento en la formación} \\
\midrule
\endhead
\midrule
\multicolumn{2}{r}{\footnotesize Continúa en la página siguiente}\\
\endfoot
\bottomrule
\endlastfoot
Administración & Incluido en el rol administrativo base protegido. \\
\addlinespace[3pt]
Analista de BI & Incluir según la responsabilidad acordada. \\
\addlinespace[3pt]
Gerencia comercial & No necesario para el recorrido de gerencia propuesto. \\
\addlinespace[3pt]
\end{longtable}
\endgroup

\begingroup
\setbox0=\vbox{\hsize=\linewidth
{\normalfont\normalsize\sffamily\bfseries\normalsize\color{Teal}\hyphenpenalty=10000\exhyphenpenalty=10000 Administrar catálogo analítico\par}

Permite crear, modificar, habilitar y retirar guías analíticas por dominio.

\textbf{Referencia técnica:} \texttt{\seqsplit{copilot.catalog.write}}\par
\par
\begingroup
\fontsize{10}{12.5}\selectfont
\setlength{\tabcolsep}{5pt}
\renewcommand{\arraystretch}{1.15}
\begin{tabular}{@{}>{\RaggedRight\arraybackslash}p{0.27\linewidth}>{\RaggedRight\arraybackslash}p{0.66\linewidth}@{}}
\toprule
\textbf{Perfil} & \textbf{Tratamiento en la formación} \\
\midrule
Administración & Incluido en el rol administrativo base protegido. \\
\bottomrule
\end{tabular}
\endgroup
}
\Needspace{\dimexpr\ht0+\dp0+48pt\relax}
\endgroup
\subsection{Administrar catálogo analítico}
Permite crear, modificar, habilitar y retirar guías analíticas por dominio.

\textbf{Referencia técnica:} \texttt{\seqsplit{copilot.catalog.write}}\par
\begingroup
\fontsize{10}{12.5}\selectfont
\setlength{\tabcolsep}{5pt}
\renewcommand{\arraystretch}{1.15}
\begin{longtable}{@{}>{\RaggedRight\arraybackslash}p{0.27\linewidth}>{\RaggedRight\arraybackslash}p{0.66\linewidth}@{}}
\toprule
\textbf{Perfil} & \textbf{Tratamiento en la formación} \\
\midrule
\endfirsthead
\toprule
\textbf{Perfil} & \textbf{Tratamiento en la formación} \\
\midrule
\endhead
\midrule
\multicolumn{2}{r}{\footnotesize Continúa en la página siguiente}\\
\endfoot
\bottomrule
\endlastfoot
Administración & Incluido en el rol administrativo base protegido. \\
\addlinespace[3pt]
Analista de BI & Incluir según la responsabilidad acordada. \\
\addlinespace[3pt]
Gerencia comercial & No necesario para el recorrido de gerencia propuesto. \\
\addlinespace[3pt]
\end{longtable}
\endgroup

\begingroup
\setbox0=\vbox{\hsize=\linewidth
{\normalfont\normalsize\sffamily\bfseries\normalsize\color{Teal}\hyphenpenalty=10000\exhyphenpenalty=10000 Consultar materializaciones\par}

Permite revisar propuestas elegibles, planes y evidencia de ejecuciones ETL.

\textbf{Referencia técnica:} \texttt{\seqsplit{etl.executions.read}}\par
\par
\begingroup
\fontsize{10}{12.5}\selectfont
\setlength{\tabcolsep}{5pt}
\renewcommand{\arraystretch}{1.15}
\begin{tabular}{@{}>{\RaggedRight\arraybackslash}p{0.27\linewidth}>{\RaggedRight\arraybackslash}p{0.66\linewidth}@{}}
\toprule
\textbf{Perfil} & \textbf{Tratamiento en la formación} \\
\midrule
Administración & Incluido en el rol administrativo base protegido. \\
\bottomrule
\end{tabular}
\endgroup
}
\Needspace{\dimexpr\ht0+\dp0+48pt\relax}
\endgroup
\subsection{Consultar materializaciones}
Permite revisar propuestas elegibles, planes y evidencia de ejecuciones ETL.

\textbf{Referencia técnica:} \texttt{\seqsplit{etl.executions.read}}\par
\begingroup
\fontsize{10}{12.5}\selectfont
\setlength{\tabcolsep}{5pt}
\renewcommand{\arraystretch}{1.15}
\begin{longtable}{@{}>{\RaggedRight\arraybackslash}p{0.27\linewidth}>{\RaggedRight\arraybackslash}p{0.66\linewidth}@{}}
\toprule
\textbf{Perfil} & \textbf{Tratamiento en la formación} \\
\midrule
\endfirsthead
\toprule
\textbf{Perfil} & \textbf{Tratamiento en la formación} \\
\midrule
\endhead
\midrule
\multicolumn{2}{r}{\footnotesize Continúa en la página siguiente}\\
\endfoot
\bottomrule
\endlastfoot
Administración & Incluido en el rol administrativo base protegido. \\
\addlinespace[3pt]
Analista de BI & Incluir según la responsabilidad acordada. \\
\addlinespace[3pt]
Gerencia comercial & No necesario para el recorrido de gerencia propuesto. \\
\addlinespace[3pt]
\end{longtable}
\endgroup

\begingroup
\setbox0=\vbox{\hsize=\linewidth
{\normalfont\normalsize\sffamily\bfseries\normalsize\color{Teal}\hyphenpenalty=10000\exhyphenpenalty=10000 Preparar materializaciones\par}

Permite seleccionar una propuesta y preparar una ejecución ETL controlada.

\textbf{Referencia técnica:} \texttt{\seqsplit{etl.executions.write}}\par
\par
\begingroup
\fontsize{10}{12.5}\selectfont
\setlength{\tabcolsep}{5pt}
\renewcommand{\arraystretch}{1.15}
\begin{tabular}{@{}>{\RaggedRight\arraybackslash}p{0.27\linewidth}>{\RaggedRight\arraybackslash}p{0.66\linewidth}@{}}
\toprule
\textbf{Perfil} & \textbf{Tratamiento en la formación} \\
\midrule
Administración & Incluido en el rol administrativo base protegido. \\
\bottomrule
\end{tabular}
\endgroup
}
\Needspace{\dimexpr\ht0+\dp0+48pt\relax}
\endgroup
\subsection{Preparar materializaciones}
Permite seleccionar una propuesta y preparar una ejecución ETL controlada.

\textbf{Referencia técnica:} \texttt{\seqsplit{etl.executions.write}}\par
\begingroup
\fontsize{10}{12.5}\selectfont
\setlength{\tabcolsep}{5pt}
\renewcommand{\arraystretch}{1.15}
\begin{longtable}{@{}>{\RaggedRight\arraybackslash}p{0.27\linewidth}>{\RaggedRight\arraybackslash}p{0.66\linewidth}@{}}
\toprule
\textbf{Perfil} & \textbf{Tratamiento en la formación} \\
\midrule
\endfirsthead
\toprule
\textbf{Perfil} & \textbf{Tratamiento en la formación} \\
\midrule
\endhead
\midrule
\multicolumn{2}{r}{\footnotesize Continúa en la página siguiente}\\
\endfoot
\bottomrule
\endlastfoot
Administración & Incluido en el rol administrativo base protegido. \\
\addlinespace[3pt]
Analista de BI & Incluir según la responsabilidad acordada. \\
\addlinespace[3pt]
Gerencia comercial & No necesario para el recorrido de gerencia propuesto. \\
\addlinespace[3pt]
\end{longtable}
\endgroup

\begingroup
\setbox0=\vbox{\hsize=\linewidth
{\normalfont\normalsize\sffamily\bfseries\normalsize\color{Teal}\hyphenpenalty=10000\exhyphenpenalty=10000 Consultar análisis visual\par}

Permite explorar KPI, visualizaciones e insights sobre un datamart conciliado.

\textbf{Referencia técnica:} \texttt{\seqsplit{analytics.dashboard.read}}\par
\par
\begingroup
\fontsize{10}{12.5}\selectfont
\setlength{\tabcolsep}{5pt}
\renewcommand{\arraystretch}{1.15}
\begin{tabular}{@{}>{\RaggedRight\arraybackslash}p{0.27\linewidth}>{\RaggedRight\arraybackslash}p{0.66\linewidth}@{}}
\toprule
\textbf{Perfil} & \textbf{Tratamiento en la formación} \\
\midrule
Administración & Incluido en el rol administrativo base protegido. \\
\bottomrule
\end{tabular}
\endgroup
}
\Needspace{\dimexpr\ht0+\dp0+48pt\relax}
\endgroup
\subsection{Consultar análisis visual}
Permite explorar KPI, visualizaciones e insights sobre un datamart conciliado.

\textbf{Referencia técnica:} \texttt{\seqsplit{analytics.dashboard.read}}\par
\begingroup
\fontsize{10}{12.5}\selectfont
\setlength{\tabcolsep}{5pt}
\renewcommand{\arraystretch}{1.15}
\begin{longtable}{@{}>{\RaggedRight\arraybackslash}p{0.27\linewidth}>{\RaggedRight\arraybackslash}p{0.66\linewidth}@{}}
\toprule
\textbf{Perfil} & \textbf{Tratamiento en la formación} \\
\midrule
\endfirsthead
\toprule
\textbf{Perfil} & \textbf{Tratamiento en la formación} \\
\midrule
\endhead
\midrule
\multicolumn{2}{r}{\footnotesize Continúa en la página siguiente}\\
\endfoot
\bottomrule
\endlastfoot
Administración & Incluido en el rol administrativo base protegido. \\
\addlinespace[3pt]
Analista de BI & Incluir según la responsabilidad acordada. \\
\addlinespace[3pt]
Gerencia comercial & Incluir para el recorrido de consulta propuesto. \\
\addlinespace[3pt]
\end{longtable}
\endgroup

\begingroup
\setbox0=\vbox{\hsize=\linewidth
{\normalfont\normalsize\sffamily\bfseries\normalsize\color{Teal}\hyphenpenalty=10000\exhyphenpenalty=10000 Exportar reportes analíticos\par}

Permite descargar en PDF o Excel la selección analítica visible y autorizada.

\textbf{Referencia técnica:} \texttt{\seqsplit{reports.analytics.export}}\par
\par
\begingroup
\fontsize{10}{12.5}\selectfont
\setlength{\tabcolsep}{5pt}
\renewcommand{\arraystretch}{1.15}
\begin{tabular}{@{}>{\RaggedRight\arraybackslash}p{0.27\linewidth}>{\RaggedRight\arraybackslash}p{0.66\linewidth}@{}}
\toprule
\textbf{Perfil} & \textbf{Tratamiento en la formación} \\
\midrule
Administración & Incluido en el rol administrativo base protegido. \\
\bottomrule
\end{tabular}
\endgroup
}
\Needspace{\dimexpr\ht0+\dp0+48pt\relax}
\endgroup
\subsection{Exportar reportes analíticos}
Permite descargar en PDF o Excel la selección analítica visible y autorizada.

\textbf{Referencia técnica:} \texttt{\seqsplit{reports.analytics.export}}\par
\begingroup
\fontsize{10}{12.5}\selectfont
\setlength{\tabcolsep}{5pt}
\renewcommand{\arraystretch}{1.15}
\begin{longtable}{@{}>{\RaggedRight\arraybackslash}p{0.27\linewidth}>{\RaggedRight\arraybackslash}p{0.66\linewidth}@{}}
\toprule
\textbf{Perfil} & \textbf{Tratamiento en la formación} \\
\midrule
\endfirsthead
\toprule
\textbf{Perfil} & \textbf{Tratamiento en la formación} \\
\midrule
\endhead
\midrule
\multicolumn{2}{r}{\footnotesize Continúa en la página siguiente}\\
\endfoot
\bottomrule
\endlastfoot
Administración & Incluido en el rol administrativo base protegido. \\
\addlinespace[3pt]
Analista de BI & Incluir según la responsabilidad acordada. \\
\addlinespace[3pt]
Gerencia comercial & Incluir para el recorrido de consulta propuesto. \\
\addlinespace[3pt]
\end{longtable}
\endgroup

\begingroup
\setbox0=\vbox{\hsize=\linewidth
{\normalfont\normalsize\sffamily\bfseries\normalsize\color{Teal}\hyphenpenalty=10000\exhyphenpenalty=10000 Consultar auditoría\par}

Permite consultar eventos de auditoría.

\textbf{Referencia técnica:} \texttt{\seqsplit{audit.read}}\par
\par
\begingroup
\fontsize{10}{12.5}\selectfont
\setlength{\tabcolsep}{5pt}
\renewcommand{\arraystretch}{1.15}
\begin{tabular}{@{}>{\RaggedRight\arraybackslash}p{0.27\linewidth}>{\RaggedRight\arraybackslash}p{0.66\linewidth}@{}}
\toprule
\textbf{Perfil} & \textbf{Tratamiento en la formación} \\
\midrule
Administración & Incluido en el rol administrativo base protegido. \\
\bottomrule
\end{tabular}
\endgroup
}
\Needspace{\dimexpr\ht0+\dp0+48pt\relax}
\endgroup
\subsection{Consultar auditoría}
Permite consultar eventos de auditoría.

\textbf{Referencia técnica:} \texttt{\seqsplit{audit.read}}\par
\begingroup
\fontsize{10}{12.5}\selectfont
\setlength{\tabcolsep}{5pt}
\renewcommand{\arraystretch}{1.15}
\begin{longtable}{@{}>{\RaggedRight\arraybackslash}p{0.27\linewidth}>{\RaggedRight\arraybackslash}p{0.66\linewidth}@{}}
\toprule
\textbf{Perfil} & \textbf{Tratamiento en la formación} \\
\midrule
\endfirsthead
\toprule
\textbf{Perfil} & \textbf{Tratamiento en la formación} \\
\midrule
\endhead
\midrule
\multicolumn{2}{r}{\footnotesize Continúa en la página siguiente}\\
\endfoot
\bottomrule
\endlastfoot
Administración & Incluido en el rol administrativo base protegido. \\
\addlinespace[3pt]
Analista de BI & No necesario para el recorrido propuesto. \\
\addlinespace[3pt]
Gerencia comercial & No necesario para el recorrido de gerencia propuesto. \\
\addlinespace[3pt]
\end{longtable}
\endgroup

\begingroup
\setbox0=\vbox{\hsize=\linewidth
{\normalfont\normalsize\RaggedRight\sffamily\bfseries\large\color{Azul}\hyphenpenalty=10000\exhyphenpenalty=10000 Dependencias que deben comprobarse\par}

\par
\begingroup
\fontsize{10}{12.5}\selectfont
\setlength{\tabcolsep}{5pt}
\renewcommand{\arraystretch}{1.15}
\begin{tabular}{@{}>{\RaggedRight\arraybackslash}p{0.27\linewidth}>{\RaggedRight\arraybackslash}p{0.66\linewidth}@{}}
\toprule
\textbf{Operación} & \textbf{Permisos y condición} \\
\midrule
Seleccionar fuente y operar pantallas con contexto & Consultar metadatos (metadata.read). También permite abrir el Explorador en consulta. \\
\bottomrule
\end{tabular}
\endgroup
}
\Needspace{\dimexpr\ht0+\dp0+48pt\relax}
\endgroup
\section{Dependencias que deben comprobarse}
\begingroup
\fontsize{10}{12.5}\selectfont
\setlength{\tabcolsep}{5pt}
\renewcommand{\arraystretch}{1.15}
\begin{longtable}{@{}>{\RaggedRight\arraybackslash}p{0.27\linewidth}>{\RaggedRight\arraybackslash}p{0.66\linewidth}@{}}
\toprule
\textbf{Operación} & \textbf{Permisos y condición} \\
\midrule
\endfirsthead
\toprule
\textbf{Operación} & \textbf{Permisos y condición} \\
\midrule
\endhead
\midrule
\multicolumn{2}{r}{\footnotesize Continúa en la página siguiente}\\
\endfoot
\bottomrule
\endlastfoot
Seleccionar fuente y operar pantallas con contexto & Consultar metadatos (metadata.read). También permite abrir el Explorador en consulta. \\
\addlinespace[3pt]
Editar usuarios no protegidos mediante la UI & Administrar usuarios y Consultar roles para cargar las opciones. \\
\addlinespace[3pt]
Crear o editar roles y menús mediante la UI & Administrar el recurso y Consultar permisos para cargar las casillas. \\
\addlinespace[3pt]
Mostrar cobertura técnica del catálogo & La configuración usa Consultar catálogo analítico; la cobertura consulta el catálogo técnico mediante Consultar propuestas BI. La falta de esa consulta puede dejar la cobertura sin mostrar. \\
\addlinespace[3pt]
Exportar PDF o Excel & Consultar análisis visual y Exportar reportes analíticos conjuntamente. \\
\addlinespace[3pt]
Revisar o generar una propuesta & Consulta de propuestas más el permiso de generar o revisar correspondiente; no se confunden. \\
\addlinespace[3pt]
\end{longtable}
\endgroup

\section{Precisiones respecto de documentación anterior}
\subsection{Catálogo de preguntas}
Cuatro orientaciones base definidas por el perfil; no hay generación LLM del editor de catálogo. Dos orientaciones adicionales de AW fueron decisiones humanas históricas.

Explicar su procedencia y su independencia por fuente.

\subsection{Restauración del catálogo}
El endpoint reset aplica default\_\allowbreak{}needs\_\allowbreak{}catalog\_\allowbreak{}configuration a la fuente actual.

Advertir que sustituye personalizaciones; no confundirlo con recuperar la última edición.

\subsection{Paginación}
App.tsx usa const pageSize = 10 en los listados auditados aunque UI\_\allowbreak{}PAGE\_\allowbreak{}SIZE se guarda.

No prometer efecto inmediato de ese parámetro; verificar antes de una demostración.

\subsection{Usuarios y auditoría}
La UI no expone buscador de usuarios ni filtros de auditoría por fecha/actor en este corte.

Documentar únicamente la paginación disponible.

\subsection{Permisos y menús}
No se crean contratos técnicos o pantallas libres desde el navegador.

Separar edición de roles/etiquetas de ampliaciones del producto.

\subsection{Pruebas por rol}
La consulta de una configuración no acredita un recorrido autenticado bajo ese perfil.

Proponer la verificación sin afirmar que se ejecutó durante esta lectura de código.

\section{Origen de los contenidos y de los ejemplos}
Esta edición integra la lectura de código y especificaciones, las capturas históricas del 18 de noviembre y las comprobaciones operativas del 25 de noviembre. Los casos 135/23 y 136/24 son recorridos nuevos comprobados con OpenAI Sol. Luna sólo tiene comprobación mínima de conexión y JSON. Los ejercicios para que el grupo practique siguen distinguiéndose de esas pruebas realizadas.

No presentar ejercicios propuestos como pruebas realizadas.

No garantizar generación idéntica o éxito de todos los escenarios de un LLM.

No incluir credenciales o tokens.

Mantener distinción pedido/factura, rol de entidad, moneda e identidad.

Poner los resultados históricos junto a su fecha y expediente.

El material base del 18 de noviembre se redactó sin nuevas propuestas ni ETL. La campaña adicional del 25 de noviembre sí incluyó configuración OpenAI, correcciones, propuestas aprobadas y ETL nuevos, documentados en sus unidades. Copiar y compilar este manual no vuelve a ejecutar esas operaciones. Las capturas actuales son vistas reales estrechas, conservadas sin estirar; no se presentan como nuevas capturas de escritorio completo.

\section{Reproducir una práctica en otro equipo}
Un portable limpio no debe contener la API key privada ni todos los expedientes locales de quien preparó el proyecto. La administración registra su propia credencial, comprueba el proveedor y configura las fuentes. Los números 135, 136, 23 y 24 son referencias de esta instalación, no valores que deban forzarse en las bases. La guía de traslado vigente acompaña la entrega y debe usarse para puertos, credenciales locales y encendido.

La política de ahorro consiste en elegir un modelo económico compatible, reducir clics repetidos y reutilizar evidencia conservada. No se recortan metadatos verificados, requisitos ni validaciones para abaratar una propuesta: esa pérdida de contexto puede dejar fuera un resultado obligatorio.

\section{Fuentes para consultar modelos y facturación}
Las tarifas se consultaron el 25 de noviembre de 2026 en la documentación oficial de OpenAI. Pueden cambiar; verifique modelo y disponibilidad antes de una demostración.

\url{https://developers.openai.com/api/docs/models/gpt-6-luna}\par\smallskip
\url{https://developers.openai.com/api/docs/models/gpt-5.4-mini}\par\smallskip
\url{https://developers.openai.com/api/docs/models/gpt-6.1-terra}\par\smallskip
\url{https://developers.openai.com/api/docs/models/gpt-6.2-sol}\par\smallskip
\url{https://developers.openai.com/api/docs/models/gpt-6.2-astra}\par\smallskip

\end{document}
